\documentclass{article}
\usepackage[paperwidth=16cm,paperheight=14cm,margin=0cm]{geometry}
\usepackage[utf8]{inputenc}\usepackage[T1]{fontenc}
\usepackage{mathpazo}\usepackage{amsmath,amssymb}
\usepackage[most]{tcolorbox}
\pagestyle{empty}\setlength{\parindent}{0pt}

\definecolor{csui}{HTML}{1F6FA8}
\definecolor{cdeg}{HTML}{6A4C93}
\definecolor{cder}{HTML}{2E7D6B}
\definecolor{ctri}{HTML}{C77F0A}
\definecolor{cpsc}{HTML}{C0503A}
\definecolor{ccvp}{HTML}{8A5A2B}
\begin{document}
\begin{tcolorbox}[enhanced,width=16cm,height=14cm,colback=white,colframe=csui,boxrule=2pt,arc=0mm,
left=6mm,right=6mm,top=0mm,bottom=4mm,sharp corners,
overlay={\fill[csui] (frame.north west) rectangle ([yshift=-1.25cm]frame.north east);
\node[anchor=west,text=white,font=\sffamily\bfseries\LARGE] at ([xshift=5mm,yshift=-0.62cm]frame.north west) {Question 1};
\node[anchor=east,text=white,font=\sffamily\itshape\large] at ([xshift=-5mm,yshift=-0.62cm]frame.north east) {Suites};}]
\vspace*{1.75cm}
\begin{minipage}[t][6.9cm][c]{\linewidth}\LARGE\centering
\begin{minipage}{0.95\linewidth}\raggedright\linespread{1.25}\selectfont $(u_n)$ est la suite arithmétique de premier terme $u_0=3$ et de raison $2$. Que vaut $u_{10}$~?\end{minipage}
\end{minipage}

\vfill
\begin{tcolorbox}[enhanced,nobeforeafter,width=7.1cm,height=2.1cm,valign=center,colback=csui!6,colframe=csui,arc=2.5mm,boxrule=0.9pt,left=3mm]{\Large\bfseries\color{csui}a)}\quad\begin{minipage}[c]{5.3cm}\raggedright\Large $23$\end{minipage}\end{tcolorbox}\hfill\begin{tcolorbox}[enhanced,nobeforeafter,width=7.1cm,height=2.1cm,valign=center,colback=csui!6,colframe=csui,arc=2.5mm,boxrule=0.9pt,left=3mm]{\Large\bfseries\color{csui}b)}\quad\begin{minipage}[c]{5.3cm}\raggedright\Large $25$\end{minipage}\end{tcolorbox}\par\vspace{2.5mm}
\begin{tcolorbox}[enhanced,nobeforeafter,width=7.1cm,height=2.1cm,valign=center,colback=csui!6,colframe=csui,arc=2.5mm,boxrule=0.9pt,left=3mm]{\Large\bfseries\color{csui}c)}\quad\begin{minipage}[c]{5.3cm}\raggedright\Large $20$\end{minipage}\end{tcolorbox}\hfill\begin{tcolorbox}[enhanced,nobeforeafter,width=7.1cm,height=2.1cm,valign=center,colback=csui!6,colframe=csui,arc=2.5mm,boxrule=0.9pt,left=3mm]{\Large\bfseries\color{csui}d)}\quad\begin{minipage}[c]{5.3cm}\raggedright\Large $21$\end{minipage}\end{tcolorbox}
\end{tcolorbox}
\newpage
\begin{tcolorbox}[enhanced,width=16cm,height=14cm,colback=white,colframe=csui,boxrule=2pt,arc=0mm,
left=6mm,right=6mm,top=0mm,bottom=4mm,sharp corners,
overlay={\fill[csui] (frame.north west) rectangle ([yshift=-1.25cm]frame.north east);
\node[anchor=west,text=white,font=\sffamily\bfseries\LARGE] at ([xshift=5mm,yshift=-0.62cm]frame.north west) {Question 2};
\node[anchor=east,text=white,font=\sffamily\itshape\large] at ([xshift=-5mm,yshift=-0.62cm]frame.north east) {Suites};}]
\vspace*{1.75cm}
\begin{minipage}[t][6.9cm][c]{\linewidth}\LARGE\centering
\begin{minipage}{0.95\linewidth}\raggedright\linespread{1.25}\selectfont $(u_n)$ est la suite arithmétique de premier terme $u_0=5$ et de raison $-3$. Que vaut $u_{8}$~?\end{minipage}
\end{minipage}

\vfill
\begin{tcolorbox}[enhanced,nobeforeafter,width=7.1cm,height=2.1cm,valign=center,colback=csui!6,colframe=csui,arc=2.5mm,boxrule=0.9pt,left=3mm]{\Large\bfseries\color{csui}a)}\quad\begin{minipage}[c]{5.3cm}\raggedright\Large $-22$\end{minipage}\end{tcolorbox}\hfill\begin{tcolorbox}[enhanced,nobeforeafter,width=7.1cm,height=2.1cm,valign=center,colback=csui!6,colframe=csui,arc=2.5mm,boxrule=0.9pt,left=3mm]{\Large\bfseries\color{csui}b)}\quad\begin{minipage}[c]{5.3cm}\raggedright\Large $-19$\end{minipage}\end{tcolorbox}\par\vspace{2.5mm}
\begin{tcolorbox}[enhanced,nobeforeafter,width=7.1cm,height=2.1cm,valign=center,colback=csui!6,colframe=csui,arc=2.5mm,boxrule=0.9pt,left=3mm]{\Large\bfseries\color{csui}c)}\quad\begin{minipage}[c]{5.3cm}\raggedright\Large $-16$\end{minipage}\end{tcolorbox}\hfill\begin{tcolorbox}[enhanced,nobeforeafter,width=7.1cm,height=2.1cm,valign=center,colback=csui!6,colframe=csui,arc=2.5mm,boxrule=0.9pt,left=3mm]{\Large\bfseries\color{csui}d)}\quad\begin{minipage}[c]{5.3cm}\raggedright\Large $-24$\end{minipage}\end{tcolorbox}
\end{tcolorbox}
\newpage
\begin{tcolorbox}[enhanced,width=16cm,height=14cm,colback=white,colframe=csui,boxrule=2pt,arc=0mm,
left=6mm,right=6mm,top=0mm,bottom=4mm,sharp corners,
overlay={\fill[csui] (frame.north west) rectangle ([yshift=-1.25cm]frame.north east);
\node[anchor=west,text=white,font=\sffamily\bfseries\LARGE] at ([xshift=5mm,yshift=-0.62cm]frame.north west) {Question 3};
\node[anchor=east,text=white,font=\sffamily\itshape\large] at ([xshift=-5mm,yshift=-0.62cm]frame.north east) {Suites};}]
\vspace*{1.75cm}
\begin{minipage}[t][6.9cm][c]{\linewidth}\LARGE\centering
\begin{minipage}{0.95\linewidth}\raggedright\linespread{1.25}\selectfont $(u_n)$ est la suite géométrique de premier terme $u_0=3$ et de raison $2$. Que vaut $u_{4}$~?\end{minipage}
\end{minipage}

\vfill
\begin{tcolorbox}[enhanced,nobeforeafter,width=7.1cm,height=2.1cm,valign=center,colback=csui!6,colframe=csui,arc=2.5mm,boxrule=0.9pt,left=3mm]{\Large\bfseries\color{csui}a)}\quad\begin{minipage}[c]{5.3cm}\raggedright\Large $11$\end{minipage}\end{tcolorbox}\hfill\begin{tcolorbox}[enhanced,nobeforeafter,width=7.1cm,height=2.1cm,valign=center,colback=csui!6,colframe=csui,arc=2.5mm,boxrule=0.9pt,left=3mm]{\Large\bfseries\color{csui}b)}\quad\begin{minipage}[c]{5.3cm}\raggedright\Large $24$\end{minipage}\end{tcolorbox}\par\vspace{2.5mm}
\begin{tcolorbox}[enhanced,nobeforeafter,width=7.1cm,height=2.1cm,valign=center,colback=csui!6,colframe=csui,arc=2.5mm,boxrule=0.9pt,left=3mm]{\Large\bfseries\color{csui}c)}\quad\begin{minipage}[c]{5.3cm}\raggedright\Large $96$\end{minipage}\end{tcolorbox}\hfill\begin{tcolorbox}[enhanced,nobeforeafter,width=7.1cm,height=2.1cm,valign=center,colback=csui!6,colframe=csui,arc=2.5mm,boxrule=0.9pt,left=3mm]{\Large\bfseries\color{csui}d)}\quad\begin{minipage}[c]{5.3cm}\raggedright\Large $48$\end{minipage}\end{tcolorbox}
\end{tcolorbox}
\newpage
\begin{tcolorbox}[enhanced,width=16cm,height=14cm,colback=white,colframe=csui,boxrule=2pt,arc=0mm,
left=6mm,right=6mm,top=0mm,bottom=4mm,sharp corners,
overlay={\fill[csui] (frame.north west) rectangle ([yshift=-1.25cm]frame.north east);
\node[anchor=west,text=white,font=\sffamily\bfseries\LARGE] at ([xshift=5mm,yshift=-0.62cm]frame.north west) {Question 4};
\node[anchor=east,text=white,font=\sffamily\itshape\large] at ([xshift=-5mm,yshift=-0.62cm]frame.north east) {Suites};}]
\vspace*{1.75cm}
\begin{minipage}[t][6.9cm][c]{\linewidth}\LARGE\centering
\begin{minipage}{0.95\linewidth}\raggedright\linespread{1.25}\selectfont $(u_n)$ est la suite géométrique de premier terme $u_0=5$ et de raison $3$. Que vaut $u_{3}$~?\end{minipage}
\end{minipage}

\vfill
\begin{tcolorbox}[enhanced,nobeforeafter,width=7.1cm,height=2.1cm,valign=center,colback=csui!6,colframe=csui,arc=2.5mm,boxrule=0.9pt,left=3mm]{\Large\bfseries\color{csui}a)}\quad\begin{minipage}[c]{5.3cm}\raggedright\Large $14$\end{minipage}\end{tcolorbox}\hfill\begin{tcolorbox}[enhanced,nobeforeafter,width=7.1cm,height=2.1cm,valign=center,colback=csui!6,colframe=csui,arc=2.5mm,boxrule=0.9pt,left=3mm]{\Large\bfseries\color{csui}b)}\quad\begin{minipage}[c]{5.3cm}\raggedright\Large $405$\end{minipage}\end{tcolorbox}\par\vspace{2.5mm}
\begin{tcolorbox}[enhanced,nobeforeafter,width=7.1cm,height=2.1cm,valign=center,colback=csui!6,colframe=csui,arc=2.5mm,boxrule=0.9pt,left=3mm]{\Large\bfseries\color{csui}c)}\quad\begin{minipage}[c]{5.3cm}\raggedright\Large $45$\end{minipage}\end{tcolorbox}\hfill\begin{tcolorbox}[enhanced,nobeforeafter,width=7.1cm,height=2.1cm,valign=center,colback=csui!6,colframe=csui,arc=2.5mm,boxrule=0.9pt,left=3mm]{\Large\bfseries\color{csui}d)}\quad\begin{minipage}[c]{5.3cm}\raggedright\Large $135$\end{minipage}\end{tcolorbox}
\end{tcolorbox}
\newpage
\begin{tcolorbox}[enhanced,width=16cm,height=14cm,colback=white,colframe=csui,boxrule=2pt,arc=0mm,
left=6mm,right=6mm,top=0mm,bottom=4mm,sharp corners,
overlay={\fill[csui] (frame.north west) rectangle ([yshift=-1.25cm]frame.north east);
\node[anchor=west,text=white,font=\sffamily\bfseries\LARGE] at ([xshift=5mm,yshift=-0.62cm]frame.north west) {Question 5};
\node[anchor=east,text=white,font=\sffamily\itshape\large] at ([xshift=-5mm,yshift=-0.62cm]frame.north east) {Suites};}]
\vspace*{1.75cm}
\begin{minipage}[t][6.9cm][c]{\linewidth}\LARGE\centering
\begin{minipage}{0.95\linewidth}\raggedright\linespread{1.25}\selectfont On définit $u_0=3$ et, pour tout entier $n$, $u_{n+1}=2u_n+1$. Que vaut $u_2$~?\end{minipage}
\end{minipage}

\vfill
\begin{tcolorbox}[enhanced,nobeforeafter,width=7.1cm,height=2.1cm,valign=center,colback=csui!6,colframe=csui,arc=2.5mm,boxrule=0.9pt,left=3mm]{\Large\bfseries\color{csui}a)}\quad\begin{minipage}[c]{5.3cm}\raggedright\Large $8$\end{minipage}\end{tcolorbox}\hfill\begin{tcolorbox}[enhanced,nobeforeafter,width=7.1cm,height=2.1cm,valign=center,colback=csui!6,colframe=csui,arc=2.5mm,boxrule=0.9pt,left=3mm]{\Large\bfseries\color{csui}b)}\quad\begin{minipage}[c]{5.3cm}\raggedright\Large $7$\end{minipage}\end{tcolorbox}\par\vspace{2.5mm}
\begin{tcolorbox}[enhanced,nobeforeafter,width=7.1cm,height=2.1cm,valign=center,colback=csui!6,colframe=csui,arc=2.5mm,boxrule=0.9pt,left=3mm]{\Large\bfseries\color{csui}c)}\quad\begin{minipage}[c]{5.3cm}\raggedright\Large $31$\end{minipage}\end{tcolorbox}\hfill\begin{tcolorbox}[enhanced,nobeforeafter,width=7.1cm,height=2.1cm,valign=center,colback=csui!6,colframe=csui,arc=2.5mm,boxrule=0.9pt,left=3mm]{\Large\bfseries\color{csui}d)}\quad\begin{minipage}[c]{5.3cm}\raggedright\Large $15$\end{minipage}\end{tcolorbox}
\end{tcolorbox}
\newpage
\begin{tcolorbox}[enhanced,width=16cm,height=14cm,colback=white,colframe=csui,boxrule=2pt,arc=0mm,
left=6mm,right=6mm,top=0mm,bottom=4mm,sharp corners,
overlay={\fill[csui] (frame.north west) rectangle ([yshift=-1.25cm]frame.north east);
\node[anchor=west,text=white,font=\sffamily\bfseries\LARGE] at ([xshift=5mm,yshift=-0.62cm]frame.north west) {Question 6};
\node[anchor=east,text=white,font=\sffamily\itshape\large] at ([xshift=-5mm,yshift=-0.62cm]frame.north east) {Suites};}]
\vspace*{1.75cm}
\begin{minipage}[t][6.9cm][c]{\linewidth}\LARGE\centering
\begin{minipage}{0.95\linewidth}\raggedright\linespread{1.25}\selectfont On définit $u_0=2$ et, pour tout entier $n$, $u_{n+1}=3u_n-2$. Que vaut $u_2$~?\end{minipage}
\end{minipage}

\vfill
\begin{tcolorbox}[enhanced,nobeforeafter,width=7.1cm,height=2.1cm,valign=center,colback=csui!6,colframe=csui,arc=2.5mm,boxrule=0.9pt,left=3mm]{\Large\bfseries\color{csui}a)}\quad\begin{minipage}[c]{5.3cm}\raggedright\Large $28$\end{minipage}\end{tcolorbox}\hfill\begin{tcolorbox}[enhanced,nobeforeafter,width=7.1cm,height=2.1cm,valign=center,colback=csui!6,colframe=csui,arc=2.5mm,boxrule=0.9pt,left=3mm]{\Large\bfseries\color{csui}b)}\quad\begin{minipage}[c]{5.3cm}\raggedright\Large $4$\end{minipage}\end{tcolorbox}\par\vspace{2.5mm}
\begin{tcolorbox}[enhanced,nobeforeafter,width=7.1cm,height=2.1cm,valign=center,colback=csui!6,colframe=csui,arc=2.5mm,boxrule=0.9pt,left=3mm]{\Large\bfseries\color{csui}c)}\quad\begin{minipage}[c]{5.3cm}\raggedright\Large $10$\end{minipage}\end{tcolorbox}\hfill\begin{tcolorbox}[enhanced,nobeforeafter,width=7.1cm,height=2.1cm,valign=center,colback=csui!6,colframe=csui,arc=2.5mm,boxrule=0.9pt,left=3mm]{\Large\bfseries\color{csui}d)}\quad\begin{minipage}[c]{5.3cm}\raggedright\Large $2$\end{minipage}\end{tcolorbox}
\end{tcolorbox}
\newpage
\begin{tcolorbox}[enhanced,width=16cm,height=14cm,colback=white,colframe=csui,boxrule=2pt,arc=0mm,
left=6mm,right=6mm,top=0mm,bottom=4mm,sharp corners,
overlay={\fill[csui] (frame.north west) rectangle ([yshift=-1.25cm]frame.north east);
\node[anchor=west,text=white,font=\sffamily\bfseries\LARGE] at ([xshift=5mm,yshift=-0.62cm]frame.north west) {Question 7};
\node[anchor=east,text=white,font=\sffamily\itshape\large] at ([xshift=-5mm,yshift=-0.62cm]frame.north east) {Suites};}]
\vspace*{1.75cm}
\begin{minipage}[t][6.9cm][c]{\linewidth}\LARGE\centering
\begin{minipage}{0.95\linewidth}\raggedright\linespread{1.25}\selectfont Que vaut la somme $1+2+3+\dots+10$~?\end{minipage}
\end{minipage}

\vfill
\begin{tcolorbox}[enhanced,nobeforeafter,width=7.1cm,height=2.1cm,valign=center,colback=csui!6,colframe=csui,arc=2.5mm,boxrule=0.9pt,left=3mm]{\Large\bfseries\color{csui}a)}\quad\begin{minipage}[c]{5.3cm}\raggedright\Large $45$\end{minipage}\end{tcolorbox}\hfill\begin{tcolorbox}[enhanced,nobeforeafter,width=7.1cm,height=2.1cm,valign=center,colback=csui!6,colframe=csui,arc=2.5mm,boxrule=0.9pt,left=3mm]{\Large\bfseries\color{csui}b)}\quad\begin{minipage}[c]{5.3cm}\raggedright\Large $110$\end{minipage}\end{tcolorbox}\par\vspace{2.5mm}
\begin{tcolorbox}[enhanced,nobeforeafter,width=7.1cm,height=2.1cm,valign=center,colback=csui!6,colframe=csui,arc=2.5mm,boxrule=0.9pt,left=3mm]{\Large\bfseries\color{csui}c)}\quad\begin{minipage}[c]{5.3cm}\raggedright\Large $100$\end{minipage}\end{tcolorbox}\hfill\begin{tcolorbox}[enhanced,nobeforeafter,width=7.1cm,height=2.1cm,valign=center,colback=csui!6,colframe=csui,arc=2.5mm,boxrule=0.9pt,left=3mm]{\Large\bfseries\color{csui}d)}\quad\begin{minipage}[c]{5.3cm}\raggedright\Large $55$\end{minipage}\end{tcolorbox}
\end{tcolorbox}
\newpage
\begin{tcolorbox}[enhanced,width=16cm,height=14cm,colback=white,colframe=csui,boxrule=2pt,arc=0mm,
left=6mm,right=6mm,top=0mm,bottom=4mm,sharp corners,
overlay={\fill[csui] (frame.north west) rectangle ([yshift=-1.25cm]frame.north east);
\node[anchor=west,text=white,font=\sffamily\bfseries\LARGE] at ([xshift=5mm,yshift=-0.62cm]frame.north west) {Question 8};
\node[anchor=east,text=white,font=\sffamily\itshape\large] at ([xshift=-5mm,yshift=-0.62cm]frame.north east) {Suites};}]
\vspace*{1.75cm}
\begin{minipage}[t][6.9cm][c]{\linewidth}\LARGE\centering
\begin{minipage}{0.95\linewidth}\raggedright\linespread{1.25}\selectfont Que vaut la somme $1+2+3+\dots+20$~?\end{minipage}
\end{minipage}

\vfill
\begin{tcolorbox}[enhanced,nobeforeafter,width=7.1cm,height=2.1cm,valign=center,colback=csui!6,colframe=csui,arc=2.5mm,boxrule=0.9pt,left=3mm]{\Large\bfseries\color{csui}a)}\quad\begin{minipage}[c]{5.3cm}\raggedright\Large $190$\end{minipage}\end{tcolorbox}\hfill\begin{tcolorbox}[enhanced,nobeforeafter,width=7.1cm,height=2.1cm,valign=center,colback=csui!6,colframe=csui,arc=2.5mm,boxrule=0.9pt,left=3mm]{\Large\bfseries\color{csui}b)}\quad\begin{minipage}[c]{5.3cm}\raggedright\Large $400$\end{minipage}\end{tcolorbox}\par\vspace{2.5mm}
\begin{tcolorbox}[enhanced,nobeforeafter,width=7.1cm,height=2.1cm,valign=center,colback=csui!6,colframe=csui,arc=2.5mm,boxrule=0.9pt,left=3mm]{\Large\bfseries\color{csui}c)}\quad\begin{minipage}[c]{5.3cm}\raggedright\Large $420$\end{minipage}\end{tcolorbox}\hfill\begin{tcolorbox}[enhanced,nobeforeafter,width=7.1cm,height=2.1cm,valign=center,colback=csui!6,colframe=csui,arc=2.5mm,boxrule=0.9pt,left=3mm]{\Large\bfseries\color{csui}d)}\quad\begin{minipage}[c]{5.3cm}\raggedright\Large $210$\end{minipage}\end{tcolorbox}
\end{tcolorbox}
\newpage
\begin{tcolorbox}[enhanced,width=16cm,height=14cm,colback=white,colframe=csui,boxrule=2pt,arc=0mm,
left=6mm,right=6mm,top=0mm,bottom=4mm,sharp corners,
overlay={\fill[csui] (frame.north west) rectangle ([yshift=-1.25cm]frame.north east);
\node[anchor=west,text=white,font=\sffamily\bfseries\LARGE] at ([xshift=5mm,yshift=-0.62cm]frame.north west) {Question 9};
\node[anchor=east,text=white,font=\sffamily\itshape\large] at ([xshift=-5mm,yshift=-0.62cm]frame.north east) {Suites};}]
\vspace*{1.75cm}
\begin{minipage}[t][6.9cm][c]{\linewidth}\LARGE\centering
\begin{minipage}{0.95\linewidth}\raggedright\linespread{1.25}\selectfont $(u_n)$ est une suite arithmétique avec $u_{3}=7$ et $u_{8}=22$. Quelle est sa raison~?\end{minipage}
\end{minipage}

\vfill
\begin{tcolorbox}[enhanced,nobeforeafter,width=7.1cm,height=2.1cm,valign=center,colback=csui!6,colframe=csui,arc=2.5mm,boxrule=0.9pt,left=3mm]{\Large\bfseries\color{csui}a)}\quad\begin{minipage}[c]{5.3cm}\raggedright\Large $\dfrac{15}{11}$\end{minipage}\end{tcolorbox}\hfill\begin{tcolorbox}[enhanced,nobeforeafter,width=7.1cm,height=2.1cm,valign=center,colback=csui!6,colframe=csui,arc=2.5mm,boxrule=0.9pt,left=3mm]{\Large\bfseries\color{csui}b)}\quad\begin{minipage}[c]{5.3cm}\raggedright\Large $15$\end{minipage}\end{tcolorbox}\par\vspace{2.5mm}
\begin{tcolorbox}[enhanced,nobeforeafter,width=7.1cm,height=2.1cm,valign=center,colback=csui!6,colframe=csui,arc=2.5mm,boxrule=0.9pt,left=3mm]{\Large\bfseries\color{csui}c)}\quad\begin{minipage}[c]{5.3cm}\raggedright\Large $3$\end{minipage}\end{tcolorbox}\hfill\begin{tcolorbox}[enhanced,nobeforeafter,width=7.1cm,height=2.1cm,valign=center,colback=csui!6,colframe=csui,arc=2.5mm,boxrule=0.9pt,left=3mm]{\Large\bfseries\color{csui}d)}\quad\begin{minipage}[c]{5.3cm}\raggedright\Large $-3$\end{minipage}\end{tcolorbox}
\end{tcolorbox}
\newpage
\begin{tcolorbox}[enhanced,width=16cm,height=14cm,colback=white,colframe=csui,boxrule=2pt,arc=0mm,
left=6mm,right=6mm,top=0mm,bottom=4mm,sharp corners,
overlay={\fill[csui] (frame.north west) rectangle ([yshift=-1.25cm]frame.north east);
\node[anchor=west,text=white,font=\sffamily\bfseries\LARGE] at ([xshift=5mm,yshift=-0.62cm]frame.north west) {Question 10};
\node[anchor=east,text=white,font=\sffamily\itshape\large] at ([xshift=-5mm,yshift=-0.62cm]frame.north east) {Suites};}]
\vspace*{1.75cm}
\begin{minipage}[t][6.9cm][c]{\linewidth}\LARGE\centering
\begin{minipage}{0.95\linewidth}\raggedright\linespread{1.25}\selectfont $(u_n)$ est une suite arithmétique avec $u_{2}=10$ et $u_{6}=-2$. Quelle est sa raison~?\end{minipage}
\end{minipage}

\vfill
\begin{tcolorbox}[enhanced,nobeforeafter,width=7.1cm,height=2.1cm,valign=center,colback=csui!6,colframe=csui,arc=2.5mm,boxrule=0.9pt,left=3mm]{\Large\bfseries\color{csui}a)}\quad\begin{minipage}[c]{5.3cm}\raggedright\Large $-12$\end{minipage}\end{tcolorbox}\hfill\begin{tcolorbox}[enhanced,nobeforeafter,width=7.1cm,height=2.1cm,valign=center,colback=csui!6,colframe=csui,arc=2.5mm,boxrule=0.9pt,left=3mm]{\Large\bfseries\color{csui}b)}\quad\begin{minipage}[c]{5.3cm}\raggedright\Large $3$\end{minipage}\end{tcolorbox}\par\vspace{2.5mm}
\begin{tcolorbox}[enhanced,nobeforeafter,width=7.1cm,height=2.1cm,valign=center,colback=csui!6,colframe=csui,arc=2.5mm,boxrule=0.9pt,left=3mm]{\Large\bfseries\color{csui}c)}\quad\begin{minipage}[c]{5.3cm}\raggedright\Large $-\dfrac{3}{2}$\end{minipage}\end{tcolorbox}\hfill\begin{tcolorbox}[enhanced,nobeforeafter,width=7.1cm,height=2.1cm,valign=center,colback=csui!6,colframe=csui,arc=2.5mm,boxrule=0.9pt,left=3mm]{\Large\bfseries\color{csui}d)}\quad\begin{minipage}[c]{5.3cm}\raggedright\Large $-3$\end{minipage}\end{tcolorbox}
\end{tcolorbox}
\newpage
\begin{tcolorbox}[enhanced,width=16cm,height=14cm,colback=white,colframe=csui,boxrule=2pt,arc=0mm,
left=6mm,right=6mm,top=0mm,bottom=4mm,sharp corners,
overlay={\fill[csui] (frame.north west) rectangle ([yshift=-1.25cm]frame.north east);
\node[anchor=west,text=white,font=\sffamily\bfseries\LARGE] at ([xshift=5mm,yshift=-0.62cm]frame.north west) {Question 11};
\node[anchor=east,text=white,font=\sffamily\itshape\large] at ([xshift=-5mm,yshift=-0.62cm]frame.north east) {Suites};}]
\vspace*{1.75cm}
\begin{minipage}[t][6.9cm][c]{\linewidth}\LARGE\centering
\begin{minipage}{0.95\linewidth}\raggedright\linespread{1.25}\selectfont Que vaut la somme $1+2+2^2+\dots+2^{6}$~?\end{minipage}
\end{minipage}

\vfill
\begin{tcolorbox}[enhanced,nobeforeafter,width=7.1cm,height=2.1cm,valign=center,colback=csui!6,colframe=csui,arc=2.5mm,boxrule=0.9pt,left=3mm]{\Large\bfseries\color{csui}a)}\quad\begin{minipage}[c]{5.3cm}\raggedright\Large $127$\end{minipage}\end{tcolorbox}\hfill\begin{tcolorbox}[enhanced,nobeforeafter,width=7.1cm,height=2.1cm,valign=center,colback=csui!6,colframe=csui,arc=2.5mm,boxrule=0.9pt,left=3mm]{\Large\bfseries\color{csui}b)}\quad\begin{minipage}[c]{5.3cm}\raggedright\Large $63$\end{minipage}\end{tcolorbox}\par\vspace{2.5mm}
\begin{tcolorbox}[enhanced,nobeforeafter,width=7.1cm,height=2.1cm,valign=center,colback=csui!6,colframe=csui,arc=2.5mm,boxrule=0.9pt,left=3mm]{\Large\bfseries\color{csui}c)}\quad\begin{minipage}[c]{5.3cm}\raggedright\Large $64$\end{minipage}\end{tcolorbox}\hfill\begin{tcolorbox}[enhanced,nobeforeafter,width=7.1cm,height=2.1cm,valign=center,colback=csui!6,colframe=csui,arc=2.5mm,boxrule=0.9pt,left=3mm]{\Large\bfseries\color{csui}d)}\quad\begin{minipage}[c]{5.3cm}\raggedright\Large $128$\end{minipage}\end{tcolorbox}
\end{tcolorbox}
\newpage
\begin{tcolorbox}[enhanced,width=16cm,height=14cm,colback=white,colframe=csui,boxrule=2pt,arc=0mm,
left=6mm,right=6mm,top=0mm,bottom=4mm,sharp corners,
overlay={\fill[csui] (frame.north west) rectangle ([yshift=-1.25cm]frame.north east);
\node[anchor=west,text=white,font=\sffamily\bfseries\LARGE] at ([xshift=5mm,yshift=-0.62cm]frame.north west) {Question 12};
\node[anchor=east,text=white,font=\sffamily\itshape\large] at ([xshift=-5mm,yshift=-0.62cm]frame.north east) {Suites};}]
\vspace*{1.75cm}
\begin{minipage}[t][6.9cm][c]{\linewidth}\LARGE\centering
\begin{minipage}{0.95\linewidth}\raggedright\linespread{1.25}\selectfont Que vaut la somme $1+3+3^2+\dots+3^{4}$~?\end{minipage}
\end{minipage}

\vfill
\begin{tcolorbox}[enhanced,nobeforeafter,width=7.1cm,height=2.1cm,valign=center,colback=csui!6,colframe=csui,arc=2.5mm,boxrule=0.9pt,left=3mm]{\Large\bfseries\color{csui}a)}\quad\begin{minipage}[c]{5.3cm}\raggedright\Large $243$\end{minipage}\end{tcolorbox}\hfill\begin{tcolorbox}[enhanced,nobeforeafter,width=7.1cm,height=2.1cm,valign=center,colback=csui!6,colframe=csui,arc=2.5mm,boxrule=0.9pt,left=3mm]{\Large\bfseries\color{csui}b)}\quad\begin{minipage}[c]{5.3cm}\raggedright\Large $121$\end{minipage}\end{tcolorbox}\par\vspace{2.5mm}
\begin{tcolorbox}[enhanced,nobeforeafter,width=7.1cm,height=2.1cm,valign=center,colback=csui!6,colframe=csui,arc=2.5mm,boxrule=0.9pt,left=3mm]{\Large\bfseries\color{csui}c)}\quad\begin{minipage}[c]{5.3cm}\raggedright\Large $242$\end{minipage}\end{tcolorbox}\hfill\begin{tcolorbox}[enhanced,nobeforeafter,width=7.1cm,height=2.1cm,valign=center,colback=csui!6,colframe=csui,arc=2.5mm,boxrule=0.9pt,left=3mm]{\Large\bfseries\color{csui}d)}\quad\begin{minipage}[c]{5.3cm}\raggedright\Large $40$\end{minipage}\end{tcolorbox}
\end{tcolorbox}
\newpage
\begin{tcolorbox}[enhanced,width=16cm,height=14cm,colback=white,colframe=csui,boxrule=2pt,arc=0mm,
left=6mm,right=6mm,top=0mm,bottom=4mm,sharp corners,
overlay={\fill[csui] (frame.north west) rectangle ([yshift=-1.25cm]frame.north east);
\node[anchor=west,text=white,font=\sffamily\bfseries\LARGE] at ([xshift=5mm,yshift=-0.62cm]frame.north west) {Question 13};
\node[anchor=east,text=white,font=\sffamily\itshape\large] at ([xshift=-5mm,yshift=-0.62cm]frame.north east) {Suites};}]
\vspace*{1.75cm}
\begin{minipage}[t][6.9cm][c]{\linewidth}\LARGE\centering
\begin{minipage}{0.95\linewidth}\raggedright\linespread{1.25}\selectfont Quel est le sens de variation de la suite définie pour tout entier naturel $n$ par $u_n=5-3n$~?\end{minipage}
\end{minipage}

\vfill
\begin{tcolorbox}[enhanced,nobeforeafter,width=7.1cm,height=2.1cm,valign=center,colback=csui!6,colframe=csui,arc=2.5mm,boxrule=0.9pt,left=3mm]{\Large\bfseries\color{csui}a)}\quad\begin{minipage}[c]{5.3cm}\raggedright\large ni croissante ni décroissante\end{minipage}\end{tcolorbox}\hfill\begin{tcolorbox}[enhanced,nobeforeafter,width=7.1cm,height=2.1cm,valign=center,colback=csui!6,colframe=csui,arc=2.5mm,boxrule=0.9pt,left=3mm]{\Large\bfseries\color{csui}b)}\quad\begin{minipage}[c]{5.3cm}\raggedright\Large constante\end{minipage}\end{tcolorbox}\par\vspace{2.5mm}
\begin{tcolorbox}[enhanced,nobeforeafter,width=7.1cm,height=2.1cm,valign=center,colback=csui!6,colframe=csui,arc=2.5mm,boxrule=0.9pt,left=3mm]{\Large\bfseries\color{csui}c)}\quad\begin{minipage}[c]{5.3cm}\raggedright\Large décroissante\end{minipage}\end{tcolorbox}\hfill\begin{tcolorbox}[enhanced,nobeforeafter,width=7.1cm,height=2.1cm,valign=center,colback=csui!6,colframe=csui,arc=2.5mm,boxrule=0.9pt,left=3mm]{\Large\bfseries\color{csui}d)}\quad\begin{minipage}[c]{5.3cm}\raggedright\Large croissante\end{minipage}\end{tcolorbox}
\end{tcolorbox}
\newpage
\begin{tcolorbox}[enhanced,width=16cm,height=14cm,colback=white,colframe=csui,boxrule=2pt,arc=0mm,
left=6mm,right=6mm,top=0mm,bottom=4mm,sharp corners,
overlay={\fill[csui] (frame.north west) rectangle ([yshift=-1.25cm]frame.north east);
\node[anchor=west,text=white,font=\sffamily\bfseries\LARGE] at ([xshift=5mm,yshift=-0.62cm]frame.north west) {Question 14};
\node[anchor=east,text=white,font=\sffamily\itshape\large] at ([xshift=-5mm,yshift=-0.62cm]frame.north east) {Suites};}]
\vspace*{1.75cm}
\begin{minipage}[t][6.9cm][c]{\linewidth}\LARGE\centering
\begin{minipage}{0.95\linewidth}\raggedright\linespread{1.25}\selectfont Quel est le sens de variation de la suite définie pour tout entier naturel $n$ par $u_n=2n+7$~?\end{minipage}
\end{minipage}

\vfill
\begin{tcolorbox}[enhanced,nobeforeafter,width=7.1cm,height=2.1cm,valign=center,colback=csui!6,colframe=csui,arc=2.5mm,boxrule=0.9pt,left=3mm]{\Large\bfseries\color{csui}a)}\quad\begin{minipage}[c]{5.3cm}\raggedright\large ni croissante ni décroissante\end{minipage}\end{tcolorbox}\hfill\begin{tcolorbox}[enhanced,nobeforeafter,width=7.1cm,height=2.1cm,valign=center,colback=csui!6,colframe=csui,arc=2.5mm,boxrule=0.9pt,left=3mm]{\Large\bfseries\color{csui}b)}\quad\begin{minipage}[c]{5.3cm}\raggedright\Large décroissante\end{minipage}\end{tcolorbox}\par\vspace{2.5mm}
\begin{tcolorbox}[enhanced,nobeforeafter,width=7.1cm,height=2.1cm,valign=center,colback=csui!6,colframe=csui,arc=2.5mm,boxrule=0.9pt,left=3mm]{\Large\bfseries\color{csui}c)}\quad\begin{minipage}[c]{5.3cm}\raggedright\Large croissante\end{minipage}\end{tcolorbox}\hfill\begin{tcolorbox}[enhanced,nobeforeafter,width=7.1cm,height=2.1cm,valign=center,colback=csui!6,colframe=csui,arc=2.5mm,boxrule=0.9pt,left=3mm]{\Large\bfseries\color{csui}d)}\quad\begin{minipage}[c]{5.3cm}\raggedright\Large constante\end{minipage}\end{tcolorbox}
\end{tcolorbox}
\newpage
\begin{tcolorbox}[enhanced,width=16cm,height=14cm,colback=white,colframe=csui,boxrule=2pt,arc=0mm,
left=6mm,right=6mm,top=0mm,bottom=4mm,sharp corners,
overlay={\fill[csui] (frame.north west) rectangle ([yshift=-1.25cm]frame.north east);
\node[anchor=west,text=white,font=\sffamily\bfseries\LARGE] at ([xshift=5mm,yshift=-0.62cm]frame.north west) {Question 15};
\node[anchor=east,text=white,font=\sffamily\itshape\large] at ([xshift=-5mm,yshift=-0.62cm]frame.north east) {Suites};}]
\vspace*{1.75cm}
\begin{minipage}[t][6.9cm][c]{\linewidth}\LARGE\centering
\begin{minipage}{0.95\linewidth}\raggedright\linespread{1.25}\selectfont Pour tout entier naturel $n$, $u_n=n^2-3n+2$. Que vaut $u_{5}$~?\end{minipage}
\end{minipage}

\vfill
\begin{tcolorbox}[enhanced,nobeforeafter,width=7.1cm,height=2.1cm,valign=center,colback=csui!6,colframe=csui,arc=2.5mm,boxrule=0.9pt,left=3mm]{\Large\bfseries\color{csui}a)}\quad\begin{minipage}[c]{5.3cm}\raggedright\Large $20$\end{minipage}\end{tcolorbox}\hfill\begin{tcolorbox}[enhanced,nobeforeafter,width=7.1cm,height=2.1cm,valign=center,colback=csui!6,colframe=csui,arc=2.5mm,boxrule=0.9pt,left=3mm]{\Large\bfseries\color{csui}b)}\quad\begin{minipage}[c]{5.3cm}\raggedright\Large $-12$\end{minipage}\end{tcolorbox}\par\vspace{2.5mm}
\begin{tcolorbox}[enhanced,nobeforeafter,width=7.1cm,height=2.1cm,valign=center,colback=csui!6,colframe=csui,arc=2.5mm,boxrule=0.9pt,left=3mm]{\Large\bfseries\color{csui}c)}\quad\begin{minipage}[c]{5.3cm}\raggedright\Large $6$\end{minipage}\end{tcolorbox}\hfill\begin{tcolorbox}[enhanced,nobeforeafter,width=7.1cm,height=2.1cm,valign=center,colback=csui!6,colframe=csui,arc=2.5mm,boxrule=0.9pt,left=3mm]{\Large\bfseries\color{csui}d)}\quad\begin{minipage}[c]{5.3cm}\raggedright\Large $12$\end{minipage}\end{tcolorbox}
\end{tcolorbox}
\newpage
\begin{tcolorbox}[enhanced,width=16cm,height=14cm,colback=white,colframe=csui,boxrule=2pt,arc=0mm,
left=6mm,right=6mm,top=0mm,bottom=4mm,sharp corners,
overlay={\fill[csui] (frame.north west) rectangle ([yshift=-1.25cm]frame.north east);
\node[anchor=west,text=white,font=\sffamily\bfseries\LARGE] at ([xshift=5mm,yshift=-0.62cm]frame.north west) {Question 16};
\node[anchor=east,text=white,font=\sffamily\itshape\large] at ([xshift=-5mm,yshift=-0.62cm]frame.north east) {Suites};}]
\vspace*{1.75cm}
\begin{minipage}[t][6.9cm][c]{\linewidth}\LARGE\centering
\begin{minipage}{0.95\linewidth}\raggedright\linespread{1.25}\selectfont Pour tout entier naturel $n$, $u_n=2n^2+n-4$. Que vaut $u_{3}$~?\end{minipage}
\end{minipage}

\vfill
\begin{tcolorbox}[enhanced,nobeforeafter,width=7.1cm,height=2.1cm,valign=center,colback=csui!6,colframe=csui,arc=2.5mm,boxrule=0.9pt,left=3mm]{\Large\bfseries\color{csui}a)}\quad\begin{minipage}[c]{5.3cm}\raggedright\Large $32$\end{minipage}\end{tcolorbox}\hfill\begin{tcolorbox}[enhanced,nobeforeafter,width=7.1cm,height=2.1cm,valign=center,colback=csui!6,colframe=csui,arc=2.5mm,boxrule=0.9pt,left=3mm]{\Large\bfseries\color{csui}b)}\quad\begin{minipage}[c]{5.3cm}\raggedright\Large $-17$\end{minipage}\end{tcolorbox}\par\vspace{2.5mm}
\begin{tcolorbox}[enhanced,nobeforeafter,width=7.1cm,height=2.1cm,valign=center,colback=csui!6,colframe=csui,arc=2.5mm,boxrule=0.9pt,left=3mm]{\Large\bfseries\color{csui}c)}\quad\begin{minipage}[c]{5.3cm}\raggedright\Large $17$\end{minipage}\end{tcolorbox}\hfill\begin{tcolorbox}[enhanced,nobeforeafter,width=7.1cm,height=2.1cm,valign=center,colback=csui!6,colframe=csui,arc=2.5mm,boxrule=0.9pt,left=3mm]{\Large\bfseries\color{csui}d)}\quad\begin{minipage}[c]{5.3cm}\raggedright\Large $6$\end{minipage}\end{tcolorbox}
\end{tcolorbox}
\newpage
\begin{tcolorbox}[enhanced,width=16cm,height=14cm,colback=white,colframe=cdeg,boxrule=2pt,arc=0mm,
left=6mm,right=6mm,top=0mm,bottom=4mm,sharp corners,
overlay={\fill[cdeg] (frame.north west) rectangle ([yshift=-1.25cm]frame.north east);
\node[anchor=west,text=white,font=\sffamily\bfseries\LARGE] at ([xshift=5mm,yshift=-0.62cm]frame.north west) {Question 17};
\node[anchor=east,text=white,font=\sffamily\itshape\large] at ([xshift=-5mm,yshift=-0.62cm]frame.north east) {Second degré};}]
\vspace*{1.75cm}
\begin{minipage}[t][6.9cm][c]{\linewidth}\LARGE\centering
\begin{minipage}{0.95\linewidth}\raggedright\linespread{1.25}\selectfont Quel est le discriminant du trinôme $x^2-5x+6$~?\end{minipage}
\end{minipage}

\vfill
\begin{tcolorbox}[enhanced,nobeforeafter,width=7.1cm,height=2.1cm,valign=center,colback=cdeg!6,colframe=cdeg,arc=2.5mm,boxrule=0.9pt,left=3mm]{\Large\bfseries\color{cdeg}a)}\quad\begin{minipage}[c]{5.3cm}\raggedright\Large $-29$\end{minipage}\end{tcolorbox}\hfill\begin{tcolorbox}[enhanced,nobeforeafter,width=7.1cm,height=2.1cm,valign=center,colback=cdeg!6,colframe=cdeg,arc=2.5mm,boxrule=0.9pt,left=3mm]{\Large\bfseries\color{cdeg}b)}\quad\begin{minipage}[c]{5.3cm}\raggedright\Large $1$\end{minipage}\end{tcolorbox}\par\vspace{2.5mm}
\begin{tcolorbox}[enhanced,nobeforeafter,width=7.1cm,height=2.1cm,valign=center,colback=cdeg!6,colframe=cdeg,arc=2.5mm,boxrule=0.9pt,left=3mm]{\Large\bfseries\color{cdeg}c)}\quad\begin{minipage}[c]{5.3cm}\raggedright\Large $49$\end{minipage}\end{tcolorbox}\hfill\begin{tcolorbox}[enhanced,nobeforeafter,width=7.1cm,height=2.1cm,valign=center,colback=cdeg!6,colframe=cdeg,arc=2.5mm,boxrule=0.9pt,left=3mm]{\Large\bfseries\color{cdeg}d)}\quad\begin{minipage}[c]{5.3cm}\raggedright\Large $-1$\end{minipage}\end{tcolorbox}
\end{tcolorbox}
\newpage
\begin{tcolorbox}[enhanced,width=16cm,height=14cm,colback=white,colframe=cdeg,boxrule=2pt,arc=0mm,
left=6mm,right=6mm,top=0mm,bottom=4mm,sharp corners,
overlay={\fill[cdeg] (frame.north west) rectangle ([yshift=-1.25cm]frame.north east);
\node[anchor=west,text=white,font=\sffamily\bfseries\LARGE] at ([xshift=5mm,yshift=-0.62cm]frame.north west) {Question 18};
\node[anchor=east,text=white,font=\sffamily\itshape\large] at ([xshift=-5mm,yshift=-0.62cm]frame.north east) {Second degré};}]
\vspace*{1.75cm}
\begin{minipage}[t][6.9cm][c]{\linewidth}\LARGE\centering
\begin{minipage}{0.95\linewidth}\raggedright\linespread{1.25}\selectfont Quel est le discriminant du trinôme $2x^2+3x-2$~?\end{minipage}
\end{minipage}

\vfill
\begin{tcolorbox}[enhanced,nobeforeafter,width=7.1cm,height=2.1cm,valign=center,colback=cdeg!6,colframe=cdeg,arc=2.5mm,boxrule=0.9pt,left=3mm]{\Large\bfseries\color{cdeg}a)}\quad\begin{minipage}[c]{5.3cm}\raggedright\Large $19$\end{minipage}\end{tcolorbox}\hfill\begin{tcolorbox}[enhanced,nobeforeafter,width=7.1cm,height=2.1cm,valign=center,colback=cdeg!6,colframe=cdeg,arc=2.5mm,boxrule=0.9pt,left=3mm]{\Large\bfseries\color{cdeg}b)}\quad\begin{minipage}[c]{5.3cm}\raggedright\Large $-7$\end{minipage}\end{tcolorbox}\par\vspace{2.5mm}
\begin{tcolorbox}[enhanced,nobeforeafter,width=7.1cm,height=2.1cm,valign=center,colback=cdeg!6,colframe=cdeg,arc=2.5mm,boxrule=0.9pt,left=3mm]{\Large\bfseries\color{cdeg}c)}\quad\begin{minipage}[c]{5.3cm}\raggedright\Large $-25$\end{minipage}\end{tcolorbox}\hfill\begin{tcolorbox}[enhanced,nobeforeafter,width=7.1cm,height=2.1cm,valign=center,colback=cdeg!6,colframe=cdeg,arc=2.5mm,boxrule=0.9pt,left=3mm]{\Large\bfseries\color{cdeg}d)}\quad\begin{minipage}[c]{5.3cm}\raggedright\Large $25$\end{minipage}\end{tcolorbox}
\end{tcolorbox}
\newpage
\begin{tcolorbox}[enhanced,width=16cm,height=14cm,colback=white,colframe=cdeg,boxrule=2pt,arc=0mm,
left=6mm,right=6mm,top=0mm,bottom=4mm,sharp corners,
overlay={\fill[cdeg] (frame.north west) rectangle ([yshift=-1.25cm]frame.north east);
\node[anchor=west,text=white,font=\sffamily\bfseries\LARGE] at ([xshift=5mm,yshift=-0.62cm]frame.north west) {Question 19};
\node[anchor=east,text=white,font=\sffamily\itshape\large] at ([xshift=-5mm,yshift=-0.62cm]frame.north east) {Second degré};}]
\vspace*{1.75cm}
\begin{minipage}[t][6.9cm][c]{\linewidth}\LARGE\centering
\begin{minipage}{0.95\linewidth}\raggedright\linespread{1.25}\selectfont Combien de racines réelles le trinôme $x^2+2x+5$ possède-t-il~?\end{minipage}
\end{minipage}

\vfill
\begin{tcolorbox}[enhanced,nobeforeafter,width=7.1cm,height=2.1cm,valign=center,colback=cdeg!6,colframe=cdeg,arc=2.5mm,boxrule=0.9pt,left=3mm]{\Large\bfseries\color{cdeg}a)}\quad\begin{minipage}[c]{5.3cm}\raggedright\Large une infinité\end{minipage}\end{tcolorbox}\hfill\begin{tcolorbox}[enhanced,nobeforeafter,width=7.1cm,height=2.1cm,valign=center,colback=cdeg!6,colframe=cdeg,arc=2.5mm,boxrule=0.9pt,left=3mm]{\Large\bfseries\color{cdeg}b)}\quad\begin{minipage}[c]{5.3cm}\raggedright\Large $1$\end{minipage}\end{tcolorbox}\par\vspace{2.5mm}
\begin{tcolorbox}[enhanced,nobeforeafter,width=7.1cm,height=2.1cm,valign=center,colback=cdeg!6,colframe=cdeg,arc=2.5mm,boxrule=0.9pt,left=3mm]{\Large\bfseries\color{cdeg}c)}\quad\begin{minipage}[c]{5.3cm}\raggedright\Large $2$\end{minipage}\end{tcolorbox}\hfill\begin{tcolorbox}[enhanced,nobeforeafter,width=7.1cm,height=2.1cm,valign=center,colback=cdeg!6,colframe=cdeg,arc=2.5mm,boxrule=0.9pt,left=3mm]{\Large\bfseries\color{cdeg}d)}\quad\begin{minipage}[c]{5.3cm}\raggedright\Large $0$\end{minipage}\end{tcolorbox}
\end{tcolorbox}
\newpage
\begin{tcolorbox}[enhanced,width=16cm,height=14cm,colback=white,colframe=cdeg,boxrule=2pt,arc=0mm,
left=6mm,right=6mm,top=0mm,bottom=4mm,sharp corners,
overlay={\fill[cdeg] (frame.north west) rectangle ([yshift=-1.25cm]frame.north east);
\node[anchor=west,text=white,font=\sffamily\bfseries\LARGE] at ([xshift=5mm,yshift=-0.62cm]frame.north west) {Question 20};
\node[anchor=east,text=white,font=\sffamily\itshape\large] at ([xshift=-5mm,yshift=-0.62cm]frame.north east) {Second degré};}]
\vspace*{1.75cm}
\begin{minipage}[t][6.9cm][c]{\linewidth}\LARGE\centering
\begin{minipage}{0.95\linewidth}\raggedright\linespread{1.25}\selectfont Combien de racines réelles le trinôme $x^2-6x+9$ possède-t-il~?\end{minipage}
\end{minipage}

\vfill
\begin{tcolorbox}[enhanced,nobeforeafter,width=7.1cm,height=2.1cm,valign=center,colback=cdeg!6,colframe=cdeg,arc=2.5mm,boxrule=0.9pt,left=3mm]{\Large\bfseries\color{cdeg}a)}\quad\begin{minipage}[c]{5.3cm}\raggedright\Large $2$\end{minipage}\end{tcolorbox}\hfill\begin{tcolorbox}[enhanced,nobeforeafter,width=7.1cm,height=2.1cm,valign=center,colback=cdeg!6,colframe=cdeg,arc=2.5mm,boxrule=0.9pt,left=3mm]{\Large\bfseries\color{cdeg}b)}\quad\begin{minipage}[c]{5.3cm}\raggedright\Large $1$\end{minipage}\end{tcolorbox}\par\vspace{2.5mm}
\begin{tcolorbox}[enhanced,nobeforeafter,width=7.1cm,height=2.1cm,valign=center,colback=cdeg!6,colframe=cdeg,arc=2.5mm,boxrule=0.9pt,left=3mm]{\Large\bfseries\color{cdeg}c)}\quad\begin{minipage}[c]{5.3cm}\raggedright\Large $0$\end{minipage}\end{tcolorbox}\hfill\begin{tcolorbox}[enhanced,nobeforeafter,width=7.1cm,height=2.1cm,valign=center,colback=cdeg!6,colframe=cdeg,arc=2.5mm,boxrule=0.9pt,left=3mm]{\Large\bfseries\color{cdeg}d)}\quad\begin{minipage}[c]{5.3cm}\raggedright\Large une infinité\end{minipage}\end{tcolorbox}
\end{tcolorbox}
\newpage
\begin{tcolorbox}[enhanced,width=16cm,height=14cm,colback=white,colframe=cdeg,boxrule=2pt,arc=0mm,
left=6mm,right=6mm,top=0mm,bottom=4mm,sharp corners,
overlay={\fill[cdeg] (frame.north west) rectangle ([yshift=-1.25cm]frame.north east);
\node[anchor=west,text=white,font=\sffamily\bfseries\LARGE] at ([xshift=5mm,yshift=-0.62cm]frame.north west) {Question 21};
\node[anchor=east,text=white,font=\sffamily\itshape\large] at ([xshift=-5mm,yshift=-0.62cm]frame.north east) {Second degré};}]
\vspace*{1.75cm}
\begin{minipage}[t][6.9cm][c]{\linewidth}\LARGE\centering
\begin{minipage}{0.95\linewidth}\raggedright\linespread{1.25}\selectfont Quelles sont les racines du trinôme $x^2-5x+6$~?\end{minipage}
\end{minipage}

\vfill
\begin{tcolorbox}[enhanced,nobeforeafter,width=7.1cm,height=2.1cm,valign=center,colback=cdeg!6,colframe=cdeg,arc=2.5mm,boxrule=0.9pt,left=3mm]{\Large\bfseries\color{cdeg}a)}\quad\begin{minipage}[c]{5.3cm}\raggedright\Large $\{2\,;\,3\}$\end{minipage}\end{tcolorbox}\hfill\begin{tcolorbox}[enhanced,nobeforeafter,width=7.1cm,height=2.1cm,valign=center,colback=cdeg!6,colframe=cdeg,arc=2.5mm,boxrule=0.9pt,left=3mm]{\Large\bfseries\color{cdeg}b)}\quad\begin{minipage}[c]{5.3cm}\raggedright\Large $\{5\,;\,6\}$\end{minipage}\end{tcolorbox}\par\vspace{2.5mm}
\begin{tcolorbox}[enhanced,nobeforeafter,width=7.1cm,height=2.1cm,valign=center,colback=cdeg!6,colframe=cdeg,arc=2.5mm,boxrule=0.9pt,left=3mm]{\Large\bfseries\color{cdeg}c)}\quad\begin{minipage}[c]{5.3cm}\raggedright\Large $\{-3\,;\,2\}$\end{minipage}\end{tcolorbox}\hfill\begin{tcolorbox}[enhanced,nobeforeafter,width=7.1cm,height=2.1cm,valign=center,colback=cdeg!6,colframe=cdeg,arc=2.5mm,boxrule=0.9pt,left=3mm]{\Large\bfseries\color{cdeg}d)}\quad\begin{minipage}[c]{5.3cm}\raggedright\Large $\{-3\,;\,-2\}$\end{minipage}\end{tcolorbox}
\end{tcolorbox}
\newpage
\begin{tcolorbox}[enhanced,width=16cm,height=14cm,colback=white,colframe=cdeg,boxrule=2pt,arc=0mm,
left=6mm,right=6mm,top=0mm,bottom=4mm,sharp corners,
overlay={\fill[cdeg] (frame.north west) rectangle ([yshift=-1.25cm]frame.north east);
\node[anchor=west,text=white,font=\sffamily\bfseries\LARGE] at ([xshift=5mm,yshift=-0.62cm]frame.north west) {Question 22};
\node[anchor=east,text=white,font=\sffamily\itshape\large] at ([xshift=-5mm,yshift=-0.62cm]frame.north east) {Second degré};}]
\vspace*{1.75cm}
\begin{minipage}[t][6.9cm][c]{\linewidth}\LARGE\centering
\begin{minipage}{0.95\linewidth}\raggedright\linespread{1.25}\selectfont Quelles sont les racines du trinôme $x^2-3x-4$~?\end{minipage}
\end{minipage}

\vfill
\begin{tcolorbox}[enhanced,nobeforeafter,width=7.1cm,height=2.1cm,valign=center,colback=cdeg!6,colframe=cdeg,arc=2.5mm,boxrule=0.9pt,left=3mm]{\Large\bfseries\color{cdeg}a)}\quad\begin{minipage}[c]{5.3cm}\raggedright\Large $\{-4\,;\,3\}$\end{minipage}\end{tcolorbox}\hfill\begin{tcolorbox}[enhanced,nobeforeafter,width=7.1cm,height=2.1cm,valign=center,colback=cdeg!6,colframe=cdeg,arc=2.5mm,boxrule=0.9pt,left=3mm]{\Large\bfseries\color{cdeg}b)}\quad\begin{minipage}[c]{5.3cm}\raggedright\Large $\{-4\,;\,-1\}$\end{minipage}\end{tcolorbox}\par\vspace{2.5mm}
\begin{tcolorbox}[enhanced,nobeforeafter,width=7.1cm,height=2.1cm,valign=center,colback=cdeg!6,colframe=cdeg,arc=2.5mm,boxrule=0.9pt,left=3mm]{\Large\bfseries\color{cdeg}c)}\quad\begin{minipage}[c]{5.3cm}\raggedright\Large $\{-4\,;\,1\}$\end{minipage}\end{tcolorbox}\hfill\begin{tcolorbox}[enhanced,nobeforeafter,width=7.1cm,height=2.1cm,valign=center,colback=cdeg!6,colframe=cdeg,arc=2.5mm,boxrule=0.9pt,left=3mm]{\Large\bfseries\color{cdeg}d)}\quad\begin{minipage}[c]{5.3cm}\raggedright\Large $\{-1\,;\,4\}$\end{minipage}\end{tcolorbox}
\end{tcolorbox}
\newpage
\begin{tcolorbox}[enhanced,width=16cm,height=14cm,colback=white,colframe=cdeg,boxrule=2pt,arc=0mm,
left=6mm,right=6mm,top=0mm,bottom=4mm,sharp corners,
overlay={\fill[cdeg] (frame.north west) rectangle ([yshift=-1.25cm]frame.north east);
\node[anchor=west,text=white,font=\sffamily\bfseries\LARGE] at ([xshift=5mm,yshift=-0.62cm]frame.north west) {Question 23};
\node[anchor=east,text=white,font=\sffamily\itshape\large] at ([xshift=-5mm,yshift=-0.62cm]frame.north east) {Second degré};}]
\vspace*{1.75cm}
\begin{minipage}[t][6.9cm][c]{\linewidth}\LARGE\centering
\begin{minipage}{0.95\linewidth}\raggedright\linespread{1.25}\selectfont Quelle est la forme canonique de $x^2+6x+5$~?\end{minipage}
\end{minipage}

\vfill
\begin{tcolorbox}[enhanced,nobeforeafter,width=7.1cm,height=2.1cm,valign=center,colback=cdeg!6,colframe=cdeg,arc=2.5mm,boxrule=0.9pt,left=3mm]{\Large\bfseries\color{cdeg}a)}\quad\begin{minipage}[c]{5.3cm}\raggedright\Large $(x+6)^2-4$\end{minipage}\end{tcolorbox}\hfill\begin{tcolorbox}[enhanced,nobeforeafter,width=7.1cm,height=2.1cm,valign=center,colback=cdeg!6,colframe=cdeg,arc=2.5mm,boxrule=0.9pt,left=3mm]{\Large\bfseries\color{cdeg}b)}\quad\begin{minipage}[c]{5.3cm}\raggedright\Large $(x+3)^2-4$\end{minipage}\end{tcolorbox}\par\vspace{2.5mm}
\begin{tcolorbox}[enhanced,nobeforeafter,width=7.1cm,height=2.1cm,valign=center,colback=cdeg!6,colframe=cdeg,arc=2.5mm,boxrule=0.9pt,left=3mm]{\Large\bfseries\color{cdeg}c)}\quad\begin{minipage}[c]{5.3cm}\raggedright\Large $(x-3)^2-4$\end{minipage}\end{tcolorbox}\hfill\begin{tcolorbox}[enhanced,nobeforeafter,width=7.1cm,height=2.1cm,valign=center,colback=cdeg!6,colframe=cdeg,arc=2.5mm,boxrule=0.9pt,left=3mm]{\Large\bfseries\color{cdeg}d)}\quad\begin{minipage}[c]{5.3cm}\raggedright\Large $(x+3)^2+5$\end{minipage}\end{tcolorbox}
\end{tcolorbox}
\newpage
\begin{tcolorbox}[enhanced,width=16cm,height=14cm,colback=white,colframe=cdeg,boxrule=2pt,arc=0mm,
left=6mm,right=6mm,top=0mm,bottom=4mm,sharp corners,
overlay={\fill[cdeg] (frame.north west) rectangle ([yshift=-1.25cm]frame.north east);
\node[anchor=west,text=white,font=\sffamily\bfseries\LARGE] at ([xshift=5mm,yshift=-0.62cm]frame.north west) {Question 24};
\node[anchor=east,text=white,font=\sffamily\itshape\large] at ([xshift=-5mm,yshift=-0.62cm]frame.north east) {Second degré};}]
\vspace*{1.75cm}
\begin{minipage}[t][6.9cm][c]{\linewidth}\LARGE\centering
\begin{minipage}{0.95\linewidth}\raggedright\linespread{1.25}\selectfont Quelle est la forme canonique de $x^2-4x+1$~?\end{minipage}
\end{minipage}

\vfill
\begin{tcolorbox}[enhanced,nobeforeafter,width=7.1cm,height=2.1cm,valign=center,colback=cdeg!6,colframe=cdeg,arc=2.5mm,boxrule=0.9pt,left=3mm]{\Large\bfseries\color{cdeg}a)}\quad\begin{minipage}[c]{5.3cm}\raggedright\Large $(x+2)^2-3$\end{minipage}\end{tcolorbox}\hfill\begin{tcolorbox}[enhanced,nobeforeafter,width=7.1cm,height=2.1cm,valign=center,colback=cdeg!6,colframe=cdeg,arc=2.5mm,boxrule=0.9pt,left=3mm]{\Large\bfseries\color{cdeg}b)}\quad\begin{minipage}[c]{5.3cm}\raggedright\Large $(x-2)^2-3$\end{minipage}\end{tcolorbox}\par\vspace{2.5mm}
\begin{tcolorbox}[enhanced,nobeforeafter,width=7.1cm,height=2.1cm,valign=center,colback=cdeg!6,colframe=cdeg,arc=2.5mm,boxrule=0.9pt,left=3mm]{\Large\bfseries\color{cdeg}c)}\quad\begin{minipage}[c]{5.3cm}\raggedright\Large $(x-4)^2-3$\end{minipage}\end{tcolorbox}\hfill\begin{tcolorbox}[enhanced,nobeforeafter,width=7.1cm,height=2.1cm,valign=center,colback=cdeg!6,colframe=cdeg,arc=2.5mm,boxrule=0.9pt,left=3mm]{\Large\bfseries\color{cdeg}d)}\quad\begin{minipage}[c]{5.3cm}\raggedright\Large $(x-2)^2+1$\end{minipage}\end{tcolorbox}
\end{tcolorbox}
\newpage
\begin{tcolorbox}[enhanced,width=16cm,height=14cm,colback=white,colframe=cdeg,boxrule=2pt,arc=0mm,
left=6mm,right=6mm,top=0mm,bottom=4mm,sharp corners,
overlay={\fill[cdeg] (frame.north west) rectangle ([yshift=-1.25cm]frame.north east);
\node[anchor=west,text=white,font=\sffamily\bfseries\LARGE] at ([xshift=5mm,yshift=-0.62cm]frame.north west) {Question 25};
\node[anchor=east,text=white,font=\sffamily\itshape\large] at ([xshift=-5mm,yshift=-0.62cm]frame.north east) {Second degré};}]
\vspace*{1.75cm}
\begin{minipage}[t][6.9cm][c]{\linewidth}\LARGE\centering
\begin{minipage}{0.95\linewidth}\raggedright\linespread{1.25}\selectfont Quelles sont les coordonnées du sommet de la parabole d'équation $y=x^2-4x+1$~?\end{minipage}
\end{minipage}

\vfill
\begin{tcolorbox}[enhanced,nobeforeafter,width=7.1cm,height=2.1cm,valign=center,colback=cdeg!6,colframe=cdeg,arc=2.5mm,boxrule=0.9pt,left=3mm]{\Large\bfseries\color{cdeg}a)}\quad\begin{minipage}[c]{5.3cm}\raggedright\Large $(2\,;\,3)$\end{minipage}\end{tcolorbox}\hfill\begin{tcolorbox}[enhanced,nobeforeafter,width=7.1cm,height=2.1cm,valign=center,colback=cdeg!6,colframe=cdeg,arc=2.5mm,boxrule=0.9pt,left=3mm]{\Large\bfseries\color{cdeg}b)}\quad\begin{minipage}[c]{5.3cm}\raggedright\Large $(-2\,;\,-3)$\end{minipage}\end{tcolorbox}\par\vspace{2.5mm}
\begin{tcolorbox}[enhanced,nobeforeafter,width=7.1cm,height=2.1cm,valign=center,colback=cdeg!6,colframe=cdeg,arc=2.5mm,boxrule=0.9pt,left=3mm]{\Large\bfseries\color{cdeg}c)}\quad\begin{minipage}[c]{5.3cm}\raggedright\Large $(2\,;\,1)$\end{minipage}\end{tcolorbox}\hfill\begin{tcolorbox}[enhanced,nobeforeafter,width=7.1cm,height=2.1cm,valign=center,colback=cdeg!6,colframe=cdeg,arc=2.5mm,boxrule=0.9pt,left=3mm]{\Large\bfseries\color{cdeg}d)}\quad\begin{minipage}[c]{5.3cm}\raggedright\Large $(2\,;\,-3)$\end{minipage}\end{tcolorbox}
\end{tcolorbox}
\newpage
\begin{tcolorbox}[enhanced,width=16cm,height=14cm,colback=white,colframe=cdeg,boxrule=2pt,arc=0mm,
left=6mm,right=6mm,top=0mm,bottom=4mm,sharp corners,
overlay={\fill[cdeg] (frame.north west) rectangle ([yshift=-1.25cm]frame.north east);
\node[anchor=west,text=white,font=\sffamily\bfseries\LARGE] at ([xshift=5mm,yshift=-0.62cm]frame.north west) {Question 26};
\node[anchor=east,text=white,font=\sffamily\itshape\large] at ([xshift=-5mm,yshift=-0.62cm]frame.north east) {Second degré};}]
\vspace*{1.75cm}
\begin{minipage}[t][6.9cm][c]{\linewidth}\LARGE\centering
\begin{minipage}{0.95\linewidth}\raggedright\linespread{1.25}\selectfont Quelles sont les coordonnées du sommet de la parabole d'équation $y=2x^2+8x+3$~?\end{minipage}
\end{minipage}

\vfill
\begin{tcolorbox}[enhanced,nobeforeafter,width=7.1cm,height=2.1cm,valign=center,colback=cdeg!6,colframe=cdeg,arc=2.5mm,boxrule=0.9pt,left=3mm]{\Large\bfseries\color{cdeg}a)}\quad\begin{minipage}[c]{5.3cm}\raggedright\Large $(2\,;\,-5)$\end{minipage}\end{tcolorbox}\hfill\begin{tcolorbox}[enhanced,nobeforeafter,width=7.1cm,height=2.1cm,valign=center,colback=cdeg!6,colframe=cdeg,arc=2.5mm,boxrule=0.9pt,left=3mm]{\Large\bfseries\color{cdeg}b)}\quad\begin{minipage}[c]{5.3cm}\raggedright\Large $(-2\,;\,-5)$\end{minipage}\end{tcolorbox}\par\vspace{2.5mm}
\begin{tcolorbox}[enhanced,nobeforeafter,width=7.1cm,height=2.1cm,valign=center,colback=cdeg!6,colframe=cdeg,arc=2.5mm,boxrule=0.9pt,left=3mm]{\Large\bfseries\color{cdeg}c)}\quad\begin{minipage}[c]{5.3cm}\raggedright\Large $(-2\,;\,5)$\end{minipage}\end{tcolorbox}\hfill\begin{tcolorbox}[enhanced,nobeforeafter,width=7.1cm,height=2.1cm,valign=center,colback=cdeg!6,colframe=cdeg,arc=2.5mm,boxrule=0.9pt,left=3mm]{\Large\bfseries\color{cdeg}d)}\quad\begin{minipage}[c]{5.3cm}\raggedright\Large $(-2\,;\,3)$\end{minipage}\end{tcolorbox}
\end{tcolorbox}
\newpage
\begin{tcolorbox}[enhanced,width=16cm,height=14cm,colback=white,colframe=cdeg,boxrule=2pt,arc=0mm,
left=6mm,right=6mm,top=0mm,bottom=4mm,sharp corners,
overlay={\fill[cdeg] (frame.north west) rectangle ([yshift=-1.25cm]frame.north east);
\node[anchor=west,text=white,font=\sffamily\bfseries\LARGE] at ([xshift=5mm,yshift=-0.62cm]frame.north west) {Question 27};
\node[anchor=east,text=white,font=\sffamily\itshape\large] at ([xshift=-5mm,yshift=-0.62cm]frame.north east) {Second degré};}]
\vspace*{1.75cm}
\begin{minipage}[t][6.9cm][c]{\linewidth}\LARGE\centering
\begin{minipage}{0.95\linewidth}\raggedright\linespread{1.25}\selectfont Quel est l'ensemble des solutions de l'inéquation $x^2-4<0$~?\end{minipage}
\end{minipage}

\vfill
\begin{tcolorbox}[enhanced,nobeforeafter,width=7.1cm,height=2.1cm,valign=center,colback=cdeg!6,colframe=cdeg,arc=2.5mm,boxrule=0.9pt,left=3mm]{\Large\bfseries\color{cdeg}a)}\quad\begin{minipage}[c]{5.3cm}\raggedright\Large $[-2\,;\,2]$\end{minipage}\end{tcolorbox}\hfill\begin{tcolorbox}[enhanced,nobeforeafter,width=7.1cm,height=2.1cm,valign=center,colback=cdeg!6,colframe=cdeg,arc=2.5mm,boxrule=0.9pt,left=3mm]{\Large\bfseries\color{cdeg}b)}\quad\begin{minipage}[c]{5.3cm}\raggedright\Large $]-2\,;\,2[$\end{minipage}\end{tcolorbox}\par\vspace{2.5mm}
\begin{tcolorbox}[enhanced,nobeforeafter,width=7.1cm,height=2.1cm,valign=center,colback=cdeg!6,colframe=cdeg,arc=2.5mm,boxrule=0.9pt,left=3mm]{\Large\bfseries\color{cdeg}c)}\quad\begin{minipage}[c]{5.3cm}\raggedright\large $]-\infty\,;\,-2]\cup[2\,;\,+\infty[$\end{minipage}\end{tcolorbox}\hfill\begin{tcolorbox}[enhanced,nobeforeafter,width=7.1cm,height=2.1cm,valign=center,colback=cdeg!6,colframe=cdeg,arc=2.5mm,boxrule=0.9pt,left=3mm]{\Large\bfseries\color{cdeg}d)}\quad\begin{minipage}[c]{5.3cm}\raggedright\large $]-\infty\,;\,-2[\cup]2\,;\,+\infty[$\end{minipage}\end{tcolorbox}
\end{tcolorbox}
\newpage
\begin{tcolorbox}[enhanced,width=16cm,height=14cm,colback=white,colframe=cdeg,boxrule=2pt,arc=0mm,
left=6mm,right=6mm,top=0mm,bottom=4mm,sharp corners,
overlay={\fill[cdeg] (frame.north west) rectangle ([yshift=-1.25cm]frame.north east);
\node[anchor=west,text=white,font=\sffamily\bfseries\LARGE] at ([xshift=5mm,yshift=-0.62cm]frame.north west) {Question 28};
\node[anchor=east,text=white,font=\sffamily\itshape\large] at ([xshift=-5mm,yshift=-0.62cm]frame.north east) {Second degré};}]
\vspace*{1.75cm}
\begin{minipage}[t][6.9cm][c]{\linewidth}\LARGE\centering
\begin{minipage}{0.95\linewidth}\raggedright\linespread{1.25}\selectfont Quel est l'ensemble des solutions de l'inéquation $x^2-6x+5>0$~?\end{minipage}
\end{minipage}

\vfill
\begin{tcolorbox}[enhanced,nobeforeafter,width=7.1cm,height=2.1cm,valign=center,colback=cdeg!6,colframe=cdeg,arc=2.5mm,boxrule=0.9pt,left=3mm]{\Large\bfseries\color{cdeg}a)}\quad\begin{minipage}[c]{5.3cm}\raggedright\large $]-\infty\,;\,1[\cup]5\,;\,+\infty[$\end{minipage}\end{tcolorbox}\hfill\begin{tcolorbox}[enhanced,nobeforeafter,width=7.1cm,height=2.1cm,valign=center,colback=cdeg!6,colframe=cdeg,arc=2.5mm,boxrule=0.9pt,left=3mm]{\Large\bfseries\color{cdeg}b)}\quad\begin{minipage}[c]{5.3cm}\raggedright\Large $]1\,;\,5[$\end{minipage}\end{tcolorbox}\par\vspace{2.5mm}
\begin{tcolorbox}[enhanced,nobeforeafter,width=7.1cm,height=2.1cm,valign=center,colback=cdeg!6,colframe=cdeg,arc=2.5mm,boxrule=0.9pt,left=3mm]{\Large\bfseries\color{cdeg}c)}\quad\begin{minipage}[c]{5.3cm}\raggedright\large $]-\infty\,;\,1]\cup[5\,;\,+\infty[$\end{minipage}\end{tcolorbox}\hfill\begin{tcolorbox}[enhanced,nobeforeafter,width=7.1cm,height=2.1cm,valign=center,colback=cdeg!6,colframe=cdeg,arc=2.5mm,boxrule=0.9pt,left=3mm]{\Large\bfseries\color{cdeg}d)}\quad\begin{minipage}[c]{5.3cm}\raggedright\Large $[1\,;\,5]$\end{minipage}\end{tcolorbox}
\end{tcolorbox}
\newpage
\begin{tcolorbox}[enhanced,width=16cm,height=14cm,colback=white,colframe=cdeg,boxrule=2pt,arc=0mm,
left=6mm,right=6mm,top=0mm,bottom=4mm,sharp corners,
overlay={\fill[cdeg] (frame.north west) rectangle ([yshift=-1.25cm]frame.north east);
\node[anchor=west,text=white,font=\sffamily\bfseries\LARGE] at ([xshift=5mm,yshift=-0.62cm]frame.north west) {Question 29};
\node[anchor=east,text=white,font=\sffamily\itshape\large] at ([xshift=-5mm,yshift=-0.62cm]frame.north east) {Second degré};}]
\vspace*{1.75cm}
\begin{minipage}[t][6.9cm][c]{\linewidth}\LARGE\centering
\begin{minipage}{0.95\linewidth}\raggedright\linespread{1.25}\selectfont Le trinôme $x^2-7x+10$ a deux racines. Que vaut leur somme~?\end{minipage}
\end{minipage}

\vfill
\begin{tcolorbox}[enhanced,nobeforeafter,width=7.1cm,height=2.1cm,valign=center,colback=cdeg!6,colframe=cdeg,arc=2.5mm,boxrule=0.9pt,left=3mm]{\Large\bfseries\color{cdeg}a)}\quad\begin{minipage}[c]{5.3cm}\raggedright\Large $-7$\end{minipage}\end{tcolorbox}\hfill\begin{tcolorbox}[enhanced,nobeforeafter,width=7.1cm,height=2.1cm,valign=center,colback=cdeg!6,colframe=cdeg,arc=2.5mm,boxrule=0.9pt,left=3mm]{\Large\bfseries\color{cdeg}b)}\quad\begin{minipage}[c]{5.3cm}\raggedright\Large $7$\end{minipage}\end{tcolorbox}\par\vspace{2.5mm}
\begin{tcolorbox}[enhanced,nobeforeafter,width=7.1cm,height=2.1cm,valign=center,colback=cdeg!6,colframe=cdeg,arc=2.5mm,boxrule=0.9pt,left=3mm]{\Large\bfseries\color{cdeg}c)}\quad\begin{minipage}[c]{5.3cm}\raggedright\Large $14$\end{minipage}\end{tcolorbox}\hfill\begin{tcolorbox}[enhanced,nobeforeafter,width=7.1cm,height=2.1cm,valign=center,colback=cdeg!6,colframe=cdeg,arc=2.5mm,boxrule=0.9pt,left=3mm]{\Large\bfseries\color{cdeg}d)}\quad\begin{minipage}[c]{5.3cm}\raggedright\Large $10$\end{minipage}\end{tcolorbox}
\end{tcolorbox}
\newpage
\begin{tcolorbox}[enhanced,width=16cm,height=14cm,colback=white,colframe=cdeg,boxrule=2pt,arc=0mm,
left=6mm,right=6mm,top=0mm,bottom=4mm,sharp corners,
overlay={\fill[cdeg] (frame.north west) rectangle ([yshift=-1.25cm]frame.north east);
\node[anchor=west,text=white,font=\sffamily\bfseries\LARGE] at ([xshift=5mm,yshift=-0.62cm]frame.north west) {Question 30};
\node[anchor=east,text=white,font=\sffamily\itshape\large] at ([xshift=-5mm,yshift=-0.62cm]frame.north east) {Second degré};}]
\vspace*{1.75cm}
\begin{minipage}[t][6.9cm][c]{\linewidth}\LARGE\centering
\begin{minipage}{0.95\linewidth}\raggedright\linespread{1.25}\selectfont Le trinôme $2x^2-6x+3$ a deux racines. Que vaut leur produit~?\end{minipage}
\end{minipage}

\vfill
\begin{tcolorbox}[enhanced,nobeforeafter,width=7.1cm,height=2.1cm,valign=center,colback=cdeg!6,colframe=cdeg,arc=2.5mm,boxrule=0.9pt,left=3mm]{\Large\bfseries\color{cdeg}a)}\quad\begin{minipage}[c]{5.3cm}\raggedright\Large $-\dfrac{3}{2}$\end{minipage}\end{tcolorbox}\hfill\begin{tcolorbox}[enhanced,nobeforeafter,width=7.1cm,height=2.1cm,valign=center,colback=cdeg!6,colframe=cdeg,arc=2.5mm,boxrule=0.9pt,left=3mm]{\Large\bfseries\color{cdeg}b)}\quad\begin{minipage}[c]{5.3cm}\raggedright\Large $-3$\end{minipage}\end{tcolorbox}\par\vspace{2.5mm}
\begin{tcolorbox}[enhanced,nobeforeafter,width=7.1cm,height=2.1cm,valign=center,colback=cdeg!6,colframe=cdeg,arc=2.5mm,boxrule=0.9pt,left=3mm]{\Large\bfseries\color{cdeg}c)}\quad\begin{minipage}[c]{5.3cm}\raggedright\Large $\dfrac{3}{2}$\end{minipage}\end{tcolorbox}\hfill\begin{tcolorbox}[enhanced,nobeforeafter,width=7.1cm,height=2.1cm,valign=center,colback=cdeg!6,colframe=cdeg,arc=2.5mm,boxrule=0.9pt,left=3mm]{\Large\bfseries\color{cdeg}d)}\quad\begin{minipage}[c]{5.3cm}\raggedright\Large $3$\end{minipage}\end{tcolorbox}
\end{tcolorbox}
\newpage
\begin{tcolorbox}[enhanced,width=16cm,height=14cm,colback=white,colframe=cdeg,boxrule=2pt,arc=0mm,
left=6mm,right=6mm,top=0mm,bottom=4mm,sharp corners,
overlay={\fill[cdeg] (frame.north west) rectangle ([yshift=-1.25cm]frame.north east);
\node[anchor=west,text=white,font=\sffamily\bfseries\LARGE] at ([xshift=5mm,yshift=-0.62cm]frame.north west) {Question 31};
\node[anchor=east,text=white,font=\sffamily\itshape\large] at ([xshift=-5mm,yshift=-0.62cm]frame.north east) {Second degré};}]
\vspace*{1.75cm}
\begin{minipage}[t][6.9cm][c]{\linewidth}\LARGE\centering
\begin{minipage}{0.95\linewidth}\raggedright\linespread{1.25}\selectfont Le trinôme $2x^2-2x-12$ a pour racines $-2$ et $3$. Quelle est sa forme factorisée~?\end{minipage}
\end{minipage}

\vfill
\begin{tcolorbox}[enhanced,nobeforeafter,width=7.1cm,height=2.1cm,valign=center,colback=cdeg!6,colframe=cdeg,arc=2.5mm,boxrule=0.9pt,left=3mm]{\Large\bfseries\color{cdeg}a)}\quad\begin{minipage}[c]{5.3cm}\raggedright\Large $2(x+2)(x-3)$\end{minipage}\end{tcolorbox}\hfill\begin{tcolorbox}[enhanced,nobeforeafter,width=7.1cm,height=2.1cm,valign=center,colback=cdeg!6,colframe=cdeg,arc=2.5mm,boxrule=0.9pt,left=3mm]{\Large\bfseries\color{cdeg}b)}\quad\begin{minipage}[c]{5.3cm}\raggedright\Large $-(x+2)(x-3)$\end{minipage}\end{tcolorbox}\par\vspace{2.5mm}
\begin{tcolorbox}[enhanced,nobeforeafter,width=7.1cm,height=2.1cm,valign=center,colback=cdeg!6,colframe=cdeg,arc=2.5mm,boxrule=0.9pt,left=3mm]{\Large\bfseries\color{cdeg}c)}\quad\begin{minipage}[c]{5.3cm}\raggedright\Large $(x+2)(x-3)$\end{minipage}\end{tcolorbox}\hfill\begin{tcolorbox}[enhanced,nobeforeafter,width=7.1cm,height=2.1cm,valign=center,colback=cdeg!6,colframe=cdeg,arc=2.5mm,boxrule=0.9pt,left=3mm]{\Large\bfseries\color{cdeg}d)}\quad\begin{minipage}[c]{5.3cm}\raggedright\Large $2(x-2)(x+3)$\end{minipage}\end{tcolorbox}
\end{tcolorbox}
\newpage
\begin{tcolorbox}[enhanced,width=16cm,height=14cm,colback=white,colframe=cdeg,boxrule=2pt,arc=0mm,
left=6mm,right=6mm,top=0mm,bottom=4mm,sharp corners,
overlay={\fill[cdeg] (frame.north west) rectangle ([yshift=-1.25cm]frame.north east);
\node[anchor=west,text=white,font=\sffamily\bfseries\LARGE] at ([xshift=5mm,yshift=-0.62cm]frame.north west) {Question 32};
\node[anchor=east,text=white,font=\sffamily\itshape\large] at ([xshift=-5mm,yshift=-0.62cm]frame.north east) {Second degré};}]
\vspace*{1.75cm}
\begin{minipage}[t][6.9cm][c]{\linewidth}\LARGE\centering
\begin{minipage}{0.95\linewidth}\raggedright\linespread{1.25}\selectfont Le trinôme $x^2-5x+4$ a pour racines $1$ et $4$. Quelle est sa forme factorisée~?\end{minipage}
\end{minipage}

\vfill
\begin{tcolorbox}[enhanced,nobeforeafter,width=7.1cm,height=2.1cm,valign=center,colback=cdeg!6,colframe=cdeg,arc=2.5mm,boxrule=0.9pt,left=3mm]{\Large\bfseries\color{cdeg}a)}\quad\begin{minipage}[c]{5.3cm}\raggedright\Large $(x+1)(x+4)$\end{minipage}\end{tcolorbox}\hfill\begin{tcolorbox}[enhanced,nobeforeafter,width=7.1cm,height=2.1cm,valign=center,colback=cdeg!6,colframe=cdeg,arc=2.5mm,boxrule=0.9pt,left=3mm]{\Large\bfseries\color{cdeg}b)}\quad\begin{minipage}[c]{5.3cm}\raggedright\Large $2(x-1)(x-4)$\end{minipage}\end{tcolorbox}\par\vspace{2.5mm}
\begin{tcolorbox}[enhanced,nobeforeafter,width=7.1cm,height=2.1cm,valign=center,colback=cdeg!6,colframe=cdeg,arc=2.5mm,boxrule=0.9pt,left=3mm]{\Large\bfseries\color{cdeg}c)}\quad\begin{minipage}[c]{5.3cm}\raggedright\Large $-(x-1)(x-4)$\end{minipage}\end{tcolorbox}\hfill\begin{tcolorbox}[enhanced,nobeforeafter,width=7.1cm,height=2.1cm,valign=center,colback=cdeg!6,colframe=cdeg,arc=2.5mm,boxrule=0.9pt,left=3mm]{\Large\bfseries\color{cdeg}d)}\quad\begin{minipage}[c]{5.3cm}\raggedright\Large $(x-1)(x-4)$\end{minipage}\end{tcolorbox}
\end{tcolorbox}
\newpage
\begin{tcolorbox}[enhanced,width=16cm,height=14cm,colback=white,colframe=cder,boxrule=2pt,arc=0mm,
left=6mm,right=6mm,top=0mm,bottom=4mm,sharp corners,
overlay={\fill[cder] (frame.north west) rectangle ([yshift=-1.25cm]frame.north east);
\node[anchor=west,text=white,font=\sffamily\bfseries\LARGE] at ([xshift=5mm,yshift=-0.62cm]frame.north west) {Question 33};
\node[anchor=east,text=white,font=\sffamily\itshape\large] at ([xshift=-5mm,yshift=-0.62cm]frame.north east) {Dérivation};}]
\vspace*{1.75cm}
\begin{minipage}[t][6.9cm][c]{\linewidth}\LARGE\centering
\begin{minipage}{0.95\linewidth}\raggedright\linespread{1.25}\selectfont Soit $f(x)=x^3-2x^2+3x+5$. Que vaut $f'(2)$~?\end{minipage}
\end{minipage}

\vfill
\begin{tcolorbox}[enhanced,nobeforeafter,width=7.1cm,height=2.1cm,valign=center,colback=cder!6,colframe=cder,arc=2.5mm,boxrule=0.9pt,left=3mm]{\Large\bfseries\color{cder}a)}\quad\begin{minipage}[c]{5.3cm}\raggedright\Large $7$\end{minipage}\end{tcolorbox}\hfill\begin{tcolorbox}[enhanced,nobeforeafter,width=7.1cm,height=2.1cm,valign=center,colback=cder!6,colframe=cder,arc=2.5mm,boxrule=0.9pt,left=3mm]{\Large\bfseries\color{cder}b)}\quad\begin{minipage}[c]{5.3cm}\raggedright\Large $3$\end{minipage}\end{tcolorbox}\par\vspace{2.5mm}
\begin{tcolorbox}[enhanced,nobeforeafter,width=7.1cm,height=2.1cm,valign=center,colback=cder!6,colframe=cder,arc=2.5mm,boxrule=0.9pt,left=3mm]{\Large\bfseries\color{cder}c)}\quad\begin{minipage}[c]{5.3cm}\raggedright\Large $11$\end{minipage}\end{tcolorbox}\hfill\begin{tcolorbox}[enhanced,nobeforeafter,width=7.1cm,height=2.1cm,valign=center,colback=cder!6,colframe=cder,arc=2.5mm,boxrule=0.9pt,left=3mm]{\Large\bfseries\color{cder}d)}\quad\begin{minipage}[c]{5.3cm}\raggedright\Large $8$\end{minipage}\end{tcolorbox}
\end{tcolorbox}
\newpage
\begin{tcolorbox}[enhanced,width=16cm,height=14cm,colback=white,colframe=cder,boxrule=2pt,arc=0mm,
left=6mm,right=6mm,top=0mm,bottom=4mm,sharp corners,
overlay={\fill[cder] (frame.north west) rectangle ([yshift=-1.25cm]frame.north east);
\node[anchor=west,text=white,font=\sffamily\bfseries\LARGE] at ([xshift=5mm,yshift=-0.62cm]frame.north west) {Question 34};
\node[anchor=east,text=white,font=\sffamily\itshape\large] at ([xshift=-5mm,yshift=-0.62cm]frame.north east) {Dérivation};}]
\vspace*{1.75cm}
\begin{minipage}[t][6.9cm][c]{\linewidth}\LARGE\centering
\begin{minipage}{0.95\linewidth}\raggedright\linespread{1.25}\selectfont Soit $f(x)=2x^3-5x+1$. Que vaut $f'(1)$~?\end{minipage}
\end{minipage}

\vfill
\begin{tcolorbox}[enhanced,nobeforeafter,width=7.1cm,height=2.1cm,valign=center,colback=cder!6,colframe=cder,arc=2.5mm,boxrule=0.9pt,left=3mm]{\Large\bfseries\color{cder}a)}\quad\begin{minipage}[c]{5.3cm}\raggedright\Large $2$\end{minipage}\end{tcolorbox}\hfill\begin{tcolorbox}[enhanced,nobeforeafter,width=7.1cm,height=2.1cm,valign=center,colback=cder!6,colframe=cder,arc=2.5mm,boxrule=0.9pt,left=3mm]{\Large\bfseries\color{cder}b)}\quad\begin{minipage}[c]{5.3cm}\raggedright\Large $1$\end{minipage}\end{tcolorbox}\par\vspace{2.5mm}
\begin{tcolorbox}[enhanced,nobeforeafter,width=7.1cm,height=2.1cm,valign=center,colback=cder!6,colframe=cder,arc=2.5mm,boxrule=0.9pt,left=3mm]{\Large\bfseries\color{cder}c)}\quad\begin{minipage}[c]{5.3cm}\raggedright\Large $-3$\end{minipage}\end{tcolorbox}\hfill\begin{tcolorbox}[enhanced,nobeforeafter,width=7.1cm,height=2.1cm,valign=center,colback=cder!6,colframe=cder,arc=2.5mm,boxrule=0.9pt,left=3mm]{\Large\bfseries\color{cder}d)}\quad\begin{minipage}[c]{5.3cm}\raggedright\Large $-2$\end{minipage}\end{tcolorbox}
\end{tcolorbox}
\newpage
\begin{tcolorbox}[enhanced,width=16cm,height=14cm,colback=white,colframe=cder,boxrule=2pt,arc=0mm,
left=6mm,right=6mm,top=0mm,bottom=4mm,sharp corners,
overlay={\fill[cder] (frame.north west) rectangle ([yshift=-1.25cm]frame.north east);
\node[anchor=west,text=white,font=\sffamily\bfseries\LARGE] at ([xshift=5mm,yshift=-0.62cm]frame.north west) {Question 35};
\node[anchor=east,text=white,font=\sffamily\itshape\large] at ([xshift=-5mm,yshift=-0.62cm]frame.north east) {Dérivation};}]
\vspace*{1.75cm}
\begin{minipage}[t][6.9cm][c]{\linewidth}\LARGE\centering
\begin{minipage}{0.95\linewidth}\raggedright\linespread{1.25}\selectfont Quelle est l'équation de la tangente à la courbe de $f(x)=x^2-2x+3$ au point d'abscisse $1$~?\end{minipage}
\end{minipage}

\vfill
\begin{tcolorbox}[enhanced,nobeforeafter,width=7.1cm,height=2.1cm,valign=center,colback=cder!6,colframe=cder,arc=2.5mm,boxrule=0.9pt,left=3mm]{\Large\bfseries\color{cder}a)}\quad\begin{minipage}[c]{5.3cm}\raggedright\Large $y=3$\end{minipage}\end{tcolorbox}\hfill\begin{tcolorbox}[enhanced,nobeforeafter,width=7.1cm,height=2.1cm,valign=center,colback=cder!6,colframe=cder,arc=2.5mm,boxrule=0.9pt,left=3mm]{\Large\bfseries\color{cder}b)}\quad\begin{minipage}[c]{5.3cm}\raggedright\Large $y=2x$\end{minipage}\end{tcolorbox}\par\vspace{2.5mm}
\begin{tcolorbox}[enhanced,nobeforeafter,width=7.1cm,height=2.1cm,valign=center,colback=cder!6,colframe=cder,arc=2.5mm,boxrule=0.9pt,left=3mm]{\Large\bfseries\color{cder}c)}\quad\begin{minipage}[c]{5.3cm}\raggedright\Large $y=2x-2$\end{minipage}\end{tcolorbox}\hfill\begin{tcolorbox}[enhanced,nobeforeafter,width=7.1cm,height=2.1cm,valign=center,colback=cder!6,colframe=cder,arc=2.5mm,boxrule=0.9pt,left=3mm]{\Large\bfseries\color{cder}d)}\quad\begin{minipage}[c]{5.3cm}\raggedright\Large $y=2$\end{minipage}\end{tcolorbox}
\end{tcolorbox}
\newpage
\begin{tcolorbox}[enhanced,width=16cm,height=14cm,colback=white,colframe=cder,boxrule=2pt,arc=0mm,
left=6mm,right=6mm,top=0mm,bottom=4mm,sharp corners,
overlay={\fill[cder] (frame.north west) rectangle ([yshift=-1.25cm]frame.north east);
\node[anchor=west,text=white,font=\sffamily\bfseries\LARGE] at ([xshift=5mm,yshift=-0.62cm]frame.north west) {Question 36};
\node[anchor=east,text=white,font=\sffamily\itshape\large] at ([xshift=-5mm,yshift=-0.62cm]frame.north east) {Dérivation};}]
\vspace*{1.75cm}
\begin{minipage}[t][6.9cm][c]{\linewidth}\LARGE\centering
\begin{minipage}{0.95\linewidth}\raggedright\linespread{1.25}\selectfont Quelle est l'équation de la tangente à la courbe de $f(x)=2x^2+x$ au point d'abscisse $-1$~?\end{minipage}
\end{minipage}

\vfill
\begin{tcolorbox}[enhanced,nobeforeafter,width=7.1cm,height=2.1cm,valign=center,colback=cder!6,colframe=cder,arc=2.5mm,boxrule=0.9pt,left=3mm]{\Large\bfseries\color{cder}a)}\quad\begin{minipage}[c]{5.3cm}\raggedright\Large $y=-3x+4$\end{minipage}\end{tcolorbox}\hfill\begin{tcolorbox}[enhanced,nobeforeafter,width=7.1cm,height=2.1cm,valign=center,colback=cder!6,colframe=cder,arc=2.5mm,boxrule=0.9pt,left=3mm]{\Large\bfseries\color{cder}b)}\quad\begin{minipage}[c]{5.3cm}\raggedright\Large $y=-3x+1$\end{minipage}\end{tcolorbox}\par\vspace{2.5mm}
\begin{tcolorbox}[enhanced,nobeforeafter,width=7.1cm,height=2.1cm,valign=center,colback=cder!6,colframe=cder,arc=2.5mm,boxrule=0.9pt,left=3mm]{\Large\bfseries\color{cder}c)}\quad\begin{minipage}[c]{5.3cm}\raggedright\Large $y=x-2$\end{minipage}\end{tcolorbox}\hfill\begin{tcolorbox}[enhanced,nobeforeafter,width=7.1cm,height=2.1cm,valign=center,colback=cder!6,colframe=cder,arc=2.5mm,boxrule=0.9pt,left=3mm]{\Large\bfseries\color{cder}d)}\quad\begin{minipage}[c]{5.3cm}\raggedright\Large $y=-3x-2$\end{minipage}\end{tcolorbox}
\end{tcolorbox}
\newpage
\begin{tcolorbox}[enhanced,width=16cm,height=14cm,colback=white,colframe=cder,boxrule=2pt,arc=0mm,
left=6mm,right=6mm,top=0mm,bottom=4mm,sharp corners,
overlay={\fill[cder] (frame.north west) rectangle ([yshift=-1.25cm]frame.north east);
\node[anchor=west,text=white,font=\sffamily\bfseries\LARGE] at ([xshift=5mm,yshift=-0.62cm]frame.north west) {Question 37};
\node[anchor=east,text=white,font=\sffamily\itshape\large] at ([xshift=-5mm,yshift=-0.62cm]frame.north east) {Dérivation};}]
\vspace*{1.75cm}
\begin{minipage}[t][6.9cm][c]{\linewidth}\LARGE\centering
\begin{minipage}{0.95\linewidth}\raggedright\linespread{1.25}\selectfont Soit $f(x)=(2x+1)(3x-1)$. Que vaut $f'(1)$~?\end{minipage}
\end{minipage}

\vfill
\begin{tcolorbox}[enhanced,nobeforeafter,width=7.1cm,height=2.1cm,valign=center,colback=cder!6,colframe=cder,arc=2.5mm,boxrule=0.9pt,left=3mm]{\Large\bfseries\color{cder}a)}\quad\begin{minipage}[c]{5.3cm}\raggedright\Large $4$\end{minipage}\end{tcolorbox}\hfill\begin{tcolorbox}[enhanced,nobeforeafter,width=7.1cm,height=2.1cm,valign=center,colback=cder!6,colframe=cder,arc=2.5mm,boxrule=0.9pt,left=3mm]{\Large\bfseries\color{cder}b)}\quad\begin{minipage}[c]{5.3cm}\raggedright\Large $15$\end{minipage}\end{tcolorbox}\par\vspace{2.5mm}
\begin{tcolorbox}[enhanced,nobeforeafter,width=7.1cm,height=2.1cm,valign=center,colback=cder!6,colframe=cder,arc=2.5mm,boxrule=0.9pt,left=3mm]{\Large\bfseries\color{cder}c)}\quad\begin{minipage}[c]{5.3cm}\raggedright\Large $6$\end{minipage}\end{tcolorbox}\hfill\begin{tcolorbox}[enhanced,nobeforeafter,width=7.1cm,height=2.1cm,valign=center,colback=cder!6,colframe=cder,arc=2.5mm,boxrule=0.9pt,left=3mm]{\Large\bfseries\color{cder}d)}\quad\begin{minipage}[c]{5.3cm}\raggedright\Large $13$\end{minipage}\end{tcolorbox}
\end{tcolorbox}
\newpage
\begin{tcolorbox}[enhanced,width=16cm,height=14cm,colback=white,colframe=cder,boxrule=2pt,arc=0mm,
left=6mm,right=6mm,top=0mm,bottom=4mm,sharp corners,
overlay={\fill[cder] (frame.north west) rectangle ([yshift=-1.25cm]frame.north east);
\node[anchor=west,text=white,font=\sffamily\bfseries\LARGE] at ([xshift=5mm,yshift=-0.62cm]frame.north west) {Question 38};
\node[anchor=east,text=white,font=\sffamily\itshape\large] at ([xshift=-5mm,yshift=-0.62cm]frame.north east) {Dérivation};}]
\vspace*{1.75cm}
\begin{minipage}[t][6.9cm][c]{\linewidth}\LARGE\centering
\begin{minipage}{0.95\linewidth}\raggedright\linespread{1.25}\selectfont Soit $f(x)=(x-3)(2x+5)$. Que vaut $f'(0)$~?\end{minipage}
\end{minipage}

\vfill
\begin{tcolorbox}[enhanced,nobeforeafter,width=7.1cm,height=2.1cm,valign=center,colback=cder!6,colframe=cder,arc=2.5mm,boxrule=0.9pt,left=3mm]{\Large\bfseries\color{cder}a)}\quad\begin{minipage}[c]{5.3cm}\raggedright\Large $2$\end{minipage}\end{tcolorbox}\hfill\begin{tcolorbox}[enhanced,nobeforeafter,width=7.1cm,height=2.1cm,valign=center,colback=cder!6,colframe=cder,arc=2.5mm,boxrule=0.9pt,left=3mm]{\Large\bfseries\color{cder}b)}\quad\begin{minipage}[c]{5.3cm}\raggedright\Large $-1$\end{minipage}\end{tcolorbox}\par\vspace{2.5mm}
\begin{tcolorbox}[enhanced,nobeforeafter,width=7.1cm,height=2.1cm,valign=center,colback=cder!6,colframe=cder,arc=2.5mm,boxrule=0.9pt,left=3mm]{\Large\bfseries\color{cder}c)}\quad\begin{minipage}[c]{5.3cm}\raggedright\Large $5$\end{minipage}\end{tcolorbox}\hfill\begin{tcolorbox}[enhanced,nobeforeafter,width=7.1cm,height=2.1cm,valign=center,colback=cder!6,colframe=cder,arc=2.5mm,boxrule=0.9pt,left=3mm]{\Large\bfseries\color{cder}d)}\quad\begin{minipage}[c]{5.3cm}\raggedright\Large $-15$\end{minipage}\end{tcolorbox}
\end{tcolorbox}
\newpage
\begin{tcolorbox}[enhanced,width=16cm,height=14cm,colback=white,colframe=cder,boxrule=2pt,arc=0mm,
left=6mm,right=6mm,top=0mm,bottom=4mm,sharp corners,
overlay={\fill[cder] (frame.north west) rectangle ([yshift=-1.25cm]frame.north east);
\node[anchor=west,text=white,font=\sffamily\bfseries\LARGE] at ([xshift=5mm,yshift=-0.62cm]frame.north west) {Question 39};
\node[anchor=east,text=white,font=\sffamily\itshape\large] at ([xshift=-5mm,yshift=-0.62cm]frame.north east) {Dérivation};}]
\vspace*{1.75cm}
\begin{minipage}[t][6.9cm][c]{\linewidth}\LARGE\centering
\begin{minipage}{0.95\linewidth}\raggedright\linespread{1.25}\selectfont Soit $f(x)=\dfrac{3}{x}$ sur $\mathbb{R}^*$. Que vaut $f'(1)$~?\end{minipage}
\end{minipage}

\vfill
\begin{tcolorbox}[enhanced,nobeforeafter,width=7.1cm,height=2.1cm,valign=center,colback=cder!6,colframe=cder,arc=2.5mm,boxrule=0.9pt,left=3mm]{\Large\bfseries\color{cder}a)}\quad\begin{minipage}[c]{5.3cm}\raggedright\Large $-2$\end{minipage}\end{tcolorbox}\hfill\begin{tcolorbox}[enhanced,nobeforeafter,width=7.1cm,height=2.1cm,valign=center,colback=cder!6,colframe=cder,arc=2.5mm,boxrule=0.9pt,left=3mm]{\Large\bfseries\color{cder}b)}\quad\begin{minipage}[c]{5.3cm}\raggedright\Large $3$\end{minipage}\end{tcolorbox}\par\vspace{2.5mm}
\begin{tcolorbox}[enhanced,nobeforeafter,width=7.1cm,height=2.1cm,valign=center,colback=cder!6,colframe=cder,arc=2.5mm,boxrule=0.9pt,left=3mm]{\Large\bfseries\color{cder}c)}\quad\begin{minipage}[c]{5.3cm}\raggedright\Large $-\dfrac{3}{2}$\end{minipage}\end{tcolorbox}\hfill\begin{tcolorbox}[enhanced,nobeforeafter,width=7.1cm,height=2.1cm,valign=center,colback=cder!6,colframe=cder,arc=2.5mm,boxrule=0.9pt,left=3mm]{\Large\bfseries\color{cder}d)}\quad\begin{minipage}[c]{5.3cm}\raggedright\Large $-3$\end{minipage}\end{tcolorbox}
\end{tcolorbox}
\newpage
\begin{tcolorbox}[enhanced,width=16cm,height=14cm,colback=white,colframe=cder,boxrule=2pt,arc=0mm,
left=6mm,right=6mm,top=0mm,bottom=4mm,sharp corners,
overlay={\fill[cder] (frame.north west) rectangle ([yshift=-1.25cm]frame.north east);
\node[anchor=west,text=white,font=\sffamily\bfseries\LARGE] at ([xshift=5mm,yshift=-0.62cm]frame.north west) {Question 40};
\node[anchor=east,text=white,font=\sffamily\itshape\large] at ([xshift=-5mm,yshift=-0.62cm]frame.north east) {Dérivation};}]
\vspace*{1.75cm}
\begin{minipage}[t][6.9cm][c]{\linewidth}\LARGE\centering
\begin{minipage}{0.95\linewidth}\raggedright\linespread{1.25}\selectfont Soit $f(x)=\dfrac{2}{x}$ sur $\mathbb{R}^*$. Que vaut $f'(-1)$~?\end{minipage}
\end{minipage}

\vfill
\begin{tcolorbox}[enhanced,nobeforeafter,width=7.1cm,height=2.1cm,valign=center,colback=cder!6,colframe=cder,arc=2.5mm,boxrule=0.9pt,left=3mm]{\Large\bfseries\color{cder}a)}\quad\begin{minipage}[c]{5.3cm}\raggedright\Large $-2$\end{minipage}\end{tcolorbox}\hfill\begin{tcolorbox}[enhanced,nobeforeafter,width=7.1cm,height=2.1cm,valign=center,colback=cder!6,colframe=cder,arc=2.5mm,boxrule=0.9pt,left=3mm]{\Large\bfseries\color{cder}b)}\quad\begin{minipage}[c]{5.3cm}\raggedright\Large $2$\end{minipage}\end{tcolorbox}\par\vspace{2.5mm}
\begin{tcolorbox}[enhanced,nobeforeafter,width=7.1cm,height=2.1cm,valign=center,colback=cder!6,colframe=cder,arc=2.5mm,boxrule=0.9pt,left=3mm]{\Large\bfseries\color{cder}c)}\quad\begin{minipage}[c]{5.3cm}\raggedright\Large $1$\end{minipage}\end{tcolorbox}\hfill\begin{tcolorbox}[enhanced,nobeforeafter,width=7.1cm,height=2.1cm,valign=center,colback=cder!6,colframe=cder,arc=2.5mm,boxrule=0.9pt,left=3mm]{\Large\bfseries\color{cder}d)}\quad\begin{minipage}[c]{5.3cm}\raggedright\Large $-1$\end{minipage}\end{tcolorbox}
\end{tcolorbox}
\newpage
\begin{tcolorbox}[enhanced,width=16cm,height=14cm,colback=white,colframe=cder,boxrule=2pt,arc=0mm,
left=6mm,right=6mm,top=0mm,bottom=4mm,sharp corners,
overlay={\fill[cder] (frame.north west) rectangle ([yshift=-1.25cm]frame.north east);
\node[anchor=west,text=white,font=\sffamily\bfseries\LARGE] at ([xshift=5mm,yshift=-0.62cm]frame.north west) {Question 41};
\node[anchor=east,text=white,font=\sffamily\itshape\large] at ([xshift=-5mm,yshift=-0.62cm]frame.north east) {Dérivation};}]
\vspace*{1.75cm}
\begin{minipage}[t][6.9cm][c]{\linewidth}\LARGE\centering
\begin{minipage}{0.95\linewidth}\raggedright\linespread{1.25}\selectfont Soit $f(x)=\sqrt{x}$. Que vaut $f'(4)$~?\end{minipage}
\end{minipage}

\vfill
\begin{tcolorbox}[enhanced,nobeforeafter,width=7.1cm,height=2.1cm,valign=center,colback=cder!6,colframe=cder,arc=2.5mm,boxrule=0.9pt,left=3mm]{\Large\bfseries\color{cder}a)}\quad\begin{minipage}[c]{5.3cm}\raggedright\Large $1$\end{minipage}\end{tcolorbox}\hfill\begin{tcolorbox}[enhanced,nobeforeafter,width=7.1cm,height=2.1cm,valign=center,colback=cder!6,colframe=cder,arc=2.5mm,boxrule=0.9pt,left=3mm]{\Large\bfseries\color{cder}b)}\quad\begin{minipage}[c]{5.3cm}\raggedright\Large $\dfrac{1}{4}$\end{minipage}\end{tcolorbox}\par\vspace{2.5mm}
\begin{tcolorbox}[enhanced,nobeforeafter,width=7.1cm,height=2.1cm,valign=center,colback=cder!6,colframe=cder,arc=2.5mm,boxrule=0.9pt,left=3mm]{\Large\bfseries\color{cder}c)}\quad\begin{minipage}[c]{5.3cm}\raggedright\Large $\dfrac{1}{2}$\end{minipage}\end{tcolorbox}\hfill\begin{tcolorbox}[enhanced,nobeforeafter,width=7.1cm,height=2.1cm,valign=center,colback=cder!6,colframe=cder,arc=2.5mm,boxrule=0.9pt,left=3mm]{\Large\bfseries\color{cder}d)}\quad\begin{minipage}[c]{5.3cm}\raggedright\Large $4$\end{minipage}\end{tcolorbox}
\end{tcolorbox}
\newpage
\begin{tcolorbox}[enhanced,width=16cm,height=14cm,colback=white,colframe=cder,boxrule=2pt,arc=0mm,
left=6mm,right=6mm,top=0mm,bottom=4mm,sharp corners,
overlay={\fill[cder] (frame.north west) rectangle ([yshift=-1.25cm]frame.north east);
\node[anchor=west,text=white,font=\sffamily\bfseries\LARGE] at ([xshift=5mm,yshift=-0.62cm]frame.north west) {Question 42};
\node[anchor=east,text=white,font=\sffamily\itshape\large] at ([xshift=-5mm,yshift=-0.62cm]frame.north east) {Dérivation};}]
\vspace*{1.75cm}
\begin{minipage}[t][6.9cm][c]{\linewidth}\LARGE\centering
\begin{minipage}{0.95\linewidth}\raggedright\linespread{1.25}\selectfont Soit $f(x)=\sqrt{x}$. Que vaut $f'(9)$~?\end{minipage}
\end{minipage}

\vfill
\begin{tcolorbox}[enhanced,nobeforeafter,width=7.1cm,height=2.1cm,valign=center,colback=cder!6,colframe=cder,arc=2.5mm,boxrule=0.9pt,left=3mm]{\Large\bfseries\color{cder}a)}\quad\begin{minipage}[c]{5.3cm}\raggedright\Large $\dfrac{1}{3}$\end{minipage}\end{tcolorbox}\hfill\begin{tcolorbox}[enhanced,nobeforeafter,width=7.1cm,height=2.1cm,valign=center,colback=cder!6,colframe=cder,arc=2.5mm,boxrule=0.9pt,left=3mm]{\Large\bfseries\color{cder}b)}\quad\begin{minipage}[c]{5.3cm}\raggedright\Large $\dfrac{1}{6}$\end{minipage}\end{tcolorbox}\par\vspace{2.5mm}
\begin{tcolorbox}[enhanced,nobeforeafter,width=7.1cm,height=2.1cm,valign=center,colback=cder!6,colframe=cder,arc=2.5mm,boxrule=0.9pt,left=3mm]{\Large\bfseries\color{cder}c)}\quad\begin{minipage}[c]{5.3cm}\raggedright\Large $\dfrac{3}{2}$\end{minipage}\end{tcolorbox}\hfill\begin{tcolorbox}[enhanced,nobeforeafter,width=7.1cm,height=2.1cm,valign=center,colback=cder!6,colframe=cder,arc=2.5mm,boxrule=0.9pt,left=3mm]{\Large\bfseries\color{cder}d)}\quad\begin{minipage}[c]{5.3cm}\raggedright\Large $6$\end{minipage}\end{tcolorbox}
\end{tcolorbox}
\newpage
\begin{tcolorbox}[enhanced,width=16cm,height=14cm,colback=white,colframe=cder,boxrule=2pt,arc=0mm,
left=6mm,right=6mm,top=0mm,bottom=4mm,sharp corners,
overlay={\fill[cder] (frame.north west) rectangle ([yshift=-1.25cm]frame.north east);
\node[anchor=west,text=white,font=\sffamily\bfseries\LARGE] at ([xshift=5mm,yshift=-0.62cm]frame.north west) {Question 43};
\node[anchor=east,text=white,font=\sffamily\itshape\large] at ([xshift=-5mm,yshift=-0.62cm]frame.north east) {Dérivation};}]
\vspace*{1.75cm}
\begin{minipage}[t][6.9cm][c]{\linewidth}\LARGE\centering
\begin{minipage}{0.95\linewidth}\raggedright\linespread{1.25}\selectfont Soit $f(x)=(2x+3)e^x$. Que vaut $f'(0)$~?\end{minipage}
\end{minipage}

\vfill
\begin{tcolorbox}[enhanced,nobeforeafter,width=7.1cm,height=2.1cm,valign=center,colback=cder!6,colframe=cder,arc=2.5mm,boxrule=0.9pt,left=3mm]{\Large\bfseries\color{cder}a)}\quad\begin{minipage}[c]{5.3cm}\raggedright\Large $2$\end{minipage}\end{tcolorbox}\hfill\begin{tcolorbox}[enhanced,nobeforeafter,width=7.1cm,height=2.1cm,valign=center,colback=cder!6,colframe=cder,arc=2.5mm,boxrule=0.9pt,left=3mm]{\Large\bfseries\color{cder}b)}\quad\begin{minipage}[c]{5.3cm}\raggedright\Large $3$\end{minipage}\end{tcolorbox}\par\vspace{2.5mm}
\begin{tcolorbox}[enhanced,nobeforeafter,width=7.1cm,height=2.1cm,valign=center,colback=cder!6,colframe=cder,arc=2.5mm,boxrule=0.9pt,left=3mm]{\Large\bfseries\color{cder}c)}\quad\begin{minipage}[c]{5.3cm}\raggedright\Large $5$\end{minipage}\end{tcolorbox}\hfill\begin{tcolorbox}[enhanced,nobeforeafter,width=7.1cm,height=2.1cm,valign=center,colback=cder!6,colframe=cder,arc=2.5mm,boxrule=0.9pt,left=3mm]{\Large\bfseries\color{cder}d)}\quad\begin{minipage}[c]{5.3cm}\raggedright\Large $6$\end{minipage}\end{tcolorbox}
\end{tcolorbox}
\newpage
\begin{tcolorbox}[enhanced,width=16cm,height=14cm,colback=white,colframe=cder,boxrule=2pt,arc=0mm,
left=6mm,right=6mm,top=0mm,bottom=4mm,sharp corners,
overlay={\fill[cder] (frame.north west) rectangle ([yshift=-1.25cm]frame.north east);
\node[anchor=west,text=white,font=\sffamily\bfseries\LARGE] at ([xshift=5mm,yshift=-0.62cm]frame.north west) {Question 44};
\node[anchor=east,text=white,font=\sffamily\itshape\large] at ([xshift=-5mm,yshift=-0.62cm]frame.north east) {Dérivation};}]
\vspace*{1.75cm}
\begin{minipage}[t][6.9cm][c]{\linewidth}\LARGE\centering
\begin{minipage}{0.95\linewidth}\raggedright\linespread{1.25}\selectfont Soit $f(x)=(x-4)e^x$. Que vaut $f'(0)$~?\end{minipage}
\end{minipage}

\vfill
\begin{tcolorbox}[enhanced,nobeforeafter,width=7.1cm,height=2.1cm,valign=center,colback=cder!6,colframe=cder,arc=2.5mm,boxrule=0.9pt,left=3mm]{\Large\bfseries\color{cder}a)}\quad\begin{minipage}[c]{5.3cm}\raggedright\Large $-2$\end{minipage}\end{tcolorbox}\hfill\begin{tcolorbox}[enhanced,nobeforeafter,width=7.1cm,height=2.1cm,valign=center,colback=cder!6,colframe=cder,arc=2.5mm,boxrule=0.9pt,left=3mm]{\Large\bfseries\color{cder}b)}\quad\begin{minipage}[c]{5.3cm}\raggedright\Large $-4$\end{minipage}\end{tcolorbox}\par\vspace{2.5mm}
\begin{tcolorbox}[enhanced,nobeforeafter,width=7.1cm,height=2.1cm,valign=center,colback=cder!6,colframe=cder,arc=2.5mm,boxrule=0.9pt,left=3mm]{\Large\bfseries\color{cder}c)}\quad\begin{minipage}[c]{5.3cm}\raggedright\Large $1$\end{minipage}\end{tcolorbox}\hfill\begin{tcolorbox}[enhanced,nobeforeafter,width=7.1cm,height=2.1cm,valign=center,colback=cder!6,colframe=cder,arc=2.5mm,boxrule=0.9pt,left=3mm]{\Large\bfseries\color{cder}d)}\quad\begin{minipage}[c]{5.3cm}\raggedright\Large $-3$\end{minipage}\end{tcolorbox}
\end{tcolorbox}
\newpage
\begin{tcolorbox}[enhanced,width=16cm,height=14cm,colback=white,colframe=cder,boxrule=2pt,arc=0mm,
left=6mm,right=6mm,top=0mm,bottom=4mm,sharp corners,
overlay={\fill[cder] (frame.north west) rectangle ([yshift=-1.25cm]frame.north east);
\node[anchor=west,text=white,font=\sffamily\bfseries\LARGE] at ([xshift=5mm,yshift=-0.62cm]frame.north west) {Question 45};
\node[anchor=east,text=white,font=\sffamily\itshape\large] at ([xshift=-5mm,yshift=-0.62cm]frame.north east) {Dérivation};}]
\vspace*{1.75cm}
\begin{minipage}[t][6.9cm][c]{\linewidth}\LARGE\centering
\begin{minipage}{0.95\linewidth}\raggedright\linespread{1.25}\selectfont La fonction $f$ est dérivable sur $\mathbb{R}$ et $f'(x)=2x-6$. Sur quel intervalle $f$ est-elle décroissante~?\end{minipage}
\end{minipage}

\vfill
\begin{tcolorbox}[enhanced,nobeforeafter,width=7.1cm,height=2.1cm,valign=center,colback=cder!6,colframe=cder,arc=2.5mm,boxrule=0.9pt,left=3mm]{\Large\bfseries\color{cder}a)}\quad\begin{minipage}[c]{5.3cm}\raggedright\Large $[3\,;\,+\infty[$\end{minipage}\end{tcolorbox}\hfill\begin{tcolorbox}[enhanced,nobeforeafter,width=7.1cm,height=2.1cm,valign=center,colback=cder!6,colframe=cder,arc=2.5mm,boxrule=0.9pt,left=3mm]{\Large\bfseries\color{cder}b)}\quad\begin{minipage}[c]{5.3cm}\raggedright\Large $]-\infty\,;\,-3]$\end{minipage}\end{tcolorbox}\par\vspace{2.5mm}
\begin{tcolorbox}[enhanced,nobeforeafter,width=7.1cm,height=2.1cm,valign=center,colback=cder!6,colframe=cder,arc=2.5mm,boxrule=0.9pt,left=3mm]{\Large\bfseries\color{cder}c)}\quad\begin{minipage}[c]{5.3cm}\raggedright\Large $[-3\,;\,+\infty[$\end{minipage}\end{tcolorbox}\hfill\begin{tcolorbox}[enhanced,nobeforeafter,width=7.1cm,height=2.1cm,valign=center,colback=cder!6,colframe=cder,arc=2.5mm,boxrule=0.9pt,left=3mm]{\Large\bfseries\color{cder}d)}\quad\begin{minipage}[c]{5.3cm}\raggedright\Large $]-\infty\,;\,3]$\end{minipage}\end{tcolorbox}
\end{tcolorbox}
\newpage
\begin{tcolorbox}[enhanced,width=16cm,height=14cm,colback=white,colframe=cder,boxrule=2pt,arc=0mm,
left=6mm,right=6mm,top=0mm,bottom=4mm,sharp corners,
overlay={\fill[cder] (frame.north west) rectangle ([yshift=-1.25cm]frame.north east);
\node[anchor=west,text=white,font=\sffamily\bfseries\LARGE] at ([xshift=5mm,yshift=-0.62cm]frame.north west) {Question 46};
\node[anchor=east,text=white,font=\sffamily\itshape\large] at ([xshift=-5mm,yshift=-0.62cm]frame.north east) {Dérivation};}]
\vspace*{1.75cm}
\begin{minipage}[t][6.9cm][c]{\linewidth}\LARGE\centering
\begin{minipage}{0.95\linewidth}\raggedright\linespread{1.25}\selectfont La fonction $f$ est dérivable sur $\mathbb{R}$ et $f'(x)=-3x+6$. Sur quel intervalle $f$ est-elle décroissante~?\end{minipage}
\end{minipage}

\vfill
\begin{tcolorbox}[enhanced,nobeforeafter,width=7.1cm,height=2.1cm,valign=center,colback=cder!6,colframe=cder,arc=2.5mm,boxrule=0.9pt,left=3mm]{\Large\bfseries\color{cder}a)}\quad\begin{minipage}[c]{5.3cm}\raggedright\Large $]-\infty\,;\,-2]$\end{minipage}\end{tcolorbox}\hfill\begin{tcolorbox}[enhanced,nobeforeafter,width=7.1cm,height=2.1cm,valign=center,colback=cder!6,colframe=cder,arc=2.5mm,boxrule=0.9pt,left=3mm]{\Large\bfseries\color{cder}b)}\quad\begin{minipage}[c]{5.3cm}\raggedright\Large $]-\infty\,;\,2]$\end{minipage}\end{tcolorbox}\par\vspace{2.5mm}
\begin{tcolorbox}[enhanced,nobeforeafter,width=7.1cm,height=2.1cm,valign=center,colback=cder!6,colframe=cder,arc=2.5mm,boxrule=0.9pt,left=3mm]{\Large\bfseries\color{cder}c)}\quad\begin{minipage}[c]{5.3cm}\raggedright\Large $[2\,;\,+\infty[$\end{minipage}\end{tcolorbox}\hfill\begin{tcolorbox}[enhanced,nobeforeafter,width=7.1cm,height=2.1cm,valign=center,colback=cder!6,colframe=cder,arc=2.5mm,boxrule=0.9pt,left=3mm]{\Large\bfseries\color{cder}d)}\quad\begin{minipage}[c]{5.3cm}\raggedright\Large $[-2\,;\,+\infty[$\end{minipage}\end{tcolorbox}
\end{tcolorbox}
\newpage
\begin{tcolorbox}[enhanced,width=16cm,height=14cm,colback=white,colframe=cder,boxrule=2pt,arc=0mm,
left=6mm,right=6mm,top=0mm,bottom=4mm,sharp corners,
overlay={\fill[cder] (frame.north west) rectangle ([yshift=-1.25cm]frame.north east);
\node[anchor=west,text=white,font=\sffamily\bfseries\LARGE] at ([xshift=5mm,yshift=-0.62cm]frame.north west) {Question 47};
\node[anchor=east,text=white,font=\sffamily\itshape\large] at ([xshift=-5mm,yshift=-0.62cm]frame.north east) {Dérivation};}]
\vspace*{1.75cm}
\begin{minipage}[t][6.9cm][c]{\linewidth}\LARGE\centering
\begin{minipage}{0.95\linewidth}\raggedright\linespread{1.25}\selectfont Soit $f(x)=x^2-4x+1$. Quel est le minimum de $f$ sur $\mathbb{R}$~?\end{minipage}
\end{minipage}

\vfill
\begin{tcolorbox}[enhanced,nobeforeafter,width=7.1cm,height=2.1cm,valign=center,colback=cder!6,colframe=cder,arc=2.5mm,boxrule=0.9pt,left=3mm]{\Large\bfseries\color{cder}a)}\quad\begin{minipage}[c]{5.3cm}\raggedright\Large $2$\end{minipage}\end{tcolorbox}\hfill\begin{tcolorbox}[enhanced,nobeforeafter,width=7.1cm,height=2.1cm,valign=center,colback=cder!6,colframe=cder,arc=2.5mm,boxrule=0.9pt,left=3mm]{\Large\bfseries\color{cder}b)}\quad\begin{minipage}[c]{5.3cm}\raggedright\Large $1$\end{minipage}\end{tcolorbox}\par\vspace{2.5mm}
\begin{tcolorbox}[enhanced,nobeforeafter,width=7.1cm,height=2.1cm,valign=center,colback=cder!6,colframe=cder,arc=2.5mm,boxrule=0.9pt,left=3mm]{\Large\bfseries\color{cder}c)}\quad\begin{minipage}[c]{5.3cm}\raggedright\Large $3$\end{minipage}\end{tcolorbox}\hfill\begin{tcolorbox}[enhanced,nobeforeafter,width=7.1cm,height=2.1cm,valign=center,colback=cder!6,colframe=cder,arc=2.5mm,boxrule=0.9pt,left=3mm]{\Large\bfseries\color{cder}d)}\quad\begin{minipage}[c]{5.3cm}\raggedright\Large $-3$\end{minipage}\end{tcolorbox}
\end{tcolorbox}
\newpage
\begin{tcolorbox}[enhanced,width=16cm,height=14cm,colback=white,colframe=cder,boxrule=2pt,arc=0mm,
left=6mm,right=6mm,top=0mm,bottom=4mm,sharp corners,
overlay={\fill[cder] (frame.north west) rectangle ([yshift=-1.25cm]frame.north east);
\node[anchor=west,text=white,font=\sffamily\bfseries\LARGE] at ([xshift=5mm,yshift=-0.62cm]frame.north west) {Question 48};
\node[anchor=east,text=white,font=\sffamily\itshape\large] at ([xshift=-5mm,yshift=-0.62cm]frame.north east) {Dérivation};}]
\vspace*{1.75cm}
\begin{minipage}[t][6.9cm][c]{\linewidth}\LARGE\centering
\begin{minipage}{0.95\linewidth}\raggedright\linespread{1.25}\selectfont Soit $f(x)=-x^2+6x-5$. Quel est le maximum de $f$ sur $\mathbb{R}$~?\end{minipage}
\end{minipage}

\vfill
\begin{tcolorbox}[enhanced,nobeforeafter,width=7.1cm,height=2.1cm,valign=center,colback=cder!6,colframe=cder,arc=2.5mm,boxrule=0.9pt,left=3mm]{\Large\bfseries\color{cder}a)}\quad\begin{minipage}[c]{5.3cm}\raggedright\Large $4$\end{minipage}\end{tcolorbox}\hfill\begin{tcolorbox}[enhanced,nobeforeafter,width=7.1cm,height=2.1cm,valign=center,colback=cder!6,colframe=cder,arc=2.5mm,boxrule=0.9pt,left=3mm]{\Large\bfseries\color{cder}b)}\quad\begin{minipage}[c]{5.3cm}\raggedright\Large $3$\end{minipage}\end{tcolorbox}\par\vspace{2.5mm}
\begin{tcolorbox}[enhanced,nobeforeafter,width=7.1cm,height=2.1cm,valign=center,colback=cder!6,colframe=cder,arc=2.5mm,boxrule=0.9pt,left=3mm]{\Large\bfseries\color{cder}c)}\quad\begin{minipage}[c]{5.3cm}\raggedright\Large $-4$\end{minipage}\end{tcolorbox}\hfill\begin{tcolorbox}[enhanced,nobeforeafter,width=7.1cm,height=2.1cm,valign=center,colback=cder!6,colframe=cder,arc=2.5mm,boxrule=0.9pt,left=3mm]{\Large\bfseries\color{cder}d)}\quad\begin{minipage}[c]{5.3cm}\raggedright\Large $-5$\end{minipage}\end{tcolorbox}
\end{tcolorbox}
\newpage
\begin{tcolorbox}[enhanced,width=16cm,height=14cm,colback=white,colframe=ctri,boxrule=2pt,arc=0mm,
left=6mm,right=6mm,top=0mm,bottom=4mm,sharp corners,
overlay={\fill[ctri] (frame.north west) rectangle ([yshift=-1.25cm]frame.north east);
\node[anchor=west,text=white,font=\sffamily\bfseries\LARGE] at ([xshift=5mm,yshift=-0.62cm]frame.north west) {Question 49};
\node[anchor=east,text=white,font=\sffamily\itshape\large] at ([xshift=-5mm,yshift=-0.62cm]frame.north east) {Trigonométrie et exponentielle};}]
\vspace*{1.75cm}
\begin{minipage}[t][6.9cm][c]{\linewidth}\LARGE\centering
\begin{minipage}{0.95\linewidth}\raggedright\linespread{1.25}\selectfont Quelle est la mesure en radians d'un angle de $60^\circ$~?\end{minipage}
\end{minipage}

\vfill
\begin{tcolorbox}[enhanced,nobeforeafter,width=7.1cm,height=2.1cm,valign=center,colback=ctri!6,colframe=ctri,arc=2.5mm,boxrule=0.9pt,left=3mm]{\Large\bfseries\color{ctri}a)}\quad\begin{minipage}[c]{5.3cm}\raggedright\Large $\dfrac{\pi}{6}$\end{minipage}\end{tcolorbox}\hfill\begin{tcolorbox}[enhanced,nobeforeafter,width=7.1cm,height=2.1cm,valign=center,colback=ctri!6,colframe=ctri,arc=2.5mm,boxrule=0.9pt,left=3mm]{\Large\bfseries\color{ctri}b)}\quad\begin{minipage}[c]{5.3cm}\raggedright\Large $\dfrac{\pi}{3}$\end{minipage}\end{tcolorbox}\par\vspace{2.5mm}
\begin{tcolorbox}[enhanced,nobeforeafter,width=7.1cm,height=2.1cm,valign=center,colback=ctri!6,colframe=ctri,arc=2.5mm,boxrule=0.9pt,left=3mm]{\Large\bfseries\color{ctri}c)}\quad\begin{minipage}[c]{5.3cm}\raggedright\Large $3\pi$\end{minipage}\end{tcolorbox}\hfill\begin{tcolorbox}[enhanced,nobeforeafter,width=7.1cm,height=2.1cm,valign=center,colback=ctri!6,colframe=ctri,arc=2.5mm,boxrule=0.9pt,left=3mm]{\Large\bfseries\color{ctri}d)}\quad\begin{minipage}[c]{5.3cm}\raggedright\Large $\dfrac{2\pi}{3}$\end{minipage}\end{tcolorbox}
\end{tcolorbox}
\newpage
\begin{tcolorbox}[enhanced,width=16cm,height=14cm,colback=white,colframe=ctri,boxrule=2pt,arc=0mm,
left=6mm,right=6mm,top=0mm,bottom=4mm,sharp corners,
overlay={\fill[ctri] (frame.north west) rectangle ([yshift=-1.25cm]frame.north east);
\node[anchor=west,text=white,font=\sffamily\bfseries\LARGE] at ([xshift=5mm,yshift=-0.62cm]frame.north west) {Question 50};
\node[anchor=east,text=white,font=\sffamily\itshape\large] at ([xshift=-5mm,yshift=-0.62cm]frame.north east) {Trigonométrie et exponentielle};}]
\vspace*{1.75cm}
\begin{minipage}[t][6.9cm][c]{\linewidth}\LARGE\centering
\begin{minipage}{0.95\linewidth}\raggedright\linespread{1.25}\selectfont Quelle est la mesure en radians d'un angle de $45^\circ$~?\end{minipage}
\end{minipage}

\vfill
\begin{tcolorbox}[enhanced,nobeforeafter,width=7.1cm,height=2.1cm,valign=center,colback=ctri!6,colframe=ctri,arc=2.5mm,boxrule=0.9pt,left=3mm]{\Large\bfseries\color{ctri}a)}\quad\begin{minipage}[c]{5.3cm}\raggedright\Large $\dfrac{\pi}{4}$\end{minipage}\end{tcolorbox}\hfill\begin{tcolorbox}[enhanced,nobeforeafter,width=7.1cm,height=2.1cm,valign=center,colback=ctri!6,colframe=ctri,arc=2.5mm,boxrule=0.9pt,left=3mm]{\Large\bfseries\color{ctri}b)}\quad\begin{minipage}[c]{5.3cm}\raggedright\Large $\dfrac{\pi}{8}$\end{minipage}\end{tcolorbox}\par\vspace{2.5mm}
\begin{tcolorbox}[enhanced,nobeforeafter,width=7.1cm,height=2.1cm,valign=center,colback=ctri!6,colframe=ctri,arc=2.5mm,boxrule=0.9pt,left=3mm]{\Large\bfseries\color{ctri}c)}\quad\begin{minipage}[c]{5.3cm}\raggedright\Large $4\pi$\end{minipage}\end{tcolorbox}\hfill\begin{tcolorbox}[enhanced,nobeforeafter,width=7.1cm,height=2.1cm,valign=center,colback=ctri!6,colframe=ctri,arc=2.5mm,boxrule=0.9pt,left=3mm]{\Large\bfseries\color{ctri}d)}\quad\begin{minipage}[c]{5.3cm}\raggedright\Large $\dfrac{\pi}{2}$\end{minipage}\end{tcolorbox}
\end{tcolorbox}
\newpage
\begin{tcolorbox}[enhanced,width=16cm,height=14cm,colback=white,colframe=ctri,boxrule=2pt,arc=0mm,
left=6mm,right=6mm,top=0mm,bottom=4mm,sharp corners,
overlay={\fill[ctri] (frame.north west) rectangle ([yshift=-1.25cm]frame.north east);
\node[anchor=west,text=white,font=\sffamily\bfseries\LARGE] at ([xshift=5mm,yshift=-0.62cm]frame.north west) {Question 51};
\node[anchor=east,text=white,font=\sffamily\itshape\large] at ([xshift=-5mm,yshift=-0.62cm]frame.north east) {Trigonométrie et exponentielle};}]
\vspace*{1.75cm}
\begin{minipage}[t][6.9cm][c]{\linewidth}\LARGE\centering
\begin{minipage}{0.95\linewidth}\raggedright\linespread{1.25}\selectfont Quelle est la mesure en degrés d'un angle de $\dfrac{\pi}{5}$ radians~?\end{minipage}
\end{minipage}

\vfill
\begin{tcolorbox}[enhanced,nobeforeafter,width=7.1cm,height=2.1cm,valign=center,colback=ctri!6,colframe=ctri,arc=2.5mm,boxrule=0.9pt,left=3mm]{\Large\bfseries\color{ctri}a)}\quad\begin{minipage}[c]{5.3cm}\raggedright\Large $72^\circ$\end{minipage}\end{tcolorbox}\hfill\begin{tcolorbox}[enhanced,nobeforeafter,width=7.1cm,height=2.1cm,valign=center,colback=ctri!6,colframe=ctri,arc=2.5mm,boxrule=0.9pt,left=3mm]{\Large\bfseries\color{ctri}b)}\quad\begin{minipage}[c]{5.3cm}\raggedright\Large $36^\circ$\end{minipage}\end{tcolorbox}\par\vspace{2.5mm}
\begin{tcolorbox}[enhanced,nobeforeafter,width=7.1cm,height=2.1cm,valign=center,colback=ctri!6,colframe=ctri,arc=2.5mm,boxrule=0.9pt,left=3mm]{\Large\bfseries\color{ctri}c)}\quad\begin{minipage}[c]{5.3cm}\raggedright\Large $324^\circ$\end{minipage}\end{tcolorbox}\hfill\begin{tcolorbox}[enhanced,nobeforeafter,width=7.1cm,height=2.1cm,valign=center,colback=ctri!6,colframe=ctri,arc=2.5mm,boxrule=0.9pt,left=3mm]{\Large\bfseries\color{ctri}d)}\quad\begin{minipage}[c]{5.3cm}\raggedright\Large $18^\circ$\end{minipage}\end{tcolorbox}
\end{tcolorbox}
\newpage
\begin{tcolorbox}[enhanced,width=16cm,height=14cm,colback=white,colframe=ctri,boxrule=2pt,arc=0mm,
left=6mm,right=6mm,top=0mm,bottom=4mm,sharp corners,
overlay={\fill[ctri] (frame.north west) rectangle ([yshift=-1.25cm]frame.north east);
\node[anchor=west,text=white,font=\sffamily\bfseries\LARGE] at ([xshift=5mm,yshift=-0.62cm]frame.north west) {Question 52};
\node[anchor=east,text=white,font=\sffamily\itshape\large] at ([xshift=-5mm,yshift=-0.62cm]frame.north east) {Trigonométrie et exponentielle};}]
\vspace*{1.75cm}
\begin{minipage}[t][6.9cm][c]{\linewidth}\LARGE\centering
\begin{minipage}{0.95\linewidth}\raggedright\linespread{1.25}\selectfont Quelle est la mesure en degrés d'un angle de $\dfrac{2\pi}{3}$ radians~?\end{minipage}
\end{minipage}

\vfill
\begin{tcolorbox}[enhanced,nobeforeafter,width=7.1cm,height=2.1cm,valign=center,colback=ctri!6,colframe=ctri,arc=2.5mm,boxrule=0.9pt,left=3mm]{\Large\bfseries\color{ctri}a)}\quad\begin{minipage}[c]{5.3cm}\raggedright\Large $150^\circ$\end{minipage}\end{tcolorbox}\hfill\begin{tcolorbox}[enhanced,nobeforeafter,width=7.1cm,height=2.1cm,valign=center,colback=ctri!6,colframe=ctri,arc=2.5mm,boxrule=0.9pt,left=3mm]{\Large\bfseries\color{ctri}b)}\quad\begin{minipage}[c]{5.3cm}\raggedright\Large $60^\circ$\end{minipage}\end{tcolorbox}\par\vspace{2.5mm}
\begin{tcolorbox}[enhanced,nobeforeafter,width=7.1cm,height=2.1cm,valign=center,colback=ctri!6,colframe=ctri,arc=2.5mm,boxrule=0.9pt,left=3mm]{\Large\bfseries\color{ctri}c)}\quad\begin{minipage}[c]{5.3cm}\raggedright\Large $240^\circ$\end{minipage}\end{tcolorbox}\hfill\begin{tcolorbox}[enhanced,nobeforeafter,width=7.1cm,height=2.1cm,valign=center,colback=ctri!6,colframe=ctri,arc=2.5mm,boxrule=0.9pt,left=3mm]{\Large\bfseries\color{ctri}d)}\quad\begin{minipage}[c]{5.3cm}\raggedright\Large $120^\circ$\end{minipage}\end{tcolorbox}
\end{tcolorbox}
\newpage
\begin{tcolorbox}[enhanced,width=16cm,height=14cm,colback=white,colframe=ctri,boxrule=2pt,arc=0mm,
left=6mm,right=6mm,top=0mm,bottom=4mm,sharp corners,
overlay={\fill[ctri] (frame.north west) rectangle ([yshift=-1.25cm]frame.north east);
\node[anchor=west,text=white,font=\sffamily\bfseries\LARGE] at ([xshift=5mm,yshift=-0.62cm]frame.north west) {Question 53};
\node[anchor=east,text=white,font=\sffamily\itshape\large] at ([xshift=-5mm,yshift=-0.62cm]frame.north east) {Trigonométrie et exponentielle};}]
\vspace*{1.75cm}
\begin{minipage}[t][6.9cm][c]{\linewidth}\LARGE\centering
\begin{minipage}{0.95\linewidth}\raggedright\linespread{1.25}\selectfont Que vaut $\cos\left(\dfrac{\pi}{3}\right)$~?\end{minipage}
\end{minipage}

\vfill
\begin{tcolorbox}[enhanced,nobeforeafter,width=7.1cm,height=2.1cm,valign=center,colback=ctri!6,colframe=ctri,arc=2.5mm,boxrule=0.9pt,left=3mm]{\Large\bfseries\color{ctri}a)}\quad\begin{minipage}[c]{5.3cm}\raggedright\Large $\dfrac12$\end{minipage}\end{tcolorbox}\hfill\begin{tcolorbox}[enhanced,nobeforeafter,width=7.1cm,height=2.1cm,valign=center,colback=ctri!6,colframe=ctri,arc=2.5mm,boxrule=0.9pt,left=3mm]{\Large\bfseries\color{ctri}b)}\quad\begin{minipage}[c]{5.3cm}\raggedright\Large $-\dfrac12$\end{minipage}\end{tcolorbox}\par\vspace{2.5mm}
\begin{tcolorbox}[enhanced,nobeforeafter,width=7.1cm,height=2.1cm,valign=center,colback=ctri!6,colframe=ctri,arc=2.5mm,boxrule=0.9pt,left=3mm]{\Large\bfseries\color{ctri}c)}\quad\begin{minipage}[c]{5.3cm}\raggedright\Large $\dfrac{\sqrt2}{2}$\end{minipage}\end{tcolorbox}\hfill\begin{tcolorbox}[enhanced,nobeforeafter,width=7.1cm,height=2.1cm,valign=center,colback=ctri!6,colframe=ctri,arc=2.5mm,boxrule=0.9pt,left=3mm]{\Large\bfseries\color{ctri}d)}\quad\begin{minipage}[c]{5.3cm}\raggedright\Large $0$\end{minipage}\end{tcolorbox}
\end{tcolorbox}
\newpage
\begin{tcolorbox}[enhanced,width=16cm,height=14cm,colback=white,colframe=ctri,boxrule=2pt,arc=0mm,
left=6mm,right=6mm,top=0mm,bottom=4mm,sharp corners,
overlay={\fill[ctri] (frame.north west) rectangle ([yshift=-1.25cm]frame.north east);
\node[anchor=west,text=white,font=\sffamily\bfseries\LARGE] at ([xshift=5mm,yshift=-0.62cm]frame.north west) {Question 54};
\node[anchor=east,text=white,font=\sffamily\itshape\large] at ([xshift=-5mm,yshift=-0.62cm]frame.north east) {Trigonométrie et exponentielle};}]
\vspace*{1.75cm}
\begin{minipage}[t][6.9cm][c]{\linewidth}\LARGE\centering
\begin{minipage}{0.95\linewidth}\raggedright\linespread{1.25}\selectfont Que vaut $\cos\left(\dfrac{\pi}{4}\right)$~?\end{minipage}
\end{minipage}

\vfill
\begin{tcolorbox}[enhanced,nobeforeafter,width=7.1cm,height=2.1cm,valign=center,colback=ctri!6,colframe=ctri,arc=2.5mm,boxrule=0.9pt,left=3mm]{\Large\bfseries\color{ctri}a)}\quad\begin{minipage}[c]{5.3cm}\raggedright\Large $-\dfrac{\sqrt2}{2}$\end{minipage}\end{tcolorbox}\hfill\begin{tcolorbox}[enhanced,nobeforeafter,width=7.1cm,height=2.1cm,valign=center,colback=ctri!6,colframe=ctri,arc=2.5mm,boxrule=0.9pt,left=3mm]{\Large\bfseries\color{ctri}b)}\quad\begin{minipage}[c]{5.3cm}\raggedright\Large $\dfrac{\sqrt3}{2}$\end{minipage}\end{tcolorbox}\par\vspace{2.5mm}
\begin{tcolorbox}[enhanced,nobeforeafter,width=7.1cm,height=2.1cm,valign=center,colback=ctri!6,colframe=ctri,arc=2.5mm,boxrule=0.9pt,left=3mm]{\Large\bfseries\color{ctri}c)}\quad\begin{minipage}[c]{5.3cm}\raggedright\Large $\dfrac12$\end{minipage}\end{tcolorbox}\hfill\begin{tcolorbox}[enhanced,nobeforeafter,width=7.1cm,height=2.1cm,valign=center,colback=ctri!6,colframe=ctri,arc=2.5mm,boxrule=0.9pt,left=3mm]{\Large\bfseries\color{ctri}d)}\quad\begin{minipage}[c]{5.3cm}\raggedright\Large $\dfrac{\sqrt2}{2}$\end{minipage}\end{tcolorbox}
\end{tcolorbox}
\newpage
\begin{tcolorbox}[enhanced,width=16cm,height=14cm,colback=white,colframe=ctri,boxrule=2pt,arc=0mm,
left=6mm,right=6mm,top=0mm,bottom=4mm,sharp corners,
overlay={\fill[ctri] (frame.north west) rectangle ([yshift=-1.25cm]frame.north east);
\node[anchor=west,text=white,font=\sffamily\bfseries\LARGE] at ([xshift=5mm,yshift=-0.62cm]frame.north west) {Question 55};
\node[anchor=east,text=white,font=\sffamily\itshape\large] at ([xshift=-5mm,yshift=-0.62cm]frame.north east) {Trigonométrie et exponentielle};}]
\vspace*{1.75cm}
\begin{minipage}[t][6.9cm][c]{\linewidth}\LARGE\centering
\begin{minipage}{0.95\linewidth}\raggedright\linespread{1.25}\selectfont Que vaut $\sin\left(\dfrac{\pi}{6}\right)$~?\end{minipage}
\end{minipage}

\vfill
\begin{tcolorbox}[enhanced,nobeforeafter,width=7.1cm,height=2.1cm,valign=center,colback=ctri!6,colframe=ctri,arc=2.5mm,boxrule=0.9pt,left=3mm]{\Large\bfseries\color{ctri}a)}\quad\begin{minipage}[c]{5.3cm}\raggedright\Large $\dfrac{\sqrt2}{2}$\end{minipage}\end{tcolorbox}\hfill\begin{tcolorbox}[enhanced,nobeforeafter,width=7.1cm,height=2.1cm,valign=center,colback=ctri!6,colframe=ctri,arc=2.5mm,boxrule=0.9pt,left=3mm]{\Large\bfseries\color{ctri}b)}\quad\begin{minipage}[c]{5.3cm}\raggedright\Large $0$\end{minipage}\end{tcolorbox}\par\vspace{2.5mm}
\begin{tcolorbox}[enhanced,nobeforeafter,width=7.1cm,height=2.1cm,valign=center,colback=ctri!6,colframe=ctri,arc=2.5mm,boxrule=0.9pt,left=3mm]{\Large\bfseries\color{ctri}c)}\quad\begin{minipage}[c]{5.3cm}\raggedright\Large $-\dfrac12$\end{minipage}\end{tcolorbox}\hfill\begin{tcolorbox}[enhanced,nobeforeafter,width=7.1cm,height=2.1cm,valign=center,colback=ctri!6,colframe=ctri,arc=2.5mm,boxrule=0.9pt,left=3mm]{\Large\bfseries\color{ctri}d)}\quad\begin{minipage}[c]{5.3cm}\raggedright\Large $\dfrac12$\end{minipage}\end{tcolorbox}
\end{tcolorbox}
\newpage
\begin{tcolorbox}[enhanced,width=16cm,height=14cm,colback=white,colframe=ctri,boxrule=2pt,arc=0mm,
left=6mm,right=6mm,top=0mm,bottom=4mm,sharp corners,
overlay={\fill[ctri] (frame.north west) rectangle ([yshift=-1.25cm]frame.north east);
\node[anchor=west,text=white,font=\sffamily\bfseries\LARGE] at ([xshift=5mm,yshift=-0.62cm]frame.north west) {Question 56};
\node[anchor=east,text=white,font=\sffamily\itshape\large] at ([xshift=-5mm,yshift=-0.62cm]frame.north east) {Trigonométrie et exponentielle};}]
\vspace*{1.75cm}
\begin{minipage}[t][6.9cm][c]{\linewidth}\LARGE\centering
\begin{minipage}{0.95\linewidth}\raggedright\linespread{1.25}\selectfont Que vaut $\sin\left(\dfrac{\pi}{4}\right)$~?\end{minipage}
\end{minipage}

\vfill
\begin{tcolorbox}[enhanced,nobeforeafter,width=7.1cm,height=2.1cm,valign=center,colback=ctri!6,colframe=ctri,arc=2.5mm,boxrule=0.9pt,left=3mm]{\Large\bfseries\color{ctri}a)}\quad\begin{minipage}[c]{5.3cm}\raggedright\Large $\dfrac12$\end{minipage}\end{tcolorbox}\hfill\begin{tcolorbox}[enhanced,nobeforeafter,width=7.1cm,height=2.1cm,valign=center,colback=ctri!6,colframe=ctri,arc=2.5mm,boxrule=0.9pt,left=3mm]{\Large\bfseries\color{ctri}b)}\quad\begin{minipage}[c]{5.3cm}\raggedright\Large $\dfrac{\sqrt3}{2}$\end{minipage}\end{tcolorbox}\par\vspace{2.5mm}
\begin{tcolorbox}[enhanced,nobeforeafter,width=7.1cm,height=2.1cm,valign=center,colback=ctri!6,colframe=ctri,arc=2.5mm,boxrule=0.9pt,left=3mm]{\Large\bfseries\color{ctri}c)}\quad\begin{minipage}[c]{5.3cm}\raggedright\Large $\dfrac{\sqrt2}{2}$\end{minipage}\end{tcolorbox}\hfill\begin{tcolorbox}[enhanced,nobeforeafter,width=7.1cm,height=2.1cm,valign=center,colback=ctri!6,colframe=ctri,arc=2.5mm,boxrule=0.9pt,left=3mm]{\Large\bfseries\color{ctri}d)}\quad\begin{minipage}[c]{5.3cm}\raggedright\Large $-\dfrac{\sqrt2}{2}$\end{minipage}\end{tcolorbox}
\end{tcolorbox}
\newpage
\begin{tcolorbox}[enhanced,width=16cm,height=14cm,colback=white,colframe=ctri,boxrule=2pt,arc=0mm,
left=6mm,right=6mm,top=0mm,bottom=4mm,sharp corners,
overlay={\fill[ctri] (frame.north west) rectangle ([yshift=-1.25cm]frame.north east);
\node[anchor=west,text=white,font=\sffamily\bfseries\LARGE] at ([xshift=5mm,yshift=-0.62cm]frame.north west) {Question 57};
\node[anchor=east,text=white,font=\sffamily\itshape\large] at ([xshift=-5mm,yshift=-0.62cm]frame.north east) {Trigonométrie et exponentielle};}]
\vspace*{1.75cm}
\begin{minipage}[t][6.9cm][c]{\linewidth}\LARGE\centering
\begin{minipage}{0.95\linewidth}\raggedright\linespread{1.25}\selectfont On sait que $\cos x=0{,}6$. Que vaut $\cos(\pi-x)$~?\end{minipage}
\end{minipage}

\vfill
\begin{tcolorbox}[enhanced,nobeforeafter,width=7.1cm,height=2.1cm,valign=center,colback=ctri!6,colframe=ctri,arc=2.5mm,boxrule=0.9pt,left=3mm]{\Large\bfseries\color{ctri}a)}\quad\begin{minipage}[c]{5.3cm}\raggedright\Large $0{,}64$\end{minipage}\end{tcolorbox}\hfill\begin{tcolorbox}[enhanced,nobeforeafter,width=7.1cm,height=2.1cm,valign=center,colback=ctri!6,colframe=ctri,arc=2.5mm,boxrule=0.9pt,left=3mm]{\Large\bfseries\color{ctri}b)}\quad\begin{minipage}[c]{5.3cm}\raggedright\Large $-0{,}6$\end{minipage}\end{tcolorbox}\par\vspace{2.5mm}
\begin{tcolorbox}[enhanced,nobeforeafter,width=7.1cm,height=2.1cm,valign=center,colback=ctri!6,colframe=ctri,arc=2.5mm,boxrule=0.9pt,left=3mm]{\Large\bfseries\color{ctri}c)}\quad\begin{minipage}[c]{5.3cm}\raggedright\Large $0{,}4$\end{minipage}\end{tcolorbox}\hfill\begin{tcolorbox}[enhanced,nobeforeafter,width=7.1cm,height=2.1cm,valign=center,colback=ctri!6,colframe=ctri,arc=2.5mm,boxrule=0.9pt,left=3mm]{\Large\bfseries\color{ctri}d)}\quad\begin{minipage}[c]{5.3cm}\raggedright\Large $0{,}6$\end{minipage}\end{tcolorbox}
\end{tcolorbox}
\newpage
\begin{tcolorbox}[enhanced,width=16cm,height=14cm,colback=white,colframe=ctri,boxrule=2pt,arc=0mm,
left=6mm,right=6mm,top=0mm,bottom=4mm,sharp corners,
overlay={\fill[ctri] (frame.north west) rectangle ([yshift=-1.25cm]frame.north east);
\node[anchor=west,text=white,font=\sffamily\bfseries\LARGE] at ([xshift=5mm,yshift=-0.62cm]frame.north west) {Question 58};
\node[anchor=east,text=white,font=\sffamily\itshape\large] at ([xshift=-5mm,yshift=-0.62cm]frame.north east) {Trigonométrie et exponentielle};}]
\vspace*{1.75cm}
\begin{minipage}[t][6.9cm][c]{\linewidth}\LARGE\centering
\begin{minipage}{0.95\linewidth}\raggedright\linespread{1.25}\selectfont On sait que $\cos x=0{,}8$. Que vaut $\cos(-x)$~?\end{minipage}
\end{minipage}

\vfill
\begin{tcolorbox}[enhanced,nobeforeafter,width=7.1cm,height=2.1cm,valign=center,colback=ctri!6,colframe=ctri,arc=2.5mm,boxrule=0.9pt,left=3mm]{\Large\bfseries\color{ctri}a)}\quad\begin{minipage}[c]{5.3cm}\raggedright\Large $0{,}36$\end{minipage}\end{tcolorbox}\hfill\begin{tcolorbox}[enhanced,nobeforeafter,width=7.1cm,height=2.1cm,valign=center,colback=ctri!6,colframe=ctri,arc=2.5mm,boxrule=0.9pt,left=3mm]{\Large\bfseries\color{ctri}b)}\quad\begin{minipage}[c]{5.3cm}\raggedright\Large $-0{,}8$\end{minipage}\end{tcolorbox}\par\vspace{2.5mm}
\begin{tcolorbox}[enhanced,nobeforeafter,width=7.1cm,height=2.1cm,valign=center,colback=ctri!6,colframe=ctri,arc=2.5mm,boxrule=0.9pt,left=3mm]{\Large\bfseries\color{ctri}c)}\quad\begin{minipage}[c]{5.3cm}\raggedright\Large $0{,}2$\end{minipage}\end{tcolorbox}\hfill\begin{tcolorbox}[enhanced,nobeforeafter,width=7.1cm,height=2.1cm,valign=center,colback=ctri!6,colframe=ctri,arc=2.5mm,boxrule=0.9pt,left=3mm]{\Large\bfseries\color{ctri}d)}\quad\begin{minipage}[c]{5.3cm}\raggedright\Large $0{,}8$\end{minipage}\end{tcolorbox}
\end{tcolorbox}
\newpage
\begin{tcolorbox}[enhanced,width=16cm,height=14cm,colback=white,colframe=ctri,boxrule=2pt,arc=0mm,
left=6mm,right=6mm,top=0mm,bottom=4mm,sharp corners,
overlay={\fill[ctri] (frame.north west) rectangle ([yshift=-1.25cm]frame.north east);
\node[anchor=west,text=white,font=\sffamily\bfseries\LARGE] at ([xshift=5mm,yshift=-0.62cm]frame.north west) {Question 59};
\node[anchor=east,text=white,font=\sffamily\itshape\large] at ([xshift=-5mm,yshift=-0.62cm]frame.north east) {Trigonométrie et exponentielle};}]
\vspace*{1.75cm}
\begin{minipage}[t][6.9cm][c]{\linewidth}\LARGE\centering
\begin{minipage}{0.95\linewidth}\raggedright\linespread{1.25}\selectfont Quelle est la mesure principale de l'angle $\dfrac{17\pi}{6}$~?\end{minipage}
\end{minipage}

\vfill
\begin{tcolorbox}[enhanced,nobeforeafter,width=7.1cm,height=2.1cm,valign=center,colback=ctri!6,colframe=ctri,arc=2.5mm,boxrule=0.9pt,left=3mm]{\Large\bfseries\color{ctri}a)}\quad\begin{minipage}[c]{5.3cm}\raggedright\Large $-\dfrac{5\pi}{6}$\end{minipage}\end{tcolorbox}\hfill\begin{tcolorbox}[enhanced,nobeforeafter,width=7.1cm,height=2.1cm,valign=center,colback=ctri!6,colframe=ctri,arc=2.5mm,boxrule=0.9pt,left=3mm]{\Large\bfseries\color{ctri}b)}\quad\begin{minipage}[c]{5.3cm}\raggedright\Large $\dfrac{5\pi}{12}$\end{minipage}\end{tcolorbox}\par\vspace{2.5mm}
\begin{tcolorbox}[enhanced,nobeforeafter,width=7.1cm,height=2.1cm,valign=center,colback=ctri!6,colframe=ctri,arc=2.5mm,boxrule=0.9pt,left=3mm]{\Large\bfseries\color{ctri}c)}\quad\begin{minipage}[c]{5.3cm}\raggedright\Large $-\dfrac{\pi}{6}$\end{minipage}\end{tcolorbox}\hfill\begin{tcolorbox}[enhanced,nobeforeafter,width=7.1cm,height=2.1cm,valign=center,colback=ctri!6,colframe=ctri,arc=2.5mm,boxrule=0.9pt,left=3mm]{\Large\bfseries\color{ctri}d)}\quad\begin{minipage}[c]{5.3cm}\raggedright\Large $\dfrac{5\pi}{6}$\end{minipage}\end{tcolorbox}
\end{tcolorbox}
\newpage
\begin{tcolorbox}[enhanced,width=16cm,height=14cm,colback=white,colframe=ctri,boxrule=2pt,arc=0mm,
left=6mm,right=6mm,top=0mm,bottom=4mm,sharp corners,
overlay={\fill[ctri] (frame.north west) rectangle ([yshift=-1.25cm]frame.north east);
\node[anchor=west,text=white,font=\sffamily\bfseries\LARGE] at ([xshift=5mm,yshift=-0.62cm]frame.north west) {Question 60};
\node[anchor=east,text=white,font=\sffamily\itshape\large] at ([xshift=-5mm,yshift=-0.62cm]frame.north east) {Trigonométrie et exponentielle};}]
\vspace*{1.75cm}
\begin{minipage}[t][6.9cm][c]{\linewidth}\LARGE\centering
\begin{minipage}{0.95\linewidth}\raggedright\linespread{1.25}\selectfont Quelle est la mesure principale de l'angle $\dfrac{7\pi}{3}$~?\end{minipage}
\end{minipage}

\vfill
\begin{tcolorbox}[enhanced,nobeforeafter,width=7.1cm,height=2.1cm,valign=center,colback=ctri!6,colframe=ctri,arc=2.5mm,boxrule=0.9pt,left=3mm]{\Large\bfseries\color{ctri}a)}\quad\begin{minipage}[c]{5.3cm}\raggedright\Large $\dfrac{\pi}{6}$\end{minipage}\end{tcolorbox}\hfill\begin{tcolorbox}[enhanced,nobeforeafter,width=7.1cm,height=2.1cm,valign=center,colback=ctri!6,colframe=ctri,arc=2.5mm,boxrule=0.9pt,left=3mm]{\Large\bfseries\color{ctri}b)}\quad\begin{minipage}[c]{5.3cm}\raggedright\Large $-\dfrac{2\pi}{3}$\end{minipage}\end{tcolorbox}\par\vspace{2.5mm}
\begin{tcolorbox}[enhanced,nobeforeafter,width=7.1cm,height=2.1cm,valign=center,colback=ctri!6,colframe=ctri,arc=2.5mm,boxrule=0.9pt,left=3mm]{\Large\bfseries\color{ctri}c)}\quad\begin{minipage}[c]{5.3cm}\raggedright\Large $-\dfrac{\pi}{3}$\end{minipage}\end{tcolorbox}\hfill\begin{tcolorbox}[enhanced,nobeforeafter,width=7.1cm,height=2.1cm,valign=center,colback=ctri!6,colframe=ctri,arc=2.5mm,boxrule=0.9pt,left=3mm]{\Large\bfseries\color{ctri}d)}\quad\begin{minipage}[c]{5.3cm}\raggedright\Large $\dfrac{\pi}{3}$\end{minipage}\end{tcolorbox}
\end{tcolorbox}
\newpage
\begin{tcolorbox}[enhanced,width=16cm,height=14cm,colback=white,colframe=ctri,boxrule=2pt,arc=0mm,
left=6mm,right=6mm,top=0mm,bottom=4mm,sharp corners,
overlay={\fill[ctri] (frame.north west) rectangle ([yshift=-1.25cm]frame.north east);
\node[anchor=west,text=white,font=\sffamily\bfseries\LARGE] at ([xshift=5mm,yshift=-0.62cm]frame.north west) {Question 61};
\node[anchor=east,text=white,font=\sffamily\itshape\large] at ([xshift=-5mm,yshift=-0.62cm]frame.north east) {Trigonométrie et exponentielle};}]
\vspace*{1.75cm}
\begin{minipage}[t][6.9cm][c]{\linewidth}\LARGE\centering
\begin{minipage}{0.95\linewidth}\raggedright\linespread{1.25}\selectfont On sait que $x\in[0\,;\,\pi]$ et $\cos x=\dfrac{3}{5}$. Que vaut $\sin x$~?\end{minipage}
\end{minipage}

\vfill
\begin{tcolorbox}[enhanced,nobeforeafter,width=7.1cm,height=2.1cm,valign=center,colback=ctri!6,colframe=ctri,arc=2.5mm,boxrule=0.9pt,left=3mm]{\Large\bfseries\color{ctri}a)}\quad\begin{minipage}[c]{5.3cm}\raggedright\Large $\dfrac{2}{5}$\end{minipage}\end{tcolorbox}\hfill\begin{tcolorbox}[enhanced,nobeforeafter,width=7.1cm,height=2.1cm,valign=center,colback=ctri!6,colframe=ctri,arc=2.5mm,boxrule=0.9pt,left=3mm]{\Large\bfseries\color{ctri}b)}\quad\begin{minipage}[c]{5.3cm}\raggedright\Large $-\dfrac{4}{5}$\end{minipage}\end{tcolorbox}\par\vspace{2.5mm}
\begin{tcolorbox}[enhanced,nobeforeafter,width=7.1cm,height=2.1cm,valign=center,colback=ctri!6,colframe=ctri,arc=2.5mm,boxrule=0.9pt,left=3mm]{\Large\bfseries\color{ctri}c)}\quad\begin{minipage}[c]{5.3cm}\raggedright\Large $\dfrac{4}{5}$\end{minipage}\end{tcolorbox}\hfill\begin{tcolorbox}[enhanced,nobeforeafter,width=7.1cm,height=2.1cm,valign=center,colback=ctri!6,colframe=ctri,arc=2.5mm,boxrule=0.9pt,left=3mm]{\Large\bfseries\color{ctri}d)}\quad\begin{minipage}[c]{5.3cm}\raggedright\Large $\dfrac{16}{25}$\end{minipage}\end{tcolorbox}
\end{tcolorbox}
\newpage
\begin{tcolorbox}[enhanced,width=16cm,height=14cm,colback=white,colframe=ctri,boxrule=2pt,arc=0mm,
left=6mm,right=6mm,top=0mm,bottom=4mm,sharp corners,
overlay={\fill[ctri] (frame.north west) rectangle ([yshift=-1.25cm]frame.north east);
\node[anchor=west,text=white,font=\sffamily\bfseries\LARGE] at ([xshift=5mm,yshift=-0.62cm]frame.north west) {Question 62};
\node[anchor=east,text=white,font=\sffamily\itshape\large] at ([xshift=-5mm,yshift=-0.62cm]frame.north east) {Trigonométrie et exponentielle};}]
\vspace*{1.75cm}
\begin{minipage}[t][6.9cm][c]{\linewidth}\LARGE\centering
\begin{minipage}{0.95\linewidth}\raggedright\linespread{1.25}\selectfont On sait que $x\in[0\,;\,\pi]$ et $\cos x=\dfrac{4}{5}$. Que vaut $\sin x$~?\end{minipage}
\end{minipage}

\vfill
\begin{tcolorbox}[enhanced,nobeforeafter,width=7.1cm,height=2.1cm,valign=center,colback=ctri!6,colframe=ctri,arc=2.5mm,boxrule=0.9pt,left=3mm]{\Large\bfseries\color{ctri}a)}\quad\begin{minipage}[c]{5.3cm}\raggedright\Large $-\dfrac{3}{5}$\end{minipage}\end{tcolorbox}\hfill\begin{tcolorbox}[enhanced,nobeforeafter,width=7.1cm,height=2.1cm,valign=center,colback=ctri!6,colframe=ctri,arc=2.5mm,boxrule=0.9pt,left=3mm]{\Large\bfseries\color{ctri}b)}\quad\begin{minipage}[c]{5.3cm}\raggedright\Large $\dfrac{3}{5}$\end{minipage}\end{tcolorbox}\par\vspace{2.5mm}
\begin{tcolorbox}[enhanced,nobeforeafter,width=7.1cm,height=2.1cm,valign=center,colback=ctri!6,colframe=ctri,arc=2.5mm,boxrule=0.9pt,left=3mm]{\Large\bfseries\color{ctri}c)}\quad\begin{minipage}[c]{5.3cm}\raggedright\Large $\dfrac{9}{25}$\end{minipage}\end{tcolorbox}\hfill\begin{tcolorbox}[enhanced,nobeforeafter,width=7.1cm,height=2.1cm,valign=center,colback=ctri!6,colframe=ctri,arc=2.5mm,boxrule=0.9pt,left=3mm]{\Large\bfseries\color{ctri}d)}\quad\begin{minipage}[c]{5.3cm}\raggedright\Large $\dfrac{1}{5}$\end{minipage}\end{tcolorbox}
\end{tcolorbox}
\newpage
\begin{tcolorbox}[enhanced,width=16cm,height=14cm,colback=white,colframe=ctri,boxrule=2pt,arc=0mm,
left=6mm,right=6mm,top=0mm,bottom=4mm,sharp corners,
overlay={\fill[ctri] (frame.north west) rectangle ([yshift=-1.25cm]frame.north east);
\node[anchor=west,text=white,font=\sffamily\bfseries\LARGE] at ([xshift=5mm,yshift=-0.62cm]frame.north west) {Question 63};
\node[anchor=east,text=white,font=\sffamily\itshape\large] at ([xshift=-5mm,yshift=-0.62cm]frame.north east) {Trigonométrie et exponentielle};}]
\vspace*{1.75cm}
\begin{minipage}[t][6.9cm][c]{\linewidth}\LARGE\centering
\begin{minipage}{0.95\linewidth}\raggedright\linespread{1.25}\selectfont Simplifier $\left(e^{2}\right)^{3}\times e^{-1}$.\end{minipage}
\end{minipage}

\vfill
\begin{tcolorbox}[enhanced,nobeforeafter,width=7.1cm,height=2.1cm,valign=center,colback=ctri!6,colframe=ctri,arc=2.5mm,boxrule=0.9pt,left=3mm]{\Large\bfseries\color{ctri}a)}\quad\begin{minipage}[c]{5.3cm}\raggedright\Large $e^{4}$\end{minipage}\end{tcolorbox}\hfill\begin{tcolorbox}[enhanced,nobeforeafter,width=7.1cm,height=2.1cm,valign=center,colback=ctri!6,colframe=ctri,arc=2.5mm,boxrule=0.9pt,left=3mm]{\Large\bfseries\color{ctri}b)}\quad\begin{minipage}[c]{5.3cm}\raggedright\Large $e^{7}$\end{minipage}\end{tcolorbox}\par\vspace{2.5mm}
\begin{tcolorbox}[enhanced,nobeforeafter,width=7.1cm,height=2.1cm,valign=center,colback=ctri!6,colframe=ctri,arc=2.5mm,boxrule=0.9pt,left=3mm]{\Large\bfseries\color{ctri}c)}\quad\begin{minipage}[c]{5.3cm}\raggedright\Large $e^{5}$\end{minipage}\end{tcolorbox}\hfill\begin{tcolorbox}[enhanced,nobeforeafter,width=7.1cm,height=2.1cm,valign=center,colback=ctri!6,colframe=ctri,arc=2.5mm,boxrule=0.9pt,left=3mm]{\Large\bfseries\color{ctri}d)}\quad\begin{minipage}[c]{5.3cm}\raggedright\Large $e^{-5}$\end{minipage}\end{tcolorbox}
\end{tcolorbox}
\newpage
\begin{tcolorbox}[enhanced,width=16cm,height=14cm,colback=white,colframe=ctri,boxrule=2pt,arc=0mm,
left=6mm,right=6mm,top=0mm,bottom=4mm,sharp corners,
overlay={\fill[ctri] (frame.north west) rectangle ([yshift=-1.25cm]frame.north east);
\node[anchor=west,text=white,font=\sffamily\bfseries\LARGE] at ([xshift=5mm,yshift=-0.62cm]frame.north west) {Question 64};
\node[anchor=east,text=white,font=\sffamily\itshape\large] at ([xshift=-5mm,yshift=-0.62cm]frame.north east) {Trigonométrie et exponentielle};}]
\vspace*{1.75cm}
\begin{minipage}[t][6.9cm][c]{\linewidth}\LARGE\centering
\begin{minipage}{0.95\linewidth}\raggedright\linespread{1.25}\selectfont Simplifier $\left(e^{-1}\right)^{4}\times e^{6}$.\end{minipage}
\end{minipage}

\vfill
\begin{tcolorbox}[enhanced,nobeforeafter,width=7.1cm,height=2.1cm,valign=center,colback=ctri!6,colframe=ctri,arc=2.5mm,boxrule=0.9pt,left=3mm]{\Large\bfseries\color{ctri}a)}\quad\begin{minipage}[c]{5.3cm}\raggedright\Large $e^{7}$\end{minipage}\end{tcolorbox}\hfill\begin{tcolorbox}[enhanced,nobeforeafter,width=7.1cm,height=2.1cm,valign=center,colback=ctri!6,colframe=ctri,arc=2.5mm,boxrule=0.9pt,left=3mm]{\Large\bfseries\color{ctri}b)}\quad\begin{minipage}[c]{5.3cm}\raggedright\Large $e^{2}$\end{minipage}\end{tcolorbox}\par\vspace{2.5mm}
\begin{tcolorbox}[enhanced,nobeforeafter,width=7.1cm,height=2.1cm,valign=center,colback=ctri!6,colframe=ctri,arc=2.5mm,boxrule=0.9pt,left=3mm]{\Large\bfseries\color{ctri}c)}\quad\begin{minipage}[c]{5.3cm}\raggedright\Large $e^{9}$\end{minipage}\end{tcolorbox}\hfill\begin{tcolorbox}[enhanced,nobeforeafter,width=7.1cm,height=2.1cm,valign=center,colback=ctri!6,colframe=ctri,arc=2.5mm,boxrule=0.9pt,left=3mm]{\Large\bfseries\color{ctri}d)}\quad\begin{minipage}[c]{5.3cm}\raggedright\Large $e^{-10}$\end{minipage}\end{tcolorbox}
\end{tcolorbox}
\newpage
\begin{tcolorbox}[enhanced,width=16cm,height=14cm,colback=white,colframe=cpsc,boxrule=2pt,arc=0mm,
left=6mm,right=6mm,top=0mm,bottom=4mm,sharp corners,
overlay={\fill[cpsc] (frame.north west) rectangle ([yshift=-1.25cm]frame.north east);
\node[anchor=west,text=white,font=\sffamily\bfseries\LARGE] at ([xshift=5mm,yshift=-0.62cm]frame.north west) {Question 65};
\node[anchor=east,text=white,font=\sffamily\itshape\large] at ([xshift=-5mm,yshift=-0.62cm]frame.north east) {Produit scalaire et géométrie repérée};}]
\vspace*{1.75cm}
\begin{minipage}[t][6.9cm][c]{\linewidth}\LARGE\centering
\begin{minipage}{0.95\linewidth}\raggedright\linespread{1.25}\selectfont Dans un repère orthonormé, que vaut $\vec u\cdot\vec v$ pour $\vec u\begin{pmatrix}2\\3\end{pmatrix}$ et $\vec v\begin{pmatrix}4\\-1\end{pmatrix}$~?\end{minipage}
\end{minipage}

\vfill
\begin{tcolorbox}[enhanced,nobeforeafter,width=7.1cm,height=2.1cm,valign=center,colback=cpsc!6,colframe=cpsc,arc=2.5mm,boxrule=0.9pt,left=3mm]{\Large\bfseries\color{cpsc}a)}\quad\begin{minipage}[c]{5.3cm}\raggedright\Large $11$\end{minipage}\end{tcolorbox}\hfill\begin{tcolorbox}[enhanced,nobeforeafter,width=7.1cm,height=2.1cm,valign=center,colback=cpsc!6,colframe=cpsc,arc=2.5mm,boxrule=0.9pt,left=3mm]{\Large\bfseries\color{cpsc}b)}\quad\begin{minipage}[c]{5.3cm}\raggedright\Large $10$\end{minipage}\end{tcolorbox}\par\vspace{2.5mm}
\begin{tcolorbox}[enhanced,nobeforeafter,width=7.1cm,height=2.1cm,valign=center,colback=cpsc!6,colframe=cpsc,arc=2.5mm,boxrule=0.9pt,left=3mm]{\Large\bfseries\color{cpsc}c)}\quad\begin{minipage}[c]{5.3cm}\raggedright\Large $5$\end{minipage}\end{tcolorbox}\hfill\begin{tcolorbox}[enhanced,nobeforeafter,width=7.1cm,height=2.1cm,valign=center,colback=cpsc!6,colframe=cpsc,arc=2.5mm,boxrule=0.9pt,left=3mm]{\Large\bfseries\color{cpsc}d)}\quad\begin{minipage}[c]{5.3cm}\raggedright\Large $-14$\end{minipage}\end{tcolorbox}
\end{tcolorbox}
\newpage
\begin{tcolorbox}[enhanced,width=16cm,height=14cm,colback=white,colframe=cpsc,boxrule=2pt,arc=0mm,
left=6mm,right=6mm,top=0mm,bottom=4mm,sharp corners,
overlay={\fill[cpsc] (frame.north west) rectangle ([yshift=-1.25cm]frame.north east);
\node[anchor=west,text=white,font=\sffamily\bfseries\LARGE] at ([xshift=5mm,yshift=-0.62cm]frame.north west) {Question 66};
\node[anchor=east,text=white,font=\sffamily\itshape\large] at ([xshift=-5mm,yshift=-0.62cm]frame.north east) {Produit scalaire et géométrie repérée};}]
\vspace*{1.75cm}
\begin{minipage}[t][6.9cm][c]{\linewidth}\LARGE\centering
\begin{minipage}{0.95\linewidth}\raggedright\linespread{1.25}\selectfont Dans un repère orthonormé, que vaut $\vec u\cdot\vec v$ pour $\vec u\begin{pmatrix}-1\\5\end{pmatrix}$ et $\vec v\begin{pmatrix}3\\2\end{pmatrix}$~?\end{minipage}
\end{minipage}

\vfill
\begin{tcolorbox}[enhanced,nobeforeafter,width=7.1cm,height=2.1cm,valign=center,colback=cpsc!6,colframe=cpsc,arc=2.5mm,boxrule=0.9pt,left=3mm]{\Large\bfseries\color{cpsc}a)}\quad\begin{minipage}[c]{5.3cm}\raggedright\Large $13$\end{minipage}\end{tcolorbox}\hfill\begin{tcolorbox}[enhanced,nobeforeafter,width=7.1cm,height=2.1cm,valign=center,colback=cpsc!6,colframe=cpsc,arc=2.5mm,boxrule=0.9pt,left=3mm]{\Large\bfseries\color{cpsc}b)}\quad\begin{minipage}[c]{5.3cm}\raggedright\Large $-17$\end{minipage}\end{tcolorbox}\par\vspace{2.5mm}
\begin{tcolorbox}[enhanced,nobeforeafter,width=7.1cm,height=2.1cm,valign=center,colback=cpsc!6,colframe=cpsc,arc=2.5mm,boxrule=0.9pt,left=3mm]{\Large\bfseries\color{cpsc}c)}\quad\begin{minipage}[c]{5.3cm}\raggedright\Large $7$\end{minipage}\end{tcolorbox}\hfill\begin{tcolorbox}[enhanced,nobeforeafter,width=7.1cm,height=2.1cm,valign=center,colback=cpsc!6,colframe=cpsc,arc=2.5mm,boxrule=0.9pt,left=3mm]{\Large\bfseries\color{cpsc}d)}\quad\begin{minipage}[c]{5.3cm}\raggedright\Large $-13$\end{minipage}\end{tcolorbox}
\end{tcolorbox}
\newpage
\begin{tcolorbox}[enhanced,width=16cm,height=14cm,colback=white,colframe=cpsc,boxrule=2pt,arc=0mm,
left=6mm,right=6mm,top=0mm,bottom=4mm,sharp corners,
overlay={\fill[cpsc] (frame.north west) rectangle ([yshift=-1.25cm]frame.north east);
\node[anchor=west,text=white,font=\sffamily\bfseries\LARGE] at ([xshift=5mm,yshift=-0.62cm]frame.north west) {Question 67};
\node[anchor=east,text=white,font=\sffamily\itshape\large] at ([xshift=-5mm,yshift=-0.62cm]frame.north east) {Produit scalaire et géométrie repérée};}]
\vspace*{1.75cm}
\begin{minipage}[t][6.9cm][c]{\linewidth}\LARGE\centering
\begin{minipage}{0.95\linewidth}\raggedright\linespread{1.25}\selectfont $\|\vec u\|=3$, $\|\vec v\|=4$ et l'angle entre $\vec u$ et $\vec v$ mesure $60^\circ$. Que vaut $\vec u\cdot\vec v$~?\end{minipage}
\end{minipage}

\vfill
\begin{tcolorbox}[enhanced,nobeforeafter,width=7.1cm,height=2.1cm,valign=center,colback=cpsc!6,colframe=cpsc,arc=2.5mm,boxrule=0.9pt,left=3mm]{\Large\bfseries\color{cpsc}a)}\quad\begin{minipage}[c]{5.3cm}\raggedright\Large $6$\end{minipage}\end{tcolorbox}\hfill\begin{tcolorbox}[enhanced,nobeforeafter,width=7.1cm,height=2.1cm,valign=center,colback=cpsc!6,colframe=cpsc,arc=2.5mm,boxrule=0.9pt,left=3mm]{\Large\bfseries\color{cpsc}b)}\quad\begin{minipage}[c]{5.3cm}\raggedright\Large $-6$\end{minipage}\end{tcolorbox}\par\vspace{2.5mm}
\begin{tcolorbox}[enhanced,nobeforeafter,width=7.1cm,height=2.1cm,valign=center,colback=cpsc!6,colframe=cpsc,arc=2.5mm,boxrule=0.9pt,left=3mm]{\Large\bfseries\color{cpsc}c)}\quad\begin{minipage}[c]{5.3cm}\raggedright\Large $7$\end{minipage}\end{tcolorbox}\hfill\begin{tcolorbox}[enhanced,nobeforeafter,width=7.1cm,height=2.1cm,valign=center,colback=cpsc!6,colframe=cpsc,arc=2.5mm,boxrule=0.9pt,left=3mm]{\Large\bfseries\color{cpsc}d)}\quad\begin{minipage}[c]{5.3cm}\raggedright\Large $12$\end{minipage}\end{tcolorbox}
\end{tcolorbox}
\newpage
\begin{tcolorbox}[enhanced,width=16cm,height=14cm,colback=white,colframe=cpsc,boxrule=2pt,arc=0mm,
left=6mm,right=6mm,top=0mm,bottom=4mm,sharp corners,
overlay={\fill[cpsc] (frame.north west) rectangle ([yshift=-1.25cm]frame.north east);
\node[anchor=west,text=white,font=\sffamily\bfseries\LARGE] at ([xshift=5mm,yshift=-0.62cm]frame.north west) {Question 68};
\node[anchor=east,text=white,font=\sffamily\itshape\large] at ([xshift=-5mm,yshift=-0.62cm]frame.north east) {Produit scalaire et géométrie repérée};}]
\vspace*{1.75cm}
\begin{minipage}[t][6.9cm][c]{\linewidth}\LARGE\centering
\begin{minipage}{0.95\linewidth}\raggedright\linespread{1.25}\selectfont $\|\vec u\|=2$, $\|\vec v\|=5$ et l'angle entre $\vec u$ et $\vec v$ mesure $120^\circ$. Que vaut $\vec u\cdot\vec v$~?\end{minipage}
\end{minipage}

\vfill
\begin{tcolorbox}[enhanced,nobeforeafter,width=7.1cm,height=2.1cm,valign=center,colback=cpsc!6,colframe=cpsc,arc=2.5mm,boxrule=0.9pt,left=3mm]{\Large\bfseries\color{cpsc}a)}\quad\begin{minipage}[c]{5.3cm}\raggedright\Large $5$\end{minipage}\end{tcolorbox}\hfill\begin{tcolorbox}[enhanced,nobeforeafter,width=7.1cm,height=2.1cm,valign=center,colback=cpsc!6,colframe=cpsc,arc=2.5mm,boxrule=0.9pt,left=3mm]{\Large\bfseries\color{cpsc}b)}\quad\begin{minipage}[c]{5.3cm}\raggedright\Large $10$\end{minipage}\end{tcolorbox}\par\vspace{2.5mm}
\begin{tcolorbox}[enhanced,nobeforeafter,width=7.1cm,height=2.1cm,valign=center,colback=cpsc!6,colframe=cpsc,arc=2.5mm,boxrule=0.9pt,left=3mm]{\Large\bfseries\color{cpsc}c)}\quad\begin{minipage}[c]{5.3cm}\raggedright\Large $-5$\end{minipage}\end{tcolorbox}\hfill\begin{tcolorbox}[enhanced,nobeforeafter,width=7.1cm,height=2.1cm,valign=center,colback=cpsc!6,colframe=cpsc,arc=2.5mm,boxrule=0.9pt,left=3mm]{\Large\bfseries\color{cpsc}d)}\quad\begin{minipage}[c]{5.3cm}\raggedright\Large $7$\end{minipage}\end{tcolorbox}
\end{tcolorbox}
\newpage
\begin{tcolorbox}[enhanced,width=16cm,height=14cm,colback=white,colframe=cpsc,boxrule=2pt,arc=0mm,
left=6mm,right=6mm,top=0mm,bottom=4mm,sharp corners,
overlay={\fill[cpsc] (frame.north west) rectangle ([yshift=-1.25cm]frame.north east);
\node[anchor=west,text=white,font=\sffamily\bfseries\LARGE] at ([xshift=5mm,yshift=-0.62cm]frame.north west) {Question 69};
\node[anchor=east,text=white,font=\sffamily\itshape\large] at ([xshift=-5mm,yshift=-0.62cm]frame.north east) {Produit scalaire et géométrie repérée};}]
\vspace*{1.75cm}
\begin{minipage}[t][6.9cm][c]{\linewidth}\LARGE\centering
\begin{minipage}{0.95\linewidth}\raggedright\linespread{1.25}\selectfont Pour quelle valeur de $m$ les vecteurs $\vec u\begin{pmatrix}2\\m\end{pmatrix}$ et $\vec v\begin{pmatrix}3\\-6\end{pmatrix}$ sont-ils orthogonaux~?\end{minipage}
\end{minipage}

\vfill
\begin{tcolorbox}[enhanced,nobeforeafter,width=7.1cm,height=2.1cm,valign=center,colback=cpsc!6,colframe=cpsc,arc=2.5mm,boxrule=0.9pt,left=3mm]{\Large\bfseries\color{cpsc}a)}\quad\begin{minipage}[c]{5.3cm}\raggedright\Large $m=1$\end{minipage}\end{tcolorbox}\hfill\begin{tcolorbox}[enhanced,nobeforeafter,width=7.1cm,height=2.1cm,valign=center,colback=cpsc!6,colframe=cpsc,arc=2.5mm,boxrule=0.9pt,left=3mm]{\Large\bfseries\color{cpsc}b)}\quad\begin{minipage}[c]{5.3cm}\raggedright\Large $m=2$\end{minipage}\end{tcolorbox}\par\vspace{2.5mm}
\begin{tcolorbox}[enhanced,nobeforeafter,width=7.1cm,height=2.1cm,valign=center,colback=cpsc!6,colframe=cpsc,arc=2.5mm,boxrule=0.9pt,left=3mm]{\Large\bfseries\color{cpsc}c)}\quad\begin{minipage}[c]{5.3cm}\raggedright\Large $m=-1$\end{minipage}\end{tcolorbox}\hfill\begin{tcolorbox}[enhanced,nobeforeafter,width=7.1cm,height=2.1cm,valign=center,colback=cpsc!6,colframe=cpsc,arc=2.5mm,boxrule=0.9pt,left=3mm]{\Large\bfseries\color{cpsc}d)}\quad\begin{minipage}[c]{5.3cm}\raggedright\Large $m=-6$\end{minipage}\end{tcolorbox}
\end{tcolorbox}
\newpage
\begin{tcolorbox}[enhanced,width=16cm,height=14cm,colback=white,colframe=cpsc,boxrule=2pt,arc=0mm,
left=6mm,right=6mm,top=0mm,bottom=4mm,sharp corners,
overlay={\fill[cpsc] (frame.north west) rectangle ([yshift=-1.25cm]frame.north east);
\node[anchor=west,text=white,font=\sffamily\bfseries\LARGE] at ([xshift=5mm,yshift=-0.62cm]frame.north west) {Question 70};
\node[anchor=east,text=white,font=\sffamily\itshape\large] at ([xshift=-5mm,yshift=-0.62cm]frame.north east) {Produit scalaire et géométrie repérée};}]
\vspace*{1.75cm}
\begin{minipage}[t][6.9cm][c]{\linewidth}\LARGE\centering
\begin{minipage}{0.95\linewidth}\raggedright\linespread{1.25}\selectfont Pour quelle valeur de $m$ les vecteurs $\vec u\begin{pmatrix}4\\m\end{pmatrix}$ et $\vec v\begin{pmatrix}1\\2\end{pmatrix}$ sont-ils orthogonaux~?\end{minipage}
\end{minipage}

\vfill
\begin{tcolorbox}[enhanced,nobeforeafter,width=7.1cm,height=2.1cm,valign=center,colback=cpsc!6,colframe=cpsc,arc=2.5mm,boxrule=0.9pt,left=3mm]{\Large\bfseries\color{cpsc}a)}\quad\begin{minipage}[c]{5.3cm}\raggedright\Large $m=2$\end{minipage}\end{tcolorbox}\hfill\begin{tcolorbox}[enhanced,nobeforeafter,width=7.1cm,height=2.1cm,valign=center,colback=cpsc!6,colframe=cpsc,arc=2.5mm,boxrule=0.9pt,left=3mm]{\Large\bfseries\color{cpsc}b)}\quad\begin{minipage}[c]{5.3cm}\raggedright\Large $m=-2$\end{minipage}\end{tcolorbox}\par\vspace{2.5mm}
\begin{tcolorbox}[enhanced,nobeforeafter,width=7.1cm,height=2.1cm,valign=center,colback=cpsc!6,colframe=cpsc,arc=2.5mm,boxrule=0.9pt,left=3mm]{\Large\bfseries\color{cpsc}c)}\quad\begin{minipage}[c]{5.3cm}\raggedright\Large $m=-4$\end{minipage}\end{tcolorbox}\hfill\begin{tcolorbox}[enhanced,nobeforeafter,width=7.1cm,height=2.1cm,valign=center,colback=cpsc!6,colframe=cpsc,arc=2.5mm,boxrule=0.9pt,left=3mm]{\Large\bfseries\color{cpsc}d)}\quad\begin{minipage}[c]{5.3cm}\raggedright\Large $m=-\dfrac{1}{2}$\end{minipage}\end{tcolorbox}
\end{tcolorbox}
\newpage
\begin{tcolorbox}[enhanced,width=16cm,height=14cm,colback=white,colframe=cpsc,boxrule=2pt,arc=0mm,
left=6mm,right=6mm,top=0mm,bottom=4mm,sharp corners,
overlay={\fill[cpsc] (frame.north west) rectangle ([yshift=-1.25cm]frame.north east);
\node[anchor=west,text=white,font=\sffamily\bfseries\LARGE] at ([xshift=5mm,yshift=-0.62cm]frame.north west) {Question 71};
\node[anchor=east,text=white,font=\sffamily\itshape\large] at ([xshift=-5mm,yshift=-0.62cm]frame.north east) {Produit scalaire et géométrie repérée};}]
\vspace*{1.75cm}
\begin{minipage}[t][6.9cm][c]{\linewidth}\LARGE\centering
\begin{minipage}{0.95\linewidth}\raggedright\linespread{1.25}\selectfont Dans le triangle $ABC$, $AB=3$, $AC=5$ et $\widehat{BAC}=60^\circ$. Combien mesure $BC$~?\end{minipage}
\end{minipage}

\vfill
\begin{tcolorbox}[enhanced,nobeforeafter,width=7.1cm,height=2.1cm,valign=center,colback=cpsc!6,colframe=cpsc,arc=2.5mm,boxrule=0.9pt,left=3mm]{\Large\bfseries\color{cpsc}a)}\quad\begin{minipage}[c]{5.3cm}\raggedright\Large $\sqrt{19}$\end{minipage}\end{tcolorbox}\hfill\begin{tcolorbox}[enhanced,nobeforeafter,width=7.1cm,height=2.1cm,valign=center,colback=cpsc!6,colframe=cpsc,arc=2.5mm,boxrule=0.9pt,left=3mm]{\Large\bfseries\color{cpsc}b)}\quad\begin{minipage}[c]{5.3cm}\raggedright\Large $\sqrt{34}$\end{minipage}\end{tcolorbox}\par\vspace{2.5mm}
\begin{tcolorbox}[enhanced,nobeforeafter,width=7.1cm,height=2.1cm,valign=center,colback=cpsc!6,colframe=cpsc,arc=2.5mm,boxrule=0.9pt,left=3mm]{\Large\bfseries\color{cpsc}c)}\quad\begin{minipage}[c]{5.3cm}\raggedright\Large $19$\end{minipage}\end{tcolorbox}\hfill\begin{tcolorbox}[enhanced,nobeforeafter,width=7.1cm,height=2.1cm,valign=center,colback=cpsc!6,colframe=cpsc,arc=2.5mm,boxrule=0.9pt,left=3mm]{\Large\bfseries\color{cpsc}d)}\quad\begin{minipage}[c]{5.3cm}\raggedright\Large $7$\end{minipage}\end{tcolorbox}
\end{tcolorbox}
\newpage
\begin{tcolorbox}[enhanced,width=16cm,height=14cm,colback=white,colframe=cpsc,boxrule=2pt,arc=0mm,
left=6mm,right=6mm,top=0mm,bottom=4mm,sharp corners,
overlay={\fill[cpsc] (frame.north west) rectangle ([yshift=-1.25cm]frame.north east);
\node[anchor=west,text=white,font=\sffamily\bfseries\LARGE] at ([xshift=5mm,yshift=-0.62cm]frame.north west) {Question 72};
\node[anchor=east,text=white,font=\sffamily\itshape\large] at ([xshift=-5mm,yshift=-0.62cm]frame.north east) {Produit scalaire et géométrie repérée};}]
\vspace*{1.75cm}
\begin{minipage}[t][6.9cm][c]{\linewidth}\LARGE\centering
\begin{minipage}{0.95\linewidth}\raggedright\linespread{1.25}\selectfont Dans le triangle $ABC$, $AB=4$, $AC=6$ et $\widehat{BAC}=60^\circ$. Combien mesure $BC$~?\end{minipage}
\end{minipage}

\vfill
\begin{tcolorbox}[enhanced,nobeforeafter,width=7.1cm,height=2.1cm,valign=center,colback=cpsc!6,colframe=cpsc,arc=2.5mm,boxrule=0.9pt,left=3mm]{\Large\bfseries\color{cpsc}a)}\quad\begin{minipage}[c]{5.3cm}\raggedright\Large $2\sqrt{19}$\end{minipage}\end{tcolorbox}\hfill\begin{tcolorbox}[enhanced,nobeforeafter,width=7.1cm,height=2.1cm,valign=center,colback=cpsc!6,colframe=cpsc,arc=2.5mm,boxrule=0.9pt,left=3mm]{\Large\bfseries\color{cpsc}b)}\quad\begin{minipage}[c]{5.3cm}\raggedright\Large $28$\end{minipage}\end{tcolorbox}\par\vspace{2.5mm}
\begin{tcolorbox}[enhanced,nobeforeafter,width=7.1cm,height=2.1cm,valign=center,colback=cpsc!6,colframe=cpsc,arc=2.5mm,boxrule=0.9pt,left=3mm]{\Large\bfseries\color{cpsc}c)}\quad\begin{minipage}[c]{5.3cm}\raggedright\Large $2\sqrt{13}$\end{minipage}\end{tcolorbox}\hfill\begin{tcolorbox}[enhanced,nobeforeafter,width=7.1cm,height=2.1cm,valign=center,colback=cpsc!6,colframe=cpsc,arc=2.5mm,boxrule=0.9pt,left=3mm]{\Large\bfseries\color{cpsc}d)}\quad\begin{minipage}[c]{5.3cm}\raggedright\Large $2\sqrt{7}$\end{minipage}\end{tcolorbox}
\end{tcolorbox}
\newpage
\begin{tcolorbox}[enhanced,width=16cm,height=14cm,colback=white,colframe=cpsc,boxrule=2pt,arc=0mm,
left=6mm,right=6mm,top=0mm,bottom=4mm,sharp corners,
overlay={\fill[cpsc] (frame.north west) rectangle ([yshift=-1.25cm]frame.north east);
\node[anchor=west,text=white,font=\sffamily\bfseries\LARGE] at ([xshift=5mm,yshift=-0.62cm]frame.north west) {Question 73};
\node[anchor=east,text=white,font=\sffamily\itshape\large] at ([xshift=-5mm,yshift=-0.62cm]frame.north east) {Produit scalaire et géométrie repérée};}]
\vspace*{1.75cm}
\begin{minipage}[t][6.9cm][c]{\linewidth}\LARGE\centering
\begin{minipage}{0.95\linewidth}\raggedright\linespread{1.25}\selectfont Lequel de ces vecteurs est un vecteur normal à la droite d'équation $2x-3y+5=0$~?\end{minipage}
\end{minipage}

\vfill
\begin{tcolorbox}[enhanced,nobeforeafter,width=7.1cm,height=2.1cm,valign=center,colback=cpsc!6,colframe=cpsc,arc=2.5mm,boxrule=0.9pt,left=3mm]{\Large\bfseries\color{cpsc}a)}\quad\begin{minipage}[c]{5.3cm}\raggedright\large $\begin{pmatrix}2\\5\end{pmatrix}$\end{minipage}\end{tcolorbox}\hfill\begin{tcolorbox}[enhanced,nobeforeafter,width=7.1cm,height=2.1cm,valign=center,colback=cpsc!6,colframe=cpsc,arc=2.5mm,boxrule=0.9pt,left=3mm]{\Large\bfseries\color{cpsc}b)}\quad\begin{minipage}[c]{5.3cm}\raggedright\large $\begin{pmatrix}-3\\2\end{pmatrix}$\end{minipage}\end{tcolorbox}\par\vspace{2.5mm}
\begin{tcolorbox}[enhanced,nobeforeafter,width=7.1cm,height=2.1cm,valign=center,colback=cpsc!6,colframe=cpsc,arc=2.5mm,boxrule=0.9pt,left=3mm]{\Large\bfseries\color{cpsc}c)}\quad\begin{minipage}[c]{5.3cm}\raggedright\large $\begin{pmatrix}2\\-3\end{pmatrix}$\end{minipage}\end{tcolorbox}\hfill\begin{tcolorbox}[enhanced,nobeforeafter,width=7.1cm,height=2.1cm,valign=center,colback=cpsc!6,colframe=cpsc,arc=2.5mm,boxrule=0.9pt,left=3mm]{\Large\bfseries\color{cpsc}d)}\quad\begin{minipage}[c]{5.3cm}\raggedright\large $\begin{pmatrix}3\\2\end{pmatrix}$\end{minipage}\end{tcolorbox}
\end{tcolorbox}
\newpage
\begin{tcolorbox}[enhanced,width=16cm,height=14cm,colback=white,colframe=cpsc,boxrule=2pt,arc=0mm,
left=6mm,right=6mm,top=0mm,bottom=4mm,sharp corners,
overlay={\fill[cpsc] (frame.north west) rectangle ([yshift=-1.25cm]frame.north east);
\node[anchor=west,text=white,font=\sffamily\bfseries\LARGE] at ([xshift=5mm,yshift=-0.62cm]frame.north west) {Question 74};
\node[anchor=east,text=white,font=\sffamily\itshape\large] at ([xshift=-5mm,yshift=-0.62cm]frame.north east) {Produit scalaire et géométrie repérée};}]
\vspace*{1.75cm}
\begin{minipage}[t][6.9cm][c]{\linewidth}\LARGE\centering
\begin{minipage}{0.95\linewidth}\raggedright\linespread{1.25}\selectfont Lequel de ces vecteurs est un vecteur normal à la droite d'équation $x+4y-2=0$~?\end{minipage}
\end{minipage}

\vfill
\begin{tcolorbox}[enhanced,nobeforeafter,width=7.1cm,height=2.1cm,valign=center,colback=cpsc!6,colframe=cpsc,arc=2.5mm,boxrule=0.9pt,left=3mm]{\Large\bfseries\color{cpsc}a)}\quad\begin{minipage}[c]{5.3cm}\raggedright\large $\begin{pmatrix}4\\1\end{pmatrix}$\end{minipage}\end{tcolorbox}\hfill\begin{tcolorbox}[enhanced,nobeforeafter,width=7.1cm,height=2.1cm,valign=center,colback=cpsc!6,colframe=cpsc,arc=2.5mm,boxrule=0.9pt,left=3mm]{\Large\bfseries\color{cpsc}b)}\quad\begin{minipage}[c]{5.3cm}\raggedright\large $\begin{pmatrix}1\\4\end{pmatrix}$\end{minipage}\end{tcolorbox}\par\vspace{2.5mm}
\begin{tcolorbox}[enhanced,nobeforeafter,width=7.1cm,height=2.1cm,valign=center,colback=cpsc!6,colframe=cpsc,arc=2.5mm,boxrule=0.9pt,left=3mm]{\Large\bfseries\color{cpsc}c)}\quad\begin{minipage}[c]{5.3cm}\raggedright\large $\begin{pmatrix}1\\-2\end{pmatrix}$\end{minipage}\end{tcolorbox}\hfill\begin{tcolorbox}[enhanced,nobeforeafter,width=7.1cm,height=2.1cm,valign=center,colback=cpsc!6,colframe=cpsc,arc=2.5mm,boxrule=0.9pt,left=3mm]{\Large\bfseries\color{cpsc}d)}\quad\begin{minipage}[c]{5.3cm}\raggedright\large $\begin{pmatrix}-4\\1\end{pmatrix}$\end{minipage}\end{tcolorbox}
\end{tcolorbox}
\newpage
\begin{tcolorbox}[enhanced,width=16cm,height=14cm,colback=white,colframe=cpsc,boxrule=2pt,arc=0mm,
left=6mm,right=6mm,top=0mm,bottom=4mm,sharp corners,
overlay={\fill[cpsc] (frame.north west) rectangle ([yshift=-1.25cm]frame.north east);
\node[anchor=west,text=white,font=\sffamily\bfseries\LARGE] at ([xshift=5mm,yshift=-0.62cm]frame.north west) {Question 75};
\node[anchor=east,text=white,font=\sffamily\itshape\large] at ([xshift=-5mm,yshift=-0.62cm]frame.north east) {Produit scalaire et géométrie repérée};}]
\vspace*{1.75cm}
\begin{minipage}[t][6.9cm][c]{\linewidth}\LARGE\centering
\begin{minipage}{0.95\linewidth}\raggedright\linespread{1.25}\selectfont Quelle est une équation cartésienne de la droite passant par $A(1\,;\,3)$ et de vecteur normal $\vec n\begin{pmatrix}2\\-1\end{pmatrix}$~?\end{minipage}
\end{minipage}

\vfill
\begin{tcolorbox}[enhanced,nobeforeafter,width=7.1cm,height=2.1cm,valign=center,colback=cpsc!6,colframe=cpsc,arc=2.5mm,boxrule=0.9pt,left=3mm]{\Large\bfseries\color{cpsc}a)}\quad\begin{minipage}[c]{5.3cm}\raggedright\Large $2x-y+1=0$\end{minipage}\end{tcolorbox}\hfill\begin{tcolorbox}[enhanced,nobeforeafter,width=7.1cm,height=2.1cm,valign=center,colback=cpsc!6,colframe=cpsc,arc=2.5mm,boxrule=0.9pt,left=3mm]{\Large\bfseries\color{cpsc}b)}\quad\begin{minipage}[c]{5.3cm}\raggedright\Large $2x-y-1=0$\end{minipage}\end{tcolorbox}\par\vspace{2.5mm}
\begin{tcolorbox}[enhanced,nobeforeafter,width=7.1cm,height=2.1cm,valign=center,colback=cpsc!6,colframe=cpsc,arc=2.5mm,boxrule=0.9pt,left=3mm]{\Large\bfseries\color{cpsc}c)}\quad\begin{minipage}[c]{5.3cm}\raggedright\Large $x+2y-7=0$\end{minipage}\end{tcolorbox}\hfill\begin{tcolorbox}[enhanced,nobeforeafter,width=7.1cm,height=2.1cm,valign=center,colback=cpsc!6,colframe=cpsc,arc=2.5mm,boxrule=0.9pt,left=3mm]{\Large\bfseries\color{cpsc}d)}\quad\begin{minipage}[c]{5.3cm}\raggedright\Large $x+3y+1=0$\end{minipage}\end{tcolorbox}
\end{tcolorbox}
\newpage
\begin{tcolorbox}[enhanced,width=16cm,height=14cm,colback=white,colframe=cpsc,boxrule=2pt,arc=0mm,
left=6mm,right=6mm,top=0mm,bottom=4mm,sharp corners,
overlay={\fill[cpsc] (frame.north west) rectangle ([yshift=-1.25cm]frame.north east);
\node[anchor=west,text=white,font=\sffamily\bfseries\LARGE] at ([xshift=5mm,yshift=-0.62cm]frame.north west) {Question 76};
\node[anchor=east,text=white,font=\sffamily\itshape\large] at ([xshift=-5mm,yshift=-0.62cm]frame.north east) {Produit scalaire et géométrie repérée};}]
\vspace*{1.75cm}
\begin{minipage}[t][6.9cm][c]{\linewidth}\LARGE\centering
\begin{minipage}{0.95\linewidth}\raggedright\linespread{1.25}\selectfont Quelle est une équation cartésienne de la droite passant par $A(-1\,;\,2)$ et de vecteur normal $\vec n\begin{pmatrix}3\\4\end{pmatrix}$~?\end{minipage}
\end{minipage}

\vfill
\begin{tcolorbox}[enhanced,nobeforeafter,width=7.1cm,height=2.1cm,valign=center,colback=cpsc!6,colframe=cpsc,arc=2.5mm,boxrule=0.9pt,left=3mm]{\Large\bfseries\color{cpsc}a)}\quad\begin{minipage}[c]{5.3cm}\raggedright\Large $-4x+3y-10=0$\end{minipage}\end{tcolorbox}\hfill\begin{tcolorbox}[enhanced,nobeforeafter,width=7.1cm,height=2.1cm,valign=center,colback=cpsc!6,colframe=cpsc,arc=2.5mm,boxrule=0.9pt,left=3mm]{\Large\bfseries\color{cpsc}b)}\quad\begin{minipage}[c]{5.3cm}\raggedright\Large $3x+4y-5=0$\end{minipage}\end{tcolorbox}\par\vspace{2.5mm}
\begin{tcolorbox}[enhanced,nobeforeafter,width=7.1cm,height=2.1cm,valign=center,colback=cpsc!6,colframe=cpsc,arc=2.5mm,boxrule=0.9pt,left=3mm]{\Large\bfseries\color{cpsc}c)}\quad\begin{minipage}[c]{5.3cm}\raggedright\Large $3x+4y+5=0$\end{minipage}\end{tcolorbox}\hfill\begin{tcolorbox}[enhanced,nobeforeafter,width=7.1cm,height=2.1cm,valign=center,colback=cpsc!6,colframe=cpsc,arc=2.5mm,boxrule=0.9pt,left=3mm]{\Large\bfseries\color{cpsc}d)}\quad\begin{minipage}[c]{5.3cm}\raggedright\Large $-x+2y-5=0$\end{minipage}\end{tcolorbox}
\end{tcolorbox}
\newpage
\begin{tcolorbox}[enhanced,width=16cm,height=14cm,colback=white,colframe=cpsc,boxrule=2pt,arc=0mm,
left=6mm,right=6mm,top=0mm,bottom=4mm,sharp corners,
overlay={\fill[cpsc] (frame.north west) rectangle ([yshift=-1.25cm]frame.north east);
\node[anchor=west,text=white,font=\sffamily\bfseries\LARGE] at ([xshift=5mm,yshift=-0.62cm]frame.north west) {Question 77};
\node[anchor=east,text=white,font=\sffamily\itshape\large] at ([xshift=-5mm,yshift=-0.62cm]frame.north east) {Produit scalaire et géométrie repérée};}]
\vspace*{1.75cm}
\begin{minipage}[t][6.9cm][c]{\linewidth}\LARGE\centering
\begin{minipage}{0.95\linewidth}\raggedright\linespread{1.25}\selectfont Quels sont le centre et le rayon du cercle d'équation $(x-2)^2+(y+1)^2=9$~?\end{minipage}
\end{minipage}

\vfill
\begin{tcolorbox}[enhanced,nobeforeafter,width=7.1cm,height=2.1cm,valign=center,colback=cpsc!6,colframe=cpsc,arc=2.5mm,boxrule=0.9pt,left=3mm]{\Large\bfseries\color{cpsc}a)}\quad\begin{minipage}[c]{5.3cm}\raggedright\large $\Omega(2\,;\,-1)\text{, rayon }9$\end{minipage}\end{tcolorbox}\hfill\begin{tcolorbox}[enhanced,nobeforeafter,width=7.1cm,height=2.1cm,valign=center,colback=cpsc!6,colframe=cpsc,arc=2.5mm,boxrule=0.9pt,left=3mm]{\Large\bfseries\color{cpsc}b)}\quad\begin{minipage}[c]{5.3cm}\raggedright\large $\Omega(-2\,;\,1)\text{, rayon }3$\end{minipage}\end{tcolorbox}\par\vspace{2.5mm}
\begin{tcolorbox}[enhanced,nobeforeafter,width=7.1cm,height=2.1cm,valign=center,colback=cpsc!6,colframe=cpsc,arc=2.5mm,boxrule=0.9pt,left=3mm]{\Large\bfseries\color{cpsc}c)}\quad\begin{minipage}[c]{5.3cm}\raggedright\large $\Omega(2\,;\,-1)\text{, rayon }3$\end{minipage}\end{tcolorbox}\hfill\begin{tcolorbox}[enhanced,nobeforeafter,width=7.1cm,height=2.1cm,valign=center,colback=cpsc!6,colframe=cpsc,arc=2.5mm,boxrule=0.9pt,left=3mm]{\Large\bfseries\color{cpsc}d)}\quad\begin{minipage}[c]{5.3cm}\raggedright\large $\Omega(-1\,;\,2)\text{, rayon }3$\end{minipage}\end{tcolorbox}
\end{tcolorbox}
\newpage
\begin{tcolorbox}[enhanced,width=16cm,height=14cm,colback=white,colframe=cpsc,boxrule=2pt,arc=0mm,
left=6mm,right=6mm,top=0mm,bottom=4mm,sharp corners,
overlay={\fill[cpsc] (frame.north west) rectangle ([yshift=-1.25cm]frame.north east);
\node[anchor=west,text=white,font=\sffamily\bfseries\LARGE] at ([xshift=5mm,yshift=-0.62cm]frame.north west) {Question 78};
\node[anchor=east,text=white,font=\sffamily\itshape\large] at ([xshift=-5mm,yshift=-0.62cm]frame.north east) {Produit scalaire et géométrie repérée};}]
\vspace*{1.75cm}
\begin{minipage}[t][6.9cm][c]{\linewidth}\LARGE\centering
\begin{minipage}{0.95\linewidth}\raggedright\linespread{1.25}\selectfont Quels sont le centre et le rayon du cercle d'équation $(x+3)^2+(y-4)^2=25$~?\end{minipage}
\end{minipage}

\vfill
\begin{tcolorbox}[enhanced,nobeforeafter,width=7.1cm,height=2.1cm,valign=center,colback=cpsc!6,colframe=cpsc,arc=2.5mm,boxrule=0.9pt,left=3mm]{\Large\bfseries\color{cpsc}a)}\quad\begin{minipage}[c]{5.3cm}\raggedright\large $\Omega(3\,;\,-4)\text{, rayon }5$\end{minipage}\end{tcolorbox}\hfill\begin{tcolorbox}[enhanced,nobeforeafter,width=7.1cm,height=2.1cm,valign=center,colback=cpsc!6,colframe=cpsc,arc=2.5mm,boxrule=0.9pt,left=3mm]{\Large\bfseries\color{cpsc}b)}\quad\begin{minipage}[c]{5.3cm}\raggedright\large $\Omega(4\,;\,-3)\text{, rayon }5$\end{minipage}\end{tcolorbox}\par\vspace{2.5mm}
\begin{tcolorbox}[enhanced,nobeforeafter,width=7.1cm,height=2.1cm,valign=center,colback=cpsc!6,colframe=cpsc,arc=2.5mm,boxrule=0.9pt,left=3mm]{\Large\bfseries\color{cpsc}c)}\quad\begin{minipage}[c]{5.3cm}\raggedright\large $\Omega(-3\,;\,4)\text{, rayon }5$\end{minipage}\end{tcolorbox}\hfill\begin{tcolorbox}[enhanced,nobeforeafter,width=7.1cm,height=2.1cm,valign=center,colback=cpsc!6,colframe=cpsc,arc=2.5mm,boxrule=0.9pt,left=3mm]{\Large\bfseries\color{cpsc}d)}\quad\begin{minipage}[c]{5.3cm}\raggedright\large $\Omega(-3\,;\,4)\text{, rayon }25$\end{minipage}\end{tcolorbox}
\end{tcolorbox}
\newpage
\begin{tcolorbox}[enhanced,width=16cm,height=14cm,colback=white,colframe=cpsc,boxrule=2pt,arc=0mm,
left=6mm,right=6mm,top=0mm,bottom=4mm,sharp corners,
overlay={\fill[cpsc] (frame.north west) rectangle ([yshift=-1.25cm]frame.north east);
\node[anchor=west,text=white,font=\sffamily\bfseries\LARGE] at ([xshift=5mm,yshift=-0.62cm]frame.north west) {Question 79};
\node[anchor=east,text=white,font=\sffamily\itshape\large] at ([xshift=-5mm,yshift=-0.62cm]frame.north east) {Produit scalaire et géométrie repérée};}]
\vspace*{1.75cm}
\begin{minipage}[t][6.9cm][c]{\linewidth}\LARGE\centering
\begin{minipage}{0.95\linewidth}\raggedright\linespread{1.25}\selectfont On a $\|\vec u\|=3$, $\|\vec v\|=4$ et $\vec u\cdot\vec v=2$. Que vaut $\|\vec u+\vec v\|^2$~?\end{minipage}
\end{minipage}

\vfill
\begin{tcolorbox}[enhanced,nobeforeafter,width=7.1cm,height=2.1cm,valign=center,colback=cpsc!6,colframe=cpsc,arc=2.5mm,boxrule=0.9pt,left=3mm]{\Large\bfseries\color{cpsc}a)}\quad\begin{minipage}[c]{5.3cm}\raggedright\Large $25$\end{minipage}\end{tcolorbox}\hfill\begin{tcolorbox}[enhanced,nobeforeafter,width=7.1cm,height=2.1cm,valign=center,colback=cpsc!6,colframe=cpsc,arc=2.5mm,boxrule=0.9pt,left=3mm]{\Large\bfseries\color{cpsc}b)}\quad\begin{minipage}[c]{5.3cm}\raggedright\Large $29$\end{minipage}\end{tcolorbox}\par\vspace{2.5mm}
\begin{tcolorbox}[enhanced,nobeforeafter,width=7.1cm,height=2.1cm,valign=center,colback=cpsc!6,colframe=cpsc,arc=2.5mm,boxrule=0.9pt,left=3mm]{\Large\bfseries\color{cpsc}c)}\quad\begin{minipage}[c]{5.3cm}\raggedright\Large $49$\end{minipage}\end{tcolorbox}\hfill\begin{tcolorbox}[enhanced,nobeforeafter,width=7.1cm,height=2.1cm,valign=center,colback=cpsc!6,colframe=cpsc,arc=2.5mm,boxrule=0.9pt,left=3mm]{\Large\bfseries\color{cpsc}d)}\quad\begin{minipage}[c]{5.3cm}\raggedright\Large $27$\end{minipage}\end{tcolorbox}
\end{tcolorbox}
\newpage
\begin{tcolorbox}[enhanced,width=16cm,height=14cm,colback=white,colframe=cpsc,boxrule=2pt,arc=0mm,
left=6mm,right=6mm,top=0mm,bottom=4mm,sharp corners,
overlay={\fill[cpsc] (frame.north west) rectangle ([yshift=-1.25cm]frame.north east);
\node[anchor=west,text=white,font=\sffamily\bfseries\LARGE] at ([xshift=5mm,yshift=-0.62cm]frame.north west) {Question 80};
\node[anchor=east,text=white,font=\sffamily\itshape\large] at ([xshift=-5mm,yshift=-0.62cm]frame.north east) {Produit scalaire et géométrie repérée};}]
\vspace*{1.75cm}
\begin{minipage}[t][6.9cm][c]{\linewidth}\LARGE\centering
\begin{minipage}{0.95\linewidth}\raggedright\linespread{1.25}\selectfont On a $\|\vec u\|=2$, $\|\vec v\|=5$ et $\vec u\cdot\vec v=-3$. Que vaut $\|\vec u+\vec v\|^2$~?\end{minipage}
\end{minipage}

\vfill
\begin{tcolorbox}[enhanced,nobeforeafter,width=7.1cm,height=2.1cm,valign=center,colback=cpsc!6,colframe=cpsc,arc=2.5mm,boxrule=0.9pt,left=3mm]{\Large\bfseries\color{cpsc}a)}\quad\begin{minipage}[c]{5.3cm}\raggedright\Large $49$\end{minipage}\end{tcolorbox}\hfill\begin{tcolorbox}[enhanced,nobeforeafter,width=7.1cm,height=2.1cm,valign=center,colback=cpsc!6,colframe=cpsc,arc=2.5mm,boxrule=0.9pt,left=3mm]{\Large\bfseries\color{cpsc}b)}\quad\begin{minipage}[c]{5.3cm}\raggedright\Large $26$\end{minipage}\end{tcolorbox}\par\vspace{2.5mm}
\begin{tcolorbox}[enhanced,nobeforeafter,width=7.1cm,height=2.1cm,valign=center,colback=cpsc!6,colframe=cpsc,arc=2.5mm,boxrule=0.9pt,left=3mm]{\Large\bfseries\color{cpsc}c)}\quad\begin{minipage}[c]{5.3cm}\raggedright\Large $23$\end{minipage}\end{tcolorbox}\hfill\begin{tcolorbox}[enhanced,nobeforeafter,width=7.1cm,height=2.1cm,valign=center,colback=cpsc!6,colframe=cpsc,arc=2.5mm,boxrule=0.9pt,left=3mm]{\Large\bfseries\color{cpsc}d)}\quad\begin{minipage}[c]{5.3cm}\raggedright\Large $29$\end{minipage}\end{tcolorbox}
\end{tcolorbox}
\newpage
\begin{tcolorbox}[enhanced,width=16cm,height=14cm,colback=white,colframe=ccvp,boxrule=2pt,arc=0mm,
left=6mm,right=6mm,top=0mm,bottom=4mm,sharp corners,
overlay={\fill[ccvp] (frame.north west) rectangle ([yshift=-1.25cm]frame.north east);
\node[anchor=west,text=white,font=\sffamily\bfseries\LARGE] at ([xshift=5mm,yshift=-0.62cm]frame.north west) {Question 81};
\node[anchor=east,text=white,font=\sffamily\itshape\large] at ([xshift=-5mm,yshift=-0.62cm]frame.north east) {Probabilités conditionnelles et variables aléatoires};}]
\vspace*{1.75cm}
\begin{minipage}[t][6.9cm][c]{\linewidth}\LARGE\centering
\begin{minipage}{0.95\linewidth}\raggedright\linespread{1.25}\selectfont On a $P(A)=0{,}5$ et $P(A\cap B)=0{,}2$. Que vaut $P_A(B)$~?\end{minipage}
\end{minipage}

\vfill
\begin{tcolorbox}[enhanced,nobeforeafter,width=7.1cm,height=2.1cm,valign=center,colback=ccvp!6,colframe=ccvp,arc=2.5mm,boxrule=0.9pt,left=3mm]{\Large\bfseries\color{ccvp}a)}\quad\begin{minipage}[c]{5.3cm}\raggedright\Large $0{,}7$\end{minipage}\end{tcolorbox}\hfill\begin{tcolorbox}[enhanced,nobeforeafter,width=7.1cm,height=2.1cm,valign=center,colback=ccvp!6,colframe=ccvp,arc=2.5mm,boxrule=0.9pt,left=3mm]{\Large\bfseries\color{ccvp}b)}\quad\begin{minipage}[c]{5.3cm}\raggedright\Large $0{,}1$\end{minipage}\end{tcolorbox}\par\vspace{2.5mm}
\begin{tcolorbox}[enhanced,nobeforeafter,width=7.1cm,height=2.1cm,valign=center,colback=ccvp!6,colframe=ccvp,arc=2.5mm,boxrule=0.9pt,left=3mm]{\Large\bfseries\color{ccvp}c)}\quad\begin{minipage}[c]{5.3cm}\raggedright\Large $0{,}3$\end{minipage}\end{tcolorbox}\hfill\begin{tcolorbox}[enhanced,nobeforeafter,width=7.1cm,height=2.1cm,valign=center,colback=ccvp!6,colframe=ccvp,arc=2.5mm,boxrule=0.9pt,left=3mm]{\Large\bfseries\color{ccvp}d)}\quad\begin{minipage}[c]{5.3cm}\raggedright\Large $0{,}4$\end{minipage}\end{tcolorbox}
\end{tcolorbox}
\newpage
\begin{tcolorbox}[enhanced,width=16cm,height=14cm,colback=white,colframe=ccvp,boxrule=2pt,arc=0mm,
left=6mm,right=6mm,top=0mm,bottom=4mm,sharp corners,
overlay={\fill[ccvp] (frame.north west) rectangle ([yshift=-1.25cm]frame.north east);
\node[anchor=west,text=white,font=\sffamily\bfseries\LARGE] at ([xshift=5mm,yshift=-0.62cm]frame.north west) {Question 82};
\node[anchor=east,text=white,font=\sffamily\itshape\large] at ([xshift=-5mm,yshift=-0.62cm]frame.north east) {Probabilités conditionnelles et variables aléatoires};}]
\vspace*{1.75cm}
\begin{minipage}[t][6.9cm][c]{\linewidth}\LARGE\centering
\begin{minipage}{0.95\linewidth}\raggedright\linespread{1.25}\selectfont On a $P(A)=0{,}4$ et $P(A\cap B)=0{,}1$. Que vaut $P_A(B)$~?\end{minipage}
\end{minipage}

\vfill
\begin{tcolorbox}[enhanced,nobeforeafter,width=7.1cm,height=2.1cm,valign=center,colback=ccvp!6,colframe=ccvp,arc=2.5mm,boxrule=0.9pt,left=3mm]{\Large\bfseries\color{ccvp}a)}\quad\begin{minipage}[c]{5.3cm}\raggedright\Large $0{,}25$\end{minipage}\end{tcolorbox}\hfill\begin{tcolorbox}[enhanced,nobeforeafter,width=7.1cm,height=2.1cm,valign=center,colback=ccvp!6,colframe=ccvp,arc=2.5mm,boxrule=0.9pt,left=3mm]{\Large\bfseries\color{ccvp}b)}\quad\begin{minipage}[c]{5.3cm}\raggedright\Large $0{,}5$\end{minipage}\end{tcolorbox}\par\vspace{2.5mm}
\begin{tcolorbox}[enhanced,nobeforeafter,width=7.1cm,height=2.1cm,valign=center,colback=ccvp!6,colframe=ccvp,arc=2.5mm,boxrule=0.9pt,left=3mm]{\Large\bfseries\color{ccvp}c)}\quad\begin{minipage}[c]{5.3cm}\raggedright\Large $0{,}04$\end{minipage}\end{tcolorbox}\hfill\begin{tcolorbox}[enhanced,nobeforeafter,width=7.1cm,height=2.1cm,valign=center,colback=ccvp!6,colframe=ccvp,arc=2.5mm,boxrule=0.9pt,left=3mm]{\Large\bfseries\color{ccvp}d)}\quad\begin{minipage}[c]{5.3cm}\raggedright\Large $0{,}3$\end{minipage}\end{tcolorbox}
\end{tcolorbox}
\newpage
\begin{tcolorbox}[enhanced,width=16cm,height=14cm,colback=white,colframe=ccvp,boxrule=2pt,arc=0mm,
left=6mm,right=6mm,top=0mm,bottom=4mm,sharp corners,
overlay={\fill[ccvp] (frame.north west) rectangle ([yshift=-1.25cm]frame.north east);
\node[anchor=west,text=white,font=\sffamily\bfseries\LARGE] at ([xshift=5mm,yshift=-0.62cm]frame.north west) {Question 83};
\node[anchor=east,text=white,font=\sffamily\itshape\large] at ([xshift=-5mm,yshift=-0.62cm]frame.north east) {Probabilités conditionnelles et variables aléatoires};}]
\vspace*{1.75cm}
\begin{minipage}[t][6.9cm][c]{\linewidth}\LARGE\centering
\begin{minipage}{0.95\linewidth}\raggedright\linespread{1.25}\selectfont On a $P(A)=0{,}6$ et $P_A(B)=0{,}5$. Que vaut $P(A\cap B)$~?\end{minipage}
\end{minipage}

\vfill
\begin{tcolorbox}[enhanced,nobeforeafter,width=7.1cm,height=2.1cm,valign=center,colback=ccvp!6,colframe=ccvp,arc=2.5mm,boxrule=0.9pt,left=3mm]{\Large\bfseries\color{ccvp}a)}\quad\begin{minipage}[c]{5.3cm}\raggedright\Large $0{,}3$\end{minipage}\end{tcolorbox}\hfill\begin{tcolorbox}[enhanced,nobeforeafter,width=7.1cm,height=2.1cm,valign=center,colback=ccvp!6,colframe=ccvp,arc=2.5mm,boxrule=0.9pt,left=3mm]{\Large\bfseries\color{ccvp}b)}\quad\begin{minipage}[c]{5.3cm}\raggedright\Large $0{,}1$\end{minipage}\end{tcolorbox}\par\vspace{2.5mm}
\begin{tcolorbox}[enhanced,nobeforeafter,width=7.1cm,height=2.1cm,valign=center,colback=ccvp!6,colframe=ccvp,arc=2.5mm,boxrule=0.9pt,left=3mm]{\Large\bfseries\color{ccvp}c)}\quad\begin{minipage}[c]{5.3cm}\raggedright\Large $0{,}5$\end{minipage}\end{tcolorbox}\hfill\begin{tcolorbox}[enhanced,nobeforeafter,width=7.1cm,height=2.1cm,valign=center,colback=ccvp!6,colframe=ccvp,arc=2.5mm,boxrule=0.9pt,left=3mm]{\Large\bfseries\color{ccvp}d)}\quad\begin{minipage}[c]{5.3cm}\raggedright\Large $\dfrac{5}{6}$\end{minipage}\end{tcolorbox}
\end{tcolorbox}
\newpage
\begin{tcolorbox}[enhanced,width=16cm,height=14cm,colback=white,colframe=ccvp,boxrule=2pt,arc=0mm,
left=6mm,right=6mm,top=0mm,bottom=4mm,sharp corners,
overlay={\fill[ccvp] (frame.north west) rectangle ([yshift=-1.25cm]frame.north east);
\node[anchor=west,text=white,font=\sffamily\bfseries\LARGE] at ([xshift=5mm,yshift=-0.62cm]frame.north west) {Question 84};
\node[anchor=east,text=white,font=\sffamily\itshape\large] at ([xshift=-5mm,yshift=-0.62cm]frame.north east) {Probabilités conditionnelles et variables aléatoires};}]
\vspace*{1.75cm}
\begin{minipage}[t][6.9cm][c]{\linewidth}\LARGE\centering
\begin{minipage}{0.95\linewidth}\raggedright\linespread{1.25}\selectfont On a $P(A)=0{,}3$ et $P_A(B)=0{,}4$. Que vaut $P(A\cap B)$~?\end{minipage}
\end{minipage}

\vfill
\begin{tcolorbox}[enhanced,nobeforeafter,width=7.1cm,height=2.1cm,valign=center,colback=ccvp!6,colframe=ccvp,arc=2.5mm,boxrule=0.9pt,left=3mm]{\Large\bfseries\color{ccvp}a)}\quad\begin{minipage}[c]{5.3cm}\raggedright\Large $0{,}18$\end{minipage}\end{tcolorbox}\hfill\begin{tcolorbox}[enhanced,nobeforeafter,width=7.1cm,height=2.1cm,valign=center,colback=ccvp!6,colframe=ccvp,arc=2.5mm,boxrule=0.9pt,left=3mm]{\Large\bfseries\color{ccvp}b)}\quad\begin{minipage}[c]{5.3cm}\raggedright\Large $0{,}7$\end{minipage}\end{tcolorbox}\par\vspace{2.5mm}
\begin{tcolorbox}[enhanced,nobeforeafter,width=7.1cm,height=2.1cm,valign=center,colback=ccvp!6,colframe=ccvp,arc=2.5mm,boxrule=0.9pt,left=3mm]{\Large\bfseries\color{ccvp}c)}\quad\begin{minipage}[c]{5.3cm}\raggedright\Large $0{,}12$\end{minipage}\end{tcolorbox}\hfill\begin{tcolorbox}[enhanced,nobeforeafter,width=7.1cm,height=2.1cm,valign=center,colback=ccvp!6,colframe=ccvp,arc=2.5mm,boxrule=0.9pt,left=3mm]{\Large\bfseries\color{ccvp}d)}\quad\begin{minipage}[c]{5.3cm}\raggedright\Large $0{,}88$\end{minipage}\end{tcolorbox}
\end{tcolorbox}
\newpage
\begin{tcolorbox}[enhanced,width=16cm,height=14cm,colback=white,colframe=ccvp,boxrule=2pt,arc=0mm,
left=6mm,right=6mm,top=0mm,bottom=4mm,sharp corners,
overlay={\fill[ccvp] (frame.north west) rectangle ([yshift=-1.25cm]frame.north east);
\node[anchor=west,text=white,font=\sffamily\bfseries\LARGE] at ([xshift=5mm,yshift=-0.62cm]frame.north west) {Question 85};
\node[anchor=east,text=white,font=\sffamily\itshape\large] at ([xshift=-5mm,yshift=-0.62cm]frame.north east) {Probabilités conditionnelles et variables aléatoires};}]
\vspace*{1.75cm}
\begin{minipage}[t][6.9cm][c]{\linewidth}\LARGE\centering
\begin{minipage}{0.95\linewidth}\raggedright\linespread{1.25}\selectfont On a $P(A)=0{,}3$, $P_A(B)=0{,}5$ et $P_{\overline{A}}(B)=0{,}2$. Que vaut $P(B)$~?\end{minipage}
\end{minipage}

\vfill
\begin{tcolorbox}[enhanced,nobeforeafter,width=7.1cm,height=2.1cm,valign=center,colback=ccvp!6,colframe=ccvp,arc=2.5mm,boxrule=0.9pt,left=3mm]{\Large\bfseries\color{ccvp}a)}\quad\begin{minipage}[c]{5.3cm}\raggedright\Large $0{,}15$\end{minipage}\end{tcolorbox}\hfill\begin{tcolorbox}[enhanced,nobeforeafter,width=7.1cm,height=2.1cm,valign=center,colback=ccvp!6,colframe=ccvp,arc=2.5mm,boxrule=0.9pt,left=3mm]{\Large\bfseries\color{ccvp}b)}\quad\begin{minipage}[c]{5.3cm}\raggedright\Large $0{,}21$\end{minipage}\end{tcolorbox}\par\vspace{2.5mm}
\begin{tcolorbox}[enhanced,nobeforeafter,width=7.1cm,height=2.1cm,valign=center,colback=ccvp!6,colframe=ccvp,arc=2.5mm,boxrule=0.9pt,left=3mm]{\Large\bfseries\color{ccvp}c)}\quad\begin{minipage}[c]{5.3cm}\raggedright\Large $0{,}7$\end{minipage}\end{tcolorbox}\hfill\begin{tcolorbox}[enhanced,nobeforeafter,width=7.1cm,height=2.1cm,valign=center,colback=ccvp!6,colframe=ccvp,arc=2.5mm,boxrule=0.9pt,left=3mm]{\Large\bfseries\color{ccvp}d)}\quad\begin{minipage}[c]{5.3cm}\raggedright\Large $0{,}29$\end{minipage}\end{tcolorbox}
\end{tcolorbox}
\newpage
\begin{tcolorbox}[enhanced,width=16cm,height=14cm,colback=white,colframe=ccvp,boxrule=2pt,arc=0mm,
left=6mm,right=6mm,top=0mm,bottom=4mm,sharp corners,
overlay={\fill[ccvp] (frame.north west) rectangle ([yshift=-1.25cm]frame.north east);
\node[anchor=west,text=white,font=\sffamily\bfseries\LARGE] at ([xshift=5mm,yshift=-0.62cm]frame.north west) {Question 86};
\node[anchor=east,text=white,font=\sffamily\itshape\large] at ([xshift=-5mm,yshift=-0.62cm]frame.north east) {Probabilités conditionnelles et variables aléatoires};}]
\vspace*{1.75cm}
\begin{minipage}[t][6.9cm][c]{\linewidth}\LARGE\centering
\begin{minipage}{0.95\linewidth}\raggedright\linespread{1.25}\selectfont On a $P(A)=0{,}6$, $P_A(B)=0{,}4$ et $P_{\overline{A}}(B)=0{,}1$. Que vaut $P(B)$~?\end{minipage}
\end{minipage}

\vfill
\begin{tcolorbox}[enhanced,nobeforeafter,width=7.1cm,height=2.1cm,valign=center,colback=ccvp!6,colframe=ccvp,arc=2.5mm,boxrule=0.9pt,left=3mm]{\Large\bfseries\color{ccvp}a)}\quad\begin{minipage}[c]{5.3cm}\raggedright\Large $0{,}3$\end{minipage}\end{tcolorbox}\hfill\begin{tcolorbox}[enhanced,nobeforeafter,width=7.1cm,height=2.1cm,valign=center,colback=ccvp!6,colframe=ccvp,arc=2.5mm,boxrule=0.9pt,left=3mm]{\Large\bfseries\color{ccvp}b)}\quad\begin{minipage}[c]{5.3cm}\raggedright\Large $0{,}28$\end{minipage}\end{tcolorbox}\par\vspace{2.5mm}
\begin{tcolorbox}[enhanced,nobeforeafter,width=7.1cm,height=2.1cm,valign=center,colback=ccvp!6,colframe=ccvp,arc=2.5mm,boxrule=0.9pt,left=3mm]{\Large\bfseries\color{ccvp}c)}\quad\begin{minipage}[c]{5.3cm}\raggedright\Large $0{,}24$\end{minipage}\end{tcolorbox}\hfill\begin{tcolorbox}[enhanced,nobeforeafter,width=7.1cm,height=2.1cm,valign=center,colback=ccvp!6,colframe=ccvp,arc=2.5mm,boxrule=0.9pt,left=3mm]{\Large\bfseries\color{ccvp}d)}\quad\begin{minipage}[c]{5.3cm}\raggedright\Large $0{,}5$\end{minipage}\end{tcolorbox}
\end{tcolorbox}
\newpage
\begin{tcolorbox}[enhanced,width=16cm,height=14cm,colback=white,colframe=ccvp,boxrule=2pt,arc=0mm,
left=6mm,right=6mm,top=0mm,bottom=4mm,sharp corners,
overlay={\fill[ccvp] (frame.north west) rectangle ([yshift=-1.25cm]frame.north east);
\node[anchor=west,text=white,font=\sffamily\bfseries\LARGE] at ([xshift=5mm,yshift=-0.62cm]frame.north west) {Question 87};
\node[anchor=east,text=white,font=\sffamily\itshape\large] at ([xshift=-5mm,yshift=-0.62cm]frame.north east) {Probabilités conditionnelles et variables aléatoires};}]
\vspace*{1.75cm}
\begin{minipage}[t][6.9cm][c]{\linewidth}\LARGE\centering
\begin{minipage}{0.95\linewidth}\raggedright\linespread{1.25}\selectfont On a $P(A)=0{,}5$, $P(B)=0{,}4$ et $P(A\cap B)=0{,}2$. Les événements $A$ et $B$ sont-ils indépendants~?\end{minipage}
\end{minipage}

\vfill
\begin{tcolorbox}[enhanced,nobeforeafter,width=7.1cm,height=2.1cm,valign=center,colback=ccvp!6,colframe=ccvp,arc=2.5mm,boxrule=0.9pt,left=3mm]{\Large\bfseries\color{ccvp}a)}\quad\begin{minipage}[c]{5.3cm}\raggedright\Large oui\end{minipage}\end{tcolorbox}\hfill\begin{tcolorbox}[enhanced,nobeforeafter,width=7.1cm,height=2.1cm,valign=center,colback=ccvp!6,colframe=ccvp,arc=2.5mm,boxrule=0.9pt,left=3mm]{\Large\bfseries\color{ccvp}b)}\quad\begin{minipage}[c]{5.3cm}\raggedright\large seulement s'ils sont incompatibles\end{minipage}\end{tcolorbox}\par\vspace{2.5mm}
\begin{tcolorbox}[enhanced,nobeforeafter,width=7.1cm,height=2.1cm,valign=center,colback=ccvp!6,colframe=ccvp,arc=2.5mm,boxrule=0.9pt,left=3mm]{\Large\bfseries\color{ccvp}c)}\quad\begin{minipage}[c]{5.3cm}\raggedright\Large on ne peut pas savoir\end{minipage}\end{tcolorbox}\hfill\begin{tcolorbox}[enhanced,nobeforeafter,width=7.1cm,height=2.1cm,valign=center,colback=ccvp!6,colframe=ccvp,arc=2.5mm,boxrule=0.9pt,left=3mm]{\Large\bfseries\color{ccvp}d)}\quad\begin{minipage}[c]{5.3cm}\raggedright\Large non\end{minipage}\end{tcolorbox}
\end{tcolorbox}
\newpage
\begin{tcolorbox}[enhanced,width=16cm,height=14cm,colback=white,colframe=ccvp,boxrule=2pt,arc=0mm,
left=6mm,right=6mm,top=0mm,bottom=4mm,sharp corners,
overlay={\fill[ccvp] (frame.north west) rectangle ([yshift=-1.25cm]frame.north east);
\node[anchor=west,text=white,font=\sffamily\bfseries\LARGE] at ([xshift=5mm,yshift=-0.62cm]frame.north west) {Question 88};
\node[anchor=east,text=white,font=\sffamily\itshape\large] at ([xshift=-5mm,yshift=-0.62cm]frame.north east) {Probabilités conditionnelles et variables aléatoires};}]
\vspace*{1.75cm}
\begin{minipage}[t][6.9cm][c]{\linewidth}\LARGE\centering
\begin{minipage}{0.95\linewidth}\raggedright\linespread{1.25}\selectfont On a $P(A)=0{,}3$, $P(B)=0{,}6$ et $P(A\cap B)=0{,}2$. Les événements $A$ et $B$ sont-ils indépendants~?\end{minipage}
\end{minipage}

\vfill
\begin{tcolorbox}[enhanced,nobeforeafter,width=7.1cm,height=2.1cm,valign=center,colback=ccvp!6,colframe=ccvp,arc=2.5mm,boxrule=0.9pt,left=3mm]{\Large\bfseries\color{ccvp}a)}\quad\begin{minipage}[c]{5.3cm}\raggedright\large seulement s'ils sont incompatibles\end{minipage}\end{tcolorbox}\hfill\begin{tcolorbox}[enhanced,nobeforeafter,width=7.1cm,height=2.1cm,valign=center,colback=ccvp!6,colframe=ccvp,arc=2.5mm,boxrule=0.9pt,left=3mm]{\Large\bfseries\color{ccvp}b)}\quad\begin{minipage}[c]{5.3cm}\raggedright\Large oui\end{minipage}\end{tcolorbox}\par\vspace{2.5mm}
\begin{tcolorbox}[enhanced,nobeforeafter,width=7.1cm,height=2.1cm,valign=center,colback=ccvp!6,colframe=ccvp,arc=2.5mm,boxrule=0.9pt,left=3mm]{\Large\bfseries\color{ccvp}c)}\quad\begin{minipage}[c]{5.3cm}\raggedright\Large non\end{minipage}\end{tcolorbox}\hfill\begin{tcolorbox}[enhanced,nobeforeafter,width=7.1cm,height=2.1cm,valign=center,colback=ccvp!6,colframe=ccvp,arc=2.5mm,boxrule=0.9pt,left=3mm]{\Large\bfseries\color{ccvp}d)}\quad\begin{minipage}[c]{5.3cm}\raggedright\Large on ne peut pas savoir\end{minipage}\end{tcolorbox}
\end{tcolorbox}
\newpage
\begin{tcolorbox}[enhanced,width=16cm,height=14cm,colback=white,colframe=ccvp,boxrule=2pt,arc=0mm,
left=6mm,right=6mm,top=0mm,bottom=4mm,sharp corners,
overlay={\fill[ccvp] (frame.north west) rectangle ([yshift=-1.25cm]frame.north east);
\node[anchor=west,text=white,font=\sffamily\bfseries\LARGE] at ([xshift=5mm,yshift=-0.62cm]frame.north west) {Question 89};
\node[anchor=east,text=white,font=\sffamily\itshape\large] at ([xshift=-5mm,yshift=-0.62cm]frame.north east) {Probabilités conditionnelles et variables aléatoires};}]
\vspace*{1.75cm}
\begin{minipage}[t][6.9cm][c]{\linewidth}\Large\centering
\begin{minipage}{0.95\linewidth}\raggedright\linespread{1.25}\selectfont Voici la loi de probabilité d'une variable aléatoire $X$. Que vaut $E(X)$~? $\begin{array}{|l|c|c|c|}\hline x_i & -2 & 1 & 5\\\hline P(X=x_i) & 0{,}3 & 0{,}5 & 0{,}2\\\hline\end{array}$\end{minipage}
\end{minipage}

\vfill
\begin{tcolorbox}[enhanced,nobeforeafter,width=7.1cm,height=2.1cm,valign=center,colback=ccvp!6,colframe=ccvp,arc=2.5mm,boxrule=0.9pt,left=3mm]{\Large\bfseries\color{ccvp}a)}\quad\begin{minipage}[c]{5.3cm}\raggedright\Large $1{,}9$\end{minipage}\end{tcolorbox}\hfill\begin{tcolorbox}[enhanced,nobeforeafter,width=7.1cm,height=2.1cm,valign=center,colback=ccvp!6,colframe=ccvp,arc=2.5mm,boxrule=0.9pt,left=3mm]{\Large\bfseries\color{ccvp}b)}\quad\begin{minipage}[c]{5.3cm}\raggedright\Large $0{,}9$\end{minipage}\end{tcolorbox}\par\vspace{2.5mm}
\begin{tcolorbox}[enhanced,nobeforeafter,width=7.1cm,height=2.1cm,valign=center,colback=ccvp!6,colframe=ccvp,arc=2.5mm,boxrule=0.9pt,left=3mm]{\Large\bfseries\color{ccvp}c)}\quad\begin{minipage}[c]{5.3cm}\raggedright\Large $\dfrac{4}{3}$\end{minipage}\end{tcolorbox}\hfill\begin{tcolorbox}[enhanced,nobeforeafter,width=7.1cm,height=2.1cm,valign=center,colback=ccvp!6,colframe=ccvp,arc=2.5mm,boxrule=0.9pt,left=3mm]{\Large\bfseries\color{ccvp}d)}\quad\begin{minipage}[c]{5.3cm}\raggedright\Large $-0{,}9$\end{minipage}\end{tcolorbox}
\end{tcolorbox}
\newpage
\begin{tcolorbox}[enhanced,width=16cm,height=14cm,colback=white,colframe=ccvp,boxrule=2pt,arc=0mm,
left=6mm,right=6mm,top=0mm,bottom=4mm,sharp corners,
overlay={\fill[ccvp] (frame.north west) rectangle ([yshift=-1.25cm]frame.north east);
\node[anchor=west,text=white,font=\sffamily\bfseries\LARGE] at ([xshift=5mm,yshift=-0.62cm]frame.north west) {Question 90};
\node[anchor=east,text=white,font=\sffamily\itshape\large] at ([xshift=-5mm,yshift=-0.62cm]frame.north east) {Probabilités conditionnelles et variables aléatoires};}]
\vspace*{1.75cm}
\begin{minipage}[t][6.9cm][c]{\linewidth}\Large\centering
\begin{minipage}{0.95\linewidth}\raggedright\linespread{1.25}\selectfont Voici la loi de probabilité d'une variable aléatoire $X$. Que vaut $E(X)$~? $\begin{array}{|l|c|c|c|}\hline x_i & 0 & 2 & 4\\\hline P(X=x_i) & 0{,}25 & 0{,}5 & 0{,}25\\\hline\end{array}$\end{minipage}
\end{minipage}

\vfill
\begin{tcolorbox}[enhanced,nobeforeafter,width=7.1cm,height=2.1cm,valign=center,colback=ccvp!6,colframe=ccvp,arc=2.5mm,boxrule=0.9pt,left=3mm]{\Large\bfseries\color{ccvp}a)}\quad\begin{minipage}[c]{5.3cm}\raggedright\Large $4$\end{minipage}\end{tcolorbox}\hfill\begin{tcolorbox}[enhanced,nobeforeafter,width=7.1cm,height=2.1cm,valign=center,colback=ccvp!6,colframe=ccvp,arc=2.5mm,boxrule=0.9pt,left=3mm]{\Large\bfseries\color{ccvp}b)}\quad\begin{minipage}[c]{5.3cm}\raggedright\Large $2$\end{minipage}\end{tcolorbox}\par\vspace{2.5mm}
\begin{tcolorbox}[enhanced,nobeforeafter,width=7.1cm,height=2.1cm,valign=center,colback=ccvp!6,colframe=ccvp,arc=2.5mm,boxrule=0.9pt,left=3mm]{\Large\bfseries\color{ccvp}c)}\quad\begin{minipage}[c]{5.3cm}\raggedright\Large $-2$\end{minipage}\end{tcolorbox}\hfill\begin{tcolorbox}[enhanced,nobeforeafter,width=7.1cm,height=2.1cm,valign=center,colback=ccvp!6,colframe=ccvp,arc=2.5mm,boxrule=0.9pt,left=3mm]{\Large\bfseries\color{ccvp}d)}\quad\begin{minipage}[c]{5.3cm}\raggedright\Large $3$\end{minipage}\end{tcolorbox}
\end{tcolorbox}
\newpage
\begin{tcolorbox}[enhanced,width=16cm,height=14cm,colback=white,colframe=ccvp,boxrule=2pt,arc=0mm,
left=6mm,right=6mm,top=0mm,bottom=4mm,sharp corners,
overlay={\fill[ccvp] (frame.north west) rectangle ([yshift=-1.25cm]frame.north east);
\node[anchor=west,text=white,font=\sffamily\bfseries\LARGE] at ([xshift=5mm,yshift=-0.62cm]frame.north west) {Question 91};
\node[anchor=east,text=white,font=\sffamily\itshape\large] at ([xshift=-5mm,yshift=-0.62cm]frame.north east) {Probabilités conditionnelles et variables aléatoires};}]
\vspace*{1.75cm}
\begin{minipage}[t][6.9cm][c]{\linewidth}\Large\centering
\begin{minipage}{0.95\linewidth}\raggedright\linespread{1.25}\selectfont Voici la loi de probabilité d'une variable aléatoire $X$. Que vaut sa variance $V(X)$~? $\begin{array}{|l|c|c|}\hline x_i & 0 & 2\\\hline P(X=x_i) & 0{,}5 & 0{,}5\\\hline\end{array}$\end{minipage}
\end{minipage}

\vfill
\begin{tcolorbox}[enhanced,nobeforeafter,width=7.1cm,height=2.1cm,valign=center,colback=ccvp!6,colframe=ccvp,arc=2.5mm,boxrule=0.9pt,left=3mm]{\Large\bfseries\color{ccvp}a)}\quad\begin{minipage}[c]{5.3cm}\raggedright\Large $1$\end{minipage}\end{tcolorbox}\hfill\begin{tcolorbox}[enhanced,nobeforeafter,width=7.1cm,height=2.1cm,valign=center,colback=ccvp!6,colframe=ccvp,arc=2.5mm,boxrule=0.9pt,left=3mm]{\Large\bfseries\color{ccvp}b)}\quad\begin{minipage}[c]{5.3cm}\raggedright\Large $0$\end{minipage}\end{tcolorbox}\par\vspace{2.5mm}
\begin{tcolorbox}[enhanced,nobeforeafter,width=7.1cm,height=2.1cm,valign=center,colback=ccvp!6,colframe=ccvp,arc=2.5mm,boxrule=0.9pt,left=3mm]{\Large\bfseries\color{ccvp}c)}\quad\begin{minipage}[c]{5.3cm}\raggedright\Large $4$\end{minipage}\end{tcolorbox}\hfill\begin{tcolorbox}[enhanced,nobeforeafter,width=7.1cm,height=2.1cm,valign=center,colback=ccvp!6,colframe=ccvp,arc=2.5mm,boxrule=0.9pt,left=3mm]{\Large\bfseries\color{ccvp}d)}\quad\begin{minipage}[c]{5.3cm}\raggedright\Large $2$\end{minipage}\end{tcolorbox}
\end{tcolorbox}
\newpage
\begin{tcolorbox}[enhanced,width=16cm,height=14cm,colback=white,colframe=ccvp,boxrule=2pt,arc=0mm,
left=6mm,right=6mm,top=0mm,bottom=4mm,sharp corners,
overlay={\fill[ccvp] (frame.north west) rectangle ([yshift=-1.25cm]frame.north east);
\node[anchor=west,text=white,font=\sffamily\bfseries\LARGE] at ([xshift=5mm,yshift=-0.62cm]frame.north west) {Question 92};
\node[anchor=east,text=white,font=\sffamily\itshape\large] at ([xshift=-5mm,yshift=-0.62cm]frame.north east) {Probabilités conditionnelles et variables aléatoires};}]
\vspace*{1.75cm}
\begin{minipage}[t][6.9cm][c]{\linewidth}\Large\centering
\begin{minipage}{0.95\linewidth}\raggedright\linespread{1.25}\selectfont Voici la loi de probabilité d'une variable aléatoire $X$. Que vaut sa variance $V(X)$~? $\begin{array}{|l|c|c|}\hline x_i & 0 & 4\\\hline P(X=x_i) & 0{,}25 & 0{,}75\\\hline\end{array}$\end{minipage}
\end{minipage}

\vfill
\begin{tcolorbox}[enhanced,nobeforeafter,width=7.1cm,height=2.1cm,valign=center,colback=ccvp!6,colframe=ccvp,arc=2.5mm,boxrule=0.9pt,left=3mm]{\Large\bfseries\color{ccvp}a)}\quad\begin{minipage}[c]{5.3cm}\raggedright\Large $6$\end{minipage}\end{tcolorbox}\hfill\begin{tcolorbox}[enhanced,nobeforeafter,width=7.1cm,height=2.1cm,valign=center,colback=ccvp!6,colframe=ccvp,arc=2.5mm,boxrule=0.9pt,left=3mm]{\Large\bfseries\color{ccvp}b)}\quad\begin{minipage}[c]{5.3cm}\raggedright\Large $3$\end{minipage}\end{tcolorbox}\par\vspace{2.5mm}
\begin{tcolorbox}[enhanced,nobeforeafter,width=7.1cm,height=2.1cm,valign=center,colback=ccvp!6,colframe=ccvp,arc=2.5mm,boxrule=0.9pt,left=3mm]{\Large\bfseries\color{ccvp}c)}\quad\begin{minipage}[c]{5.3cm}\raggedright\Large $12$\end{minipage}\end{tcolorbox}\hfill\begin{tcolorbox}[enhanced,nobeforeafter,width=7.1cm,height=2.1cm,valign=center,colback=ccvp!6,colframe=ccvp,arc=2.5mm,boxrule=0.9pt,left=3mm]{\Large\bfseries\color{ccvp}d)}\quad\begin{minipage}[c]{5.3cm}\raggedright\Large $4$\end{minipage}\end{tcolorbox}
\end{tcolorbox}
\newpage
\begin{tcolorbox}[enhanced,width=16cm,height=14cm,colback=white,colframe=ccvp,boxrule=2pt,arc=0mm,
left=6mm,right=6mm,top=0mm,bottom=4mm,sharp corners,
overlay={\fill[ccvp] (frame.north west) rectangle ([yshift=-1.25cm]frame.north east);
\node[anchor=west,text=white,font=\sffamily\bfseries\LARGE] at ([xshift=5mm,yshift=-0.62cm]frame.north west) {Question 93};
\node[anchor=east,text=white,font=\sffamily\itshape\large] at ([xshift=-5mm,yshift=-0.62cm]frame.north east) {Probabilités conditionnelles et variables aléatoires};}]
\vspace*{1.75cm}
\begin{minipage}[t][6.9cm][c]{\linewidth}\Large\centering
\begin{minipage}{0.95\linewidth}\raggedright\linespread{1.25}\selectfont On choisit au hasard une fille de ce groupe. Quelle est la probabilité que ce soit un demi-pensionnaire~? $\begin{array}{|l|c|c|}\hline & \text{Demi-pens.} & \text{Externes}\\\hline \text{Filles} & 12 & 8\\\hline \text{Garçons} & 10 & 20\\\hline\end{array}$\end{minipage}
\end{minipage}

\vfill
\begin{tcolorbox}[enhanced,nobeforeafter,width=7.1cm,height=2.1cm,valign=center,colback=ccvp!6,colframe=ccvp,arc=2.5mm,boxrule=0.9pt,left=3mm]{\Large\bfseries\color{ccvp}a)}\quad\begin{minipage}[c]{5.3cm}\raggedright\Large $\dfrac{6}{25}$\end{minipage}\end{tcolorbox}\hfill\begin{tcolorbox}[enhanced,nobeforeafter,width=7.1cm,height=2.1cm,valign=center,colback=ccvp!6,colframe=ccvp,arc=2.5mm,boxrule=0.9pt,left=3mm]{\Large\bfseries\color{ccvp}b)}\quad\begin{minipage}[c]{5.3cm}\raggedright\Large $\dfrac{3}{5}$\end{minipage}\end{tcolorbox}\par\vspace{2.5mm}
\begin{tcolorbox}[enhanced,nobeforeafter,width=7.1cm,height=2.1cm,valign=center,colback=ccvp!6,colframe=ccvp,arc=2.5mm,boxrule=0.9pt,left=3mm]{\Large\bfseries\color{ccvp}c)}\quad\begin{minipage}[c]{5.3cm}\raggedright\Large $\dfrac{2}{5}$\end{minipage}\end{tcolorbox}\hfill\begin{tcolorbox}[enhanced,nobeforeafter,width=7.1cm,height=2.1cm,valign=center,colback=ccvp!6,colframe=ccvp,arc=2.5mm,boxrule=0.9pt,left=3mm]{\Large\bfseries\color{ccvp}d)}\quad\begin{minipage}[c]{5.3cm}\raggedright\Large $\dfrac{6}{11}$\end{minipage}\end{tcolorbox}
\end{tcolorbox}
\newpage
\begin{tcolorbox}[enhanced,width=16cm,height=14cm,colback=white,colframe=ccvp,boxrule=2pt,arc=0mm,
left=6mm,right=6mm,top=0mm,bottom=4mm,sharp corners,
overlay={\fill[ccvp] (frame.north west) rectangle ([yshift=-1.25cm]frame.north east);
\node[anchor=west,text=white,font=\sffamily\bfseries\LARGE] at ([xshift=5mm,yshift=-0.62cm]frame.north west) {Question 94};
\node[anchor=east,text=white,font=\sffamily\itshape\large] at ([xshift=-5mm,yshift=-0.62cm]frame.north east) {Probabilités conditionnelles et variables aléatoires};}]
\vspace*{1.75cm}
\begin{minipage}[t][6.9cm][c]{\linewidth}\Large\centering
\begin{minipage}{0.95\linewidth}\raggedright\linespread{1.25}\selectfont On choisit au hasard un garçon de ce groupe. Quelle est la probabilité que ce soit un demi-pensionnaire~? $\begin{array}{|l|c|c|}\hline & \text{Demi-pens.} & \text{Externes}\\\hline \text{Filles} & 15 & 5\\\hline \text{Garçons} & 9 & 11\\\hline\end{array}$\end{minipage}
\end{minipage}

\vfill
\begin{tcolorbox}[enhanced,nobeforeafter,width=7.1cm,height=2.1cm,valign=center,colback=ccvp!6,colframe=ccvp,arc=2.5mm,boxrule=0.9pt,left=3mm]{\Large\bfseries\color{ccvp}a)}\quad\begin{minipage}[c]{5.3cm}\raggedright\Large $\dfrac{9}{20}$\end{minipage}\end{tcolorbox}\hfill\begin{tcolorbox}[enhanced,nobeforeafter,width=7.1cm,height=2.1cm,valign=center,colback=ccvp!6,colframe=ccvp,arc=2.5mm,boxrule=0.9pt,left=3mm]{\Large\bfseries\color{ccvp}b)}\quad\begin{minipage}[c]{5.3cm}\raggedright\Large $\dfrac{11}{20}$\end{minipage}\end{tcolorbox}\par\vspace{2.5mm}
\begin{tcolorbox}[enhanced,nobeforeafter,width=7.1cm,height=2.1cm,valign=center,colback=ccvp!6,colframe=ccvp,arc=2.5mm,boxrule=0.9pt,left=3mm]{\Large\bfseries\color{ccvp}c)}\quad\begin{minipage}[c]{5.3cm}\raggedright\Large $\dfrac{9}{40}$\end{minipage}\end{tcolorbox}\hfill\begin{tcolorbox}[enhanced,nobeforeafter,width=7.1cm,height=2.1cm,valign=center,colback=ccvp!6,colframe=ccvp,arc=2.5mm,boxrule=0.9pt,left=3mm]{\Large\bfseries\color{ccvp}d)}\quad\begin{minipage}[c]{5.3cm}\raggedright\Large $\dfrac{3}{8}$\end{minipage}\end{tcolorbox}
\end{tcolorbox}
\newpage
\begin{tcolorbox}[enhanced,width=16cm,height=14cm,colback=white,colframe=ccvp,boxrule=2pt,arc=0mm,
left=6mm,right=6mm,top=0mm,bottom=4mm,sharp corners,
overlay={\fill[ccvp] (frame.north west) rectangle ([yshift=-1.25cm]frame.north east);
\node[anchor=west,text=white,font=\sffamily\bfseries\LARGE] at ([xshift=5mm,yshift=-0.62cm]frame.north west) {Question 95};
\node[anchor=east,text=white,font=\sffamily\itshape\large] at ([xshift=-5mm,yshift=-0.62cm]frame.north east) {Probabilités conditionnelles et variables aléatoires};}]
\vspace*{1.75cm}
\begin{minipage}[t][6.9cm][c]{\linewidth}\Large\centering
\begin{minipage}{0.95\linewidth}\raggedright\linespread{1.25}\selectfont Voici la loi de probabilité d'une variable aléatoire $X$. Que vaut $P(X\geqslant 2)$~? $\begin{array}{|l|c|c|c|c|c|}\hline x_i & 0 & 1 & 2 & 3 & 4\\\hline P(X=x_i) & 0{,}1 & 0{,}2 & 0{,}3 & 0{,}25 & 0{,}15\\\hline\end{array}$\end{minipage}
\end{minipage}

\vfill
\begin{tcolorbox}[enhanced,nobeforeafter,width=7.1cm,height=2.1cm,valign=center,colback=ccvp!6,colframe=ccvp,arc=2.5mm,boxrule=0.9pt,left=3mm]{\Large\bfseries\color{ccvp}a)}\quad\begin{minipage}[c]{5.3cm}\raggedright\Large $0{,}3$\end{minipage}\end{tcolorbox}\hfill\begin{tcolorbox}[enhanced,nobeforeafter,width=7.1cm,height=2.1cm,valign=center,colback=ccvp!6,colframe=ccvp,arc=2.5mm,boxrule=0.9pt,left=3mm]{\Large\bfseries\color{ccvp}b)}\quad\begin{minipage}[c]{5.3cm}\raggedright\Large $0{,}4$\end{minipage}\end{tcolorbox}\par\vspace{2.5mm}
\begin{tcolorbox}[enhanced,nobeforeafter,width=7.1cm,height=2.1cm,valign=center,colback=ccvp!6,colframe=ccvp,arc=2.5mm,boxrule=0.9pt,left=3mm]{\Large\bfseries\color{ccvp}c)}\quad\begin{minipage}[c]{5.3cm}\raggedright\Large $0{,}9$\end{minipage}\end{tcolorbox}\hfill\begin{tcolorbox}[enhanced,nobeforeafter,width=7.1cm,height=2.1cm,valign=center,colback=ccvp!6,colframe=ccvp,arc=2.5mm,boxrule=0.9pt,left=3mm]{\Large\bfseries\color{ccvp}d)}\quad\begin{minipage}[c]{5.3cm}\raggedright\Large $0{,}7$\end{minipage}\end{tcolorbox}
\end{tcolorbox}
\newpage
\begin{tcolorbox}[enhanced,width=16cm,height=14cm,colback=white,colframe=ccvp,boxrule=2pt,arc=0mm,
left=6mm,right=6mm,top=0mm,bottom=4mm,sharp corners,
overlay={\fill[ccvp] (frame.north west) rectangle ([yshift=-1.25cm]frame.north east);
\node[anchor=west,text=white,font=\sffamily\bfseries\LARGE] at ([xshift=5mm,yshift=-0.62cm]frame.north west) {Question 96};
\node[anchor=east,text=white,font=\sffamily\itshape\large] at ([xshift=-5mm,yshift=-0.62cm]frame.north east) {Probabilités conditionnelles et variables aléatoires};}]
\vspace*{1.75cm}
\begin{minipage}[t][6.9cm][c]{\linewidth}\Large\centering
\begin{minipage}{0.95\linewidth}\raggedright\linespread{1.25}\selectfont Voici la loi de probabilité d'une variable aléatoire $X$. Que vaut $P(X\geqslant 3)$~? $\begin{array}{|l|c|c|c|c|c|}\hline x_i & 0 & 1 & 2 & 3 & 4\\\hline P(X=x_i) & 0{,}2 & 0{,}2 & 0{,}2 & 0{,}2 & 0{,}2\\\hline\end{array}$\end{minipage}
\end{minipage}

\vfill
\begin{tcolorbox}[enhanced,nobeforeafter,width=7.1cm,height=2.1cm,valign=center,colback=ccvp!6,colframe=ccvp,arc=2.5mm,boxrule=0.9pt,left=3mm]{\Large\bfseries\color{ccvp}a)}\quad\begin{minipage}[c]{5.3cm}\raggedright\Large $0{,}6$\end{minipage}\end{tcolorbox}\hfill\begin{tcolorbox}[enhanced,nobeforeafter,width=7.1cm,height=2.1cm,valign=center,colback=ccvp!6,colframe=ccvp,arc=2.5mm,boxrule=0.9pt,left=3mm]{\Large\bfseries\color{ccvp}b)}\quad\begin{minipage}[c]{5.3cm}\raggedright\Large $0{,}4$\end{minipage}\end{tcolorbox}\par\vspace{2.5mm}
\begin{tcolorbox}[enhanced,nobeforeafter,width=7.1cm,height=2.1cm,valign=center,colback=ccvp!6,colframe=ccvp,arc=2.5mm,boxrule=0.9pt,left=3mm]{\Large\bfseries\color{ccvp}c)}\quad\begin{minipage}[c]{5.3cm}\raggedright\Large $0{,}3$\end{minipage}\end{tcolorbox}\hfill\begin{tcolorbox}[enhanced,nobeforeafter,width=7.1cm,height=2.1cm,valign=center,colback=ccvp!6,colframe=ccvp,arc=2.5mm,boxrule=0.9pt,left=3mm]{\Large\bfseries\color{ccvp}d)}\quad\begin{minipage}[c]{5.3cm}\raggedright\Large $0{,}2$\end{minipage}\end{tcolorbox}
\end{tcolorbox}
\newpage
\begin{tcolorbox}[enhanced,width=16cm,height=14cm,colback=white,colframe=csui,boxrule=2pt,arc=0mm,
left=6mm,right=6mm,top=0mm,bottom=4mm,sharp corners,
overlay={\fill[csui] (frame.north west) rectangle ([yshift=-1.25cm]frame.north east);
\node[anchor=west,text=white,font=\sffamily\bfseries\LARGE] at ([xshift=5mm,yshift=-0.62cm]frame.north west) {Question 1 B};
\node[anchor=east,text=white,font=\sffamily\itshape\large] at ([xshift=-5mm,yshift=-0.62cm]frame.north east) {Suites};}]
\vspace*{1.75cm}
\begin{minipage}[t][6.9cm][c]{\linewidth}\LARGE\centering
\begin{minipage}{0.95\linewidth}\raggedright\linespread{1.25}\selectfont $(u_n)$ est la suite arithmétique de premier terme $u_0=-4$ et de raison $5$. Que vaut $u_{12}$~?\end{minipage}
\end{minipage}

\vfill
\begin{tcolorbox}[enhanced,nobeforeafter,width=7.1cm,height=2.1cm,valign=center,colback=csui!6,colframe=csui,arc=2.5mm,boxrule=0.9pt,left=3mm]{\Large\bfseries\color{csui}a)}\quad\begin{minipage}[c]{5.3cm}\raggedright\Large $51$\end{minipage}\end{tcolorbox}\hfill\begin{tcolorbox}[enhanced,nobeforeafter,width=7.1cm,height=2.1cm,valign=center,colback=csui!6,colframe=csui,arc=2.5mm,boxrule=0.9pt,left=3mm]{\Large\bfseries\color{csui}b)}\quad\begin{minipage}[c]{5.3cm}\raggedright\Large $61$\end{minipage}\end{tcolorbox}\par\vspace{2.5mm}
\begin{tcolorbox}[enhanced,nobeforeafter,width=7.1cm,height=2.1cm,valign=center,colback=csui!6,colframe=csui,arc=2.5mm,boxrule=0.9pt,left=3mm]{\Large\bfseries\color{csui}c)}\quad\begin{minipage}[c]{5.3cm}\raggedright\Large $60$\end{minipage}\end{tcolorbox}\hfill\begin{tcolorbox}[enhanced,nobeforeafter,width=7.1cm,height=2.1cm,valign=center,colback=csui!6,colframe=csui,arc=2.5mm,boxrule=0.9pt,left=3mm]{\Large\bfseries\color{csui}d)}\quad\begin{minipage}[c]{5.3cm}\raggedright\Large $56$\end{minipage}\end{tcolorbox}
\end{tcolorbox}
\newpage
\begin{tcolorbox}[enhanced,width=16cm,height=14cm,colback=white,colframe=csui,boxrule=2pt,arc=0mm,
left=6mm,right=6mm,top=0mm,bottom=4mm,sharp corners,
overlay={\fill[csui] (frame.north west) rectangle ([yshift=-1.25cm]frame.north east);
\node[anchor=west,text=white,font=\sffamily\bfseries\LARGE] at ([xshift=5mm,yshift=-0.62cm]frame.north west) {Question 2 B};
\node[anchor=east,text=white,font=\sffamily\itshape\large] at ([xshift=-5mm,yshift=-0.62cm]frame.north east) {Suites};}]
\vspace*{1.75cm}
\begin{minipage}[t][6.9cm][c]{\linewidth}\LARGE\centering
\begin{minipage}{0.95\linewidth}\raggedright\linespread{1.25}\selectfont $(u_n)$ est la suite arithmétique de premier terme $u_0=7$ et de raison $4$. Que vaut $u_{15}$~?\end{minipage}
\end{minipage}

\vfill
\begin{tcolorbox}[enhanced,nobeforeafter,width=7.1cm,height=2.1cm,valign=center,colback=csui!6,colframe=csui,arc=2.5mm,boxrule=0.9pt,left=3mm]{\Large\bfseries\color{csui}a)}\quad\begin{minipage}[c]{5.3cm}\raggedright\Large $67$\end{minipage}\end{tcolorbox}\hfill\begin{tcolorbox}[enhanced,nobeforeafter,width=7.1cm,height=2.1cm,valign=center,colback=csui!6,colframe=csui,arc=2.5mm,boxrule=0.9pt,left=3mm]{\Large\bfseries\color{csui}b)}\quad\begin{minipage}[c]{5.3cm}\raggedright\Large $71$\end{minipage}\end{tcolorbox}\par\vspace{2.5mm}
\begin{tcolorbox}[enhanced,nobeforeafter,width=7.1cm,height=2.1cm,valign=center,colback=csui!6,colframe=csui,arc=2.5mm,boxrule=0.9pt,left=3mm]{\Large\bfseries\color{csui}c)}\quad\begin{minipage}[c]{5.3cm}\raggedright\Large $63$\end{minipage}\end{tcolorbox}\hfill\begin{tcolorbox}[enhanced,nobeforeafter,width=7.1cm,height=2.1cm,valign=center,colback=csui!6,colframe=csui,arc=2.5mm,boxrule=0.9pt,left=3mm]{\Large\bfseries\color{csui}d)}\quad\begin{minipage}[c]{5.3cm}\raggedright\Large $60$\end{minipage}\end{tcolorbox}
\end{tcolorbox}
\newpage
\begin{tcolorbox}[enhanced,width=16cm,height=14cm,colback=white,colframe=csui,boxrule=2pt,arc=0mm,
left=6mm,right=6mm,top=0mm,bottom=4mm,sharp corners,
overlay={\fill[csui] (frame.north west) rectangle ([yshift=-1.25cm]frame.north east);
\node[anchor=west,text=white,font=\sffamily\bfseries\LARGE] at ([xshift=5mm,yshift=-0.62cm]frame.north west) {Question 3 B};
\node[anchor=east,text=white,font=\sffamily\itshape\large] at ([xshift=-5mm,yshift=-0.62cm]frame.north east) {Suites};}]
\vspace*{1.75cm}
\begin{minipage}[t][6.9cm][c]{\linewidth}\LARGE\centering
\begin{minipage}{0.95\linewidth}\raggedright\linespread{1.25}\selectfont $(u_n)$ est la suite géométrique de premier terme $u_0=2$ et de raison $-2$. Que vaut $u_{5}$~?\end{minipage}
\end{minipage}

\vfill
\begin{tcolorbox}[enhanced,nobeforeafter,width=7.1cm,height=2.1cm,valign=center,colback=csui!6,colframe=csui,arc=2.5mm,boxrule=0.9pt,left=3mm]{\Large\bfseries\color{csui}a)}\quad\begin{minipage}[c]{5.3cm}\raggedright\Large $-8$\end{minipage}\end{tcolorbox}\hfill\begin{tcolorbox}[enhanced,nobeforeafter,width=7.1cm,height=2.1cm,valign=center,colback=csui!6,colframe=csui,arc=2.5mm,boxrule=0.9pt,left=3mm]{\Large\bfseries\color{csui}b)}\quad\begin{minipage}[c]{5.3cm}\raggedright\Large $32$\end{minipage}\end{tcolorbox}\par\vspace{2.5mm}
\begin{tcolorbox}[enhanced,nobeforeafter,width=7.1cm,height=2.1cm,valign=center,colback=csui!6,colframe=csui,arc=2.5mm,boxrule=0.9pt,left=3mm]{\Large\bfseries\color{csui}c)}\quad\begin{minipage}[c]{5.3cm}\raggedright\Large $-64$\end{minipage}\end{tcolorbox}\hfill\begin{tcolorbox}[enhanced,nobeforeafter,width=7.1cm,height=2.1cm,valign=center,colback=csui!6,colframe=csui,arc=2.5mm,boxrule=0.9pt,left=3mm]{\Large\bfseries\color{csui}d)}\quad\begin{minipage}[c]{5.3cm}\raggedright\Large $128$\end{minipage}\end{tcolorbox}
\end{tcolorbox}
\newpage
\begin{tcolorbox}[enhanced,width=16cm,height=14cm,colback=white,colframe=csui,boxrule=2pt,arc=0mm,
left=6mm,right=6mm,top=0mm,bottom=4mm,sharp corners,
overlay={\fill[csui] (frame.north west) rectangle ([yshift=-1.25cm]frame.north east);
\node[anchor=west,text=white,font=\sffamily\bfseries\LARGE] at ([xshift=5mm,yshift=-0.62cm]frame.north west) {Question 4 B};
\node[anchor=east,text=white,font=\sffamily\itshape\large] at ([xshift=-5mm,yshift=-0.62cm]frame.north east) {Suites};}]
\vspace*{1.75cm}
\begin{minipage}[t][6.9cm][c]{\linewidth}\LARGE\centering
\begin{minipage}{0.95\linewidth}\raggedright\linespread{1.25}\selectfont $(u_n)$ est la suite géométrique de premier terme $u_0=1$ et de raison $4$. Que vaut $u_{4}$~?\end{minipage}
\end{minipage}

\vfill
\begin{tcolorbox}[enhanced,nobeforeafter,width=7.1cm,height=2.1cm,valign=center,colback=csui!6,colframe=csui,arc=2.5mm,boxrule=0.9pt,left=3mm]{\Large\bfseries\color{csui}a)}\quad\begin{minipage}[c]{5.3cm}\raggedright\Large $1\,024$\end{minipage}\end{tcolorbox}\hfill\begin{tcolorbox}[enhanced,nobeforeafter,width=7.1cm,height=2.1cm,valign=center,colback=csui!6,colframe=csui,arc=2.5mm,boxrule=0.9pt,left=3mm]{\Large\bfseries\color{csui}b)}\quad\begin{minipage}[c]{5.3cm}\raggedright\Large $256$\end{minipage}\end{tcolorbox}\par\vspace{2.5mm}
\begin{tcolorbox}[enhanced,nobeforeafter,width=7.1cm,height=2.1cm,valign=center,colback=csui!6,colframe=csui,arc=2.5mm,boxrule=0.9pt,left=3mm]{\Large\bfseries\color{csui}c)}\quad\begin{minipage}[c]{5.3cm}\raggedright\Large $64$\end{minipage}\end{tcolorbox}\hfill\begin{tcolorbox}[enhanced,nobeforeafter,width=7.1cm,height=2.1cm,valign=center,colback=csui!6,colframe=csui,arc=2.5mm,boxrule=0.9pt,left=3mm]{\Large\bfseries\color{csui}d)}\quad\begin{minipage}[c]{5.3cm}\raggedright\Large $17$\end{minipage}\end{tcolorbox}
\end{tcolorbox}
\newpage
\begin{tcolorbox}[enhanced,width=16cm,height=14cm,colback=white,colframe=csui,boxrule=2pt,arc=0mm,
left=6mm,right=6mm,top=0mm,bottom=4mm,sharp corners,
overlay={\fill[csui] (frame.north west) rectangle ([yshift=-1.25cm]frame.north east);
\node[anchor=west,text=white,font=\sffamily\bfseries\LARGE] at ([xshift=5mm,yshift=-0.62cm]frame.north west) {Question 5 B};
\node[anchor=east,text=white,font=\sffamily\itshape\large] at ([xshift=-5mm,yshift=-0.62cm]frame.north east) {Suites};}]
\vspace*{1.75cm}
\begin{minipage}[t][6.9cm][c]{\linewidth}\LARGE\centering
\begin{minipage}{0.95\linewidth}\raggedright\linespread{1.25}\selectfont On définit $u_0=12$ et, pour tout entier $n$, $u_{n+1}=\dfrac{1}{2}u_n+4$. Que vaut $u_2$~?\end{minipage}
\end{minipage}

\vfill
\begin{tcolorbox}[enhanced,nobeforeafter,width=7.1cm,height=2.1cm,valign=center,colback=csui!6,colframe=csui,arc=2.5mm,boxrule=0.9pt,left=3mm]{\Large\bfseries\color{csui}a)}\quad\begin{minipage}[c]{5.3cm}\raggedright\Large $9$\end{minipage}\end{tcolorbox}\hfill\begin{tcolorbox}[enhanced,nobeforeafter,width=7.1cm,height=2.1cm,valign=center,colback=csui!6,colframe=csui,arc=2.5mm,boxrule=0.9pt,left=3mm]{\Large\bfseries\color{csui}b)}\quad\begin{minipage}[c]{5.3cm}\raggedright\Large $\dfrac{17}{2}$\end{minipage}\end{tcolorbox}\par\vspace{2.5mm}
\begin{tcolorbox}[enhanced,nobeforeafter,width=7.1cm,height=2.1cm,valign=center,colback=csui!6,colframe=csui,arc=2.5mm,boxrule=0.9pt,left=3mm]{\Large\bfseries\color{csui}c)}\quad\begin{minipage}[c]{5.3cm}\raggedright\Large $14$\end{minipage}\end{tcolorbox}\hfill\begin{tcolorbox}[enhanced,nobeforeafter,width=7.1cm,height=2.1cm,valign=center,colback=csui!6,colframe=csui,arc=2.5mm,boxrule=0.9pt,left=3mm]{\Large\bfseries\color{csui}d)}\quad\begin{minipage}[c]{5.3cm}\raggedright\Large $10$\end{minipage}\end{tcolorbox}
\end{tcolorbox}
\newpage
\begin{tcolorbox}[enhanced,width=16cm,height=14cm,colback=white,colframe=csui,boxrule=2pt,arc=0mm,
left=6mm,right=6mm,top=0mm,bottom=4mm,sharp corners,
overlay={\fill[csui] (frame.north west) rectangle ([yshift=-1.25cm]frame.north east);
\node[anchor=west,text=white,font=\sffamily\bfseries\LARGE] at ([xshift=5mm,yshift=-0.62cm]frame.north west) {Question 6 B};
\node[anchor=east,text=white,font=\sffamily\itshape\large] at ([xshift=-5mm,yshift=-0.62cm]frame.north east) {Suites};}]
\vspace*{1.75cm}
\begin{minipage}[t][6.9cm][c]{\linewidth}\LARGE\centering
\begin{minipage}{0.95\linewidth}\raggedright\linespread{1.25}\selectfont On définit $u_0=2$ et, pour tout entier $n$, $u_{n+1}=-2u_n+3$. Que vaut $u_2$~?\end{minipage}
\end{minipage}

\vfill
\begin{tcolorbox}[enhanced,nobeforeafter,width=7.1cm,height=2.1cm,valign=center,colback=csui!6,colframe=csui,arc=2.5mm,boxrule=0.9pt,left=3mm]{\Large\bfseries\color{csui}a)}\quad\begin{minipage}[c]{5.3cm}\raggedright\Large $-7$\end{minipage}\end{tcolorbox}\hfill\begin{tcolorbox}[enhanced,nobeforeafter,width=7.1cm,height=2.1cm,valign=center,colback=csui!6,colframe=csui,arc=2.5mm,boxrule=0.9pt,left=3mm]{\Large\bfseries\color{csui}b)}\quad\begin{minipage}[c]{5.3cm}\raggedright\Large $2$\end{minipage}\end{tcolorbox}\par\vspace{2.5mm}
\begin{tcolorbox}[enhanced,nobeforeafter,width=7.1cm,height=2.1cm,valign=center,colback=csui!6,colframe=csui,arc=2.5mm,boxrule=0.9pt,left=3mm]{\Large\bfseries\color{csui}c)}\quad\begin{minipage}[c]{5.3cm}\raggedright\Large $-1$\end{minipage}\end{tcolorbox}\hfill\begin{tcolorbox}[enhanced,nobeforeafter,width=7.1cm,height=2.1cm,valign=center,colback=csui!6,colframe=csui,arc=2.5mm,boxrule=0.9pt,left=3mm]{\Large\bfseries\color{csui}d)}\quad\begin{minipage}[c]{5.3cm}\raggedright\Large $5$\end{minipage}\end{tcolorbox}
\end{tcolorbox}
\newpage
\begin{tcolorbox}[enhanced,width=16cm,height=14cm,colback=white,colframe=csui,boxrule=2pt,arc=0mm,
left=6mm,right=6mm,top=0mm,bottom=4mm,sharp corners,
overlay={\fill[csui] (frame.north west) rectangle ([yshift=-1.25cm]frame.north east);
\node[anchor=west,text=white,font=\sffamily\bfseries\LARGE] at ([xshift=5mm,yshift=-0.62cm]frame.north west) {Question 7 B};
\node[anchor=east,text=white,font=\sffamily\itshape\large] at ([xshift=-5mm,yshift=-0.62cm]frame.north east) {Suites};}]
\vspace*{1.75cm}
\begin{minipage}[t][6.9cm][c]{\linewidth}\LARGE\centering
\begin{minipage}{0.95\linewidth}\raggedright\linespread{1.25}\selectfont Que vaut la somme $1+2+3+\dots+50$~?\end{minipage}
\end{minipage}

\vfill
\begin{tcolorbox}[enhanced,nobeforeafter,width=7.1cm,height=2.1cm,valign=center,colback=csui!6,colframe=csui,arc=2.5mm,boxrule=0.9pt,left=3mm]{\Large\bfseries\color{csui}a)}\quad\begin{minipage}[c]{5.3cm}\raggedright\Large $1\,275$\end{minipage}\end{tcolorbox}\hfill\begin{tcolorbox}[enhanced,nobeforeafter,width=7.1cm,height=2.1cm,valign=center,colback=csui!6,colframe=csui,arc=2.5mm,boxrule=0.9pt,left=3mm]{\Large\bfseries\color{csui}b)}\quad\begin{minipage}[c]{5.3cm}\raggedright\Large $2\,550$\end{minipage}\end{tcolorbox}\par\vspace{2.5mm}
\begin{tcolorbox}[enhanced,nobeforeafter,width=7.1cm,height=2.1cm,valign=center,colback=csui!6,colframe=csui,arc=2.5mm,boxrule=0.9pt,left=3mm]{\Large\bfseries\color{csui}c)}\quad\begin{minipage}[c]{5.3cm}\raggedright\Large $2\,500$\end{minipage}\end{tcolorbox}\hfill\begin{tcolorbox}[enhanced,nobeforeafter,width=7.1cm,height=2.1cm,valign=center,colback=csui!6,colframe=csui,arc=2.5mm,boxrule=0.9pt,left=3mm]{\Large\bfseries\color{csui}d)}\quad\begin{minipage}[c]{5.3cm}\raggedright\Large $1\,225$\end{minipage}\end{tcolorbox}
\end{tcolorbox}
\newpage
\begin{tcolorbox}[enhanced,width=16cm,height=14cm,colback=white,colframe=csui,boxrule=2pt,arc=0mm,
left=6mm,right=6mm,top=0mm,bottom=4mm,sharp corners,
overlay={\fill[csui] (frame.north west) rectangle ([yshift=-1.25cm]frame.north east);
\node[anchor=west,text=white,font=\sffamily\bfseries\LARGE] at ([xshift=5mm,yshift=-0.62cm]frame.north west) {Question 8 B};
\node[anchor=east,text=white,font=\sffamily\itshape\large] at ([xshift=-5mm,yshift=-0.62cm]frame.north east) {Suites};}]
\vspace*{1.75cm}
\begin{minipage}[t][6.9cm][c]{\linewidth}\LARGE\centering
\begin{minipage}{0.95\linewidth}\raggedright\linespread{1.25}\selectfont Que vaut la somme $1+2+3+\dots+100$~?\end{minipage}
\end{minipage}

\vfill
\begin{tcolorbox}[enhanced,nobeforeafter,width=7.1cm,height=2.1cm,valign=center,colback=csui!6,colframe=csui,arc=2.5mm,boxrule=0.9pt,left=3mm]{\Large\bfseries\color{csui}a)}\quad\begin{minipage}[c]{5.3cm}\raggedright\Large $4\,950$\end{minipage}\end{tcolorbox}\hfill\begin{tcolorbox}[enhanced,nobeforeafter,width=7.1cm,height=2.1cm,valign=center,colback=csui!6,colframe=csui,arc=2.5mm,boxrule=0.9pt,left=3mm]{\Large\bfseries\color{csui}b)}\quad\begin{minipage}[c]{5.3cm}\raggedright\Large $5\,050$\end{minipage}\end{tcolorbox}\par\vspace{2.5mm}
\begin{tcolorbox}[enhanced,nobeforeafter,width=7.1cm,height=2.1cm,valign=center,colback=csui!6,colframe=csui,arc=2.5mm,boxrule=0.9pt,left=3mm]{\Large\bfseries\color{csui}c)}\quad\begin{minipage}[c]{5.3cm}\raggedright\Large $10\,100$\end{minipage}\end{tcolorbox}\hfill\begin{tcolorbox}[enhanced,nobeforeafter,width=7.1cm,height=2.1cm,valign=center,colback=csui!6,colframe=csui,arc=2.5mm,boxrule=0.9pt,left=3mm]{\Large\bfseries\color{csui}d)}\quad\begin{minipage}[c]{5.3cm}\raggedright\Large $10\,000$\end{minipage}\end{tcolorbox}
\end{tcolorbox}
\newpage
\begin{tcolorbox}[enhanced,width=16cm,height=14cm,colback=white,colframe=csui,boxrule=2pt,arc=0mm,
left=6mm,right=6mm,top=0mm,bottom=4mm,sharp corners,
overlay={\fill[csui] (frame.north west) rectangle ([yshift=-1.25cm]frame.north east);
\node[anchor=west,text=white,font=\sffamily\bfseries\LARGE] at ([xshift=5mm,yshift=-0.62cm]frame.north west) {Question 9 B};
\node[anchor=east,text=white,font=\sffamily\itshape\large] at ([xshift=-5mm,yshift=-0.62cm]frame.north east) {Suites};}]
\vspace*{1.75cm}
\begin{minipage}[t][6.9cm][c]{\linewidth}\LARGE\centering
\begin{minipage}{0.95\linewidth}\raggedright\linespread{1.25}\selectfont $(u_n)$ est une suite arithmétique avec $u_{1}=4$ et $u_{5}=14$. Quelle est sa raison~?\end{minipage}
\end{minipage}

\vfill
\begin{tcolorbox}[enhanced,nobeforeafter,width=7.1cm,height=2.1cm,valign=center,colback=csui!6,colframe=csui,arc=2.5mm,boxrule=0.9pt,left=3mm]{\Large\bfseries\color{csui}a)}\quad\begin{minipage}[c]{5.3cm}\raggedright\Large $-\dfrac{5}{2}$\end{minipage}\end{tcolorbox}\hfill\begin{tcolorbox}[enhanced,nobeforeafter,width=7.1cm,height=2.1cm,valign=center,colback=csui!6,colframe=csui,arc=2.5mm,boxrule=0.9pt,left=3mm]{\Large\bfseries\color{csui}b)}\quad\begin{minipage}[c]{5.3cm}\raggedright\Large $\dfrac{5}{3}$\end{minipage}\end{tcolorbox}\par\vspace{2.5mm}
\begin{tcolorbox}[enhanced,nobeforeafter,width=7.1cm,height=2.1cm,valign=center,colback=csui!6,colframe=csui,arc=2.5mm,boxrule=0.9pt,left=3mm]{\Large\bfseries\color{csui}c)}\quad\begin{minipage}[c]{5.3cm}\raggedright\Large $10$\end{minipage}\end{tcolorbox}\hfill\begin{tcolorbox}[enhanced,nobeforeafter,width=7.1cm,height=2.1cm,valign=center,colback=csui!6,colframe=csui,arc=2.5mm,boxrule=0.9pt,left=3mm]{\Large\bfseries\color{csui}d)}\quad\begin{minipage}[c]{5.3cm}\raggedright\Large $\dfrac{5}{2}$\end{minipage}\end{tcolorbox}
\end{tcolorbox}
\newpage
\begin{tcolorbox}[enhanced,width=16cm,height=14cm,colback=white,colframe=csui,boxrule=2pt,arc=0mm,
left=6mm,right=6mm,top=0mm,bottom=4mm,sharp corners,
overlay={\fill[csui] (frame.north west) rectangle ([yshift=-1.25cm]frame.north east);
\node[anchor=west,text=white,font=\sffamily\bfseries\LARGE] at ([xshift=5mm,yshift=-0.62cm]frame.north west) {Question 10 B};
\node[anchor=east,text=white,font=\sffamily\itshape\large] at ([xshift=-5mm,yshift=-0.62cm]frame.north east) {Suites};}]
\vspace*{1.75cm}
\begin{minipage}[t][6.9cm][c]{\linewidth}\LARGE\centering
\begin{minipage}{0.95\linewidth}\raggedright\linespread{1.25}\selectfont $(u_n)$ est une suite arithmétique avec $u_{0}=-3$ et $u_{4}=13$. Quelle est sa raison~?\end{minipage}
\end{minipage}

\vfill
\begin{tcolorbox}[enhanced,nobeforeafter,width=7.1cm,height=2.1cm,valign=center,colback=csui!6,colframe=csui,arc=2.5mm,boxrule=0.9pt,left=3mm]{\Large\bfseries\color{csui}a)}\quad\begin{minipage}[c]{5.3cm}\raggedright\Large $\dfrac{1}{4}$\end{minipage}\end{tcolorbox}\hfill\begin{tcolorbox}[enhanced,nobeforeafter,width=7.1cm,height=2.1cm,valign=center,colback=csui!6,colframe=csui,arc=2.5mm,boxrule=0.9pt,left=3mm]{\Large\bfseries\color{csui}b)}\quad\begin{minipage}[c]{5.3cm}\raggedright\Large $16$\end{minipage}\end{tcolorbox}\par\vspace{2.5mm}
\begin{tcolorbox}[enhanced,nobeforeafter,width=7.1cm,height=2.1cm,valign=center,colback=csui!6,colframe=csui,arc=2.5mm,boxrule=0.9pt,left=3mm]{\Large\bfseries\color{csui}c)}\quad\begin{minipage}[c]{5.3cm}\raggedright\Large $-4$\end{minipage}\end{tcolorbox}\hfill\begin{tcolorbox}[enhanced,nobeforeafter,width=7.1cm,height=2.1cm,valign=center,colback=csui!6,colframe=csui,arc=2.5mm,boxrule=0.9pt,left=3mm]{\Large\bfseries\color{csui}d)}\quad\begin{minipage}[c]{5.3cm}\raggedright\Large $4$\end{minipage}\end{tcolorbox}
\end{tcolorbox}
\newpage
\begin{tcolorbox}[enhanced,width=16cm,height=14cm,colback=white,colframe=csui,boxrule=2pt,arc=0mm,
left=6mm,right=6mm,top=0mm,bottom=4mm,sharp corners,
overlay={\fill[csui] (frame.north west) rectangle ([yshift=-1.25cm]frame.north east);
\node[anchor=west,text=white,font=\sffamily\bfseries\LARGE] at ([xshift=5mm,yshift=-0.62cm]frame.north west) {Question 11 B};
\node[anchor=east,text=white,font=\sffamily\itshape\large] at ([xshift=-5mm,yshift=-0.62cm]frame.north east) {Suites};}]
\vspace*{1.75cm}
\begin{minipage}[t][6.9cm][c]{\linewidth}\LARGE\centering
\begin{minipage}{0.95\linewidth}\raggedright\linespread{1.25}\selectfont Que vaut la somme $1+2+2^2+\dots+2^{9}$~?\end{minipage}
\end{minipage}

\vfill
\begin{tcolorbox}[enhanced,nobeforeafter,width=7.1cm,height=2.1cm,valign=center,colback=csui!6,colframe=csui,arc=2.5mm,boxrule=0.9pt,left=3mm]{\Large\bfseries\color{csui}a)}\quad\begin{minipage}[c]{5.3cm}\raggedright\Large $512$\end{minipage}\end{tcolorbox}\hfill\begin{tcolorbox}[enhanced,nobeforeafter,width=7.1cm,height=2.1cm,valign=center,colback=csui!6,colframe=csui,arc=2.5mm,boxrule=0.9pt,left=3mm]{\Large\bfseries\color{csui}b)}\quad\begin{minipage}[c]{5.3cm}\raggedright\Large $1\,023$\end{minipage}\end{tcolorbox}\par\vspace{2.5mm}
\begin{tcolorbox}[enhanced,nobeforeafter,width=7.1cm,height=2.1cm,valign=center,colback=csui!6,colframe=csui,arc=2.5mm,boxrule=0.9pt,left=3mm]{\Large\bfseries\color{csui}c)}\quad\begin{minipage}[c]{5.3cm}\raggedright\Large $1\,024$\end{minipage}\end{tcolorbox}\hfill\begin{tcolorbox}[enhanced,nobeforeafter,width=7.1cm,height=2.1cm,valign=center,colback=csui!6,colframe=csui,arc=2.5mm,boxrule=0.9pt,left=3mm]{\Large\bfseries\color{csui}d)}\quad\begin{minipage}[c]{5.3cm}\raggedright\Large $511$\end{minipage}\end{tcolorbox}
\end{tcolorbox}
\newpage
\begin{tcolorbox}[enhanced,width=16cm,height=14cm,colback=white,colframe=csui,boxrule=2pt,arc=0mm,
left=6mm,right=6mm,top=0mm,bottom=4mm,sharp corners,
overlay={\fill[csui] (frame.north west) rectangle ([yshift=-1.25cm]frame.north east);
\node[anchor=west,text=white,font=\sffamily\bfseries\LARGE] at ([xshift=5mm,yshift=-0.62cm]frame.north west) {Question 12 B};
\node[anchor=east,text=white,font=\sffamily\itshape\large] at ([xshift=-5mm,yshift=-0.62cm]frame.north east) {Suites};}]
\vspace*{1.75cm}
\begin{minipage}[t][6.9cm][c]{\linewidth}\LARGE\centering
\begin{minipage}{0.95\linewidth}\raggedright\linespread{1.25}\selectfont Que vaut la somme $1+5+5^2+\dots+5^{3}$~?\end{minipage}
\end{minipage}

\vfill
\begin{tcolorbox}[enhanced,nobeforeafter,width=7.1cm,height=2.1cm,valign=center,colback=csui!6,colframe=csui,arc=2.5mm,boxrule=0.9pt,left=3mm]{\Large\bfseries\color{csui}a)}\quad\begin{minipage}[c]{5.3cm}\raggedright\Large $31$\end{minipage}\end{tcolorbox}\hfill\begin{tcolorbox}[enhanced,nobeforeafter,width=7.1cm,height=2.1cm,valign=center,colback=csui!6,colframe=csui,arc=2.5mm,boxrule=0.9pt,left=3mm]{\Large\bfseries\color{csui}b)}\quad\begin{minipage}[c]{5.3cm}\raggedright\Large $625$\end{minipage}\end{tcolorbox}\par\vspace{2.5mm}
\begin{tcolorbox}[enhanced,nobeforeafter,width=7.1cm,height=2.1cm,valign=center,colback=csui!6,colframe=csui,arc=2.5mm,boxrule=0.9pt,left=3mm]{\Large\bfseries\color{csui}c)}\quad\begin{minipage}[c]{5.3cm}\raggedright\Large $156$\end{minipage}\end{tcolorbox}\hfill\begin{tcolorbox}[enhanced,nobeforeafter,width=7.1cm,height=2.1cm,valign=center,colback=csui!6,colframe=csui,arc=2.5mm,boxrule=0.9pt,left=3mm]{\Large\bfseries\color{csui}d)}\quad\begin{minipage}[c]{5.3cm}\raggedright\Large $624$\end{minipage}\end{tcolorbox}
\end{tcolorbox}
\newpage
\begin{tcolorbox}[enhanced,width=16cm,height=14cm,colback=white,colframe=csui,boxrule=2pt,arc=0mm,
left=6mm,right=6mm,top=0mm,bottom=4mm,sharp corners,
overlay={\fill[csui] (frame.north west) rectangle ([yshift=-1.25cm]frame.north east);
\node[anchor=west,text=white,font=\sffamily\bfseries\LARGE] at ([xshift=5mm,yshift=-0.62cm]frame.north west) {Question 13 B};
\node[anchor=east,text=white,font=\sffamily\itshape\large] at ([xshift=-5mm,yshift=-0.62cm]frame.north east) {Suites};}]
\vspace*{1.75cm}
\begin{minipage}[t][6.9cm][c]{\linewidth}\LARGE\centering
\begin{minipage}{0.95\linewidth}\raggedright\linespread{1.25}\selectfont Quel est le sens de variation de la suite définie pour tout entier naturel $n$ par $u_n=3\times2^n$~?\end{minipage}
\end{minipage}

\vfill
\begin{tcolorbox}[enhanced,nobeforeafter,width=7.1cm,height=2.1cm,valign=center,colback=csui!6,colframe=csui,arc=2.5mm,boxrule=0.9pt,left=3mm]{\Large\bfseries\color{csui}a)}\quad\begin{minipage}[c]{5.3cm}\raggedright\Large croissante\end{minipage}\end{tcolorbox}\hfill\begin{tcolorbox}[enhanced,nobeforeafter,width=7.1cm,height=2.1cm,valign=center,colback=csui!6,colframe=csui,arc=2.5mm,boxrule=0.9pt,left=3mm]{\Large\bfseries\color{csui}b)}\quad\begin{minipage}[c]{5.3cm}\raggedright\large ni croissante ni décroissante\end{minipage}\end{tcolorbox}\par\vspace{2.5mm}
\begin{tcolorbox}[enhanced,nobeforeafter,width=7.1cm,height=2.1cm,valign=center,colback=csui!6,colframe=csui,arc=2.5mm,boxrule=0.9pt,left=3mm]{\Large\bfseries\color{csui}c)}\quad\begin{minipage}[c]{5.3cm}\raggedright\Large constante\end{minipage}\end{tcolorbox}\hfill\begin{tcolorbox}[enhanced,nobeforeafter,width=7.1cm,height=2.1cm,valign=center,colback=csui!6,colframe=csui,arc=2.5mm,boxrule=0.9pt,left=3mm]{\Large\bfseries\color{csui}d)}\quad\begin{minipage}[c]{5.3cm}\raggedright\Large décroissante\end{minipage}\end{tcolorbox}
\end{tcolorbox}
\newpage
\begin{tcolorbox}[enhanced,width=16cm,height=14cm,colback=white,colframe=csui,boxrule=2pt,arc=0mm,
left=6mm,right=6mm,top=0mm,bottom=4mm,sharp corners,
overlay={\fill[csui] (frame.north west) rectangle ([yshift=-1.25cm]frame.north east);
\node[anchor=west,text=white,font=\sffamily\bfseries\LARGE] at ([xshift=5mm,yshift=-0.62cm]frame.north west) {Question 14 B};
\node[anchor=east,text=white,font=\sffamily\itshape\large] at ([xshift=-5mm,yshift=-0.62cm]frame.north east) {Suites};}]
\vspace*{1.75cm}
\begin{minipage}[t][6.9cm][c]{\linewidth}\LARGE\centering
\begin{minipage}{0.95\linewidth}\raggedright\linespread{1.25}\selectfont Quel est le sens de variation de la suite définie pour tout entier naturel $n$ par $u_n=4\times0{,}5^n$~?\end{minipage}
\end{minipage}

\vfill
\begin{tcolorbox}[enhanced,nobeforeafter,width=7.1cm,height=2.1cm,valign=center,colback=csui!6,colframe=csui,arc=2.5mm,boxrule=0.9pt,left=3mm]{\Large\bfseries\color{csui}a)}\quad\begin{minipage}[c]{5.3cm}\raggedright\large ni croissante ni décroissante\end{minipage}\end{tcolorbox}\hfill\begin{tcolorbox}[enhanced,nobeforeafter,width=7.1cm,height=2.1cm,valign=center,colback=csui!6,colframe=csui,arc=2.5mm,boxrule=0.9pt,left=3mm]{\Large\bfseries\color{csui}b)}\quad\begin{minipage}[c]{5.3cm}\raggedright\Large croissante\end{minipage}\end{tcolorbox}\par\vspace{2.5mm}
\begin{tcolorbox}[enhanced,nobeforeafter,width=7.1cm,height=2.1cm,valign=center,colback=csui!6,colframe=csui,arc=2.5mm,boxrule=0.9pt,left=3mm]{\Large\bfseries\color{csui}c)}\quad\begin{minipage}[c]{5.3cm}\raggedright\Large décroissante\end{minipage}\end{tcolorbox}\hfill\begin{tcolorbox}[enhanced,nobeforeafter,width=7.1cm,height=2.1cm,valign=center,colback=csui!6,colframe=csui,arc=2.5mm,boxrule=0.9pt,left=3mm]{\Large\bfseries\color{csui}d)}\quad\begin{minipage}[c]{5.3cm}\raggedright\Large constante\end{minipage}\end{tcolorbox}
\end{tcolorbox}
\newpage
\begin{tcolorbox}[enhanced,width=16cm,height=14cm,colback=white,colframe=csui,boxrule=2pt,arc=0mm,
left=6mm,right=6mm,top=0mm,bottom=4mm,sharp corners,
overlay={\fill[csui] (frame.north west) rectangle ([yshift=-1.25cm]frame.north east);
\node[anchor=west,text=white,font=\sffamily\bfseries\LARGE] at ([xshift=5mm,yshift=-0.62cm]frame.north west) {Question 15 B};
\node[anchor=east,text=white,font=\sffamily\itshape\large] at ([xshift=-5mm,yshift=-0.62cm]frame.north east) {Suites};}]
\vspace*{1.75cm}
\begin{minipage}[t][6.9cm][c]{\linewidth}\LARGE\centering
\begin{minipage}{0.95\linewidth}\raggedright\linespread{1.25}\selectfont Pour tout entier naturel $n$, $u_n=-n^2+6n$. Que vaut $u_{4}$~?\end{minipage}
\end{minipage}

\vfill
\begin{tcolorbox}[enhanced,nobeforeafter,width=7.1cm,height=2.1cm,valign=center,colback=csui!6,colframe=csui,arc=2.5mm,boxrule=0.9pt,left=3mm]{\Large\bfseries\color{csui}a)}\quad\begin{minipage}[c]{5.3cm}\raggedright\Large $5$\end{minipage}\end{tcolorbox}\hfill\begin{tcolorbox}[enhanced,nobeforeafter,width=7.1cm,height=2.1cm,valign=center,colback=csui!6,colframe=csui,arc=2.5mm,boxrule=0.9pt,left=3mm]{\Large\bfseries\color{csui}b)}\quad\begin{minipage}[c]{5.3cm}\raggedright\Large $9$\end{minipage}\end{tcolorbox}\par\vspace{2.5mm}
\begin{tcolorbox}[enhanced,nobeforeafter,width=7.1cm,height=2.1cm,valign=center,colback=csui!6,colframe=csui,arc=2.5mm,boxrule=0.9pt,left=3mm]{\Large\bfseries\color{csui}c)}\quad\begin{minipage}[c]{5.3cm}\raggedright\Large $-8$\end{minipage}\end{tcolorbox}\hfill\begin{tcolorbox}[enhanced,nobeforeafter,width=7.1cm,height=2.1cm,valign=center,colback=csui!6,colframe=csui,arc=2.5mm,boxrule=0.9pt,left=3mm]{\Large\bfseries\color{csui}d)}\quad\begin{minipage}[c]{5.3cm}\raggedright\Large $8$\end{minipage}\end{tcolorbox}
\end{tcolorbox}
\newpage
\begin{tcolorbox}[enhanced,width=16cm,height=14cm,colback=white,colframe=csui,boxrule=2pt,arc=0mm,
left=6mm,right=6mm,top=0mm,bottom=4mm,sharp corners,
overlay={\fill[csui] (frame.north west) rectangle ([yshift=-1.25cm]frame.north east);
\node[anchor=west,text=white,font=\sffamily\bfseries\LARGE] at ([xshift=5mm,yshift=-0.62cm]frame.north west) {Question 16 B};
\node[anchor=east,text=white,font=\sffamily\itshape\large] at ([xshift=-5mm,yshift=-0.62cm]frame.north east) {Suites};}]
\vspace*{1.75cm}
\begin{minipage}[t][6.9cm][c]{\linewidth}\LARGE\centering
\begin{minipage}{0.95\linewidth}\raggedright\linespread{1.25}\selectfont Pour tout entier naturel $n$, $u_n=3n^2-2n+1$. Que vaut $u_{2}$~?\end{minipage}
\end{minipage}

\vfill
\begin{tcolorbox}[enhanced,nobeforeafter,width=7.1cm,height=2.1cm,valign=center,colback=csui!6,colframe=csui,arc=2.5mm,boxrule=0.9pt,left=3mm]{\Large\bfseries\color{csui}a)}\quad\begin{minipage}[c]{5.3cm}\raggedright\Large $9$\end{minipage}\end{tcolorbox}\hfill\begin{tcolorbox}[enhanced,nobeforeafter,width=7.1cm,height=2.1cm,valign=center,colback=csui!6,colframe=csui,arc=2.5mm,boxrule=0.9pt,left=3mm]{\Large\bfseries\color{csui}b)}\quad\begin{minipage}[c]{5.3cm}\raggedright\Large $-9$\end{minipage}\end{tcolorbox}\par\vspace{2.5mm}
\begin{tcolorbox}[enhanced,nobeforeafter,width=7.1cm,height=2.1cm,valign=center,colback=csui!6,colframe=csui,arc=2.5mm,boxrule=0.9pt,left=3mm]{\Large\bfseries\color{csui}c)}\quad\begin{minipage}[c]{5.3cm}\raggedright\Large $2$\end{minipage}\end{tcolorbox}\hfill\begin{tcolorbox}[enhanced,nobeforeafter,width=7.1cm,height=2.1cm,valign=center,colback=csui!6,colframe=csui,arc=2.5mm,boxrule=0.9pt,left=3mm]{\Large\bfseries\color{csui}d)}\quad\begin{minipage}[c]{5.3cm}\raggedright\Large $22$\end{minipage}\end{tcolorbox}
\end{tcolorbox}
\newpage
\begin{tcolorbox}[enhanced,width=16cm,height=14cm,colback=white,colframe=cdeg,boxrule=2pt,arc=0mm,
left=6mm,right=6mm,top=0mm,bottom=4mm,sharp corners,
overlay={\fill[cdeg] (frame.north west) rectangle ([yshift=-1.25cm]frame.north east);
\node[anchor=west,text=white,font=\sffamily\bfseries\LARGE] at ([xshift=5mm,yshift=-0.62cm]frame.north west) {Question 17 B};
\node[anchor=east,text=white,font=\sffamily\itshape\large] at ([xshift=-5mm,yshift=-0.62cm]frame.north east) {Second degré};}]
\vspace*{1.75cm}
\begin{minipage}[t][6.9cm][c]{\linewidth}\LARGE\centering
\begin{minipage}{0.95\linewidth}\raggedright\linespread{1.25}\selectfont Quel est le discriminant du trinôme $x^2+4x+4$~?\end{minipage}
\end{minipage}

\vfill
\begin{tcolorbox}[enhanced,nobeforeafter,width=7.1cm,height=2.1cm,valign=center,colback=cdeg!6,colframe=cdeg,arc=2.5mm,boxrule=0.9pt,left=3mm]{\Large\bfseries\color{cdeg}a)}\quad\begin{minipage}[c]{5.3cm}\raggedright\Large $1$\end{minipage}\end{tcolorbox}\hfill\begin{tcolorbox}[enhanced,nobeforeafter,width=7.1cm,height=2.1cm,valign=center,colback=cdeg!6,colframe=cdeg,arc=2.5mm,boxrule=0.9pt,left=3mm]{\Large\bfseries\color{cdeg}b)}\quad\begin{minipage}[c]{5.3cm}\raggedright\Large $-12$\end{minipage}\end{tcolorbox}\par\vspace{2.5mm}
\begin{tcolorbox}[enhanced,nobeforeafter,width=7.1cm,height=2.1cm,valign=center,colback=cdeg!6,colframe=cdeg,arc=2.5mm,boxrule=0.9pt,left=3mm]{\Large\bfseries\color{cdeg}c)}\quad\begin{minipage}[c]{5.3cm}\raggedright\Large $0$\end{minipage}\end{tcolorbox}\hfill\begin{tcolorbox}[enhanced,nobeforeafter,width=7.1cm,height=2.1cm,valign=center,colback=cdeg!6,colframe=cdeg,arc=2.5mm,boxrule=0.9pt,left=3mm]{\Large\bfseries\color{cdeg}d)}\quad\begin{minipage}[c]{5.3cm}\raggedright\Large $32$\end{minipage}\end{tcolorbox}
\end{tcolorbox}
\newpage
\begin{tcolorbox}[enhanced,width=16cm,height=14cm,colback=white,colframe=cdeg,boxrule=2pt,arc=0mm,
left=6mm,right=6mm,top=0mm,bottom=4mm,sharp corners,
overlay={\fill[cdeg] (frame.north west) rectangle ([yshift=-1.25cm]frame.north east);
\node[anchor=west,text=white,font=\sffamily\bfseries\LARGE] at ([xshift=5mm,yshift=-0.62cm]frame.north west) {Question 18 B};
\node[anchor=east,text=white,font=\sffamily\itshape\large] at ([xshift=-5mm,yshift=-0.62cm]frame.north east) {Second degré};}]
\vspace*{1.75cm}
\begin{minipage}[t][6.9cm][c]{\linewidth}\LARGE\centering
\begin{minipage}{0.95\linewidth}\raggedright\linespread{1.25}\selectfont Quel est le discriminant du trinôme $3x^2-2x+5$~?\end{minipage}
\end{minipage}

\vfill
\begin{tcolorbox}[enhanced,nobeforeafter,width=7.1cm,height=2.1cm,valign=center,colback=cdeg!6,colframe=cdeg,arc=2.5mm,boxrule=0.9pt,left=3mm]{\Large\bfseries\color{cdeg}a)}\quad\begin{minipage}[c]{5.3cm}\raggedright\Large $56$\end{minipage}\end{tcolorbox}\hfill\begin{tcolorbox}[enhanced,nobeforeafter,width=7.1cm,height=2.1cm,valign=center,colback=cdeg!6,colframe=cdeg,arc=2.5mm,boxrule=0.9pt,left=3mm]{\Large\bfseries\color{cdeg}b)}\quad\begin{minipage}[c]{5.3cm}\raggedright\Large $64$\end{minipage}\end{tcolorbox}\par\vspace{2.5mm}
\begin{tcolorbox}[enhanced,nobeforeafter,width=7.1cm,height=2.1cm,valign=center,colback=cdeg!6,colframe=cdeg,arc=2.5mm,boxrule=0.9pt,left=3mm]{\Large\bfseries\color{cdeg}c)}\quad\begin{minipage}[c]{5.3cm}\raggedright\Large $-56$\end{minipage}\end{tcolorbox}\hfill\begin{tcolorbox}[enhanced,nobeforeafter,width=7.1cm,height=2.1cm,valign=center,colback=cdeg!6,colframe=cdeg,arc=2.5mm,boxrule=0.9pt,left=3mm]{\Large\bfseries\color{cdeg}d)}\quad\begin{minipage}[c]{5.3cm}\raggedright\Large $-62$\end{minipage}\end{tcolorbox}
\end{tcolorbox}
\newpage
\begin{tcolorbox}[enhanced,width=16cm,height=14cm,colback=white,colframe=cdeg,boxrule=2pt,arc=0mm,
left=6mm,right=6mm,top=0mm,bottom=4mm,sharp corners,
overlay={\fill[cdeg] (frame.north west) rectangle ([yshift=-1.25cm]frame.north east);
\node[anchor=west,text=white,font=\sffamily\bfseries\LARGE] at ([xshift=5mm,yshift=-0.62cm]frame.north west) {Question 19 B};
\node[anchor=east,text=white,font=\sffamily\itshape\large] at ([xshift=-5mm,yshift=-0.62cm]frame.north east) {Second degré};}]
\vspace*{1.75cm}
\begin{minipage}[t][6.9cm][c]{\linewidth}\LARGE\centering
\begin{minipage}{0.95\linewidth}\raggedright\linespread{1.25}\selectfont Combien de racines réelles le trinôme $2x^2+5x+1$ possède-t-il~?\end{minipage}
\end{minipage}

\vfill
\begin{tcolorbox}[enhanced,nobeforeafter,width=7.1cm,height=2.1cm,valign=center,colback=cdeg!6,colframe=cdeg,arc=2.5mm,boxrule=0.9pt,left=3mm]{\Large\bfseries\color{cdeg}a)}\quad\begin{minipage}[c]{5.3cm}\raggedright\Large $0$\end{minipage}\end{tcolorbox}\hfill\begin{tcolorbox}[enhanced,nobeforeafter,width=7.1cm,height=2.1cm,valign=center,colback=cdeg!6,colframe=cdeg,arc=2.5mm,boxrule=0.9pt,left=3mm]{\Large\bfseries\color{cdeg}b)}\quad\begin{minipage}[c]{5.3cm}\raggedright\Large une infinité\end{minipage}\end{tcolorbox}\par\vspace{2.5mm}
\begin{tcolorbox}[enhanced,nobeforeafter,width=7.1cm,height=2.1cm,valign=center,colback=cdeg!6,colframe=cdeg,arc=2.5mm,boxrule=0.9pt,left=3mm]{\Large\bfseries\color{cdeg}c)}\quad\begin{minipage}[c]{5.3cm}\raggedright\Large $1$\end{minipage}\end{tcolorbox}\hfill\begin{tcolorbox}[enhanced,nobeforeafter,width=7.1cm,height=2.1cm,valign=center,colback=cdeg!6,colframe=cdeg,arc=2.5mm,boxrule=0.9pt,left=3mm]{\Large\bfseries\color{cdeg}d)}\quad\begin{minipage}[c]{5.3cm}\raggedright\Large $2$\end{minipage}\end{tcolorbox}
\end{tcolorbox}
\newpage
\begin{tcolorbox}[enhanced,width=16cm,height=14cm,colback=white,colframe=cdeg,boxrule=2pt,arc=0mm,
left=6mm,right=6mm,top=0mm,bottom=4mm,sharp corners,
overlay={\fill[cdeg] (frame.north west) rectangle ([yshift=-1.25cm]frame.north east);
\node[anchor=west,text=white,font=\sffamily\bfseries\LARGE] at ([xshift=5mm,yshift=-0.62cm]frame.north west) {Question 20 B};
\node[anchor=east,text=white,font=\sffamily\itshape\large] at ([xshift=-5mm,yshift=-0.62cm]frame.north east) {Second degré};}]
\vspace*{1.75cm}
\begin{minipage}[t][6.9cm][c]{\linewidth}\LARGE\centering
\begin{minipage}{0.95\linewidth}\raggedright\linespread{1.25}\selectfont Combien de racines réelles le trinôme $-3x^2+2x+4$ possède-t-il~?\end{minipage}
\end{minipage}

\vfill
\begin{tcolorbox}[enhanced,nobeforeafter,width=7.1cm,height=2.1cm,valign=center,colback=cdeg!6,colframe=cdeg,arc=2.5mm,boxrule=0.9pt,left=3mm]{\Large\bfseries\color{cdeg}a)}\quad\begin{minipage}[c]{5.3cm}\raggedright\Large $0$\end{minipage}\end{tcolorbox}\hfill\begin{tcolorbox}[enhanced,nobeforeafter,width=7.1cm,height=2.1cm,valign=center,colback=cdeg!6,colframe=cdeg,arc=2.5mm,boxrule=0.9pt,left=3mm]{\Large\bfseries\color{cdeg}b)}\quad\begin{minipage}[c]{5.3cm}\raggedright\Large une infinité\end{minipage}\end{tcolorbox}\par\vspace{2.5mm}
\begin{tcolorbox}[enhanced,nobeforeafter,width=7.1cm,height=2.1cm,valign=center,colback=cdeg!6,colframe=cdeg,arc=2.5mm,boxrule=0.9pt,left=3mm]{\Large\bfseries\color{cdeg}c)}\quad\begin{minipage}[c]{5.3cm}\raggedright\Large $1$\end{minipage}\end{tcolorbox}\hfill\begin{tcolorbox}[enhanced,nobeforeafter,width=7.1cm,height=2.1cm,valign=center,colback=cdeg!6,colframe=cdeg,arc=2.5mm,boxrule=0.9pt,left=3mm]{\Large\bfseries\color{cdeg}d)}\quad\begin{minipage}[c]{5.3cm}\raggedright\Large $2$\end{minipage}\end{tcolorbox}
\end{tcolorbox}
\newpage
\begin{tcolorbox}[enhanced,width=16cm,height=14cm,colback=white,colframe=cdeg,boxrule=2pt,arc=0mm,
left=6mm,right=6mm,top=0mm,bottom=4mm,sharp corners,
overlay={\fill[cdeg] (frame.north west) rectangle ([yshift=-1.25cm]frame.north east);
\node[anchor=west,text=white,font=\sffamily\bfseries\LARGE] at ([xshift=5mm,yshift=-0.62cm]frame.north west) {Question 21 B};
\node[anchor=east,text=white,font=\sffamily\itshape\large] at ([xshift=-5mm,yshift=-0.62cm]frame.north east) {Second degré};}]
\vspace*{1.75cm}
\begin{minipage}[t][6.9cm][c]{\linewidth}\LARGE\centering
\begin{minipage}{0.95\linewidth}\raggedright\linespread{1.25}\selectfont Quelles sont les racines du trinôme $2x^2+4x-6$~?\end{minipage}
\end{minipage}

\vfill
\begin{tcolorbox}[enhanced,nobeforeafter,width=7.1cm,height=2.1cm,valign=center,colback=cdeg!6,colframe=cdeg,arc=2.5mm,boxrule=0.9pt,left=3mm]{\Large\bfseries\color{cdeg}a)}\quad\begin{minipage}[c]{5.3cm}\raggedright\Large $\{-3\,;\,1\}$\end{minipage}\end{tcolorbox}\hfill\begin{tcolorbox}[enhanced,nobeforeafter,width=7.1cm,height=2.1cm,valign=center,colback=cdeg!6,colframe=cdeg,arc=2.5mm,boxrule=0.9pt,left=3mm]{\Large\bfseries\color{cdeg}b)}\quad\begin{minipage}[c]{5.3cm}\raggedright\Large $\{-1\,;\,3\}$\end{minipage}\end{tcolorbox}\par\vspace{2.5mm}
\begin{tcolorbox}[enhanced,nobeforeafter,width=7.1cm,height=2.1cm,valign=center,colback=cdeg!6,colframe=cdeg,arc=2.5mm,boxrule=0.9pt,left=3mm]{\Large\bfseries\color{cdeg}c)}\quad\begin{minipage}[c]{5.3cm}\raggedright\Large $\{1\,;\,3\}$\end{minipage}\end{tcolorbox}\hfill\begin{tcolorbox}[enhanced,nobeforeafter,width=7.1cm,height=2.1cm,valign=center,colback=cdeg!6,colframe=cdeg,arc=2.5mm,boxrule=0.9pt,left=3mm]{\Large\bfseries\color{cdeg}d)}\quad\begin{minipage}[c]{5.3cm}\raggedright\Large $\{-3\,;\,-2\}$\end{minipage}\end{tcolorbox}
\end{tcolorbox}
\newpage
\begin{tcolorbox}[enhanced,width=16cm,height=14cm,colback=white,colframe=cdeg,boxrule=2pt,arc=0mm,
left=6mm,right=6mm,top=0mm,bottom=4mm,sharp corners,
overlay={\fill[cdeg] (frame.north west) rectangle ([yshift=-1.25cm]frame.north east);
\node[anchor=west,text=white,font=\sffamily\bfseries\LARGE] at ([xshift=5mm,yshift=-0.62cm]frame.north west) {Question 22 B};
\node[anchor=east,text=white,font=\sffamily\itshape\large] at ([xshift=-5mm,yshift=-0.62cm]frame.north east) {Second degré};}]
\vspace*{1.75cm}
\begin{minipage}[t][6.9cm][c]{\linewidth}\LARGE\centering
\begin{minipage}{0.95\linewidth}\raggedright\linespread{1.25}\selectfont Quelles sont les racines du trinôme $x^2+7x+10$~?\end{minipage}
\end{minipage}

\vfill
\begin{tcolorbox}[enhanced,nobeforeafter,width=7.1cm,height=2.1cm,valign=center,colback=cdeg!6,colframe=cdeg,arc=2.5mm,boxrule=0.9pt,left=3mm]{\Large\bfseries\color{cdeg}a)}\quad\begin{minipage}[c]{5.3cm}\raggedright\Large $\{-5\,;\,-2\}$\end{minipage}\end{tcolorbox}\hfill\begin{tcolorbox}[enhanced,nobeforeafter,width=7.1cm,height=2.1cm,valign=center,colback=cdeg!6,colframe=cdeg,arc=2.5mm,boxrule=0.9pt,left=3mm]{\Large\bfseries\color{cdeg}b)}\quad\begin{minipage}[c]{5.3cm}\raggedright\Large $\{-5\,;\,2\}$\end{minipage}\end{tcolorbox}\par\vspace{2.5mm}
\begin{tcolorbox}[enhanced,nobeforeafter,width=7.1cm,height=2.1cm,valign=center,colback=cdeg!6,colframe=cdeg,arc=2.5mm,boxrule=0.9pt,left=3mm]{\Large\bfseries\color{cdeg}c)}\quad\begin{minipage}[c]{5.3cm}\raggedright\Large $\{-7\,;\,10\}$\end{minipage}\end{tcolorbox}\hfill\begin{tcolorbox}[enhanced,nobeforeafter,width=7.1cm,height=2.1cm,valign=center,colback=cdeg!6,colframe=cdeg,arc=2.5mm,boxrule=0.9pt,left=3mm]{\Large\bfseries\color{cdeg}d)}\quad\begin{minipage}[c]{5.3cm}\raggedright\Large $\{2\,;\,5\}$\end{minipage}\end{tcolorbox}
\end{tcolorbox}
\newpage
\begin{tcolorbox}[enhanced,width=16cm,height=14cm,colback=white,colframe=cdeg,boxrule=2pt,arc=0mm,
left=6mm,right=6mm,top=0mm,bottom=4mm,sharp corners,
overlay={\fill[cdeg] (frame.north west) rectangle ([yshift=-1.25cm]frame.north east);
\node[anchor=west,text=white,font=\sffamily\bfseries\LARGE] at ([xshift=5mm,yshift=-0.62cm]frame.north west) {Question 23 B};
\node[anchor=east,text=white,font=\sffamily\itshape\large] at ([xshift=-5mm,yshift=-0.62cm]frame.north east) {Second degré};}]
\vspace*{1.75cm}
\begin{minipage}[t][6.9cm][c]{\linewidth}\LARGE\centering
\begin{minipage}{0.95\linewidth}\raggedright\linespread{1.25}\selectfont Quelle est la forme canonique de $x^2+2x-3$~?\end{minipage}
\end{minipage}

\vfill
\begin{tcolorbox}[enhanced,nobeforeafter,width=7.1cm,height=2.1cm,valign=center,colback=cdeg!6,colframe=cdeg,arc=2.5mm,boxrule=0.9pt,left=3mm]{\Large\bfseries\color{cdeg}a)}\quad\begin{minipage}[c]{5.3cm}\raggedright\Large $(x+1)^2-3$\end{minipage}\end{tcolorbox}\hfill\begin{tcolorbox}[enhanced,nobeforeafter,width=7.1cm,height=2.1cm,valign=center,colback=cdeg!6,colframe=cdeg,arc=2.5mm,boxrule=0.9pt,left=3mm]{\Large\bfseries\color{cdeg}b)}\quad\begin{minipage}[c]{5.3cm}\raggedright\Large $(x+2)^2-4$\end{minipage}\end{tcolorbox}\par\vspace{2.5mm}
\begin{tcolorbox}[enhanced,nobeforeafter,width=7.1cm,height=2.1cm,valign=center,colback=cdeg!6,colframe=cdeg,arc=2.5mm,boxrule=0.9pt,left=3mm]{\Large\bfseries\color{cdeg}c)}\quad\begin{minipage}[c]{5.3cm}\raggedright\Large $(x-1)^2-4$\end{minipage}\end{tcolorbox}\hfill\begin{tcolorbox}[enhanced,nobeforeafter,width=7.1cm,height=2.1cm,valign=center,colback=cdeg!6,colframe=cdeg,arc=2.5mm,boxrule=0.9pt,left=3mm]{\Large\bfseries\color{cdeg}d)}\quad\begin{minipage}[c]{5.3cm}\raggedright\Large $(x+1)^2-4$\end{minipage}\end{tcolorbox}
\end{tcolorbox}
\newpage
\begin{tcolorbox}[enhanced,width=16cm,height=14cm,colback=white,colframe=cdeg,boxrule=2pt,arc=0mm,
left=6mm,right=6mm,top=0mm,bottom=4mm,sharp corners,
overlay={\fill[cdeg] (frame.north west) rectangle ([yshift=-1.25cm]frame.north east);
\node[anchor=west,text=white,font=\sffamily\bfseries\LARGE] at ([xshift=5mm,yshift=-0.62cm]frame.north west) {Question 24 B};
\node[anchor=east,text=white,font=\sffamily\itshape\large] at ([xshift=-5mm,yshift=-0.62cm]frame.north east) {Second degré};}]
\vspace*{1.75cm}
\begin{minipage}[t][6.9cm][c]{\linewidth}\LARGE\centering
\begin{minipage}{0.95\linewidth}\raggedright\linespread{1.25}\selectfont Quelle est la forme canonique de $x^2-10x+30$~?\end{minipage}
\end{minipage}

\vfill
\begin{tcolorbox}[enhanced,nobeforeafter,width=7.1cm,height=2.1cm,valign=center,colback=cdeg!6,colframe=cdeg,arc=2.5mm,boxrule=0.9pt,left=3mm]{\Large\bfseries\color{cdeg}a)}\quad\begin{minipage}[c]{5.3cm}\raggedright\Large $(x-10)^2+5$\end{minipage}\end{tcolorbox}\hfill\begin{tcolorbox}[enhanced,nobeforeafter,width=7.1cm,height=2.1cm,valign=center,colback=cdeg!6,colframe=cdeg,arc=2.5mm,boxrule=0.9pt,left=3mm]{\Large\bfseries\color{cdeg}b)}\quad\begin{minipage}[c]{5.3cm}\raggedright\Large $(x+5)^2+5$\end{minipage}\end{tcolorbox}\par\vspace{2.5mm}
\begin{tcolorbox}[enhanced,nobeforeafter,width=7.1cm,height=2.1cm,valign=center,colback=cdeg!6,colframe=cdeg,arc=2.5mm,boxrule=0.9pt,left=3mm]{\Large\bfseries\color{cdeg}c)}\quad\begin{minipage}[c]{5.3cm}\raggedright\Large $(x-5)^2+30$\end{minipage}\end{tcolorbox}\hfill\begin{tcolorbox}[enhanced,nobeforeafter,width=7.1cm,height=2.1cm,valign=center,colback=cdeg!6,colframe=cdeg,arc=2.5mm,boxrule=0.9pt,left=3mm]{\Large\bfseries\color{cdeg}d)}\quad\begin{minipage}[c]{5.3cm}\raggedright\Large $(x-5)^2+5$\end{minipage}\end{tcolorbox}
\end{tcolorbox}
\newpage
\begin{tcolorbox}[enhanced,width=16cm,height=14cm,colback=white,colframe=cdeg,boxrule=2pt,arc=0mm,
left=6mm,right=6mm,top=0mm,bottom=4mm,sharp corners,
overlay={\fill[cdeg] (frame.north west) rectangle ([yshift=-1.25cm]frame.north east);
\node[anchor=west,text=white,font=\sffamily\bfseries\LARGE] at ([xshift=5mm,yshift=-0.62cm]frame.north west) {Question 25 B};
\node[anchor=east,text=white,font=\sffamily\itshape\large] at ([xshift=-5mm,yshift=-0.62cm]frame.north east) {Second degré};}]
\vspace*{1.75cm}
\begin{minipage}[t][6.9cm][c]{\linewidth}\LARGE\centering
\begin{minipage}{0.95\linewidth}\raggedright\linespread{1.25}\selectfont Quelles sont les coordonnées du sommet de la parabole d'équation $y=-x^2+6x-5$~?\end{minipage}
\end{minipage}

\vfill
\begin{tcolorbox}[enhanced,nobeforeafter,width=7.1cm,height=2.1cm,valign=center,colback=cdeg!6,colframe=cdeg,arc=2.5mm,boxrule=0.9pt,left=3mm]{\Large\bfseries\color{cdeg}a)}\quad\begin{minipage}[c]{5.3cm}\raggedright\Large $(3\,;\,4)$\end{minipage}\end{tcolorbox}\hfill\begin{tcolorbox}[enhanced,nobeforeafter,width=7.1cm,height=2.1cm,valign=center,colback=cdeg!6,colframe=cdeg,arc=2.5mm,boxrule=0.9pt,left=3mm]{\Large\bfseries\color{cdeg}b)}\quad\begin{minipage}[c]{5.3cm}\raggedright\Large $(-3\,;\,4)$\end{minipage}\end{tcolorbox}\par\vspace{2.5mm}
\begin{tcolorbox}[enhanced,nobeforeafter,width=7.1cm,height=2.1cm,valign=center,colback=cdeg!6,colframe=cdeg,arc=2.5mm,boxrule=0.9pt,left=3mm]{\Large\bfseries\color{cdeg}c)}\quad\begin{minipage}[c]{5.3cm}\raggedright\Large $(3\,;\,-5)$\end{minipage}\end{tcolorbox}\hfill\begin{tcolorbox}[enhanced,nobeforeafter,width=7.1cm,height=2.1cm,valign=center,colback=cdeg!6,colframe=cdeg,arc=2.5mm,boxrule=0.9pt,left=3mm]{\Large\bfseries\color{cdeg}d)}\quad\begin{minipage}[c]{5.3cm}\raggedright\Large $(3\,;\,-4)$\end{minipage}\end{tcolorbox}
\end{tcolorbox}
\newpage
\begin{tcolorbox}[enhanced,width=16cm,height=14cm,colback=white,colframe=cdeg,boxrule=2pt,arc=0mm,
left=6mm,right=6mm,top=0mm,bottom=4mm,sharp corners,
overlay={\fill[cdeg] (frame.north west) rectangle ([yshift=-1.25cm]frame.north east);
\node[anchor=west,text=white,font=\sffamily\bfseries\LARGE] at ([xshift=5mm,yshift=-0.62cm]frame.north west) {Question 26 B};
\node[anchor=east,text=white,font=\sffamily\itshape\large] at ([xshift=-5mm,yshift=-0.62cm]frame.north east) {Second degré};}]
\vspace*{1.75cm}
\begin{minipage}[t][6.9cm][c]{\linewidth}\LARGE\centering
\begin{minipage}{0.95\linewidth}\raggedright\linespread{1.25}\selectfont Quelles sont les coordonnées du sommet de la parabole d'équation $y=x^2+2x-8$~?\end{minipage}
\end{minipage}

\vfill
\begin{tcolorbox}[enhanced,nobeforeafter,width=7.1cm,height=2.1cm,valign=center,colback=cdeg!6,colframe=cdeg,arc=2.5mm,boxrule=0.9pt,left=3mm]{\Large\bfseries\color{cdeg}a)}\quad\begin{minipage}[c]{5.3cm}\raggedright\Large $(1\,;\,-9)$\end{minipage}\end{tcolorbox}\hfill\begin{tcolorbox}[enhanced,nobeforeafter,width=7.1cm,height=2.1cm,valign=center,colback=cdeg!6,colframe=cdeg,arc=2.5mm,boxrule=0.9pt,left=3mm]{\Large\bfseries\color{cdeg}b)}\quad\begin{minipage}[c]{5.3cm}\raggedright\Large $(-1\,;\,9)$\end{minipage}\end{tcolorbox}\par\vspace{2.5mm}
\begin{tcolorbox}[enhanced,nobeforeafter,width=7.1cm,height=2.1cm,valign=center,colback=cdeg!6,colframe=cdeg,arc=2.5mm,boxrule=0.9pt,left=3mm]{\Large\bfseries\color{cdeg}c)}\quad\begin{minipage}[c]{5.3cm}\raggedright\Large $(-1\,;\,-8)$\end{minipage}\end{tcolorbox}\hfill\begin{tcolorbox}[enhanced,nobeforeafter,width=7.1cm,height=2.1cm,valign=center,colback=cdeg!6,colframe=cdeg,arc=2.5mm,boxrule=0.9pt,left=3mm]{\Large\bfseries\color{cdeg}d)}\quad\begin{minipage}[c]{5.3cm}\raggedright\Large $(-1\,;\,-9)$\end{minipage}\end{tcolorbox}
\end{tcolorbox}
\newpage
\begin{tcolorbox}[enhanced,width=16cm,height=14cm,colback=white,colframe=cdeg,boxrule=2pt,arc=0mm,
left=6mm,right=6mm,top=0mm,bottom=4mm,sharp corners,
overlay={\fill[cdeg] (frame.north west) rectangle ([yshift=-1.25cm]frame.north east);
\node[anchor=west,text=white,font=\sffamily\bfseries\LARGE] at ([xshift=5mm,yshift=-0.62cm]frame.north west) {Question 27 B};
\node[anchor=east,text=white,font=\sffamily\itshape\large] at ([xshift=-5mm,yshift=-0.62cm]frame.north east) {Second degré};}]
\vspace*{1.75cm}
\begin{minipage}[t][6.9cm][c]{\linewidth}\LARGE\centering
\begin{minipage}{0.95\linewidth}\raggedright\linespread{1.25}\selectfont Quel est l'ensemble des solutions de l'inéquation $-x^2+3x>0$~?\end{minipage}
\end{minipage}

\vfill
\begin{tcolorbox}[enhanced,nobeforeafter,width=7.1cm,height=2.1cm,valign=center,colback=cdeg!6,colframe=cdeg,arc=2.5mm,boxrule=0.9pt,left=3mm]{\Large\bfseries\color{cdeg}a)}\quad\begin{minipage}[c]{5.3cm}\raggedright\large $]-\infty\,;\,0]\cup[3\,;\,+\infty[$\end{minipage}\end{tcolorbox}\hfill\begin{tcolorbox}[enhanced,nobeforeafter,width=7.1cm,height=2.1cm,valign=center,colback=cdeg!6,colframe=cdeg,arc=2.5mm,boxrule=0.9pt,left=3mm]{\Large\bfseries\color{cdeg}b)}\quad\begin{minipage}[c]{5.3cm}\raggedright\Large $]0\,;\,3[$\end{minipage}\end{tcolorbox}\par\vspace{2.5mm}
\begin{tcolorbox}[enhanced,nobeforeafter,width=7.1cm,height=2.1cm,valign=center,colback=cdeg!6,colframe=cdeg,arc=2.5mm,boxrule=0.9pt,left=3mm]{\Large\bfseries\color{cdeg}c)}\quad\begin{minipage}[c]{5.3cm}\raggedright\large $]-\infty\,;\,0[\cup]3\,;\,+\infty[$\end{minipage}\end{tcolorbox}\hfill\begin{tcolorbox}[enhanced,nobeforeafter,width=7.1cm,height=2.1cm,valign=center,colback=cdeg!6,colframe=cdeg,arc=2.5mm,boxrule=0.9pt,left=3mm]{\Large\bfseries\color{cdeg}d)}\quad\begin{minipage}[c]{5.3cm}\raggedright\Large $[0\,;\,3]$\end{minipage}\end{tcolorbox}
\end{tcolorbox}
\newpage
\begin{tcolorbox}[enhanced,width=16cm,height=14cm,colback=white,colframe=cdeg,boxrule=2pt,arc=0mm,
left=6mm,right=6mm,top=0mm,bottom=4mm,sharp corners,
overlay={\fill[cdeg] (frame.north west) rectangle ([yshift=-1.25cm]frame.north east);
\node[anchor=west,text=white,font=\sffamily\bfseries\LARGE] at ([xshift=5mm,yshift=-0.62cm]frame.north west) {Question 28 B};
\node[anchor=east,text=white,font=\sffamily\itshape\large] at ([xshift=-5mm,yshift=-0.62cm]frame.north east) {Second degré};}]
\vspace*{1.75cm}
\begin{minipage}[t][6.9cm][c]{\linewidth}\LARGE\centering
\begin{minipage}{0.95\linewidth}\raggedright\linespread{1.25}\selectfont Quel est l'ensemble des solutions de l'inéquation $2x^2+4x-6\leqslant 0$~?\end{minipage}
\end{minipage}

\vfill
\begin{tcolorbox}[enhanced,nobeforeafter,width=7.1cm,height=2.1cm,valign=center,colback=cdeg!6,colframe=cdeg,arc=2.5mm,boxrule=0.9pt,left=3mm]{\Large\bfseries\color{cdeg}a)}\quad\begin{minipage}[c]{5.3cm}\raggedright\large $]-\infty\,;\,-3[\cup]1\,;\,+\infty[$\end{minipage}\end{tcolorbox}\hfill\begin{tcolorbox}[enhanced,nobeforeafter,width=7.1cm,height=2.1cm,valign=center,colback=cdeg!6,colframe=cdeg,arc=2.5mm,boxrule=0.9pt,left=3mm]{\Large\bfseries\color{cdeg}b)}\quad\begin{minipage}[c]{5.3cm}\raggedright\Large $]-3\,;\,1[$\end{minipage}\end{tcolorbox}\par\vspace{2.5mm}
\begin{tcolorbox}[enhanced,nobeforeafter,width=7.1cm,height=2.1cm,valign=center,colback=cdeg!6,colframe=cdeg,arc=2.5mm,boxrule=0.9pt,left=3mm]{\Large\bfseries\color{cdeg}c)}\quad\begin{minipage}[c]{5.3cm}\raggedright\large $]-\infty\,;\,-3]\cup[1\,;\,+\infty[$\end{minipage}\end{tcolorbox}\hfill\begin{tcolorbox}[enhanced,nobeforeafter,width=7.1cm,height=2.1cm,valign=center,colback=cdeg!6,colframe=cdeg,arc=2.5mm,boxrule=0.9pt,left=3mm]{\Large\bfseries\color{cdeg}d)}\quad\begin{minipage}[c]{5.3cm}\raggedright\Large $[-3\,;\,1]$\end{minipage}\end{tcolorbox}
\end{tcolorbox}
\newpage
\begin{tcolorbox}[enhanced,width=16cm,height=14cm,colback=white,colframe=cdeg,boxrule=2pt,arc=0mm,
left=6mm,right=6mm,top=0mm,bottom=4mm,sharp corners,
overlay={\fill[cdeg] (frame.north west) rectangle ([yshift=-1.25cm]frame.north east);
\node[anchor=west,text=white,font=\sffamily\bfseries\LARGE] at ([xshift=5mm,yshift=-0.62cm]frame.north west) {Question 29 B};
\node[anchor=east,text=white,font=\sffamily\itshape\large] at ([xshift=-5mm,yshift=-0.62cm]frame.north east) {Second degré};}]
\vspace*{1.75cm}
\begin{minipage}[t][6.9cm][c]{\linewidth}\LARGE\centering
\begin{minipage}{0.95\linewidth}\raggedright\linespread{1.25}\selectfont Le trinôme $x^2+4x-5$ a deux racines. Que vaut leur somme~?\end{minipage}
\end{minipage}

\vfill
\begin{tcolorbox}[enhanced,nobeforeafter,width=7.1cm,height=2.1cm,valign=center,colback=cdeg!6,colframe=cdeg,arc=2.5mm,boxrule=0.9pt,left=3mm]{\Large\bfseries\color{cdeg}a)}\quad\begin{minipage}[c]{5.3cm}\raggedright\Large $-4$\end{minipage}\end{tcolorbox}\hfill\begin{tcolorbox}[enhanced,nobeforeafter,width=7.1cm,height=2.1cm,valign=center,colback=cdeg!6,colframe=cdeg,arc=2.5mm,boxrule=0.9pt,left=3mm]{\Large\bfseries\color{cdeg}b)}\quad\begin{minipage}[c]{5.3cm}\raggedright\Large $-5$\end{minipage}\end{tcolorbox}\par\vspace{2.5mm}
\begin{tcolorbox}[enhanced,nobeforeafter,width=7.1cm,height=2.1cm,valign=center,colback=cdeg!6,colframe=cdeg,arc=2.5mm,boxrule=0.9pt,left=3mm]{\Large\bfseries\color{cdeg}c)}\quad\begin{minipage}[c]{5.3cm}\raggedright\Large $-8$\end{minipage}\end{tcolorbox}\hfill\begin{tcolorbox}[enhanced,nobeforeafter,width=7.1cm,height=2.1cm,valign=center,colback=cdeg!6,colframe=cdeg,arc=2.5mm,boxrule=0.9pt,left=3mm]{\Large\bfseries\color{cdeg}d)}\quad\begin{minipage}[c]{5.3cm}\raggedright\Large $4$\end{minipage}\end{tcolorbox}
\end{tcolorbox}
\newpage
\begin{tcolorbox}[enhanced,width=16cm,height=14cm,colback=white,colframe=cdeg,boxrule=2pt,arc=0mm,
left=6mm,right=6mm,top=0mm,bottom=4mm,sharp corners,
overlay={\fill[cdeg] (frame.north west) rectangle ([yshift=-1.25cm]frame.north east);
\node[anchor=west,text=white,font=\sffamily\bfseries\LARGE] at ([xshift=5mm,yshift=-0.62cm]frame.north west) {Question 30 B};
\node[anchor=east,text=white,font=\sffamily\itshape\large] at ([xshift=-5mm,yshift=-0.62cm]frame.north east) {Second degré};}]
\vspace*{1.75cm}
\begin{minipage}[t][6.9cm][c]{\linewidth}\LARGE\centering
\begin{minipage}{0.95\linewidth}\raggedright\linespread{1.25}\selectfont Le trinôme $3x^2-9x+2$ a deux racines. Que vaut leur produit~?\end{minipage}
\end{minipage}

\vfill
\begin{tcolorbox}[enhanced,nobeforeafter,width=7.1cm,height=2.1cm,valign=center,colback=cdeg!6,colframe=cdeg,arc=2.5mm,boxrule=0.9pt,left=3mm]{\Large\bfseries\color{cdeg}a)}\quad\begin{minipage}[c]{5.3cm}\raggedright\Large $\dfrac{2}{3}$\end{minipage}\end{tcolorbox}\hfill\begin{tcolorbox}[enhanced,nobeforeafter,width=7.1cm,height=2.1cm,valign=center,colback=cdeg!6,colframe=cdeg,arc=2.5mm,boxrule=0.9pt,left=3mm]{\Large\bfseries\color{cdeg}b)}\quad\begin{minipage}[c]{5.3cm}\raggedright\Large $3$\end{minipage}\end{tcolorbox}\par\vspace{2.5mm}
\begin{tcolorbox}[enhanced,nobeforeafter,width=7.1cm,height=2.1cm,valign=center,colback=cdeg!6,colframe=cdeg,arc=2.5mm,boxrule=0.9pt,left=3mm]{\Large\bfseries\color{cdeg}c)}\quad\begin{minipage}[c]{5.3cm}\raggedright\Large $-3$\end{minipage}\end{tcolorbox}\hfill\begin{tcolorbox}[enhanced,nobeforeafter,width=7.1cm,height=2.1cm,valign=center,colback=cdeg!6,colframe=cdeg,arc=2.5mm,boxrule=0.9pt,left=3mm]{\Large\bfseries\color{cdeg}d)}\quad\begin{minipage}[c]{5.3cm}\raggedright\Large $-\dfrac{2}{3}$\end{minipage}\end{tcolorbox}
\end{tcolorbox}
\newpage
\begin{tcolorbox}[enhanced,width=16cm,height=14cm,colback=white,colframe=cdeg,boxrule=2pt,arc=0mm,
left=6mm,right=6mm,top=0mm,bottom=4mm,sharp corners,
overlay={\fill[cdeg] (frame.north west) rectangle ([yshift=-1.25cm]frame.north east);
\node[anchor=west,text=white,font=\sffamily\bfseries\LARGE] at ([xshift=5mm,yshift=-0.62cm]frame.north west) {Question 31 B};
\node[anchor=east,text=white,font=\sffamily\itshape\large] at ([xshift=-5mm,yshift=-0.62cm]frame.north east) {Second degré};}]
\vspace*{1.75cm}
\begin{minipage}[t][6.9cm][c]{\linewidth}\LARGE\centering
\begin{minipage}{0.95\linewidth}\raggedright\linespread{1.25}\selectfont Le trinôme $3x^2-3x-6$ a pour racines $-1$ et $2$. Quelle est sa forme factorisée~?\end{minipage}
\end{minipage}

\vfill
\begin{tcolorbox}[enhanced,nobeforeafter,width=7.1cm,height=2.1cm,valign=center,colback=cdeg!6,colframe=cdeg,arc=2.5mm,boxrule=0.9pt,left=3mm]{\Large\bfseries\color{cdeg}a)}\quad\begin{minipage}[c]{5.3cm}\raggedright\Large $3(x+1)(x-2)$\end{minipage}\end{tcolorbox}\hfill\begin{tcolorbox}[enhanced,nobeforeafter,width=7.1cm,height=2.1cm,valign=center,colback=cdeg!6,colframe=cdeg,arc=2.5mm,boxrule=0.9pt,left=3mm]{\Large\bfseries\color{cdeg}b)}\quad\begin{minipage}[c]{5.3cm}\raggedright\Large $(x+1)(x-2)$\end{minipage}\end{tcolorbox}\par\vspace{2.5mm}
\begin{tcolorbox}[enhanced,nobeforeafter,width=7.1cm,height=2.1cm,valign=center,colback=cdeg!6,colframe=cdeg,arc=2.5mm,boxrule=0.9pt,left=3mm]{\Large\bfseries\color{cdeg}c)}\quad\begin{minipage}[c]{5.3cm}\raggedright\Large $3(x-1)(x+2)$\end{minipage}\end{tcolorbox}\hfill\begin{tcolorbox}[enhanced,nobeforeafter,width=7.1cm,height=2.1cm,valign=center,colback=cdeg!6,colframe=cdeg,arc=2.5mm,boxrule=0.9pt,left=3mm]{\Large\bfseries\color{cdeg}d)}\quad\begin{minipage}[c]{5.3cm}\raggedright\Large $-(x+1)(x-2)$\end{minipage}\end{tcolorbox}
\end{tcolorbox}
\newpage
\begin{tcolorbox}[enhanced,width=16cm,height=14cm,colback=white,colframe=cdeg,boxrule=2pt,arc=0mm,
left=6mm,right=6mm,top=0mm,bottom=4mm,sharp corners,
overlay={\fill[cdeg] (frame.north west) rectangle ([yshift=-1.25cm]frame.north east);
\node[anchor=west,text=white,font=\sffamily\bfseries\LARGE] at ([xshift=5mm,yshift=-0.62cm]frame.north west) {Question 32 B};
\node[anchor=east,text=white,font=\sffamily\itshape\large] at ([xshift=-5mm,yshift=-0.62cm]frame.north east) {Second degré};}]
\vspace*{1.75cm}
\begin{minipage}[t][6.9cm][c]{\linewidth}\LARGE\centering
\begin{minipage}{0.95\linewidth}\raggedright\linespread{1.25}\selectfont Le trinôme $-x^2+5x$ a pour racines $0$ et $5$. Quelle est sa forme factorisée~?\end{minipage}
\end{minipage}

\vfill
\begin{tcolorbox}[enhanced,nobeforeafter,width=7.1cm,height=2.1cm,valign=center,colback=cdeg!6,colframe=cdeg,arc=2.5mm,boxrule=0.9pt,left=3mm]{\Large\bfseries\color{cdeg}a)}\quad\begin{minipage}[c]{5.3cm}\raggedright\Large $-x(x-6)$\end{minipage}\end{tcolorbox}\hfill\begin{tcolorbox}[enhanced,nobeforeafter,width=7.1cm,height=2.1cm,valign=center,colback=cdeg!6,colframe=cdeg,arc=2.5mm,boxrule=0.9pt,left=3mm]{\Large\bfseries\color{cdeg}b)}\quad\begin{minipage}[c]{5.3cm}\raggedright\Large $-x(x-5)$\end{minipage}\end{tcolorbox}\par\vspace{2.5mm}
\begin{tcolorbox}[enhanced,nobeforeafter,width=7.1cm,height=2.1cm,valign=center,colback=cdeg!6,colframe=cdeg,arc=2.5mm,boxrule=0.9pt,left=3mm]{\Large\bfseries\color{cdeg}c)}\quad\begin{minipage}[c]{5.3cm}\raggedright\Large $x(x-5)$\end{minipage}\end{tcolorbox}\hfill\begin{tcolorbox}[enhanced,nobeforeafter,width=7.1cm,height=2.1cm,valign=center,colback=cdeg!6,colframe=cdeg,arc=2.5mm,boxrule=0.9pt,left=3mm]{\Large\bfseries\color{cdeg}d)}\quad\begin{minipage}[c]{5.3cm}\raggedright\Large $-x(x+5)$\end{minipage}\end{tcolorbox}
\end{tcolorbox}
\newpage
\begin{tcolorbox}[enhanced,width=16cm,height=14cm,colback=white,colframe=cder,boxrule=2pt,arc=0mm,
left=6mm,right=6mm,top=0mm,bottom=4mm,sharp corners,
overlay={\fill[cder] (frame.north west) rectangle ([yshift=-1.25cm]frame.north east);
\node[anchor=west,text=white,font=\sffamily\bfseries\LARGE] at ([xshift=5mm,yshift=-0.62cm]frame.north west) {Question 33 B};
\node[anchor=east,text=white,font=\sffamily\itshape\large] at ([xshift=-5mm,yshift=-0.62cm]frame.north east) {Dérivation};}]
\vspace*{1.75cm}
\begin{minipage}[t][6.9cm][c]{\linewidth}\LARGE\centering
\begin{minipage}{0.95\linewidth}\raggedright\linespread{1.25}\selectfont Soit $f(x)=-x^3+3x^2+4$. Que vaut $f'(3)$~?\end{minipage}
\end{minipage}

\vfill
\begin{tcolorbox}[enhanced,nobeforeafter,width=7.1cm,height=2.1cm,valign=center,colback=cder!6,colframe=cder,arc=2.5mm,boxrule=0.9pt,left=3mm]{\Large\bfseries\color{cder}a)}\quad\begin{minipage}[c]{5.3cm}\raggedright\Large $-18$\end{minipage}\end{tcolorbox}\hfill\begin{tcolorbox}[enhanced,nobeforeafter,width=7.1cm,height=2.1cm,valign=center,colback=cder!6,colframe=cder,arc=2.5mm,boxrule=0.9pt,left=3mm]{\Large\bfseries\color{cder}b)}\quad\begin{minipage}[c]{5.3cm}\raggedright\Large $4$\end{minipage}\end{tcolorbox}\par\vspace{2.5mm}
\begin{tcolorbox}[enhanced,nobeforeafter,width=7.1cm,height=2.1cm,valign=center,colback=cder!6,colframe=cder,arc=2.5mm,boxrule=0.9pt,left=3mm]{\Large\bfseries\color{cder}c)}\quad\begin{minipage}[c]{5.3cm}\raggedright\Large $0$\end{minipage}\end{tcolorbox}\hfill\begin{tcolorbox}[enhanced,nobeforeafter,width=7.1cm,height=2.1cm,valign=center,colback=cder!6,colframe=cder,arc=2.5mm,boxrule=0.9pt,left=3mm]{\Large\bfseries\color{cder}d)}\quad\begin{minipage}[c]{5.3cm}\raggedright\Large $-9$\end{minipage}\end{tcolorbox}
\end{tcolorbox}
\newpage
\begin{tcolorbox}[enhanced,width=16cm,height=14cm,colback=white,colframe=cder,boxrule=2pt,arc=0mm,
left=6mm,right=6mm,top=0mm,bottom=4mm,sharp corners,
overlay={\fill[cder] (frame.north west) rectangle ([yshift=-1.25cm]frame.north east);
\node[anchor=west,text=white,font=\sffamily\bfseries\LARGE] at ([xshift=5mm,yshift=-0.62cm]frame.north west) {Question 34 B};
\node[anchor=east,text=white,font=\sffamily\itshape\large] at ([xshift=-5mm,yshift=-0.62cm]frame.north east) {Dérivation};}]
\vspace*{1.75cm}
\begin{minipage}[t][6.9cm][c]{\linewidth}\LARGE\centering
\begin{minipage}{0.95\linewidth}\raggedright\linespread{1.25}\selectfont Soit $f(x)=x^3+x^2+x+1$. Que vaut $f'(-1)$~?\end{minipage}
\end{minipage}

\vfill
\begin{tcolorbox}[enhanced,nobeforeafter,width=7.1cm,height=2.1cm,valign=center,colback=cder!6,colframe=cder,arc=2.5mm,boxrule=0.9pt,left=3mm]{\Large\bfseries\color{cder}a)}\quad\begin{minipage}[c]{5.3cm}\raggedright\Large $2$\end{minipage}\end{tcolorbox}\hfill\begin{tcolorbox}[enhanced,nobeforeafter,width=7.1cm,height=2.1cm,valign=center,colback=cder!6,colframe=cder,arc=2.5mm,boxrule=0.9pt,left=3mm]{\Large\bfseries\color{cder}b)}\quad\begin{minipage}[c]{5.3cm}\raggedright\Large $0$\end{minipage}\end{tcolorbox}\par\vspace{2.5mm}
\begin{tcolorbox}[enhanced,nobeforeafter,width=7.1cm,height=2.1cm,valign=center,colback=cder!6,colframe=cder,arc=2.5mm,boxrule=0.9pt,left=3mm]{\Large\bfseries\color{cder}c)}\quad\begin{minipage}[c]{5.3cm}\raggedright\Large $1$\end{minipage}\end{tcolorbox}\hfill\begin{tcolorbox}[enhanced,nobeforeafter,width=7.1cm,height=2.1cm,valign=center,colback=cder!6,colframe=cder,arc=2.5mm,boxrule=0.9pt,left=3mm]{\Large\bfseries\color{cder}d)}\quad\begin{minipage}[c]{5.3cm}\raggedright\Large $3$\end{minipage}\end{tcolorbox}
\end{tcolorbox}
\newpage
\begin{tcolorbox}[enhanced,width=16cm,height=14cm,colback=white,colframe=cder,boxrule=2pt,arc=0mm,
left=6mm,right=6mm,top=0mm,bottom=4mm,sharp corners,
overlay={\fill[cder] (frame.north west) rectangle ([yshift=-1.25cm]frame.north east);
\node[anchor=west,text=white,font=\sffamily\bfseries\LARGE] at ([xshift=5mm,yshift=-0.62cm]frame.north west) {Question 35 B};
\node[anchor=east,text=white,font=\sffamily\itshape\large] at ([xshift=-5mm,yshift=-0.62cm]frame.north east) {Dérivation};}]
\vspace*{1.75cm}
\begin{minipage}[t][6.9cm][c]{\linewidth}\LARGE\centering
\begin{minipage}{0.95\linewidth}\raggedright\linespread{1.25}\selectfont Quelle est l'équation de la tangente à la courbe de $f(x)=x^2-4$ au point d'abscisse $2$~?\end{minipage}
\end{minipage}

\vfill
\begin{tcolorbox}[enhanced,nobeforeafter,width=7.1cm,height=2.1cm,valign=center,colback=cder!6,colframe=cder,arc=2.5mm,boxrule=0.9pt,left=3mm]{\Large\bfseries\color{cder}a)}\quad\begin{minipage}[c]{5.3cm}\raggedright\Large $y=4x$\end{minipage}\end{tcolorbox}\hfill\begin{tcolorbox}[enhanced,nobeforeafter,width=7.1cm,height=2.1cm,valign=center,colback=cder!6,colframe=cder,arc=2.5mm,boxrule=0.9pt,left=3mm]{\Large\bfseries\color{cder}b)}\quad\begin{minipage}[c]{5.3cm}\raggedright\Large $y=4$\end{minipage}\end{tcolorbox}\par\vspace{2.5mm}
\begin{tcolorbox}[enhanced,nobeforeafter,width=7.1cm,height=2.1cm,valign=center,colback=cder!6,colframe=cder,arc=2.5mm,boxrule=0.9pt,left=3mm]{\Large\bfseries\color{cder}c)}\quad\begin{minipage}[c]{5.3cm}\raggedright\Large $y=4x+8$\end{minipage}\end{tcolorbox}\hfill\begin{tcolorbox}[enhanced,nobeforeafter,width=7.1cm,height=2.1cm,valign=center,colback=cder!6,colframe=cder,arc=2.5mm,boxrule=0.9pt,left=3mm]{\Large\bfseries\color{cder}d)}\quad\begin{minipage}[c]{5.3cm}\raggedright\Large $y=4x-8$\end{minipage}\end{tcolorbox}
\end{tcolorbox}
\newpage
\begin{tcolorbox}[enhanced,width=16cm,height=14cm,colback=white,colframe=cder,boxrule=2pt,arc=0mm,
left=6mm,right=6mm,top=0mm,bottom=4mm,sharp corners,
overlay={\fill[cder] (frame.north west) rectangle ([yshift=-1.25cm]frame.north east);
\node[anchor=west,text=white,font=\sffamily\bfseries\LARGE] at ([xshift=5mm,yshift=-0.62cm]frame.north west) {Question 36 B};
\node[anchor=east,text=white,font=\sffamily\itshape\large] at ([xshift=-5mm,yshift=-0.62cm]frame.north east) {Dérivation};}]
\vspace*{1.75cm}
\begin{minipage}[t][6.9cm][c]{\linewidth}\LARGE\centering
\begin{minipage}{0.95\linewidth}\raggedright\linespread{1.25}\selectfont Quelle est l'équation de la tangente à la courbe de $f(x)=-x^2+4x+1$ au point d'abscisse $1$~?\end{minipage}
\end{minipage}

\vfill
\begin{tcolorbox}[enhanced,nobeforeafter,width=7.1cm,height=2.1cm,valign=center,colback=cder!6,colframe=cder,arc=2.5mm,boxrule=0.9pt,left=3mm]{\Large\bfseries\color{cder}a)}\quad\begin{minipage}[c]{5.3cm}\raggedright\Large $y=2x+4$\end{minipage}\end{tcolorbox}\hfill\begin{tcolorbox}[enhanced,nobeforeafter,width=7.1cm,height=2.1cm,valign=center,colback=cder!6,colframe=cder,arc=2.5mm,boxrule=0.9pt,left=3mm]{\Large\bfseries\color{cder}b)}\quad\begin{minipage}[c]{5.3cm}\raggedright\Large $y=2x+6$\end{minipage}\end{tcolorbox}\par\vspace{2.5mm}
\begin{tcolorbox}[enhanced,nobeforeafter,width=7.1cm,height=2.1cm,valign=center,colback=cder!6,colframe=cder,arc=2.5mm,boxrule=0.9pt,left=3mm]{\Large\bfseries\color{cder}c)}\quad\begin{minipage}[c]{5.3cm}\raggedright\Large $y=2x+2$\end{minipage}\end{tcolorbox}\hfill\begin{tcolorbox}[enhanced,nobeforeafter,width=7.1cm,height=2.1cm,valign=center,colback=cder!6,colframe=cder,arc=2.5mm,boxrule=0.9pt,left=3mm]{\Large\bfseries\color{cder}d)}\quad\begin{minipage}[c]{5.3cm}\raggedright\Large $y=4x-2$\end{minipage}\end{tcolorbox}
\end{tcolorbox}
\newpage
\begin{tcolorbox}[enhanced,width=16cm,height=14cm,colback=white,colframe=cder,boxrule=2pt,arc=0mm,
left=6mm,right=6mm,top=0mm,bottom=4mm,sharp corners,
overlay={\fill[cder] (frame.north west) rectangle ([yshift=-1.25cm]frame.north east);
\node[anchor=west,text=white,font=\sffamily\bfseries\LARGE] at ([xshift=5mm,yshift=-0.62cm]frame.north west) {Question 37 B};
\node[anchor=east,text=white,font=\sffamily\itshape\large] at ([xshift=-5mm,yshift=-0.62cm]frame.north east) {Dérivation};}]
\vspace*{1.75cm}
\begin{minipage}[t][6.9cm][c]{\linewidth}\LARGE\centering
\begin{minipage}{0.95\linewidth}\raggedright\linespread{1.25}\selectfont Soit $f(x)=(-x+2)(4x+1)$. Que vaut $f'(2)$~?\end{minipage}
\end{minipage}

\vfill
\begin{tcolorbox}[enhanced,nobeforeafter,width=7.1cm,height=2.1cm,valign=center,colback=cder!6,colframe=cder,arc=2.5mm,boxrule=0.9pt,left=3mm]{\Large\bfseries\color{cder}a)}\quad\begin{minipage}[c]{5.3cm}\raggedright\Large $-7$\end{minipage}\end{tcolorbox}\hfill\begin{tcolorbox}[enhanced,nobeforeafter,width=7.1cm,height=2.1cm,valign=center,colback=cder!6,colframe=cder,arc=2.5mm,boxrule=0.9pt,left=3mm]{\Large\bfseries\color{cder}b)}\quad\begin{minipage}[c]{5.3cm}\raggedright\Large $0$\end{minipage}\end{tcolorbox}\par\vspace{2.5mm}
\begin{tcolorbox}[enhanced,nobeforeafter,width=7.1cm,height=2.1cm,valign=center,colback=cder!6,colframe=cder,arc=2.5mm,boxrule=0.9pt,left=3mm]{\Large\bfseries\color{cder}c)}\quad\begin{minipage}[c]{5.3cm}\raggedright\Large $-4$\end{minipage}\end{tcolorbox}\hfill\begin{tcolorbox}[enhanced,nobeforeafter,width=7.1cm,height=2.1cm,valign=center,colback=cder!6,colframe=cder,arc=2.5mm,boxrule=0.9pt,left=3mm]{\Large\bfseries\color{cder}d)}\quad\begin{minipage}[c]{5.3cm}\raggedright\Large $-9$\end{minipage}\end{tcolorbox}
\end{tcolorbox}
\newpage
\begin{tcolorbox}[enhanced,width=16cm,height=14cm,colback=white,colframe=cder,boxrule=2pt,arc=0mm,
left=6mm,right=6mm,top=0mm,bottom=4mm,sharp corners,
overlay={\fill[cder] (frame.north west) rectangle ([yshift=-1.25cm]frame.north east);
\node[anchor=west,text=white,font=\sffamily\bfseries\LARGE] at ([xshift=5mm,yshift=-0.62cm]frame.north west) {Question 38 B};
\node[anchor=east,text=white,font=\sffamily\itshape\large] at ([xshift=-5mm,yshift=-0.62cm]frame.north east) {Dérivation};}]
\vspace*{1.75cm}
\begin{minipage}[t][6.9cm][c]{\linewidth}\LARGE\centering
\begin{minipage}{0.95\linewidth}\raggedright\linespread{1.25}\selectfont Soit $f(x)=(3x)(x-2)$. Que vaut $f'(1)$~?\end{minipage}
\end{minipage}

\vfill
\begin{tcolorbox}[enhanced,nobeforeafter,width=7.1cm,height=2.1cm,valign=center,colback=cder!6,colframe=cder,arc=2.5mm,boxrule=0.9pt,left=3mm]{\Large\bfseries\color{cder}a)}\quad\begin{minipage}[c]{5.3cm}\raggedright\Large $-3$\end{minipage}\end{tcolorbox}\hfill\begin{tcolorbox}[enhanced,nobeforeafter,width=7.1cm,height=2.1cm,valign=center,colback=cder!6,colframe=cder,arc=2.5mm,boxrule=0.9pt,left=3mm]{\Large\bfseries\color{cder}b)}\quad\begin{minipage}[c]{5.3cm}\raggedright\Large $2$\end{minipage}\end{tcolorbox}\par\vspace{2.5mm}
\begin{tcolorbox}[enhanced,nobeforeafter,width=7.1cm,height=2.1cm,valign=center,colback=cder!6,colframe=cder,arc=2.5mm,boxrule=0.9pt,left=3mm]{\Large\bfseries\color{cder}c)}\quad\begin{minipage}[c]{5.3cm}\raggedright\Large $0$\end{minipage}\end{tcolorbox}\hfill\begin{tcolorbox}[enhanced,nobeforeafter,width=7.1cm,height=2.1cm,valign=center,colback=cder!6,colframe=cder,arc=2.5mm,boxrule=0.9pt,left=3mm]{\Large\bfseries\color{cder}d)}\quad\begin{minipage}[c]{5.3cm}\raggedright\Large $3$\end{minipage}\end{tcolorbox}
\end{tcolorbox}
\newpage
\begin{tcolorbox}[enhanced,width=16cm,height=14cm,colback=white,colframe=cder,boxrule=2pt,arc=0mm,
left=6mm,right=6mm,top=0mm,bottom=4mm,sharp corners,
overlay={\fill[cder] (frame.north west) rectangle ([yshift=-1.25cm]frame.north east);
\node[anchor=west,text=white,font=\sffamily\bfseries\LARGE] at ([xshift=5mm,yshift=-0.62cm]frame.north west) {Question 39 B};
\node[anchor=east,text=white,font=\sffamily\itshape\large] at ([xshift=-5mm,yshift=-0.62cm]frame.north east) {Dérivation};}]
\vspace*{1.75cm}
\begin{minipage}[t][6.9cm][c]{\linewidth}\LARGE\centering
\begin{minipage}{0.95\linewidth}\raggedright\linespread{1.25}\selectfont Soit $f(x)=\dfrac{5}{x}$ sur $\mathbb{R}^*$. Que vaut $f'(2)$~?\end{minipage}
\end{minipage}

\vfill
\begin{tcolorbox}[enhanced,nobeforeafter,width=7.1cm,height=2.1cm,valign=center,colback=cder!6,colframe=cder,arc=2.5mm,boxrule=0.9pt,left=3mm]{\Large\bfseries\color{cder}a)}\quad\begin{minipage}[c]{5.3cm}\raggedright\Large $-\dfrac{5}{2}$\end{minipage}\end{tcolorbox}\hfill\begin{tcolorbox}[enhanced,nobeforeafter,width=7.1cm,height=2.1cm,valign=center,colback=cder!6,colframe=cder,arc=2.5mm,boxrule=0.9pt,left=3mm]{\Large\bfseries\color{cder}b)}\quad\begin{minipage}[c]{5.3cm}\raggedright\Large $-\dfrac{5}{4}$\end{minipage}\end{tcolorbox}\par\vspace{2.5mm}
\begin{tcolorbox}[enhanced,nobeforeafter,width=7.1cm,height=2.1cm,valign=center,colback=cder!6,colframe=cder,arc=2.5mm,boxrule=0.9pt,left=3mm]{\Large\bfseries\color{cder}c)}\quad\begin{minipage}[c]{5.3cm}\raggedright\Large $10$\end{minipage}\end{tcolorbox}\hfill\begin{tcolorbox}[enhanced,nobeforeafter,width=7.1cm,height=2.1cm,valign=center,colback=cder!6,colframe=cder,arc=2.5mm,boxrule=0.9pt,left=3mm]{\Large\bfseries\color{cder}d)}\quad\begin{minipage}[c]{5.3cm}\raggedright\Large $\dfrac{5}{4}$\end{minipage}\end{tcolorbox}
\end{tcolorbox}
\newpage
\begin{tcolorbox}[enhanced,width=16cm,height=14cm,colback=white,colframe=cder,boxrule=2pt,arc=0mm,
left=6mm,right=6mm,top=0mm,bottom=4mm,sharp corners,
overlay={\fill[cder] (frame.north west) rectangle ([yshift=-1.25cm]frame.north east);
\node[anchor=west,text=white,font=\sffamily\bfseries\LARGE] at ([xshift=5mm,yshift=-0.62cm]frame.north west) {Question 40 B};
\node[anchor=east,text=white,font=\sffamily\itshape\large] at ([xshift=-5mm,yshift=-0.62cm]frame.north east) {Dérivation};}]
\vspace*{1.75cm}
\begin{minipage}[t][6.9cm][c]{\linewidth}\LARGE\centering
\begin{minipage}{0.95\linewidth}\raggedright\linespread{1.25}\selectfont Soit $f(x)=\dfrac{-4}{x}$ sur $\mathbb{R}^*$. Que vaut $f'(2)$~?\end{minipage}
\end{minipage}

\vfill
\begin{tcolorbox}[enhanced,nobeforeafter,width=7.1cm,height=2.1cm,valign=center,colback=cder!6,colframe=cder,arc=2.5mm,boxrule=0.9pt,left=3mm]{\Large\bfseries\color{cder}a)}\quad\begin{minipage}[c]{5.3cm}\raggedright\Large $1$\end{minipage}\end{tcolorbox}\hfill\begin{tcolorbox}[enhanced,nobeforeafter,width=7.1cm,height=2.1cm,valign=center,colback=cder!6,colframe=cder,arc=2.5mm,boxrule=0.9pt,left=3mm]{\Large\bfseries\color{cder}b)}\quad\begin{minipage}[c]{5.3cm}\raggedright\Large $-1$\end{minipage}\end{tcolorbox}\par\vspace{2.5mm}
\begin{tcolorbox}[enhanced,nobeforeafter,width=7.1cm,height=2.1cm,valign=center,colback=cder!6,colframe=cder,arc=2.5mm,boxrule=0.9pt,left=3mm]{\Large\bfseries\color{cder}c)}\quad\begin{minipage}[c]{5.3cm}\raggedright\Large $2$\end{minipage}\end{tcolorbox}\hfill\begin{tcolorbox}[enhanced,nobeforeafter,width=7.1cm,height=2.1cm,valign=center,colback=cder!6,colframe=cder,arc=2.5mm,boxrule=0.9pt,left=3mm]{\Large\bfseries\color{cder}d)}\quad\begin{minipage}[c]{5.3cm}\raggedright\Large $-8$\end{minipage}\end{tcolorbox}
\end{tcolorbox}
\newpage
\begin{tcolorbox}[enhanced,width=16cm,height=14cm,colback=white,colframe=cder,boxrule=2pt,arc=0mm,
left=6mm,right=6mm,top=0mm,bottom=4mm,sharp corners,
overlay={\fill[cder] (frame.north west) rectangle ([yshift=-1.25cm]frame.north east);
\node[anchor=west,text=white,font=\sffamily\bfseries\LARGE] at ([xshift=5mm,yshift=-0.62cm]frame.north west) {Question 41 B};
\node[anchor=east,text=white,font=\sffamily\itshape\large] at ([xshift=-5mm,yshift=-0.62cm]frame.north east) {Dérivation};}]
\vspace*{1.75cm}
\begin{minipage}[t][6.9cm][c]{\linewidth}\LARGE\centering
\begin{minipage}{0.95\linewidth}\raggedright\linespread{1.25}\selectfont Soit $f(x)=\sqrt{x}$. Que vaut $f'(16)$~?\end{minipage}
\end{minipage}

\vfill
\begin{tcolorbox}[enhanced,nobeforeafter,width=7.1cm,height=2.1cm,valign=center,colback=cder!6,colframe=cder,arc=2.5mm,boxrule=0.9pt,left=3mm]{\Large\bfseries\color{cder}a)}\quad\begin{minipage}[c]{5.3cm}\raggedright\Large $8$\end{minipage}\end{tcolorbox}\hfill\begin{tcolorbox}[enhanced,nobeforeafter,width=7.1cm,height=2.1cm,valign=center,colback=cder!6,colframe=cder,arc=2.5mm,boxrule=0.9pt,left=3mm]{\Large\bfseries\color{cder}b)}\quad\begin{minipage}[c]{5.3cm}\raggedright\Large $2$\end{minipage}\end{tcolorbox}\par\vspace{2.5mm}
\begin{tcolorbox}[enhanced,nobeforeafter,width=7.1cm,height=2.1cm,valign=center,colback=cder!6,colframe=cder,arc=2.5mm,boxrule=0.9pt,left=3mm]{\Large\bfseries\color{cder}c)}\quad\begin{minipage}[c]{5.3cm}\raggedright\Large $\dfrac{1}{4}$\end{minipage}\end{tcolorbox}\hfill\begin{tcolorbox}[enhanced,nobeforeafter,width=7.1cm,height=2.1cm,valign=center,colback=cder!6,colframe=cder,arc=2.5mm,boxrule=0.9pt,left=3mm]{\Large\bfseries\color{cder}d)}\quad\begin{minipage}[c]{5.3cm}\raggedright\Large $\dfrac{1}{8}$\end{minipage}\end{tcolorbox}
\end{tcolorbox}
\newpage
\begin{tcolorbox}[enhanced,width=16cm,height=14cm,colback=white,colframe=cder,boxrule=2pt,arc=0mm,
left=6mm,right=6mm,top=0mm,bottom=4mm,sharp corners,
overlay={\fill[cder] (frame.north west) rectangle ([yshift=-1.25cm]frame.north east);
\node[anchor=west,text=white,font=\sffamily\bfseries\LARGE] at ([xshift=5mm,yshift=-0.62cm]frame.north west) {Question 42 B};
\node[anchor=east,text=white,font=\sffamily\itshape\large] at ([xshift=-5mm,yshift=-0.62cm]frame.north east) {Dérivation};}]
\vspace*{1.75cm}
\begin{minipage}[t][6.9cm][c]{\linewidth}\LARGE\centering
\begin{minipage}{0.95\linewidth}\raggedright\linespread{1.25}\selectfont Soit $f(x)=\sqrt{x}$. Que vaut $f'(1)$~?\end{minipage}
\end{minipage}

\vfill
\begin{tcolorbox}[enhanced,nobeforeafter,width=7.1cm,height=2.1cm,valign=center,colback=cder!6,colframe=cder,arc=2.5mm,boxrule=0.9pt,left=3mm]{\Large\bfseries\color{cder}a)}\quad\begin{minipage}[c]{5.3cm}\raggedright\Large $2$\end{minipage}\end{tcolorbox}\hfill\begin{tcolorbox}[enhanced,nobeforeafter,width=7.1cm,height=2.1cm,valign=center,colback=cder!6,colframe=cder,arc=2.5mm,boxrule=0.9pt,left=3mm]{\Large\bfseries\color{cder}b)}\quad\begin{minipage}[c]{5.3cm}\raggedright\Large $\dfrac{1}{2}$\end{minipage}\end{tcolorbox}\par\vspace{2.5mm}
\begin{tcolorbox}[enhanced,nobeforeafter,width=7.1cm,height=2.1cm,valign=center,colback=cder!6,colframe=cder,arc=2.5mm,boxrule=0.9pt,left=3mm]{\Large\bfseries\color{cder}c)}\quad\begin{minipage}[c]{5.3cm}\raggedright\Large $\dfrac{1}{4}$\end{minipage}\end{tcolorbox}\hfill\begin{tcolorbox}[enhanced,nobeforeafter,width=7.1cm,height=2.1cm,valign=center,colback=cder!6,colframe=cder,arc=2.5mm,boxrule=0.9pt,left=3mm]{\Large\bfseries\color{cder}d)}\quad\begin{minipage}[c]{5.3cm}\raggedright\Large $1$\end{minipage}\end{tcolorbox}
\end{tcolorbox}
\newpage
\begin{tcolorbox}[enhanced,width=16cm,height=14cm,colback=white,colframe=cder,boxrule=2pt,arc=0mm,
left=6mm,right=6mm,top=0mm,bottom=4mm,sharp corners,
overlay={\fill[cder] (frame.north west) rectangle ([yshift=-1.25cm]frame.north east);
\node[anchor=west,text=white,font=\sffamily\bfseries\LARGE] at ([xshift=5mm,yshift=-0.62cm]frame.north west) {Question 43 B};
\node[anchor=east,text=white,font=\sffamily\itshape\large] at ([xshift=-5mm,yshift=-0.62cm]frame.north east) {Dérivation};}]
\vspace*{1.75cm}
\begin{minipage}[t][6.9cm][c]{\linewidth}\LARGE\centering
\begin{minipage}{0.95\linewidth}\raggedright\linespread{1.25}\selectfont Soit $f(x)=(-3x+1)e^x$. Que vaut $f'(0)$~?\end{minipage}
\end{minipage}

\vfill
\begin{tcolorbox}[enhanced,nobeforeafter,width=7.1cm,height=2.1cm,valign=center,colback=cder!6,colframe=cder,arc=2.5mm,boxrule=0.9pt,left=3mm]{\Large\bfseries\color{cder}a)}\quad\begin{minipage}[c]{5.3cm}\raggedright\Large $-2$\end{minipage}\end{tcolorbox}\hfill\begin{tcolorbox}[enhanced,nobeforeafter,width=7.1cm,height=2.1cm,valign=center,colback=cder!6,colframe=cder,arc=2.5mm,boxrule=0.9pt,left=3mm]{\Large\bfseries\color{cder}b)}\quad\begin{minipage}[c]{5.3cm}\raggedright\Large $-3$\end{minipage}\end{tcolorbox}\par\vspace{2.5mm}
\begin{tcolorbox}[enhanced,nobeforeafter,width=7.1cm,height=2.1cm,valign=center,colback=cder!6,colframe=cder,arc=2.5mm,boxrule=0.9pt,left=3mm]{\Large\bfseries\color{cder}c)}\quad\begin{minipage}[c]{5.3cm}\raggedright\Large $1$\end{minipage}\end{tcolorbox}\hfill\begin{tcolorbox}[enhanced,nobeforeafter,width=7.1cm,height=2.1cm,valign=center,colback=cder!6,colframe=cder,arc=2.5mm,boxrule=0.9pt,left=3mm]{\Large\bfseries\color{cder}d)}\quad\begin{minipage}[c]{5.3cm}\raggedright\Large $-1$\end{minipage}\end{tcolorbox}
\end{tcolorbox}
\newpage
\begin{tcolorbox}[enhanced,width=16cm,height=14cm,colback=white,colframe=cder,boxrule=2pt,arc=0mm,
left=6mm,right=6mm,top=0mm,bottom=4mm,sharp corners,
overlay={\fill[cder] (frame.north west) rectangle ([yshift=-1.25cm]frame.north east);
\node[anchor=west,text=white,font=\sffamily\bfseries\LARGE] at ([xshift=5mm,yshift=-0.62cm]frame.north west) {Question 44 B};
\node[anchor=east,text=white,font=\sffamily\itshape\large] at ([xshift=-5mm,yshift=-0.62cm]frame.north east) {Dérivation};}]
\vspace*{1.75cm}
\begin{minipage}[t][6.9cm][c]{\linewidth}\LARGE\centering
\begin{minipage}{0.95\linewidth}\raggedright\linespread{1.25}\selectfont Soit $f(x)=(4x)e^x$. Que vaut $f'(0)$~?\end{minipage}
\end{minipage}

\vfill
\begin{tcolorbox}[enhanced,nobeforeafter,width=7.1cm,height=2.1cm,valign=center,colback=cder!6,colframe=cder,arc=2.5mm,boxrule=0.9pt,left=3mm]{\Large\bfseries\color{cder}a)}\quad\begin{minipage}[c]{5.3cm}\raggedright\Large $0$\end{minipage}\end{tcolorbox}\hfill\begin{tcolorbox}[enhanced,nobeforeafter,width=7.1cm,height=2.1cm,valign=center,colback=cder!6,colframe=cder,arc=2.5mm,boxrule=0.9pt,left=3mm]{\Large\bfseries\color{cder}b)}\quad\begin{minipage}[c]{5.3cm}\raggedright\Large $5$\end{minipage}\end{tcolorbox}\par\vspace{2.5mm}
\begin{tcolorbox}[enhanced,nobeforeafter,width=7.1cm,height=2.1cm,valign=center,colback=cder!6,colframe=cder,arc=2.5mm,boxrule=0.9pt,left=3mm]{\Large\bfseries\color{cder}c)}\quad\begin{minipage}[c]{5.3cm}\raggedright\Large $4$\end{minipage}\end{tcolorbox}\hfill\begin{tcolorbox}[enhanced,nobeforeafter,width=7.1cm,height=2.1cm,valign=center,colback=cder!6,colframe=cder,arc=2.5mm,boxrule=0.9pt,left=3mm]{\Large\bfseries\color{cder}d)}\quad\begin{minipage}[c]{5.3cm}\raggedright\Large $3$\end{minipage}\end{tcolorbox}
\end{tcolorbox}
\newpage
\begin{tcolorbox}[enhanced,width=16cm,height=14cm,colback=white,colframe=cder,boxrule=2pt,arc=0mm,
left=6mm,right=6mm,top=0mm,bottom=4mm,sharp corners,
overlay={\fill[cder] (frame.north west) rectangle ([yshift=-1.25cm]frame.north east);
\node[anchor=west,text=white,font=\sffamily\bfseries\LARGE] at ([xshift=5mm,yshift=-0.62cm]frame.north west) {Question 45 B};
\node[anchor=east,text=white,font=\sffamily\itshape\large] at ([xshift=-5mm,yshift=-0.62cm]frame.north east) {Dérivation};}]
\vspace*{1.75cm}
\begin{minipage}[t][6.9cm][c]{\linewidth}\LARGE\centering
\begin{minipage}{0.95\linewidth}\raggedright\linespread{1.25}\selectfont La fonction $f$ est dérivable sur $\mathbb{R}$ et $f'(x)=4x+8$. Sur quel intervalle $f$ est-elle décroissante~?\end{minipage}
\end{minipage}

\vfill
\begin{tcolorbox}[enhanced,nobeforeafter,width=7.1cm,height=2.1cm,valign=center,colback=cder!6,colframe=cder,arc=2.5mm,boxrule=0.9pt,left=3mm]{\Large\bfseries\color{cder}a)}\quad\begin{minipage}[c]{5.3cm}\raggedright\Large $[-2\,;\,+\infty[$\end{minipage}\end{tcolorbox}\hfill\begin{tcolorbox}[enhanced,nobeforeafter,width=7.1cm,height=2.1cm,valign=center,colback=cder!6,colframe=cder,arc=2.5mm,boxrule=0.9pt,left=3mm]{\Large\bfseries\color{cder}b)}\quad\begin{minipage}[c]{5.3cm}\raggedright\Large $[2\,;\,+\infty[$\end{minipage}\end{tcolorbox}\par\vspace{2.5mm}
\begin{tcolorbox}[enhanced,nobeforeafter,width=7.1cm,height=2.1cm,valign=center,colback=cder!6,colframe=cder,arc=2.5mm,boxrule=0.9pt,left=3mm]{\Large\bfseries\color{cder}c)}\quad\begin{minipage}[c]{5.3cm}\raggedright\Large $]-\infty\,;\,2]$\end{minipage}\end{tcolorbox}\hfill\begin{tcolorbox}[enhanced,nobeforeafter,width=7.1cm,height=2.1cm,valign=center,colback=cder!6,colframe=cder,arc=2.5mm,boxrule=0.9pt,left=3mm]{\Large\bfseries\color{cder}d)}\quad\begin{minipage}[c]{5.3cm}\raggedright\Large $]-\infty\,;\,-2]$\end{minipage}\end{tcolorbox}
\end{tcolorbox}
\newpage
\begin{tcolorbox}[enhanced,width=16cm,height=14cm,colback=white,colframe=cder,boxrule=2pt,arc=0mm,
left=6mm,right=6mm,top=0mm,bottom=4mm,sharp corners,
overlay={\fill[cder] (frame.north west) rectangle ([yshift=-1.25cm]frame.north east);
\node[anchor=west,text=white,font=\sffamily\bfseries\LARGE] at ([xshift=5mm,yshift=-0.62cm]frame.north west) {Question 46 B};
\node[anchor=east,text=white,font=\sffamily\itshape\large] at ([xshift=-5mm,yshift=-0.62cm]frame.north east) {Dérivation};}]
\vspace*{1.75cm}
\begin{minipage}[t][6.9cm][c]{\linewidth}\LARGE\centering
\begin{minipage}{0.95\linewidth}\raggedright\linespread{1.25}\selectfont La fonction $f$ est dérivable sur $\mathbb{R}$ et $f'(x)=-x-5$. Sur quel intervalle $f$ est-elle décroissante~?\end{minipage}
\end{minipage}

\vfill
\begin{tcolorbox}[enhanced,nobeforeafter,width=7.1cm,height=2.1cm,valign=center,colback=cder!6,colframe=cder,arc=2.5mm,boxrule=0.9pt,left=3mm]{\Large\bfseries\color{cder}a)}\quad\begin{minipage}[c]{5.3cm}\raggedright\Large $[-5\,;\,+\infty[$\end{minipage}\end{tcolorbox}\hfill\begin{tcolorbox}[enhanced,nobeforeafter,width=7.1cm,height=2.1cm,valign=center,colback=cder!6,colframe=cder,arc=2.5mm,boxrule=0.9pt,left=3mm]{\Large\bfseries\color{cder}b)}\quad\begin{minipage}[c]{5.3cm}\raggedright\Large $[5\,;\,+\infty[$\end{minipage}\end{tcolorbox}\par\vspace{2.5mm}
\begin{tcolorbox}[enhanced,nobeforeafter,width=7.1cm,height=2.1cm,valign=center,colback=cder!6,colframe=cder,arc=2.5mm,boxrule=0.9pt,left=3mm]{\Large\bfseries\color{cder}c)}\quad\begin{minipage}[c]{5.3cm}\raggedright\Large $]-\infty\,;\,5]$\end{minipage}\end{tcolorbox}\hfill\begin{tcolorbox}[enhanced,nobeforeafter,width=7.1cm,height=2.1cm,valign=center,colback=cder!6,colframe=cder,arc=2.5mm,boxrule=0.9pt,left=3mm]{\Large\bfseries\color{cder}d)}\quad\begin{minipage}[c]{5.3cm}\raggedright\Large $]-\infty\,;\,-5]$\end{minipage}\end{tcolorbox}
\end{tcolorbox}
\newpage
\begin{tcolorbox}[enhanced,width=16cm,height=14cm,colback=white,colframe=cder,boxrule=2pt,arc=0mm,
left=6mm,right=6mm,top=0mm,bottom=4mm,sharp corners,
overlay={\fill[cder] (frame.north west) rectangle ([yshift=-1.25cm]frame.north east);
\node[anchor=west,text=white,font=\sffamily\bfseries\LARGE] at ([xshift=5mm,yshift=-0.62cm]frame.north west) {Question 47 B};
\node[anchor=east,text=white,font=\sffamily\itshape\large] at ([xshift=-5mm,yshift=-0.62cm]frame.north east) {Dérivation};}]
\vspace*{1.75cm}
\begin{minipage}[t][6.9cm][c]{\linewidth}\LARGE\centering
\begin{minipage}{0.95\linewidth}\raggedright\linespread{1.25}\selectfont Soit $f(x)=2x^2+8x+3$. Quel est le minimum de $f$ sur $\mathbb{R}$~?\end{minipage}
\end{minipage}

\vfill
\begin{tcolorbox}[enhanced,nobeforeafter,width=7.1cm,height=2.1cm,valign=center,colback=cder!6,colframe=cder,arc=2.5mm,boxrule=0.9pt,left=3mm]{\Large\bfseries\color{cder}a)}\quad\begin{minipage}[c]{5.3cm}\raggedright\Large $-5$\end{minipage}\end{tcolorbox}\hfill\begin{tcolorbox}[enhanced,nobeforeafter,width=7.1cm,height=2.1cm,valign=center,colback=cder!6,colframe=cder,arc=2.5mm,boxrule=0.9pt,left=3mm]{\Large\bfseries\color{cder}b)}\quad\begin{minipage}[c]{5.3cm}\raggedright\Large $5$\end{minipage}\end{tcolorbox}\par\vspace{2.5mm}
\begin{tcolorbox}[enhanced,nobeforeafter,width=7.1cm,height=2.1cm,valign=center,colback=cder!6,colframe=cder,arc=2.5mm,boxrule=0.9pt,left=3mm]{\Large\bfseries\color{cder}c)}\quad\begin{minipage}[c]{5.3cm}\raggedright\Large $-2$\end{minipage}\end{tcolorbox}\hfill\begin{tcolorbox}[enhanced,nobeforeafter,width=7.1cm,height=2.1cm,valign=center,colback=cder!6,colframe=cder,arc=2.5mm,boxrule=0.9pt,left=3mm]{\Large\bfseries\color{cder}d)}\quad\begin{minipage}[c]{5.3cm}\raggedright\Large $3$\end{minipage}\end{tcolorbox}
\end{tcolorbox}
\newpage
\begin{tcolorbox}[enhanced,width=16cm,height=14cm,colback=white,colframe=cder,boxrule=2pt,arc=0mm,
left=6mm,right=6mm,top=0mm,bottom=4mm,sharp corners,
overlay={\fill[cder] (frame.north west) rectangle ([yshift=-1.25cm]frame.north east);
\node[anchor=west,text=white,font=\sffamily\bfseries\LARGE] at ([xshift=5mm,yshift=-0.62cm]frame.north west) {Question 48 B};
\node[anchor=east,text=white,font=\sffamily\itshape\large] at ([xshift=-5mm,yshift=-0.62cm]frame.north east) {Dérivation};}]
\vspace*{1.75cm}
\begin{minipage}[t][6.9cm][c]{\linewidth}\LARGE\centering
\begin{minipage}{0.95\linewidth}\raggedright\linespread{1.25}\selectfont Soit $f(x)=-2x^2+4x+1$. Quel est le maximum de $f$ sur $\mathbb{R}$~?\end{minipage}
\end{minipage}

\vfill
\begin{tcolorbox}[enhanced,nobeforeafter,width=7.1cm,height=2.1cm,valign=center,colback=cder!6,colframe=cder,arc=2.5mm,boxrule=0.9pt,left=3mm]{\Large\bfseries\color{cder}a)}\quad\begin{minipage}[c]{5.3cm}\raggedright\Large $-3$\end{minipage}\end{tcolorbox}\hfill\begin{tcolorbox}[enhanced,nobeforeafter,width=7.1cm,height=2.1cm,valign=center,colback=cder!6,colframe=cder,arc=2.5mm,boxrule=0.9pt,left=3mm]{\Large\bfseries\color{cder}b)}\quad\begin{minipage}[c]{5.3cm}\raggedright\Large $5$\end{minipage}\end{tcolorbox}\par\vspace{2.5mm}
\begin{tcolorbox}[enhanced,nobeforeafter,width=7.1cm,height=2.1cm,valign=center,colback=cder!6,colframe=cder,arc=2.5mm,boxrule=0.9pt,left=3mm]{\Large\bfseries\color{cder}c)}\quad\begin{minipage}[c]{5.3cm}\raggedright\Large $3$\end{minipage}\end{tcolorbox}\hfill\begin{tcolorbox}[enhanced,nobeforeafter,width=7.1cm,height=2.1cm,valign=center,colback=cder!6,colframe=cder,arc=2.5mm,boxrule=0.9pt,left=3mm]{\Large\bfseries\color{cder}d)}\quad\begin{minipage}[c]{5.3cm}\raggedright\Large $1$\end{minipage}\end{tcolorbox}
\end{tcolorbox}
\newpage
\begin{tcolorbox}[enhanced,width=16cm,height=14cm,colback=white,colframe=ctri,boxrule=2pt,arc=0mm,
left=6mm,right=6mm,top=0mm,bottom=4mm,sharp corners,
overlay={\fill[ctri] (frame.north west) rectangle ([yshift=-1.25cm]frame.north east);
\node[anchor=west,text=white,font=\sffamily\bfseries\LARGE] at ([xshift=5mm,yshift=-0.62cm]frame.north west) {Question 49 B};
\node[anchor=east,text=white,font=\sffamily\itshape\large] at ([xshift=-5mm,yshift=-0.62cm]frame.north east) {Trigonométrie et exponentielle};}]
\vspace*{1.75cm}
\begin{minipage}[t][6.9cm][c]{\linewidth}\LARGE\centering
\begin{minipage}{0.95\linewidth}\raggedright\linespread{1.25}\selectfont Quelle est la mesure en radians d'un angle de $30^\circ$~?\end{minipage}
\end{minipage}

\vfill
\begin{tcolorbox}[enhanced,nobeforeafter,width=7.1cm,height=2.1cm,valign=center,colback=ctri!6,colframe=ctri,arc=2.5mm,boxrule=0.9pt,left=3mm]{\Large\bfseries\color{ctri}a)}\quad\begin{minipage}[c]{5.3cm}\raggedright\Large $6\pi$\end{minipage}\end{tcolorbox}\hfill\begin{tcolorbox}[enhanced,nobeforeafter,width=7.1cm,height=2.1cm,valign=center,colback=ctri!6,colframe=ctri,arc=2.5mm,boxrule=0.9pt,left=3mm]{\Large\bfseries\color{ctri}b)}\quad\begin{minipage}[c]{5.3cm}\raggedright\Large $\dfrac{\pi}{12}$\end{minipage}\end{tcolorbox}\par\vspace{2.5mm}
\begin{tcolorbox}[enhanced,nobeforeafter,width=7.1cm,height=2.1cm,valign=center,colback=ctri!6,colframe=ctri,arc=2.5mm,boxrule=0.9pt,left=3mm]{\Large\bfseries\color{ctri}c)}\quad\begin{minipage}[c]{5.3cm}\raggedright\Large $\dfrac{\pi}{6}$\end{minipage}\end{tcolorbox}\hfill\begin{tcolorbox}[enhanced,nobeforeafter,width=7.1cm,height=2.1cm,valign=center,colback=ctri!6,colframe=ctri,arc=2.5mm,boxrule=0.9pt,left=3mm]{\Large\bfseries\color{ctri}d)}\quad\begin{minipage}[c]{5.3cm}\raggedright\Large $\dfrac{\pi}{3}$\end{minipage}\end{tcolorbox}
\end{tcolorbox}
\newpage
\begin{tcolorbox}[enhanced,width=16cm,height=14cm,colback=white,colframe=ctri,boxrule=2pt,arc=0mm,
left=6mm,right=6mm,top=0mm,bottom=4mm,sharp corners,
overlay={\fill[ctri] (frame.north west) rectangle ([yshift=-1.25cm]frame.north east);
\node[anchor=west,text=white,font=\sffamily\bfseries\LARGE] at ([xshift=5mm,yshift=-0.62cm]frame.north west) {Question 50 B};
\node[anchor=east,text=white,font=\sffamily\itshape\large] at ([xshift=-5mm,yshift=-0.62cm]frame.north east) {Trigonométrie et exponentielle};}]
\vspace*{1.75cm}
\begin{minipage}[t][6.9cm][c]{\linewidth}\LARGE\centering
\begin{minipage}{0.95\linewidth}\raggedright\linespread{1.25}\selectfont Quelle est la mesure en radians d'un angle de $120^\circ$~?\end{minipage}
\end{minipage}

\vfill
\begin{tcolorbox}[enhanced,nobeforeafter,width=7.1cm,height=2.1cm,valign=center,colback=ctri!6,colframe=ctri,arc=2.5mm,boxrule=0.9pt,left=3mm]{\Large\bfseries\color{ctri}a)}\quad\begin{minipage}[c]{5.3cm}\raggedright\Large $\dfrac{4\pi}{3}$\end{minipage}\end{tcolorbox}\hfill\begin{tcolorbox}[enhanced,nobeforeafter,width=7.1cm,height=2.1cm,valign=center,colback=ctri!6,colframe=ctri,arc=2.5mm,boxrule=0.9pt,left=3mm]{\Large\bfseries\color{ctri}b)}\quad\begin{minipage}[c]{5.3cm}\raggedright\Large $\dfrac{2\pi}{3}$\end{minipage}\end{tcolorbox}\par\vspace{2.5mm}
\begin{tcolorbox}[enhanced,nobeforeafter,width=7.1cm,height=2.1cm,valign=center,colback=ctri!6,colframe=ctri,arc=2.5mm,boxrule=0.9pt,left=3mm]{\Large\bfseries\color{ctri}c)}\quad\begin{minipage}[c]{5.3cm}\raggedright\Large $\dfrac{\pi}{3}$\end{minipage}\end{tcolorbox}\hfill\begin{tcolorbox}[enhanced,nobeforeafter,width=7.1cm,height=2.1cm,valign=center,colback=ctri!6,colframe=ctri,arc=2.5mm,boxrule=0.9pt,left=3mm]{\Large\bfseries\color{ctri}d)}\quad\begin{minipage}[c]{5.3cm}\raggedright\Large $\dfrac{3\pi}{2}$\end{minipage}\end{tcolorbox}
\end{tcolorbox}
\newpage
\begin{tcolorbox}[enhanced,width=16cm,height=14cm,colback=white,colframe=ctri,boxrule=2pt,arc=0mm,
left=6mm,right=6mm,top=0mm,bottom=4mm,sharp corners,
overlay={\fill[ctri] (frame.north west) rectangle ([yshift=-1.25cm]frame.north east);
\node[anchor=west,text=white,font=\sffamily\bfseries\LARGE] at ([xshift=5mm,yshift=-0.62cm]frame.north west) {Question 51 B};
\node[anchor=east,text=white,font=\sffamily\itshape\large] at ([xshift=-5mm,yshift=-0.62cm]frame.north east) {Trigonométrie et exponentielle};}]
\vspace*{1.75cm}
\begin{minipage}[t][6.9cm][c]{\linewidth}\LARGE\centering
\begin{minipage}{0.95\linewidth}\raggedright\linespread{1.25}\selectfont Quelle est la mesure en degrés d'un angle de $\dfrac{3\pi}{2}$ radians~?\end{minipage}
\end{minipage}

\vfill
\begin{tcolorbox}[enhanced,nobeforeafter,width=7.1cm,height=2.1cm,valign=center,colback=ctri!6,colframe=ctri,arc=2.5mm,boxrule=0.9pt,left=3mm]{\Large\bfseries\color{ctri}a)}\quad\begin{minipage}[c]{5.3cm}\raggedright\Large $270^\circ$\end{minipage}\end{tcolorbox}\hfill\begin{tcolorbox}[enhanced,nobeforeafter,width=7.1cm,height=2.1cm,valign=center,colback=ctri!6,colframe=ctri,arc=2.5mm,boxrule=0.9pt,left=3mm]{\Large\bfseries\color{ctri}b)}\quad\begin{minipage}[c]{5.3cm}\raggedright\Large $540^\circ$\end{minipage}\end{tcolorbox}\par\vspace{2.5mm}
\begin{tcolorbox}[enhanced,nobeforeafter,width=7.1cm,height=2.1cm,valign=center,colback=ctri!6,colframe=ctri,arc=2.5mm,boxrule=0.9pt,left=3mm]{\Large\bfseries\color{ctri}c)}\quad\begin{minipage}[c]{5.3cm}\raggedright\Large $90^\circ$\end{minipage}\end{tcolorbox}\hfill\begin{tcolorbox}[enhanced,nobeforeafter,width=7.1cm,height=2.1cm,valign=center,colback=ctri!6,colframe=ctri,arc=2.5mm,boxrule=0.9pt,left=3mm]{\Large\bfseries\color{ctri}d)}\quad\begin{minipage}[c]{5.3cm}\raggedright\Large $135^\circ$\end{minipage}\end{tcolorbox}
\end{tcolorbox}
\newpage
\begin{tcolorbox}[enhanced,width=16cm,height=14cm,colback=white,colframe=ctri,boxrule=2pt,arc=0mm,
left=6mm,right=6mm,top=0mm,bottom=4mm,sharp corners,
overlay={\fill[ctri] (frame.north west) rectangle ([yshift=-1.25cm]frame.north east);
\node[anchor=west,text=white,font=\sffamily\bfseries\LARGE] at ([xshift=5mm,yshift=-0.62cm]frame.north west) {Question 52 B};
\node[anchor=east,text=white,font=\sffamily\itshape\large] at ([xshift=-5mm,yshift=-0.62cm]frame.north east) {Trigonométrie et exponentielle};}]
\vspace*{1.75cm}
\begin{minipage}[t][6.9cm][c]{\linewidth}\LARGE\centering
\begin{minipage}{0.95\linewidth}\raggedright\linespread{1.25}\selectfont Quelle est la mesure en degrés d'un angle de $\dfrac{\pi}{12}$ radians~?\end{minipage}
\end{minipage}

\vfill
\begin{tcolorbox}[enhanced,nobeforeafter,width=7.1cm,height=2.1cm,valign=center,colback=ctri!6,colframe=ctri,arc=2.5mm,boxrule=0.9pt,left=3mm]{\Large\bfseries\color{ctri}a)}\quad\begin{minipage}[c]{5.3cm}\raggedright\Large $30^\circ$\end{minipage}\end{tcolorbox}\hfill\begin{tcolorbox}[enhanced,nobeforeafter,width=7.1cm,height=2.1cm,valign=center,colback=ctri!6,colframe=ctri,arc=2.5mm,boxrule=0.9pt,left=3mm]{\Large\bfseries\color{ctri}b)}\quad\begin{minipage}[c]{5.3cm}\raggedright\Large $7{,}5^\circ$\end{minipage}\end{tcolorbox}\par\vspace{2.5mm}
\begin{tcolorbox}[enhanced,nobeforeafter,width=7.1cm,height=2.1cm,valign=center,colback=ctri!6,colframe=ctri,arc=2.5mm,boxrule=0.9pt,left=3mm]{\Large\bfseries\color{ctri}c)}\quad\begin{minipage}[c]{5.3cm}\raggedright\Large $15^\circ$\end{minipage}\end{tcolorbox}\hfill\begin{tcolorbox}[enhanced,nobeforeafter,width=7.1cm,height=2.1cm,valign=center,colback=ctri!6,colframe=ctri,arc=2.5mm,boxrule=0.9pt,left=3mm]{\Large\bfseries\color{ctri}d)}\quad\begin{minipage}[c]{5.3cm}\raggedright\Large $345^\circ$\end{minipage}\end{tcolorbox}
\end{tcolorbox}
\newpage
\begin{tcolorbox}[enhanced,width=16cm,height=14cm,colback=white,colframe=ctri,boxrule=2pt,arc=0mm,
left=6mm,right=6mm,top=0mm,bottom=4mm,sharp corners,
overlay={\fill[ctri] (frame.north west) rectangle ([yshift=-1.25cm]frame.north east);
\node[anchor=west,text=white,font=\sffamily\bfseries\LARGE] at ([xshift=5mm,yshift=-0.62cm]frame.north west) {Question 53 B};
\node[anchor=east,text=white,font=\sffamily\itshape\large] at ([xshift=-5mm,yshift=-0.62cm]frame.north east) {Trigonométrie et exponentielle};}]
\vspace*{1.75cm}
\begin{minipage}[t][6.9cm][c]{\linewidth}\LARGE\centering
\begin{minipage}{0.95\linewidth}\raggedright\linespread{1.25}\selectfont Que vaut $\cos\left(\dfrac{\pi}{6}\right)$~?\end{minipage}
\end{minipage}

\vfill
\begin{tcolorbox}[enhanced,nobeforeafter,width=7.1cm,height=2.1cm,valign=center,colback=ctri!6,colframe=ctri,arc=2.5mm,boxrule=0.9pt,left=3mm]{\Large\bfseries\color{ctri}a)}\quad\begin{minipage}[c]{5.3cm}\raggedright\Large $1$\end{minipage}\end{tcolorbox}\hfill\begin{tcolorbox}[enhanced,nobeforeafter,width=7.1cm,height=2.1cm,valign=center,colback=ctri!6,colframe=ctri,arc=2.5mm,boxrule=0.9pt,left=3mm]{\Large\bfseries\color{ctri}b)}\quad\begin{minipage}[c]{5.3cm}\raggedright\Large $-\dfrac{\sqrt3}{2}$\end{minipage}\end{tcolorbox}\par\vspace{2.5mm}
\begin{tcolorbox}[enhanced,nobeforeafter,width=7.1cm,height=2.1cm,valign=center,colback=ctri!6,colframe=ctri,arc=2.5mm,boxrule=0.9pt,left=3mm]{\Large\bfseries\color{ctri}c)}\quad\begin{minipage}[c]{5.3cm}\raggedright\Large $\dfrac{\sqrt2}{2}$\end{minipage}\end{tcolorbox}\hfill\begin{tcolorbox}[enhanced,nobeforeafter,width=7.1cm,height=2.1cm,valign=center,colback=ctri!6,colframe=ctri,arc=2.5mm,boxrule=0.9pt,left=3mm]{\Large\bfseries\color{ctri}d)}\quad\begin{minipage}[c]{5.3cm}\raggedright\Large $\dfrac{\sqrt3}{2}$\end{minipage}\end{tcolorbox}
\end{tcolorbox}
\newpage
\begin{tcolorbox}[enhanced,width=16cm,height=14cm,colback=white,colframe=ctri,boxrule=2pt,arc=0mm,
left=6mm,right=6mm,top=0mm,bottom=4mm,sharp corners,
overlay={\fill[ctri] (frame.north west) rectangle ([yshift=-1.25cm]frame.north east);
\node[anchor=west,text=white,font=\sffamily\bfseries\LARGE] at ([xshift=5mm,yshift=-0.62cm]frame.north west) {Question 54 B};
\node[anchor=east,text=white,font=\sffamily\itshape\large] at ([xshift=-5mm,yshift=-0.62cm]frame.north east) {Trigonométrie et exponentielle};}]
\vspace*{1.75cm}
\begin{minipage}[t][6.9cm][c]{\linewidth}\LARGE\centering
\begin{minipage}{0.95\linewidth}\raggedright\linespread{1.25}\selectfont Que vaut $\cos\left(\dfrac{2\pi}{3}\right)$~?\end{minipage}
\end{minipage}

\vfill
\begin{tcolorbox}[enhanced,nobeforeafter,width=7.1cm,height=2.1cm,valign=center,colback=ctri!6,colframe=ctri,arc=2.5mm,boxrule=0.9pt,left=3mm]{\Large\bfseries\color{ctri}a)}\quad\begin{minipage}[c]{5.3cm}\raggedright\Large $\dfrac12$\end{minipage}\end{tcolorbox}\hfill\begin{tcolorbox}[enhanced,nobeforeafter,width=7.1cm,height=2.1cm,valign=center,colback=ctri!6,colframe=ctri,arc=2.5mm,boxrule=0.9pt,left=3mm]{\Large\bfseries\color{ctri}b)}\quad\begin{minipage}[c]{5.3cm}\raggedright\Large $0$\end{minipage}\end{tcolorbox}\par\vspace{2.5mm}
\begin{tcolorbox}[enhanced,nobeforeafter,width=7.1cm,height=2.1cm,valign=center,colback=ctri!6,colframe=ctri,arc=2.5mm,boxrule=0.9pt,left=3mm]{\Large\bfseries\color{ctri}c)}\quad\begin{minipage}[c]{5.3cm}\raggedright\Large $-\dfrac12$\end{minipage}\end{tcolorbox}\hfill\begin{tcolorbox}[enhanced,nobeforeafter,width=7.1cm,height=2.1cm,valign=center,colback=ctri!6,colframe=ctri,arc=2.5mm,boxrule=0.9pt,left=3mm]{\Large\bfseries\color{ctri}d)}\quad\begin{minipage}[c]{5.3cm}\raggedright\Large $-\dfrac{\sqrt2}{2}$\end{minipage}\end{tcolorbox}
\end{tcolorbox}
\newpage
\begin{tcolorbox}[enhanced,width=16cm,height=14cm,colback=white,colframe=ctri,boxrule=2pt,arc=0mm,
left=6mm,right=6mm,top=0mm,bottom=4mm,sharp corners,
overlay={\fill[ctri] (frame.north west) rectangle ([yshift=-1.25cm]frame.north east);
\node[anchor=west,text=white,font=\sffamily\bfseries\LARGE] at ([xshift=5mm,yshift=-0.62cm]frame.north west) {Question 55 B};
\node[anchor=east,text=white,font=\sffamily\itshape\large] at ([xshift=-5mm,yshift=-0.62cm]frame.north east) {Trigonométrie et exponentielle};}]
\vspace*{1.75cm}
\begin{minipage}[t][6.9cm][c]{\linewidth}\LARGE\centering
\begin{minipage}{0.95\linewidth}\raggedright\linespread{1.25}\selectfont Que vaut $\sin\left(\dfrac{\pi}{3}\right)$~?\end{minipage}
\end{minipage}

\vfill
\begin{tcolorbox}[enhanced,nobeforeafter,width=7.1cm,height=2.1cm,valign=center,colback=ctri!6,colframe=ctri,arc=2.5mm,boxrule=0.9pt,left=3mm]{\Large\bfseries\color{ctri}a)}\quad\begin{minipage}[c]{5.3cm}\raggedright\Large $\dfrac{\sqrt2}{2}$\end{minipage}\end{tcolorbox}\hfill\begin{tcolorbox}[enhanced,nobeforeafter,width=7.1cm,height=2.1cm,valign=center,colback=ctri!6,colframe=ctri,arc=2.5mm,boxrule=0.9pt,left=3mm]{\Large\bfseries\color{ctri}b)}\quad\begin{minipage}[c]{5.3cm}\raggedright\Large $\dfrac{\sqrt3}{2}$\end{minipage}\end{tcolorbox}\par\vspace{2.5mm}
\begin{tcolorbox}[enhanced,nobeforeafter,width=7.1cm,height=2.1cm,valign=center,colback=ctri!6,colframe=ctri,arc=2.5mm,boxrule=0.9pt,left=3mm]{\Large\bfseries\color{ctri}c)}\quad\begin{minipage}[c]{5.3cm}\raggedright\Large $-\dfrac{\sqrt3}{2}$\end{minipage}\end{tcolorbox}\hfill\begin{tcolorbox}[enhanced,nobeforeafter,width=7.1cm,height=2.1cm,valign=center,colback=ctri!6,colframe=ctri,arc=2.5mm,boxrule=0.9pt,left=3mm]{\Large\bfseries\color{ctri}d)}\quad\begin{minipage}[c]{5.3cm}\raggedright\Large $1$\end{minipage}\end{tcolorbox}
\end{tcolorbox}
\newpage
\begin{tcolorbox}[enhanced,width=16cm,height=14cm,colback=white,colframe=ctri,boxrule=2pt,arc=0mm,
left=6mm,right=6mm,top=0mm,bottom=4mm,sharp corners,
overlay={\fill[ctri] (frame.north west) rectangle ([yshift=-1.25cm]frame.north east);
\node[anchor=west,text=white,font=\sffamily\bfseries\LARGE] at ([xshift=5mm,yshift=-0.62cm]frame.north west) {Question 56 B};
\node[anchor=east,text=white,font=\sffamily\itshape\large] at ([xshift=-5mm,yshift=-0.62cm]frame.north east) {Trigonométrie et exponentielle};}]
\vspace*{1.75cm}
\begin{minipage}[t][6.9cm][c]{\linewidth}\LARGE\centering
\begin{minipage}{0.95\linewidth}\raggedright\linespread{1.25}\selectfont Que vaut $\sin\left(\dfrac{\pi}{2}\right)$~?\end{minipage}
\end{minipage}

\vfill
\begin{tcolorbox}[enhanced,nobeforeafter,width=7.1cm,height=2.1cm,valign=center,colback=ctri!6,colframe=ctri,arc=2.5mm,boxrule=0.9pt,left=3mm]{\Large\bfseries\color{ctri}a)}\quad\begin{minipage}[c]{5.3cm}\raggedright\Large $\dfrac{\sqrt2}{2}$\end{minipage}\end{tcolorbox}\hfill\begin{tcolorbox}[enhanced,nobeforeafter,width=7.1cm,height=2.1cm,valign=center,colback=ctri!6,colframe=ctri,arc=2.5mm,boxrule=0.9pt,left=3mm]{\Large\bfseries\color{ctri}b)}\quad\begin{minipage}[c]{5.3cm}\raggedright\Large $-1$\end{minipage}\end{tcolorbox}\par\vspace{2.5mm}
\begin{tcolorbox}[enhanced,nobeforeafter,width=7.1cm,height=2.1cm,valign=center,colback=ctri!6,colframe=ctri,arc=2.5mm,boxrule=0.9pt,left=3mm]{\Large\bfseries\color{ctri}c)}\quad\begin{minipage}[c]{5.3cm}\raggedright\Large $1$\end{minipage}\end{tcolorbox}\hfill\begin{tcolorbox}[enhanced,nobeforeafter,width=7.1cm,height=2.1cm,valign=center,colback=ctri!6,colframe=ctri,arc=2.5mm,boxrule=0.9pt,left=3mm]{\Large\bfseries\color{ctri}d)}\quad\begin{minipage}[c]{5.3cm}\raggedright\Large $\dfrac{\sqrt3}{2}$\end{minipage}\end{tcolorbox}
\end{tcolorbox}
\newpage
\begin{tcolorbox}[enhanced,width=16cm,height=14cm,colback=white,colframe=ctri,boxrule=2pt,arc=0mm,
left=6mm,right=6mm,top=0mm,bottom=4mm,sharp corners,
overlay={\fill[ctri] (frame.north west) rectangle ([yshift=-1.25cm]frame.north east);
\node[anchor=west,text=white,font=\sffamily\bfseries\LARGE] at ([xshift=5mm,yshift=-0.62cm]frame.north west) {Question 57 B};
\node[anchor=east,text=white,font=\sffamily\itshape\large] at ([xshift=-5mm,yshift=-0.62cm]frame.north east) {Trigonométrie et exponentielle};}]
\vspace*{1.75cm}
\begin{minipage}[t][6.9cm][c]{\linewidth}\LARGE\centering
\begin{minipage}{0.95\linewidth}\raggedright\linespread{1.25}\selectfont On sait que $\sin x=0{,}3$. Que vaut $\sin(-x)$~?\end{minipage}
\end{minipage}

\vfill
\begin{tcolorbox}[enhanced,nobeforeafter,width=7.1cm,height=2.1cm,valign=center,colback=ctri!6,colframe=ctri,arc=2.5mm,boxrule=0.9pt,left=3mm]{\Large\bfseries\color{ctri}a)}\quad\begin{minipage}[c]{5.3cm}\raggedright\Large $0{,}91$\end{minipage}\end{tcolorbox}\hfill\begin{tcolorbox}[enhanced,nobeforeafter,width=7.1cm,height=2.1cm,valign=center,colback=ctri!6,colframe=ctri,arc=2.5mm,boxrule=0.9pt,left=3mm]{\Large\bfseries\color{ctri}b)}\quad\begin{minipage}[c]{5.3cm}\raggedright\Large $0{,}7$\end{minipage}\end{tcolorbox}\par\vspace{2.5mm}
\begin{tcolorbox}[enhanced,nobeforeafter,width=7.1cm,height=2.1cm,valign=center,colback=ctri!6,colframe=ctri,arc=2.5mm,boxrule=0.9pt,left=3mm]{\Large\bfseries\color{ctri}c)}\quad\begin{minipage}[c]{5.3cm}\raggedright\Large $0{,}3$\end{minipage}\end{tcolorbox}\hfill\begin{tcolorbox}[enhanced,nobeforeafter,width=7.1cm,height=2.1cm,valign=center,colback=ctri!6,colframe=ctri,arc=2.5mm,boxrule=0.9pt,left=3mm]{\Large\bfseries\color{ctri}d)}\quad\begin{minipage}[c]{5.3cm}\raggedright\Large $-0{,}3$\end{minipage}\end{tcolorbox}
\end{tcolorbox}
\newpage
\begin{tcolorbox}[enhanced,width=16cm,height=14cm,colback=white,colframe=ctri,boxrule=2pt,arc=0mm,
left=6mm,right=6mm,top=0mm,bottom=4mm,sharp corners,
overlay={\fill[ctri] (frame.north west) rectangle ([yshift=-1.25cm]frame.north east);
\node[anchor=west,text=white,font=\sffamily\bfseries\LARGE] at ([xshift=5mm,yshift=-0.62cm]frame.north west) {Question 58 B};
\node[anchor=east,text=white,font=\sffamily\itshape\large] at ([xshift=-5mm,yshift=-0.62cm]frame.north east) {Trigonométrie et exponentielle};}]
\vspace*{1.75cm}
\begin{minipage}[t][6.9cm][c]{\linewidth}\LARGE\centering
\begin{minipage}{0.95\linewidth}\raggedright\linespread{1.25}\selectfont On sait que $\cos x=-0{,}4$. Que vaut $\cos(\pi+x)$~?\end{minipage}
\end{minipage}

\vfill
\begin{tcolorbox}[enhanced,nobeforeafter,width=7.1cm,height=2.1cm,valign=center,colback=ctri!6,colframe=ctri,arc=2.5mm,boxrule=0.9pt,left=3mm]{\Large\bfseries\color{ctri}a)}\quad\begin{minipage}[c]{5.3cm}\raggedright\Large $0{,}84$\end{minipage}\end{tcolorbox}\hfill\begin{tcolorbox}[enhanced,nobeforeafter,width=7.1cm,height=2.1cm,valign=center,colback=ctri!6,colframe=ctri,arc=2.5mm,boxrule=0.9pt,left=3mm]{\Large\bfseries\color{ctri}b)}\quad\begin{minipage}[c]{5.3cm}\raggedright\Large $0{,}4$\end{minipage}\end{tcolorbox}\par\vspace{2.5mm}
\begin{tcolorbox}[enhanced,nobeforeafter,width=7.1cm,height=2.1cm,valign=center,colback=ctri!6,colframe=ctri,arc=2.5mm,boxrule=0.9pt,left=3mm]{\Large\bfseries\color{ctri}c)}\quad\begin{minipage}[c]{5.3cm}\raggedright\Large $-0{,}4$\end{minipage}\end{tcolorbox}\hfill\begin{tcolorbox}[enhanced,nobeforeafter,width=7.1cm,height=2.1cm,valign=center,colback=ctri!6,colframe=ctri,arc=2.5mm,boxrule=0.9pt,left=3mm]{\Large\bfseries\color{ctri}d)}\quad\begin{minipage}[c]{5.3cm}\raggedright\Large $1{,}4$\end{minipage}\end{tcolorbox}
\end{tcolorbox}
\newpage
\begin{tcolorbox}[enhanced,width=16cm,height=14cm,colback=white,colframe=ctri,boxrule=2pt,arc=0mm,
left=6mm,right=6mm,top=0mm,bottom=4mm,sharp corners,
overlay={\fill[ctri] (frame.north west) rectangle ([yshift=-1.25cm]frame.north east);
\node[anchor=west,text=white,font=\sffamily\bfseries\LARGE] at ([xshift=5mm,yshift=-0.62cm]frame.north west) {Question 59 B};
\node[anchor=east,text=white,font=\sffamily\itshape\large] at ([xshift=-5mm,yshift=-0.62cm]frame.north east) {Trigonométrie et exponentielle};}]
\vspace*{1.75cm}
\begin{minipage}[t][6.9cm][c]{\linewidth}\LARGE\centering
\begin{minipage}{0.95\linewidth}\raggedright\linespread{1.25}\selectfont Quelle est la mesure principale de l'angle $-\dfrac{5\pi}{4}$~?\end{minipage}
\end{minipage}

\vfill
\begin{tcolorbox}[enhanced,nobeforeafter,width=7.1cm,height=2.1cm,valign=center,colback=ctri!6,colframe=ctri,arc=2.5mm,boxrule=0.9pt,left=3mm]{\Large\bfseries\color{ctri}a)}\quad\begin{minipage}[c]{5.3cm}\raggedright\Large $\dfrac{3\pi}{8}$\end{minipage}\end{tcolorbox}\hfill\begin{tcolorbox}[enhanced,nobeforeafter,width=7.1cm,height=2.1cm,valign=center,colback=ctri!6,colframe=ctri,arc=2.5mm,boxrule=0.9pt,left=3mm]{\Large\bfseries\color{ctri}b)}\quad\begin{minipage}[c]{5.3cm}\raggedright\Large $\dfrac{3\pi}{4}$\end{minipage}\end{tcolorbox}\par\vspace{2.5mm}
\begin{tcolorbox}[enhanced,nobeforeafter,width=7.1cm,height=2.1cm,valign=center,colback=ctri!6,colframe=ctri,arc=2.5mm,boxrule=0.9pt,left=3mm]{\Large\bfseries\color{ctri}c)}\quad\begin{minipage}[c]{5.3cm}\raggedright\Large $-\dfrac{3\pi}{4}$\end{minipage}\end{tcolorbox}\hfill\begin{tcolorbox}[enhanced,nobeforeafter,width=7.1cm,height=2.1cm,valign=center,colback=ctri!6,colframe=ctri,arc=2.5mm,boxrule=0.9pt,left=3mm]{\Large\bfseries\color{ctri}d)}\quad\begin{minipage}[c]{5.3cm}\raggedright\Large $-\dfrac{\pi}{4}$\end{minipage}\end{tcolorbox}
\end{tcolorbox}
\newpage
\begin{tcolorbox}[enhanced,width=16cm,height=14cm,colback=white,colframe=ctri,boxrule=2pt,arc=0mm,
left=6mm,right=6mm,top=0mm,bottom=4mm,sharp corners,
overlay={\fill[ctri] (frame.north west) rectangle ([yshift=-1.25cm]frame.north east);
\node[anchor=west,text=white,font=\sffamily\bfseries\LARGE] at ([xshift=5mm,yshift=-0.62cm]frame.north west) {Question 60 B};
\node[anchor=east,text=white,font=\sffamily\itshape\large] at ([xshift=-5mm,yshift=-0.62cm]frame.north east) {Trigonométrie et exponentielle};}]
\vspace*{1.75cm}
\begin{minipage}[t][6.9cm][c]{\linewidth}\LARGE\centering
\begin{minipage}{0.95\linewidth}\raggedright\linespread{1.25}\selectfont Quelle est la mesure principale de l'angle $\dfrac{9\pi}{2}$~?\end{minipage}
\end{minipage}

\vfill
\begin{tcolorbox}[enhanced,nobeforeafter,width=7.1cm,height=2.1cm,valign=center,colback=ctri!6,colframe=ctri,arc=2.5mm,boxrule=0.9pt,left=3mm]{\Large\bfseries\color{ctri}a)}\quad\begin{minipage}[c]{5.3cm}\raggedright\Large $\dfrac{\pi}{4}$\end{minipage}\end{tcolorbox}\hfill\begin{tcolorbox}[enhanced,nobeforeafter,width=7.1cm,height=2.1cm,valign=center,colback=ctri!6,colframe=ctri,arc=2.5mm,boxrule=0.9pt,left=3mm]{\Large\bfseries\color{ctri}b)}\quad\begin{minipage}[c]{5.3cm}\raggedright\Large $\dfrac{\pi}{2}$\end{minipage}\end{tcolorbox}\par\vspace{2.5mm}
\begin{tcolorbox}[enhanced,nobeforeafter,width=7.1cm,height=2.1cm,valign=center,colback=ctri!6,colframe=ctri,arc=2.5mm,boxrule=0.9pt,left=3mm]{\Large\bfseries\color{ctri}c)}\quad\begin{minipage}[c]{5.3cm}\raggedright\Large $\dfrac{5\pi}{6}$\end{minipage}\end{tcolorbox}\hfill\begin{tcolorbox}[enhanced,nobeforeafter,width=7.1cm,height=2.1cm,valign=center,colback=ctri!6,colframe=ctri,arc=2.5mm,boxrule=0.9pt,left=3mm]{\Large\bfseries\color{ctri}d)}\quad\begin{minipage}[c]{5.3cm}\raggedright\Large $-\dfrac{\pi}{2}$\end{minipage}\end{tcolorbox}
\end{tcolorbox}
\newpage
\begin{tcolorbox}[enhanced,width=16cm,height=14cm,colback=white,colframe=ctri,boxrule=2pt,arc=0mm,
left=6mm,right=6mm,top=0mm,bottom=4mm,sharp corners,
overlay={\fill[ctri] (frame.north west) rectangle ([yshift=-1.25cm]frame.north east);
\node[anchor=west,text=white,font=\sffamily\bfseries\LARGE] at ([xshift=5mm,yshift=-0.62cm]frame.north west) {Question 61 B};
\node[anchor=east,text=white,font=\sffamily\itshape\large] at ([xshift=-5mm,yshift=-0.62cm]frame.north east) {Trigonométrie et exponentielle};}]
\vspace*{1.75cm}
\begin{minipage}[t][6.9cm][c]{\linewidth}\LARGE\centering
\begin{minipage}{0.95\linewidth}\raggedright\linespread{1.25}\selectfont On sait que $x\in[0\,;\,\pi]$ et $\cos x=\dfrac{5}{13}$. Que vaut $\sin x$~?\end{minipage}
\end{minipage}

\vfill
\begin{tcolorbox}[enhanced,nobeforeafter,width=7.1cm,height=2.1cm,valign=center,colback=ctri!6,colframe=ctri,arc=2.5mm,boxrule=0.9pt,left=3mm]{\Large\bfseries\color{ctri}a)}\quad\begin{minipage}[c]{5.3cm}\raggedright\Large $\dfrac{12}{13}$\end{minipage}\end{tcolorbox}\hfill\begin{tcolorbox}[enhanced,nobeforeafter,width=7.1cm,height=2.1cm,valign=center,colback=ctri!6,colframe=ctri,arc=2.5mm,boxrule=0.9pt,left=3mm]{\Large\bfseries\color{ctri}b)}\quad\begin{minipage}[c]{5.3cm}\raggedright\Large $\dfrac{144}{169}$\end{minipage}\end{tcolorbox}\par\vspace{2.5mm}
\begin{tcolorbox}[enhanced,nobeforeafter,width=7.1cm,height=2.1cm,valign=center,colback=ctri!6,colframe=ctri,arc=2.5mm,boxrule=0.9pt,left=3mm]{\Large\bfseries\color{ctri}c)}\quad\begin{minipage}[c]{5.3cm}\raggedright\Large $\dfrac{8}{13}$\end{minipage}\end{tcolorbox}\hfill\begin{tcolorbox}[enhanced,nobeforeafter,width=7.1cm,height=2.1cm,valign=center,colback=ctri!6,colframe=ctri,arc=2.5mm,boxrule=0.9pt,left=3mm]{\Large\bfseries\color{ctri}d)}\quad\begin{minipage}[c]{5.3cm}\raggedright\Large $-\dfrac{12}{13}$\end{minipage}\end{tcolorbox}
\end{tcolorbox}
\newpage
\begin{tcolorbox}[enhanced,width=16cm,height=14cm,colback=white,colframe=ctri,boxrule=2pt,arc=0mm,
left=6mm,right=6mm,top=0mm,bottom=4mm,sharp corners,
overlay={\fill[ctri] (frame.north west) rectangle ([yshift=-1.25cm]frame.north east);
\node[anchor=west,text=white,font=\sffamily\bfseries\LARGE] at ([xshift=5mm,yshift=-0.62cm]frame.north west) {Question 62 B};
\node[anchor=east,text=white,font=\sffamily\itshape\large] at ([xshift=-5mm,yshift=-0.62cm]frame.north east) {Trigonométrie et exponentielle};}]
\vspace*{1.75cm}
\begin{minipage}[t][6.9cm][c]{\linewidth}\LARGE\centering
\begin{minipage}{0.95\linewidth}\raggedright\linespread{1.25}\selectfont On sait que $x\in[0\,;\,\pi]$ et $\cos x=-\dfrac{8}{17}$. Que vaut $\sin x$~?\end{minipage}
\end{minipage}

\vfill
\begin{tcolorbox}[enhanced,nobeforeafter,width=7.1cm,height=2.1cm,valign=center,colback=ctri!6,colframe=ctri,arc=2.5mm,boxrule=0.9pt,left=3mm]{\Large\bfseries\color{ctri}a)}\quad\begin{minipage}[c]{5.3cm}\raggedright\Large $-\dfrac{15}{17}$\end{minipage}\end{tcolorbox}\hfill\begin{tcolorbox}[enhanced,nobeforeafter,width=7.1cm,height=2.1cm,valign=center,colback=ctri!6,colframe=ctri,arc=2.5mm,boxrule=0.9pt,left=3mm]{\Large\bfseries\color{ctri}b)}\quad\begin{minipage}[c]{5.3cm}\raggedright\Large $\dfrac{15}{17}$\end{minipage}\end{tcolorbox}\par\vspace{2.5mm}
\begin{tcolorbox}[enhanced,nobeforeafter,width=7.1cm,height=2.1cm,valign=center,colback=ctri!6,colframe=ctri,arc=2.5mm,boxrule=0.9pt,left=3mm]{\Large\bfseries\color{ctri}c)}\quad\begin{minipage}[c]{5.3cm}\raggedright\Large $\dfrac{225}{289}$\end{minipage}\end{tcolorbox}\hfill\begin{tcolorbox}[enhanced,nobeforeafter,width=7.1cm,height=2.1cm,valign=center,colback=ctri!6,colframe=ctri,arc=2.5mm,boxrule=0.9pt,left=3mm]{\Large\bfseries\color{ctri}d)}\quad\begin{minipage}[c]{5.3cm}\raggedright\Large $\dfrac{25}{17}$\end{minipage}\end{tcolorbox}
\end{tcolorbox}
\newpage
\begin{tcolorbox}[enhanced,width=16cm,height=14cm,colback=white,colframe=ctri,boxrule=2pt,arc=0mm,
left=6mm,right=6mm,top=0mm,bottom=4mm,sharp corners,
overlay={\fill[ctri] (frame.north west) rectangle ([yshift=-1.25cm]frame.north east);
\node[anchor=west,text=white,font=\sffamily\bfseries\LARGE] at ([xshift=5mm,yshift=-0.62cm]frame.north west) {Question 63 B};
\node[anchor=east,text=white,font=\sffamily\itshape\large] at ([xshift=-5mm,yshift=-0.62cm]frame.north east) {Trigonométrie et exponentielle};}]
\vspace*{1.75cm}
\begin{minipage}[t][6.9cm][c]{\linewidth}\LARGE\centering
\begin{minipage}{0.95\linewidth}\raggedright\linespread{1.25}\selectfont Simplifier $\left(e^{3}\right)^{2}\times e^{-7}$.\end{minipage}
\end{minipage}

\vfill
\begin{tcolorbox}[enhanced,nobeforeafter,width=7.1cm,height=2.1cm,valign=center,colback=ctri!6,colframe=ctri,arc=2.5mm,boxrule=0.9pt,left=3mm]{\Large\bfseries\color{ctri}a)}\quad\begin{minipage}[c]{5.3cm}\raggedright\Large $e^{-2}$\end{minipage}\end{tcolorbox}\hfill\begin{tcolorbox}[enhanced,nobeforeafter,width=7.1cm,height=2.1cm,valign=center,colback=ctri!6,colframe=ctri,arc=2.5mm,boxrule=0.9pt,left=3mm]{\Large\bfseries\color{ctri}b)}\quad\begin{minipage}[c]{5.3cm}\raggedright\Large $e^{-15}$\end{minipage}\end{tcolorbox}\par\vspace{2.5mm}
\begin{tcolorbox}[enhanced,nobeforeafter,width=7.1cm,height=2.1cm,valign=center,colback=ctri!6,colframe=ctri,arc=2.5mm,boxrule=0.9pt,left=3mm]{\Large\bfseries\color{ctri}c)}\quad\begin{minipage}[c]{5.3cm}\raggedright\Large $e^{-1}$\end{minipage}\end{tcolorbox}\hfill\begin{tcolorbox}[enhanced,nobeforeafter,width=7.1cm,height=2.1cm,valign=center,colback=ctri!6,colframe=ctri,arc=2.5mm,boxrule=0.9pt,left=3mm]{\Large\bfseries\color{ctri}d)}\quad\begin{minipage}[c]{5.3cm}\raggedright\Large $e^{2}$\end{minipage}\end{tcolorbox}
\end{tcolorbox}
\newpage
\begin{tcolorbox}[enhanced,width=16cm,height=14cm,colback=white,colframe=ctri,boxrule=2pt,arc=0mm,
left=6mm,right=6mm,top=0mm,bottom=4mm,sharp corners,
overlay={\fill[ctri] (frame.north west) rectangle ([yshift=-1.25cm]frame.north east);
\node[anchor=west,text=white,font=\sffamily\bfseries\LARGE] at ([xshift=5mm,yshift=-0.62cm]frame.north west) {Question 64 B};
\node[anchor=east,text=white,font=\sffamily\itshape\large] at ([xshift=-5mm,yshift=-0.62cm]frame.north east) {Trigonométrie et exponentielle};}]
\vspace*{1.75cm}
\begin{minipage}[t][6.9cm][c]{\linewidth}\LARGE\centering
\begin{minipage}{0.95\linewidth}\raggedright\linespread{1.25}\selectfont Simplifier $\left(e^{1}\right)^{5}\times e^{-5}$.\end{minipage}
\end{minipage}

\vfill
\begin{tcolorbox}[enhanced,nobeforeafter,width=7.1cm,height=2.1cm,valign=center,colback=ctri!6,colframe=ctri,arc=2.5mm,boxrule=0.9pt,left=3mm]{\Large\bfseries\color{ctri}a)}\quad\begin{minipage}[c]{5.3cm}\raggedright\Large $1$\end{minipage}\end{tcolorbox}\hfill\begin{tcolorbox}[enhanced,nobeforeafter,width=7.1cm,height=2.1cm,valign=center,colback=ctri!6,colframe=ctri,arc=2.5mm,boxrule=0.9pt,left=3mm]{\Large\bfseries\color{ctri}b)}\quad\begin{minipage}[c]{5.3cm}\raggedright\Large $e$\end{minipage}\end{tcolorbox}\par\vspace{2.5mm}
\begin{tcolorbox}[enhanced,nobeforeafter,width=7.1cm,height=2.1cm,valign=center,colback=ctri!6,colframe=ctri,arc=2.5mm,boxrule=0.9pt,left=3mm]{\Large\bfseries\color{ctri}c)}\quad\begin{minipage}[c]{5.3cm}\raggedright\Large $e^{-4}$\end{minipage}\end{tcolorbox}\hfill\begin{tcolorbox}[enhanced,nobeforeafter,width=7.1cm,height=2.1cm,valign=center,colback=ctri!6,colframe=ctri,arc=2.5mm,boxrule=0.9pt,left=3mm]{\Large\bfseries\color{ctri}d)}\quad\begin{minipage}[c]{5.3cm}\raggedright\Large $e^{2}$\end{minipage}\end{tcolorbox}
\end{tcolorbox}
\newpage
\begin{tcolorbox}[enhanced,width=16cm,height=14cm,colback=white,colframe=cpsc,boxrule=2pt,arc=0mm,
left=6mm,right=6mm,top=0mm,bottom=4mm,sharp corners,
overlay={\fill[cpsc] (frame.north west) rectangle ([yshift=-1.25cm]frame.north east);
\node[anchor=west,text=white,font=\sffamily\bfseries\LARGE] at ([xshift=5mm,yshift=-0.62cm]frame.north west) {Question 65 B};
\node[anchor=east,text=white,font=\sffamily\itshape\large] at ([xshift=-5mm,yshift=-0.62cm]frame.north east) {Produit scalaire et géométrie repérée};}]
\vspace*{1.75cm}
\begin{minipage}[t][6.9cm][c]{\linewidth}\LARGE\centering
\begin{minipage}{0.95\linewidth}\raggedright\linespread{1.25}\selectfont Dans un repère orthonormé, que vaut $\vec u\cdot\vec v$ pour $\vec u\begin{pmatrix}4\\-2\end{pmatrix}$ et $\vec v\begin{pmatrix}1\\2\end{pmatrix}$~?\end{minipage}
\end{minipage}

\vfill
\begin{tcolorbox}[enhanced,nobeforeafter,width=7.1cm,height=2.1cm,valign=center,colback=cpsc!6,colframe=cpsc,arc=2.5mm,boxrule=0.9pt,left=3mm]{\Large\bfseries\color{cpsc}a)}\quad\begin{minipage}[c]{5.3cm}\raggedright\Large $6$\end{minipage}\end{tcolorbox}\hfill\begin{tcolorbox}[enhanced,nobeforeafter,width=7.1cm,height=2.1cm,valign=center,colback=cpsc!6,colframe=cpsc,arc=2.5mm,boxrule=0.9pt,left=3mm]{\Large\bfseries\color{cpsc}b)}\quad\begin{minipage}[c]{5.3cm}\raggedright\Large $8$\end{minipage}\end{tcolorbox}\par\vspace{2.5mm}
\begin{tcolorbox}[enhanced,nobeforeafter,width=7.1cm,height=2.1cm,valign=center,colback=cpsc!6,colframe=cpsc,arc=2.5mm,boxrule=0.9pt,left=3mm]{\Large\bfseries\color{cpsc}c)}\quad\begin{minipage}[c]{5.3cm}\raggedright\Large $0$\end{minipage}\end{tcolorbox}\hfill\begin{tcolorbox}[enhanced,nobeforeafter,width=7.1cm,height=2.1cm,valign=center,colback=cpsc!6,colframe=cpsc,arc=2.5mm,boxrule=0.9pt,left=3mm]{\Large\bfseries\color{cpsc}d)}\quad\begin{minipage}[c]{5.3cm}\raggedright\Large $10$\end{minipage}\end{tcolorbox}
\end{tcolorbox}
\newpage
\begin{tcolorbox}[enhanced,width=16cm,height=14cm,colback=white,colframe=cpsc,boxrule=2pt,arc=0mm,
left=6mm,right=6mm,top=0mm,bottom=4mm,sharp corners,
overlay={\fill[cpsc] (frame.north west) rectangle ([yshift=-1.25cm]frame.north east);
\node[anchor=west,text=white,font=\sffamily\bfseries\LARGE] at ([xshift=5mm,yshift=-0.62cm]frame.north west) {Question 66 B};
\node[anchor=east,text=white,font=\sffamily\itshape\large] at ([xshift=-5mm,yshift=-0.62cm]frame.north east) {Produit scalaire et géométrie repérée};}]
\vspace*{1.75cm}
\begin{minipage}[t][6.9cm][c]{\linewidth}\LARGE\centering
\begin{minipage}{0.95\linewidth}\raggedright\linespread{1.25}\selectfont Dans un repère orthonormé, que vaut $\vec u\cdot\vec v$ pour $\vec u\begin{pmatrix}3\\1\end{pmatrix}$ et $\vec v\begin{pmatrix}-2\\5\end{pmatrix}$~?\end{minipage}
\end{minipage}

\vfill
\begin{tcolorbox}[enhanced,nobeforeafter,width=7.1cm,height=2.1cm,valign=center,colback=cpsc!6,colframe=cpsc,arc=2.5mm,boxrule=0.9pt,left=3mm]{\Large\bfseries\color{cpsc}a)}\quad\begin{minipage}[c]{5.3cm}\raggedright\Large $-1$\end{minipage}\end{tcolorbox}\hfill\begin{tcolorbox}[enhanced,nobeforeafter,width=7.1cm,height=2.1cm,valign=center,colback=cpsc!6,colframe=cpsc,arc=2.5mm,boxrule=0.9pt,left=3mm]{\Large\bfseries\color{cpsc}b)}\quad\begin{minipage}[c]{5.3cm}\raggedright\Large $13$\end{minipage}\end{tcolorbox}\par\vspace{2.5mm}
\begin{tcolorbox}[enhanced,nobeforeafter,width=7.1cm,height=2.1cm,valign=center,colback=cpsc!6,colframe=cpsc,arc=2.5mm,boxrule=0.9pt,left=3mm]{\Large\bfseries\color{cpsc}c)}\quad\begin{minipage}[c]{5.3cm}\raggedright\Large $17$\end{minipage}\end{tcolorbox}\hfill\begin{tcolorbox}[enhanced,nobeforeafter,width=7.1cm,height=2.1cm,valign=center,colback=cpsc!6,colframe=cpsc,arc=2.5mm,boxrule=0.9pt,left=3mm]{\Large\bfseries\color{cpsc}d)}\quad\begin{minipage}[c]{5.3cm}\raggedright\Large $-11$\end{minipage}\end{tcolorbox}
\end{tcolorbox}
\newpage
\begin{tcolorbox}[enhanced,width=16cm,height=14cm,colback=white,colframe=cpsc,boxrule=2pt,arc=0mm,
left=6mm,right=6mm,top=0mm,bottom=4mm,sharp corners,
overlay={\fill[cpsc] (frame.north west) rectangle ([yshift=-1.25cm]frame.north east);
\node[anchor=west,text=white,font=\sffamily\bfseries\LARGE] at ([xshift=5mm,yshift=-0.62cm]frame.north west) {Question 67 B};
\node[anchor=east,text=white,font=\sffamily\itshape\large] at ([xshift=-5mm,yshift=-0.62cm]frame.north east) {Produit scalaire et géométrie repérée};}]
\vspace*{1.75cm}
\begin{minipage}[t][6.9cm][c]{\linewidth}\LARGE\centering
\begin{minipage}{0.95\linewidth}\raggedright\linespread{1.25}\selectfont $\|\vec u\|=6$, $\|\vec v\|=1$ et l'angle entre $\vec u$ et $\vec v$ mesure $90^\circ$. Que vaut $\vec u\cdot\vec v$~?\end{minipage}
\end{minipage}

\vfill
\begin{tcolorbox}[enhanced,nobeforeafter,width=7.1cm,height=2.1cm,valign=center,colback=cpsc!6,colframe=cpsc,arc=2.5mm,boxrule=0.9pt,left=3mm]{\Large\bfseries\color{cpsc}a)}\quad\begin{minipage}[c]{5.3cm}\raggedright\Large $6$\end{minipage}\end{tcolorbox}\hfill\begin{tcolorbox}[enhanced,nobeforeafter,width=7.1cm,height=2.1cm,valign=center,colback=cpsc!6,colframe=cpsc,arc=2.5mm,boxrule=0.9pt,left=3mm]{\Large\bfseries\color{cpsc}b)}\quad\begin{minipage}[c]{5.3cm}\raggedright\Large $0$\end{minipage}\end{tcolorbox}\par\vspace{2.5mm}
\begin{tcolorbox}[enhanced,nobeforeafter,width=7.1cm,height=2.1cm,valign=center,colback=cpsc!6,colframe=cpsc,arc=2.5mm,boxrule=0.9pt,left=3mm]{\Large\bfseries\color{cpsc}c)}\quad\begin{minipage}[c]{5.3cm}\raggedright\Large $3$\end{minipage}\end{tcolorbox}\hfill\begin{tcolorbox}[enhanced,nobeforeafter,width=7.1cm,height=2.1cm,valign=center,colback=cpsc!6,colframe=cpsc,arc=2.5mm,boxrule=0.9pt,left=3mm]{\Large\bfseries\color{cpsc}d)}\quad\begin{minipage}[c]{5.3cm}\raggedright\Large $7$\end{minipage}\end{tcolorbox}
\end{tcolorbox}
\newpage
\begin{tcolorbox}[enhanced,width=16cm,height=14cm,colback=white,colframe=cpsc,boxrule=2pt,arc=0mm,
left=6mm,right=6mm,top=0mm,bottom=4mm,sharp corners,
overlay={\fill[cpsc] (frame.north west) rectangle ([yshift=-1.25cm]frame.north east);
\node[anchor=west,text=white,font=\sffamily\bfseries\LARGE] at ([xshift=5mm,yshift=-0.62cm]frame.north west) {Question 68 B};
\node[anchor=east,text=white,font=\sffamily\itshape\large] at ([xshift=-5mm,yshift=-0.62cm]frame.north east) {Produit scalaire et géométrie repérée};}]
\vspace*{1.75cm}
\begin{minipage}[t][6.9cm][c]{\linewidth}\LARGE\centering
\begin{minipage}{0.95\linewidth}\raggedright\linespread{1.25}\selectfont $\|\vec u\|=4$, $\|\vec v\|=4$ et l'angle entre $\vec u$ et $\vec v$ mesure $0^\circ$. Que vaut $\vec u\cdot\vec v$~?\end{minipage}
\end{minipage}

\vfill
\begin{tcolorbox}[enhanced,nobeforeafter,width=7.1cm,height=2.1cm,valign=center,colback=cpsc!6,colframe=cpsc,arc=2.5mm,boxrule=0.9pt,left=3mm]{\Large\bfseries\color{cpsc}a)}\quad\begin{minipage}[c]{5.3cm}\raggedright\Large $16$\end{minipage}\end{tcolorbox}\hfill\begin{tcolorbox}[enhanced,nobeforeafter,width=7.1cm,height=2.1cm,valign=center,colback=cpsc!6,colframe=cpsc,arc=2.5mm,boxrule=0.9pt,left=3mm]{\Large\bfseries\color{cpsc}b)}\quad\begin{minipage}[c]{5.3cm}\raggedright\Large $-16$\end{minipage}\end{tcolorbox}\par\vspace{2.5mm}
\begin{tcolorbox}[enhanced,nobeforeafter,width=7.1cm,height=2.1cm,valign=center,colback=cpsc!6,colframe=cpsc,arc=2.5mm,boxrule=0.9pt,left=3mm]{\Large\bfseries\color{cpsc}c)}\quad\begin{minipage}[c]{5.3cm}\raggedright\Large $17$\end{minipage}\end{tcolorbox}\hfill\begin{tcolorbox}[enhanced,nobeforeafter,width=7.1cm,height=2.1cm,valign=center,colback=cpsc!6,colframe=cpsc,arc=2.5mm,boxrule=0.9pt,left=3mm]{\Large\bfseries\color{cpsc}d)}\quad\begin{minipage}[c]{5.3cm}\raggedright\Large $8$\end{minipage}\end{tcolorbox}
\end{tcolorbox}
\newpage
\begin{tcolorbox}[enhanced,width=16cm,height=14cm,colback=white,colframe=cpsc,boxrule=2pt,arc=0mm,
left=6mm,right=6mm,top=0mm,bottom=4mm,sharp corners,
overlay={\fill[cpsc] (frame.north west) rectangle ([yshift=-1.25cm]frame.north east);
\node[anchor=west,text=white,font=\sffamily\bfseries\LARGE] at ([xshift=5mm,yshift=-0.62cm]frame.north west) {Question 69 B};
\node[anchor=east,text=white,font=\sffamily\itshape\large] at ([xshift=-5mm,yshift=-0.62cm]frame.north east) {Produit scalaire et géométrie repérée};}]
\vspace*{1.75cm}
\begin{minipage}[t][6.9cm][c]{\linewidth}\LARGE\centering
\begin{minipage}{0.95\linewidth}\raggedright\linespread{1.25}\selectfont Pour quelle valeur de $m$ les vecteurs $\vec u\begin{pmatrix}-3\\m\end{pmatrix}$ et $\vec v\begin{pmatrix}2\\3\end{pmatrix}$ sont-ils orthogonaux~?\end{minipage}
\end{minipage}

\vfill
\begin{tcolorbox}[enhanced,nobeforeafter,width=7.1cm,height=2.1cm,valign=center,colback=cpsc!6,colframe=cpsc,arc=2.5mm,boxrule=0.9pt,left=3mm]{\Large\bfseries\color{cpsc}a)}\quad\begin{minipage}[c]{5.3cm}\raggedright\Large $m=6$\end{minipage}\end{tcolorbox}\hfill\begin{tcolorbox}[enhanced,nobeforeafter,width=7.1cm,height=2.1cm,valign=center,colback=cpsc!6,colframe=cpsc,arc=2.5mm,boxrule=0.9pt,left=3mm]{\Large\bfseries\color{cpsc}b)}\quad\begin{minipage}[c]{5.3cm}\raggedright\Large $m=-2$\end{minipage}\end{tcolorbox}\par\vspace{2.5mm}
\begin{tcolorbox}[enhanced,nobeforeafter,width=7.1cm,height=2.1cm,valign=center,colback=cpsc!6,colframe=cpsc,arc=2.5mm,boxrule=0.9pt,left=3mm]{\Large\bfseries\color{cpsc}c)}\quad\begin{minipage}[c]{5.3cm}\raggedright\Large $m=\dfrac{1}{2}$\end{minipage}\end{tcolorbox}\hfill\begin{tcolorbox}[enhanced,nobeforeafter,width=7.1cm,height=2.1cm,valign=center,colback=cpsc!6,colframe=cpsc,arc=2.5mm,boxrule=0.9pt,left=3mm]{\Large\bfseries\color{cpsc}d)}\quad\begin{minipage}[c]{5.3cm}\raggedright\Large $m=2$\end{minipage}\end{tcolorbox}
\end{tcolorbox}
\newpage
\begin{tcolorbox}[enhanced,width=16cm,height=14cm,colback=white,colframe=cpsc,boxrule=2pt,arc=0mm,
left=6mm,right=6mm,top=0mm,bottom=4mm,sharp corners,
overlay={\fill[cpsc] (frame.north west) rectangle ([yshift=-1.25cm]frame.north east);
\node[anchor=west,text=white,font=\sffamily\bfseries\LARGE] at ([xshift=5mm,yshift=-0.62cm]frame.north west) {Question 70 B};
\node[anchor=east,text=white,font=\sffamily\itshape\large] at ([xshift=-5mm,yshift=-0.62cm]frame.north east) {Produit scalaire et géométrie repérée};}]
\vspace*{1.75cm}
\begin{minipage}[t][6.9cm][c]{\linewidth}\LARGE\centering
\begin{minipage}{0.95\linewidth}\raggedright\linespread{1.25}\selectfont Pour quelle valeur de $m$ les vecteurs $\vec u\begin{pmatrix}5\\m\end{pmatrix}$ et $\vec v\begin{pmatrix}2\\-4\end{pmatrix}$ sont-ils orthogonaux~?\end{minipage}
\end{minipage}

\vfill
\begin{tcolorbox}[enhanced,nobeforeafter,width=7.1cm,height=2.1cm,valign=center,colback=cpsc!6,colframe=cpsc,arc=2.5mm,boxrule=0.9pt,left=3mm]{\Large\bfseries\color{cpsc}a)}\quad\begin{minipage}[c]{5.3cm}\raggedright\Large $m=\dfrac{5}{2}$\end{minipage}\end{tcolorbox}\hfill\begin{tcolorbox}[enhanced,nobeforeafter,width=7.1cm,height=2.1cm,valign=center,colback=cpsc!6,colframe=cpsc,arc=2.5mm,boxrule=0.9pt,left=3mm]{\Large\bfseries\color{cpsc}b)}\quad\begin{minipage}[c]{5.3cm}\raggedright\Large $m=-10$\end{minipage}\end{tcolorbox}\par\vspace{2.5mm}
\begin{tcolorbox}[enhanced,nobeforeafter,width=7.1cm,height=2.1cm,valign=center,colback=cpsc!6,colframe=cpsc,arc=2.5mm,boxrule=0.9pt,left=3mm]{\Large\bfseries\color{cpsc}c)}\quad\begin{minipage}[c]{5.3cm}\raggedright\Large $m=\dfrac{2}{5}$\end{minipage}\end{tcolorbox}\hfill\begin{tcolorbox}[enhanced,nobeforeafter,width=7.1cm,height=2.1cm,valign=center,colback=cpsc!6,colframe=cpsc,arc=2.5mm,boxrule=0.9pt,left=3mm]{\Large\bfseries\color{cpsc}d)}\quad\begin{minipage}[c]{5.3cm}\raggedright\Large $m=-\dfrac{5}{2}$\end{minipage}\end{tcolorbox}
\end{tcolorbox}
\newpage
\begin{tcolorbox}[enhanced,width=16cm,height=14cm,colback=white,colframe=cpsc,boxrule=2pt,arc=0mm,
left=6mm,right=6mm,top=0mm,bottom=4mm,sharp corners,
overlay={\fill[cpsc] (frame.north west) rectangle ([yshift=-1.25cm]frame.north east);
\node[anchor=west,text=white,font=\sffamily\bfseries\LARGE] at ([xshift=5mm,yshift=-0.62cm]frame.north west) {Question 71 B};
\node[anchor=east,text=white,font=\sffamily\itshape\large] at ([xshift=-5mm,yshift=-0.62cm]frame.north east) {Produit scalaire et géométrie repérée};}]
\vspace*{1.75cm}
\begin{minipage}[t][6.9cm][c]{\linewidth}\LARGE\centering
\begin{minipage}{0.95\linewidth}\raggedright\linespread{1.25}\selectfont Dans le triangle $ABC$, $AB=4$, $AC=5$ et $\widehat{BAC}=120^\circ$. Combien mesure $BC$~?\end{minipage}
\end{minipage}

\vfill
\begin{tcolorbox}[enhanced,nobeforeafter,width=7.1cm,height=2.1cm,valign=center,colback=cpsc!6,colframe=cpsc,arc=2.5mm,boxrule=0.9pt,left=3mm]{\Large\bfseries\color{cpsc}a)}\quad\begin{minipage}[c]{5.3cm}\raggedright\Large $\sqrt{41}$\end{minipage}\end{tcolorbox}\hfill\begin{tcolorbox}[enhanced,nobeforeafter,width=7.1cm,height=2.1cm,valign=center,colback=cpsc!6,colframe=cpsc,arc=2.5mm,boxrule=0.9pt,left=3mm]{\Large\bfseries\color{cpsc}b)}\quad\begin{minipage}[c]{5.3cm}\raggedright\Large $\sqrt{21}$\end{minipage}\end{tcolorbox}\par\vspace{2.5mm}
\begin{tcolorbox}[enhanced,nobeforeafter,width=7.1cm,height=2.1cm,valign=center,colback=cpsc!6,colframe=cpsc,arc=2.5mm,boxrule=0.9pt,left=3mm]{\Large\bfseries\color{cpsc}c)}\quad\begin{minipage}[c]{5.3cm}\raggedright\Large $61$\end{minipage}\end{tcolorbox}\hfill\begin{tcolorbox}[enhanced,nobeforeafter,width=7.1cm,height=2.1cm,valign=center,colback=cpsc!6,colframe=cpsc,arc=2.5mm,boxrule=0.9pt,left=3mm]{\Large\bfseries\color{cpsc}d)}\quad\begin{minipage}[c]{5.3cm}\raggedright\Large $\sqrt{61}$\end{minipage}\end{tcolorbox}
\end{tcolorbox}
\newpage
\begin{tcolorbox}[enhanced,width=16cm,height=14cm,colback=white,colframe=cpsc,boxrule=2pt,arc=0mm,
left=6mm,right=6mm,top=0mm,bottom=4mm,sharp corners,
overlay={\fill[cpsc] (frame.north west) rectangle ([yshift=-1.25cm]frame.north east);
\node[anchor=west,text=white,font=\sffamily\bfseries\LARGE] at ([xshift=5mm,yshift=-0.62cm]frame.north west) {Question 72 B};
\node[anchor=east,text=white,font=\sffamily\itshape\large] at ([xshift=-5mm,yshift=-0.62cm]frame.north east) {Produit scalaire et géométrie repérée};}]
\vspace*{1.75cm}
\begin{minipage}[t][6.9cm][c]{\linewidth}\LARGE\centering
\begin{minipage}{0.95\linewidth}\raggedright\linespread{1.25}\selectfont Dans le triangle $ABC$, $AB=5$, $AC=8$ et $\widehat{BAC}=60^\circ$. Combien mesure $BC$~?\end{minipage}
\end{minipage}

\vfill
\begin{tcolorbox}[enhanced,nobeforeafter,width=7.1cm,height=2.1cm,valign=center,colback=cpsc!6,colframe=cpsc,arc=2.5mm,boxrule=0.9pt,left=3mm]{\Large\bfseries\color{cpsc}a)}\quad\begin{minipage}[c]{5.3cm}\raggedright\Large $\sqrt{89}$\end{minipage}\end{tcolorbox}\hfill\begin{tcolorbox}[enhanced,nobeforeafter,width=7.1cm,height=2.1cm,valign=center,colback=cpsc!6,colframe=cpsc,arc=2.5mm,boxrule=0.9pt,left=3mm]{\Large\bfseries\color{cpsc}b)}\quad\begin{minipage}[c]{5.3cm}\raggedright\Large $\sqrt{129}$\end{minipage}\end{tcolorbox}\par\vspace{2.5mm}
\begin{tcolorbox}[enhanced,nobeforeafter,width=7.1cm,height=2.1cm,valign=center,colback=cpsc!6,colframe=cpsc,arc=2.5mm,boxrule=0.9pt,left=3mm]{\Large\bfseries\color{cpsc}c)}\quad\begin{minipage}[c]{5.3cm}\raggedright\Large $49$\end{minipage}\end{tcolorbox}\hfill\begin{tcolorbox}[enhanced,nobeforeafter,width=7.1cm,height=2.1cm,valign=center,colback=cpsc!6,colframe=cpsc,arc=2.5mm,boxrule=0.9pt,left=3mm]{\Large\bfseries\color{cpsc}d)}\quad\begin{minipage}[c]{5.3cm}\raggedright\Large $7$\end{minipage}\end{tcolorbox}
\end{tcolorbox}
\newpage
\begin{tcolorbox}[enhanced,width=16cm,height=14cm,colback=white,colframe=cpsc,boxrule=2pt,arc=0mm,
left=6mm,right=6mm,top=0mm,bottom=4mm,sharp corners,
overlay={\fill[cpsc] (frame.north west) rectangle ([yshift=-1.25cm]frame.north east);
\node[anchor=west,text=white,font=\sffamily\bfseries\LARGE] at ([xshift=5mm,yshift=-0.62cm]frame.north west) {Question 73 B};
\node[anchor=east,text=white,font=\sffamily\itshape\large] at ([xshift=-5mm,yshift=-0.62cm]frame.north east) {Produit scalaire et géométrie repérée};}]
\vspace*{1.75cm}
\begin{minipage}[t][6.9cm][c]{\linewidth}\LARGE\centering
\begin{minipage}{0.95\linewidth}\raggedright\linespread{1.25}\selectfont Lequel de ces vecteurs est un vecteur normal à la droite d'équation $-5x+y+3=0$~?\end{minipage}
\end{minipage}

\vfill
\begin{tcolorbox}[enhanced,nobeforeafter,width=7.1cm,height=2.1cm,valign=center,colback=cpsc!6,colframe=cpsc,arc=2.5mm,boxrule=0.9pt,left=3mm]{\Large\bfseries\color{cpsc}a)}\quad\begin{minipage}[c]{5.3cm}\raggedright\large $\begin{pmatrix}-1\\-5\end{pmatrix}$\end{minipage}\end{tcolorbox}\hfill\begin{tcolorbox}[enhanced,nobeforeafter,width=7.1cm,height=2.1cm,valign=center,colback=cpsc!6,colframe=cpsc,arc=2.5mm,boxrule=0.9pt,left=3mm]{\Large\bfseries\color{cpsc}b)}\quad\begin{minipage}[c]{5.3cm}\raggedright\large $\begin{pmatrix}-5\\1\end{pmatrix}$\end{minipage}\end{tcolorbox}\par\vspace{2.5mm}
\begin{tcolorbox}[enhanced,nobeforeafter,width=7.1cm,height=2.1cm,valign=center,colback=cpsc!6,colframe=cpsc,arc=2.5mm,boxrule=0.9pt,left=3mm]{\Large\bfseries\color{cpsc}c)}\quad\begin{minipage}[c]{5.3cm}\raggedright\large $\begin{pmatrix}1\\-5\end{pmatrix}$\end{minipage}\end{tcolorbox}\hfill\begin{tcolorbox}[enhanced,nobeforeafter,width=7.1cm,height=2.1cm,valign=center,colback=cpsc!6,colframe=cpsc,arc=2.5mm,boxrule=0.9pt,left=3mm]{\Large\bfseries\color{cpsc}d)}\quad\begin{minipage}[c]{5.3cm}\raggedright\large $\begin{pmatrix}-5\\3\end{pmatrix}$\end{minipage}\end{tcolorbox}
\end{tcolorbox}
\newpage
\begin{tcolorbox}[enhanced,width=16cm,height=14cm,colback=white,colframe=cpsc,boxrule=2pt,arc=0mm,
left=6mm,right=6mm,top=0mm,bottom=4mm,sharp corners,
overlay={\fill[cpsc] (frame.north west) rectangle ([yshift=-1.25cm]frame.north east);
\node[anchor=west,text=white,font=\sffamily\bfseries\LARGE] at ([xshift=5mm,yshift=-0.62cm]frame.north west) {Question 74 B};
\node[anchor=east,text=white,font=\sffamily\itshape\large] at ([xshift=-5mm,yshift=-0.62cm]frame.north east) {Produit scalaire et géométrie repérée};}]
\vspace*{1.75cm}
\begin{minipage}[t][6.9cm][c]{\linewidth}\LARGE\centering
\begin{minipage}{0.95\linewidth}\raggedright\linespread{1.25}\selectfont Lequel de ces vecteurs est un vecteur normal à la droite d'équation $3x+3y-1=0$~?\end{minipage}
\end{minipage}

\vfill
\begin{tcolorbox}[enhanced,nobeforeafter,width=7.1cm,height=2.1cm,valign=center,colback=cpsc!6,colframe=cpsc,arc=2.5mm,boxrule=0.9pt,left=3mm]{\Large\bfseries\color{cpsc}a)}\quad\begin{minipage}[c]{5.3cm}\raggedright\large $\begin{pmatrix}3\\3\end{pmatrix}$\end{minipage}\end{tcolorbox}\hfill\begin{tcolorbox}[enhanced,nobeforeafter,width=7.1cm,height=2.1cm,valign=center,colback=cpsc!6,colframe=cpsc,arc=2.5mm,boxrule=0.9pt,left=3mm]{\Large\bfseries\color{cpsc}b)}\quad\begin{minipage}[c]{5.3cm}\raggedright\large $\begin{pmatrix}3\\-1\end{pmatrix}$\end{minipage}\end{tcolorbox}\par\vspace{2.5mm}
\begin{tcolorbox}[enhanced,nobeforeafter,width=7.1cm,height=2.1cm,valign=center,colback=cpsc!6,colframe=cpsc,arc=2.5mm,boxrule=0.9pt,left=3mm]{\Large\bfseries\color{cpsc}c)}\quad\begin{minipage}[c]{5.3cm}\raggedright\large $\begin{pmatrix}3\\-3\end{pmatrix}$\end{minipage}\end{tcolorbox}\hfill\begin{tcolorbox}[enhanced,nobeforeafter,width=7.1cm,height=2.1cm,valign=center,colback=cpsc!6,colframe=cpsc,arc=2.5mm,boxrule=0.9pt,left=3mm]{\Large\bfseries\color{cpsc}d)}\quad\begin{minipage}[c]{5.3cm}\raggedright\large $\begin{pmatrix}-3\\3\end{pmatrix}$\end{minipage}\end{tcolorbox}
\end{tcolorbox}
\newpage
\begin{tcolorbox}[enhanced,width=16cm,height=14cm,colback=white,colframe=cpsc,boxrule=2pt,arc=0mm,
left=6mm,right=6mm,top=0mm,bottom=4mm,sharp corners,
overlay={\fill[cpsc] (frame.north west) rectangle ([yshift=-1.25cm]frame.north east);
\node[anchor=west,text=white,font=\sffamily\bfseries\LARGE] at ([xshift=5mm,yshift=-0.62cm]frame.north west) {Question 75 B};
\node[anchor=east,text=white,font=\sffamily\itshape\large] at ([xshift=-5mm,yshift=-0.62cm]frame.north east) {Produit scalaire et géométrie repérée};}]
\vspace*{1.75cm}
\begin{minipage}[t][6.9cm][c]{\linewidth}\LARGE\centering
\begin{minipage}{0.95\linewidth}\raggedright\linespread{1.25}\selectfont Quelle est une équation cartésienne de la droite passant par $A(2\,;\,-5)$ et de vecteur normal $\vec n\begin{pmatrix}1\\1\end{pmatrix}$~?\end{minipage}
\end{minipage}

\vfill
\begin{tcolorbox}[enhanced,nobeforeafter,width=7.1cm,height=2.1cm,valign=center,colback=cpsc!6,colframe=cpsc,arc=2.5mm,boxrule=0.9pt,left=3mm]{\Large\bfseries\color{cpsc}a)}\quad\begin{minipage}[c]{5.3cm}\raggedright\Large $x+y+3=0$\end{minipage}\end{tcolorbox}\hfill\begin{tcolorbox}[enhanced,nobeforeafter,width=7.1cm,height=2.1cm,valign=center,colback=cpsc!6,colframe=cpsc,arc=2.5mm,boxrule=0.9pt,left=3mm]{\Large\bfseries\color{cpsc}b)}\quad\begin{minipage}[c]{5.3cm}\raggedright\Large $x+y-3=0$\end{minipage}\end{tcolorbox}\par\vspace{2.5mm}
\begin{tcolorbox}[enhanced,nobeforeafter,width=7.1cm,height=2.1cm,valign=center,colback=cpsc!6,colframe=cpsc,arc=2.5mm,boxrule=0.9pt,left=3mm]{\Large\bfseries\color{cpsc}c)}\quad\begin{minipage}[c]{5.3cm}\raggedright\Large $-x+y+7=0$\end{minipage}\end{tcolorbox}\hfill\begin{tcolorbox}[enhanced,nobeforeafter,width=7.1cm,height=2.1cm,valign=center,colback=cpsc!6,colframe=cpsc,arc=2.5mm,boxrule=0.9pt,left=3mm]{\Large\bfseries\color{cpsc}d)}\quad\begin{minipage}[c]{5.3cm}\raggedright\Large $2x-5y+3=0$\end{minipage}\end{tcolorbox}
\end{tcolorbox}
\newpage
\begin{tcolorbox}[enhanced,width=16cm,height=14cm,colback=white,colframe=cpsc,boxrule=2pt,arc=0mm,
left=6mm,right=6mm,top=0mm,bottom=4mm,sharp corners,
overlay={\fill[cpsc] (frame.north west) rectangle ([yshift=-1.25cm]frame.north east);
\node[anchor=west,text=white,font=\sffamily\bfseries\LARGE] at ([xshift=5mm,yshift=-0.62cm]frame.north west) {Question 76 B};
\node[anchor=east,text=white,font=\sffamily\itshape\large] at ([xshift=-5mm,yshift=-0.62cm]frame.north east) {Produit scalaire et géométrie repérée};}]
\vspace*{1.75cm}
\begin{minipage}[t][6.9cm][c]{\linewidth}\LARGE\centering
\begin{minipage}{0.95\linewidth}\raggedright\linespread{1.25}\selectfont Quelle est une équation cartésienne de la droite passant par $A(0\,;\,1)$ et de vecteur normal $\vec n\begin{pmatrix}-2\\5\end{pmatrix}$~?\end{minipage}
\end{minipage}

\vfill
\begin{tcolorbox}[enhanced,nobeforeafter,width=7.1cm,height=2.1cm,valign=center,colback=cpsc!6,colframe=cpsc,arc=2.5mm,boxrule=0.9pt,left=3mm]{\Large\bfseries\color{cpsc}a)}\quad\begin{minipage}[c]{5.3cm}\raggedright\Large $-2x+5y+5=0$\end{minipage}\end{tcolorbox}\hfill\begin{tcolorbox}[enhanced,nobeforeafter,width=7.1cm,height=2.1cm,valign=center,colback=cpsc!6,colframe=cpsc,arc=2.5mm,boxrule=0.9pt,left=3mm]{\Large\bfseries\color{cpsc}b)}\quad\begin{minipage}[c]{5.3cm}\raggedright\Large $-2x+5y-5=0$\end{minipage}\end{tcolorbox}\par\vspace{2.5mm}
\begin{tcolorbox}[enhanced,nobeforeafter,width=7.1cm,height=2.1cm,valign=center,colback=cpsc!6,colframe=cpsc,arc=2.5mm,boxrule=0.9pt,left=3mm]{\Large\bfseries\color{cpsc}c)}\quad\begin{minipage}[c]{5.3cm}\raggedright\Large $y-5=0$\end{minipage}\end{tcolorbox}\hfill\begin{tcolorbox}[enhanced,nobeforeafter,width=7.1cm,height=2.1cm,valign=center,colback=cpsc!6,colframe=cpsc,arc=2.5mm,boxrule=0.9pt,left=3mm]{\Large\bfseries\color{cpsc}d)}\quad\begin{minipage}[c]{5.3cm}\raggedright\Large $-5x-2y+2=0$\end{minipage}\end{tcolorbox}
\end{tcolorbox}
\newpage
\begin{tcolorbox}[enhanced,width=16cm,height=14cm,colback=white,colframe=cpsc,boxrule=2pt,arc=0mm,
left=6mm,right=6mm,top=0mm,bottom=4mm,sharp corners,
overlay={\fill[cpsc] (frame.north west) rectangle ([yshift=-1.25cm]frame.north east);
\node[anchor=west,text=white,font=\sffamily\bfseries\LARGE] at ([xshift=5mm,yshift=-0.62cm]frame.north west) {Question 77 B};
\node[anchor=east,text=white,font=\sffamily\itshape\large] at ([xshift=-5mm,yshift=-0.62cm]frame.north east) {Produit scalaire et géométrie repérée};}]
\vspace*{1.75cm}
\begin{minipage}[t][6.9cm][c]{\linewidth}\LARGE\centering
\begin{minipage}{0.95\linewidth}\raggedright\linespread{1.25}\selectfont Quels sont le centre et le rayon du cercle d'équation $x^2+(y-5)^2=16$~?\end{minipage}
\end{minipage}

\vfill
\begin{tcolorbox}[enhanced,nobeforeafter,width=7.1cm,height=2.1cm,valign=center,colback=cpsc!6,colframe=cpsc,arc=2.5mm,boxrule=0.9pt,left=3mm]{\Large\bfseries\color{cpsc}a)}\quad\begin{minipage}[c]{5.3cm}\raggedright\large $\Omega(0\,;\,-5)\text{, rayon }4$\end{minipage}\end{tcolorbox}\hfill\begin{tcolorbox}[enhanced,nobeforeafter,width=7.1cm,height=2.1cm,valign=center,colback=cpsc!6,colframe=cpsc,arc=2.5mm,boxrule=0.9pt,left=3mm]{\Large\bfseries\color{cpsc}b)}\quad\begin{minipage}[c]{5.3cm}\raggedright\large $\Omega(5\,;\,0)\text{, rayon }4$\end{minipage}\end{tcolorbox}\par\vspace{2.5mm}
\begin{tcolorbox}[enhanced,nobeforeafter,width=7.1cm,height=2.1cm,valign=center,colback=cpsc!6,colframe=cpsc,arc=2.5mm,boxrule=0.9pt,left=3mm]{\Large\bfseries\color{cpsc}c)}\quad\begin{minipage}[c]{5.3cm}\raggedright\large $\Omega(0\,;\,5)\text{, rayon }4$\end{minipage}\end{tcolorbox}\hfill\begin{tcolorbox}[enhanced,nobeforeafter,width=7.1cm,height=2.1cm,valign=center,colback=cpsc!6,colframe=cpsc,arc=2.5mm,boxrule=0.9pt,left=3mm]{\Large\bfseries\color{cpsc}d)}\quad\begin{minipage}[c]{5.3cm}\raggedright\large $\Omega(0\,;\,5)\text{, rayon }16$\end{minipage}\end{tcolorbox}
\end{tcolorbox}
\newpage
\begin{tcolorbox}[enhanced,width=16cm,height=14cm,colback=white,colframe=cpsc,boxrule=2pt,arc=0mm,
left=6mm,right=6mm,top=0mm,bottom=4mm,sharp corners,
overlay={\fill[cpsc] (frame.north west) rectangle ([yshift=-1.25cm]frame.north east);
\node[anchor=west,text=white,font=\sffamily\bfseries\LARGE] at ([xshift=5mm,yshift=-0.62cm]frame.north west) {Question 78 B};
\node[anchor=east,text=white,font=\sffamily\itshape\large] at ([xshift=-5mm,yshift=-0.62cm]frame.north east) {Produit scalaire et géométrie repérée};}]
\vspace*{1.75cm}
\begin{minipage}[t][6.9cm][c]{\linewidth}\LARGE\centering
\begin{minipage}{0.95\linewidth}\raggedright\linespread{1.25}\selectfont Quels sont le centre et le rayon du cercle d'équation $(x-1)^2+(y-1)^2=2$~?\end{minipage}
\end{minipage}

\vfill
\begin{tcolorbox}[enhanced,nobeforeafter,width=7.1cm,height=2.1cm,valign=center,colback=cpsc!6,colframe=cpsc,arc=2.5mm,boxrule=0.9pt,left=3mm]{\Large\bfseries\color{cpsc}a)}\quad\begin{minipage}[c]{5.3cm}\raggedright\large $\Omega(1\,;\,-1)\text{, rayon }\sqrt{2}$\end{minipage}\end{tcolorbox}\hfill\begin{tcolorbox}[enhanced,nobeforeafter,width=7.1cm,height=2.1cm,valign=center,colback=cpsc!6,colframe=cpsc,arc=2.5mm,boxrule=0.9pt,left=3mm]{\Large\bfseries\color{cpsc}b)}\quad\begin{minipage}[c]{5.3cm}\raggedright\large $\Omega(1\,;\,1)\text{, rayon }\sqrt{2}$\end{minipage}\end{tcolorbox}\par\vspace{2.5mm}
\begin{tcolorbox}[enhanced,nobeforeafter,width=7.1cm,height=2.1cm,valign=center,colback=cpsc!6,colframe=cpsc,arc=2.5mm,boxrule=0.9pt,left=3mm]{\Large\bfseries\color{cpsc}c)}\quad\begin{minipage}[c]{5.3cm}\raggedright\large $\Omega(1\,;\,1)\text{, rayon }2$\end{minipage}\end{tcolorbox}\hfill\begin{tcolorbox}[enhanced,nobeforeafter,width=7.1cm,height=2.1cm,valign=center,colback=cpsc!6,colframe=cpsc,arc=2.5mm,boxrule=0.9pt,left=3mm]{\Large\bfseries\color{cpsc}d)}\quad\begin{minipage}[c]{5.3cm}\raggedright\large $\Omega(-1\,;\,-1)\text{, rayon }\sqrt{2}$\end{minipage}\end{tcolorbox}
\end{tcolorbox}
\newpage
\begin{tcolorbox}[enhanced,width=16cm,height=14cm,colback=white,colframe=cpsc,boxrule=2pt,arc=0mm,
left=6mm,right=6mm,top=0mm,bottom=4mm,sharp corners,
overlay={\fill[cpsc] (frame.north west) rectangle ([yshift=-1.25cm]frame.north east);
\node[anchor=west,text=white,font=\sffamily\bfseries\LARGE] at ([xshift=5mm,yshift=-0.62cm]frame.north west) {Question 79 B};
\node[anchor=east,text=white,font=\sffamily\itshape\large] at ([xshift=-5mm,yshift=-0.62cm]frame.north east) {Produit scalaire et géométrie repérée};}]
\vspace*{1.75cm}
\begin{minipage}[t][6.9cm][c]{\linewidth}\LARGE\centering
\begin{minipage}{0.95\linewidth}\raggedright\linespread{1.25}\selectfont On a $\|\vec u\|=1$, $\|\vec v\|=1$ et $\vec u\cdot\vec v=0$. Que vaut $\|\vec u+\vec v\|^2$~?\end{minipage}
\end{minipage}

\vfill
\begin{tcolorbox}[enhanced,nobeforeafter,width=7.1cm,height=2.1cm,valign=center,colback=cpsc!6,colframe=cpsc,arc=2.5mm,boxrule=0.9pt,left=3mm]{\Large\bfseries\color{cpsc}a)}\quad\begin{minipage}[c]{5.3cm}\raggedright\Large $2$\end{minipage}\end{tcolorbox}\hfill\begin{tcolorbox}[enhanced,nobeforeafter,width=7.1cm,height=2.1cm,valign=center,colback=cpsc!6,colframe=cpsc,arc=2.5mm,boxrule=0.9pt,left=3mm]{\Large\bfseries\color{cpsc}b)}\quad\begin{minipage}[c]{5.3cm}\raggedright\Large $1$\end{minipage}\end{tcolorbox}\par\vspace{2.5mm}
\begin{tcolorbox}[enhanced,nobeforeafter,width=7.1cm,height=2.1cm,valign=center,colback=cpsc!6,colframe=cpsc,arc=2.5mm,boxrule=0.9pt,left=3mm]{\Large\bfseries\color{cpsc}c)}\quad\begin{minipage}[c]{5.3cm}\raggedright\Large $3$\end{minipage}\end{tcolorbox}\hfill\begin{tcolorbox}[enhanced,nobeforeafter,width=7.1cm,height=2.1cm,valign=center,colback=cpsc!6,colframe=cpsc,arc=2.5mm,boxrule=0.9pt,left=3mm]{\Large\bfseries\color{cpsc}d)}\quad\begin{minipage}[c]{5.3cm}\raggedright\Large $4$\end{minipage}\end{tcolorbox}
\end{tcolorbox}
\newpage
\begin{tcolorbox}[enhanced,width=16cm,height=14cm,colback=white,colframe=cpsc,boxrule=2pt,arc=0mm,
left=6mm,right=6mm,top=0mm,bottom=4mm,sharp corners,
overlay={\fill[cpsc] (frame.north west) rectangle ([yshift=-1.25cm]frame.north east);
\node[anchor=west,text=white,font=\sffamily\bfseries\LARGE] at ([xshift=5mm,yshift=-0.62cm]frame.north west) {Question 80 B};
\node[anchor=east,text=white,font=\sffamily\itshape\large] at ([xshift=-5mm,yshift=-0.62cm]frame.north east) {Produit scalaire et géométrie repérée};}]
\vspace*{1.75cm}
\begin{minipage}[t][6.9cm][c]{\linewidth}\LARGE\centering
\begin{minipage}{0.95\linewidth}\raggedright\linespread{1.25}\selectfont On a $\|\vec u\|=6$, $\|\vec v\|=2$ et $\vec u\cdot\vec v=-5$. Que vaut $\|\vec u+\vec v\|^2$~?\end{minipage}
\end{minipage}

\vfill
\begin{tcolorbox}[enhanced,nobeforeafter,width=7.1cm,height=2.1cm,valign=center,colback=cpsc!6,colframe=cpsc,arc=2.5mm,boxrule=0.9pt,left=3mm]{\Large\bfseries\color{cpsc}a)}\quad\begin{minipage}[c]{5.3cm}\raggedright\Large $64$\end{minipage}\end{tcolorbox}\hfill\begin{tcolorbox}[enhanced,nobeforeafter,width=7.1cm,height=2.1cm,valign=center,colback=cpsc!6,colframe=cpsc,arc=2.5mm,boxrule=0.9pt,left=3mm]{\Large\bfseries\color{cpsc}b)}\quad\begin{minipage}[c]{5.3cm}\raggedright\Large $35$\end{minipage}\end{tcolorbox}\par\vspace{2.5mm}
\begin{tcolorbox}[enhanced,nobeforeafter,width=7.1cm,height=2.1cm,valign=center,colback=cpsc!6,colframe=cpsc,arc=2.5mm,boxrule=0.9pt,left=3mm]{\Large\bfseries\color{cpsc}c)}\quad\begin{minipage}[c]{5.3cm}\raggedright\Large $30$\end{minipage}\end{tcolorbox}\hfill\begin{tcolorbox}[enhanced,nobeforeafter,width=7.1cm,height=2.1cm,valign=center,colback=cpsc!6,colframe=cpsc,arc=2.5mm,boxrule=0.9pt,left=3mm]{\Large\bfseries\color{cpsc}d)}\quad\begin{minipage}[c]{5.3cm}\raggedright\Large $40$\end{minipage}\end{tcolorbox}
\end{tcolorbox}
\newpage
\begin{tcolorbox}[enhanced,width=16cm,height=14cm,colback=white,colframe=ccvp,boxrule=2pt,arc=0mm,
left=6mm,right=6mm,top=0mm,bottom=4mm,sharp corners,
overlay={\fill[ccvp] (frame.north west) rectangle ([yshift=-1.25cm]frame.north east);
\node[anchor=west,text=white,font=\sffamily\bfseries\LARGE] at ([xshift=5mm,yshift=-0.62cm]frame.north west) {Question 81 B};
\node[anchor=east,text=white,font=\sffamily\itshape\large] at ([xshift=-5mm,yshift=-0.62cm]frame.north east) {Probabilités conditionnelles et variables aléatoires};}]
\vspace*{1.75cm}
\begin{minipage}[t][6.9cm][c]{\linewidth}\LARGE\centering
\begin{minipage}{0.95\linewidth}\raggedright\linespread{1.25}\selectfont On a $P(A)=0{,}8$ et $P(A\cap B)=0{,}6$. Que vaut $P_A(B)$~?\end{minipage}
\end{minipage}

\vfill
\begin{tcolorbox}[enhanced,nobeforeafter,width=7.1cm,height=2.1cm,valign=center,colback=ccvp!6,colframe=ccvp,arc=2.5mm,boxrule=0.9pt,left=3mm]{\Large\bfseries\color{ccvp}a)}\quad\begin{minipage}[c]{5.3cm}\raggedright\Large $0{,}48$\end{minipage}\end{tcolorbox}\hfill\begin{tcolorbox}[enhanced,nobeforeafter,width=7.1cm,height=2.1cm,valign=center,colback=ccvp!6,colframe=ccvp,arc=2.5mm,boxrule=0.9pt,left=3mm]{\Large\bfseries\color{ccvp}b)}\quad\begin{minipage}[c]{5.3cm}\raggedright\Large $0{,}75$\end{minipage}\end{tcolorbox}\par\vspace{2.5mm}
\begin{tcolorbox}[enhanced,nobeforeafter,width=7.1cm,height=2.1cm,valign=center,colback=ccvp!6,colframe=ccvp,arc=2.5mm,boxrule=0.9pt,left=3mm]{\Large\bfseries\color{ccvp}c)}\quad\begin{minipage}[c]{5.3cm}\raggedright\Large $0{,}2$\end{minipage}\end{tcolorbox}\hfill\begin{tcolorbox}[enhanced,nobeforeafter,width=7.1cm,height=2.1cm,valign=center,colback=ccvp!6,colframe=ccvp,arc=2.5mm,boxrule=0.9pt,left=3mm]{\Large\bfseries\color{ccvp}d)}\quad\begin{minipage}[c]{5.3cm}\raggedright\Large $0{,}4$\end{minipage}\end{tcolorbox}
\end{tcolorbox}
\newpage
\begin{tcolorbox}[enhanced,width=16cm,height=14cm,colback=white,colframe=ccvp,boxrule=2pt,arc=0mm,
left=6mm,right=6mm,top=0mm,bottom=4mm,sharp corners,
overlay={\fill[ccvp] (frame.north west) rectangle ([yshift=-1.25cm]frame.north east);
\node[anchor=west,text=white,font=\sffamily\bfseries\LARGE] at ([xshift=5mm,yshift=-0.62cm]frame.north west) {Question 82 B};
\node[anchor=east,text=white,font=\sffamily\itshape\large] at ([xshift=-5mm,yshift=-0.62cm]frame.north east) {Probabilités conditionnelles et variables aléatoires};}]
\vspace*{1.75cm}
\begin{minipage}[t][6.9cm][c]{\linewidth}\LARGE\centering
\begin{minipage}{0.95\linewidth}\raggedright\linespread{1.25}\selectfont On a $P(A)=0{,}25$ et $P(A\cap B)=0{,}05$. Que vaut $P_A(B)$~?\end{minipage}
\end{minipage}

\vfill
\begin{tcolorbox}[enhanced,nobeforeafter,width=7.1cm,height=2.1cm,valign=center,colback=ccvp!6,colframe=ccvp,arc=2.5mm,boxrule=0.9pt,left=3mm]{\Large\bfseries\color{ccvp}a)}\quad\begin{minipage}[c]{5.3cm}\raggedright\Large $0{,}3$\end{minipage}\end{tcolorbox}\hfill\begin{tcolorbox}[enhanced,nobeforeafter,width=7.1cm,height=2.1cm,valign=center,colback=ccvp!6,colframe=ccvp,arc=2.5mm,boxrule=0.9pt,left=3mm]{\Large\bfseries\color{ccvp}b)}\quad\begin{minipage}[c]{5.3cm}\raggedright\Large $0{,}2$\end{minipage}\end{tcolorbox}\par\vspace{2.5mm}
\begin{tcolorbox}[enhanced,nobeforeafter,width=7.1cm,height=2.1cm,valign=center,colback=ccvp!6,colframe=ccvp,arc=2.5mm,boxrule=0.9pt,left=3mm]{\Large\bfseries\color{ccvp}c)}\quad\begin{minipage}[c]{5.3cm}\raggedright\Large $0{,}05$\end{minipage}\end{tcolorbox}\hfill\begin{tcolorbox}[enhanced,nobeforeafter,width=7.1cm,height=2.1cm,valign=center,colback=ccvp!6,colframe=ccvp,arc=2.5mm,boxrule=0.9pt,left=3mm]{\Large\bfseries\color{ccvp}d)}\quad\begin{minipage}[c]{5.3cm}\raggedright\Large $0{,}0125$\end{minipage}\end{tcolorbox}
\end{tcolorbox}
\newpage
\begin{tcolorbox}[enhanced,width=16cm,height=14cm,colback=white,colframe=ccvp,boxrule=2pt,arc=0mm,
left=6mm,right=6mm,top=0mm,bottom=4mm,sharp corners,
overlay={\fill[ccvp] (frame.north west) rectangle ([yshift=-1.25cm]frame.north east);
\node[anchor=west,text=white,font=\sffamily\bfseries\LARGE] at ([xshift=5mm,yshift=-0.62cm]frame.north west) {Question 83 B};
\node[anchor=east,text=white,font=\sffamily\itshape\large] at ([xshift=-5mm,yshift=-0.62cm]frame.north east) {Probabilités conditionnelles et variables aléatoires};}]
\vspace*{1.75cm}
\begin{minipage}[t][6.9cm][c]{\linewidth}\LARGE\centering
\begin{minipage}{0.95\linewidth}\raggedright\linespread{1.25}\selectfont On a $P(A)=0{,}8$ et $P_A(B)=0{,}25$. Que vaut $P(A\cap B)$~?\end{minipage}
\end{minipage}

\vfill
\begin{tcolorbox}[enhanced,nobeforeafter,width=7.1cm,height=2.1cm,valign=center,colback=ccvp!6,colframe=ccvp,arc=2.5mm,boxrule=0.9pt,left=3mm]{\Large\bfseries\color{ccvp}a)}\quad\begin{minipage}[c]{5.3cm}\raggedright\Large $0{,}2$\end{minipage}\end{tcolorbox}\hfill\begin{tcolorbox}[enhanced,nobeforeafter,width=7.1cm,height=2.1cm,valign=center,colback=ccvp!6,colframe=ccvp,arc=2.5mm,boxrule=0.9pt,left=3mm]{\Large\bfseries\color{ccvp}b)}\quad\begin{minipage}[c]{5.3cm}\raggedright\Large $0{,}05$\end{minipage}\end{tcolorbox}\par\vspace{2.5mm}
\begin{tcolorbox}[enhanced,nobeforeafter,width=7.1cm,height=2.1cm,valign=center,colback=ccvp!6,colframe=ccvp,arc=2.5mm,boxrule=0.9pt,left=3mm]{\Large\bfseries\color{ccvp}c)}\quad\begin{minipage}[c]{5.3cm}\raggedright\Large $0{,}3125$\end{minipage}\end{tcolorbox}\hfill\begin{tcolorbox}[enhanced,nobeforeafter,width=7.1cm,height=2.1cm,valign=center,colback=ccvp!6,colframe=ccvp,arc=2.5mm,boxrule=0.9pt,left=3mm]{\Large\bfseries\color{ccvp}d)}\quad\begin{minipage}[c]{5.3cm}\raggedright\Large $0{,}6$\end{minipage}\end{tcolorbox}
\end{tcolorbox}
\newpage
\begin{tcolorbox}[enhanced,width=16cm,height=14cm,colback=white,colframe=ccvp,boxrule=2pt,arc=0mm,
left=6mm,right=6mm,top=0mm,bottom=4mm,sharp corners,
overlay={\fill[ccvp] (frame.north west) rectangle ([yshift=-1.25cm]frame.north east);
\node[anchor=west,text=white,font=\sffamily\bfseries\LARGE] at ([xshift=5mm,yshift=-0.62cm]frame.north west) {Question 84 B};
\node[anchor=east,text=white,font=\sffamily\itshape\large] at ([xshift=-5mm,yshift=-0.62cm]frame.north east) {Probabilités conditionnelles et variables aléatoires};}]
\vspace*{1.75cm}
\begin{minipage}[t][6.9cm][c]{\linewidth}\LARGE\centering
\begin{minipage}{0.95\linewidth}\raggedright\linespread{1.25}\selectfont On a $P(A)=0{,}45$ et $P_A(B)=0{,}2$. Que vaut $P(A\cap B)$~?\end{minipage}
\end{minipage}

\vfill
\begin{tcolorbox}[enhanced,nobeforeafter,width=7.1cm,height=2.1cm,valign=center,colback=ccvp!6,colframe=ccvp,arc=2.5mm,boxrule=0.9pt,left=3mm]{\Large\bfseries\color{ccvp}a)}\quad\begin{minipage}[c]{5.3cm}\raggedright\Large $\dfrac{4}{9}$\end{minipage}\end{tcolorbox}\hfill\begin{tcolorbox}[enhanced,nobeforeafter,width=7.1cm,height=2.1cm,valign=center,colback=ccvp!6,colframe=ccvp,arc=2.5mm,boxrule=0.9pt,left=3mm]{\Large\bfseries\color{ccvp}b)}\quad\begin{minipage}[c]{5.3cm}\raggedright\Large $0{,}36$\end{minipage}\end{tcolorbox}\par\vspace{2.5mm}
\begin{tcolorbox}[enhanced,nobeforeafter,width=7.1cm,height=2.1cm,valign=center,colback=ccvp!6,colframe=ccvp,arc=2.5mm,boxrule=0.9pt,left=3mm]{\Large\bfseries\color{ccvp}c)}\quad\begin{minipage}[c]{5.3cm}\raggedright\Large $0{,}09$\end{minipage}\end{tcolorbox}\hfill\begin{tcolorbox}[enhanced,nobeforeafter,width=7.1cm,height=2.1cm,valign=center,colback=ccvp!6,colframe=ccvp,arc=2.5mm,boxrule=0.9pt,left=3mm]{\Large\bfseries\color{ccvp}d)}\quad\begin{minipage}[c]{5.3cm}\raggedright\Large $0{,}65$\end{minipage}\end{tcolorbox}
\end{tcolorbox}
\newpage
\begin{tcolorbox}[enhanced,width=16cm,height=14cm,colback=white,colframe=ccvp,boxrule=2pt,arc=0mm,
left=6mm,right=6mm,top=0mm,bottom=4mm,sharp corners,
overlay={\fill[ccvp] (frame.north west) rectangle ([yshift=-1.25cm]frame.north east);
\node[anchor=west,text=white,font=\sffamily\bfseries\LARGE] at ([xshift=5mm,yshift=-0.62cm]frame.north west) {Question 85 B};
\node[anchor=east,text=white,font=\sffamily\itshape\large] at ([xshift=-5mm,yshift=-0.62cm]frame.north east) {Probabilités conditionnelles et variables aléatoires};}]
\vspace*{1.75cm}
\begin{minipage}[t][6.9cm][c]{\linewidth}\LARGE\centering
\begin{minipage}{0.95\linewidth}\raggedright\linespread{1.25}\selectfont On a $P(A)=0{,}5$, $P_A(B)=0{,}8$ et $P_{\overline{A}}(B)=0{,}4$. Que vaut $P(B)$~?\end{minipage}
\end{minipage}

\vfill
\begin{tcolorbox}[enhanced,nobeforeafter,width=7.1cm,height=2.1cm,valign=center,colback=ccvp!6,colframe=ccvp,arc=2.5mm,boxrule=0.9pt,left=3mm]{\Large\bfseries\color{ccvp}a)}\quad\begin{minipage}[c]{5.3cm}\raggedright\Large $0{,}7$\end{minipage}\end{tcolorbox}\hfill\begin{tcolorbox}[enhanced,nobeforeafter,width=7.1cm,height=2.1cm,valign=center,colback=ccvp!6,colframe=ccvp,arc=2.5mm,boxrule=0.9pt,left=3mm]{\Large\bfseries\color{ccvp}b)}\quad\begin{minipage}[c]{5.3cm}\raggedright\Large $0{,}4$\end{minipage}\end{tcolorbox}\par\vspace{2.5mm}
\begin{tcolorbox}[enhanced,nobeforeafter,width=7.1cm,height=2.1cm,valign=center,colback=ccvp!6,colframe=ccvp,arc=2.5mm,boxrule=0.9pt,left=3mm]{\Large\bfseries\color{ccvp}c)}\quad\begin{minipage}[c]{5.3cm}\raggedright\Large $0{,}32$\end{minipage}\end{tcolorbox}\hfill\begin{tcolorbox}[enhanced,nobeforeafter,width=7.1cm,height=2.1cm,valign=center,colback=ccvp!6,colframe=ccvp,arc=2.5mm,boxrule=0.9pt,left=3mm]{\Large\bfseries\color{ccvp}d)}\quad\begin{minipage}[c]{5.3cm}\raggedright\Large $0{,}6$\end{minipage}\end{tcolorbox}
\end{tcolorbox}
\newpage
\begin{tcolorbox}[enhanced,width=16cm,height=14cm,colback=white,colframe=ccvp,boxrule=2pt,arc=0mm,
left=6mm,right=6mm,top=0mm,bottom=4mm,sharp corners,
overlay={\fill[ccvp] (frame.north west) rectangle ([yshift=-1.25cm]frame.north east);
\node[anchor=west,text=white,font=\sffamily\bfseries\LARGE] at ([xshift=5mm,yshift=-0.62cm]frame.north west) {Question 86 B};
\node[anchor=east,text=white,font=\sffamily\itshape\large] at ([xshift=-5mm,yshift=-0.62cm]frame.north east) {Probabilités conditionnelles et variables aléatoires};}]
\vspace*{1.75cm}
\begin{minipage}[t][6.9cm][c]{\linewidth}\LARGE\centering
\begin{minipage}{0.95\linewidth}\raggedright\linespread{1.25}\selectfont On a $P(A)=0{,}2$, $P_A(B)=0{,}9$ et $P_{\overline{A}}(B)=0{,}3$. Que vaut $P(B)$~?\end{minipage}
\end{minipage}

\vfill
\begin{tcolorbox}[enhanced,nobeforeafter,width=7.1cm,height=2.1cm,valign=center,colback=ccvp!6,colframe=ccvp,arc=2.5mm,boxrule=0.9pt,left=3mm]{\Large\bfseries\color{ccvp}a)}\quad\begin{minipage}[c]{5.3cm}\raggedright\Large $0{,}18$\end{minipage}\end{tcolorbox}\hfill\begin{tcolorbox}[enhanced,nobeforeafter,width=7.1cm,height=2.1cm,valign=center,colback=ccvp!6,colframe=ccvp,arc=2.5mm,boxrule=0.9pt,left=3mm]{\Large\bfseries\color{ccvp}b)}\quad\begin{minipage}[c]{5.3cm}\raggedright\Large $0{,}27$\end{minipage}\end{tcolorbox}\par\vspace{2.5mm}
\begin{tcolorbox}[enhanced,nobeforeafter,width=7.1cm,height=2.1cm,valign=center,colback=ccvp!6,colframe=ccvp,arc=2.5mm,boxrule=0.9pt,left=3mm]{\Large\bfseries\color{ccvp}c)}\quad\begin{minipage}[c]{5.3cm}\raggedright\Large $0{,}42$\end{minipage}\end{tcolorbox}\hfill\begin{tcolorbox}[enhanced,nobeforeafter,width=7.1cm,height=2.1cm,valign=center,colback=ccvp!6,colframe=ccvp,arc=2.5mm,boxrule=0.9pt,left=3mm]{\Large\bfseries\color{ccvp}d)}\quad\begin{minipage}[c]{5.3cm}\raggedright\Large $0{,}24$\end{minipage}\end{tcolorbox}
\end{tcolorbox}
\newpage
\begin{tcolorbox}[enhanced,width=16cm,height=14cm,colback=white,colframe=ccvp,boxrule=2pt,arc=0mm,
left=6mm,right=6mm,top=0mm,bottom=4mm,sharp corners,
overlay={\fill[ccvp] (frame.north west) rectangle ([yshift=-1.25cm]frame.north east);
\node[anchor=west,text=white,font=\sffamily\bfseries\LARGE] at ([xshift=5mm,yshift=-0.62cm]frame.north west) {Question 87 B};
\node[anchor=east,text=white,font=\sffamily\itshape\large] at ([xshift=-5mm,yshift=-0.62cm]frame.north east) {Probabilités conditionnelles et variables aléatoires};}]
\vspace*{1.75cm}
\begin{minipage}[t][6.9cm][c]{\linewidth}\LARGE\centering
\begin{minipage}{0.95\linewidth}\raggedright\linespread{1.25}\selectfont On a $P(A)=0{,}25$, $P(B)=0{,}8$ et $P(A\cap B)=0{,}2$. Les événements $A$ et $B$ sont-ils indépendants~?\end{minipage}
\end{minipage}

\vfill
\begin{tcolorbox}[enhanced,nobeforeafter,width=7.1cm,height=2.1cm,valign=center,colback=ccvp!6,colframe=ccvp,arc=2.5mm,boxrule=0.9pt,left=3mm]{\Large\bfseries\color{ccvp}a)}\quad\begin{minipage}[c]{5.3cm}\raggedright\Large on ne peut pas savoir\end{minipage}\end{tcolorbox}\hfill\begin{tcolorbox}[enhanced,nobeforeafter,width=7.1cm,height=2.1cm,valign=center,colback=ccvp!6,colframe=ccvp,arc=2.5mm,boxrule=0.9pt,left=3mm]{\Large\bfseries\color{ccvp}b)}\quad\begin{minipage}[c]{5.3cm}\raggedright\Large non\end{minipage}\end{tcolorbox}\par\vspace{2.5mm}
\begin{tcolorbox}[enhanced,nobeforeafter,width=7.1cm,height=2.1cm,valign=center,colback=ccvp!6,colframe=ccvp,arc=2.5mm,boxrule=0.9pt,left=3mm]{\Large\bfseries\color{ccvp}c)}\quad\begin{minipage}[c]{5.3cm}\raggedright\Large oui\end{minipage}\end{tcolorbox}\hfill\begin{tcolorbox}[enhanced,nobeforeafter,width=7.1cm,height=2.1cm,valign=center,colback=ccvp!6,colframe=ccvp,arc=2.5mm,boxrule=0.9pt,left=3mm]{\Large\bfseries\color{ccvp}d)}\quad\begin{minipage}[c]{5.3cm}\raggedright\large seulement s'ils sont incompatibles\end{minipage}\end{tcolorbox}
\end{tcolorbox}
\newpage
\begin{tcolorbox}[enhanced,width=16cm,height=14cm,colback=white,colframe=ccvp,boxrule=2pt,arc=0mm,
left=6mm,right=6mm,top=0mm,bottom=4mm,sharp corners,
overlay={\fill[ccvp] (frame.north west) rectangle ([yshift=-1.25cm]frame.north east);
\node[anchor=west,text=white,font=\sffamily\bfseries\LARGE] at ([xshift=5mm,yshift=-0.62cm]frame.north west) {Question 88 B};
\node[anchor=east,text=white,font=\sffamily\itshape\large] at ([xshift=-5mm,yshift=-0.62cm]frame.north east) {Probabilités conditionnelles et variables aléatoires};}]
\vspace*{1.75cm}
\begin{minipage}[t][6.9cm][c]{\linewidth}\LARGE\centering
\begin{minipage}{0.95\linewidth}\raggedright\linespread{1.25}\selectfont On a $P(A)=0{,}7$, $P(B)=0{,}5$ et $P(A\cap B)=0{,}3$. Les événements $A$ et $B$ sont-ils indépendants~?\end{minipage}
\end{minipage}

\vfill
\begin{tcolorbox}[enhanced,nobeforeafter,width=7.1cm,height=2.1cm,valign=center,colback=ccvp!6,colframe=ccvp,arc=2.5mm,boxrule=0.9pt,left=3mm]{\Large\bfseries\color{ccvp}a)}\quad\begin{minipage}[c]{5.3cm}\raggedright\Large oui\end{minipage}\end{tcolorbox}\hfill\begin{tcolorbox}[enhanced,nobeforeafter,width=7.1cm,height=2.1cm,valign=center,colback=ccvp!6,colframe=ccvp,arc=2.5mm,boxrule=0.9pt,left=3mm]{\Large\bfseries\color{ccvp}b)}\quad\begin{minipage}[c]{5.3cm}\raggedright\Large on ne peut pas savoir\end{minipage}\end{tcolorbox}\par\vspace{2.5mm}
\begin{tcolorbox}[enhanced,nobeforeafter,width=7.1cm,height=2.1cm,valign=center,colback=ccvp!6,colframe=ccvp,arc=2.5mm,boxrule=0.9pt,left=3mm]{\Large\bfseries\color{ccvp}c)}\quad\begin{minipage}[c]{5.3cm}\raggedright\Large non\end{minipage}\end{tcolorbox}\hfill\begin{tcolorbox}[enhanced,nobeforeafter,width=7.1cm,height=2.1cm,valign=center,colback=ccvp!6,colframe=ccvp,arc=2.5mm,boxrule=0.9pt,left=3mm]{\Large\bfseries\color{ccvp}d)}\quad\begin{minipage}[c]{5.3cm}\raggedright\large seulement s'ils sont incompatibles\end{minipage}\end{tcolorbox}
\end{tcolorbox}
\newpage
\begin{tcolorbox}[enhanced,width=16cm,height=14cm,colback=white,colframe=ccvp,boxrule=2pt,arc=0mm,
left=6mm,right=6mm,top=0mm,bottom=4mm,sharp corners,
overlay={\fill[ccvp] (frame.north west) rectangle ([yshift=-1.25cm]frame.north east);
\node[anchor=west,text=white,font=\sffamily\bfseries\LARGE] at ([xshift=5mm,yshift=-0.62cm]frame.north west) {Question 89 B};
\node[anchor=east,text=white,font=\sffamily\itshape\large] at ([xshift=-5mm,yshift=-0.62cm]frame.north east) {Probabilités conditionnelles et variables aléatoires};}]
\vspace*{1.75cm}
\begin{minipage}[t][6.9cm][c]{\linewidth}\Large\centering
\begin{minipage}{0.95\linewidth}\raggedright\linespread{1.25}\selectfont Voici la loi de probabilité d'une variable aléatoire $X$. Que vaut $E(X)$~? $\begin{array}{|l|c|c|}\hline x_i & -1 & 3\\\hline P(X=x_i) & 0{,}4 & 0{,}6\\\hline\end{array}$\end{minipage}
\end{minipage}

\vfill
\begin{tcolorbox}[enhanced,nobeforeafter,width=7.1cm,height=2.1cm,valign=center,colback=ccvp!6,colframe=ccvp,arc=2.5mm,boxrule=0.9pt,left=3mm]{\Large\bfseries\color{ccvp}a)}\quad\begin{minipage}[c]{5.3cm}\raggedright\Large $1{,}4$\end{minipage}\end{tcolorbox}\hfill\begin{tcolorbox}[enhanced,nobeforeafter,width=7.1cm,height=2.1cm,valign=center,colback=ccvp!6,colframe=ccvp,arc=2.5mm,boxrule=0.9pt,left=3mm]{\Large\bfseries\color{ccvp}b)}\quad\begin{minipage}[c]{5.3cm}\raggedright\Large $1$\end{minipage}\end{tcolorbox}\par\vspace{2.5mm}
\begin{tcolorbox}[enhanced,nobeforeafter,width=7.1cm,height=2.1cm,valign=center,colback=ccvp!6,colframe=ccvp,arc=2.5mm,boxrule=0.9pt,left=3mm]{\Large\bfseries\color{ccvp}c)}\quad\begin{minipage}[c]{5.3cm}\raggedright\Large $-1{,}4$\end{minipage}\end{tcolorbox}\hfill\begin{tcolorbox}[enhanced,nobeforeafter,width=7.1cm,height=2.1cm,valign=center,colback=ccvp!6,colframe=ccvp,arc=2.5mm,boxrule=0.9pt,left=3mm]{\Large\bfseries\color{ccvp}d)}\quad\begin{minipage}[c]{5.3cm}\raggedright\Large $2{,}4$\end{minipage}\end{tcolorbox}
\end{tcolorbox}
\newpage
\begin{tcolorbox}[enhanced,width=16cm,height=14cm,colback=white,colframe=ccvp,boxrule=2pt,arc=0mm,
left=6mm,right=6mm,top=0mm,bottom=4mm,sharp corners,
overlay={\fill[ccvp] (frame.north west) rectangle ([yshift=-1.25cm]frame.north east);
\node[anchor=west,text=white,font=\sffamily\bfseries\LARGE] at ([xshift=5mm,yshift=-0.62cm]frame.north west) {Question 90 B};
\node[anchor=east,text=white,font=\sffamily\itshape\large] at ([xshift=-5mm,yshift=-0.62cm]frame.north east) {Probabilités conditionnelles et variables aléatoires};}]
\vspace*{1.75cm}
\begin{minipage}[t][6.9cm][c]{\linewidth}\Large\centering
\begin{minipage}{0.95\linewidth}\raggedright\linespread{1.25}\selectfont Voici la loi de probabilité d'une variable aléatoire $X$. Que vaut $E(X)$~? $\begin{array}{|l|c|c|c|}\hline x_i & 1 & 2 & 3\\\hline P(X=x_i) & 0{,}2 & 0{,}3 & 0{,}5\\\hline\end{array}$\end{minipage}
\end{minipage}

\vfill
\begin{tcolorbox}[enhanced,nobeforeafter,width=7.1cm,height=2.1cm,valign=center,colback=ccvp!6,colframe=ccvp,arc=2.5mm,boxrule=0.9pt,left=3mm]{\Large\bfseries\color{ccvp}a)}\quad\begin{minipage}[c]{5.3cm}\raggedright\Large $3{,}3$\end{minipage}\end{tcolorbox}\hfill\begin{tcolorbox}[enhanced,nobeforeafter,width=7.1cm,height=2.1cm,valign=center,colback=ccvp!6,colframe=ccvp,arc=2.5mm,boxrule=0.9pt,left=3mm]{\Large\bfseries\color{ccvp}b)}\quad\begin{minipage}[c]{5.3cm}\raggedright\Large $2$\end{minipage}\end{tcolorbox}\par\vspace{2.5mm}
\begin{tcolorbox}[enhanced,nobeforeafter,width=7.1cm,height=2.1cm,valign=center,colback=ccvp!6,colframe=ccvp,arc=2.5mm,boxrule=0.9pt,left=3mm]{\Large\bfseries\color{ccvp}c)}\quad\begin{minipage}[c]{5.3cm}\raggedright\Large $-2{,}3$\end{minipage}\end{tcolorbox}\hfill\begin{tcolorbox}[enhanced,nobeforeafter,width=7.1cm,height=2.1cm,valign=center,colback=ccvp!6,colframe=ccvp,arc=2.5mm,boxrule=0.9pt,left=3mm]{\Large\bfseries\color{ccvp}d)}\quad\begin{minipage}[c]{5.3cm}\raggedright\Large $2{,}3$\end{minipage}\end{tcolorbox}
\end{tcolorbox}
\newpage
\begin{tcolorbox}[enhanced,width=16cm,height=14cm,colback=white,colframe=ccvp,boxrule=2pt,arc=0mm,
left=6mm,right=6mm,top=0mm,bottom=4mm,sharp corners,
overlay={\fill[ccvp] (frame.north west) rectangle ([yshift=-1.25cm]frame.north east);
\node[anchor=west,text=white,font=\sffamily\bfseries\LARGE] at ([xshift=5mm,yshift=-0.62cm]frame.north west) {Question 91 B};
\node[anchor=east,text=white,font=\sffamily\itshape\large] at ([xshift=-5mm,yshift=-0.62cm]frame.north east) {Probabilités conditionnelles et variables aléatoires};}]
\vspace*{1.75cm}
\begin{minipage}[t][6.9cm][c]{\linewidth}\Large\centering
\begin{minipage}{0.95\linewidth}\raggedright\linespread{1.25}\selectfont Voici la loi de probabilité d'une variable aléatoire $X$. Que vaut sa variance $V(X)$~? $\begin{array}{|l|c|c|}\hline x_i & -2 & 2\\\hline P(X=x_i) & 0{,}5 & 0{,}5\\\hline\end{array}$\end{minipage}
\end{minipage}

\vfill
\begin{tcolorbox}[enhanced,nobeforeafter,width=7.1cm,height=2.1cm,valign=center,colback=ccvp!6,colframe=ccvp,arc=2.5mm,boxrule=0.9pt,left=3mm]{\Large\bfseries\color{ccvp}a)}\quad\begin{minipage}[c]{5.3cm}\raggedright\Large $5$\end{minipage}\end{tcolorbox}\hfill\begin{tcolorbox}[enhanced,nobeforeafter,width=7.1cm,height=2.1cm,valign=center,colback=ccvp!6,colframe=ccvp,arc=2.5mm,boxrule=0.9pt,left=3mm]{\Large\bfseries\color{ccvp}b)}\quad\begin{minipage}[c]{5.3cm}\raggedright\Large $0$\end{minipage}\end{tcolorbox}\par\vspace{2.5mm}
\begin{tcolorbox}[enhanced,nobeforeafter,width=7.1cm,height=2.1cm,valign=center,colback=ccvp!6,colframe=ccvp,arc=2.5mm,boxrule=0.9pt,left=3mm]{\Large\bfseries\color{ccvp}c)}\quad\begin{minipage}[c]{5.3cm}\raggedright\Large $4$\end{minipage}\end{tcolorbox}\hfill\begin{tcolorbox}[enhanced,nobeforeafter,width=7.1cm,height=2.1cm,valign=center,colback=ccvp!6,colframe=ccvp,arc=2.5mm,boxrule=0.9pt,left=3mm]{\Large\bfseries\color{ccvp}d)}\quad\begin{minipage}[c]{5.3cm}\raggedright\Large $6$\end{minipage}\end{tcolorbox}
\end{tcolorbox}
\newpage
\begin{tcolorbox}[enhanced,width=16cm,height=14cm,colback=white,colframe=ccvp,boxrule=2pt,arc=0mm,
left=6mm,right=6mm,top=0mm,bottom=4mm,sharp corners,
overlay={\fill[ccvp] (frame.north west) rectangle ([yshift=-1.25cm]frame.north east);
\node[anchor=west,text=white,font=\sffamily\bfseries\LARGE] at ([xshift=5mm,yshift=-0.62cm]frame.north west) {Question 92 B};
\node[anchor=east,text=white,font=\sffamily\itshape\large] at ([xshift=-5mm,yshift=-0.62cm]frame.north east) {Probabilités conditionnelles et variables aléatoires};}]
\vspace*{1.75cm}
\begin{minipage}[t][6.9cm][c]{\linewidth}\Large\centering
\begin{minipage}{0.95\linewidth}\raggedright\linespread{1.25}\selectfont Voici la loi de probabilité d'une variable aléatoire $X$. Que vaut sa variance $V(X)$~? $\begin{array}{|l|c|c|}\hline x_i & 0 & 10\\\hline P(X=x_i) & 0{,}6 & 0{,}4\\\hline\end{array}$\end{minipage}
\end{minipage}

\vfill
\begin{tcolorbox}[enhanced,nobeforeafter,width=7.1cm,height=2.1cm,valign=center,colback=ccvp!6,colframe=ccvp,arc=2.5mm,boxrule=0.9pt,left=3mm]{\Large\bfseries\color{ccvp}a)}\quad\begin{minipage}[c]{5.3cm}\raggedright\Large $24$\end{minipage}\end{tcolorbox}\hfill\begin{tcolorbox}[enhanced,nobeforeafter,width=7.1cm,height=2.1cm,valign=center,colback=ccvp!6,colframe=ccvp,arc=2.5mm,boxrule=0.9pt,left=3mm]{\Large\bfseries\color{ccvp}b)}\quad\begin{minipage}[c]{5.3cm}\raggedright\Large $4$\end{minipage}\end{tcolorbox}\par\vspace{2.5mm}
\begin{tcolorbox}[enhanced,nobeforeafter,width=7.1cm,height=2.1cm,valign=center,colback=ccvp!6,colframe=ccvp,arc=2.5mm,boxrule=0.9pt,left=3mm]{\Large\bfseries\color{ccvp}c)}\quad\begin{minipage}[c]{5.3cm}\raggedright\Large $25$\end{minipage}\end{tcolorbox}\hfill\begin{tcolorbox}[enhanced,nobeforeafter,width=7.1cm,height=2.1cm,valign=center,colback=ccvp!6,colframe=ccvp,arc=2.5mm,boxrule=0.9pt,left=3mm]{\Large\bfseries\color{ccvp}d)}\quad\begin{minipage}[c]{5.3cm}\raggedright\Large $40$\end{minipage}\end{tcolorbox}
\end{tcolorbox}
\newpage
\begin{tcolorbox}[enhanced,width=16cm,height=14cm,colback=white,colframe=ccvp,boxrule=2pt,arc=0mm,
left=6mm,right=6mm,top=0mm,bottom=4mm,sharp corners,
overlay={\fill[ccvp] (frame.north west) rectangle ([yshift=-1.25cm]frame.north east);
\node[anchor=west,text=white,font=\sffamily\bfseries\LARGE] at ([xshift=5mm,yshift=-0.62cm]frame.north west) {Question 93 B};
\node[anchor=east,text=white,font=\sffamily\itshape\large] at ([xshift=-5mm,yshift=-0.62cm]frame.north east) {Probabilités conditionnelles et variables aléatoires};}]
\vspace*{1.75cm}
\begin{minipage}[t][6.9cm][c]{\linewidth}\Large\centering
\begin{minipage}{0.95\linewidth}\raggedright\linespread{1.25}\selectfont On choisit au hasard une fille de ce groupe. Quelle est la probabilité que ce soit un demi-pensionnaire~? $\begin{array}{|l|c|c|}\hline & \text{Demi-pens.} & \text{Externes}\\\hline \text{Filles} & 6 & 14\\\hline \text{Garçons} & 18 & 12\\\hline\end{array}$\end{minipage}
\end{minipage}

\vfill
\begin{tcolorbox}[enhanced,nobeforeafter,width=7.1cm,height=2.1cm,valign=center,colback=ccvp!6,colframe=ccvp,arc=2.5mm,boxrule=0.9pt,left=3mm]{\Large\bfseries\color{ccvp}a)}\quad\begin{minipage}[c]{5.3cm}\raggedright\Large $\dfrac{7}{10}$\end{minipage}\end{tcolorbox}\hfill\begin{tcolorbox}[enhanced,nobeforeafter,width=7.1cm,height=2.1cm,valign=center,colback=ccvp!6,colframe=ccvp,arc=2.5mm,boxrule=0.9pt,left=3mm]{\Large\bfseries\color{ccvp}b)}\quad\begin{minipage}[c]{5.3cm}\raggedright\Large $\dfrac{1}{4}$\end{minipage}\end{tcolorbox}\par\vspace{2.5mm}
\begin{tcolorbox}[enhanced,nobeforeafter,width=7.1cm,height=2.1cm,valign=center,colback=ccvp!6,colframe=ccvp,arc=2.5mm,boxrule=0.9pt,left=3mm]{\Large\bfseries\color{ccvp}c)}\quad\begin{minipage}[c]{5.3cm}\raggedright\Large $\dfrac{3}{10}$\end{minipage}\end{tcolorbox}\hfill\begin{tcolorbox}[enhanced,nobeforeafter,width=7.1cm,height=2.1cm,valign=center,colback=ccvp!6,colframe=ccvp,arc=2.5mm,boxrule=0.9pt,left=3mm]{\Large\bfseries\color{ccvp}d)}\quad\begin{minipage}[c]{5.3cm}\raggedright\Large $\dfrac{3}{25}$\end{minipage}\end{tcolorbox}
\end{tcolorbox}
\newpage
\begin{tcolorbox}[enhanced,width=16cm,height=14cm,colback=white,colframe=ccvp,boxrule=2pt,arc=0mm,
left=6mm,right=6mm,top=0mm,bottom=4mm,sharp corners,
overlay={\fill[ccvp] (frame.north west) rectangle ([yshift=-1.25cm]frame.north east);
\node[anchor=west,text=white,font=\sffamily\bfseries\LARGE] at ([xshift=5mm,yshift=-0.62cm]frame.north west) {Question 94 B};
\node[anchor=east,text=white,font=\sffamily\itshape\large] at ([xshift=-5mm,yshift=-0.62cm]frame.north east) {Probabilités conditionnelles et variables aléatoires};}]
\vspace*{1.75cm}
\begin{minipage}[t][6.9cm][c]{\linewidth}\Large\centering
\begin{minipage}{0.95\linewidth}\raggedright\linespread{1.25}\selectfont On choisit au hasard un garçon de ce groupe. Quelle est la probabilité que ce soit un demi-pensionnaire~? $\begin{array}{|l|c|c|}\hline & \text{Demi-pens.} & \text{Externes}\\\hline \text{Filles} & 20 & 10\\\hline \text{Garçons} & 5 & 15\\\hline\end{array}$\end{minipage}
\end{minipage}

\vfill
\begin{tcolorbox}[enhanced,nobeforeafter,width=7.1cm,height=2.1cm,valign=center,colback=ccvp!6,colframe=ccvp,arc=2.5mm,boxrule=0.9pt,left=3mm]{\Large\bfseries\color{ccvp}a)}\quad\begin{minipage}[c]{5.3cm}\raggedright\Large $\dfrac{1}{4}$\end{minipage}\end{tcolorbox}\hfill\begin{tcolorbox}[enhanced,nobeforeafter,width=7.1cm,height=2.1cm,valign=center,colback=ccvp!6,colframe=ccvp,arc=2.5mm,boxrule=0.9pt,left=3mm]{\Large\bfseries\color{ccvp}b)}\quad\begin{minipage}[c]{5.3cm}\raggedright\Large $\dfrac{1}{5}$\end{minipage}\end{tcolorbox}\par\vspace{2.5mm}
\begin{tcolorbox}[enhanced,nobeforeafter,width=7.1cm,height=2.1cm,valign=center,colback=ccvp!6,colframe=ccvp,arc=2.5mm,boxrule=0.9pt,left=3mm]{\Large\bfseries\color{ccvp}c)}\quad\begin{minipage}[c]{5.3cm}\raggedright\Large $\dfrac{3}{4}$\end{minipage}\end{tcolorbox}\hfill\begin{tcolorbox}[enhanced,nobeforeafter,width=7.1cm,height=2.1cm,valign=center,colback=ccvp!6,colframe=ccvp,arc=2.5mm,boxrule=0.9pt,left=3mm]{\Large\bfseries\color{ccvp}d)}\quad\begin{minipage}[c]{5.3cm}\raggedright\Large $\dfrac{1}{10}$\end{minipage}\end{tcolorbox}
\end{tcolorbox}
\newpage
\begin{tcolorbox}[enhanced,width=16cm,height=14cm,colback=white,colframe=ccvp,boxrule=2pt,arc=0mm,
left=6mm,right=6mm,top=0mm,bottom=4mm,sharp corners,
overlay={\fill[ccvp] (frame.north west) rectangle ([yshift=-1.25cm]frame.north east);
\node[anchor=west,text=white,font=\sffamily\bfseries\LARGE] at ([xshift=5mm,yshift=-0.62cm]frame.north west) {Question 95 B};
\node[anchor=east,text=white,font=\sffamily\itshape\large] at ([xshift=-5mm,yshift=-0.62cm]frame.north east) {Probabilités conditionnelles et variables aléatoires};}]
\vspace*{1.75cm}
\begin{minipage}[t][6.9cm][c]{\linewidth}\Large\centering
\begin{minipage}{0.95\linewidth}\raggedright\linespread{1.25}\selectfont Voici la loi de probabilité d'une variable aléatoire $X$. Que vaut $P(X\geqslant 1)$~? $\begin{array}{|l|c|c|c|c|c|}\hline x_i & 0 & 1 & 2 & 3 & 4\\\hline P(X=x_i) & 0{,}05 & 0{,}15 & 0{,}4 & 0{,}3 & 0{,}1\\\hline\end{array}$\end{minipage}
\end{minipage}

\vfill
\begin{tcolorbox}[enhanced,nobeforeafter,width=7.1cm,height=2.1cm,valign=center,colback=ccvp!6,colframe=ccvp,arc=2.5mm,boxrule=0.9pt,left=3mm]{\Large\bfseries\color{ccvp}a)}\quad\begin{minipage}[c]{5.3cm}\raggedright\Large $0{,}95$\end{minipage}\end{tcolorbox}\hfill\begin{tcolorbox}[enhanced,nobeforeafter,width=7.1cm,height=2.1cm,valign=center,colback=ccvp!6,colframe=ccvp,arc=2.5mm,boxrule=0.9pt,left=3mm]{\Large\bfseries\color{ccvp}b)}\quad\begin{minipage}[c]{5.3cm}\raggedright\Large $0{,}15$\end{minipage}\end{tcolorbox}\par\vspace{2.5mm}
\begin{tcolorbox}[enhanced,nobeforeafter,width=7.1cm,height=2.1cm,valign=center,colback=ccvp!6,colframe=ccvp,arc=2.5mm,boxrule=0.9pt,left=3mm]{\Large\bfseries\color{ccvp}c)}\quad\begin{minipage}[c]{5.3cm}\raggedright\Large $0{,}8$\end{minipage}\end{tcolorbox}\hfill\begin{tcolorbox}[enhanced,nobeforeafter,width=7.1cm,height=2.1cm,valign=center,colback=ccvp!6,colframe=ccvp,arc=2.5mm,boxrule=0.9pt,left=3mm]{\Large\bfseries\color{ccvp}d)}\quad\begin{minipage}[c]{5.3cm}\raggedright\Large $0{,}05$\end{minipage}\end{tcolorbox}
\end{tcolorbox}
\newpage
\begin{tcolorbox}[enhanced,width=16cm,height=14cm,colback=white,colframe=ccvp,boxrule=2pt,arc=0mm,
left=6mm,right=6mm,top=0mm,bottom=4mm,sharp corners,
overlay={\fill[ccvp] (frame.north west) rectangle ([yshift=-1.25cm]frame.north east);
\node[anchor=west,text=white,font=\sffamily\bfseries\LARGE] at ([xshift=5mm,yshift=-0.62cm]frame.north west) {Question 96 B};
\node[anchor=east,text=white,font=\sffamily\itshape\large] at ([xshift=-5mm,yshift=-0.62cm]frame.north east) {Probabilités conditionnelles et variables aléatoires};}]
\vspace*{1.75cm}
\begin{minipage}[t][6.9cm][c]{\linewidth}\Large\centering
\begin{minipage}{0.95\linewidth}\raggedright\linespread{1.25}\selectfont Voici la loi de probabilité d'une variable aléatoire $X$. Que vaut $P(X\geqslant 3)$~? $\begin{array}{|l|c|c|c|c|c|}\hline x_i & 0 & 1 & 2 & 3 & 4\\\hline P(X=x_i) & 0{,}3 & 0{,}1 & 0{,}1 & 0{,}35 & 0{,}15\\\hline\end{array}$\end{minipage}
\end{minipage}

\vfill
\begin{tcolorbox}[enhanced,nobeforeafter,width=7.1cm,height=2.1cm,valign=center,colback=ccvp!6,colframe=ccvp,arc=2.5mm,boxrule=0.9pt,left=3mm]{\Large\bfseries\color{ccvp}a)}\quad\begin{minipage}[c]{5.3cm}\raggedright\Large $0{,}6$\end{minipage}\end{tcolorbox}\hfill\begin{tcolorbox}[enhanced,nobeforeafter,width=7.1cm,height=2.1cm,valign=center,colback=ccvp!6,colframe=ccvp,arc=2.5mm,boxrule=0.9pt,left=3mm]{\Large\bfseries\color{ccvp}b)}\quad\begin{minipage}[c]{5.3cm}\raggedright\Large $0{,}15$\end{minipage}\end{tcolorbox}\par\vspace{2.5mm}
\begin{tcolorbox}[enhanced,nobeforeafter,width=7.1cm,height=2.1cm,valign=center,colback=ccvp!6,colframe=ccvp,arc=2.5mm,boxrule=0.9pt,left=3mm]{\Large\bfseries\color{ccvp}c)}\quad\begin{minipage}[c]{5.3cm}\raggedright\Large $0{,}5$\end{minipage}\end{tcolorbox}\hfill\begin{tcolorbox}[enhanced,nobeforeafter,width=7.1cm,height=2.1cm,valign=center,colback=ccvp!6,colframe=ccvp,arc=2.5mm,boxrule=0.9pt,left=3mm]{\Large\bfseries\color{ccvp}d)}\quad\begin{minipage}[c]{5.3cm}\raggedright\Large $0{,}35$\end{minipage}\end{tcolorbox}
\end{tcolorbox}
\newpage
\begin{tcolorbox}[enhanced,width=16cm,height=14cm,colback=white,colframe=csui,boxrule=2pt,arc=0mm,
left=6mm,right=6mm,top=0mm,bottom=4mm,sharp corners,
overlay={\fill[csui] (frame.north west) rectangle ([yshift=-1.25cm]frame.north east);
\node[anchor=west,text=white,font=\sffamily\bfseries\LARGE] at ([xshift=5mm,yshift=-0.62cm]frame.north west) {Question 1 C};
\node[anchor=east,text=white,font=\sffamily\itshape\large] at ([xshift=-5mm,yshift=-0.62cm]frame.north east) {Suites};}]
\vspace*{1.75cm}
\begin{minipage}[t][6.9cm][c]{\linewidth}\LARGE\centering
\begin{minipage}{0.95\linewidth}\raggedright\linespread{1.25}\selectfont $(u_n)$ est la suite arithmétique de premier terme $u_0=2$ et de raison $-4$. Que vaut $u_{9}$~?\end{minipage}
\end{minipage}

\vfill
\begin{tcolorbox}[enhanced,nobeforeafter,width=7.1cm,height=2.1cm,valign=center,colback=csui!6,colframe=csui,arc=2.5mm,boxrule=0.9pt,left=3mm]{\Large\bfseries\color{csui}a)}\quad\begin{minipage}[c]{5.3cm}\raggedright\Large $-30$\end{minipage}\end{tcolorbox}\hfill\begin{tcolorbox}[enhanced,nobeforeafter,width=7.1cm,height=2.1cm,valign=center,colback=csui!6,colframe=csui,arc=2.5mm,boxrule=0.9pt,left=3mm]{\Large\bfseries\color{csui}b)}\quad\begin{minipage}[c]{5.3cm}\raggedright\Large $-36$\end{minipage}\end{tcolorbox}\par\vspace{2.5mm}
\begin{tcolorbox}[enhanced,nobeforeafter,width=7.1cm,height=2.1cm,valign=center,colback=csui!6,colframe=csui,arc=2.5mm,boxrule=0.9pt,left=3mm]{\Large\bfseries\color{csui}c)}\quad\begin{minipage}[c]{5.3cm}\raggedright\Large $-38$\end{minipage}\end{tcolorbox}\hfill\begin{tcolorbox}[enhanced,nobeforeafter,width=7.1cm,height=2.1cm,valign=center,colback=csui!6,colframe=csui,arc=2.5mm,boxrule=0.9pt,left=3mm]{\Large\bfseries\color{csui}d)}\quad\begin{minipage}[c]{5.3cm}\raggedright\Large $-34$\end{minipage}\end{tcolorbox}
\end{tcolorbox}
\newpage
\begin{tcolorbox}[enhanced,width=16cm,height=14cm,colback=white,colframe=csui,boxrule=2pt,arc=0mm,
left=6mm,right=6mm,top=0mm,bottom=4mm,sharp corners,
overlay={\fill[csui] (frame.north west) rectangle ([yshift=-1.25cm]frame.north east);
\node[anchor=west,text=white,font=\sffamily\bfseries\LARGE] at ([xshift=5mm,yshift=-0.62cm]frame.north west) {Question 2 C};
\node[anchor=east,text=white,font=\sffamily\itshape\large] at ([xshift=-5mm,yshift=-0.62cm]frame.north east) {Suites};}]
\vspace*{1.75cm}
\begin{minipage}[t][6.9cm][c]{\linewidth}\LARGE\centering
\begin{minipage}{0.95\linewidth}\raggedright\linespread{1.25}\selectfont $(u_n)$ est la suite arithmétique de premier terme $u_0=-6$ et de raison $3$. Que vaut $u_{11}$~?\end{minipage}
\end{minipage}

\vfill
\begin{tcolorbox}[enhanced,nobeforeafter,width=7.1cm,height=2.1cm,valign=center,colback=csui!6,colframe=csui,arc=2.5mm,boxrule=0.9pt,left=3mm]{\Large\bfseries\color{csui}a)}\quad\begin{minipage}[c]{5.3cm}\raggedright\Large $24$\end{minipage}\end{tcolorbox}\hfill\begin{tcolorbox}[enhanced,nobeforeafter,width=7.1cm,height=2.1cm,valign=center,colback=csui!6,colframe=csui,arc=2.5mm,boxrule=0.9pt,left=3mm]{\Large\bfseries\color{csui}b)}\quad\begin{minipage}[c]{5.3cm}\raggedright\Large $33$\end{minipage}\end{tcolorbox}\par\vspace{2.5mm}
\begin{tcolorbox}[enhanced,nobeforeafter,width=7.1cm,height=2.1cm,valign=center,colback=csui!6,colframe=csui,arc=2.5mm,boxrule=0.9pt,left=3mm]{\Large\bfseries\color{csui}c)}\quad\begin{minipage}[c]{5.3cm}\raggedright\Large $30$\end{minipage}\end{tcolorbox}\hfill\begin{tcolorbox}[enhanced,nobeforeafter,width=7.1cm,height=2.1cm,valign=center,colback=csui!6,colframe=csui,arc=2.5mm,boxrule=0.9pt,left=3mm]{\Large\bfseries\color{csui}d)}\quad\begin{minipage}[c]{5.3cm}\raggedright\Large $27$\end{minipage}\end{tcolorbox}
\end{tcolorbox}
\newpage
\begin{tcolorbox}[enhanced,width=16cm,height=14cm,colback=white,colframe=csui,boxrule=2pt,arc=0mm,
left=6mm,right=6mm,top=0mm,bottom=4mm,sharp corners,
overlay={\fill[csui] (frame.north west) rectangle ([yshift=-1.25cm]frame.north east);
\node[anchor=west,text=white,font=\sffamily\bfseries\LARGE] at ([xshift=5mm,yshift=-0.62cm]frame.north west) {Question 3 C};
\node[anchor=east,text=white,font=\sffamily\itshape\large] at ([xshift=-5mm,yshift=-0.62cm]frame.north east) {Suites};}]
\vspace*{1.75cm}
\begin{minipage}[t][6.9cm][c]{\linewidth}\LARGE\centering
\begin{minipage}{0.95\linewidth}\raggedright\linespread{1.25}\selectfont $(u_n)$ est la suite géométrique de premier terme $u_0=64$ et de raison $\dfrac{1}{2}$. Que vaut $u_{3}$~?\end{minipage}
\end{minipage}

\vfill
\begin{tcolorbox}[enhanced,nobeforeafter,width=7.1cm,height=2.1cm,valign=center,colback=csui!6,colframe=csui,arc=2.5mm,boxrule=0.9pt,left=3mm]{\Large\bfseries\color{csui}a)}\quad\begin{minipage}[c]{5.3cm}\raggedright\Large $8$\end{minipage}\end{tcolorbox}\hfill\begin{tcolorbox}[enhanced,nobeforeafter,width=7.1cm,height=2.1cm,valign=center,colback=csui!6,colframe=csui,arc=2.5mm,boxrule=0.9pt,left=3mm]{\Large\bfseries\color{csui}b)}\quad\begin{minipage}[c]{5.3cm}\raggedright\Large $4$\end{minipage}\end{tcolorbox}\par\vspace{2.5mm}
\begin{tcolorbox}[enhanced,nobeforeafter,width=7.1cm,height=2.1cm,valign=center,colback=csui!6,colframe=csui,arc=2.5mm,boxrule=0.9pt,left=3mm]{\Large\bfseries\color{csui}c)}\quad\begin{minipage}[c]{5.3cm}\raggedright\Large $16$\end{minipage}\end{tcolorbox}\hfill\begin{tcolorbox}[enhanced,nobeforeafter,width=7.1cm,height=2.1cm,valign=center,colback=csui!6,colframe=csui,arc=2.5mm,boxrule=0.9pt,left=3mm]{\Large\bfseries\color{csui}d)}\quad\begin{minipage}[c]{5.3cm}\raggedright\Large $\dfrac{131}{2}$\end{minipage}\end{tcolorbox}
\end{tcolorbox}
\newpage
\begin{tcolorbox}[enhanced,width=16cm,height=14cm,colback=white,colframe=csui,boxrule=2pt,arc=0mm,
left=6mm,right=6mm,top=0mm,bottom=4mm,sharp corners,
overlay={\fill[csui] (frame.north west) rectangle ([yshift=-1.25cm]frame.north east);
\node[anchor=west,text=white,font=\sffamily\bfseries\LARGE] at ([xshift=5mm,yshift=-0.62cm]frame.north west) {Question 4 C};
\node[anchor=east,text=white,font=\sffamily\itshape\large] at ([xshift=-5mm,yshift=-0.62cm]frame.north east) {Suites};}]
\vspace*{1.75cm}
\begin{minipage}[t][6.9cm][c]{\linewidth}\LARGE\centering
\begin{minipage}{0.95\linewidth}\raggedright\linespread{1.25}\selectfont $(u_n)$ est la suite géométrique de premier terme $u_0=-3$ et de raison $2$. Que vaut $u_{5}$~?\end{minipage}
\end{minipage}

\vfill
\begin{tcolorbox}[enhanced,nobeforeafter,width=7.1cm,height=2.1cm,valign=center,colback=csui!6,colframe=csui,arc=2.5mm,boxrule=0.9pt,left=3mm]{\Large\bfseries\color{csui}a)}\quad\begin{minipage}[c]{5.3cm}\raggedright\Large $-96$\end{minipage}\end{tcolorbox}\hfill\begin{tcolorbox}[enhanced,nobeforeafter,width=7.1cm,height=2.1cm,valign=center,colback=csui!6,colframe=csui,arc=2.5mm,boxrule=0.9pt,left=3mm]{\Large\bfseries\color{csui}b)}\quad\begin{minipage}[c]{5.3cm}\raggedright\Large $-48$\end{minipage}\end{tcolorbox}\par\vspace{2.5mm}
\begin{tcolorbox}[enhanced,nobeforeafter,width=7.1cm,height=2.1cm,valign=center,colback=csui!6,colframe=csui,arc=2.5mm,boxrule=0.9pt,left=3mm]{\Large\bfseries\color{csui}c)}\quad\begin{minipage}[c]{5.3cm}\raggedright\Large $-192$\end{minipage}\end{tcolorbox}\hfill\begin{tcolorbox}[enhanced,nobeforeafter,width=7.1cm,height=2.1cm,valign=center,colback=csui!6,colframe=csui,arc=2.5mm,boxrule=0.9pt,left=3mm]{\Large\bfseries\color{csui}d)}\quad\begin{minipage}[c]{5.3cm}\raggedright\Large $7$\end{minipage}\end{tcolorbox}
\end{tcolorbox}
\newpage
\begin{tcolorbox}[enhanced,width=16cm,height=14cm,colback=white,colframe=csui,boxrule=2pt,arc=0mm,
left=6mm,right=6mm,top=0mm,bottom=4mm,sharp corners,
overlay={\fill[csui] (frame.north west) rectangle ([yshift=-1.25cm]frame.north east);
\node[anchor=west,text=white,font=\sffamily\bfseries\LARGE] at ([xshift=5mm,yshift=-0.62cm]frame.north west) {Question 5 C};
\node[anchor=east,text=white,font=\sffamily\itshape\large] at ([xshift=-5mm,yshift=-0.62cm]frame.north east) {Suites};}]
\vspace*{1.75cm}
\begin{minipage}[t][6.9cm][c]{\linewidth}\LARGE\centering
\begin{minipage}{0.95\linewidth}\raggedright\linespread{1.25}\selectfont On définit $u_0=2$ et, pour tout entier $n$, $u_{n+1}=4u_n-5$. Que vaut $u_2$~?\end{minipage}
\end{minipage}

\vfill
\begin{tcolorbox}[enhanced,nobeforeafter,width=7.1cm,height=2.1cm,valign=center,colback=csui!6,colframe=csui,arc=2.5mm,boxrule=0.9pt,left=3mm]{\Large\bfseries\color{csui}a)}\quad\begin{minipage}[c]{5.3cm}\raggedright\Large $-2$\end{minipage}\end{tcolorbox}\hfill\begin{tcolorbox}[enhanced,nobeforeafter,width=7.1cm,height=2.1cm,valign=center,colback=csui!6,colframe=csui,arc=2.5mm,boxrule=0.9pt,left=3mm]{\Large\bfseries\color{csui}b)}\quad\begin{minipage}[c]{5.3cm}\raggedright\Large $7$\end{minipage}\end{tcolorbox}\par\vspace{2.5mm}
\begin{tcolorbox}[enhanced,nobeforeafter,width=7.1cm,height=2.1cm,valign=center,colback=csui!6,colframe=csui,arc=2.5mm,boxrule=0.9pt,left=3mm]{\Large\bfseries\color{csui}c)}\quad\begin{minipage}[c]{5.3cm}\raggedright\Large $23$\end{minipage}\end{tcolorbox}\hfill\begin{tcolorbox}[enhanced,nobeforeafter,width=7.1cm,height=2.1cm,valign=center,colback=csui!6,colframe=csui,arc=2.5mm,boxrule=0.9pt,left=3mm]{\Large\bfseries\color{csui}d)}\quad\begin{minipage}[c]{5.3cm}\raggedright\Large $3$\end{minipage}\end{tcolorbox}
\end{tcolorbox}
\newpage
\begin{tcolorbox}[enhanced,width=16cm,height=14cm,colback=white,colframe=csui,boxrule=2pt,arc=0mm,
left=6mm,right=6mm,top=0mm,bottom=4mm,sharp corners,
overlay={\fill[csui] (frame.north west) rectangle ([yshift=-1.25cm]frame.north east);
\node[anchor=west,text=white,font=\sffamily\bfseries\LARGE] at ([xshift=5mm,yshift=-0.62cm]frame.north west) {Question 6 C};
\node[anchor=east,text=white,font=\sffamily\itshape\large] at ([xshift=-5mm,yshift=-0.62cm]frame.north east) {Suites};}]
\vspace*{1.75cm}
\begin{minipage}[t][6.9cm][c]{\linewidth}\LARGE\centering
\begin{minipage}{0.95\linewidth}\raggedright\linespread{1.25}\selectfont On définit $u_0=5$ et, pour tout entier $n$, $u_{n+1}=2u_n-3$. Que vaut $u_2$~?\end{minipage}
\end{minipage}

\vfill
\begin{tcolorbox}[enhanced,nobeforeafter,width=7.1cm,height=2.1cm,valign=center,colback=csui!6,colframe=csui,arc=2.5mm,boxrule=0.9pt,left=3mm]{\Large\bfseries\color{csui}a)}\quad\begin{minipage}[c]{5.3cm}\raggedright\Large $7$\end{minipage}\end{tcolorbox}\hfill\begin{tcolorbox}[enhanced,nobeforeafter,width=7.1cm,height=2.1cm,valign=center,colback=csui!6,colframe=csui,arc=2.5mm,boxrule=0.9pt,left=3mm]{\Large\bfseries\color{csui}b)}\quad\begin{minipage}[c]{5.3cm}\raggedright\Large $4$\end{minipage}\end{tcolorbox}\par\vspace{2.5mm}
\begin{tcolorbox}[enhanced,nobeforeafter,width=7.1cm,height=2.1cm,valign=center,colback=csui!6,colframe=csui,arc=2.5mm,boxrule=0.9pt,left=3mm]{\Large\bfseries\color{csui}c)}\quad\begin{minipage}[c]{5.3cm}\raggedright\Large $11$\end{minipage}\end{tcolorbox}\hfill\begin{tcolorbox}[enhanced,nobeforeafter,width=7.1cm,height=2.1cm,valign=center,colback=csui!6,colframe=csui,arc=2.5mm,boxrule=0.9pt,left=3mm]{\Large\bfseries\color{csui}d)}\quad\begin{minipage}[c]{5.3cm}\raggedright\Large $19$\end{minipage}\end{tcolorbox}
\end{tcolorbox}
\newpage
\begin{tcolorbox}[enhanced,width=16cm,height=14cm,colback=white,colframe=csui,boxrule=2pt,arc=0mm,
left=6mm,right=6mm,top=0mm,bottom=4mm,sharp corners,
overlay={\fill[csui] (frame.north west) rectangle ([yshift=-1.25cm]frame.north east);
\node[anchor=west,text=white,font=\sffamily\bfseries\LARGE] at ([xshift=5mm,yshift=-0.62cm]frame.north west) {Question 7 C};
\node[anchor=east,text=white,font=\sffamily\itshape\large] at ([xshift=-5mm,yshift=-0.62cm]frame.north east) {Suites};}]
\vspace*{1.75cm}
\begin{minipage}[t][6.9cm][c]{\linewidth}\LARGE\centering
\begin{minipage}{0.95\linewidth}\raggedright\linespread{1.25}\selectfont Que vaut la somme $1+2+3+\dots+15$~?\end{minipage}
\end{minipage}

\vfill
\begin{tcolorbox}[enhanced,nobeforeafter,width=7.1cm,height=2.1cm,valign=center,colback=csui!6,colframe=csui,arc=2.5mm,boxrule=0.9pt,left=3mm]{\Large\bfseries\color{csui}a)}\quad\begin{minipage}[c]{5.3cm}\raggedright\Large $120$\end{minipage}\end{tcolorbox}\hfill\begin{tcolorbox}[enhanced,nobeforeafter,width=7.1cm,height=2.1cm,valign=center,colback=csui!6,colframe=csui,arc=2.5mm,boxrule=0.9pt,left=3mm]{\Large\bfseries\color{csui}b)}\quad\begin{minipage}[c]{5.3cm}\raggedright\Large $240$\end{minipage}\end{tcolorbox}\par\vspace{2.5mm}
\begin{tcolorbox}[enhanced,nobeforeafter,width=7.1cm,height=2.1cm,valign=center,colback=csui!6,colframe=csui,arc=2.5mm,boxrule=0.9pt,left=3mm]{\Large\bfseries\color{csui}c)}\quad\begin{minipage}[c]{5.3cm}\raggedright\Large $105$\end{minipage}\end{tcolorbox}\hfill\begin{tcolorbox}[enhanced,nobeforeafter,width=7.1cm,height=2.1cm,valign=center,colback=csui!6,colframe=csui,arc=2.5mm,boxrule=0.9pt,left=3mm]{\Large\bfseries\color{csui}d)}\quad\begin{minipage}[c]{5.3cm}\raggedright\Large $225$\end{minipage}\end{tcolorbox}
\end{tcolorbox}
\newpage
\begin{tcolorbox}[enhanced,width=16cm,height=14cm,colback=white,colframe=csui,boxrule=2pt,arc=0mm,
left=6mm,right=6mm,top=0mm,bottom=4mm,sharp corners,
overlay={\fill[csui] (frame.north west) rectangle ([yshift=-1.25cm]frame.north east);
\node[anchor=west,text=white,font=\sffamily\bfseries\LARGE] at ([xshift=5mm,yshift=-0.62cm]frame.north west) {Question 8 C};
\node[anchor=east,text=white,font=\sffamily\itshape\large] at ([xshift=-5mm,yshift=-0.62cm]frame.north east) {Suites};}]
\vspace*{1.75cm}
\begin{minipage}[t][6.9cm][c]{\linewidth}\LARGE\centering
\begin{minipage}{0.95\linewidth}\raggedright\linespread{1.25}\selectfont Que vaut la somme $1+2+3+\dots+30$~?\end{minipage}
\end{minipage}

\vfill
\begin{tcolorbox}[enhanced,nobeforeafter,width=7.1cm,height=2.1cm,valign=center,colback=csui!6,colframe=csui,arc=2.5mm,boxrule=0.9pt,left=3mm]{\Large\bfseries\color{csui}a)}\quad\begin{minipage}[c]{5.3cm}\raggedright\Large $435$\end{minipage}\end{tcolorbox}\hfill\begin{tcolorbox}[enhanced,nobeforeafter,width=7.1cm,height=2.1cm,valign=center,colback=csui!6,colframe=csui,arc=2.5mm,boxrule=0.9pt,left=3mm]{\Large\bfseries\color{csui}b)}\quad\begin{minipage}[c]{5.3cm}\raggedright\Large $900$\end{minipage}\end{tcolorbox}\par\vspace{2.5mm}
\begin{tcolorbox}[enhanced,nobeforeafter,width=7.1cm,height=2.1cm,valign=center,colback=csui!6,colframe=csui,arc=2.5mm,boxrule=0.9pt,left=3mm]{\Large\bfseries\color{csui}c)}\quad\begin{minipage}[c]{5.3cm}\raggedright\Large $465$\end{minipage}\end{tcolorbox}\hfill\begin{tcolorbox}[enhanced,nobeforeafter,width=7.1cm,height=2.1cm,valign=center,colback=csui!6,colframe=csui,arc=2.5mm,boxrule=0.9pt,left=3mm]{\Large\bfseries\color{csui}d)}\quad\begin{minipage}[c]{5.3cm}\raggedright\Large $930$\end{minipage}\end{tcolorbox}
\end{tcolorbox}
\newpage
\begin{tcolorbox}[enhanced,width=16cm,height=14cm,colback=white,colframe=csui,boxrule=2pt,arc=0mm,
left=6mm,right=6mm,top=0mm,bottom=4mm,sharp corners,
overlay={\fill[csui] (frame.north west) rectangle ([yshift=-1.25cm]frame.north east);
\node[anchor=west,text=white,font=\sffamily\bfseries\LARGE] at ([xshift=5mm,yshift=-0.62cm]frame.north west) {Question 9 C};
\node[anchor=east,text=white,font=\sffamily\itshape\large] at ([xshift=-5mm,yshift=-0.62cm]frame.north east) {Suites};}]
\vspace*{1.75cm}
\begin{minipage}[t][6.9cm][c]{\linewidth}\LARGE\centering
\begin{minipage}{0.95\linewidth}\raggedright\linespread{1.25}\selectfont $(u_n)$ est une suite arithmétique avec $u_{5}=20$ et $u_{10}=10$. Quelle est sa raison~?\end{minipage}
\end{minipage}

\vfill
\begin{tcolorbox}[enhanced,nobeforeafter,width=7.1cm,height=2.1cm,valign=center,colback=csui!6,colframe=csui,arc=2.5mm,boxrule=0.9pt,left=3mm]{\Large\bfseries\color{csui}a)}\quad\begin{minipage}[c]{5.3cm}\raggedright\Large $-2$\end{minipage}\end{tcolorbox}\hfill\begin{tcolorbox}[enhanced,nobeforeafter,width=7.1cm,height=2.1cm,valign=center,colback=csui!6,colframe=csui,arc=2.5mm,boxrule=0.9pt,left=3mm]{\Large\bfseries\color{csui}b)}\quad\begin{minipage}[c]{5.3cm}\raggedright\Large $-10$\end{minipage}\end{tcolorbox}\par\vspace{2.5mm}
\begin{tcolorbox}[enhanced,nobeforeafter,width=7.1cm,height=2.1cm,valign=center,colback=csui!6,colframe=csui,arc=2.5mm,boxrule=0.9pt,left=3mm]{\Large\bfseries\color{csui}c)}\quad\begin{minipage}[c]{5.3cm}\raggedright\Large $-\dfrac{2}{3}$\end{minipage}\end{tcolorbox}\hfill\begin{tcolorbox}[enhanced,nobeforeafter,width=7.1cm,height=2.1cm,valign=center,colback=csui!6,colframe=csui,arc=2.5mm,boxrule=0.9pt,left=3mm]{\Large\bfseries\color{csui}d)}\quad\begin{minipage}[c]{5.3cm}\raggedright\Large $2$\end{minipage}\end{tcolorbox}
\end{tcolorbox}
\newpage
\begin{tcolorbox}[enhanced,width=16cm,height=14cm,colback=white,colframe=csui,boxrule=2pt,arc=0mm,
left=6mm,right=6mm,top=0mm,bottom=4mm,sharp corners,
overlay={\fill[csui] (frame.north west) rectangle ([yshift=-1.25cm]frame.north east);
\node[anchor=west,text=white,font=\sffamily\bfseries\LARGE] at ([xshift=5mm,yshift=-0.62cm]frame.north west) {Question 10 C};
\node[anchor=east,text=white,font=\sffamily\itshape\large] at ([xshift=-5mm,yshift=-0.62cm]frame.north east) {Suites};}]
\vspace*{1.75cm}
\begin{minipage}[t][6.9cm][c]{\linewidth}\LARGE\centering
\begin{minipage}{0.95\linewidth}\raggedright\linespread{1.25}\selectfont $(u_n)$ est une suite arithmétique avec $u_{3}=1$ et $u_{9}=19$. Quelle est sa raison~?\end{minipage}
\end{minipage}

\vfill
\begin{tcolorbox}[enhanced,nobeforeafter,width=7.1cm,height=2.1cm,valign=center,colback=csui!6,colframe=csui,arc=2.5mm,boxrule=0.9pt,left=3mm]{\Large\bfseries\color{csui}a)}\quad\begin{minipage}[c]{5.3cm}\raggedright\Large $-3$\end{minipage}\end{tcolorbox}\hfill\begin{tcolorbox}[enhanced,nobeforeafter,width=7.1cm,height=2.1cm,valign=center,colback=csui!6,colframe=csui,arc=2.5mm,boxrule=0.9pt,left=3mm]{\Large\bfseries\color{csui}b)}\quad\begin{minipage}[c]{5.3cm}\raggedright\Large $3$\end{minipage}\end{tcolorbox}\par\vspace{2.5mm}
\begin{tcolorbox}[enhanced,nobeforeafter,width=7.1cm,height=2.1cm,valign=center,colback=csui!6,colframe=csui,arc=2.5mm,boxrule=0.9pt,left=3mm]{\Large\bfseries\color{csui}c)}\quad\begin{minipage}[c]{5.3cm}\raggedright\Large $18$\end{minipage}\end{tcolorbox}\hfill\begin{tcolorbox}[enhanced,nobeforeafter,width=7.1cm,height=2.1cm,valign=center,colback=csui!6,colframe=csui,arc=2.5mm,boxrule=0.9pt,left=3mm]{\Large\bfseries\color{csui}d)}\quad\begin{minipage}[c]{5.3cm}\raggedright\Large $\dfrac{3}{2}$\end{minipage}\end{tcolorbox}
\end{tcolorbox}
\newpage
\begin{tcolorbox}[enhanced,width=16cm,height=14cm,colback=white,colframe=csui,boxrule=2pt,arc=0mm,
left=6mm,right=6mm,top=0mm,bottom=4mm,sharp corners,
overlay={\fill[csui] (frame.north west) rectangle ([yshift=-1.25cm]frame.north east);
\node[anchor=west,text=white,font=\sffamily\bfseries\LARGE] at ([xshift=5mm,yshift=-0.62cm]frame.north west) {Question 11 C};
\node[anchor=east,text=white,font=\sffamily\itshape\large] at ([xshift=-5mm,yshift=-0.62cm]frame.north east) {Suites};}]
\vspace*{1.75cm}
\begin{minipage}[t][6.9cm][c]{\linewidth}\LARGE\centering
\begin{minipage}{0.95\linewidth}\raggedright\linespread{1.25}\selectfont Que vaut la somme $1+4+4^2+\dots+4^{4}$~?\end{minipage}
\end{minipage}

\vfill
\begin{tcolorbox}[enhanced,nobeforeafter,width=7.1cm,height=2.1cm,valign=center,colback=csui!6,colframe=csui,arc=2.5mm,boxrule=0.9pt,left=3mm]{\Large\bfseries\color{csui}a)}\quad\begin{minipage}[c]{5.3cm}\raggedright\Large $1\,024$\end{minipage}\end{tcolorbox}\hfill\begin{tcolorbox}[enhanced,nobeforeafter,width=7.1cm,height=2.1cm,valign=center,colback=csui!6,colframe=csui,arc=2.5mm,boxrule=0.9pt,left=3mm]{\Large\bfseries\color{csui}b)}\quad\begin{minipage}[c]{5.3cm}\raggedright\Large $1\,023$\end{minipage}\end{tcolorbox}\par\vspace{2.5mm}
\begin{tcolorbox}[enhanced,nobeforeafter,width=7.1cm,height=2.1cm,valign=center,colback=csui!6,colframe=csui,arc=2.5mm,boxrule=0.9pt,left=3mm]{\Large\bfseries\color{csui}c)}\quad\begin{minipage}[c]{5.3cm}\raggedright\Large $85$\end{minipage}\end{tcolorbox}\hfill\begin{tcolorbox}[enhanced,nobeforeafter,width=7.1cm,height=2.1cm,valign=center,colback=csui!6,colframe=csui,arc=2.5mm,boxrule=0.9pt,left=3mm]{\Large\bfseries\color{csui}d)}\quad\begin{minipage}[c]{5.3cm}\raggedright\Large $341$\end{minipage}\end{tcolorbox}
\end{tcolorbox}
\newpage
\begin{tcolorbox}[enhanced,width=16cm,height=14cm,colback=white,colframe=csui,boxrule=2pt,arc=0mm,
left=6mm,right=6mm,top=0mm,bottom=4mm,sharp corners,
overlay={\fill[csui] (frame.north west) rectangle ([yshift=-1.25cm]frame.north east);
\node[anchor=west,text=white,font=\sffamily\bfseries\LARGE] at ([xshift=5mm,yshift=-0.62cm]frame.north west) {Question 12 C};
\node[anchor=east,text=white,font=\sffamily\itshape\large] at ([xshift=-5mm,yshift=-0.62cm]frame.north east) {Suites};}]
\vspace*{1.75cm}
\begin{minipage}[t][6.9cm][c]{\linewidth}\LARGE\centering
\begin{minipage}{0.95\linewidth}\raggedright\linespread{1.25}\selectfont Que vaut la somme $1+3+3^2+\dots+3^{5}$~?\end{minipage}
\end{minipage}

\vfill
\begin{tcolorbox}[enhanced,nobeforeafter,width=7.1cm,height=2.1cm,valign=center,colback=csui!6,colframe=csui,arc=2.5mm,boxrule=0.9pt,left=3mm]{\Large\bfseries\color{csui}a)}\quad\begin{minipage}[c]{5.3cm}\raggedright\Large $364$\end{minipage}\end{tcolorbox}\hfill\begin{tcolorbox}[enhanced,nobeforeafter,width=7.1cm,height=2.1cm,valign=center,colback=csui!6,colframe=csui,arc=2.5mm,boxrule=0.9pt,left=3mm]{\Large\bfseries\color{csui}b)}\quad\begin{minipage}[c]{5.3cm}\raggedright\Large $728$\end{minipage}\end{tcolorbox}\par\vspace{2.5mm}
\begin{tcolorbox}[enhanced,nobeforeafter,width=7.1cm,height=2.1cm,valign=center,colback=csui!6,colframe=csui,arc=2.5mm,boxrule=0.9pt,left=3mm]{\Large\bfseries\color{csui}c)}\quad\begin{minipage}[c]{5.3cm}\raggedright\Large $121$\end{minipage}\end{tcolorbox}\hfill\begin{tcolorbox}[enhanced,nobeforeafter,width=7.1cm,height=2.1cm,valign=center,colback=csui!6,colframe=csui,arc=2.5mm,boxrule=0.9pt,left=3mm]{\Large\bfseries\color{csui}d)}\quad\begin{minipage}[c]{5.3cm}\raggedright\Large $729$\end{minipage}\end{tcolorbox}
\end{tcolorbox}
\newpage
\begin{tcolorbox}[enhanced,width=16cm,height=14cm,colback=white,colframe=csui,boxrule=2pt,arc=0mm,
left=6mm,right=6mm,top=0mm,bottom=4mm,sharp corners,
overlay={\fill[csui] (frame.north west) rectangle ([yshift=-1.25cm]frame.north east);
\node[anchor=west,text=white,font=\sffamily\bfseries\LARGE] at ([xshift=5mm,yshift=-0.62cm]frame.north west) {Question 13 C};
\node[anchor=east,text=white,font=\sffamily\itshape\large] at ([xshift=-5mm,yshift=-0.62cm]frame.north east) {Suites};}]
\vspace*{1.75cm}
\begin{minipage}[t][6.9cm][c]{\linewidth}\LARGE\centering
\begin{minipage}{0.95\linewidth}\raggedright\linespread{1.25}\selectfont Quel est le sens de variation de la suite définie pour tout entier naturel $n$ par $u_n=-2\times3^n$~?\end{minipage}
\end{minipage}

\vfill
\begin{tcolorbox}[enhanced,nobeforeafter,width=7.1cm,height=2.1cm,valign=center,colback=csui!6,colframe=csui,arc=2.5mm,boxrule=0.9pt,left=3mm]{\Large\bfseries\color{csui}a)}\quad\begin{minipage}[c]{5.3cm}\raggedright\large ni croissante ni décroissante\end{minipage}\end{tcolorbox}\hfill\begin{tcolorbox}[enhanced,nobeforeafter,width=7.1cm,height=2.1cm,valign=center,colback=csui!6,colframe=csui,arc=2.5mm,boxrule=0.9pt,left=3mm]{\Large\bfseries\color{csui}b)}\quad\begin{minipage}[c]{5.3cm}\raggedright\Large croissante\end{minipage}\end{tcolorbox}\par\vspace{2.5mm}
\begin{tcolorbox}[enhanced,nobeforeafter,width=7.1cm,height=2.1cm,valign=center,colback=csui!6,colframe=csui,arc=2.5mm,boxrule=0.9pt,left=3mm]{\Large\bfseries\color{csui}c)}\quad\begin{minipage}[c]{5.3cm}\raggedright\Large constante\end{minipage}\end{tcolorbox}\hfill\begin{tcolorbox}[enhanced,nobeforeafter,width=7.1cm,height=2.1cm,valign=center,colback=csui!6,colframe=csui,arc=2.5mm,boxrule=0.9pt,left=3mm]{\Large\bfseries\color{csui}d)}\quad\begin{minipage}[c]{5.3cm}\raggedright\Large décroissante\end{minipage}\end{tcolorbox}
\end{tcolorbox}
\newpage
\begin{tcolorbox}[enhanced,width=16cm,height=14cm,colback=white,colframe=csui,boxrule=2pt,arc=0mm,
left=6mm,right=6mm,top=0mm,bottom=4mm,sharp corners,
overlay={\fill[csui] (frame.north west) rectangle ([yshift=-1.25cm]frame.north east);
\node[anchor=west,text=white,font=\sffamily\bfseries\LARGE] at ([xshift=5mm,yshift=-0.62cm]frame.north west) {Question 14 C};
\node[anchor=east,text=white,font=\sffamily\itshape\large] at ([xshift=-5mm,yshift=-0.62cm]frame.north east) {Suites};}]
\vspace*{1.75cm}
\begin{minipage}[t][6.9cm][c]{\linewidth}\LARGE\centering
\begin{minipage}{0.95\linewidth}\raggedright\linespread{1.25}\selectfont Quel est le sens de variation de la suite définie pour tout entier naturel $n$ par $u_n=n^2+1$~?\end{minipage}
\end{minipage}

\vfill
\begin{tcolorbox}[enhanced,nobeforeafter,width=7.1cm,height=2.1cm,valign=center,colback=csui!6,colframe=csui,arc=2.5mm,boxrule=0.9pt,left=3mm]{\Large\bfseries\color{csui}a)}\quad\begin{minipage}[c]{5.3cm}\raggedright\Large décroissante\end{minipage}\end{tcolorbox}\hfill\begin{tcolorbox}[enhanced,nobeforeafter,width=7.1cm,height=2.1cm,valign=center,colback=csui!6,colframe=csui,arc=2.5mm,boxrule=0.9pt,left=3mm]{\Large\bfseries\color{csui}b)}\quad\begin{minipage}[c]{5.3cm}\raggedright\Large constante\end{minipage}\end{tcolorbox}\par\vspace{2.5mm}
\begin{tcolorbox}[enhanced,nobeforeafter,width=7.1cm,height=2.1cm,valign=center,colback=csui!6,colframe=csui,arc=2.5mm,boxrule=0.9pt,left=3mm]{\Large\bfseries\color{csui}c)}\quad\begin{minipage}[c]{5.3cm}\raggedright\Large croissante\end{minipage}\end{tcolorbox}\hfill\begin{tcolorbox}[enhanced,nobeforeafter,width=7.1cm,height=2.1cm,valign=center,colback=csui!6,colframe=csui,arc=2.5mm,boxrule=0.9pt,left=3mm]{\Large\bfseries\color{csui}d)}\quad\begin{minipage}[c]{5.3cm}\raggedright\large ni croissante ni décroissante\end{minipage}\end{tcolorbox}
\end{tcolorbox}
\newpage
\begin{tcolorbox}[enhanced,width=16cm,height=14cm,colback=white,colframe=csui,boxrule=2pt,arc=0mm,
left=6mm,right=6mm,top=0mm,bottom=4mm,sharp corners,
overlay={\fill[csui] (frame.north west) rectangle ([yshift=-1.25cm]frame.north east);
\node[anchor=west,text=white,font=\sffamily\bfseries\LARGE] at ([xshift=5mm,yshift=-0.62cm]frame.north west) {Question 15 C};
\node[anchor=east,text=white,font=\sffamily\itshape\large] at ([xshift=-5mm,yshift=-0.62cm]frame.north east) {Suites};}]
\vspace*{1.75cm}
\begin{minipage}[t][6.9cm][c]{\linewidth}\LARGE\centering
\begin{minipage}{0.95\linewidth}\raggedright\linespread{1.25}\selectfont Pour tout entier naturel $n$, $u_n=n^2-10$. Que vaut $u_{6}$~?\end{minipage}
\end{minipage}

\vfill
\begin{tcolorbox}[enhanced,nobeforeafter,width=7.1cm,height=2.1cm,valign=center,colback=csui!6,colframe=csui,arc=2.5mm,boxrule=0.9pt,left=3mm]{\Large\bfseries\color{csui}a)}\quad\begin{minipage}[c]{5.3cm}\raggedright\Large $-26$\end{minipage}\end{tcolorbox}\hfill\begin{tcolorbox}[enhanced,nobeforeafter,width=7.1cm,height=2.1cm,valign=center,colback=csui!6,colframe=csui,arc=2.5mm,boxrule=0.9pt,left=3mm]{\Large\bfseries\color{csui}b)}\quad\begin{minipage}[c]{5.3cm}\raggedright\Large $26$\end{minipage}\end{tcolorbox}\par\vspace{2.5mm}
\begin{tcolorbox}[enhanced,nobeforeafter,width=7.1cm,height=2.1cm,valign=center,colback=csui!6,colframe=csui,arc=2.5mm,boxrule=0.9pt,left=3mm]{\Large\bfseries\color{csui}c)}\quad\begin{minipage}[c]{5.3cm}\raggedright\Large $15$\end{minipage}\end{tcolorbox}\hfill\begin{tcolorbox}[enhanced,nobeforeafter,width=7.1cm,height=2.1cm,valign=center,colback=csui!6,colframe=csui,arc=2.5mm,boxrule=0.9pt,left=3mm]{\Large\bfseries\color{csui}d)}\quad\begin{minipage}[c]{5.3cm}\raggedright\Large $39$\end{minipage}\end{tcolorbox}
\end{tcolorbox}
\newpage
\begin{tcolorbox}[enhanced,width=16cm,height=14cm,colback=white,colframe=csui,boxrule=2pt,arc=0mm,
left=6mm,right=6mm,top=0mm,bottom=4mm,sharp corners,
overlay={\fill[csui] (frame.north west) rectangle ([yshift=-1.25cm]frame.north east);
\node[anchor=west,text=white,font=\sffamily\bfseries\LARGE] at ([xshift=5mm,yshift=-0.62cm]frame.north west) {Question 16 C};
\node[anchor=east,text=white,font=\sffamily\itshape\large] at ([xshift=-5mm,yshift=-0.62cm]frame.north east) {Suites};}]
\vspace*{1.75cm}
\begin{minipage}[t][6.9cm][c]{\linewidth}\LARGE\centering
\begin{minipage}{0.95\linewidth}\raggedright\linespread{1.25}\selectfont Pour tout entier naturel $n$, $u_n=2n^2-5n+3$. Que vaut $u_{4}$~?\end{minipage}
\end{minipage}

\vfill
\begin{tcolorbox}[enhanced,nobeforeafter,width=7.1cm,height=2.1cm,valign=center,colback=csui!6,colframe=csui,arc=2.5mm,boxrule=0.9pt,left=3mm]{\Large\bfseries\color{csui}a)}\quad\begin{minipage}[c]{5.3cm}\raggedright\Large $-15$\end{minipage}\end{tcolorbox}\hfill\begin{tcolorbox}[enhanced,nobeforeafter,width=7.1cm,height=2.1cm,valign=center,colback=csui!6,colframe=csui,arc=2.5mm,boxrule=0.9pt,left=3mm]{\Large\bfseries\color{csui}b)}\quad\begin{minipage}[c]{5.3cm}\raggedright\Large $15$\end{minipage}\end{tcolorbox}\par\vspace{2.5mm}
\begin{tcolorbox}[enhanced,nobeforeafter,width=7.1cm,height=2.1cm,valign=center,colback=csui!6,colframe=csui,arc=2.5mm,boxrule=0.9pt,left=3mm]{\Large\bfseries\color{csui}c)}\quad\begin{minipage}[c]{5.3cm}\raggedright\Large $28$\end{minipage}\end{tcolorbox}\hfill\begin{tcolorbox}[enhanced,nobeforeafter,width=7.1cm,height=2.1cm,valign=center,colback=csui!6,colframe=csui,arc=2.5mm,boxrule=0.9pt,left=3mm]{\Large\bfseries\color{csui}d)}\quad\begin{minipage}[c]{5.3cm}\raggedright\Large $6$\end{minipage}\end{tcolorbox}
\end{tcolorbox}
\newpage
\begin{tcolorbox}[enhanced,width=16cm,height=14cm,colback=white,colframe=cdeg,boxrule=2pt,arc=0mm,
left=6mm,right=6mm,top=0mm,bottom=4mm,sharp corners,
overlay={\fill[cdeg] (frame.north west) rectangle ([yshift=-1.25cm]frame.north east);
\node[anchor=west,text=white,font=\sffamily\bfseries\LARGE] at ([xshift=5mm,yshift=-0.62cm]frame.north west) {Question 17 C};
\node[anchor=east,text=white,font=\sffamily\itshape\large] at ([xshift=-5mm,yshift=-0.62cm]frame.north east) {Second degré};}]
\vspace*{1.75cm}
\begin{minipage}[t][6.9cm][c]{\linewidth}\LARGE\centering
\begin{minipage}{0.95\linewidth}\raggedright\linespread{1.25}\selectfont Quel est le discriminant du trinôme $-x^2+6x-9$~?\end{minipage}
\end{minipage}

\vfill
\begin{tcolorbox}[enhanced,nobeforeafter,width=7.1cm,height=2.1cm,valign=center,colback=cdeg!6,colframe=cdeg,arc=2.5mm,boxrule=0.9pt,left=3mm]{\Large\bfseries\color{cdeg}a)}\quad\begin{minipage}[c]{5.3cm}\raggedright\Large $-30$\end{minipage}\end{tcolorbox}\hfill\begin{tcolorbox}[enhanced,nobeforeafter,width=7.1cm,height=2.1cm,valign=center,colback=cdeg!6,colframe=cdeg,arc=2.5mm,boxrule=0.9pt,left=3mm]{\Large\bfseries\color{cdeg}b)}\quad\begin{minipage}[c]{5.3cm}\raggedright\Large $0$\end{minipage}\end{tcolorbox}\par\vspace{2.5mm}
\begin{tcolorbox}[enhanced,nobeforeafter,width=7.1cm,height=2.1cm,valign=center,colback=cdeg!6,colframe=cdeg,arc=2.5mm,boxrule=0.9pt,left=3mm]{\Large\bfseries\color{cdeg}c)}\quad\begin{minipage}[c]{5.3cm}\raggedright\Large $1$\end{minipage}\end{tcolorbox}\hfill\begin{tcolorbox}[enhanced,nobeforeafter,width=7.1cm,height=2.1cm,valign=center,colback=cdeg!6,colframe=cdeg,arc=2.5mm,boxrule=0.9pt,left=3mm]{\Large\bfseries\color{cdeg}d)}\quad\begin{minipage}[c]{5.3cm}\raggedright\Large $72$\end{minipage}\end{tcolorbox}
\end{tcolorbox}
\newpage
\begin{tcolorbox}[enhanced,width=16cm,height=14cm,colback=white,colframe=cdeg,boxrule=2pt,arc=0mm,
left=6mm,right=6mm,top=0mm,bottom=4mm,sharp corners,
overlay={\fill[cdeg] (frame.north west) rectangle ([yshift=-1.25cm]frame.north east);
\node[anchor=west,text=white,font=\sffamily\bfseries\LARGE] at ([xshift=5mm,yshift=-0.62cm]frame.north west) {Question 18 C};
\node[anchor=east,text=white,font=\sffamily\itshape\large] at ([xshift=-5mm,yshift=-0.62cm]frame.north east) {Second degré};}]
\vspace*{1.75cm}
\begin{minipage}[t][6.9cm][c]{\linewidth}\LARGE\centering
\begin{minipage}{0.95\linewidth}\raggedright\linespread{1.25}\selectfont Quel est le discriminant du trinôme $2x^2-7x+3$~?\end{minipage}
\end{minipage}

\vfill
\begin{tcolorbox}[enhanced,nobeforeafter,width=7.1cm,height=2.1cm,valign=center,colback=cdeg!6,colframe=cdeg,arc=2.5mm,boxrule=0.9pt,left=3mm]{\Large\bfseries\color{cdeg}a)}\quad\begin{minipage}[c]{5.3cm}\raggedright\Large $25$\end{minipage}\end{tcolorbox}\hfill\begin{tcolorbox}[enhanced,nobeforeafter,width=7.1cm,height=2.1cm,valign=center,colback=cdeg!6,colframe=cdeg,arc=2.5mm,boxrule=0.9pt,left=3mm]{\Large\bfseries\color{cdeg}b)}\quad\begin{minipage}[c]{5.3cm}\raggedright\Large $73$\end{minipage}\end{tcolorbox}\par\vspace{2.5mm}
\begin{tcolorbox}[enhanced,nobeforeafter,width=7.1cm,height=2.1cm,valign=center,colback=cdeg!6,colframe=cdeg,arc=2.5mm,boxrule=0.9pt,left=3mm]{\Large\bfseries\color{cdeg}c)}\quad\begin{minipage}[c]{5.3cm}\raggedright\Large $-31$\end{minipage}\end{tcolorbox}\hfill\begin{tcolorbox}[enhanced,nobeforeafter,width=7.1cm,height=2.1cm,valign=center,colback=cdeg!6,colframe=cdeg,arc=2.5mm,boxrule=0.9pt,left=3mm]{\Large\bfseries\color{cdeg}d)}\quad\begin{minipage}[c]{5.3cm}\raggedright\Large $-25$\end{minipage}\end{tcolorbox}
\end{tcolorbox}
\newpage
\begin{tcolorbox}[enhanced,width=16cm,height=14cm,colback=white,colframe=cdeg,boxrule=2pt,arc=0mm,
left=6mm,right=6mm,top=0mm,bottom=4mm,sharp corners,
overlay={\fill[cdeg] (frame.north west) rectangle ([yshift=-1.25cm]frame.north east);
\node[anchor=west,text=white,font=\sffamily\bfseries\LARGE] at ([xshift=5mm,yshift=-0.62cm]frame.north west) {Question 19 C};
\node[anchor=east,text=white,font=\sffamily\itshape\large] at ([xshift=-5mm,yshift=-0.62cm]frame.north east) {Second degré};}]
\vspace*{1.75cm}
\begin{minipage}[t][6.9cm][c]{\linewidth}\LARGE\centering
\begin{minipage}{0.95\linewidth}\raggedright\linespread{1.25}\selectfont Combien de racines réelles le trinôme $4x^2+4x+1$ possède-t-il~?\end{minipage}
\end{minipage}

\vfill
\begin{tcolorbox}[enhanced,nobeforeafter,width=7.1cm,height=2.1cm,valign=center,colback=cdeg!6,colframe=cdeg,arc=2.5mm,boxrule=0.9pt,left=3mm]{\Large\bfseries\color{cdeg}a)}\quad\begin{minipage}[c]{5.3cm}\raggedright\Large $0$\end{minipage}\end{tcolorbox}\hfill\begin{tcolorbox}[enhanced,nobeforeafter,width=7.1cm,height=2.1cm,valign=center,colback=cdeg!6,colframe=cdeg,arc=2.5mm,boxrule=0.9pt,left=3mm]{\Large\bfseries\color{cdeg}b)}\quad\begin{minipage}[c]{5.3cm}\raggedright\Large $1$\end{minipage}\end{tcolorbox}\par\vspace{2.5mm}
\begin{tcolorbox}[enhanced,nobeforeafter,width=7.1cm,height=2.1cm,valign=center,colback=cdeg!6,colframe=cdeg,arc=2.5mm,boxrule=0.9pt,left=3mm]{\Large\bfseries\color{cdeg}c)}\quad\begin{minipage}[c]{5.3cm}\raggedright\Large une infinité\end{minipage}\end{tcolorbox}\hfill\begin{tcolorbox}[enhanced,nobeforeafter,width=7.1cm,height=2.1cm,valign=center,colback=cdeg!6,colframe=cdeg,arc=2.5mm,boxrule=0.9pt,left=3mm]{\Large\bfseries\color{cdeg}d)}\quad\begin{minipage}[c]{5.3cm}\raggedright\Large $2$\end{minipage}\end{tcolorbox}
\end{tcolorbox}
\newpage
\begin{tcolorbox}[enhanced,width=16cm,height=14cm,colback=white,colframe=cdeg,boxrule=2pt,arc=0mm,
left=6mm,right=6mm,top=0mm,bottom=4mm,sharp corners,
overlay={\fill[cdeg] (frame.north west) rectangle ([yshift=-1.25cm]frame.north east);
\node[anchor=west,text=white,font=\sffamily\bfseries\LARGE] at ([xshift=5mm,yshift=-0.62cm]frame.north west) {Question 20 C};
\node[anchor=east,text=white,font=\sffamily\itshape\large] at ([xshift=-5mm,yshift=-0.62cm]frame.north east) {Second degré};}]
\vspace*{1.75cm}
\begin{minipage}[t][6.9cm][c]{\linewidth}\LARGE\centering
\begin{minipage}{0.95\linewidth}\raggedright\linespread{1.25}\selectfont Combien de racines réelles le trinôme $x^2+3$ possède-t-il~?\end{minipage}
\end{minipage}

\vfill
\begin{tcolorbox}[enhanced,nobeforeafter,width=7.1cm,height=2.1cm,valign=center,colback=cdeg!6,colframe=cdeg,arc=2.5mm,boxrule=0.9pt,left=3mm]{\Large\bfseries\color{cdeg}a)}\quad\begin{minipage}[c]{5.3cm}\raggedright\Large $2$\end{minipage}\end{tcolorbox}\hfill\begin{tcolorbox}[enhanced,nobeforeafter,width=7.1cm,height=2.1cm,valign=center,colback=cdeg!6,colframe=cdeg,arc=2.5mm,boxrule=0.9pt,left=3mm]{\Large\bfseries\color{cdeg}b)}\quad\begin{minipage}[c]{5.3cm}\raggedright\Large $1$\end{minipage}\end{tcolorbox}\par\vspace{2.5mm}
\begin{tcolorbox}[enhanced,nobeforeafter,width=7.1cm,height=2.1cm,valign=center,colback=cdeg!6,colframe=cdeg,arc=2.5mm,boxrule=0.9pt,left=3mm]{\Large\bfseries\color{cdeg}c)}\quad\begin{minipage}[c]{5.3cm}\raggedright\Large $0$\end{minipage}\end{tcolorbox}\hfill\begin{tcolorbox}[enhanced,nobeforeafter,width=7.1cm,height=2.1cm,valign=center,colback=cdeg!6,colframe=cdeg,arc=2.5mm,boxrule=0.9pt,left=3mm]{\Large\bfseries\color{cdeg}d)}\quad\begin{minipage}[c]{5.3cm}\raggedright\Large une infinité\end{minipage}\end{tcolorbox}
\end{tcolorbox}
\newpage
\begin{tcolorbox}[enhanced,width=16cm,height=14cm,colback=white,colframe=cdeg,boxrule=2pt,arc=0mm,
left=6mm,right=6mm,top=0mm,bottom=4mm,sharp corners,
overlay={\fill[cdeg] (frame.north west) rectangle ([yshift=-1.25cm]frame.north east);
\node[anchor=west,text=white,font=\sffamily\bfseries\LARGE] at ([xshift=5mm,yshift=-0.62cm]frame.north west) {Question 21 C};
\node[anchor=east,text=white,font=\sffamily\itshape\large] at ([xshift=-5mm,yshift=-0.62cm]frame.north east) {Second degré};}]
\vspace*{1.75cm}
\begin{minipage}[t][6.9cm][c]{\linewidth}\LARGE\centering
\begin{minipage}{0.95\linewidth}\raggedright\linespread{1.25}\selectfont Quelles sont les racines du trinôme $-x^2+8x-12$~?\end{minipage}
\end{minipage}

\vfill
\begin{tcolorbox}[enhanced,nobeforeafter,width=7.1cm,height=2.1cm,valign=center,colback=cdeg!6,colframe=cdeg,arc=2.5mm,boxrule=0.9pt,left=3mm]{\Large\bfseries\color{cdeg}a)}\quad\begin{minipage}[c]{5.3cm}\raggedright\Large $\{-6\,;\,2\}$\end{minipage}\end{tcolorbox}\hfill\begin{tcolorbox}[enhanced,nobeforeafter,width=7.1cm,height=2.1cm,valign=center,colback=cdeg!6,colframe=cdeg,arc=2.5mm,boxrule=0.9pt,left=3mm]{\Large\bfseries\color{cdeg}b)}\quad\begin{minipage}[c]{5.3cm}\raggedright\Large $\{2\,;\,6\}$\end{minipage}\end{tcolorbox}\par\vspace{2.5mm}
\begin{tcolorbox}[enhanced,nobeforeafter,width=7.1cm,height=2.1cm,valign=center,colback=cdeg!6,colframe=cdeg,arc=2.5mm,boxrule=0.9pt,left=3mm]{\Large\bfseries\color{cdeg}c)}\quad\begin{minipage}[c]{5.3cm}\raggedright\Large $\{8\,;\,12\}$\end{minipage}\end{tcolorbox}\hfill\begin{tcolorbox}[enhanced,nobeforeafter,width=7.1cm,height=2.1cm,valign=center,colback=cdeg!6,colframe=cdeg,arc=2.5mm,boxrule=0.9pt,left=3mm]{\Large\bfseries\color{cdeg}d)}\quad\begin{minipage}[c]{5.3cm}\raggedright\Large $\{-6\,;\,-2\}$\end{minipage}\end{tcolorbox}
\end{tcolorbox}
\newpage
\begin{tcolorbox}[enhanced,width=16cm,height=14cm,colback=white,colframe=cdeg,boxrule=2pt,arc=0mm,
left=6mm,right=6mm,top=0mm,bottom=4mm,sharp corners,
overlay={\fill[cdeg] (frame.north west) rectangle ([yshift=-1.25cm]frame.north east);
\node[anchor=west,text=white,font=\sffamily\bfseries\LARGE] at ([xshift=5mm,yshift=-0.62cm]frame.north west) {Question 22 C};
\node[anchor=east,text=white,font=\sffamily\itshape\large] at ([xshift=-5mm,yshift=-0.62cm]frame.north east) {Second degré};}]
\vspace*{1.75cm}
\begin{minipage}[t][6.9cm][c]{\linewidth}\LARGE\centering
\begin{minipage}{0.95\linewidth}\raggedright\linespread{1.25}\selectfont Quelles sont les racines du trinôme $3x^2-9x+6$~?\end{minipage}
\end{minipage}

\vfill
\begin{tcolorbox}[enhanced,nobeforeafter,width=7.1cm,height=2.1cm,valign=center,colback=cdeg!6,colframe=cdeg,arc=2.5mm,boxrule=0.9pt,left=3mm]{\Large\bfseries\color{cdeg}a)}\quad\begin{minipage}[c]{5.3cm}\raggedright\Large $\{1\,;\,2\}$\end{minipage}\end{tcolorbox}\hfill\begin{tcolorbox}[enhanced,nobeforeafter,width=7.1cm,height=2.1cm,valign=center,colback=cdeg!6,colframe=cdeg,arc=2.5mm,boxrule=0.9pt,left=3mm]{\Large\bfseries\color{cdeg}b)}\quad\begin{minipage}[c]{5.3cm}\raggedright\Large $\{2\,;\,3\}$\end{minipage}\end{tcolorbox}\par\vspace{2.5mm}
\begin{tcolorbox}[enhanced,nobeforeafter,width=7.1cm,height=2.1cm,valign=center,colback=cdeg!6,colframe=cdeg,arc=2.5mm,boxrule=0.9pt,left=3mm]{\Large\bfseries\color{cdeg}c)}\quad\begin{minipage}[c]{5.3cm}\raggedright\Large $\{-2\,;\,-1\}$\end{minipage}\end{tcolorbox}\hfill\begin{tcolorbox}[enhanced,nobeforeafter,width=7.1cm,height=2.1cm,valign=center,colback=cdeg!6,colframe=cdeg,arc=2.5mm,boxrule=0.9pt,left=3mm]{\Large\bfseries\color{cdeg}d)}\quad\begin{minipage}[c]{5.3cm}\raggedright\Large $\{-2\,;\,1\}$\end{minipage}\end{tcolorbox}
\end{tcolorbox}
\newpage
\begin{tcolorbox}[enhanced,width=16cm,height=14cm,colback=white,colframe=cdeg,boxrule=2pt,arc=0mm,
left=6mm,right=6mm,top=0mm,bottom=4mm,sharp corners,
overlay={\fill[cdeg] (frame.north west) rectangle ([yshift=-1.25cm]frame.north east);
\node[anchor=west,text=white,font=\sffamily\bfseries\LARGE] at ([xshift=5mm,yshift=-0.62cm]frame.north west) {Question 23 C};
\node[anchor=east,text=white,font=\sffamily\itshape\large] at ([xshift=-5mm,yshift=-0.62cm]frame.north east) {Second degré};}]
\vspace*{1.75cm}
\begin{minipage}[t][6.9cm][c]{\linewidth}\LARGE\centering
\begin{minipage}{0.95\linewidth}\raggedright\linespread{1.25}\selectfont Quelle est la forme canonique de $x^2+8x+10$~?\end{minipage}
\end{minipage}

\vfill
\begin{tcolorbox}[enhanced,nobeforeafter,width=7.1cm,height=2.1cm,valign=center,colback=cdeg!6,colframe=cdeg,arc=2.5mm,boxrule=0.9pt,left=3mm]{\Large\bfseries\color{cdeg}a)}\quad\begin{minipage}[c]{5.3cm}\raggedright\Large $(x+4)^2+10$\end{minipage}\end{tcolorbox}\hfill\begin{tcolorbox}[enhanced,nobeforeafter,width=7.1cm,height=2.1cm,valign=center,colback=cdeg!6,colframe=cdeg,arc=2.5mm,boxrule=0.9pt,left=3mm]{\Large\bfseries\color{cdeg}b)}\quad\begin{minipage}[c]{5.3cm}\raggedright\Large $(x-4)^2-6$\end{minipage}\end{tcolorbox}\par\vspace{2.5mm}
\begin{tcolorbox}[enhanced,nobeforeafter,width=7.1cm,height=2.1cm,valign=center,colback=cdeg!6,colframe=cdeg,arc=2.5mm,boxrule=0.9pt,left=3mm]{\Large\bfseries\color{cdeg}c)}\quad\begin{minipage}[c]{5.3cm}\raggedright\Large $(x+4)^2-6$\end{minipage}\end{tcolorbox}\hfill\begin{tcolorbox}[enhanced,nobeforeafter,width=7.1cm,height=2.1cm,valign=center,colback=cdeg!6,colframe=cdeg,arc=2.5mm,boxrule=0.9pt,left=3mm]{\Large\bfseries\color{cdeg}d)}\quad\begin{minipage}[c]{5.3cm}\raggedright\Large $(x+8)^2-6$\end{minipage}\end{tcolorbox}
\end{tcolorbox}
\newpage
\begin{tcolorbox}[enhanced,width=16cm,height=14cm,colback=white,colframe=cdeg,boxrule=2pt,arc=0mm,
left=6mm,right=6mm,top=0mm,bottom=4mm,sharp corners,
overlay={\fill[cdeg] (frame.north west) rectangle ([yshift=-1.25cm]frame.north east);
\node[anchor=west,text=white,font=\sffamily\bfseries\LARGE] at ([xshift=5mm,yshift=-0.62cm]frame.north west) {Question 24 C};
\node[anchor=east,text=white,font=\sffamily\itshape\large] at ([xshift=-5mm,yshift=-0.62cm]frame.north east) {Second degré};}]
\vspace*{1.75cm}
\begin{minipage}[t][6.9cm][c]{\linewidth}\LARGE\centering
\begin{minipage}{0.95\linewidth}\raggedright\linespread{1.25}\selectfont Quelle est la forme canonique de $x^2-6x+2$~?\end{minipage}
\end{minipage}

\vfill
\begin{tcolorbox}[enhanced,nobeforeafter,width=7.1cm,height=2.1cm,valign=center,colback=cdeg!6,colframe=cdeg,arc=2.5mm,boxrule=0.9pt,left=3mm]{\Large\bfseries\color{cdeg}a)}\quad\begin{minipage}[c]{5.3cm}\raggedright\Large $(x-3)^2+2$\end{minipage}\end{tcolorbox}\hfill\begin{tcolorbox}[enhanced,nobeforeafter,width=7.1cm,height=2.1cm,valign=center,colback=cdeg!6,colframe=cdeg,arc=2.5mm,boxrule=0.9pt,left=3mm]{\Large\bfseries\color{cdeg}b)}\quad\begin{minipage}[c]{5.3cm}\raggedright\Large $(x-6)^2-7$\end{minipage}\end{tcolorbox}\par\vspace{2.5mm}
\begin{tcolorbox}[enhanced,nobeforeafter,width=7.1cm,height=2.1cm,valign=center,colback=cdeg!6,colframe=cdeg,arc=2.5mm,boxrule=0.9pt,left=3mm]{\Large\bfseries\color{cdeg}c)}\quad\begin{minipage}[c]{5.3cm}\raggedright\Large $(x-3)^2-7$\end{minipage}\end{tcolorbox}\hfill\begin{tcolorbox}[enhanced,nobeforeafter,width=7.1cm,height=2.1cm,valign=center,colback=cdeg!6,colframe=cdeg,arc=2.5mm,boxrule=0.9pt,left=3mm]{\Large\bfseries\color{cdeg}d)}\quad\begin{minipage}[c]{5.3cm}\raggedright\Large $(x+3)^2-7$\end{minipage}\end{tcolorbox}
\end{tcolorbox}
\newpage
\begin{tcolorbox}[enhanced,width=16cm,height=14cm,colback=white,colframe=cdeg,boxrule=2pt,arc=0mm,
left=6mm,right=6mm,top=0mm,bottom=4mm,sharp corners,
overlay={\fill[cdeg] (frame.north west) rectangle ([yshift=-1.25cm]frame.north east);
\node[anchor=west,text=white,font=\sffamily\bfseries\LARGE] at ([xshift=5mm,yshift=-0.62cm]frame.north west) {Question 25 C};
\node[anchor=east,text=white,font=\sffamily\itshape\large] at ([xshift=-5mm,yshift=-0.62cm]frame.north east) {Second degré};}]
\vspace*{1.75cm}
\begin{minipage}[t][6.9cm][c]{\linewidth}\LARGE\centering
\begin{minipage}{0.95\linewidth}\raggedright\linespread{1.25}\selectfont Quelles sont les coordonnées du sommet de la parabole d'équation $y=3x^2-12x+7$~?\end{minipage}
\end{minipage}

\vfill
\begin{tcolorbox}[enhanced,nobeforeafter,width=7.1cm,height=2.1cm,valign=center,colback=cdeg!6,colframe=cdeg,arc=2.5mm,boxrule=0.9pt,left=3mm]{\Large\bfseries\color{cdeg}a)}\quad\begin{minipage}[c]{5.3cm}\raggedright\Large $(2\,;\,-5)$\end{minipage}\end{tcolorbox}\hfill\begin{tcolorbox}[enhanced,nobeforeafter,width=7.1cm,height=2.1cm,valign=center,colback=cdeg!6,colframe=cdeg,arc=2.5mm,boxrule=0.9pt,left=3mm]{\Large\bfseries\color{cdeg}b)}\quad\begin{minipage}[c]{5.3cm}\raggedright\Large $(2\,;\,5)$\end{minipage}\end{tcolorbox}\par\vspace{2.5mm}
\begin{tcolorbox}[enhanced,nobeforeafter,width=7.1cm,height=2.1cm,valign=center,colback=cdeg!6,colframe=cdeg,arc=2.5mm,boxrule=0.9pt,left=3mm]{\Large\bfseries\color{cdeg}c)}\quad\begin{minipage}[c]{5.3cm}\raggedright\Large $(2\,;\,7)$\end{minipage}\end{tcolorbox}\hfill\begin{tcolorbox}[enhanced,nobeforeafter,width=7.1cm,height=2.1cm,valign=center,colback=cdeg!6,colframe=cdeg,arc=2.5mm,boxrule=0.9pt,left=3mm]{\Large\bfseries\color{cdeg}d)}\quad\begin{minipage}[c]{5.3cm}\raggedright\Large $(-2\,;\,-5)$\end{minipage}\end{tcolorbox}
\end{tcolorbox}
\newpage
\begin{tcolorbox}[enhanced,width=16cm,height=14cm,colback=white,colframe=cdeg,boxrule=2pt,arc=0mm,
left=6mm,right=6mm,top=0mm,bottom=4mm,sharp corners,
overlay={\fill[cdeg] (frame.north west) rectangle ([yshift=-1.25cm]frame.north east);
\node[anchor=west,text=white,font=\sffamily\bfseries\LARGE] at ([xshift=5mm,yshift=-0.62cm]frame.north west) {Question 26 C};
\node[anchor=east,text=white,font=\sffamily\itshape\large] at ([xshift=-5mm,yshift=-0.62cm]frame.north east) {Second degré};}]
\vspace*{1.75cm}
\begin{minipage}[t][6.9cm][c]{\linewidth}\LARGE\centering
\begin{minipage}{0.95\linewidth}\raggedright\linespread{1.25}\selectfont Quelles sont les coordonnées du sommet de la parabole d'équation $y=-2x^2+4x+1$~?\end{minipage}
\end{minipage}

\vfill
\begin{tcolorbox}[enhanced,nobeforeafter,width=7.1cm,height=2.1cm,valign=center,colback=cdeg!6,colframe=cdeg,arc=2.5mm,boxrule=0.9pt,left=3mm]{\Large\bfseries\color{cdeg}a)}\quad\begin{minipage}[c]{5.3cm}\raggedright\Large $(1\,;\,-3)$\end{minipage}\end{tcolorbox}\hfill\begin{tcolorbox}[enhanced,nobeforeafter,width=7.1cm,height=2.1cm,valign=center,colback=cdeg!6,colframe=cdeg,arc=2.5mm,boxrule=0.9pt,left=3mm]{\Large\bfseries\color{cdeg}b)}\quad\begin{minipage}[c]{5.3cm}\raggedright\Large $(1\,;\,1)$\end{minipage}\end{tcolorbox}\par\vspace{2.5mm}
\begin{tcolorbox}[enhanced,nobeforeafter,width=7.1cm,height=2.1cm,valign=center,colback=cdeg!6,colframe=cdeg,arc=2.5mm,boxrule=0.9pt,left=3mm]{\Large\bfseries\color{cdeg}c)}\quad\begin{minipage}[c]{5.3cm}\raggedright\Large $(-1\,;\,3)$\end{minipage}\end{tcolorbox}\hfill\begin{tcolorbox}[enhanced,nobeforeafter,width=7.1cm,height=2.1cm,valign=center,colback=cdeg!6,colframe=cdeg,arc=2.5mm,boxrule=0.9pt,left=3mm]{\Large\bfseries\color{cdeg}d)}\quad\begin{minipage}[c]{5.3cm}\raggedright\Large $(1\,;\,3)$\end{minipage}\end{tcolorbox}
\end{tcolorbox}
\newpage
\begin{tcolorbox}[enhanced,width=16cm,height=14cm,colback=white,colframe=cdeg,boxrule=2pt,arc=0mm,
left=6mm,right=6mm,top=0mm,bottom=4mm,sharp corners,
overlay={\fill[cdeg] (frame.north west) rectangle ([yshift=-1.25cm]frame.north east);
\node[anchor=west,text=white,font=\sffamily\bfseries\LARGE] at ([xshift=5mm,yshift=-0.62cm]frame.north west) {Question 27 C};
\node[anchor=east,text=white,font=\sffamily\itshape\large] at ([xshift=-5mm,yshift=-0.62cm]frame.north east) {Second degré};}]
\vspace*{1.75cm}
\begin{minipage}[t][6.9cm][c]{\linewidth}\LARGE\centering
\begin{minipage}{0.95\linewidth}\raggedright\linespread{1.25}\selectfont Quel est l'ensemble des solutions de l'inéquation $x^2-3x-4\geqslant 0$~?\end{minipage}
\end{minipage}

\vfill
\begin{tcolorbox}[enhanced,nobeforeafter,width=7.1cm,height=2.1cm,valign=center,colback=cdeg!6,colframe=cdeg,arc=2.5mm,boxrule=0.9pt,left=3mm]{\Large\bfseries\color{cdeg}a)}\quad\begin{minipage}[c]{5.3cm}\raggedright\Large $[-1\,;\,4]$\end{minipage}\end{tcolorbox}\hfill\begin{tcolorbox}[enhanced,nobeforeafter,width=7.1cm,height=2.1cm,valign=center,colback=cdeg!6,colframe=cdeg,arc=2.5mm,boxrule=0.9pt,left=3mm]{\Large\bfseries\color{cdeg}b)}\quad\begin{minipage}[c]{5.3cm}\raggedright\Large $]-1\,;\,4[$\end{minipage}\end{tcolorbox}\par\vspace{2.5mm}
\begin{tcolorbox}[enhanced,nobeforeafter,width=7.1cm,height=2.1cm,valign=center,colback=cdeg!6,colframe=cdeg,arc=2.5mm,boxrule=0.9pt,left=3mm]{\Large\bfseries\color{cdeg}c)}\quad\begin{minipage}[c]{5.3cm}\raggedright\large $]-\infty\,;\,-1]\cup[4\,;\,+\infty[$\end{minipage}\end{tcolorbox}\hfill\begin{tcolorbox}[enhanced,nobeforeafter,width=7.1cm,height=2.1cm,valign=center,colback=cdeg!6,colframe=cdeg,arc=2.5mm,boxrule=0.9pt,left=3mm]{\Large\bfseries\color{cdeg}d)}\quad\begin{minipage}[c]{5.3cm}\raggedright\large $]-\infty\,;\,-1[\cup]4\,;\,+\infty[$\end{minipage}\end{tcolorbox}
\end{tcolorbox}
\newpage
\begin{tcolorbox}[enhanced,width=16cm,height=14cm,colback=white,colframe=cdeg,boxrule=2pt,arc=0mm,
left=6mm,right=6mm,top=0mm,bottom=4mm,sharp corners,
overlay={\fill[cdeg] (frame.north west) rectangle ([yshift=-1.25cm]frame.north east);
\node[anchor=west,text=white,font=\sffamily\bfseries\LARGE] at ([xshift=5mm,yshift=-0.62cm]frame.north west) {Question 28 C};
\node[anchor=east,text=white,font=\sffamily\itshape\large] at ([xshift=-5mm,yshift=-0.62cm]frame.north east) {Second degré};}]
\vspace*{1.75cm}
\begin{minipage}[t][6.9cm][c]{\linewidth}\LARGE\centering
\begin{minipage}{0.95\linewidth}\raggedright\linespread{1.25}\selectfont Quel est l'ensemble des solutions de l'inéquation $-2x^2-12x-10<0$~?\end{minipage}
\end{minipage}

\vfill
\begin{tcolorbox}[enhanced,nobeforeafter,width=7.1cm,height=2.1cm,valign=center,colback=cdeg!6,colframe=cdeg,arc=2.5mm,boxrule=0.9pt,left=3mm]{\Large\bfseries\color{cdeg}a)}\quad\begin{minipage}[c]{5.3cm}\raggedright\Large $[-5\,;\,-1]$\end{minipage}\end{tcolorbox}\hfill\begin{tcolorbox}[enhanced,nobeforeafter,width=7.1cm,height=2.1cm,valign=center,colback=cdeg!6,colframe=cdeg,arc=2.5mm,boxrule=0.9pt,left=3mm]{\Large\bfseries\color{cdeg}b)}\quad\begin{minipage}[c]{5.3cm}\raggedright\Large $]-5\,;\,-1[$\end{minipage}\end{tcolorbox}\par\vspace{2.5mm}
\begin{tcolorbox}[enhanced,nobeforeafter,width=7.1cm,height=2.1cm,valign=center,colback=cdeg!6,colframe=cdeg,arc=2.5mm,boxrule=0.9pt,left=3mm]{\Large\bfseries\color{cdeg}c)}\quad\begin{minipage}[c]{5.3cm}\raggedright\large $]-\infty\,;\,-5[\cup]-1\,;\,+\infty[$\end{minipage}\end{tcolorbox}\hfill\begin{tcolorbox}[enhanced,nobeforeafter,width=7.1cm,height=2.1cm,valign=center,colback=cdeg!6,colframe=cdeg,arc=2.5mm,boxrule=0.9pt,left=3mm]{\Large\bfseries\color{cdeg}d)}\quad\begin{minipage}[c]{5.3cm}\raggedright\large $]-\infty\,;\,-5]\cup[-1\,;\,+\infty[$\end{minipage}\end{tcolorbox}
\end{tcolorbox}
\newpage
\begin{tcolorbox}[enhanced,width=16cm,height=14cm,colback=white,colframe=cdeg,boxrule=2pt,arc=0mm,
left=6mm,right=6mm,top=0mm,bottom=4mm,sharp corners,
overlay={\fill[cdeg] (frame.north west) rectangle ([yshift=-1.25cm]frame.north east);
\node[anchor=west,text=white,font=\sffamily\bfseries\LARGE] at ([xshift=5mm,yshift=-0.62cm]frame.north west) {Question 29 C};
\node[anchor=east,text=white,font=\sffamily\itshape\large] at ([xshift=-5mm,yshift=-0.62cm]frame.north east) {Second degré};}]
\vspace*{1.75cm}
\begin{minipage}[t][6.9cm][c]{\linewidth}\LARGE\centering
\begin{minipage}{0.95\linewidth}\raggedright\linespread{1.25}\selectfont Le trinôme $x^2-x-12$ a deux racines. Que vaut leur produit~?\end{minipage}
\end{minipage}

\vfill
\begin{tcolorbox}[enhanced,nobeforeafter,width=7.1cm,height=2.1cm,valign=center,colback=cdeg!6,colframe=cdeg,arc=2.5mm,boxrule=0.9pt,left=3mm]{\Large\bfseries\color{cdeg}a)}\quad\begin{minipage}[c]{5.3cm}\raggedright\Large $1$\end{minipage}\end{tcolorbox}\hfill\begin{tcolorbox}[enhanced,nobeforeafter,width=7.1cm,height=2.1cm,valign=center,colback=cdeg!6,colframe=cdeg,arc=2.5mm,boxrule=0.9pt,left=3mm]{\Large\bfseries\color{cdeg}b)}\quad\begin{minipage}[c]{5.3cm}\raggedright\Large $-12$\end{minipage}\end{tcolorbox}\par\vspace{2.5mm}
\begin{tcolorbox}[enhanced,nobeforeafter,width=7.1cm,height=2.1cm,valign=center,colback=cdeg!6,colframe=cdeg,arc=2.5mm,boxrule=0.9pt,left=3mm]{\Large\bfseries\color{cdeg}c)}\quad\begin{minipage}[c]{5.3cm}\raggedright\Large $12$\end{minipage}\end{tcolorbox}\hfill\begin{tcolorbox}[enhanced,nobeforeafter,width=7.1cm,height=2.1cm,valign=center,colback=cdeg!6,colframe=cdeg,arc=2.5mm,boxrule=0.9pt,left=3mm]{\Large\bfseries\color{cdeg}d)}\quad\begin{minipage}[c]{5.3cm}\raggedright\Large $-1$\end{minipage}\end{tcolorbox}
\end{tcolorbox}
\newpage
\begin{tcolorbox}[enhanced,width=16cm,height=14cm,colback=white,colframe=cdeg,boxrule=2pt,arc=0mm,
left=6mm,right=6mm,top=0mm,bottom=4mm,sharp corners,
overlay={\fill[cdeg] (frame.north west) rectangle ([yshift=-1.25cm]frame.north east);
\node[anchor=west,text=white,font=\sffamily\bfseries\LARGE] at ([xshift=5mm,yshift=-0.62cm]frame.north west) {Question 30 C};
\node[anchor=east,text=white,font=\sffamily\itshape\large] at ([xshift=-5mm,yshift=-0.62cm]frame.north east) {Second degré};}]
\vspace*{1.75cm}
\begin{minipage}[t][6.9cm][c]{\linewidth}\LARGE\centering
\begin{minipage}{0.95\linewidth}\raggedright\linespread{1.25}\selectfont Le trinôme $2x^2+5x-3$ a deux racines. Que vaut leur somme~?\end{minipage}
\end{minipage}

\vfill
\begin{tcolorbox}[enhanced,nobeforeafter,width=7.1cm,height=2.1cm,valign=center,colback=cdeg!6,colframe=cdeg,arc=2.5mm,boxrule=0.9pt,left=3mm]{\Large\bfseries\color{cdeg}a)}\quad\begin{minipage}[c]{5.3cm}\raggedright\Large $\dfrac{5}{2}$\end{minipage}\end{tcolorbox}\hfill\begin{tcolorbox}[enhanced,nobeforeafter,width=7.1cm,height=2.1cm,valign=center,colback=cdeg!6,colframe=cdeg,arc=2.5mm,boxrule=0.9pt,left=3mm]{\Large\bfseries\color{cdeg}b)}\quad\begin{minipage}[c]{5.3cm}\raggedright\Large $-\dfrac{3}{2}$\end{minipage}\end{tcolorbox}\par\vspace{2.5mm}
\begin{tcolorbox}[enhanced,nobeforeafter,width=7.1cm,height=2.1cm,valign=center,colback=cdeg!6,colframe=cdeg,arc=2.5mm,boxrule=0.9pt,left=3mm]{\Large\bfseries\color{cdeg}c)}\quad\begin{minipage}[c]{5.3cm}\raggedright\Large $-\dfrac{5}{2}$\end{minipage}\end{tcolorbox}\hfill\begin{tcolorbox}[enhanced,nobeforeafter,width=7.1cm,height=2.1cm,valign=center,colback=cdeg!6,colframe=cdeg,arc=2.5mm,boxrule=0.9pt,left=3mm]{\Large\bfseries\color{cdeg}d)}\quad\begin{minipage}[c]{5.3cm}\raggedright\Large $-5$\end{minipage}\end{tcolorbox}
\end{tcolorbox}
\newpage
\begin{tcolorbox}[enhanced,width=16cm,height=14cm,colback=white,colframe=cdeg,boxrule=2pt,arc=0mm,
left=6mm,right=6mm,top=0mm,bottom=4mm,sharp corners,
overlay={\fill[cdeg] (frame.north west) rectangle ([yshift=-1.25cm]frame.north east);
\node[anchor=west,text=white,font=\sffamily\bfseries\LARGE] at ([xshift=5mm,yshift=-0.62cm]frame.north west) {Question 31 C};
\node[anchor=east,text=white,font=\sffamily\itshape\large] at ([xshift=-5mm,yshift=-0.62cm]frame.north east) {Second degré};}]
\vspace*{1.75cm}
\begin{minipage}[t][6.9cm][c]{\linewidth}\LARGE\centering
\begin{minipage}{0.95\linewidth}\raggedright\linespread{1.25}\selectfont Le trinôme $x^2+9x+18$ a pour racines $-3$ et $-6$. Quelle est sa forme factorisée~?\end{minipage}
\end{minipage}

\vfill
\begin{tcolorbox}[enhanced,nobeforeafter,width=7.1cm,height=2.1cm,valign=center,colback=cdeg!6,colframe=cdeg,arc=2.5mm,boxrule=0.9pt,left=3mm]{\Large\bfseries\color{cdeg}a)}\quad\begin{minipage}[c]{5.3cm}\raggedright\Large $(x-3)(x-6)$\end{minipage}\end{tcolorbox}\hfill\begin{tcolorbox}[enhanced,nobeforeafter,width=7.1cm,height=2.1cm,valign=center,colback=cdeg!6,colframe=cdeg,arc=2.5mm,boxrule=0.9pt,left=3mm]{\Large\bfseries\color{cdeg}b)}\quad\begin{minipage}[c]{5.3cm}\raggedright\Large $-(x+3)(x+6)$\end{minipage}\end{tcolorbox}\par\vspace{2.5mm}
\begin{tcolorbox}[enhanced,nobeforeafter,width=7.1cm,height=2.1cm,valign=center,colback=cdeg!6,colframe=cdeg,arc=2.5mm,boxrule=0.9pt,left=3mm]{\Large\bfseries\color{cdeg}c)}\quad\begin{minipage}[c]{5.3cm}\raggedright\Large $(x+3)(x+6)$\end{minipage}\end{tcolorbox}\hfill\begin{tcolorbox}[enhanced,nobeforeafter,width=7.1cm,height=2.1cm,valign=center,colback=cdeg!6,colframe=cdeg,arc=2.5mm,boxrule=0.9pt,left=3mm]{\Large\bfseries\color{cdeg}d)}\quad\begin{minipage}[c]{5.3cm}\raggedright\Large $2(x+3)(x+6)$\end{minipage}\end{tcolorbox}
\end{tcolorbox}
\newpage
\begin{tcolorbox}[enhanced,width=16cm,height=14cm,colback=white,colframe=cdeg,boxrule=2pt,arc=0mm,
left=6mm,right=6mm,top=0mm,bottom=4mm,sharp corners,
overlay={\fill[cdeg] (frame.north west) rectangle ([yshift=-1.25cm]frame.north east);
\node[anchor=west,text=white,font=\sffamily\bfseries\LARGE] at ([xshift=5mm,yshift=-0.62cm]frame.north west) {Question 32 C};
\node[anchor=east,text=white,font=\sffamily\itshape\large] at ([xshift=-5mm,yshift=-0.62cm]frame.north east) {Second degré};}]
\vspace*{1.75cm}
\begin{minipage}[t][6.9cm][c]{\linewidth}\LARGE\centering
\begin{minipage}{0.95\linewidth}\raggedright\linespread{1.25}\selectfont Le trinôme $2x^2-8x+6$ a pour racines $1$ et $3$. Quelle est sa forme factorisée~?\end{minipage}
\end{minipage}

\vfill
\begin{tcolorbox}[enhanced,nobeforeafter,width=7.1cm,height=2.1cm,valign=center,colback=cdeg!6,colframe=cdeg,arc=2.5mm,boxrule=0.9pt,left=3mm]{\Large\bfseries\color{cdeg}a)}\quad\begin{minipage}[c]{5.3cm}\raggedright\Large $(x-1)(x-3)$\end{minipage}\end{tcolorbox}\hfill\begin{tcolorbox}[enhanced,nobeforeafter,width=7.1cm,height=2.1cm,valign=center,colback=cdeg!6,colframe=cdeg,arc=2.5mm,boxrule=0.9pt,left=3mm]{\Large\bfseries\color{cdeg}b)}\quad\begin{minipage}[c]{5.3cm}\raggedright\Large $2(x+1)(x+3)$\end{minipage}\end{tcolorbox}\par\vspace{2.5mm}
\begin{tcolorbox}[enhanced,nobeforeafter,width=7.1cm,height=2.1cm,valign=center,colback=cdeg!6,colframe=cdeg,arc=2.5mm,boxrule=0.9pt,left=3mm]{\Large\bfseries\color{cdeg}c)}\quad\begin{minipage}[c]{5.3cm}\raggedright\Large $-(x-1)(x-3)$\end{minipage}\end{tcolorbox}\hfill\begin{tcolorbox}[enhanced,nobeforeafter,width=7.1cm,height=2.1cm,valign=center,colback=cdeg!6,colframe=cdeg,arc=2.5mm,boxrule=0.9pt,left=3mm]{\Large\bfseries\color{cdeg}d)}\quad\begin{minipage}[c]{5.3cm}\raggedright\Large $2(x-1)(x-3)$\end{minipage}\end{tcolorbox}
\end{tcolorbox}
\newpage
\begin{tcolorbox}[enhanced,width=16cm,height=14cm,colback=white,colframe=cder,boxrule=2pt,arc=0mm,
left=6mm,right=6mm,top=0mm,bottom=4mm,sharp corners,
overlay={\fill[cder] (frame.north west) rectangle ([yshift=-1.25cm]frame.north east);
\node[anchor=west,text=white,font=\sffamily\bfseries\LARGE] at ([xshift=5mm,yshift=-0.62cm]frame.north west) {Question 33 C};
\node[anchor=east,text=white,font=\sffamily\itshape\large] at ([xshift=-5mm,yshift=-0.62cm]frame.north east) {Dérivation};}]
\vspace*{1.75cm}
\begin{minipage}[t][6.9cm][c]{\linewidth}\LARGE\centering
\begin{minipage}{0.95\linewidth}\raggedright\linespread{1.25}\selectfont Soit $f(x)=3x^3-4x^2+2x$. Que vaut $f'(0)$~?\end{minipage}
\end{minipage}

\vfill
\begin{tcolorbox}[enhanced,nobeforeafter,width=7.1cm,height=2.1cm,valign=center,colback=cder!6,colframe=cder,arc=2.5mm,boxrule=0.9pt,left=3mm]{\Large\bfseries\color{cder}a)}\quad\begin{minipage}[c]{5.3cm}\raggedright\Large $3$\end{minipage}\end{tcolorbox}\hfill\begin{tcolorbox}[enhanced,nobeforeafter,width=7.1cm,height=2.1cm,valign=center,colback=cder!6,colframe=cder,arc=2.5mm,boxrule=0.9pt,left=3mm]{\Large\bfseries\color{cder}b)}\quad\begin{minipage}[c]{5.3cm}\raggedright\Large $1$\end{minipage}\end{tcolorbox}\par\vspace{2.5mm}
\begin{tcolorbox}[enhanced,nobeforeafter,width=7.1cm,height=2.1cm,valign=center,colback=cder!6,colframe=cder,arc=2.5mm,boxrule=0.9pt,left=3mm]{\Large\bfseries\color{cder}c)}\quad\begin{minipage}[c]{5.3cm}\raggedright\Large $2$\end{minipage}\end{tcolorbox}\hfill\begin{tcolorbox}[enhanced,nobeforeafter,width=7.1cm,height=2.1cm,valign=center,colback=cder!6,colframe=cder,arc=2.5mm,boxrule=0.9pt,left=3mm]{\Large\bfseries\color{cder}d)}\quad\begin{minipage}[c]{5.3cm}\raggedright\Large $0$\end{minipage}\end{tcolorbox}
\end{tcolorbox}
\newpage
\begin{tcolorbox}[enhanced,width=16cm,height=14cm,colback=white,colframe=cder,boxrule=2pt,arc=0mm,
left=6mm,right=6mm,top=0mm,bottom=4mm,sharp corners,
overlay={\fill[cder] (frame.north west) rectangle ([yshift=-1.25cm]frame.north east);
\node[anchor=west,text=white,font=\sffamily\bfseries\LARGE] at ([xshift=5mm,yshift=-0.62cm]frame.north west) {Question 34 C};
\node[anchor=east,text=white,font=\sffamily\itshape\large] at ([xshift=-5mm,yshift=-0.62cm]frame.north east) {Dérivation};}]
\vspace*{1.75cm}
\begin{minipage}[t][6.9cm][c]{\linewidth}\LARGE\centering
\begin{minipage}{0.95\linewidth}\raggedright\linespread{1.25}\selectfont Soit $f(x)=x^3-12x+4$. Que vaut $f'(2)$~?\end{minipage}
\end{minipage}

\vfill
\begin{tcolorbox}[enhanced,nobeforeafter,width=7.1cm,height=2.1cm,valign=center,colback=cder!6,colframe=cder,arc=2.5mm,boxrule=0.9pt,left=3mm]{\Large\bfseries\color{cder}a)}\quad\begin{minipage}[c]{5.3cm}\raggedright\Large $0$\end{minipage}\end{tcolorbox}\hfill\begin{tcolorbox}[enhanced,nobeforeafter,width=7.1cm,height=2.1cm,valign=center,colback=cder!6,colframe=cder,arc=2.5mm,boxrule=0.9pt,left=3mm]{\Large\bfseries\color{cder}b)}\quad\begin{minipage}[c]{5.3cm}\raggedright\Large $-12$\end{minipage}\end{tcolorbox}\par\vspace{2.5mm}
\begin{tcolorbox}[enhanced,nobeforeafter,width=7.1cm,height=2.1cm,valign=center,colback=cder!6,colframe=cder,arc=2.5mm,boxrule=0.9pt,left=3mm]{\Large\bfseries\color{cder}c)}\quad\begin{minipage}[c]{5.3cm}\raggedright\Large $1$\end{minipage}\end{tcolorbox}\hfill\begin{tcolorbox}[enhanced,nobeforeafter,width=7.1cm,height=2.1cm,valign=center,colback=cder!6,colframe=cder,arc=2.5mm,boxrule=0.9pt,left=3mm]{\Large\bfseries\color{cder}d)}\quad\begin{minipage}[c]{5.3cm}\raggedright\Large $-8$\end{minipage}\end{tcolorbox}
\end{tcolorbox}
\newpage
\begin{tcolorbox}[enhanced,width=16cm,height=14cm,colback=white,colframe=cder,boxrule=2pt,arc=0mm,
left=6mm,right=6mm,top=0mm,bottom=4mm,sharp corners,
overlay={\fill[cder] (frame.north west) rectangle ([yshift=-1.25cm]frame.north east);
\node[anchor=west,text=white,font=\sffamily\bfseries\LARGE] at ([xshift=5mm,yshift=-0.62cm]frame.north west) {Question 35 C};
\node[anchor=east,text=white,font=\sffamily\itshape\large] at ([xshift=-5mm,yshift=-0.62cm]frame.north east) {Dérivation};}]
\vspace*{1.75cm}
\begin{minipage}[t][6.9cm][c]{\linewidth}\LARGE\centering
\begin{minipage}{0.95\linewidth}\raggedright\linespread{1.25}\selectfont Quelle est l'équation de la tangente à la courbe de $f(x)=3x^2-6x+2$ au point d'abscisse $2$~?\end{minipage}
\end{minipage}

\vfill
\begin{tcolorbox}[enhanced,nobeforeafter,width=7.1cm,height=2.1cm,valign=center,colback=cder!6,colframe=cder,arc=2.5mm,boxrule=0.9pt,left=3mm]{\Large\bfseries\color{cder}a)}\quad\begin{minipage}[c]{5.3cm}\raggedright\Large $y=6x+14$\end{minipage}\end{tcolorbox}\hfill\begin{tcolorbox}[enhanced,nobeforeafter,width=7.1cm,height=2.1cm,valign=center,colback=cder!6,colframe=cder,arc=2.5mm,boxrule=0.9pt,left=3mm]{\Large\bfseries\color{cder}b)}\quad\begin{minipage}[c]{5.3cm}\raggedright\Large $y=6x+2$\end{minipage}\end{tcolorbox}\par\vspace{2.5mm}
\begin{tcolorbox}[enhanced,nobeforeafter,width=7.1cm,height=2.1cm,valign=center,colback=cder!6,colframe=cder,arc=2.5mm,boxrule=0.9pt,left=3mm]{\Large\bfseries\color{cder}c)}\quad\begin{minipage}[c]{5.3cm}\raggedright\Large $y=2x+2$\end{minipage}\end{tcolorbox}\hfill\begin{tcolorbox}[enhanced,nobeforeafter,width=7.1cm,height=2.1cm,valign=center,colback=cder!6,colframe=cder,arc=2.5mm,boxrule=0.9pt,left=3mm]{\Large\bfseries\color{cder}d)}\quad\begin{minipage}[c]{5.3cm}\raggedright\Large $y=6x-10$\end{minipage}\end{tcolorbox}
\end{tcolorbox}
\newpage
\begin{tcolorbox}[enhanced,width=16cm,height=14cm,colback=white,colframe=cder,boxrule=2pt,arc=0mm,
left=6mm,right=6mm,top=0mm,bottom=4mm,sharp corners,
overlay={\fill[cder] (frame.north west) rectangle ([yshift=-1.25cm]frame.north east);
\node[anchor=west,text=white,font=\sffamily\bfseries\LARGE] at ([xshift=5mm,yshift=-0.62cm]frame.north west) {Question 36 C};
\node[anchor=east,text=white,font=\sffamily\itshape\large] at ([xshift=-5mm,yshift=-0.62cm]frame.north east) {Dérivation};}]
\vspace*{1.75cm}
\begin{minipage}[t][6.9cm][c]{\linewidth}\LARGE\centering
\begin{minipage}{0.95\linewidth}\raggedright\linespread{1.25}\selectfont Quelle est l'équation de la tangente à la courbe de $f(x)=x^2+3x-1$ au point d'abscisse $0$~?\end{minipage}
\end{minipage}

\vfill
\begin{tcolorbox}[enhanced,nobeforeafter,width=7.1cm,height=2.1cm,valign=center,colback=cder!6,colframe=cder,arc=2.5mm,boxrule=0.9pt,left=3mm]{\Large\bfseries\color{cder}a)}\quad\begin{minipage}[c]{5.3cm}\raggedright\Large $y=3x-1$\end{minipage}\end{tcolorbox}\hfill\begin{tcolorbox}[enhanced,nobeforeafter,width=7.1cm,height=2.1cm,valign=center,colback=cder!6,colframe=cder,arc=2.5mm,boxrule=0.9pt,left=3mm]{\Large\bfseries\color{cder}b)}\quad\begin{minipage}[c]{5.3cm}\raggedright\Large $y=-x$\end{minipage}\end{tcolorbox}\par\vspace{2.5mm}
\begin{tcolorbox}[enhanced,nobeforeafter,width=7.1cm,height=2.1cm,valign=center,colback=cder!6,colframe=cder,arc=2.5mm,boxrule=0.9pt,left=3mm]{\Large\bfseries\color{cder}c)}\quad\begin{minipage}[c]{5.3cm}\raggedright\Large $y=3x$\end{minipage}\end{tcolorbox}\hfill\begin{tcolorbox}[enhanced,nobeforeafter,width=7.1cm,height=2.1cm,valign=center,colback=cder!6,colframe=cder,arc=2.5mm,boxrule=0.9pt,left=3mm]{\Large\bfseries\color{cder}d)}\quad\begin{minipage}[c]{5.3cm}\raggedright\Large $y=-x+3$\end{minipage}\end{tcolorbox}
\end{tcolorbox}
\newpage
\begin{tcolorbox}[enhanced,width=16cm,height=14cm,colback=white,colframe=cder,boxrule=2pt,arc=0mm,
left=6mm,right=6mm,top=0mm,bottom=4mm,sharp corners,
overlay={\fill[cder] (frame.north west) rectangle ([yshift=-1.25cm]frame.north east);
\node[anchor=west,text=white,font=\sffamily\bfseries\LARGE] at ([xshift=5mm,yshift=-0.62cm]frame.north west) {Question 37 C};
\node[anchor=east,text=white,font=\sffamily\itshape\large] at ([xshift=-5mm,yshift=-0.62cm]frame.north east) {Dérivation};}]
\vspace*{1.75cm}
\begin{minipage}[t][6.9cm][c]{\linewidth}\LARGE\centering
\begin{minipage}{0.95\linewidth}\raggedright\linespread{1.25}\selectfont Soit $f(x)=(2x-1)(-x+3)$. Que vaut $f'(-1)$~?\end{minipage}
\end{minipage}

\vfill
\begin{tcolorbox}[enhanced,nobeforeafter,width=7.1cm,height=2.1cm,valign=center,colback=cder!6,colframe=cder,arc=2.5mm,boxrule=0.9pt,left=3mm]{\Large\bfseries\color{cder}a)}\quad\begin{minipage}[c]{5.3cm}\raggedright\Large $11$\end{minipage}\end{tcolorbox}\hfill\begin{tcolorbox}[enhanced,nobeforeafter,width=7.1cm,height=2.1cm,valign=center,colback=cder!6,colframe=cder,arc=2.5mm,boxrule=0.9pt,left=3mm]{\Large\bfseries\color{cder}b)}\quad\begin{minipage}[c]{5.3cm}\raggedright\Large $-12$\end{minipage}\end{tcolorbox}\par\vspace{2.5mm}
\begin{tcolorbox}[enhanced,nobeforeafter,width=7.1cm,height=2.1cm,valign=center,colback=cder!6,colframe=cder,arc=2.5mm,boxrule=0.9pt,left=3mm]{\Large\bfseries\color{cder}c)}\quad\begin{minipage}[c]{5.3cm}\raggedright\Large $-2$\end{minipage}\end{tcolorbox}\hfill\begin{tcolorbox}[enhanced,nobeforeafter,width=7.1cm,height=2.1cm,valign=center,colback=cder!6,colframe=cder,arc=2.5mm,boxrule=0.9pt,left=3mm]{\Large\bfseries\color{cder}d)}\quad\begin{minipage}[c]{5.3cm}\raggedright\Large $8$\end{minipage}\end{tcolorbox}
\end{tcolorbox}
\newpage
\begin{tcolorbox}[enhanced,width=16cm,height=14cm,colback=white,colframe=cder,boxrule=2pt,arc=0mm,
left=6mm,right=6mm,top=0mm,bottom=4mm,sharp corners,
overlay={\fill[cder] (frame.north west) rectangle ([yshift=-1.25cm]frame.north east);
\node[anchor=west,text=white,font=\sffamily\bfseries\LARGE] at ([xshift=5mm,yshift=-0.62cm]frame.north west) {Question 38 C};
\node[anchor=east,text=white,font=\sffamily\itshape\large] at ([xshift=-5mm,yshift=-0.62cm]frame.north east) {Dérivation};}]
\vspace*{1.75cm}
\begin{minipage}[t][6.9cm][c]{\linewidth}\LARGE\centering
\begin{minipage}{0.95\linewidth}\raggedright\linespread{1.25}\selectfont Soit $f(x)=(x+4)(x+4)$. Que vaut $f'(-2)$~?\end{minipage}
\end{minipage}

\vfill
\begin{tcolorbox}[enhanced,nobeforeafter,width=7.1cm,height=2.1cm,valign=center,colback=cder!6,colframe=cder,arc=2.5mm,boxrule=0.9pt,left=3mm]{\Large\bfseries\color{cder}a)}\quad\begin{minipage}[c]{5.3cm}\raggedright\Large $2$\end{minipage}\end{tcolorbox}\hfill\begin{tcolorbox}[enhanced,nobeforeafter,width=7.1cm,height=2.1cm,valign=center,colback=cder!6,colframe=cder,arc=2.5mm,boxrule=0.9pt,left=3mm]{\Large\bfseries\color{cder}b)}\quad\begin{minipage}[c]{5.3cm}\raggedright\Large $6$\end{minipage}\end{tcolorbox}\par\vspace{2.5mm}
\begin{tcolorbox}[enhanced,nobeforeafter,width=7.1cm,height=2.1cm,valign=center,colback=cder!6,colframe=cder,arc=2.5mm,boxrule=0.9pt,left=3mm]{\Large\bfseries\color{cder}c)}\quad\begin{minipage}[c]{5.3cm}\raggedright\Large $1$\end{minipage}\end{tcolorbox}\hfill\begin{tcolorbox}[enhanced,nobeforeafter,width=7.1cm,height=2.1cm,valign=center,colback=cder!6,colframe=cder,arc=2.5mm,boxrule=0.9pt,left=3mm]{\Large\bfseries\color{cder}d)}\quad\begin{minipage}[c]{5.3cm}\raggedright\Large $4$\end{minipage}\end{tcolorbox}
\end{tcolorbox}
\newpage
\begin{tcolorbox}[enhanced,width=16cm,height=14cm,colback=white,colframe=cder,boxrule=2pt,arc=0mm,
left=6mm,right=6mm,top=0mm,bottom=4mm,sharp corners,
overlay={\fill[cder] (frame.north west) rectangle ([yshift=-1.25cm]frame.north east);
\node[anchor=west,text=white,font=\sffamily\bfseries\LARGE] at ([xshift=5mm,yshift=-0.62cm]frame.north west) {Question 39 C};
\node[anchor=east,text=white,font=\sffamily\itshape\large] at ([xshift=-5mm,yshift=-0.62cm]frame.north east) {Dérivation};}]
\vspace*{1.75cm}
\begin{minipage}[t][6.9cm][c]{\linewidth}\LARGE\centering
\begin{minipage}{0.95\linewidth}\raggedright\linespread{1.25}\selectfont Soit $f(x)=\dfrac{6}{x}$ sur $\mathbb{R}^*$. Que vaut $f'(3)$~?\end{minipage}
\end{minipage}

\vfill
\begin{tcolorbox}[enhanced,nobeforeafter,width=7.1cm,height=2.1cm,valign=center,colback=cder!6,colframe=cder,arc=2.5mm,boxrule=0.9pt,left=3mm]{\Large\bfseries\color{cder}a)}\quad\begin{minipage}[c]{5.3cm}\raggedright\Large $\dfrac{2}{3}$\end{minipage}\end{tcolorbox}\hfill\begin{tcolorbox}[enhanced,nobeforeafter,width=7.1cm,height=2.1cm,valign=center,colback=cder!6,colframe=cder,arc=2.5mm,boxrule=0.9pt,left=3mm]{\Large\bfseries\color{cder}b)}\quad\begin{minipage}[c]{5.3cm}\raggedright\Large $18$\end{minipage}\end{tcolorbox}\par\vspace{2.5mm}
\begin{tcolorbox}[enhanced,nobeforeafter,width=7.1cm,height=2.1cm,valign=center,colback=cder!6,colframe=cder,arc=2.5mm,boxrule=0.9pt,left=3mm]{\Large\bfseries\color{cder}c)}\quad\begin{minipage}[c]{5.3cm}\raggedright\Large $-2$\end{minipage}\end{tcolorbox}\hfill\begin{tcolorbox}[enhanced,nobeforeafter,width=7.1cm,height=2.1cm,valign=center,colback=cder!6,colframe=cder,arc=2.5mm,boxrule=0.9pt,left=3mm]{\Large\bfseries\color{cder}d)}\quad\begin{minipage}[c]{5.3cm}\raggedright\Large $-\dfrac{2}{3}$\end{minipage}\end{tcolorbox}
\end{tcolorbox}
\newpage
\begin{tcolorbox}[enhanced,width=16cm,height=14cm,colback=white,colframe=cder,boxrule=2pt,arc=0mm,
left=6mm,right=6mm,top=0mm,bottom=4mm,sharp corners,
overlay={\fill[cder] (frame.north west) rectangle ([yshift=-1.25cm]frame.north east);
\node[anchor=west,text=white,font=\sffamily\bfseries\LARGE] at ([xshift=5mm,yshift=-0.62cm]frame.north west) {Question 40 C};
\node[anchor=east,text=white,font=\sffamily\itshape\large] at ([xshift=-5mm,yshift=-0.62cm]frame.north east) {Dérivation};}]
\vspace*{1.75cm}
\begin{minipage}[t][6.9cm][c]{\linewidth}\LARGE\centering
\begin{minipage}{0.95\linewidth}\raggedright\linespread{1.25}\selectfont Soit $f(x)=\dfrac{1}{x}$ sur $\mathbb{R}^*$. Que vaut $f'(\dfrac{1}{2})$~?\end{minipage}
\end{minipage}

\vfill
\begin{tcolorbox}[enhanced,nobeforeafter,width=7.1cm,height=2.1cm,valign=center,colback=cder!6,colframe=cder,arc=2.5mm,boxrule=0.9pt,left=3mm]{\Large\bfseries\color{cder}a)}\quad\begin{minipage}[c]{5.3cm}\raggedright\Large $-4$\end{minipage}\end{tcolorbox}\hfill\begin{tcolorbox}[enhanced,nobeforeafter,width=7.1cm,height=2.1cm,valign=center,colback=cder!6,colframe=cder,arc=2.5mm,boxrule=0.9pt,left=3mm]{\Large\bfseries\color{cder}b)}\quad\begin{minipage}[c]{5.3cm}\raggedright\Large $\dfrac{1}{2}$\end{minipage}\end{tcolorbox}\par\vspace{2.5mm}
\begin{tcolorbox}[enhanced,nobeforeafter,width=7.1cm,height=2.1cm,valign=center,colback=cder!6,colframe=cder,arc=2.5mm,boxrule=0.9pt,left=3mm]{\Large\bfseries\color{cder}c)}\quad\begin{minipage}[c]{5.3cm}\raggedright\Large $-2$\end{minipage}\end{tcolorbox}\hfill\begin{tcolorbox}[enhanced,nobeforeafter,width=7.1cm,height=2.1cm,valign=center,colback=cder!6,colframe=cder,arc=2.5mm,boxrule=0.9pt,left=3mm]{\Large\bfseries\color{cder}d)}\quad\begin{minipage}[c]{5.3cm}\raggedright\Large $4$\end{minipage}\end{tcolorbox}
\end{tcolorbox}
\newpage
\begin{tcolorbox}[enhanced,width=16cm,height=14cm,colback=white,colframe=cder,boxrule=2pt,arc=0mm,
left=6mm,right=6mm,top=0mm,bottom=4mm,sharp corners,
overlay={\fill[cder] (frame.north west) rectangle ([yshift=-1.25cm]frame.north east);
\node[anchor=west,text=white,font=\sffamily\bfseries\LARGE] at ([xshift=5mm,yshift=-0.62cm]frame.north west) {Question 41 C};
\node[anchor=east,text=white,font=\sffamily\itshape\large] at ([xshift=-5mm,yshift=-0.62cm]frame.north east) {Dérivation};}]
\vspace*{1.75cm}
\begin{minipage}[t][6.9cm][c]{\linewidth}\LARGE\centering
\begin{minipage}{0.95\linewidth}\raggedright\linespread{1.25}\selectfont Soit $f(x)=\sqrt{x}$. Que vaut $f'(25)$~?\end{minipage}
\end{minipage}

\vfill
\begin{tcolorbox}[enhanced,nobeforeafter,width=7.1cm,height=2.1cm,valign=center,colback=cder!6,colframe=cder,arc=2.5mm,boxrule=0.9pt,left=3mm]{\Large\bfseries\color{cder}a)}\quad\begin{minipage}[c]{5.3cm}\raggedright\Large $\dfrac{5}{2}$\end{minipage}\end{tcolorbox}\hfill\begin{tcolorbox}[enhanced,nobeforeafter,width=7.1cm,height=2.1cm,valign=center,colback=cder!6,colframe=cder,arc=2.5mm,boxrule=0.9pt,left=3mm]{\Large\bfseries\color{cder}b)}\quad\begin{minipage}[c]{5.3cm}\raggedright\Large $10$\end{minipage}\end{tcolorbox}\par\vspace{2.5mm}
\begin{tcolorbox}[enhanced,nobeforeafter,width=7.1cm,height=2.1cm,valign=center,colback=cder!6,colframe=cder,arc=2.5mm,boxrule=0.9pt,left=3mm]{\Large\bfseries\color{cder}c)}\quad\begin{minipage}[c]{5.3cm}\raggedright\Large $\dfrac{1}{10}$\end{minipage}\end{tcolorbox}\hfill\begin{tcolorbox}[enhanced,nobeforeafter,width=7.1cm,height=2.1cm,valign=center,colback=cder!6,colframe=cder,arc=2.5mm,boxrule=0.9pt,left=3mm]{\Large\bfseries\color{cder}d)}\quad\begin{minipage}[c]{5.3cm}\raggedright\Large $\dfrac{1}{5}$\end{minipage}\end{tcolorbox}
\end{tcolorbox}
\newpage
\begin{tcolorbox}[enhanced,width=16cm,height=14cm,colback=white,colframe=cder,boxrule=2pt,arc=0mm,
left=6mm,right=6mm,top=0mm,bottom=4mm,sharp corners,
overlay={\fill[cder] (frame.north west) rectangle ([yshift=-1.25cm]frame.north east);
\node[anchor=west,text=white,font=\sffamily\bfseries\LARGE] at ([xshift=5mm,yshift=-0.62cm]frame.north west) {Question 42 C};
\node[anchor=east,text=white,font=\sffamily\itshape\large] at ([xshift=-5mm,yshift=-0.62cm]frame.north east) {Dérivation};}]
\vspace*{1.75cm}
\begin{minipage}[t][6.9cm][c]{\linewidth}\LARGE\centering
\begin{minipage}{0.95\linewidth}\raggedright\linespread{1.25}\selectfont Soit $f(x)=\sqrt{x}$. Que vaut $f'(36)$~?\end{minipage}
\end{minipage}

\vfill
\begin{tcolorbox}[enhanced,nobeforeafter,width=7.1cm,height=2.1cm,valign=center,colback=cder!6,colframe=cder,arc=2.5mm,boxrule=0.9pt,left=3mm]{\Large\bfseries\color{cder}a)}\quad\begin{minipage}[c]{5.3cm}\raggedright\Large $3$\end{minipage}\end{tcolorbox}\hfill\begin{tcolorbox}[enhanced,nobeforeafter,width=7.1cm,height=2.1cm,valign=center,colback=cder!6,colframe=cder,arc=2.5mm,boxrule=0.9pt,left=3mm]{\Large\bfseries\color{cder}b)}\quad\begin{minipage}[c]{5.3cm}\raggedright\Large $12$\end{minipage}\end{tcolorbox}\par\vspace{2.5mm}
\begin{tcolorbox}[enhanced,nobeforeafter,width=7.1cm,height=2.1cm,valign=center,colback=cder!6,colframe=cder,arc=2.5mm,boxrule=0.9pt,left=3mm]{\Large\bfseries\color{cder}c)}\quad\begin{minipage}[c]{5.3cm}\raggedright\Large $\dfrac{1}{6}$\end{minipage}\end{tcolorbox}\hfill\begin{tcolorbox}[enhanced,nobeforeafter,width=7.1cm,height=2.1cm,valign=center,colback=cder!6,colframe=cder,arc=2.5mm,boxrule=0.9pt,left=3mm]{\Large\bfseries\color{cder}d)}\quad\begin{minipage}[c]{5.3cm}\raggedright\Large $\dfrac{1}{12}$\end{minipage}\end{tcolorbox}
\end{tcolorbox}
\newpage
\begin{tcolorbox}[enhanced,width=16cm,height=14cm,colback=white,colframe=cder,boxrule=2pt,arc=0mm,
left=6mm,right=6mm,top=0mm,bottom=4mm,sharp corners,
overlay={\fill[cder] (frame.north west) rectangle ([yshift=-1.25cm]frame.north east);
\node[anchor=west,text=white,font=\sffamily\bfseries\LARGE] at ([xshift=5mm,yshift=-0.62cm]frame.north west) {Question 43 C};
\node[anchor=east,text=white,font=\sffamily\itshape\large] at ([xshift=-5mm,yshift=-0.62cm]frame.north east) {Dérivation};}]
\vspace*{1.75cm}
\begin{minipage}[t][6.9cm][c]{\linewidth}\LARGE\centering
\begin{minipage}{0.95\linewidth}\raggedright\linespread{1.25}\selectfont Soit $f(x)=(x+1)e^x$. Que vaut $f'(0)$~?\end{minipage}
\end{minipage}

\vfill
\begin{tcolorbox}[enhanced,nobeforeafter,width=7.1cm,height=2.1cm,valign=center,colback=cder!6,colframe=cder,arc=2.5mm,boxrule=0.9pt,left=3mm]{\Large\bfseries\color{cder}a)}\quad\begin{minipage}[c]{5.3cm}\raggedright\Large $3$\end{minipage}\end{tcolorbox}\hfill\begin{tcolorbox}[enhanced,nobeforeafter,width=7.1cm,height=2.1cm,valign=center,colback=cder!6,colframe=cder,arc=2.5mm,boxrule=0.9pt,left=3mm]{\Large\bfseries\color{cder}b)}\quad\begin{minipage}[c]{5.3cm}\raggedright\Large $1$\end{minipage}\end{tcolorbox}\par\vspace{2.5mm}
\begin{tcolorbox}[enhanced,nobeforeafter,width=7.1cm,height=2.1cm,valign=center,colback=cder!6,colframe=cder,arc=2.5mm,boxrule=0.9pt,left=3mm]{\Large\bfseries\color{cder}c)}\quad\begin{minipage}[c]{5.3cm}\raggedright\Large $2$\end{minipage}\end{tcolorbox}\hfill\begin{tcolorbox}[enhanced,nobeforeafter,width=7.1cm,height=2.1cm,valign=center,colback=cder!6,colframe=cder,arc=2.5mm,boxrule=0.9pt,left=3mm]{\Large\bfseries\color{cder}d)}\quad\begin{minipage}[c]{5.3cm}\raggedright\Large $4$\end{minipage}\end{tcolorbox}
\end{tcolorbox}
\newpage
\begin{tcolorbox}[enhanced,width=16cm,height=14cm,colback=white,colframe=cder,boxrule=2pt,arc=0mm,
left=6mm,right=6mm,top=0mm,bottom=4mm,sharp corners,
overlay={\fill[cder] (frame.north west) rectangle ([yshift=-1.25cm]frame.north east);
\node[anchor=west,text=white,font=\sffamily\bfseries\LARGE] at ([xshift=5mm,yshift=-0.62cm]frame.north west) {Question 44 C};
\node[anchor=east,text=white,font=\sffamily\itshape\large] at ([xshift=-5mm,yshift=-0.62cm]frame.north east) {Dérivation};}]
\vspace*{1.75cm}
\begin{minipage}[t][6.9cm][c]{\linewidth}\LARGE\centering
\begin{minipage}{0.95\linewidth}\raggedright\linespread{1.25}\selectfont Soit $f(x)=(-2x+5)e^x$. Que vaut $f'(0)$~?\end{minipage}
\end{minipage}

\vfill
\begin{tcolorbox}[enhanced,nobeforeafter,width=7.1cm,height=2.1cm,valign=center,colback=cder!6,colframe=cder,arc=2.5mm,boxrule=0.9pt,left=3mm]{\Large\bfseries\color{cder}a)}\quad\begin{minipage}[c]{5.3cm}\raggedright\Large $5$\end{minipage}\end{tcolorbox}\hfill\begin{tcolorbox}[enhanced,nobeforeafter,width=7.1cm,height=2.1cm,valign=center,colback=cder!6,colframe=cder,arc=2.5mm,boxrule=0.9pt,left=3mm]{\Large\bfseries\color{cder}b)}\quad\begin{minipage}[c]{5.3cm}\raggedright\Large $3$\end{minipage}\end{tcolorbox}\par\vspace{2.5mm}
\begin{tcolorbox}[enhanced,nobeforeafter,width=7.1cm,height=2.1cm,valign=center,colback=cder!6,colframe=cder,arc=2.5mm,boxrule=0.9pt,left=3mm]{\Large\bfseries\color{cder}c)}\quad\begin{minipage}[c]{5.3cm}\raggedright\Large $-10$\end{minipage}\end{tcolorbox}\hfill\begin{tcolorbox}[enhanced,nobeforeafter,width=7.1cm,height=2.1cm,valign=center,colback=cder!6,colframe=cder,arc=2.5mm,boxrule=0.9pt,left=3mm]{\Large\bfseries\color{cder}d)}\quad\begin{minipage}[c]{5.3cm}\raggedright\Large $-2$\end{minipage}\end{tcolorbox}
\end{tcolorbox}
\newpage
\begin{tcolorbox}[enhanced,width=16cm,height=14cm,colback=white,colframe=cder,boxrule=2pt,arc=0mm,
left=6mm,right=6mm,top=0mm,bottom=4mm,sharp corners,
overlay={\fill[cder] (frame.north west) rectangle ([yshift=-1.25cm]frame.north east);
\node[anchor=west,text=white,font=\sffamily\bfseries\LARGE] at ([xshift=5mm,yshift=-0.62cm]frame.north west) {Question 45 C};
\node[anchor=east,text=white,font=\sffamily\itshape\large] at ([xshift=-5mm,yshift=-0.62cm]frame.north east) {Dérivation};}]
\vspace*{1.75cm}
\begin{minipage}[t][6.9cm][c]{\linewidth}\LARGE\centering
\begin{minipage}{0.95\linewidth}\raggedright\linespread{1.25}\selectfont La fonction $f$ est dérivable sur $\mathbb{R}$ et $f'(x)=5x-10$. Sur quel intervalle $f$ est-elle décroissante~?\end{minipage}
\end{minipage}

\vfill
\begin{tcolorbox}[enhanced,nobeforeafter,width=7.1cm,height=2.1cm,valign=center,colback=cder!6,colframe=cder,arc=2.5mm,boxrule=0.9pt,left=3mm]{\Large\bfseries\color{cder}a)}\quad\begin{minipage}[c]{5.3cm}\raggedright\Large $[2\,;\,+\infty[$\end{minipage}\end{tcolorbox}\hfill\begin{tcolorbox}[enhanced,nobeforeafter,width=7.1cm,height=2.1cm,valign=center,colback=cder!6,colframe=cder,arc=2.5mm,boxrule=0.9pt,left=3mm]{\Large\bfseries\color{cder}b)}\quad\begin{minipage}[c]{5.3cm}\raggedright\Large $]-\infty\,;\,2]$\end{minipage}\end{tcolorbox}\par\vspace{2.5mm}
\begin{tcolorbox}[enhanced,nobeforeafter,width=7.1cm,height=2.1cm,valign=center,colback=cder!6,colframe=cder,arc=2.5mm,boxrule=0.9pt,left=3mm]{\Large\bfseries\color{cder}c)}\quad\begin{minipage}[c]{5.3cm}\raggedright\Large $[-2\,;\,+\infty[$\end{minipage}\end{tcolorbox}\hfill\begin{tcolorbox}[enhanced,nobeforeafter,width=7.1cm,height=2.1cm,valign=center,colback=cder!6,colframe=cder,arc=2.5mm,boxrule=0.9pt,left=3mm]{\Large\bfseries\color{cder}d)}\quad\begin{minipage}[c]{5.3cm}\raggedright\Large $]-\infty\,;\,-2]$\end{minipage}\end{tcolorbox}
\end{tcolorbox}
\newpage
\begin{tcolorbox}[enhanced,width=16cm,height=14cm,colback=white,colframe=cder,boxrule=2pt,arc=0mm,
left=6mm,right=6mm,top=0mm,bottom=4mm,sharp corners,
overlay={\fill[cder] (frame.north west) rectangle ([yshift=-1.25cm]frame.north east);
\node[anchor=west,text=white,font=\sffamily\bfseries\LARGE] at ([xshift=5mm,yshift=-0.62cm]frame.north west) {Question 46 C};
\node[anchor=east,text=white,font=\sffamily\itshape\large] at ([xshift=-5mm,yshift=-0.62cm]frame.north east) {Dérivation};}]
\vspace*{1.75cm}
\begin{minipage}[t][6.9cm][c]{\linewidth}\LARGE\centering
\begin{minipage}{0.95\linewidth}\raggedright\linespread{1.25}\selectfont La fonction $f$ est dérivable sur $\mathbb{R}$ et $f'(x)=-2x+1$. Sur quel intervalle $f$ est-elle décroissante~?\end{minipage}
\end{minipage}

\vfill
\begin{tcolorbox}[enhanced,nobeforeafter,width=7.1cm,height=2.1cm,valign=center,colback=cder!6,colframe=cder,arc=2.5mm,boxrule=0.9pt,left=3mm]{\Large\bfseries\color{cder}a)}\quad\begin{minipage}[c]{5.3cm}\raggedright\Large $[0{,}5\,;\,+\infty[$\end{minipage}\end{tcolorbox}\hfill\begin{tcolorbox}[enhanced,nobeforeafter,width=7.1cm,height=2.1cm,valign=center,colback=cder!6,colframe=cder,arc=2.5mm,boxrule=0.9pt,left=3mm]{\Large\bfseries\color{cder}b)}\quad\begin{minipage}[c]{5.3cm}\raggedright\Large $[-0{,}5\,;\,+\infty[$\end{minipage}\end{tcolorbox}\par\vspace{2.5mm}
\begin{tcolorbox}[enhanced,nobeforeafter,width=7.1cm,height=2.1cm,valign=center,colback=cder!6,colframe=cder,arc=2.5mm,boxrule=0.9pt,left=3mm]{\Large\bfseries\color{cder}c)}\quad\begin{minipage}[c]{5.3cm}\raggedright\Large $]-\infty\,;\,0{,}5]$\end{minipage}\end{tcolorbox}\hfill\begin{tcolorbox}[enhanced,nobeforeafter,width=7.1cm,height=2.1cm,valign=center,colback=cder!6,colframe=cder,arc=2.5mm,boxrule=0.9pt,left=3mm]{\Large\bfseries\color{cder}d)}\quad\begin{minipage}[c]{5.3cm}\raggedright\Large $]-\infty\,;\,-0{,}5]$\end{minipage}\end{tcolorbox}
\end{tcolorbox}
\newpage
\begin{tcolorbox}[enhanced,width=16cm,height=14cm,colback=white,colframe=cder,boxrule=2pt,arc=0mm,
left=6mm,right=6mm,top=0mm,bottom=4mm,sharp corners,
overlay={\fill[cder] (frame.north west) rectangle ([yshift=-1.25cm]frame.north east);
\node[anchor=west,text=white,font=\sffamily\bfseries\LARGE] at ([xshift=5mm,yshift=-0.62cm]frame.north west) {Question 47 C};
\node[anchor=east,text=white,font=\sffamily\itshape\large] at ([xshift=-5mm,yshift=-0.62cm]frame.north east) {Dérivation};}]
\vspace*{1.75cm}
\begin{minipage}[t][6.9cm][c]{\linewidth}\LARGE\centering
\begin{minipage}{0.95\linewidth}\raggedright\linespread{1.25}\selectfont Soit $f(x)=x^2+2x-8$. Quel est le minimum de $f$ sur $\mathbb{R}$~?\end{minipage}
\end{minipage}

\vfill
\begin{tcolorbox}[enhanced,nobeforeafter,width=7.1cm,height=2.1cm,valign=center,colback=cder!6,colframe=cder,arc=2.5mm,boxrule=0.9pt,left=3mm]{\Large\bfseries\color{cder}a)}\quad\begin{minipage}[c]{5.3cm}\raggedright\Large $-1$\end{minipage}\end{tcolorbox}\hfill\begin{tcolorbox}[enhanced,nobeforeafter,width=7.1cm,height=2.1cm,valign=center,colback=cder!6,colframe=cder,arc=2.5mm,boxrule=0.9pt,left=3mm]{\Large\bfseries\color{cder}b)}\quad\begin{minipage}[c]{5.3cm}\raggedright\Large $-8$\end{minipage}\end{tcolorbox}\par\vspace{2.5mm}
\begin{tcolorbox}[enhanced,nobeforeafter,width=7.1cm,height=2.1cm,valign=center,colback=cder!6,colframe=cder,arc=2.5mm,boxrule=0.9pt,left=3mm]{\Large\bfseries\color{cder}c)}\quad\begin{minipage}[c]{5.3cm}\raggedright\Large $9$\end{minipage}\end{tcolorbox}\hfill\begin{tcolorbox}[enhanced,nobeforeafter,width=7.1cm,height=2.1cm,valign=center,colback=cder!6,colframe=cder,arc=2.5mm,boxrule=0.9pt,left=3mm]{\Large\bfseries\color{cder}d)}\quad\begin{minipage}[c]{5.3cm}\raggedright\Large $-9$\end{minipage}\end{tcolorbox}
\end{tcolorbox}
\newpage
\begin{tcolorbox}[enhanced,width=16cm,height=14cm,colback=white,colframe=cder,boxrule=2pt,arc=0mm,
left=6mm,right=6mm,top=0mm,bottom=4mm,sharp corners,
overlay={\fill[cder] (frame.north west) rectangle ([yshift=-1.25cm]frame.north east);
\node[anchor=west,text=white,font=\sffamily\bfseries\LARGE] at ([xshift=5mm,yshift=-0.62cm]frame.north west) {Question 48 C};
\node[anchor=east,text=white,font=\sffamily\itshape\large] at ([xshift=-5mm,yshift=-0.62cm]frame.north east) {Dérivation};}]
\vspace*{1.75cm}
\begin{minipage}[t][6.9cm][c]{\linewidth}\LARGE\centering
\begin{minipage}{0.95\linewidth}\raggedright\linespread{1.25}\selectfont Soit $f(x)=-x^2-4x+3$. Quel est le maximum de $f$ sur $\mathbb{R}$~?\end{minipage}
\end{minipage}

\vfill
\begin{tcolorbox}[enhanced,nobeforeafter,width=7.1cm,height=2.1cm,valign=center,colback=cder!6,colframe=cder,arc=2.5mm,boxrule=0.9pt,left=3mm]{\Large\bfseries\color{cder}a)}\quad\begin{minipage}[c]{5.3cm}\raggedright\Large $-7$\end{minipage}\end{tcolorbox}\hfill\begin{tcolorbox}[enhanced,nobeforeafter,width=7.1cm,height=2.1cm,valign=center,colback=cder!6,colframe=cder,arc=2.5mm,boxrule=0.9pt,left=3mm]{\Large\bfseries\color{cder}b)}\quad\begin{minipage}[c]{5.3cm}\raggedright\Large $7$\end{minipage}\end{tcolorbox}\par\vspace{2.5mm}
\begin{tcolorbox}[enhanced,nobeforeafter,width=7.1cm,height=2.1cm,valign=center,colback=cder!6,colframe=cder,arc=2.5mm,boxrule=0.9pt,left=3mm]{\Large\bfseries\color{cder}c)}\quad\begin{minipage}[c]{5.3cm}\raggedright\Large $-2$\end{minipage}\end{tcolorbox}\hfill\begin{tcolorbox}[enhanced,nobeforeafter,width=7.1cm,height=2.1cm,valign=center,colback=cder!6,colframe=cder,arc=2.5mm,boxrule=0.9pt,left=3mm]{\Large\bfseries\color{cder}d)}\quad\begin{minipage}[c]{5.3cm}\raggedright\Large $3$\end{minipage}\end{tcolorbox}
\end{tcolorbox}
\newpage
\begin{tcolorbox}[enhanced,width=16cm,height=14cm,colback=white,colframe=ctri,boxrule=2pt,arc=0mm,
left=6mm,right=6mm,top=0mm,bottom=4mm,sharp corners,
overlay={\fill[ctri] (frame.north west) rectangle ([yshift=-1.25cm]frame.north east);
\node[anchor=west,text=white,font=\sffamily\bfseries\LARGE] at ([xshift=5mm,yshift=-0.62cm]frame.north west) {Question 49 C};
\node[anchor=east,text=white,font=\sffamily\itshape\large] at ([xshift=-5mm,yshift=-0.62cm]frame.north east) {Trigonométrie et exponentielle};}]
\vspace*{1.75cm}
\begin{minipage}[t][6.9cm][c]{\linewidth}\LARGE\centering
\begin{minipage}{0.95\linewidth}\raggedright\linespread{1.25}\selectfont Quelle est la mesure en radians d'un angle de $135^\circ$~?\end{minipage}
\end{minipage}

\vfill
\begin{tcolorbox}[enhanced,nobeforeafter,width=7.1cm,height=2.1cm,valign=center,colback=ctri!6,colframe=ctri,arc=2.5mm,boxrule=0.9pt,left=3mm]{\Large\bfseries\color{ctri}a)}\quad\begin{minipage}[c]{5.3cm}\raggedright\Large $\dfrac{3\pi}{2}$\end{minipage}\end{tcolorbox}\hfill\begin{tcolorbox}[enhanced,nobeforeafter,width=7.1cm,height=2.1cm,valign=center,colback=ctri!6,colframe=ctri,arc=2.5mm,boxrule=0.9pt,left=3mm]{\Large\bfseries\color{ctri}b)}\quad\begin{minipage}[c]{5.3cm}\raggedright\Large $\dfrac{3\pi}{8}$\end{minipage}\end{tcolorbox}\par\vspace{2.5mm}
\begin{tcolorbox}[enhanced,nobeforeafter,width=7.1cm,height=2.1cm,valign=center,colback=ctri!6,colframe=ctri,arc=2.5mm,boxrule=0.9pt,left=3mm]{\Large\bfseries\color{ctri}c)}\quad\begin{minipage}[c]{5.3cm}\raggedright\Large $\dfrac{3\pi}{4}$\end{minipage}\end{tcolorbox}\hfill\begin{tcolorbox}[enhanced,nobeforeafter,width=7.1cm,height=2.1cm,valign=center,colback=ctri!6,colframe=ctri,arc=2.5mm,boxrule=0.9pt,left=3mm]{\Large\bfseries\color{ctri}d)}\quad\begin{minipage}[c]{5.3cm}\raggedright\Large $\dfrac{4\pi}{3}$\end{minipage}\end{tcolorbox}
\end{tcolorbox}
\newpage
\begin{tcolorbox}[enhanced,width=16cm,height=14cm,colback=white,colframe=ctri,boxrule=2pt,arc=0mm,
left=6mm,right=6mm,top=0mm,bottom=4mm,sharp corners,
overlay={\fill[ctri] (frame.north west) rectangle ([yshift=-1.25cm]frame.north east);
\node[anchor=west,text=white,font=\sffamily\bfseries\LARGE] at ([xshift=5mm,yshift=-0.62cm]frame.north west) {Question 50 C};
\node[anchor=east,text=white,font=\sffamily\itshape\large] at ([xshift=-5mm,yshift=-0.62cm]frame.north east) {Trigonométrie et exponentielle};}]
\vspace*{1.75cm}
\begin{minipage}[t][6.9cm][c]{\linewidth}\LARGE\centering
\begin{minipage}{0.95\linewidth}\raggedright\linespread{1.25}\selectfont Quelle est la mesure en radians d'un angle de $150^\circ$~?\end{minipage}
\end{minipage}

\vfill
\begin{tcolorbox}[enhanced,nobeforeafter,width=7.1cm,height=2.1cm,valign=center,colback=ctri!6,colframe=ctri,arc=2.5mm,boxrule=0.9pt,left=3mm]{\Large\bfseries\color{ctri}a)}\quad\begin{minipage}[c]{5.3cm}\raggedright\Large $\dfrac{5\pi}{12}$\end{minipage}\end{tcolorbox}\hfill\begin{tcolorbox}[enhanced,nobeforeafter,width=7.1cm,height=2.1cm,valign=center,colback=ctri!6,colframe=ctri,arc=2.5mm,boxrule=0.9pt,left=3mm]{\Large\bfseries\color{ctri}b)}\quad\begin{minipage}[c]{5.3cm}\raggedright\Large $\dfrac{5\pi}{6}$\end{minipage}\end{tcolorbox}\par\vspace{2.5mm}
\begin{tcolorbox}[enhanced,nobeforeafter,width=7.1cm,height=2.1cm,valign=center,colback=ctri!6,colframe=ctri,arc=2.5mm,boxrule=0.9pt,left=3mm]{\Large\bfseries\color{ctri}c)}\quad\begin{minipage}[c]{5.3cm}\raggedright\Large $\dfrac{6\pi}{5}$\end{minipage}\end{tcolorbox}\hfill\begin{tcolorbox}[enhanced,nobeforeafter,width=7.1cm,height=2.1cm,valign=center,colback=ctri!6,colframe=ctri,arc=2.5mm,boxrule=0.9pt,left=3mm]{\Large\bfseries\color{ctri}d)}\quad\begin{minipage}[c]{5.3cm}\raggedright\Large $\dfrac{5\pi}{3}$\end{minipage}\end{tcolorbox}
\end{tcolorbox}
\newpage
\begin{tcolorbox}[enhanced,width=16cm,height=14cm,colback=white,colframe=ctri,boxrule=2pt,arc=0mm,
left=6mm,right=6mm,top=0mm,bottom=4mm,sharp corners,
overlay={\fill[ctri] (frame.north west) rectangle ([yshift=-1.25cm]frame.north east);
\node[anchor=west,text=white,font=\sffamily\bfseries\LARGE] at ([xshift=5mm,yshift=-0.62cm]frame.north west) {Question 51 C};
\node[anchor=east,text=white,font=\sffamily\itshape\large] at ([xshift=-5mm,yshift=-0.62cm]frame.north east) {Trigonométrie et exponentielle};}]
\vspace*{1.75cm}
\begin{minipage}[t][6.9cm][c]{\linewidth}\LARGE\centering
\begin{minipage}{0.95\linewidth}\raggedright\linespread{1.25}\selectfont Quelle est la mesure en degrés d'un angle de $\dfrac{5\pi}{4}$ radians~?\end{minipage}
\end{minipage}

\vfill
\begin{tcolorbox}[enhanced,nobeforeafter,width=7.1cm,height=2.1cm,valign=center,colback=ctri!6,colframe=ctri,arc=2.5mm,boxrule=0.9pt,left=3mm]{\Large\bfseries\color{ctri}a)}\quad\begin{minipage}[c]{5.3cm}\raggedright\Large $450^\circ$\end{minipage}\end{tcolorbox}\hfill\begin{tcolorbox}[enhanced,nobeforeafter,width=7.1cm,height=2.1cm,valign=center,colback=ctri!6,colframe=ctri,arc=2.5mm,boxrule=0.9pt,left=3mm]{\Large\bfseries\color{ctri}b)}\quad\begin{minipage}[c]{5.3cm}\raggedright\Large $225^\circ$\end{minipage}\end{tcolorbox}\par\vspace{2.5mm}
\begin{tcolorbox}[enhanced,nobeforeafter,width=7.1cm,height=2.1cm,valign=center,colback=ctri!6,colframe=ctri,arc=2.5mm,boxrule=0.9pt,left=3mm]{\Large\bfseries\color{ctri}c)}\quad\begin{minipage}[c]{5.3cm}\raggedright\Large $135^\circ$\end{minipage}\end{tcolorbox}\hfill\begin{tcolorbox}[enhanced,nobeforeafter,width=7.1cm,height=2.1cm,valign=center,colback=ctri!6,colframe=ctri,arc=2.5mm,boxrule=0.9pt,left=3mm]{\Large\bfseries\color{ctri}d)}\quad\begin{minipage}[c]{5.3cm}\raggedright\Large $112{,}5^\circ$\end{minipage}\end{tcolorbox}
\end{tcolorbox}
\newpage
\begin{tcolorbox}[enhanced,width=16cm,height=14cm,colback=white,colframe=ctri,boxrule=2pt,arc=0mm,
left=6mm,right=6mm,top=0mm,bottom=4mm,sharp corners,
overlay={\fill[ctri] (frame.north west) rectangle ([yshift=-1.25cm]frame.north east);
\node[anchor=west,text=white,font=\sffamily\bfseries\LARGE] at ([xshift=5mm,yshift=-0.62cm]frame.north west) {Question 52 C};
\node[anchor=east,text=white,font=\sffamily\itshape\large] at ([xshift=-5mm,yshift=-0.62cm]frame.north east) {Trigonométrie et exponentielle};}]
\vspace*{1.75cm}
\begin{minipage}[t][6.9cm][c]{\linewidth}\LARGE\centering
\begin{minipage}{0.95\linewidth}\raggedright\linespread{1.25}\selectfont Quelle est la mesure en degrés d'un angle de $\dfrac{7\pi}{6}$ radians~?\end{minipage}
\end{minipage}

\vfill
\begin{tcolorbox}[enhanced,nobeforeafter,width=7.1cm,height=2.1cm,valign=center,colback=ctri!6,colframe=ctri,arc=2.5mm,boxrule=0.9pt,left=3mm]{\Large\bfseries\color{ctri}a)}\quad\begin{minipage}[c]{5.3cm}\raggedright\Large $210^\circ$\end{minipage}\end{tcolorbox}\hfill\begin{tcolorbox}[enhanced,nobeforeafter,width=7.1cm,height=2.1cm,valign=center,colback=ctri!6,colframe=ctri,arc=2.5mm,boxrule=0.9pt,left=3mm]{\Large\bfseries\color{ctri}b)}\quad\begin{minipage}[c]{5.3cm}\raggedright\Large $150^\circ$\end{minipage}\end{tcolorbox}\par\vspace{2.5mm}
\begin{tcolorbox}[enhanced,nobeforeafter,width=7.1cm,height=2.1cm,valign=center,colback=ctri!6,colframe=ctri,arc=2.5mm,boxrule=0.9pt,left=3mm]{\Large\bfseries\color{ctri}c)}\quad\begin{minipage}[c]{5.3cm}\raggedright\Large $105^\circ$\end{minipage}\end{tcolorbox}\hfill\begin{tcolorbox}[enhanced,nobeforeafter,width=7.1cm,height=2.1cm,valign=center,colback=ctri!6,colframe=ctri,arc=2.5mm,boxrule=0.9pt,left=3mm]{\Large\bfseries\color{ctri}d)}\quad\begin{minipage}[c]{5.3cm}\raggedright\Large $420^\circ$\end{minipage}\end{tcolorbox}
\end{tcolorbox}
\newpage
\begin{tcolorbox}[enhanced,width=16cm,height=14cm,colback=white,colframe=ctri,boxrule=2pt,arc=0mm,
left=6mm,right=6mm,top=0mm,bottom=4mm,sharp corners,
overlay={\fill[ctri] (frame.north west) rectangle ([yshift=-1.25cm]frame.north east);
\node[anchor=west,text=white,font=\sffamily\bfseries\LARGE] at ([xshift=5mm,yshift=-0.62cm]frame.north west) {Question 53 C};
\node[anchor=east,text=white,font=\sffamily\itshape\large] at ([xshift=-5mm,yshift=-0.62cm]frame.north east) {Trigonométrie et exponentielle};}]
\vspace*{1.75cm}
\begin{minipage}[t][6.9cm][c]{\linewidth}\LARGE\centering
\begin{minipage}{0.95\linewidth}\raggedright\linespread{1.25}\selectfont Que vaut $\cos\left(\pi\right)$~?\end{minipage}
\end{minipage}

\vfill
\begin{tcolorbox}[enhanced,nobeforeafter,width=7.1cm,height=2.1cm,valign=center,colback=ctri!6,colframe=ctri,arc=2.5mm,boxrule=0.9pt,left=3mm]{\Large\bfseries\color{ctri}a)}\quad\begin{minipage}[c]{5.3cm}\raggedright\Large $\dfrac{\sqrt3}{2}$\end{minipage}\end{tcolorbox}\hfill\begin{tcolorbox}[enhanced,nobeforeafter,width=7.1cm,height=2.1cm,valign=center,colback=ctri!6,colframe=ctri,arc=2.5mm,boxrule=0.9pt,left=3mm]{\Large\bfseries\color{ctri}b)}\quad\begin{minipage}[c]{5.3cm}\raggedright\Large $-\dfrac{\sqrt3}{2}$\end{minipage}\end{tcolorbox}\par\vspace{2.5mm}
\begin{tcolorbox}[enhanced,nobeforeafter,width=7.1cm,height=2.1cm,valign=center,colback=ctri!6,colframe=ctri,arc=2.5mm,boxrule=0.9pt,left=3mm]{\Large\bfseries\color{ctri}c)}\quad\begin{minipage}[c]{5.3cm}\raggedright\Large $-1$\end{minipage}\end{tcolorbox}\hfill\begin{tcolorbox}[enhanced,nobeforeafter,width=7.1cm,height=2.1cm,valign=center,colback=ctri!6,colframe=ctri,arc=2.5mm,boxrule=0.9pt,left=3mm]{\Large\bfseries\color{ctri}d)}\quad\begin{minipage}[c]{5.3cm}\raggedright\Large $1$\end{minipage}\end{tcolorbox}
\end{tcolorbox}
\newpage
\begin{tcolorbox}[enhanced,width=16cm,height=14cm,colback=white,colframe=ctri,boxrule=2pt,arc=0mm,
left=6mm,right=6mm,top=0mm,bottom=4mm,sharp corners,
overlay={\fill[ctri] (frame.north west) rectangle ([yshift=-1.25cm]frame.north east);
\node[anchor=west,text=white,font=\sffamily\bfseries\LARGE] at ([xshift=5mm,yshift=-0.62cm]frame.north west) {Question 54 C};
\node[anchor=east,text=white,font=\sffamily\itshape\large] at ([xshift=-5mm,yshift=-0.62cm]frame.north east) {Trigonométrie et exponentielle};}]
\vspace*{1.75cm}
\begin{minipage}[t][6.9cm][c]{\linewidth}\LARGE\centering
\begin{minipage}{0.95\linewidth}\raggedright\linespread{1.25}\selectfont Que vaut $\cos\left(\dfrac{3\pi}{4}\right)$~?\end{minipage}
\end{minipage}

\vfill
\begin{tcolorbox}[enhanced,nobeforeafter,width=7.1cm,height=2.1cm,valign=center,colback=ctri!6,colframe=ctri,arc=2.5mm,boxrule=0.9pt,left=3mm]{\Large\bfseries\color{ctri}a)}\quad\begin{minipage}[c]{5.3cm}\raggedright\Large $-\dfrac{\sqrt3}{2}$\end{minipage}\end{tcolorbox}\hfill\begin{tcolorbox}[enhanced,nobeforeafter,width=7.1cm,height=2.1cm,valign=center,colback=ctri!6,colframe=ctri,arc=2.5mm,boxrule=0.9pt,left=3mm]{\Large\bfseries\color{ctri}b)}\quad\begin{minipage}[c]{5.3cm}\raggedright\Large $\dfrac{\sqrt2}{2}$\end{minipage}\end{tcolorbox}\par\vspace{2.5mm}
\begin{tcolorbox}[enhanced,nobeforeafter,width=7.1cm,height=2.1cm,valign=center,colback=ctri!6,colframe=ctri,arc=2.5mm,boxrule=0.9pt,left=3mm]{\Large\bfseries\color{ctri}c)}\quad\begin{minipage}[c]{5.3cm}\raggedright\Large $-\dfrac12$\end{minipage}\end{tcolorbox}\hfill\begin{tcolorbox}[enhanced,nobeforeafter,width=7.1cm,height=2.1cm,valign=center,colback=ctri!6,colframe=ctri,arc=2.5mm,boxrule=0.9pt,left=3mm]{\Large\bfseries\color{ctri}d)}\quad\begin{minipage}[c]{5.3cm}\raggedright\Large $-\dfrac{\sqrt2}{2}$\end{minipage}\end{tcolorbox}
\end{tcolorbox}
\newpage
\begin{tcolorbox}[enhanced,width=16cm,height=14cm,colback=white,colframe=ctri,boxrule=2pt,arc=0mm,
left=6mm,right=6mm,top=0mm,bottom=4mm,sharp corners,
overlay={\fill[ctri] (frame.north west) rectangle ([yshift=-1.25cm]frame.north east);
\node[anchor=west,text=white,font=\sffamily\bfseries\LARGE] at ([xshift=5mm,yshift=-0.62cm]frame.north west) {Question 55 C};
\node[anchor=east,text=white,font=\sffamily\itshape\large] at ([xshift=-5mm,yshift=-0.62cm]frame.north east) {Trigonométrie et exponentielle};}]
\vspace*{1.75cm}
\begin{minipage}[t][6.9cm][c]{\linewidth}\LARGE\centering
\begin{minipage}{0.95\linewidth}\raggedright\linespread{1.25}\selectfont Que vaut $\sin\left(-\dfrac{\pi}{3}\right)$~?\end{minipage}
\end{minipage}

\vfill
\begin{tcolorbox}[enhanced,nobeforeafter,width=7.1cm,height=2.1cm,valign=center,colback=ctri!6,colframe=ctri,arc=2.5mm,boxrule=0.9pt,left=3mm]{\Large\bfseries\color{ctri}a)}\quad\begin{minipage}[c]{5.3cm}\raggedright\Large $-1$\end{minipage}\end{tcolorbox}\hfill\begin{tcolorbox}[enhanced,nobeforeafter,width=7.1cm,height=2.1cm,valign=center,colback=ctri!6,colframe=ctri,arc=2.5mm,boxrule=0.9pt,left=3mm]{\Large\bfseries\color{ctri}b)}\quad\begin{minipage}[c]{5.3cm}\raggedright\Large $-\dfrac{\sqrt2}{2}$\end{minipage}\end{tcolorbox}\par\vspace{2.5mm}
\begin{tcolorbox}[enhanced,nobeforeafter,width=7.1cm,height=2.1cm,valign=center,colback=ctri!6,colframe=ctri,arc=2.5mm,boxrule=0.9pt,left=3mm]{\Large\bfseries\color{ctri}c)}\quad\begin{minipage}[c]{5.3cm}\raggedright\Large $-\dfrac{\sqrt3}{2}$\end{minipage}\end{tcolorbox}\hfill\begin{tcolorbox}[enhanced,nobeforeafter,width=7.1cm,height=2.1cm,valign=center,colback=ctri!6,colframe=ctri,arc=2.5mm,boxrule=0.9pt,left=3mm]{\Large\bfseries\color{ctri}d)}\quad\begin{minipage}[c]{5.3cm}\raggedright\Large $\dfrac{\sqrt3}{2}$\end{minipage}\end{tcolorbox}
\end{tcolorbox}
\newpage
\begin{tcolorbox}[enhanced,width=16cm,height=14cm,colback=white,colframe=ctri,boxrule=2pt,arc=0mm,
left=6mm,right=6mm,top=0mm,bottom=4mm,sharp corners,
overlay={\fill[ctri] (frame.north west) rectangle ([yshift=-1.25cm]frame.north east);
\node[anchor=west,text=white,font=\sffamily\bfseries\LARGE] at ([xshift=5mm,yshift=-0.62cm]frame.north west) {Question 56 C};
\node[anchor=east,text=white,font=\sffamily\itshape\large] at ([xshift=-5mm,yshift=-0.62cm]frame.north east) {Trigonométrie et exponentielle};}]
\vspace*{1.75cm}
\begin{minipage}[t][6.9cm][c]{\linewidth}\LARGE\centering
\begin{minipage}{0.95\linewidth}\raggedright\linespread{1.25}\selectfont Que vaut $\sin\left(\pi\right)$~?\end{minipage}
\end{minipage}

\vfill
\begin{tcolorbox}[enhanced,nobeforeafter,width=7.1cm,height=2.1cm,valign=center,colback=ctri!6,colframe=ctri,arc=2.5mm,boxrule=0.9pt,left=3mm]{\Large\bfseries\color{ctri}a)}\quad\begin{minipage}[c]{5.3cm}\raggedright\Large $1$\end{minipage}\end{tcolorbox}\hfill\begin{tcolorbox}[enhanced,nobeforeafter,width=7.1cm,height=2.1cm,valign=center,colback=ctri!6,colframe=ctri,arc=2.5mm,boxrule=0.9pt,left=3mm]{\Large\bfseries\color{ctri}b)}\quad\begin{minipage}[c]{5.3cm}\raggedright\Large $-\dfrac12$\end{minipage}\end{tcolorbox}\par\vspace{2.5mm}
\begin{tcolorbox}[enhanced,nobeforeafter,width=7.1cm,height=2.1cm,valign=center,colback=ctri!6,colframe=ctri,arc=2.5mm,boxrule=0.9pt,left=3mm]{\Large\bfseries\color{ctri}c)}\quad\begin{minipage}[c]{5.3cm}\raggedright\Large $0$\end{minipage}\end{tcolorbox}\hfill\begin{tcolorbox}[enhanced,nobeforeafter,width=7.1cm,height=2.1cm,valign=center,colback=ctri!6,colframe=ctri,arc=2.5mm,boxrule=0.9pt,left=3mm]{\Large\bfseries\color{ctri}d)}\quad\begin{minipage}[c]{5.3cm}\raggedright\Large $\dfrac12$\end{minipage}\end{tcolorbox}
\end{tcolorbox}
\newpage
\begin{tcolorbox}[enhanced,width=16cm,height=14cm,colback=white,colframe=ctri,boxrule=2pt,arc=0mm,
left=6mm,right=6mm,top=0mm,bottom=4mm,sharp corners,
overlay={\fill[ctri] (frame.north west) rectangle ([yshift=-1.25cm]frame.north east);
\node[anchor=west,text=white,font=\sffamily\bfseries\LARGE] at ([xshift=5mm,yshift=-0.62cm]frame.north west) {Question 57 C};
\node[anchor=east,text=white,font=\sffamily\itshape\large] at ([xshift=-5mm,yshift=-0.62cm]frame.north east) {Trigonométrie et exponentielle};}]
\vspace*{1.75cm}
\begin{minipage}[t][6.9cm][c]{\linewidth}\LARGE\centering
\begin{minipage}{0.95\linewidth}\raggedright\linespread{1.25}\selectfont On sait que $\sin x=0{,}25$. Que vaut $\sin(\pi-x)$~?\end{minipage}
\end{minipage}

\vfill
\begin{tcolorbox}[enhanced,nobeforeafter,width=7.1cm,height=2.1cm,valign=center,colback=ctri!6,colframe=ctri,arc=2.5mm,boxrule=0.9pt,left=3mm]{\Large\bfseries\color{ctri}a)}\quad\begin{minipage}[c]{5.3cm}\raggedright\Large $-0{,}25$\end{minipage}\end{tcolorbox}\hfill\begin{tcolorbox}[enhanced,nobeforeafter,width=7.1cm,height=2.1cm,valign=center,colback=ctri!6,colframe=ctri,arc=2.5mm,boxrule=0.9pt,left=3mm]{\Large\bfseries\color{ctri}b)}\quad\begin{minipage}[c]{5.3cm}\raggedright\Large $0{,}9375$\end{minipage}\end{tcolorbox}\par\vspace{2.5mm}
\begin{tcolorbox}[enhanced,nobeforeafter,width=7.1cm,height=2.1cm,valign=center,colback=ctri!6,colframe=ctri,arc=2.5mm,boxrule=0.9pt,left=3mm]{\Large\bfseries\color{ctri}c)}\quad\begin{minipage}[c]{5.3cm}\raggedright\Large $0{,}75$\end{minipage}\end{tcolorbox}\hfill\begin{tcolorbox}[enhanced,nobeforeafter,width=7.1cm,height=2.1cm,valign=center,colback=ctri!6,colframe=ctri,arc=2.5mm,boxrule=0.9pt,left=3mm]{\Large\bfseries\color{ctri}d)}\quad\begin{minipage}[c]{5.3cm}\raggedright\Large $0{,}25$\end{minipage}\end{tcolorbox}
\end{tcolorbox}
\newpage
\begin{tcolorbox}[enhanced,width=16cm,height=14cm,colback=white,colframe=ctri,boxrule=2pt,arc=0mm,
left=6mm,right=6mm,top=0mm,bottom=4mm,sharp corners,
overlay={\fill[ctri] (frame.north west) rectangle ([yshift=-1.25cm]frame.north east);
\node[anchor=west,text=white,font=\sffamily\bfseries\LARGE] at ([xshift=5mm,yshift=-0.62cm]frame.north west) {Question 58 C};
\node[anchor=east,text=white,font=\sffamily\itshape\large] at ([xshift=-5mm,yshift=-0.62cm]frame.north east) {Trigonométrie et exponentielle};}]
\vspace*{1.75cm}
\begin{minipage}[t][6.9cm][c]{\linewidth}\LARGE\centering
\begin{minipage}{0.95\linewidth}\raggedright\linespread{1.25}\selectfont On sait que $\sin x=0{,}7$. Que vaut $\sin(\pi+x)$~?\end{minipage}
\end{minipage}

\vfill
\begin{tcolorbox}[enhanced,nobeforeafter,width=7.1cm,height=2.1cm,valign=center,colback=ctri!6,colframe=ctri,arc=2.5mm,boxrule=0.9pt,left=3mm]{\Large\bfseries\color{ctri}a)}\quad\begin{minipage}[c]{5.3cm}\raggedright\Large $0{,}7$\end{minipage}\end{tcolorbox}\hfill\begin{tcolorbox}[enhanced,nobeforeafter,width=7.1cm,height=2.1cm,valign=center,colback=ctri!6,colframe=ctri,arc=2.5mm,boxrule=0.9pt,left=3mm]{\Large\bfseries\color{ctri}b)}\quad\begin{minipage}[c]{5.3cm}\raggedright\Large $-0{,}7$\end{minipage}\end{tcolorbox}\par\vspace{2.5mm}
\begin{tcolorbox}[enhanced,nobeforeafter,width=7.1cm,height=2.1cm,valign=center,colback=ctri!6,colframe=ctri,arc=2.5mm,boxrule=0.9pt,left=3mm]{\Large\bfseries\color{ctri}c)}\quad\begin{minipage}[c]{5.3cm}\raggedright\Large $0{,}51$\end{minipage}\end{tcolorbox}\hfill\begin{tcolorbox}[enhanced,nobeforeafter,width=7.1cm,height=2.1cm,valign=center,colback=ctri!6,colframe=ctri,arc=2.5mm,boxrule=0.9pt,left=3mm]{\Large\bfseries\color{ctri}d)}\quad\begin{minipage}[c]{5.3cm}\raggedright\Large $0{,}3$\end{minipage}\end{tcolorbox}
\end{tcolorbox}
\newpage
\begin{tcolorbox}[enhanced,width=16cm,height=14cm,colback=white,colframe=ctri,boxrule=2pt,arc=0mm,
left=6mm,right=6mm,top=0mm,bottom=4mm,sharp corners,
overlay={\fill[ctri] (frame.north west) rectangle ([yshift=-1.25cm]frame.north east);
\node[anchor=west,text=white,font=\sffamily\bfseries\LARGE] at ([xshift=5mm,yshift=-0.62cm]frame.north west) {Question 59 C};
\node[anchor=east,text=white,font=\sffamily\itshape\large] at ([xshift=-5mm,yshift=-0.62cm]frame.north east) {Trigonométrie et exponentielle};}]
\vspace*{1.75cm}
\begin{minipage}[t][6.9cm][c]{\linewidth}\LARGE\centering
\begin{minipage}{0.95\linewidth}\raggedright\linespread{1.25}\selectfont Quelle est la mesure principale de l'angle $\dfrac{13\pi}{4}$~?\end{minipage}
\end{minipage}

\vfill
\begin{tcolorbox}[enhanced,nobeforeafter,width=7.1cm,height=2.1cm,valign=center,colback=ctri!6,colframe=ctri,arc=2.5mm,boxrule=0.9pt,left=3mm]{\Large\bfseries\color{ctri}a)}\quad\begin{minipage}[c]{5.3cm}\raggedright\Large $-\dfrac{3\pi}{8}$\end{minipage}\end{tcolorbox}\hfill\begin{tcolorbox}[enhanced,nobeforeafter,width=7.1cm,height=2.1cm,valign=center,colback=ctri!6,colframe=ctri,arc=2.5mm,boxrule=0.9pt,left=3mm]{\Large\bfseries\color{ctri}b)}\quad\begin{minipage}[c]{5.3cm}\raggedright\Large $-\dfrac{3\pi}{4}$\end{minipage}\end{tcolorbox}\par\vspace{2.5mm}
\begin{tcolorbox}[enhanced,nobeforeafter,width=7.1cm,height=2.1cm,valign=center,colback=ctri!6,colframe=ctri,arc=2.5mm,boxrule=0.9pt,left=3mm]{\Large\bfseries\color{ctri}c)}\quad\begin{minipage}[c]{5.3cm}\raggedright\Large $\dfrac{\pi}{4}$\end{minipage}\end{tcolorbox}\hfill\begin{tcolorbox}[enhanced,nobeforeafter,width=7.1cm,height=2.1cm,valign=center,colback=ctri!6,colframe=ctri,arc=2.5mm,boxrule=0.9pt,left=3mm]{\Large\bfseries\color{ctri}d)}\quad\begin{minipage}[c]{5.3cm}\raggedright\Large $\dfrac{3\pi}{4}$\end{minipage}\end{tcolorbox}
\end{tcolorbox}
\newpage
\begin{tcolorbox}[enhanced,width=16cm,height=14cm,colback=white,colframe=ctri,boxrule=2pt,arc=0mm,
left=6mm,right=6mm,top=0mm,bottom=4mm,sharp corners,
overlay={\fill[ctri] (frame.north west) rectangle ([yshift=-1.25cm]frame.north east);
\node[anchor=west,text=white,font=\sffamily\bfseries\LARGE] at ([xshift=5mm,yshift=-0.62cm]frame.north west) {Question 60 C};
\node[anchor=east,text=white,font=\sffamily\itshape\large] at ([xshift=-5mm,yshift=-0.62cm]frame.north east) {Trigonométrie et exponentielle};}]
\vspace*{1.75cm}
\begin{minipage}[t][6.9cm][c]{\linewidth}\LARGE\centering
\begin{minipage}{0.95\linewidth}\raggedright\linespread{1.25}\selectfont Quelle est la mesure principale de l'angle $\dfrac{25\pi}{6}$~?\end{minipage}
\end{minipage}

\vfill
\begin{tcolorbox}[enhanced,nobeforeafter,width=7.1cm,height=2.1cm,valign=center,colback=ctri!6,colframe=ctri,arc=2.5mm,boxrule=0.9pt,left=3mm]{\Large\bfseries\color{ctri}a)}\quad\begin{minipage}[c]{5.3cm}\raggedright\Large $\dfrac{\pi}{12}$\end{minipage}\end{tcolorbox}\hfill\begin{tcolorbox}[enhanced,nobeforeafter,width=7.1cm,height=2.1cm,valign=center,colback=ctri!6,colframe=ctri,arc=2.5mm,boxrule=0.9pt,left=3mm]{\Large\bfseries\color{ctri}b)}\quad\begin{minipage}[c]{5.3cm}\raggedright\Large $-\dfrac{5\pi}{6}$\end{minipage}\end{tcolorbox}\par\vspace{2.5mm}
\begin{tcolorbox}[enhanced,nobeforeafter,width=7.1cm,height=2.1cm,valign=center,colback=ctri!6,colframe=ctri,arc=2.5mm,boxrule=0.9pt,left=3mm]{\Large\bfseries\color{ctri}c)}\quad\begin{minipage}[c]{5.3cm}\raggedright\Large $\dfrac{\pi}{6}$\end{minipage}\end{tcolorbox}\hfill\begin{tcolorbox}[enhanced,nobeforeafter,width=7.1cm,height=2.1cm,valign=center,colback=ctri!6,colframe=ctri,arc=2.5mm,boxrule=0.9pt,left=3mm]{\Large\bfseries\color{ctri}d)}\quad\begin{minipage}[c]{5.3cm}\raggedright\Large $-\dfrac{\pi}{6}$\end{minipage}\end{tcolorbox}
\end{tcolorbox}
\newpage
\begin{tcolorbox}[enhanced,width=16cm,height=14cm,colback=white,colframe=ctri,boxrule=2pt,arc=0mm,
left=6mm,right=6mm,top=0mm,bottom=4mm,sharp corners,
overlay={\fill[ctri] (frame.north west) rectangle ([yshift=-1.25cm]frame.north east);
\node[anchor=west,text=white,font=\sffamily\bfseries\LARGE] at ([xshift=5mm,yshift=-0.62cm]frame.north west) {Question 61 C};
\node[anchor=east,text=white,font=\sffamily\itshape\large] at ([xshift=-5mm,yshift=-0.62cm]frame.north east) {Trigonométrie et exponentielle};}]
\vspace*{1.75cm}
\begin{minipage}[t][6.9cm][c]{\linewidth}\LARGE\centering
\begin{minipage}{0.95\linewidth}\raggedright\linespread{1.25}\selectfont On sait que $x\in[0\,;\,\pi]$ et $\cos x=\dfrac{1}{2}$. Que vaut $\sin x$~?\end{minipage}
\end{minipage}

\vfill
\begin{tcolorbox}[enhanced,nobeforeafter,width=7.1cm,height=2.1cm,valign=center,colback=ctri!6,colframe=ctri,arc=2.5mm,boxrule=0.9pt,left=3mm]{\Large\bfseries\color{ctri}a)}\quad\begin{minipage}[c]{5.3cm}\raggedright\Large $\dfrac{1}{2}$\end{minipage}\end{tcolorbox}\hfill\begin{tcolorbox}[enhanced,nobeforeafter,width=7.1cm,height=2.1cm,valign=center,colback=ctri!6,colframe=ctri,arc=2.5mm,boxrule=0.9pt,left=3mm]{\Large\bfseries\color{ctri}b)}\quad\begin{minipage}[c]{5.3cm}\raggedright\Large $-\dfrac{\sqrt{3}}{2}$\end{minipage}\end{tcolorbox}\par\vspace{2.5mm}
\begin{tcolorbox}[enhanced,nobeforeafter,width=7.1cm,height=2.1cm,valign=center,colback=ctri!6,colframe=ctri,arc=2.5mm,boxrule=0.9pt,left=3mm]{\Large\bfseries\color{ctri}c)}\quad\begin{minipage}[c]{5.3cm}\raggedright\Large $\dfrac{\sqrt{3}}{2}$\end{minipage}\end{tcolorbox}\hfill\begin{tcolorbox}[enhanced,nobeforeafter,width=7.1cm,height=2.1cm,valign=center,colback=ctri!6,colframe=ctri,arc=2.5mm,boxrule=0.9pt,left=3mm]{\Large\bfseries\color{ctri}d)}\quad\begin{minipage}[c]{5.3cm}\raggedright\Large $\dfrac{3}{4}$\end{minipage}\end{tcolorbox}
\end{tcolorbox}
\newpage
\begin{tcolorbox}[enhanced,width=16cm,height=14cm,colback=white,colframe=ctri,boxrule=2pt,arc=0mm,
left=6mm,right=6mm,top=0mm,bottom=4mm,sharp corners,
overlay={\fill[ctri] (frame.north west) rectangle ([yshift=-1.25cm]frame.north east);
\node[anchor=west,text=white,font=\sffamily\bfseries\LARGE] at ([xshift=5mm,yshift=-0.62cm]frame.north west) {Question 62 C};
\node[anchor=east,text=white,font=\sffamily\itshape\large] at ([xshift=-5mm,yshift=-0.62cm]frame.north east) {Trigonométrie et exponentielle};}]
\vspace*{1.75cm}
\begin{minipage}[t][6.9cm][c]{\linewidth}\LARGE\centering
\begin{minipage}{0.95\linewidth}\raggedright\linespread{1.25}\selectfont On sait que $x\in[0\,;\,\pi]$ et $\cos x=-\dfrac{7}{25}$. Que vaut $\sin x$~?\end{minipage}
\end{minipage}

\vfill
\begin{tcolorbox}[enhanced,nobeforeafter,width=7.1cm,height=2.1cm,valign=center,colback=ctri!6,colframe=ctri,arc=2.5mm,boxrule=0.9pt,left=3mm]{\Large\bfseries\color{ctri}a)}\quad\begin{minipage}[c]{5.3cm}\raggedright\Large $-\dfrac{24}{25}$\end{minipage}\end{tcolorbox}\hfill\begin{tcolorbox}[enhanced,nobeforeafter,width=7.1cm,height=2.1cm,valign=center,colback=ctri!6,colframe=ctri,arc=2.5mm,boxrule=0.9pt,left=3mm]{\Large\bfseries\color{ctri}b)}\quad\begin{minipage}[c]{5.3cm}\raggedright\Large $\dfrac{576}{625}$\end{minipage}\end{tcolorbox}\par\vspace{2.5mm}
\begin{tcolorbox}[enhanced,nobeforeafter,width=7.1cm,height=2.1cm,valign=center,colback=ctri!6,colframe=ctri,arc=2.5mm,boxrule=0.9pt,left=3mm]{\Large\bfseries\color{ctri}c)}\quad\begin{minipage}[c]{5.3cm}\raggedright\Large $\dfrac{24}{25}$\end{minipage}\end{tcolorbox}\hfill\begin{tcolorbox}[enhanced,nobeforeafter,width=7.1cm,height=2.1cm,valign=center,colback=ctri!6,colframe=ctri,arc=2.5mm,boxrule=0.9pt,left=3mm]{\Large\bfseries\color{ctri}d)}\quad\begin{minipage}[c]{5.3cm}\raggedright\Large $\dfrac{32}{25}$\end{minipage}\end{tcolorbox}
\end{tcolorbox}
\newpage
\begin{tcolorbox}[enhanced,width=16cm,height=14cm,colback=white,colframe=ctri,boxrule=2pt,arc=0mm,
left=6mm,right=6mm,top=0mm,bottom=4mm,sharp corners,
overlay={\fill[ctri] (frame.north west) rectangle ([yshift=-1.25cm]frame.north east);
\node[anchor=west,text=white,font=\sffamily\bfseries\LARGE] at ([xshift=5mm,yshift=-0.62cm]frame.north west) {Question 63 C};
\node[anchor=east,text=white,font=\sffamily\itshape\large] at ([xshift=-5mm,yshift=-0.62cm]frame.north east) {Trigonométrie et exponentielle};}]
\vspace*{1.75cm}
\begin{minipage}[t][6.9cm][c]{\linewidth}\LARGE\centering
\begin{minipage}{0.95\linewidth}\raggedright\linespread{1.25}\selectfont Simplifier $\left(e^{4}\right)^{-1}\times e^{2}$.\end{minipage}
\end{minipage}

\vfill
\begin{tcolorbox}[enhanced,nobeforeafter,width=7.1cm,height=2.1cm,valign=center,colback=ctri!6,colframe=ctri,arc=2.5mm,boxrule=0.9pt,left=3mm]{\Large\bfseries\color{ctri}a)}\quad\begin{minipage}[c]{5.3cm}\raggedright\Large $e^{-2}$\end{minipage}\end{tcolorbox}\hfill\begin{tcolorbox}[enhanced,nobeforeafter,width=7.1cm,height=2.1cm,valign=center,colback=ctri!6,colframe=ctri,arc=2.5mm,boxrule=0.9pt,left=3mm]{\Large\bfseries\color{ctri}b)}\quad\begin{minipage}[c]{5.3cm}\raggedright\Large $e^{5}$\end{minipage}\end{tcolorbox}\par\vspace{2.5mm}
\begin{tcolorbox}[enhanced,nobeforeafter,width=7.1cm,height=2.1cm,valign=center,colback=ctri!6,colframe=ctri,arc=2.5mm,boxrule=0.9pt,left=3mm]{\Large\bfseries\color{ctri}c)}\quad\begin{minipage}[c]{5.3cm}\raggedright\Large $e^{2.25}$\end{minipage}\end{tcolorbox}\hfill\begin{tcolorbox}[enhanced,nobeforeafter,width=7.1cm,height=2.1cm,valign=center,colback=ctri!6,colframe=ctri,arc=2.5mm,boxrule=0.9pt,left=3mm]{\Large\bfseries\color{ctri}d)}\quad\begin{minipage}[c]{5.3cm}\raggedright\Large $e^{4}$\end{minipage}\end{tcolorbox}
\end{tcolorbox}
\newpage
\begin{tcolorbox}[enhanced,width=16cm,height=14cm,colback=white,colframe=ctri,boxrule=2pt,arc=0mm,
left=6mm,right=6mm,top=0mm,bottom=4mm,sharp corners,
overlay={\fill[ctri] (frame.north west) rectangle ([yshift=-1.25cm]frame.north east);
\node[anchor=west,text=white,font=\sffamily\bfseries\LARGE] at ([xshift=5mm,yshift=-0.62cm]frame.north west) {Question 64 C};
\node[anchor=east,text=white,font=\sffamily\itshape\large] at ([xshift=-5mm,yshift=-0.62cm]frame.north east) {Trigonométrie et exponentielle};}]
\vspace*{1.75cm}
\begin{minipage}[t][6.9cm][c]{\linewidth}\LARGE\centering
\begin{minipage}{0.95\linewidth}\raggedright\linespread{1.25}\selectfont Simplifier $\left(e^{-2}\right)^{-3}\times e^{1}$.\end{minipage}
\end{minipage}

\vfill
\begin{tcolorbox}[enhanced,nobeforeafter,width=7.1cm,height=2.1cm,valign=center,colback=ctri!6,colframe=ctri,arc=2.5mm,boxrule=0.9pt,left=3mm]{\Large\bfseries\color{ctri}a)}\quad\begin{minipage}[c]{5.3cm}\raggedright\Large $e^{-4}$\end{minipage}\end{tcolorbox}\hfill\begin{tcolorbox}[enhanced,nobeforeafter,width=7.1cm,height=2.1cm,valign=center,colback=ctri!6,colframe=ctri,arc=2.5mm,boxrule=0.9pt,left=3mm]{\Large\bfseries\color{ctri}b)}\quad\begin{minipage}[c]{5.3cm}\raggedright\Large $e^{4}$\end{minipage}\end{tcolorbox}\par\vspace{2.5mm}
\begin{tcolorbox}[enhanced,nobeforeafter,width=7.1cm,height=2.1cm,valign=center,colback=ctri!6,colframe=ctri,arc=2.5mm,boxrule=0.9pt,left=3mm]{\Large\bfseries\color{ctri}c)}\quad\begin{minipage}[c]{5.3cm}\raggedright\Large $e^{0.875}$\end{minipage}\end{tcolorbox}\hfill\begin{tcolorbox}[enhanced,nobeforeafter,width=7.1cm,height=2.1cm,valign=center,colback=ctri!6,colframe=ctri,arc=2.5mm,boxrule=0.9pt,left=3mm]{\Large\bfseries\color{ctri}d)}\quad\begin{minipage}[c]{5.3cm}\raggedright\Large $e^{7}$\end{minipage}\end{tcolorbox}
\end{tcolorbox}
\newpage
\begin{tcolorbox}[enhanced,width=16cm,height=14cm,colback=white,colframe=cpsc,boxrule=2pt,arc=0mm,
left=6mm,right=6mm,top=0mm,bottom=4mm,sharp corners,
overlay={\fill[cpsc] (frame.north west) rectangle ([yshift=-1.25cm]frame.north east);
\node[anchor=west,text=white,font=\sffamily\bfseries\LARGE] at ([xshift=5mm,yshift=-0.62cm]frame.north west) {Question 65 C};
\node[anchor=east,text=white,font=\sffamily\itshape\large] at ([xshift=-5mm,yshift=-0.62cm]frame.north east) {Produit scalaire et géométrie repérée};}]
\vspace*{1.75cm}
\begin{minipage}[t][6.9cm][c]{\linewidth}\Large\centering
\begin{minipage}{0.95\linewidth}\raggedright\linespread{1.25}\selectfont Dans un repère orthonormé, que vaut $\vec u\cdot\vec v$ pour $\vec u\begin{pmatrix}-3\\-4\end{pmatrix}$ et $\vec v\begin{pmatrix}2\\-1\end{pmatrix}$~?\end{minipage}
\end{minipage}

\vfill
\begin{tcolorbox}[enhanced,nobeforeafter,width=7.1cm,height=2.1cm,valign=center,colback=cpsc!6,colframe=cpsc,arc=2.5mm,boxrule=0.9pt,left=3mm]{\Large\bfseries\color{cpsc}a)}\quad\begin{minipage}[c]{5.3cm}\raggedright\Large $-5$\end{minipage}\end{tcolorbox}\hfill\begin{tcolorbox}[enhanced,nobeforeafter,width=7.1cm,height=2.1cm,valign=center,colback=cpsc!6,colframe=cpsc,arc=2.5mm,boxrule=0.9pt,left=3mm]{\Large\bfseries\color{cpsc}b)}\quad\begin{minipage}[c]{5.3cm}\raggedright\Large $-2$\end{minipage}\end{tcolorbox}\par\vspace{2.5mm}
\begin{tcolorbox}[enhanced,nobeforeafter,width=7.1cm,height=2.1cm,valign=center,colback=cpsc!6,colframe=cpsc,arc=2.5mm,boxrule=0.9pt,left=3mm]{\Large\bfseries\color{cpsc}c)}\quad\begin{minipage}[c]{5.3cm}\raggedright\Large $11$\end{minipage}\end{tcolorbox}\hfill\begin{tcolorbox}[enhanced,nobeforeafter,width=7.1cm,height=2.1cm,valign=center,colback=cpsc!6,colframe=cpsc,arc=2.5mm,boxrule=0.9pt,left=3mm]{\Large\bfseries\color{cpsc}d)}\quad\begin{minipage}[c]{5.3cm}\raggedright\Large $-10$\end{minipage}\end{tcolorbox}
\end{tcolorbox}
\newpage
\begin{tcolorbox}[enhanced,width=16cm,height=14cm,colback=white,colframe=cpsc,boxrule=2pt,arc=0mm,
left=6mm,right=6mm,top=0mm,bottom=4mm,sharp corners,
overlay={\fill[cpsc] (frame.north west) rectangle ([yshift=-1.25cm]frame.north east);
\node[anchor=west,text=white,font=\sffamily\bfseries\LARGE] at ([xshift=5mm,yshift=-0.62cm]frame.north west) {Question 66 C};
\node[anchor=east,text=white,font=\sffamily\itshape\large] at ([xshift=-5mm,yshift=-0.62cm]frame.north east) {Produit scalaire et géométrie repérée};}]
\vspace*{1.75cm}
\begin{minipage}[t][6.9cm][c]{\linewidth}\LARGE\centering
\begin{minipage}{0.95\linewidth}\raggedright\linespread{1.25}\selectfont Dans un repère orthonormé, que vaut $\vec u\cdot\vec v$ pour $\vec u\begin{pmatrix}5\\2\end{pmatrix}$ et $\vec v\begin{pmatrix}1\\3\end{pmatrix}$~?\end{minipage}
\end{minipage}

\vfill
\begin{tcolorbox}[enhanced,nobeforeafter,width=7.1cm,height=2.1cm,valign=center,colback=cpsc!6,colframe=cpsc,arc=2.5mm,boxrule=0.9pt,left=3mm]{\Large\bfseries\color{cpsc}a)}\quad\begin{minipage}[c]{5.3cm}\raggedright\Large $-1$\end{minipage}\end{tcolorbox}\hfill\begin{tcolorbox}[enhanced,nobeforeafter,width=7.1cm,height=2.1cm,valign=center,colback=cpsc!6,colframe=cpsc,arc=2.5mm,boxrule=0.9pt,left=3mm]{\Large\bfseries\color{cpsc}b)}\quad\begin{minipage}[c]{5.3cm}\raggedright\Large $13$\end{minipage}\end{tcolorbox}\par\vspace{2.5mm}
\begin{tcolorbox}[enhanced,nobeforeafter,width=7.1cm,height=2.1cm,valign=center,colback=cpsc!6,colframe=cpsc,arc=2.5mm,boxrule=0.9pt,left=3mm]{\Large\bfseries\color{cpsc}c)}\quad\begin{minipage}[c]{5.3cm}\raggedright\Large $11$\end{minipage}\end{tcolorbox}\hfill\begin{tcolorbox}[enhanced,nobeforeafter,width=7.1cm,height=2.1cm,valign=center,colback=cpsc!6,colframe=cpsc,arc=2.5mm,boxrule=0.9pt,left=3mm]{\Large\bfseries\color{cpsc}d)}\quad\begin{minipage}[c]{5.3cm}\raggedright\Large $17$\end{minipage}\end{tcolorbox}
\end{tcolorbox}
\newpage
\begin{tcolorbox}[enhanced,width=16cm,height=14cm,colback=white,colframe=cpsc,boxrule=2pt,arc=0mm,
left=6mm,right=6mm,top=0mm,bottom=4mm,sharp corners,
overlay={\fill[cpsc] (frame.north west) rectangle ([yshift=-1.25cm]frame.north east);
\node[anchor=west,text=white,font=\sffamily\bfseries\LARGE] at ([xshift=5mm,yshift=-0.62cm]frame.north west) {Question 67 C};
\node[anchor=east,text=white,font=\sffamily\itshape\large] at ([xshift=-5mm,yshift=-0.62cm]frame.north east) {Produit scalaire et géométrie repérée};}]
\vspace*{1.75cm}
\begin{minipage}[t][6.9cm][c]{\linewidth}\LARGE\centering
\begin{minipage}{0.95\linewidth}\raggedright\linespread{1.25}\selectfont $\|\vec u\|=2$, $\|\vec v\|=3$ et l'angle entre $\vec u$ et $\vec v$ mesure $180^\circ$. Que vaut $\vec u\cdot\vec v$~?\end{minipage}
\end{minipage}

\vfill
\begin{tcolorbox}[enhanced,nobeforeafter,width=7.1cm,height=2.1cm,valign=center,colback=cpsc!6,colframe=cpsc,arc=2.5mm,boxrule=0.9pt,left=3mm]{\Large\bfseries\color{cpsc}a)}\quad\begin{minipage}[c]{5.3cm}\raggedright\Large $-6$\end{minipage}\end{tcolorbox}\hfill\begin{tcolorbox}[enhanced,nobeforeafter,width=7.1cm,height=2.1cm,valign=center,colback=cpsc!6,colframe=cpsc,arc=2.5mm,boxrule=0.9pt,left=3mm]{\Large\bfseries\color{cpsc}b)}\quad\begin{minipage}[c]{5.3cm}\raggedright\Large $6$\end{minipage}\end{tcolorbox}\par\vspace{2.5mm}
\begin{tcolorbox}[enhanced,nobeforeafter,width=7.1cm,height=2.1cm,valign=center,colback=cpsc!6,colframe=cpsc,arc=2.5mm,boxrule=0.9pt,left=3mm]{\Large\bfseries\color{cpsc}c)}\quad\begin{minipage}[c]{5.3cm}\raggedright\Large $-5$\end{minipage}\end{tcolorbox}\hfill\begin{tcolorbox}[enhanced,nobeforeafter,width=7.1cm,height=2.1cm,valign=center,colback=cpsc!6,colframe=cpsc,arc=2.5mm,boxrule=0.9pt,left=3mm]{\Large\bfseries\color{cpsc}d)}\quad\begin{minipage}[c]{5.3cm}\raggedright\Large $5$\end{minipage}\end{tcolorbox}
\end{tcolorbox}
\newpage
\begin{tcolorbox}[enhanced,width=16cm,height=14cm,colback=white,colframe=cpsc,boxrule=2pt,arc=0mm,
left=6mm,right=6mm,top=0mm,bottom=4mm,sharp corners,
overlay={\fill[cpsc] (frame.north west) rectangle ([yshift=-1.25cm]frame.north east);
\node[anchor=west,text=white,font=\sffamily\bfseries\LARGE] at ([xshift=5mm,yshift=-0.62cm]frame.north west) {Question 68 C};
\node[anchor=east,text=white,font=\sffamily\itshape\large] at ([xshift=-5mm,yshift=-0.62cm]frame.north east) {Produit scalaire et géométrie repérée};}]
\vspace*{1.75cm}
\begin{minipage}[t][6.9cm][c]{\linewidth}\LARGE\centering
\begin{minipage}{0.95\linewidth}\raggedright\linespread{1.25}\selectfont $\|\vec u\|=5$, $\|\vec v\|=2$ et l'angle entre $\vec u$ et $\vec v$ mesure $60^\circ$. Que vaut $\vec u\cdot\vec v$~?\end{minipage}
\end{minipage}

\vfill
\begin{tcolorbox}[enhanced,nobeforeafter,width=7.1cm,height=2.1cm,valign=center,colback=cpsc!6,colframe=cpsc,arc=2.5mm,boxrule=0.9pt,left=3mm]{\Large\bfseries\color{cpsc}a)}\quad\begin{minipage}[c]{5.3cm}\raggedright\Large $7$\end{minipage}\end{tcolorbox}\hfill\begin{tcolorbox}[enhanced,nobeforeafter,width=7.1cm,height=2.1cm,valign=center,colback=cpsc!6,colframe=cpsc,arc=2.5mm,boxrule=0.9pt,left=3mm]{\Large\bfseries\color{cpsc}b)}\quad\begin{minipage}[c]{5.3cm}\raggedright\Large $-5$\end{minipage}\end{tcolorbox}\par\vspace{2.5mm}
\begin{tcolorbox}[enhanced,nobeforeafter,width=7.1cm,height=2.1cm,valign=center,colback=cpsc!6,colframe=cpsc,arc=2.5mm,boxrule=0.9pt,left=3mm]{\Large\bfseries\color{cpsc}c)}\quad\begin{minipage}[c]{5.3cm}\raggedright\Large $5$\end{minipage}\end{tcolorbox}\hfill\begin{tcolorbox}[enhanced,nobeforeafter,width=7.1cm,height=2.1cm,valign=center,colback=cpsc!6,colframe=cpsc,arc=2.5mm,boxrule=0.9pt,left=3mm]{\Large\bfseries\color{cpsc}d)}\quad\begin{minipage}[c]{5.3cm}\raggedright\Large $10$\end{minipage}\end{tcolorbox}
\end{tcolorbox}
\newpage
\begin{tcolorbox}[enhanced,width=16cm,height=14cm,colback=white,colframe=cpsc,boxrule=2pt,arc=0mm,
left=6mm,right=6mm,top=0mm,bottom=4mm,sharp corners,
overlay={\fill[cpsc] (frame.north west) rectangle ([yshift=-1.25cm]frame.north east);
\node[anchor=west,text=white,font=\sffamily\bfseries\LARGE] at ([xshift=5mm,yshift=-0.62cm]frame.north west) {Question 69 C};
\node[anchor=east,text=white,font=\sffamily\itshape\large] at ([xshift=-5mm,yshift=-0.62cm]frame.north east) {Produit scalaire et géométrie repérée};}]
\vspace*{1.75cm}
\begin{minipage}[t][6.9cm][c]{\linewidth}\LARGE\centering
\begin{minipage}{0.95\linewidth}\raggedright\linespread{1.25}\selectfont Pour quelle valeur de $m$ les vecteurs $\vec u\begin{pmatrix}1\\m\end{pmatrix}$ et $\vec v\begin{pmatrix}-3\\6\end{pmatrix}$ sont-ils orthogonaux~?\end{minipage}
\end{minipage}

\vfill
\begin{tcolorbox}[enhanced,nobeforeafter,width=7.1cm,height=2.1cm,valign=center,colback=cpsc!6,colframe=cpsc,arc=2.5mm,boxrule=0.9pt,left=3mm]{\Large\bfseries\color{cpsc}a)}\quad\begin{minipage}[c]{5.3cm}\raggedright\Large $m=\dfrac{1}{2}$\end{minipage}\end{tcolorbox}\hfill\begin{tcolorbox}[enhanced,nobeforeafter,width=7.1cm,height=2.1cm,valign=center,colback=cpsc!6,colframe=cpsc,arc=2.5mm,boxrule=0.9pt,left=3mm]{\Large\bfseries\color{cpsc}b)}\quad\begin{minipage}[c]{5.3cm}\raggedright\Large $m=-\dfrac{1}{2}$\end{minipage}\end{tcolorbox}\par\vspace{2.5mm}
\begin{tcolorbox}[enhanced,nobeforeafter,width=7.1cm,height=2.1cm,valign=center,colback=cpsc!6,colframe=cpsc,arc=2.5mm,boxrule=0.9pt,left=3mm]{\Large\bfseries\color{cpsc}c)}\quad\begin{minipage}[c]{5.3cm}\raggedright\Large $m=2$\end{minipage}\end{tcolorbox}\hfill\begin{tcolorbox}[enhanced,nobeforeafter,width=7.1cm,height=2.1cm,valign=center,colback=cpsc!6,colframe=cpsc,arc=2.5mm,boxrule=0.9pt,left=3mm]{\Large\bfseries\color{cpsc}d)}\quad\begin{minipage}[c]{5.3cm}\raggedright\Large $m=3$\end{minipage}\end{tcolorbox}
\end{tcolorbox}
\newpage
\begin{tcolorbox}[enhanced,width=16cm,height=14cm,colback=white,colframe=cpsc,boxrule=2pt,arc=0mm,
left=6mm,right=6mm,top=0mm,bottom=4mm,sharp corners,
overlay={\fill[cpsc] (frame.north west) rectangle ([yshift=-1.25cm]frame.north east);
\node[anchor=west,text=white,font=\sffamily\bfseries\LARGE] at ([xshift=5mm,yshift=-0.62cm]frame.north west) {Question 70 C};
\node[anchor=east,text=white,font=\sffamily\itshape\large] at ([xshift=-5mm,yshift=-0.62cm]frame.north east) {Produit scalaire et géométrie repérée};}]
\vspace*{1.75cm}
\begin{minipage}[t][6.9cm][c]{\linewidth}\LARGE\centering
\begin{minipage}{0.95\linewidth}\raggedright\linespread{1.25}\selectfont Pour quelle valeur de $m$ les vecteurs $\vec u\begin{pmatrix}6\\m\end{pmatrix}$ et $\vec v\begin{pmatrix}-1\\3\end{pmatrix}$ sont-ils orthogonaux~?\end{minipage}
\end{minipage}

\vfill
\begin{tcolorbox}[enhanced,nobeforeafter,width=7.1cm,height=2.1cm,valign=center,colback=cpsc!6,colframe=cpsc,arc=2.5mm,boxrule=0.9pt,left=3mm]{\Large\bfseries\color{cpsc}a)}\quad\begin{minipage}[c]{5.3cm}\raggedright\Large $m=6$\end{minipage}\end{tcolorbox}\hfill\begin{tcolorbox}[enhanced,nobeforeafter,width=7.1cm,height=2.1cm,valign=center,colback=cpsc!6,colframe=cpsc,arc=2.5mm,boxrule=0.9pt,left=3mm]{\Large\bfseries\color{cpsc}b)}\quad\begin{minipage}[c]{5.3cm}\raggedright\Large $m=\dfrac{1}{2}$\end{minipage}\end{tcolorbox}\par\vspace{2.5mm}
\begin{tcolorbox}[enhanced,nobeforeafter,width=7.1cm,height=2.1cm,valign=center,colback=cpsc!6,colframe=cpsc,arc=2.5mm,boxrule=0.9pt,left=3mm]{\Large\bfseries\color{cpsc}c)}\quad\begin{minipage}[c]{5.3cm}\raggedright\Large $m=-2$\end{minipage}\end{tcolorbox}\hfill\begin{tcolorbox}[enhanced,nobeforeafter,width=7.1cm,height=2.1cm,valign=center,colback=cpsc!6,colframe=cpsc,arc=2.5mm,boxrule=0.9pt,left=3mm]{\Large\bfseries\color{cpsc}d)}\quad\begin{minipage}[c]{5.3cm}\raggedright\Large $m=2$\end{minipage}\end{tcolorbox}
\end{tcolorbox}
\newpage
\begin{tcolorbox}[enhanced,width=16cm,height=14cm,colback=white,colframe=cpsc,boxrule=2pt,arc=0mm,
left=6mm,right=6mm,top=0mm,bottom=4mm,sharp corners,
overlay={\fill[cpsc] (frame.north west) rectangle ([yshift=-1.25cm]frame.north east);
\node[anchor=west,text=white,font=\sffamily\bfseries\LARGE] at ([xshift=5mm,yshift=-0.62cm]frame.north west) {Question 71 C};
\node[anchor=east,text=white,font=\sffamily\itshape\large] at ([xshift=-5mm,yshift=-0.62cm]frame.north east) {Produit scalaire et géométrie repérée};}]
\vspace*{1.75cm}
\begin{minipage}[t][6.9cm][c]{\linewidth}\LARGE\centering
\begin{minipage}{0.95\linewidth}\raggedright\linespread{1.25}\selectfont Dans le triangle $ABC$, $AB=3$, $AC=3$ et $\widehat{BAC}=120^\circ$. Combien mesure $BC$~?\end{minipage}
\end{minipage}

\vfill
\begin{tcolorbox}[enhanced,nobeforeafter,width=7.1cm,height=2.1cm,valign=center,colback=cpsc!6,colframe=cpsc,arc=2.5mm,boxrule=0.9pt,left=3mm]{\Large\bfseries\color{cpsc}a)}\quad\begin{minipage}[c]{5.3cm}\raggedright\Large $3\sqrt{2}$\end{minipage}\end{tcolorbox}\hfill\begin{tcolorbox}[enhanced,nobeforeafter,width=7.1cm,height=2.1cm,valign=center,colback=cpsc!6,colframe=cpsc,arc=2.5mm,boxrule=0.9pt,left=3mm]{\Large\bfseries\color{cpsc}b)}\quad\begin{minipage}[c]{5.3cm}\raggedright\Large $27$\end{minipage}\end{tcolorbox}\par\vspace{2.5mm}
\begin{tcolorbox}[enhanced,nobeforeafter,width=7.1cm,height=2.1cm,valign=center,colback=cpsc!6,colframe=cpsc,arc=2.5mm,boxrule=0.9pt,left=3mm]{\Large\bfseries\color{cpsc}c)}\quad\begin{minipage}[c]{5.3cm}\raggedright\Large $3\sqrt{3}$\end{minipage}\end{tcolorbox}\hfill\begin{tcolorbox}[enhanced,nobeforeafter,width=7.1cm,height=2.1cm,valign=center,colback=cpsc!6,colframe=cpsc,arc=2.5mm,boxrule=0.9pt,left=3mm]{\Large\bfseries\color{cpsc}d)}\quad\begin{minipage}[c]{5.3cm}\raggedright\Large $3$\end{minipage}\end{tcolorbox}
\end{tcolorbox}
\newpage
\begin{tcolorbox}[enhanced,width=16cm,height=14cm,colback=white,colframe=cpsc,boxrule=2pt,arc=0mm,
left=6mm,right=6mm,top=0mm,bottom=4mm,sharp corners,
overlay={\fill[cpsc] (frame.north west) rectangle ([yshift=-1.25cm]frame.north east);
\node[anchor=west,text=white,font=\sffamily\bfseries\LARGE] at ([xshift=5mm,yshift=-0.62cm]frame.north west) {Question 72 C};
\node[anchor=east,text=white,font=\sffamily\itshape\large] at ([xshift=-5mm,yshift=-0.62cm]frame.north east) {Produit scalaire et géométrie repérée};}]
\vspace*{1.75cm}
\begin{minipage}[t][6.9cm][c]{\linewidth}\LARGE\centering
\begin{minipage}{0.95\linewidth}\raggedright\linespread{1.25}\selectfont Dans le triangle $ABC$, $AB=6$, $AC=10$ et $\widehat{BAC}=120^\circ$. Combien mesure $BC$~?\end{minipage}
\end{minipage}

\vfill
\begin{tcolorbox}[enhanced,nobeforeafter,width=7.1cm,height=2.1cm,valign=center,colback=cpsc!6,colframe=cpsc,arc=2.5mm,boxrule=0.9pt,left=3mm]{\Large\bfseries\color{cpsc}a)}\quad\begin{minipage}[c]{5.3cm}\raggedright\Large $196$\end{minipage}\end{tcolorbox}\hfill\begin{tcolorbox}[enhanced,nobeforeafter,width=7.1cm,height=2.1cm,valign=center,colback=cpsc!6,colframe=cpsc,arc=2.5mm,boxrule=0.9pt,left=3mm]{\Large\bfseries\color{cpsc}b)}\quad\begin{minipage}[c]{5.3cm}\raggedright\Large $14$\end{minipage}\end{tcolorbox}\par\vspace{2.5mm}
\begin{tcolorbox}[enhanced,nobeforeafter,width=7.1cm,height=2.1cm,valign=center,colback=cpsc!6,colframe=cpsc,arc=2.5mm,boxrule=0.9pt,left=3mm]{\Large\bfseries\color{cpsc}c)}\quad\begin{minipage}[c]{5.3cm}\raggedright\Large $2\sqrt{34}$\end{minipage}\end{tcolorbox}\hfill\begin{tcolorbox}[enhanced,nobeforeafter,width=7.1cm,height=2.1cm,valign=center,colback=cpsc!6,colframe=cpsc,arc=2.5mm,boxrule=0.9pt,left=3mm]{\Large\bfseries\color{cpsc}d)}\quad\begin{minipage}[c]{5.3cm}\raggedright\Large $2\sqrt{19}$\end{minipage}\end{tcolorbox}
\end{tcolorbox}
\newpage
\begin{tcolorbox}[enhanced,width=16cm,height=14cm,colback=white,colframe=cpsc,boxrule=2pt,arc=0mm,
left=6mm,right=6mm,top=0mm,bottom=4mm,sharp corners,
overlay={\fill[cpsc] (frame.north west) rectangle ([yshift=-1.25cm]frame.north east);
\node[anchor=west,text=white,font=\sffamily\bfseries\LARGE] at ([xshift=5mm,yshift=-0.62cm]frame.north west) {Question 73 C};
\node[anchor=east,text=white,font=\sffamily\itshape\large] at ([xshift=-5mm,yshift=-0.62cm]frame.north east) {Produit scalaire et géométrie repérée};}]
\vspace*{1.75cm}
\begin{minipage}[t][6.9cm][c]{\linewidth}\LARGE\centering
\begin{minipage}{0.95\linewidth}\raggedright\linespread{1.25}\selectfont Lequel de ces vecteurs est un vecteur normal à la droite d'équation $4x-y=0$~?\end{minipage}
\end{minipage}

\vfill
\begin{tcolorbox}[enhanced,nobeforeafter,width=7.1cm,height=2.1cm,valign=center,colback=cpsc!6,colframe=cpsc,arc=2.5mm,boxrule=0.9pt,left=3mm]{\Large\bfseries\color{cpsc}a)}\quad\begin{minipage}[c]{5.3cm}\raggedright\large $\begin{pmatrix}-1\\4\end{pmatrix}$\end{minipage}\end{tcolorbox}\hfill\begin{tcolorbox}[enhanced,nobeforeafter,width=7.1cm,height=2.1cm,valign=center,colback=cpsc!6,colframe=cpsc,arc=2.5mm,boxrule=0.9pt,left=3mm]{\Large\bfseries\color{cpsc}b)}\quad\begin{minipage}[c]{5.3cm}\raggedright\large $\begin{pmatrix}4\\-1\end{pmatrix}$\end{minipage}\end{tcolorbox}\par\vspace{2.5mm}
\begin{tcolorbox}[enhanced,nobeforeafter,width=7.1cm,height=2.1cm,valign=center,colback=cpsc!6,colframe=cpsc,arc=2.5mm,boxrule=0.9pt,left=3mm]{\Large\bfseries\color{cpsc}c)}\quad\begin{minipage}[c]{5.3cm}\raggedright\large $\begin{pmatrix}4\\0\end{pmatrix}$\end{minipage}\end{tcolorbox}\hfill\begin{tcolorbox}[enhanced,nobeforeafter,width=7.1cm,height=2.1cm,valign=center,colback=cpsc!6,colframe=cpsc,arc=2.5mm,boxrule=0.9pt,left=3mm]{\Large\bfseries\color{cpsc}d)}\quad\begin{minipage}[c]{5.3cm}\raggedright\large $\begin{pmatrix}1\\4\end{pmatrix}$\end{minipage}\end{tcolorbox}
\end{tcolorbox}
\newpage
\begin{tcolorbox}[enhanced,width=16cm,height=14cm,colback=white,colframe=cpsc,boxrule=2pt,arc=0mm,
left=6mm,right=6mm,top=0mm,bottom=4mm,sharp corners,
overlay={\fill[cpsc] (frame.north west) rectangle ([yshift=-1.25cm]frame.north east);
\node[anchor=west,text=white,font=\sffamily\bfseries\LARGE] at ([xshift=5mm,yshift=-0.62cm]frame.north west) {Question 74 C};
\node[anchor=east,text=white,font=\sffamily\itshape\large] at ([xshift=-5mm,yshift=-0.62cm]frame.north east) {Produit scalaire et géométrie repérée};}]
\vspace*{1.75cm}
\begin{minipage}[t][6.9cm][c]{\linewidth}\LARGE\centering
\begin{minipage}{0.95\linewidth}\raggedright\linespread{1.25}\selectfont Lequel de ces vecteurs est un vecteur normal à la droite d'équation $x-2y+7=0$~?\end{minipage}
\end{minipage}

\vfill
\begin{tcolorbox}[enhanced,nobeforeafter,width=7.1cm,height=2.1cm,valign=center,colback=cpsc!6,colframe=cpsc,arc=2.5mm,boxrule=0.9pt,left=3mm]{\Large\bfseries\color{cpsc}a)}\quad\begin{minipage}[c]{5.3cm}\raggedright\large $\begin{pmatrix}1\\-2\end{pmatrix}$\end{minipage}\end{tcolorbox}\hfill\begin{tcolorbox}[enhanced,nobeforeafter,width=7.1cm,height=2.1cm,valign=center,colback=cpsc!6,colframe=cpsc,arc=2.5mm,boxrule=0.9pt,left=3mm]{\Large\bfseries\color{cpsc}b)}\quad\begin{minipage}[c]{5.3cm}\raggedright\large $\begin{pmatrix}1\\7\end{pmatrix}$\end{minipage}\end{tcolorbox}\par\vspace{2.5mm}
\begin{tcolorbox}[enhanced,nobeforeafter,width=7.1cm,height=2.1cm,valign=center,colback=cpsc!6,colframe=cpsc,arc=2.5mm,boxrule=0.9pt,left=3mm]{\Large\bfseries\color{cpsc}c)}\quad\begin{minipage}[c]{5.3cm}\raggedright\large $\begin{pmatrix}-2\\1\end{pmatrix}$\end{minipage}\end{tcolorbox}\hfill\begin{tcolorbox}[enhanced,nobeforeafter,width=7.1cm,height=2.1cm,valign=center,colback=cpsc!6,colframe=cpsc,arc=2.5mm,boxrule=0.9pt,left=3mm]{\Large\bfseries\color{cpsc}d)}\quad\begin{minipage}[c]{5.3cm}\raggedright\large $\begin{pmatrix}2\\1\end{pmatrix}$\end{minipage}\end{tcolorbox}
\end{tcolorbox}
\newpage
\begin{tcolorbox}[enhanced,width=16cm,height=14cm,colback=white,colframe=cpsc,boxrule=2pt,arc=0mm,
left=6mm,right=6mm,top=0mm,bottom=4mm,sharp corners,
overlay={\fill[cpsc] (frame.north west) rectangle ([yshift=-1.25cm]frame.north east);
\node[anchor=west,text=white,font=\sffamily\bfseries\LARGE] at ([xshift=5mm,yshift=-0.62cm]frame.north west) {Question 75 C};
\node[anchor=east,text=white,font=\sffamily\itshape\large] at ([xshift=-5mm,yshift=-0.62cm]frame.north east) {Produit scalaire et géométrie repérée};}]
\vspace*{1.75cm}
\begin{minipage}[t][6.9cm][c]{\linewidth}\LARGE\centering
\begin{minipage}{0.95\linewidth}\raggedright\linespread{1.25}\selectfont Quelle est une équation cartésienne de la droite passant par $A(3\,;\,3)$ et de vecteur normal $\vec n\begin{pmatrix}4\\-3\end{pmatrix}$~?\end{minipage}
\end{minipage}

\vfill
\begin{tcolorbox}[enhanced,nobeforeafter,width=7.1cm,height=2.1cm,valign=center,colback=cpsc!6,colframe=cpsc,arc=2.5mm,boxrule=0.9pt,left=3mm]{\Large\bfseries\color{cpsc}a)}\quad\begin{minipage}[c]{5.3cm}\raggedright\Large $4x-3y-3=0$\end{minipage}\end{tcolorbox}\hfill\begin{tcolorbox}[enhanced,nobeforeafter,width=7.1cm,height=2.1cm,valign=center,colback=cpsc!6,colframe=cpsc,arc=2.5mm,boxrule=0.9pt,left=3mm]{\Large\bfseries\color{cpsc}b)}\quad\begin{minipage}[c]{5.3cm}\raggedright\Large $3x+3y-3=0$\end{minipage}\end{tcolorbox}\par\vspace{2.5mm}
\begin{tcolorbox}[enhanced,nobeforeafter,width=7.1cm,height=2.1cm,valign=center,colback=cpsc!6,colframe=cpsc,arc=2.5mm,boxrule=0.9pt,left=3mm]{\Large\bfseries\color{cpsc}c)}\quad\begin{minipage}[c]{5.3cm}\raggedright\Large $3x+4y-21=0$\end{minipage}\end{tcolorbox}\hfill\begin{tcolorbox}[enhanced,nobeforeafter,width=7.1cm,height=2.1cm,valign=center,colback=cpsc!6,colframe=cpsc,arc=2.5mm,boxrule=0.9pt,left=3mm]{\Large\bfseries\color{cpsc}d)}\quad\begin{minipage}[c]{5.3cm}\raggedright\Large $4x-3y+3=0$\end{minipage}\end{tcolorbox}
\end{tcolorbox}
\newpage
\begin{tcolorbox}[enhanced,width=16cm,height=14cm,colback=white,colframe=cpsc,boxrule=2pt,arc=0mm,
left=6mm,right=6mm,top=0mm,bottom=4mm,sharp corners,
overlay={\fill[cpsc] (frame.north west) rectangle ([yshift=-1.25cm]frame.north east);
\node[anchor=west,text=white,font=\sffamily\bfseries\LARGE] at ([xshift=5mm,yshift=-0.62cm]frame.north west) {Question 76 C};
\node[anchor=east,text=white,font=\sffamily\itshape\large] at ([xshift=-5mm,yshift=-0.62cm]frame.north east) {Produit scalaire et géométrie repérée};}]
\vspace*{1.75cm}
\begin{minipage}[t][6.9cm][c]{\linewidth}\LARGE\centering
\begin{minipage}{0.95\linewidth}\raggedright\linespread{1.25}\selectfont Quelle est une équation cartésienne de la droite passant par $A(-4\,;\,-1)$ et de vecteur normal $\vec n\begin{pmatrix}1\\-2\end{pmatrix}$~?\end{minipage}
\end{minipage}

\vfill
\begin{tcolorbox}[enhanced,nobeforeafter,width=7.1cm,height=2.1cm,valign=center,colback=cpsc!6,colframe=cpsc,arc=2.5mm,boxrule=0.9pt,left=3mm]{\Large\bfseries\color{cpsc}a)}\quad\begin{minipage}[c]{5.3cm}\raggedright\Large $x-2y-2=0$\end{minipage}\end{tcolorbox}\hfill\begin{tcolorbox}[enhanced,nobeforeafter,width=7.1cm,height=2.1cm,valign=center,colback=cpsc!6,colframe=cpsc,arc=2.5mm,boxrule=0.9pt,left=3mm]{\Large\bfseries\color{cpsc}b)}\quad\begin{minipage}[c]{5.3cm}\raggedright\Large $-4x-y+2=0$\end{minipage}\end{tcolorbox}\par\vspace{2.5mm}
\begin{tcolorbox}[enhanced,nobeforeafter,width=7.1cm,height=2.1cm,valign=center,colback=cpsc!6,colframe=cpsc,arc=2.5mm,boxrule=0.9pt,left=3mm]{\Large\bfseries\color{cpsc}c)}\quad\begin{minipage}[c]{5.3cm}\raggedright\Large $2x+y+9=0$\end{minipage}\end{tcolorbox}\hfill\begin{tcolorbox}[enhanced,nobeforeafter,width=7.1cm,height=2.1cm,valign=center,colback=cpsc!6,colframe=cpsc,arc=2.5mm,boxrule=0.9pt,left=3mm]{\Large\bfseries\color{cpsc}d)}\quad\begin{minipage}[c]{5.3cm}\raggedright\Large $x-2y+2=0$\end{minipage}\end{tcolorbox}
\end{tcolorbox}
\newpage
\begin{tcolorbox}[enhanced,width=16cm,height=14cm,colback=white,colframe=cpsc,boxrule=2pt,arc=0mm,
left=6mm,right=6mm,top=0mm,bottom=4mm,sharp corners,
overlay={\fill[cpsc] (frame.north west) rectangle ([yshift=-1.25cm]frame.north east);
\node[anchor=west,text=white,font=\sffamily\bfseries\LARGE] at ([xshift=5mm,yshift=-0.62cm]frame.north west) {Question 77 C};
\node[anchor=east,text=white,font=\sffamily\itshape\large] at ([xshift=-5mm,yshift=-0.62cm]frame.north east) {Produit scalaire et géométrie repérée};}]
\vspace*{1.75cm}
\begin{minipage}[t][6.9cm][c]{\linewidth}\LARGE\centering
\begin{minipage}{0.95\linewidth}\raggedright\linespread{1.25}\selectfont Quels sont le centre et le rayon du cercle d'équation $(x+2)^2+(y+6)^2=49$~?\end{minipage}
\end{minipage}

\vfill
\begin{tcolorbox}[enhanced,nobeforeafter,width=7.1cm,height=2.1cm,valign=center,colback=cpsc!6,colframe=cpsc,arc=2.5mm,boxrule=0.9pt,left=3mm]{\Large\bfseries\color{cpsc}a)}\quad\begin{minipage}[c]{5.3cm}\raggedright\large $\Omega(-2\,;\,-6)\text{, rayon }7$\end{minipage}\end{tcolorbox}\hfill\begin{tcolorbox}[enhanced,nobeforeafter,width=7.1cm,height=2.1cm,valign=center,colback=cpsc!6,colframe=cpsc,arc=2.5mm,boxrule=0.9pt,left=3mm]{\Large\bfseries\color{cpsc}b)}\quad\begin{minipage}[c]{5.3cm}\raggedright\large $\Omega(-6\,;\,-2)\text{, rayon }7$\end{minipage}\end{tcolorbox}\par\vspace{2.5mm}
\begin{tcolorbox}[enhanced,nobeforeafter,width=7.1cm,height=2.1cm,valign=center,colback=cpsc!6,colframe=cpsc,arc=2.5mm,boxrule=0.9pt,left=3mm]{\Large\bfseries\color{cpsc}c)}\quad\begin{minipage}[c]{5.3cm}\raggedright\large $\Omega(2\,;\,6)\text{, rayon }7$\end{minipage}\end{tcolorbox}\hfill\begin{tcolorbox}[enhanced,nobeforeafter,width=7.1cm,height=2.1cm,valign=center,colback=cpsc!6,colframe=cpsc,arc=2.5mm,boxrule=0.9pt,left=3mm]{\Large\bfseries\color{cpsc}d)}\quad\begin{minipage}[c]{5.3cm}\raggedright\large $\Omega(-2\,;\,-6)\text{, rayon }49$\end{minipage}\end{tcolorbox}
\end{tcolorbox}
\newpage
\begin{tcolorbox}[enhanced,width=16cm,height=14cm,colback=white,colframe=cpsc,boxrule=2pt,arc=0mm,
left=6mm,right=6mm,top=0mm,bottom=4mm,sharp corners,
overlay={\fill[cpsc] (frame.north west) rectangle ([yshift=-1.25cm]frame.north east);
\node[anchor=west,text=white,font=\sffamily\bfseries\LARGE] at ([xshift=5mm,yshift=-0.62cm]frame.north west) {Question 78 C};
\node[anchor=east,text=white,font=\sffamily\itshape\large] at ([xshift=-5mm,yshift=-0.62cm]frame.north east) {Produit scalaire et géométrie repérée};}]
\vspace*{1.75cm}
\begin{minipage}[t][6.9cm][c]{\linewidth}\LARGE\centering
\begin{minipage}{0.95\linewidth}\raggedright\linespread{1.25}\selectfont Quels sont le centre et le rayon du cercle d'équation $(x-4)^2+y^2=12$~?\end{minipage}
\end{minipage}

\vfill
\begin{tcolorbox}[enhanced,nobeforeafter,width=7.1cm,height=2.1cm,valign=center,colback=cpsc!6,colframe=cpsc,arc=2.5mm,boxrule=0.9pt,left=3mm]{\Large\bfseries\color{cpsc}a)}\quad\begin{minipage}[c]{5.3cm}\raggedright\large $\Omega(-4\,;\,0)\text{, rayon }2\sqrt{3}$\end{minipage}\end{tcolorbox}\hfill\begin{tcolorbox}[enhanced,nobeforeafter,width=7.1cm,height=2.1cm,valign=center,colback=cpsc!6,colframe=cpsc,arc=2.5mm,boxrule=0.9pt,left=3mm]{\Large\bfseries\color{cpsc}b)}\quad\begin{minipage}[c]{5.3cm}\raggedright\large $\Omega(4\,;\,0)\text{, rayon }12$\end{minipage}\end{tcolorbox}\par\vspace{2.5mm}
\begin{tcolorbox}[enhanced,nobeforeafter,width=7.1cm,height=2.1cm,valign=center,colback=cpsc!6,colframe=cpsc,arc=2.5mm,boxrule=0.9pt,left=3mm]{\Large\bfseries\color{cpsc}c)}\quad\begin{minipage}[c]{5.3cm}\raggedright\large $\Omega(0\,;\,4)\text{, rayon }2\sqrt{3}$\end{minipage}\end{tcolorbox}\hfill\begin{tcolorbox}[enhanced,nobeforeafter,width=7.1cm,height=2.1cm,valign=center,colback=cpsc!6,colframe=cpsc,arc=2.5mm,boxrule=0.9pt,left=3mm]{\Large\bfseries\color{cpsc}d)}\quad\begin{minipage}[c]{5.3cm}\raggedright\large $\Omega(4\,;\,0)\text{, rayon }2\sqrt{3}$\end{minipage}\end{tcolorbox}
\end{tcolorbox}
\newpage
\begin{tcolorbox}[enhanced,width=16cm,height=14cm,colback=white,colframe=cpsc,boxrule=2pt,arc=0mm,
left=6mm,right=6mm,top=0mm,bottom=4mm,sharp corners,
overlay={\fill[cpsc] (frame.north west) rectangle ([yshift=-1.25cm]frame.north east);
\node[anchor=west,text=white,font=\sffamily\bfseries\LARGE] at ([xshift=5mm,yshift=-0.62cm]frame.north west) {Question 79 C};
\node[anchor=east,text=white,font=\sffamily\itshape\large] at ([xshift=-5mm,yshift=-0.62cm]frame.north east) {Produit scalaire et géométrie repérée};}]
\vspace*{1.75cm}
\begin{minipage}[t][6.9cm][c]{\linewidth}\LARGE\centering
\begin{minipage}{0.95\linewidth}\raggedright\linespread{1.25}\selectfont On a $\|\vec u\|=4$, $\|\vec v\|=3$ et $\vec u\cdot\vec v=6$. Que vaut $\|\vec u+\vec v\|^2$~?\end{minipage}
\end{minipage}

\vfill
\begin{tcolorbox}[enhanced,nobeforeafter,width=7.1cm,height=2.1cm,valign=center,colback=cpsc!6,colframe=cpsc,arc=2.5mm,boxrule=0.9pt,left=3mm]{\Large\bfseries\color{cpsc}a)}\quad\begin{minipage}[c]{5.3cm}\raggedright\Large $31$\end{minipage}\end{tcolorbox}\hfill\begin{tcolorbox}[enhanced,nobeforeafter,width=7.1cm,height=2.1cm,valign=center,colback=cpsc!6,colframe=cpsc,arc=2.5mm,boxrule=0.9pt,left=3mm]{\Large\bfseries\color{cpsc}b)}\quad\begin{minipage}[c]{5.3cm}\raggedright\Large $49$\end{minipage}\end{tcolorbox}\par\vspace{2.5mm}
\begin{tcolorbox}[enhanced,nobeforeafter,width=7.1cm,height=2.1cm,valign=center,colback=cpsc!6,colframe=cpsc,arc=2.5mm,boxrule=0.9pt,left=3mm]{\Large\bfseries\color{cpsc}c)}\quad\begin{minipage}[c]{5.3cm}\raggedright\Large $37$\end{minipage}\end{tcolorbox}\hfill\begin{tcolorbox}[enhanced,nobeforeafter,width=7.1cm,height=2.1cm,valign=center,colback=cpsc!6,colframe=cpsc,arc=2.5mm,boxrule=0.9pt,left=3mm]{\Large\bfseries\color{cpsc}d)}\quad\begin{minipage}[c]{5.3cm}\raggedright\Large $25$\end{minipage}\end{tcolorbox}
\end{tcolorbox}
\newpage
\begin{tcolorbox}[enhanced,width=16cm,height=14cm,colback=white,colframe=cpsc,boxrule=2pt,arc=0mm,
left=6mm,right=6mm,top=0mm,bottom=4mm,sharp corners,
overlay={\fill[cpsc] (frame.north west) rectangle ([yshift=-1.25cm]frame.north east);
\node[anchor=west,text=white,font=\sffamily\bfseries\LARGE] at ([xshift=5mm,yshift=-0.62cm]frame.north west) {Question 80 C};
\node[anchor=east,text=white,font=\sffamily\itshape\large] at ([xshift=-5mm,yshift=-0.62cm]frame.north east) {Produit scalaire et géométrie repérée};}]
\vspace*{1.75cm}
\begin{minipage}[t][6.9cm][c]{\linewidth}\LARGE\centering
\begin{minipage}{0.95\linewidth}\raggedright\linespread{1.25}\selectfont On a $\|\vec u\|=5$, $\|\vec v\|=5$ et $\vec u\cdot\vec v=-8$. Que vaut $\|\vec u+\vec v\|^2$~?\end{minipage}
\end{minipage}

\vfill
\begin{tcolorbox}[enhanced,nobeforeafter,width=7.1cm,height=2.1cm,valign=center,colback=cpsc!6,colframe=cpsc,arc=2.5mm,boxrule=0.9pt,left=3mm]{\Large\bfseries\color{cpsc}a)}\quad\begin{minipage}[c]{5.3cm}\raggedright\Large $100$\end{minipage}\end{tcolorbox}\hfill\begin{tcolorbox}[enhanced,nobeforeafter,width=7.1cm,height=2.1cm,valign=center,colback=cpsc!6,colframe=cpsc,arc=2.5mm,boxrule=0.9pt,left=3mm]{\Large\bfseries\color{cpsc}b)}\quad\begin{minipage}[c]{5.3cm}\raggedright\Large $42$\end{minipage}\end{tcolorbox}\par\vspace{2.5mm}
\begin{tcolorbox}[enhanced,nobeforeafter,width=7.1cm,height=2.1cm,valign=center,colback=cpsc!6,colframe=cpsc,arc=2.5mm,boxrule=0.9pt,left=3mm]{\Large\bfseries\color{cpsc}c)}\quad\begin{minipage}[c]{5.3cm}\raggedright\Large $50$\end{minipage}\end{tcolorbox}\hfill\begin{tcolorbox}[enhanced,nobeforeafter,width=7.1cm,height=2.1cm,valign=center,colback=cpsc!6,colframe=cpsc,arc=2.5mm,boxrule=0.9pt,left=3mm]{\Large\bfseries\color{cpsc}d)}\quad\begin{minipage}[c]{5.3cm}\raggedright\Large $34$\end{minipage}\end{tcolorbox}
\end{tcolorbox}
\newpage
\begin{tcolorbox}[enhanced,width=16cm,height=14cm,colback=white,colframe=ccvp,boxrule=2pt,arc=0mm,
left=6mm,right=6mm,top=0mm,bottom=4mm,sharp corners,
overlay={\fill[ccvp] (frame.north west) rectangle ([yshift=-1.25cm]frame.north east);
\node[anchor=west,text=white,font=\sffamily\bfseries\LARGE] at ([xshift=5mm,yshift=-0.62cm]frame.north west) {Question 81 C};
\node[anchor=east,text=white,font=\sffamily\itshape\large] at ([xshift=-5mm,yshift=-0.62cm]frame.north east) {Probabilités conditionnelles et variables aléatoires};}]
\vspace*{1.75cm}
\begin{minipage}[t][6.9cm][c]{\linewidth}\LARGE\centering
\begin{minipage}{0.95\linewidth}\raggedright\linespread{1.25}\selectfont On a $P(A)=0{,}6$ et $P(A\cap B)=0{,}3$. Que vaut $P_A(B)$~?\end{minipage}
\end{minipage}

\vfill
\begin{tcolorbox}[enhanced,nobeforeafter,width=7.1cm,height=2.1cm,valign=center,colback=ccvp!6,colframe=ccvp,arc=2.5mm,boxrule=0.9pt,left=3mm]{\Large\bfseries\color{ccvp}a)}\quad\begin{minipage}[c]{5.3cm}\raggedright\Large $0{,}18$\end{minipage}\end{tcolorbox}\hfill\begin{tcolorbox}[enhanced,nobeforeafter,width=7.1cm,height=2.1cm,valign=center,colback=ccvp!6,colframe=ccvp,arc=2.5mm,boxrule=0.9pt,left=3mm]{\Large\bfseries\color{ccvp}b)}\quad\begin{minipage}[c]{5.3cm}\raggedright\Large $0{,}3$\end{minipage}\end{tcolorbox}\par\vspace{2.5mm}
\begin{tcolorbox}[enhanced,nobeforeafter,width=7.1cm,height=2.1cm,valign=center,colback=ccvp!6,colframe=ccvp,arc=2.5mm,boxrule=0.9pt,left=3mm]{\Large\bfseries\color{ccvp}c)}\quad\begin{minipage}[c]{5.3cm}\raggedright\Large $0{,}9$\end{minipage}\end{tcolorbox}\hfill\begin{tcolorbox}[enhanced,nobeforeafter,width=7.1cm,height=2.1cm,valign=center,colback=ccvp!6,colframe=ccvp,arc=2.5mm,boxrule=0.9pt,left=3mm]{\Large\bfseries\color{ccvp}d)}\quad\begin{minipage}[c]{5.3cm}\raggedright\Large $0{,}5$\end{minipage}\end{tcolorbox}
\end{tcolorbox}
\newpage
\begin{tcolorbox}[enhanced,width=16cm,height=14cm,colback=white,colframe=ccvp,boxrule=2pt,arc=0mm,
left=6mm,right=6mm,top=0mm,bottom=4mm,sharp corners,
overlay={\fill[ccvp] (frame.north west) rectangle ([yshift=-1.25cm]frame.north east);
\node[anchor=west,text=white,font=\sffamily\bfseries\LARGE] at ([xshift=5mm,yshift=-0.62cm]frame.north west) {Question 82 C};
\node[anchor=east,text=white,font=\sffamily\itshape\large] at ([xshift=-5mm,yshift=-0.62cm]frame.north east) {Probabilités conditionnelles et variables aléatoires};}]
\vspace*{1.75cm}
\begin{minipage}[t][6.9cm][c]{\linewidth}\LARGE\centering
\begin{minipage}{0.95\linewidth}\raggedright\linespread{1.25}\selectfont On a $P(A)=0{,}5$ et $P(A\cap B)=0{,}35$. Que vaut $P_A(B)$~?\end{minipage}
\end{minipage}

\vfill
\begin{tcolorbox}[enhanced,nobeforeafter,width=7.1cm,height=2.1cm,valign=center,colback=ccvp!6,colframe=ccvp,arc=2.5mm,boxrule=0.9pt,left=3mm]{\Large\bfseries\color{ccvp}a)}\quad\begin{minipage}[c]{5.3cm}\raggedright\Large $0{,}175$\end{minipage}\end{tcolorbox}\hfill\begin{tcolorbox}[enhanced,nobeforeafter,width=7.1cm,height=2.1cm,valign=center,colback=ccvp!6,colframe=ccvp,arc=2.5mm,boxrule=0.9pt,left=3mm]{\Large\bfseries\color{ccvp}b)}\quad\begin{minipage}[c]{5.3cm}\raggedright\Large $0{,}15$\end{minipage}\end{tcolorbox}\par\vspace{2.5mm}
\begin{tcolorbox}[enhanced,nobeforeafter,width=7.1cm,height=2.1cm,valign=center,colback=ccvp!6,colframe=ccvp,arc=2.5mm,boxrule=0.9pt,left=3mm]{\Large\bfseries\color{ccvp}c)}\quad\begin{minipage}[c]{5.3cm}\raggedright\Large $0{,}85$\end{minipage}\end{tcolorbox}\hfill\begin{tcolorbox}[enhanced,nobeforeafter,width=7.1cm,height=2.1cm,valign=center,colback=ccvp!6,colframe=ccvp,arc=2.5mm,boxrule=0.9pt,left=3mm]{\Large\bfseries\color{ccvp}d)}\quad\begin{minipage}[c]{5.3cm}\raggedright\Large $0{,}7$\end{minipage}\end{tcolorbox}
\end{tcolorbox}
\newpage
\begin{tcolorbox}[enhanced,width=16cm,height=14cm,colback=white,colframe=ccvp,boxrule=2pt,arc=0mm,
left=6mm,right=6mm,top=0mm,bottom=4mm,sharp corners,
overlay={\fill[ccvp] (frame.north west) rectangle ([yshift=-1.25cm]frame.north east);
\node[anchor=west,text=white,font=\sffamily\bfseries\LARGE] at ([xshift=5mm,yshift=-0.62cm]frame.north west) {Question 83 C};
\node[anchor=east,text=white,font=\sffamily\itshape\large] at ([xshift=-5mm,yshift=-0.62cm]frame.north east) {Probabilités conditionnelles et variables aléatoires};}]
\vspace*{1.75cm}
\begin{minipage}[t][6.9cm][c]{\linewidth}\LARGE\centering
\begin{minipage}{0.95\linewidth}\raggedright\linespread{1.25}\selectfont On a $P(A)=0{,}7$ et $P_A(B)=0{,}9$. Que vaut $P(A\cap B)$~?\end{minipage}
\end{minipage}

\vfill
\begin{tcolorbox}[enhanced,nobeforeafter,width=7.1cm,height=2.1cm,valign=center,colback=ccvp!6,colframe=ccvp,arc=2.5mm,boxrule=0.9pt,left=3mm]{\Large\bfseries\color{ccvp}a)}\quad\begin{minipage}[c]{5.3cm}\raggedright\Large $0{,}37$\end{minipage}\end{tcolorbox}\hfill\begin{tcolorbox}[enhanced,nobeforeafter,width=7.1cm,height=2.1cm,valign=center,colback=ccvp!6,colframe=ccvp,arc=2.5mm,boxrule=0.9pt,left=3mm]{\Large\bfseries\color{ccvp}b)}\quad\begin{minipage}[c]{5.3cm}\raggedright\Large $0{,}6$\end{minipage}\end{tcolorbox}\par\vspace{2.5mm}
\begin{tcolorbox}[enhanced,nobeforeafter,width=7.1cm,height=2.1cm,valign=center,colback=ccvp!6,colframe=ccvp,arc=2.5mm,boxrule=0.9pt,left=3mm]{\Large\bfseries\color{ccvp}c)}\quad\begin{minipage}[c]{5.3cm}\raggedright\Large $0{,}63$\end{minipage}\end{tcolorbox}\hfill\begin{tcolorbox}[enhanced,nobeforeafter,width=7.1cm,height=2.1cm,valign=center,colback=ccvp!6,colframe=ccvp,arc=2.5mm,boxrule=0.9pt,left=3mm]{\Large\bfseries\color{ccvp}d)}\quad\begin{minipage}[c]{5.3cm}\raggedright\Large $0{,}07$\end{minipage}\end{tcolorbox}
\end{tcolorbox}
\newpage
\begin{tcolorbox}[enhanced,width=16cm,height=14cm,colback=white,colframe=ccvp,boxrule=2pt,arc=0mm,
left=6mm,right=6mm,top=0mm,bottom=4mm,sharp corners,
overlay={\fill[ccvp] (frame.north west) rectangle ([yshift=-1.25cm]frame.north east);
\node[anchor=west,text=white,font=\sffamily\bfseries\LARGE] at ([xshift=5mm,yshift=-0.62cm]frame.north west) {Question 84 C};
\node[anchor=east,text=white,font=\sffamily\itshape\large] at ([xshift=-5mm,yshift=-0.62cm]frame.north east) {Probabilités conditionnelles et variables aléatoires};}]
\vspace*{1.75cm}
\begin{minipage}[t][6.9cm][c]{\linewidth}\LARGE\centering
\begin{minipage}{0.95\linewidth}\raggedright\linespread{1.25}\selectfont On a $P(A)=0{,}25$ et $P_A(B)=0{,}6$. Que vaut $P(A\cap B)$~?\end{minipage}
\end{minipage}

\vfill
\begin{tcolorbox}[enhanced,nobeforeafter,width=7.1cm,height=2.1cm,valign=center,colback=ccvp!6,colframe=ccvp,arc=2.5mm,boxrule=0.9pt,left=3mm]{\Large\bfseries\color{ccvp}a)}\quad\begin{minipage}[c]{5.3cm}\raggedright\Large $0{,}85$\end{minipage}\end{tcolorbox}\hfill\begin{tcolorbox}[enhanced,nobeforeafter,width=7.1cm,height=2.1cm,valign=center,colback=ccvp!6,colframe=ccvp,arc=2.5mm,boxrule=0.9pt,left=3mm]{\Large\bfseries\color{ccvp}b)}\quad\begin{minipage}[c]{5.3cm}\raggedright\Large $0{,}15$\end{minipage}\end{tcolorbox}\par\vspace{2.5mm}
\begin{tcolorbox}[enhanced,nobeforeafter,width=7.1cm,height=2.1cm,valign=center,colback=ccvp!6,colframe=ccvp,arc=2.5mm,boxrule=0.9pt,left=3mm]{\Large\bfseries\color{ccvp}c)}\quad\begin{minipage}[c]{5.3cm}\raggedright\Large $0{,}6$\end{minipage}\end{tcolorbox}\hfill\begin{tcolorbox}[enhanced,nobeforeafter,width=7.1cm,height=2.1cm,valign=center,colback=ccvp!6,colframe=ccvp,arc=2.5mm,boxrule=0.9pt,left=3mm]{\Large\bfseries\color{ccvp}d)}\quad\begin{minipage}[c]{5.3cm}\raggedright\Large $0{,}1$\end{minipage}\end{tcolorbox}
\end{tcolorbox}
\newpage
\begin{tcolorbox}[enhanced,width=16cm,height=14cm,colback=white,colframe=ccvp,boxrule=2pt,arc=0mm,
left=6mm,right=6mm,top=0mm,bottom=4mm,sharp corners,
overlay={\fill[ccvp] (frame.north west) rectangle ([yshift=-1.25cm]frame.north east);
\node[anchor=west,text=white,font=\sffamily\bfseries\LARGE] at ([xshift=5mm,yshift=-0.62cm]frame.north west) {Question 85 C};
\node[anchor=east,text=white,font=\sffamily\itshape\large] at ([xshift=-5mm,yshift=-0.62cm]frame.north east) {Probabilités conditionnelles et variables aléatoires};}]
\vspace*{1.75cm}
\begin{minipage}[t][6.9cm][c]{\linewidth}\LARGE\centering
\begin{minipage}{0.95\linewidth}\raggedright\linespread{1.25}\selectfont On a $P(A)=0{,}4$, $P_A(B)=0{,}25$ et $P_{\overline{A}}(B)=0{,}5$. Que vaut $P(B)$~?\end{minipage}
\end{minipage}

\vfill
\begin{tcolorbox}[enhanced,nobeforeafter,width=7.1cm,height=2.1cm,valign=center,colback=ccvp!6,colframe=ccvp,arc=2.5mm,boxrule=0.9pt,left=3mm]{\Large\bfseries\color{ccvp}a)}\quad\begin{minipage}[c]{5.3cm}\raggedright\Large $0{,}3$\end{minipage}\end{tcolorbox}\hfill\begin{tcolorbox}[enhanced,nobeforeafter,width=7.1cm,height=2.1cm,valign=center,colback=ccvp!6,colframe=ccvp,arc=2.5mm,boxrule=0.9pt,left=3mm]{\Large\bfseries\color{ccvp}b)}\quad\begin{minipage}[c]{5.3cm}\raggedright\Large $0{,}4$\end{minipage}\end{tcolorbox}\par\vspace{2.5mm}
\begin{tcolorbox}[enhanced,nobeforeafter,width=7.1cm,height=2.1cm,valign=center,colback=ccvp!6,colframe=ccvp,arc=2.5mm,boxrule=0.9pt,left=3mm]{\Large\bfseries\color{ccvp}c)}\quad\begin{minipage}[c]{5.3cm}\raggedright\Large $0{,}75$\end{minipage}\end{tcolorbox}\hfill\begin{tcolorbox}[enhanced,nobeforeafter,width=7.1cm,height=2.1cm,valign=center,colback=ccvp!6,colframe=ccvp,arc=2.5mm,boxrule=0.9pt,left=3mm]{\Large\bfseries\color{ccvp}d)}\quad\begin{minipage}[c]{5.3cm}\raggedright\Large $0{,}1$\end{minipage}\end{tcolorbox}
\end{tcolorbox}
\newpage
\begin{tcolorbox}[enhanced,width=16cm,height=14cm,colback=white,colframe=ccvp,boxrule=2pt,arc=0mm,
left=6mm,right=6mm,top=0mm,bottom=4mm,sharp corners,
overlay={\fill[ccvp] (frame.north west) rectangle ([yshift=-1.25cm]frame.north east);
\node[anchor=west,text=white,font=\sffamily\bfseries\LARGE] at ([xshift=5mm,yshift=-0.62cm]frame.north west) {Question 86 C};
\node[anchor=east,text=white,font=\sffamily\itshape\large] at ([xshift=-5mm,yshift=-0.62cm]frame.north east) {Probabilités conditionnelles et variables aléatoires};}]
\vspace*{1.75cm}
\begin{minipage}[t][6.9cm][c]{\linewidth}\LARGE\centering
\begin{minipage}{0.95\linewidth}\raggedright\linespread{1.25}\selectfont On a $P(A)=0{,}7$, $P_A(B)=0{,}6$ et $P_{\overline{A}}(B)=0{,}2$. Que vaut $P(B)$~?\end{minipage}
\end{minipage}

\vfill
\begin{tcolorbox}[enhanced,nobeforeafter,width=7.1cm,height=2.1cm,valign=center,colback=ccvp!6,colframe=ccvp,arc=2.5mm,boxrule=0.9pt,left=3mm]{\Large\bfseries\color{ccvp}a)}\quad\begin{minipage}[c]{5.3cm}\raggedright\Large $0{,}56$\end{minipage}\end{tcolorbox}\hfill\begin{tcolorbox}[enhanced,nobeforeafter,width=7.1cm,height=2.1cm,valign=center,colback=ccvp!6,colframe=ccvp,arc=2.5mm,boxrule=0.9pt,left=3mm]{\Large\bfseries\color{ccvp}b)}\quad\begin{minipage}[c]{5.3cm}\raggedright\Large $0{,}48$\end{minipage}\end{tcolorbox}\par\vspace{2.5mm}
\begin{tcolorbox}[enhanced,nobeforeafter,width=7.1cm,height=2.1cm,valign=center,colback=ccvp!6,colframe=ccvp,arc=2.5mm,boxrule=0.9pt,left=3mm]{\Large\bfseries\color{ccvp}c)}\quad\begin{minipage}[c]{5.3cm}\raggedright\Large $0{,}42$\end{minipage}\end{tcolorbox}\hfill\begin{tcolorbox}[enhanced,nobeforeafter,width=7.1cm,height=2.1cm,valign=center,colback=ccvp!6,colframe=ccvp,arc=2.5mm,boxrule=0.9pt,left=3mm]{\Large\bfseries\color{ccvp}d)}\quad\begin{minipage}[c]{5.3cm}\raggedright\Large $0{,}8$\end{minipage}\end{tcolorbox}
\end{tcolorbox}
\newpage
\begin{tcolorbox}[enhanced,width=16cm,height=14cm,colback=white,colframe=ccvp,boxrule=2pt,arc=0mm,
left=6mm,right=6mm,top=0mm,bottom=4mm,sharp corners,
overlay={\fill[ccvp] (frame.north west) rectangle ([yshift=-1.25cm]frame.north east);
\node[anchor=west,text=white,font=\sffamily\bfseries\LARGE] at ([xshift=5mm,yshift=-0.62cm]frame.north west) {Question 87 C};
\node[anchor=east,text=white,font=\sffamily\itshape\large] at ([xshift=-5mm,yshift=-0.62cm]frame.north east) {Probabilités conditionnelles et variables aléatoires};}]
\vspace*{1.75cm}
\begin{minipage}[t][6.9cm][c]{\linewidth}\LARGE\centering
\begin{minipage}{0.95\linewidth}\raggedright\linespread{1.25}\selectfont On a $P(A)=0{,}6$, $P(B)=0{,}5$ et $P(A\cap B)=0{,}3$. Les événements $A$ et $B$ sont-ils indépendants~?\end{minipage}
\end{minipage}

\vfill
\begin{tcolorbox}[enhanced,nobeforeafter,width=7.1cm,height=2.1cm,valign=center,colback=ccvp!6,colframe=ccvp,arc=2.5mm,boxrule=0.9pt,left=3mm]{\Large\bfseries\color{ccvp}a)}\quad\begin{minipage}[c]{5.3cm}\raggedright\Large oui\end{minipage}\end{tcolorbox}\hfill\begin{tcolorbox}[enhanced,nobeforeafter,width=7.1cm,height=2.1cm,valign=center,colback=ccvp!6,colframe=ccvp,arc=2.5mm,boxrule=0.9pt,left=3mm]{\Large\bfseries\color{ccvp}b)}\quad\begin{minipage}[c]{5.3cm}\raggedright\Large on ne peut pas savoir\end{minipage}\end{tcolorbox}\par\vspace{2.5mm}
\begin{tcolorbox}[enhanced,nobeforeafter,width=7.1cm,height=2.1cm,valign=center,colback=ccvp!6,colframe=ccvp,arc=2.5mm,boxrule=0.9pt,left=3mm]{\Large\bfseries\color{ccvp}c)}\quad\begin{minipage}[c]{5.3cm}\raggedright\large seulement s'ils sont incompatibles\end{minipage}\end{tcolorbox}\hfill\begin{tcolorbox}[enhanced,nobeforeafter,width=7.1cm,height=2.1cm,valign=center,colback=ccvp!6,colframe=ccvp,arc=2.5mm,boxrule=0.9pt,left=3mm]{\Large\bfseries\color{ccvp}d)}\quad\begin{minipage}[c]{5.3cm}\raggedright\Large non\end{minipage}\end{tcolorbox}
\end{tcolorbox}
\newpage
\begin{tcolorbox}[enhanced,width=16cm,height=14cm,colback=white,colframe=ccvp,boxrule=2pt,arc=0mm,
left=6mm,right=6mm,top=0mm,bottom=4mm,sharp corners,
overlay={\fill[ccvp] (frame.north west) rectangle ([yshift=-1.25cm]frame.north east);
\node[anchor=west,text=white,font=\sffamily\bfseries\LARGE] at ([xshift=5mm,yshift=-0.62cm]frame.north west) {Question 88 C};
\node[anchor=east,text=white,font=\sffamily\itshape\large] at ([xshift=-5mm,yshift=-0.62cm]frame.north east) {Probabilités conditionnelles et variables aléatoires};}]
\vspace*{1.75cm}
\begin{minipage}[t][6.9cm][c]{\linewidth}\LARGE\centering
\begin{minipage}{0.95\linewidth}\raggedright\linespread{1.25}\selectfont On a $P(A)=0{,}2$, $P(B)=0{,}4$ et $P(A\cap B)=0{,}1$. Les événements $A$ et $B$ sont-ils indépendants~?\end{minipage}
\end{minipage}

\vfill
\begin{tcolorbox}[enhanced,nobeforeafter,width=7.1cm,height=2.1cm,valign=center,colback=ccvp!6,colframe=ccvp,arc=2.5mm,boxrule=0.9pt,left=3mm]{\Large\bfseries\color{ccvp}a)}\quad\begin{minipage}[c]{5.3cm}\raggedright\large seulement s'ils sont incompatibles\end{minipage}\end{tcolorbox}\hfill\begin{tcolorbox}[enhanced,nobeforeafter,width=7.1cm,height=2.1cm,valign=center,colback=ccvp!6,colframe=ccvp,arc=2.5mm,boxrule=0.9pt,left=3mm]{\Large\bfseries\color{ccvp}b)}\quad\begin{minipage}[c]{5.3cm}\raggedright\Large on ne peut pas savoir\end{minipage}\end{tcolorbox}\par\vspace{2.5mm}
\begin{tcolorbox}[enhanced,nobeforeafter,width=7.1cm,height=2.1cm,valign=center,colback=ccvp!6,colframe=ccvp,arc=2.5mm,boxrule=0.9pt,left=3mm]{\Large\bfseries\color{ccvp}c)}\quad\begin{minipage}[c]{5.3cm}\raggedright\Large oui\end{minipage}\end{tcolorbox}\hfill\begin{tcolorbox}[enhanced,nobeforeafter,width=7.1cm,height=2.1cm,valign=center,colback=ccvp!6,colframe=ccvp,arc=2.5mm,boxrule=0.9pt,left=3mm]{\Large\bfseries\color{ccvp}d)}\quad\begin{minipage}[c]{5.3cm}\raggedright\Large non\end{minipage}\end{tcolorbox}
\end{tcolorbox}
\newpage
\begin{tcolorbox}[enhanced,width=16cm,height=14cm,colback=white,colframe=ccvp,boxrule=2pt,arc=0mm,
left=6mm,right=6mm,top=0mm,bottom=4mm,sharp corners,
overlay={\fill[ccvp] (frame.north west) rectangle ([yshift=-1.25cm]frame.north east);
\node[anchor=west,text=white,font=\sffamily\bfseries\LARGE] at ([xshift=5mm,yshift=-0.62cm]frame.north west) {Question 89 C};
\node[anchor=east,text=white,font=\sffamily\itshape\large] at ([xshift=-5mm,yshift=-0.62cm]frame.north east) {Probabilités conditionnelles et variables aléatoires};}]
\vspace*{1.75cm}
\begin{minipage}[t][6.9cm][c]{\linewidth}\Large\centering
\begin{minipage}{0.95\linewidth}\raggedright\linespread{1.25}\selectfont Voici la loi de probabilité d'une variable aléatoire $X$. Que vaut $E(X)$~? $\begin{array}{|l|c|c|c|}\hline x_i & -5 & 0 & 10\\\hline P(X=x_i) & 0{,}1 & 0{,}6 & 0{,}3\\\hline\end{array}$\end{minipage}
\end{minipage}

\vfill
\begin{tcolorbox}[enhanced,nobeforeafter,width=7.1cm,height=2.1cm,valign=center,colback=ccvp!6,colframe=ccvp,arc=2.5mm,boxrule=0.9pt,left=3mm]{\Large\bfseries\color{ccvp}a)}\quad\begin{minipage}[c]{5.3cm}\raggedright\Large $2{,}5$\end{minipage}\end{tcolorbox}\hfill\begin{tcolorbox}[enhanced,nobeforeafter,width=7.1cm,height=2.1cm,valign=center,colback=ccvp!6,colframe=ccvp,arc=2.5mm,boxrule=0.9pt,left=3mm]{\Large\bfseries\color{ccvp}b)}\quad\begin{minipage}[c]{5.3cm}\raggedright\Large $3{,}5$\end{minipage}\end{tcolorbox}\par\vspace{2.5mm}
\begin{tcolorbox}[enhanced,nobeforeafter,width=7.1cm,height=2.1cm,valign=center,colback=ccvp!6,colframe=ccvp,arc=2.5mm,boxrule=0.9pt,left=3mm]{\Large\bfseries\color{ccvp}c)}\quad\begin{minipage}[c]{5.3cm}\raggedright\Large $-2{,}5$\end{minipage}\end{tcolorbox}\hfill\begin{tcolorbox}[enhanced,nobeforeafter,width=7.1cm,height=2.1cm,valign=center,colback=ccvp!6,colframe=ccvp,arc=2.5mm,boxrule=0.9pt,left=3mm]{\Large\bfseries\color{ccvp}d)}\quad\begin{minipage}[c]{5.3cm}\raggedright\Large $\dfrac{5}{3}$\end{minipage}\end{tcolorbox}
\end{tcolorbox}
\newpage
\begin{tcolorbox}[enhanced,width=16cm,height=14cm,colback=white,colframe=ccvp,boxrule=2pt,arc=0mm,
left=6mm,right=6mm,top=0mm,bottom=4mm,sharp corners,
overlay={\fill[ccvp] (frame.north west) rectangle ([yshift=-1.25cm]frame.north east);
\node[anchor=west,text=white,font=\sffamily\bfseries\LARGE] at ([xshift=5mm,yshift=-0.62cm]frame.north west) {Question 90 C};
\node[anchor=east,text=white,font=\sffamily\itshape\large] at ([xshift=-5mm,yshift=-0.62cm]frame.north east) {Probabilités conditionnelles et variables aléatoires};}]
\vspace*{1.75cm}
\begin{minipage}[t][6.9cm][c]{\linewidth}\Large\centering
\begin{minipage}{0.95\linewidth}\raggedright\linespread{1.25}\selectfont Voici la loi de probabilité d'une variable aléatoire $X$. Que vaut $E(X)$~? $\begin{array}{|l|c|c|}\hline x_i & 2 & 6\\\hline P(X=x_i) & 0{,}5 & 0{,}5\\\hline\end{array}$\end{minipage}
\end{minipage}

\vfill
\begin{tcolorbox}[enhanced,nobeforeafter,width=7.1cm,height=2.1cm,valign=center,colback=ccvp!6,colframe=ccvp,arc=2.5mm,boxrule=0.9pt,left=3mm]{\Large\bfseries\color{ccvp}a)}\quad\begin{minipage}[c]{5.3cm}\raggedright\Large $-4$\end{minipage}\end{tcolorbox}\hfill\begin{tcolorbox}[enhanced,nobeforeafter,width=7.1cm,height=2.1cm,valign=center,colback=ccvp!6,colframe=ccvp,arc=2.5mm,boxrule=0.9pt,left=3mm]{\Large\bfseries\color{ccvp}b)}\quad\begin{minipage}[c]{5.3cm}\raggedright\Large $6$\end{minipage}\end{tcolorbox}\par\vspace{2.5mm}
\begin{tcolorbox}[enhanced,nobeforeafter,width=7.1cm,height=2.1cm,valign=center,colback=ccvp!6,colframe=ccvp,arc=2.5mm,boxrule=0.9pt,left=3mm]{\Large\bfseries\color{ccvp}c)}\quad\begin{minipage}[c]{5.3cm}\raggedright\Large $4$\end{minipage}\end{tcolorbox}\hfill\begin{tcolorbox}[enhanced,nobeforeafter,width=7.1cm,height=2.1cm,valign=center,colback=ccvp!6,colframe=ccvp,arc=2.5mm,boxrule=0.9pt,left=3mm]{\Large\bfseries\color{ccvp}d)}\quad\begin{minipage}[c]{5.3cm}\raggedright\Large $5$\end{minipage}\end{tcolorbox}
\end{tcolorbox}
\newpage
\begin{tcolorbox}[enhanced,width=16cm,height=14cm,colback=white,colframe=ccvp,boxrule=2pt,arc=0mm,
left=6mm,right=6mm,top=0mm,bottom=4mm,sharp corners,
overlay={\fill[ccvp] (frame.north west) rectangle ([yshift=-1.25cm]frame.north east);
\node[anchor=west,text=white,font=\sffamily\bfseries\LARGE] at ([xshift=5mm,yshift=-0.62cm]frame.north west) {Question 91 C};
\node[anchor=east,text=white,font=\sffamily\itshape\large] at ([xshift=-5mm,yshift=-0.62cm]frame.north east) {Probabilités conditionnelles et variables aléatoires};}]
\vspace*{1.75cm}
\begin{minipage}[t][6.9cm][c]{\linewidth}\Large\centering
\begin{minipage}{0.95\linewidth}\raggedright\linespread{1.25}\selectfont Voici la loi de probabilité d'une variable aléatoire $X$. Que vaut sa variance $V(X)$~? $\begin{array}{|l|c|c|c|}\hline x_i & 1 & 2 & 3\\\hline P(X=x_i) & 0{,}25 & 0{,}5 & 0{,}25\\\hline\end{array}$\end{minipage}
\end{minipage}

\vfill
\begin{tcolorbox}[enhanced,nobeforeafter,width=7.1cm,height=2.1cm,valign=center,colback=ccvp!6,colframe=ccvp,arc=2.5mm,boxrule=0.9pt,left=3mm]{\Large\bfseries\color{ccvp}a)}\quad\begin{minipage}[c]{5.3cm}\raggedright\Large $0{,}5$\end{minipage}\end{tcolorbox}\hfill\begin{tcolorbox}[enhanced,nobeforeafter,width=7.1cm,height=2.1cm,valign=center,colback=ccvp!6,colframe=ccvp,arc=2.5mm,boxrule=0.9pt,left=3mm]{\Large\bfseries\color{ccvp}b)}\quad\begin{minipage}[c]{5.3cm}\raggedright\Large $4{,}5$\end{minipage}\end{tcolorbox}\par\vspace{2.5mm}
\begin{tcolorbox}[enhanced,nobeforeafter,width=7.1cm,height=2.1cm,valign=center,colback=ccvp!6,colframe=ccvp,arc=2.5mm,boxrule=0.9pt,left=3mm]{\Large\bfseries\color{ccvp}c)}\quad\begin{minipage}[c]{5.3cm}\raggedright\Large $2$\end{minipage}\end{tcolorbox}\hfill\begin{tcolorbox}[enhanced,nobeforeafter,width=7.1cm,height=2.1cm,valign=center,colback=ccvp!6,colframe=ccvp,arc=2.5mm,boxrule=0.9pt,left=3mm]{\Large\bfseries\color{ccvp}d)}\quad\begin{minipage}[c]{5.3cm}\raggedright\Large $1{,}5$\end{minipage}\end{tcolorbox}
\end{tcolorbox}
\newpage
\begin{tcolorbox}[enhanced,width=16cm,height=14cm,colback=white,colframe=ccvp,boxrule=2pt,arc=0mm,
left=6mm,right=6mm,top=0mm,bottom=4mm,sharp corners,
overlay={\fill[ccvp] (frame.north west) rectangle ([yshift=-1.25cm]frame.north east);
\node[anchor=west,text=white,font=\sffamily\bfseries\LARGE] at ([xshift=5mm,yshift=-0.62cm]frame.north west) {Question 92 C};
\node[anchor=east,text=white,font=\sffamily\itshape\large] at ([xshift=-5mm,yshift=-0.62cm]frame.north east) {Probabilités conditionnelles et variables aléatoires};}]
\vspace*{1.75cm}
\begin{minipage}[t][6.9cm][c]{\linewidth}\Large\centering
\begin{minipage}{0.95\linewidth}\raggedright\linespread{1.25}\selectfont Voici la loi de probabilité d'une variable aléatoire $X$. Que vaut sa variance $V(X)$~? $\begin{array}{|l|c|c|}\hline x_i & -3 & 3\\\hline P(X=x_i) & 0{,}5 & 0{,}5\\\hline\end{array}$\end{minipage}
\end{minipage}

\vfill
\begin{tcolorbox}[enhanced,nobeforeafter,width=7.1cm,height=2.1cm,valign=center,colback=ccvp!6,colframe=ccvp,arc=2.5mm,boxrule=0.9pt,left=3mm]{\Large\bfseries\color{ccvp}a)}\quad\begin{minipage}[c]{5.3cm}\raggedright\Large $11$\end{minipage}\end{tcolorbox}\hfill\begin{tcolorbox}[enhanced,nobeforeafter,width=7.1cm,height=2.1cm,valign=center,colback=ccvp!6,colframe=ccvp,arc=2.5mm,boxrule=0.9pt,left=3mm]{\Large\bfseries\color{ccvp}b)}\quad\begin{minipage}[c]{5.3cm}\raggedright\Large $10$\end{minipage}\end{tcolorbox}\par\vspace{2.5mm}
\begin{tcolorbox}[enhanced,nobeforeafter,width=7.1cm,height=2.1cm,valign=center,colback=ccvp!6,colframe=ccvp,arc=2.5mm,boxrule=0.9pt,left=3mm]{\Large\bfseries\color{ccvp}c)}\quad\begin{minipage}[c]{5.3cm}\raggedright\Large $0$\end{minipage}\end{tcolorbox}\hfill\begin{tcolorbox}[enhanced,nobeforeafter,width=7.1cm,height=2.1cm,valign=center,colback=ccvp!6,colframe=ccvp,arc=2.5mm,boxrule=0.9pt,left=3mm]{\Large\bfseries\color{ccvp}d)}\quad\begin{minipage}[c]{5.3cm}\raggedright\Large $9$\end{minipage}\end{tcolorbox}
\end{tcolorbox}
\newpage
\begin{tcolorbox}[enhanced,width=16cm,height=14cm,colback=white,colframe=ccvp,boxrule=2pt,arc=0mm,
left=6mm,right=6mm,top=0mm,bottom=4mm,sharp corners,
overlay={\fill[ccvp] (frame.north west) rectangle ([yshift=-1.25cm]frame.north east);
\node[anchor=west,text=white,font=\sffamily\bfseries\LARGE] at ([xshift=5mm,yshift=-0.62cm]frame.north west) {Question 93 C};
\node[anchor=east,text=white,font=\sffamily\itshape\large] at ([xshift=-5mm,yshift=-0.62cm]frame.north east) {Probabilités conditionnelles et variables aléatoires};}]
\vspace*{1.75cm}
\begin{minipage}[t][6.9cm][c]{\linewidth}\Large\centering
\begin{minipage}{0.95\linewidth}\raggedright\linespread{1.25}\selectfont On choisit au hasard une fille de ce groupe. Quelle est la probabilité que ce soit un demi-pensionnaire~? $\begin{array}{|l|c|c|}\hline & \text{Demi-pens.} & \text{Externes}\\\hline \text{Filles} & 9 & 3\\\hline \text{Garçons} & 7 & 21\\\hline\end{array}$\end{minipage}
\end{minipage}

\vfill
\begin{tcolorbox}[enhanced,nobeforeafter,width=7.1cm,height=2.1cm,valign=center,colback=ccvp!6,colframe=ccvp,arc=2.5mm,boxrule=0.9pt,left=3mm]{\Large\bfseries\color{ccvp}a)}\quad\begin{minipage}[c]{5.3cm}\raggedright\Large $\dfrac{3}{4}$\end{minipage}\end{tcolorbox}\hfill\begin{tcolorbox}[enhanced,nobeforeafter,width=7.1cm,height=2.1cm,valign=center,colback=ccvp!6,colframe=ccvp,arc=2.5mm,boxrule=0.9pt,left=3mm]{\Large\bfseries\color{ccvp}b)}\quad\begin{minipage}[c]{5.3cm}\raggedright\Large $\dfrac{9}{16}$\end{minipage}\end{tcolorbox}\par\vspace{2.5mm}
\begin{tcolorbox}[enhanced,nobeforeafter,width=7.1cm,height=2.1cm,valign=center,colback=ccvp!6,colframe=ccvp,arc=2.5mm,boxrule=0.9pt,left=3mm]{\Large\bfseries\color{ccvp}c)}\quad\begin{minipage}[c]{5.3cm}\raggedright\Large $\dfrac{1}{4}$\end{minipage}\end{tcolorbox}\hfill\begin{tcolorbox}[enhanced,nobeforeafter,width=7.1cm,height=2.1cm,valign=center,colback=ccvp!6,colframe=ccvp,arc=2.5mm,boxrule=0.9pt,left=3mm]{\Large\bfseries\color{ccvp}d)}\quad\begin{minipage}[c]{5.3cm}\raggedright\Large $\dfrac{9}{40}$\end{minipage}\end{tcolorbox}
\end{tcolorbox}
\newpage
\begin{tcolorbox}[enhanced,width=16cm,height=14cm,colback=white,colframe=ccvp,boxrule=2pt,arc=0mm,
left=6mm,right=6mm,top=0mm,bottom=4mm,sharp corners,
overlay={\fill[ccvp] (frame.north west) rectangle ([yshift=-1.25cm]frame.north east);
\node[anchor=west,text=white,font=\sffamily\bfseries\LARGE] at ([xshift=5mm,yshift=-0.62cm]frame.north west) {Question 94 C};
\node[anchor=east,text=white,font=\sffamily\itshape\large] at ([xshift=-5mm,yshift=-0.62cm]frame.north east) {Probabilités conditionnelles et variables aléatoires};}]
\vspace*{1.75cm}
\begin{minipage}[t][6.9cm][c]{\linewidth}\Large\centering
\begin{minipage}{0.95\linewidth}\raggedright\linespread{1.25}\selectfont On choisit au hasard un garçon de ce groupe. Quelle est la probabilité que ce soit un demi-pensionnaire~? $\begin{array}{|l|c|c|}\hline & \text{Demi-pens.} & \text{Externes}\\\hline \text{Filles} & 14 & 6\\\hline \text{Garçons} & 8 & 12\\\hline\end{array}$\end{minipage}
\end{minipage}

\vfill
\begin{tcolorbox}[enhanced,nobeforeafter,width=7.1cm,height=2.1cm,valign=center,colback=ccvp!6,colframe=ccvp,arc=2.5mm,boxrule=0.9pt,left=3mm]{\Large\bfseries\color{ccvp}a)}\quad\begin{minipage}[c]{5.3cm}\raggedright\Large $\dfrac{2}{5}$\end{minipage}\end{tcolorbox}\hfill\begin{tcolorbox}[enhanced,nobeforeafter,width=7.1cm,height=2.1cm,valign=center,colback=ccvp!6,colframe=ccvp,arc=2.5mm,boxrule=0.9pt,left=3mm]{\Large\bfseries\color{ccvp}b)}\quad\begin{minipage}[c]{5.3cm}\raggedright\Large $\dfrac{4}{11}$\end{minipage}\end{tcolorbox}\par\vspace{2.5mm}
\begin{tcolorbox}[enhanced,nobeforeafter,width=7.1cm,height=2.1cm,valign=center,colback=ccvp!6,colframe=ccvp,arc=2.5mm,boxrule=0.9pt,left=3mm]{\Large\bfseries\color{ccvp}c)}\quad\begin{minipage}[c]{5.3cm}\raggedright\Large $\dfrac{3}{5}$\end{minipage}\end{tcolorbox}\hfill\begin{tcolorbox}[enhanced,nobeforeafter,width=7.1cm,height=2.1cm,valign=center,colback=ccvp!6,colframe=ccvp,arc=2.5mm,boxrule=0.9pt,left=3mm]{\Large\bfseries\color{ccvp}d)}\quad\begin{minipage}[c]{5.3cm}\raggedright\Large $\dfrac{1}{5}$\end{minipage}\end{tcolorbox}
\end{tcolorbox}
\newpage
\begin{tcolorbox}[enhanced,width=16cm,height=14cm,colback=white,colframe=ccvp,boxrule=2pt,arc=0mm,
left=6mm,right=6mm,top=0mm,bottom=4mm,sharp corners,
overlay={\fill[ccvp] (frame.north west) rectangle ([yshift=-1.25cm]frame.north east);
\node[anchor=west,text=white,font=\sffamily\bfseries\LARGE] at ([xshift=5mm,yshift=-0.62cm]frame.north west) {Question 95 C};
\node[anchor=east,text=white,font=\sffamily\itshape\large] at ([xshift=-5mm,yshift=-0.62cm]frame.north east) {Probabilités conditionnelles et variables aléatoires};}]
\vspace*{1.75cm}
\begin{minipage}[t][6.9cm][c]{\linewidth}\Large\centering
\begin{minipage}{0.95\linewidth}\raggedright\linespread{1.25}\selectfont Voici la loi de probabilité d'une variable aléatoire $X$. Que vaut $P(X\geqslant 2)$~? $\begin{array}{|l|c|c|c|c|c|}\hline x_i & 0 & 1 & 2 & 3 & 4\\\hline P(X=x_i) & 0{,}1 & 0{,}1 & 0{,}6 & 0{,}1 & 0{,}1\\\hline\end{array}$\end{minipage}
\end{minipage}

\vfill
\begin{tcolorbox}[enhanced,nobeforeafter,width=7.1cm,height=2.1cm,valign=center,colback=ccvp!6,colframe=ccvp,arc=2.5mm,boxrule=0.9pt,left=3mm]{\Large\bfseries\color{ccvp}a)}\quad\begin{minipage}[c]{5.3cm}\raggedright\Large $0{,}2$\end{minipage}\end{tcolorbox}\hfill\begin{tcolorbox}[enhanced,nobeforeafter,width=7.1cm,height=2.1cm,valign=center,colback=ccvp!6,colframe=ccvp,arc=2.5mm,boxrule=0.9pt,left=3mm]{\Large\bfseries\color{ccvp}b)}\quad\begin{minipage}[c]{5.3cm}\raggedright\Large $0{,}8$\end{minipage}\end{tcolorbox}\par\vspace{2.5mm}
\begin{tcolorbox}[enhanced,nobeforeafter,width=7.1cm,height=2.1cm,valign=center,colback=ccvp!6,colframe=ccvp,arc=2.5mm,boxrule=0.9pt,left=3mm]{\Large\bfseries\color{ccvp}c)}\quad\begin{minipage}[c]{5.3cm}\raggedright\Large $0{,}6$\end{minipage}\end{tcolorbox}\hfill\begin{tcolorbox}[enhanced,nobeforeafter,width=7.1cm,height=2.1cm,valign=center,colback=ccvp!6,colframe=ccvp,arc=2.5mm,boxrule=0.9pt,left=3mm]{\Large\bfseries\color{ccvp}d)}\quad\begin{minipage}[c]{5.3cm}\raggedright\Large $0{,}9$\end{minipage}\end{tcolorbox}
\end{tcolorbox}
\newpage
\begin{tcolorbox}[enhanced,width=16cm,height=14cm,colback=white,colframe=ccvp,boxrule=2pt,arc=0mm,
left=6mm,right=6mm,top=0mm,bottom=4mm,sharp corners,
overlay={\fill[ccvp] (frame.north west) rectangle ([yshift=-1.25cm]frame.north east);
\node[anchor=west,text=white,font=\sffamily\bfseries\LARGE] at ([xshift=5mm,yshift=-0.62cm]frame.north west) {Question 96 C};
\node[anchor=east,text=white,font=\sffamily\itshape\large] at ([xshift=-5mm,yshift=-0.62cm]frame.north east) {Probabilités conditionnelles et variables aléatoires};}]
\vspace*{1.75cm}
\begin{minipage}[t][6.9cm][c]{\linewidth}\Large\centering
\begin{minipage}{0.95\linewidth}\raggedright\linespread{1.25}\selectfont Voici la loi de probabilité d'une variable aléatoire $X$. Que vaut $P(X\geqslant 4)$~? $\begin{array}{|l|c|c|c|c|c|}\hline x_i & 0 & 1 & 2 & 3 & 4\\\hline P(X=x_i) & 0{,}25 & 0{,}25 & 0{,}2 & 0{,}2 & 0{,}1\\\hline\end{array}$\end{minipage}
\end{minipage}

\vfill
\begin{tcolorbox}[enhanced,nobeforeafter,width=7.1cm,height=2.1cm,valign=center,colback=ccvp!6,colframe=ccvp,arc=2.5mm,boxrule=0.9pt,left=3mm]{\Large\bfseries\color{ccvp}a)}\quad\begin{minipage}[c]{5.3cm}\raggedright\Large $0{,}1$\end{minipage}\end{tcolorbox}\hfill\begin{tcolorbox}[enhanced,nobeforeafter,width=7.1cm,height=2.1cm,valign=center,colback=ccvp!6,colframe=ccvp,arc=2.5mm,boxrule=0.9pt,left=3mm]{\Large\bfseries\color{ccvp}b)}\quad\begin{minipage}[c]{5.3cm}\raggedright\Large $0{,}9$\end{minipage}\end{tcolorbox}\par\vspace{2.5mm}
\begin{tcolorbox}[enhanced,nobeforeafter,width=7.1cm,height=2.1cm,valign=center,colback=ccvp!6,colframe=ccvp,arc=2.5mm,boxrule=0.9pt,left=3mm]{\Large\bfseries\color{ccvp}c)}\quad\begin{minipage}[c]{5.3cm}\raggedright\Large $0$\end{minipage}\end{tcolorbox}\hfill\begin{tcolorbox}[enhanced,nobeforeafter,width=7.1cm,height=2.1cm,valign=center,colback=ccvp!6,colframe=ccvp,arc=2.5mm,boxrule=0.9pt,left=3mm]{\Large\bfseries\color{ccvp}d)}\quad\begin{minipage}[c]{5.3cm}\raggedright\Large $0{,}05$\end{minipage}\end{tcolorbox}
\end{tcolorbox}
\newpage
\begin{tcolorbox}[enhanced,width=16cm,height=14cm,colback=white,colframe=csui,boxrule=2pt,arc=0mm,
left=6mm,right=6mm,top=0mm,bottom=4mm,sharp corners,
overlay={\fill[csui] (frame.north west) rectangle ([yshift=-1.25cm]frame.north east);
\node[anchor=west,text=white,font=\sffamily\bfseries\LARGE] at ([xshift=5mm,yshift=-0.62cm]frame.north west) {Question 1 D};
\node[anchor=east,text=white,font=\sffamily\itshape\large] at ([xshift=-5mm,yshift=-0.62cm]frame.north east) {Suites};}]
\vspace*{1.75cm}
\begin{minipage}[t][6.9cm][c]{\linewidth}\LARGE\centering
\begin{minipage}{0.95\linewidth}\raggedright\linespread{1.25}\selectfont $(u_n)$ est la suite arithmétique de premier terme $u_0=10$ et de raison $-2$. Que vaut $u_{7}$~?\end{minipage}
\end{minipage}

\vfill
\begin{tcolorbox}[enhanced,nobeforeafter,width=7.1cm,height=2.1cm,valign=center,colback=csui!6,colframe=csui,arc=2.5mm,boxrule=0.9pt,left=3mm]{\Large\bfseries\color{csui}a)}\quad\begin{minipage}[c]{5.3cm}\raggedright\Large $-4$\end{minipage}\end{tcolorbox}\hfill\begin{tcolorbox}[enhanced,nobeforeafter,width=7.1cm,height=2.1cm,valign=center,colback=csui!6,colframe=csui,arc=2.5mm,boxrule=0.9pt,left=3mm]{\Large\bfseries\color{csui}b)}\quad\begin{minipage}[c]{5.3cm}\raggedright\Large $-6$\end{minipage}\end{tcolorbox}\par\vspace{2.5mm}
\begin{tcolorbox}[enhanced,nobeforeafter,width=7.1cm,height=2.1cm,valign=center,colback=csui!6,colframe=csui,arc=2.5mm,boxrule=0.9pt,left=3mm]{\Large\bfseries\color{csui}c)}\quad\begin{minipage}[c]{5.3cm}\raggedright\Large $-2$\end{minipage}\end{tcolorbox}\hfill\begin{tcolorbox}[enhanced,nobeforeafter,width=7.1cm,height=2.1cm,valign=center,colback=csui!6,colframe=csui,arc=2.5mm,boxrule=0.9pt,left=3mm]{\Large\bfseries\color{csui}d)}\quad\begin{minipage}[c]{5.3cm}\raggedright\Large $-14$\end{minipage}\end{tcolorbox}
\end{tcolorbox}
\newpage
\begin{tcolorbox}[enhanced,width=16cm,height=14cm,colback=white,colframe=csui,boxrule=2pt,arc=0mm,
left=6mm,right=6mm,top=0mm,bottom=4mm,sharp corners,
overlay={\fill[csui] (frame.north west) rectangle ([yshift=-1.25cm]frame.north east);
\node[anchor=west,text=white,font=\sffamily\bfseries\LARGE] at ([xshift=5mm,yshift=-0.62cm]frame.north west) {Question 2 D};
\node[anchor=east,text=white,font=\sffamily\itshape\large] at ([xshift=-5mm,yshift=-0.62cm]frame.north east) {Suites};}]
\vspace*{1.75cm}
\begin{minipage}[t][6.9cm][c]{\linewidth}\LARGE\centering
\begin{minipage}{0.95\linewidth}\raggedright\linespread{1.25}\selectfont $(u_n)$ est la suite arithmétique de premier terme $u_0=4$ et de raison $6$. Que vaut $u_{13}$~?\end{minipage}
\end{minipage}

\vfill
\begin{tcolorbox}[enhanced,nobeforeafter,width=7.1cm,height=2.1cm,valign=center,colback=csui!6,colframe=csui,arc=2.5mm,boxrule=0.9pt,left=3mm]{\Large\bfseries\color{csui}a)}\quad\begin{minipage}[c]{5.3cm}\raggedright\Large $76$\end{minipage}\end{tcolorbox}\hfill\begin{tcolorbox}[enhanced,nobeforeafter,width=7.1cm,height=2.1cm,valign=center,colback=csui!6,colframe=csui,arc=2.5mm,boxrule=0.9pt,left=3mm]{\Large\bfseries\color{csui}b)}\quad\begin{minipage}[c]{5.3cm}\raggedright\Large $88$\end{minipage}\end{tcolorbox}\par\vspace{2.5mm}
\begin{tcolorbox}[enhanced,nobeforeafter,width=7.1cm,height=2.1cm,valign=center,colback=csui!6,colframe=csui,arc=2.5mm,boxrule=0.9pt,left=3mm]{\Large\bfseries\color{csui}c)}\quad\begin{minipage}[c]{5.3cm}\raggedright\Large $82$\end{minipage}\end{tcolorbox}\hfill\begin{tcolorbox}[enhanced,nobeforeafter,width=7.1cm,height=2.1cm,valign=center,colback=csui!6,colframe=csui,arc=2.5mm,boxrule=0.9pt,left=3mm]{\Large\bfseries\color{csui}d)}\quad\begin{minipage}[c]{5.3cm}\raggedright\Large $78$\end{minipage}\end{tcolorbox}
\end{tcolorbox}
\newpage
\begin{tcolorbox}[enhanced,width=16cm,height=14cm,colback=white,colframe=csui,boxrule=2pt,arc=0mm,
left=6mm,right=6mm,top=0mm,bottom=4mm,sharp corners,
overlay={\fill[csui] (frame.north west) rectangle ([yshift=-1.25cm]frame.north east);
\node[anchor=west,text=white,font=\sffamily\bfseries\LARGE] at ([xshift=5mm,yshift=-0.62cm]frame.north west) {Question 3 D};
\node[anchor=east,text=white,font=\sffamily\itshape\large] at ([xshift=-5mm,yshift=-0.62cm]frame.north east) {Suites};}]
\vspace*{1.75cm}
\begin{minipage}[t][6.9cm][c]{\linewidth}\LARGE\centering
\begin{minipage}{0.95\linewidth}\raggedright\linespread{1.25}\selectfont $(u_n)$ est la suite géométrique de premier terme $u_0=81$ et de raison $\dfrac{1}{3}$. Que vaut $u_{4}$~?\end{minipage}
\end{minipage}

\vfill
\begin{tcolorbox}[enhanced,nobeforeafter,width=7.1cm,height=2.1cm,valign=center,colback=csui!6,colframe=csui,arc=2.5mm,boxrule=0.9pt,left=3mm]{\Large\bfseries\color{csui}a)}\quad\begin{minipage}[c]{5.3cm}\raggedright\Large $\dfrac{247}{3}$\end{minipage}\end{tcolorbox}\hfill\begin{tcolorbox}[enhanced,nobeforeafter,width=7.1cm,height=2.1cm,valign=center,colback=csui!6,colframe=csui,arc=2.5mm,boxrule=0.9pt,left=3mm]{\Large\bfseries\color{csui}b)}\quad\begin{minipage}[c]{5.3cm}\raggedright\Large $1$\end{minipage}\end{tcolorbox}\par\vspace{2.5mm}
\begin{tcolorbox}[enhanced,nobeforeafter,width=7.1cm,height=2.1cm,valign=center,colback=csui!6,colframe=csui,arc=2.5mm,boxrule=0.9pt,left=3mm]{\Large\bfseries\color{csui}c)}\quad\begin{minipage}[c]{5.3cm}\raggedright\Large $3$\end{minipage}\end{tcolorbox}\hfill\begin{tcolorbox}[enhanced,nobeforeafter,width=7.1cm,height=2.1cm,valign=center,colback=csui!6,colframe=csui,arc=2.5mm,boxrule=0.9pt,left=3mm]{\Large\bfseries\color{csui}d)}\quad\begin{minipage}[c]{5.3cm}\raggedright\Large $\dfrac{1}{3}$\end{minipage}\end{tcolorbox}
\end{tcolorbox}
\newpage
\begin{tcolorbox}[enhanced,width=16cm,height=14cm,colback=white,colframe=csui,boxrule=2pt,arc=0mm,
left=6mm,right=6mm,top=0mm,bottom=4mm,sharp corners,
overlay={\fill[csui] (frame.north west) rectangle ([yshift=-1.25cm]frame.north east);
\node[anchor=west,text=white,font=\sffamily\bfseries\LARGE] at ([xshift=5mm,yshift=-0.62cm]frame.north west) {Question 4 D};
\node[anchor=east,text=white,font=\sffamily\itshape\large] at ([xshift=-5mm,yshift=-0.62cm]frame.north east) {Suites};}]
\vspace*{1.75cm}
\begin{minipage}[t][6.9cm][c]{\linewidth}\LARGE\centering
\begin{minipage}{0.95\linewidth}\raggedright\linespread{1.25}\selectfont $(u_n)$ est la suite géométrique de premier terme $u_0=6$ et de raison $-2$. Que vaut $u_{3}$~?\end{minipage}
\end{minipage}

\vfill
\begin{tcolorbox}[enhanced,nobeforeafter,width=7.1cm,height=2.1cm,valign=center,colback=csui!6,colframe=csui,arc=2.5mm,boxrule=0.9pt,left=3mm]{\Large\bfseries\color{csui}a)}\quad\begin{minipage}[c]{5.3cm}\raggedright\Large $24$\end{minipage}\end{tcolorbox}\hfill\begin{tcolorbox}[enhanced,nobeforeafter,width=7.1cm,height=2.1cm,valign=center,colback=csui!6,colframe=csui,arc=2.5mm,boxrule=0.9pt,left=3mm]{\Large\bfseries\color{csui}b)}\quad\begin{minipage}[c]{5.3cm}\raggedright\Large $0$\end{minipage}\end{tcolorbox}\par\vspace{2.5mm}
\begin{tcolorbox}[enhanced,nobeforeafter,width=7.1cm,height=2.1cm,valign=center,colback=csui!6,colframe=csui,arc=2.5mm,boxrule=0.9pt,left=3mm]{\Large\bfseries\color{csui}c)}\quad\begin{minipage}[c]{5.3cm}\raggedright\Large $-48$\end{minipage}\end{tcolorbox}\hfill\begin{tcolorbox}[enhanced,nobeforeafter,width=7.1cm,height=2.1cm,valign=center,colback=csui!6,colframe=csui,arc=2.5mm,boxrule=0.9pt,left=3mm]{\Large\bfseries\color{csui}d)}\quad\begin{minipage}[c]{5.3cm}\raggedright\Large $96$\end{minipage}\end{tcolorbox}
\end{tcolorbox}
\newpage
\begin{tcolorbox}[enhanced,width=16cm,height=14cm,colback=white,colframe=csui,boxrule=2pt,arc=0mm,
left=6mm,right=6mm,top=0mm,bottom=4mm,sharp corners,
overlay={\fill[csui] (frame.north west) rectangle ([yshift=-1.25cm]frame.north east);
\node[anchor=west,text=white,font=\sffamily\bfseries\LARGE] at ([xshift=5mm,yshift=-0.62cm]frame.north west) {Question 5 D};
\node[anchor=east,text=white,font=\sffamily\itshape\large] at ([xshift=-5mm,yshift=-0.62cm]frame.north east) {Suites};}]
\vspace*{1.75cm}
\begin{minipage}[t][6.9cm][c]{\linewidth}\LARGE\centering
\begin{minipage}{0.95\linewidth}\raggedright\linespread{1.25}\selectfont On définit $u_0=-4$ et, pour tout entier $n$, $u_{n+1}=u_n+6$. Que vaut $u_2$~?\end{minipage}
\end{minipage}

\vfill
\begin{tcolorbox}[enhanced,nobeforeafter,width=7.1cm,height=2.1cm,valign=center,colback=csui!6,colframe=csui,arc=2.5mm,boxrule=0.9pt,left=3mm]{\Large\bfseries\color{csui}a)}\quad\begin{minipage}[c]{5.3cm}\raggedright\Large $2$\end{minipage}\end{tcolorbox}\hfill\begin{tcolorbox}[enhanced,nobeforeafter,width=7.1cm,height=2.1cm,valign=center,colback=csui!6,colframe=csui,arc=2.5mm,boxrule=0.9pt,left=3mm]{\Large\bfseries\color{csui}b)}\quad\begin{minipage}[c]{5.3cm}\raggedright\Large $14$\end{minipage}\end{tcolorbox}\par\vspace{2.5mm}
\begin{tcolorbox}[enhanced,nobeforeafter,width=7.1cm,height=2.1cm,valign=center,colback=csui!6,colframe=csui,arc=2.5mm,boxrule=0.9pt,left=3mm]{\Large\bfseries\color{csui}c)}\quad\begin{minipage}[c]{5.3cm}\raggedright\Large $8$\end{minipage}\end{tcolorbox}\hfill\begin{tcolorbox}[enhanced,nobeforeafter,width=7.1cm,height=2.1cm,valign=center,colback=csui!6,colframe=csui,arc=2.5mm,boxrule=0.9pt,left=3mm]{\Large\bfseries\color{csui}d)}\quad\begin{minipage}[c]{5.3cm}\raggedright\Large $9$\end{minipage}\end{tcolorbox}
\end{tcolorbox}
\newpage
\begin{tcolorbox}[enhanced,width=16cm,height=14cm,colback=white,colframe=csui,boxrule=2pt,arc=0mm,
left=6mm,right=6mm,top=0mm,bottom=4mm,sharp corners,
overlay={\fill[csui] (frame.north west) rectangle ([yshift=-1.25cm]frame.north east);
\node[anchor=west,text=white,font=\sffamily\bfseries\LARGE] at ([xshift=5mm,yshift=-0.62cm]frame.north west) {Question 6 D};
\node[anchor=east,text=white,font=\sffamily\itshape\large] at ([xshift=-5mm,yshift=-0.62cm]frame.north east) {Suites};}]
\vspace*{1.75cm}
\begin{minipage}[t][6.9cm][c]{\linewidth}\LARGE\centering
\begin{minipage}{0.95\linewidth}\raggedright\linespread{1.25}\selectfont On définit $u_0=3$ et, pour tout entier $n$, $u_{n+1}=-u_n+10$. Que vaut $u_2$~?\end{minipage}
\end{minipage}

\vfill
\begin{tcolorbox}[enhanced,nobeforeafter,width=7.1cm,height=2.1cm,valign=center,colback=csui!6,colframe=csui,arc=2.5mm,boxrule=0.9pt,left=3mm]{\Large\bfseries\color{csui}a)}\quad\begin{minipage}[c]{5.3cm}\raggedright\Large $7$\end{minipage}\end{tcolorbox}\hfill\begin{tcolorbox}[enhanced,nobeforeafter,width=7.1cm,height=2.1cm,valign=center,colback=csui!6,colframe=csui,arc=2.5mm,boxrule=0.9pt,left=3mm]{\Large\bfseries\color{csui}b)}\quad\begin{minipage}[c]{5.3cm}\raggedright\Large $4$\end{minipage}\end{tcolorbox}\par\vspace{2.5mm}
\begin{tcolorbox}[enhanced,nobeforeafter,width=7.1cm,height=2.1cm,valign=center,colback=csui!6,colframe=csui,arc=2.5mm,boxrule=0.9pt,left=3mm]{\Large\bfseries\color{csui}c)}\quad\begin{minipage}[c]{5.3cm}\raggedright\Large $3$\end{minipage}\end{tcolorbox}\hfill\begin{tcolorbox}[enhanced,nobeforeafter,width=7.1cm,height=2.1cm,valign=center,colback=csui!6,colframe=csui,arc=2.5mm,boxrule=0.9pt,left=3mm]{\Large\bfseries\color{csui}d)}\quad\begin{minipage}[c]{5.3cm}\raggedright\Large $17$\end{minipage}\end{tcolorbox}
\end{tcolorbox}
\newpage
\begin{tcolorbox}[enhanced,width=16cm,height=14cm,colback=white,colframe=csui,boxrule=2pt,arc=0mm,
left=6mm,right=6mm,top=0mm,bottom=4mm,sharp corners,
overlay={\fill[csui] (frame.north west) rectangle ([yshift=-1.25cm]frame.north east);
\node[anchor=west,text=white,font=\sffamily\bfseries\LARGE] at ([xshift=5mm,yshift=-0.62cm]frame.north west) {Question 7 D};
\node[anchor=east,text=white,font=\sffamily\itshape\large] at ([xshift=-5mm,yshift=-0.62cm]frame.north east) {Suites};}]
\vspace*{1.75cm}
\begin{minipage}[t][6.9cm][c]{\linewidth}\LARGE\centering
\begin{minipage}{0.95\linewidth}\raggedright\linespread{1.25}\selectfont Que vaut la somme $1+2+3+\dots+40$~?\end{minipage}
\end{minipage}

\vfill
\begin{tcolorbox}[enhanced,nobeforeafter,width=7.1cm,height=2.1cm,valign=center,colback=csui!6,colframe=csui,arc=2.5mm,boxrule=0.9pt,left=3mm]{\Large\bfseries\color{csui}a)}\quad\begin{minipage}[c]{5.3cm}\raggedright\Large $1\,640$\end{minipage}\end{tcolorbox}\hfill\begin{tcolorbox}[enhanced,nobeforeafter,width=7.1cm,height=2.1cm,valign=center,colback=csui!6,colframe=csui,arc=2.5mm,boxrule=0.9pt,left=3mm]{\Large\bfseries\color{csui}b)}\quad\begin{minipage}[c]{5.3cm}\raggedright\Large $820$\end{minipage}\end{tcolorbox}\par\vspace{2.5mm}
\begin{tcolorbox}[enhanced,nobeforeafter,width=7.1cm,height=2.1cm,valign=center,colback=csui!6,colframe=csui,arc=2.5mm,boxrule=0.9pt,left=3mm]{\Large\bfseries\color{csui}c)}\quad\begin{minipage}[c]{5.3cm}\raggedright\Large $780$\end{minipage}\end{tcolorbox}\hfill\begin{tcolorbox}[enhanced,nobeforeafter,width=7.1cm,height=2.1cm,valign=center,colback=csui!6,colframe=csui,arc=2.5mm,boxrule=0.9pt,left=3mm]{\Large\bfseries\color{csui}d)}\quad\begin{minipage}[c]{5.3cm}\raggedright\Large $1\,600$\end{minipage}\end{tcolorbox}
\end{tcolorbox}
\newpage
\begin{tcolorbox}[enhanced,width=16cm,height=14cm,colback=white,colframe=csui,boxrule=2pt,arc=0mm,
left=6mm,right=6mm,top=0mm,bottom=4mm,sharp corners,
overlay={\fill[csui] (frame.north west) rectangle ([yshift=-1.25cm]frame.north east);
\node[anchor=west,text=white,font=\sffamily\bfseries\LARGE] at ([xshift=5mm,yshift=-0.62cm]frame.north west) {Question 8 D};
\node[anchor=east,text=white,font=\sffamily\itshape\large] at ([xshift=-5mm,yshift=-0.62cm]frame.north east) {Suites};}]
\vspace*{1.75cm}
\begin{minipage}[t][6.9cm][c]{\linewidth}\LARGE\centering
\begin{minipage}{0.95\linewidth}\raggedright\linespread{1.25}\selectfont Que vaut la somme $1+2+3+\dots+12$~?\end{minipage}
\end{minipage}

\vfill
\begin{tcolorbox}[enhanced,nobeforeafter,width=7.1cm,height=2.1cm,valign=center,colback=csui!6,colframe=csui,arc=2.5mm,boxrule=0.9pt,left=3mm]{\Large\bfseries\color{csui}a)}\quad\begin{minipage}[c]{5.3cm}\raggedright\Large $144$\end{minipage}\end{tcolorbox}\hfill\begin{tcolorbox}[enhanced,nobeforeafter,width=7.1cm,height=2.1cm,valign=center,colback=csui!6,colframe=csui,arc=2.5mm,boxrule=0.9pt,left=3mm]{\Large\bfseries\color{csui}b)}\quad\begin{minipage}[c]{5.3cm}\raggedright\Large $66$\end{minipage}\end{tcolorbox}\par\vspace{2.5mm}
\begin{tcolorbox}[enhanced,nobeforeafter,width=7.1cm,height=2.1cm,valign=center,colback=csui!6,colframe=csui,arc=2.5mm,boxrule=0.9pt,left=3mm]{\Large\bfseries\color{csui}c)}\quad\begin{minipage}[c]{5.3cm}\raggedright\Large $156$\end{minipage}\end{tcolorbox}\hfill\begin{tcolorbox}[enhanced,nobeforeafter,width=7.1cm,height=2.1cm,valign=center,colback=csui!6,colframe=csui,arc=2.5mm,boxrule=0.9pt,left=3mm]{\Large\bfseries\color{csui}d)}\quad\begin{minipage}[c]{5.3cm}\raggedright\Large $78$\end{minipage}\end{tcolorbox}
\end{tcolorbox}
\newpage
\begin{tcolorbox}[enhanced,width=16cm,height=14cm,colback=white,colframe=csui,boxrule=2pt,arc=0mm,
left=6mm,right=6mm,top=0mm,bottom=4mm,sharp corners,
overlay={\fill[csui] (frame.north west) rectangle ([yshift=-1.25cm]frame.north east);
\node[anchor=west,text=white,font=\sffamily\bfseries\LARGE] at ([xshift=5mm,yshift=-0.62cm]frame.north west) {Question 9 D};
\node[anchor=east,text=white,font=\sffamily\itshape\large] at ([xshift=-5mm,yshift=-0.62cm]frame.north east) {Suites};}]
\vspace*{1.75cm}
\begin{minipage}[t][6.9cm][c]{\linewidth}\LARGE\centering
\begin{minipage}{0.95\linewidth}\raggedright\linespread{1.25}\selectfont $(u_n)$ est une suite arithmétique avec $u_{4}=9$ et $u_{6}=12$. Quelle est sa raison~?\end{minipage}
\end{minipage}

\vfill
\begin{tcolorbox}[enhanced,nobeforeafter,width=7.1cm,height=2.1cm,valign=center,colback=csui!6,colframe=csui,arc=2.5mm,boxrule=0.9pt,left=3mm]{\Large\bfseries\color{csui}a)}\quad\begin{minipage}[c]{5.3cm}\raggedright\Large $\dfrac{3}{2}$\end{minipage}\end{tcolorbox}\hfill\begin{tcolorbox}[enhanced,nobeforeafter,width=7.1cm,height=2.1cm,valign=center,colback=csui!6,colframe=csui,arc=2.5mm,boxrule=0.9pt,left=3mm]{\Large\bfseries\color{csui}b)}\quad\begin{minipage}[c]{5.3cm}\raggedright\Large $-\dfrac{3}{2}$\end{minipage}\end{tcolorbox}\par\vspace{2.5mm}
\begin{tcolorbox}[enhanced,nobeforeafter,width=7.1cm,height=2.1cm,valign=center,colback=csui!6,colframe=csui,arc=2.5mm,boxrule=0.9pt,left=3mm]{\Large\bfseries\color{csui}c)}\quad\begin{minipage}[c]{5.3cm}\raggedright\Large $3$\end{minipage}\end{tcolorbox}\hfill\begin{tcolorbox}[enhanced,nobeforeafter,width=7.1cm,height=2.1cm,valign=center,colback=csui!6,colframe=csui,arc=2.5mm,boxrule=0.9pt,left=3mm]{\Large\bfseries\color{csui}d)}\quad\begin{minipage}[c]{5.3cm}\raggedright\Large $\dfrac{3}{10}$\end{minipage}\end{tcolorbox}
\end{tcolorbox}
\newpage
\begin{tcolorbox}[enhanced,width=16cm,height=14cm,colback=white,colframe=csui,boxrule=2pt,arc=0mm,
left=6mm,right=6mm,top=0mm,bottom=4mm,sharp corners,
overlay={\fill[csui] (frame.north west) rectangle ([yshift=-1.25cm]frame.north east);
\node[anchor=west,text=white,font=\sffamily\bfseries\LARGE] at ([xshift=5mm,yshift=-0.62cm]frame.north west) {Question 10 D};
\node[anchor=east,text=white,font=\sffamily\itshape\large] at ([xshift=-5mm,yshift=-0.62cm]frame.north east) {Suites};}]
\vspace*{1.75cm}
\begin{minipage}[t][6.9cm][c]{\linewidth}\LARGE\centering
\begin{minipage}{0.95\linewidth}\raggedright\linespread{1.25}\selectfont $(u_n)$ est une suite arithmétique avec $u_{1}=0$ et $u_{6}=-15$. Quelle est sa raison~?\end{minipage}
\end{minipage}

\vfill
\begin{tcolorbox}[enhanced,nobeforeafter,width=7.1cm,height=2.1cm,valign=center,colback=csui!6,colframe=csui,arc=2.5mm,boxrule=0.9pt,left=3mm]{\Large\bfseries\color{csui}a)}\quad\begin{minipage}[c]{5.3cm}\raggedright\Large $3$\end{minipage}\end{tcolorbox}\hfill\begin{tcolorbox}[enhanced,nobeforeafter,width=7.1cm,height=2.1cm,valign=center,colback=csui!6,colframe=csui,arc=2.5mm,boxrule=0.9pt,left=3mm]{\Large\bfseries\color{csui}b)}\quad\begin{minipage}[c]{5.3cm}\raggedright\Large $-3$\end{minipage}\end{tcolorbox}\par\vspace{2.5mm}
\begin{tcolorbox}[enhanced,nobeforeafter,width=7.1cm,height=2.1cm,valign=center,colback=csui!6,colframe=csui,arc=2.5mm,boxrule=0.9pt,left=3mm]{\Large\bfseries\color{csui}c)}\quad\begin{minipage}[c]{5.3cm}\raggedright\Large $-\dfrac{15}{7}$\end{minipage}\end{tcolorbox}\hfill\begin{tcolorbox}[enhanced,nobeforeafter,width=7.1cm,height=2.1cm,valign=center,colback=csui!6,colframe=csui,arc=2.5mm,boxrule=0.9pt,left=3mm]{\Large\bfseries\color{csui}d)}\quad\begin{minipage}[c]{5.3cm}\raggedright\Large $-15$\end{minipage}\end{tcolorbox}
\end{tcolorbox}
\newpage
\begin{tcolorbox}[enhanced,width=16cm,height=14cm,colback=white,colframe=csui,boxrule=2pt,arc=0mm,
left=6mm,right=6mm,top=0mm,bottom=4mm,sharp corners,
overlay={\fill[csui] (frame.north west) rectangle ([yshift=-1.25cm]frame.north east);
\node[anchor=west,text=white,font=\sffamily\bfseries\LARGE] at ([xshift=5mm,yshift=-0.62cm]frame.north west) {Question 11 D};
\node[anchor=east,text=white,font=\sffamily\itshape\large] at ([xshift=-5mm,yshift=-0.62cm]frame.north east) {Suites};}]
\vspace*{1.75cm}
\begin{minipage}[t][6.9cm][c]{\linewidth}\LARGE\centering
\begin{minipage}{0.95\linewidth}\raggedright\linespread{1.25}\selectfont Que vaut la somme $1+10+10^2+\dots+10^{4}$~?\end{minipage}
\end{minipage}

\vfill
\begin{tcolorbox}[enhanced,nobeforeafter,width=7.1cm,height=2.1cm,valign=center,colback=csui!6,colframe=csui,arc=2.5mm,boxrule=0.9pt,left=3mm]{\Large\bfseries\color{csui}a)}\quad\begin{minipage}[c]{5.3cm}\raggedright\Large $11\,111$\end{minipage}\end{tcolorbox}\hfill\begin{tcolorbox}[enhanced,nobeforeafter,width=7.1cm,height=2.1cm,valign=center,colback=csui!6,colframe=csui,arc=2.5mm,boxrule=0.9pt,left=3mm]{\Large\bfseries\color{csui}b)}\quad\begin{minipage}[c]{5.3cm}\raggedright\Large $100\,000$\end{minipage}\end{tcolorbox}\par\vspace{2.5mm}
\begin{tcolorbox}[enhanced,nobeforeafter,width=7.1cm,height=2.1cm,valign=center,colback=csui!6,colframe=csui,arc=2.5mm,boxrule=0.9pt,left=3mm]{\Large\bfseries\color{csui}c)}\quad\begin{minipage}[c]{5.3cm}\raggedright\Large $1\,111$\end{minipage}\end{tcolorbox}\hfill\begin{tcolorbox}[enhanced,nobeforeafter,width=7.1cm,height=2.1cm,valign=center,colback=csui!6,colframe=csui,arc=2.5mm,boxrule=0.9pt,left=3mm]{\Large\bfseries\color{csui}d)}\quad\begin{minipage}[c]{5.3cm}\raggedright\Large $99\,999$\end{minipage}\end{tcolorbox}
\end{tcolorbox}
\newpage
\begin{tcolorbox}[enhanced,width=16cm,height=14cm,colback=white,colframe=csui,boxrule=2pt,arc=0mm,
left=6mm,right=6mm,top=0mm,bottom=4mm,sharp corners,
overlay={\fill[csui] (frame.north west) rectangle ([yshift=-1.25cm]frame.north east);
\node[anchor=west,text=white,font=\sffamily\bfseries\LARGE] at ([xshift=5mm,yshift=-0.62cm]frame.north west) {Question 12 D};
\node[anchor=east,text=white,font=\sffamily\itshape\large] at ([xshift=-5mm,yshift=-0.62cm]frame.north east) {Suites};}]
\vspace*{1.75cm}
\begin{minipage}[t][6.9cm][c]{\linewidth}\LARGE\centering
\begin{minipage}{0.95\linewidth}\raggedright\linespread{1.25}\selectfont Que vaut la somme $1+2+2^2+\dots+2^{4}$~?\end{minipage}
\end{minipage}

\vfill
\begin{tcolorbox}[enhanced,nobeforeafter,width=7.1cm,height=2.1cm,valign=center,colback=csui!6,colframe=csui,arc=2.5mm,boxrule=0.9pt,left=3mm]{\Large\bfseries\color{csui}a)}\quad\begin{minipage}[c]{5.3cm}\raggedright\Large $15$\end{minipage}\end{tcolorbox}\hfill\begin{tcolorbox}[enhanced,nobeforeafter,width=7.1cm,height=2.1cm,valign=center,colback=csui!6,colframe=csui,arc=2.5mm,boxrule=0.9pt,left=3mm]{\Large\bfseries\color{csui}b)}\quad\begin{minipage}[c]{5.3cm}\raggedright\Large $31$\end{minipage}\end{tcolorbox}\par\vspace{2.5mm}
\begin{tcolorbox}[enhanced,nobeforeafter,width=7.1cm,height=2.1cm,valign=center,colback=csui!6,colframe=csui,arc=2.5mm,boxrule=0.9pt,left=3mm]{\Large\bfseries\color{csui}c)}\quad\begin{minipage}[c]{5.3cm}\raggedright\Large $32$\end{minipage}\end{tcolorbox}\hfill\begin{tcolorbox}[enhanced,nobeforeafter,width=7.1cm,height=2.1cm,valign=center,colback=csui!6,colframe=csui,arc=2.5mm,boxrule=0.9pt,left=3mm]{\Large\bfseries\color{csui}d)}\quad\begin{minipage}[c]{5.3cm}\raggedright\Large $16$\end{minipage}\end{tcolorbox}
\end{tcolorbox}
\newpage
\begin{tcolorbox}[enhanced,width=16cm,height=14cm,colback=white,colframe=csui,boxrule=2pt,arc=0mm,
left=6mm,right=6mm,top=0mm,bottom=4mm,sharp corners,
overlay={\fill[csui] (frame.north west) rectangle ([yshift=-1.25cm]frame.north east);
\node[anchor=west,text=white,font=\sffamily\bfseries\LARGE] at ([xshift=5mm,yshift=-0.62cm]frame.north west) {Question 13 D};
\node[anchor=east,text=white,font=\sffamily\itshape\large] at ([xshift=-5mm,yshift=-0.62cm]frame.north east) {Suites};}]
\vspace*{1.75cm}
\begin{minipage}[t][6.9cm][c]{\linewidth}\LARGE\centering
\begin{minipage}{0.95\linewidth}\raggedright\linespread{1.25}\selectfont Quel est le sens de variation de la suite définie pour tout entier naturel $n$ par $u_n=\dfrac{1}{n+1}$~?\end{minipage}
\end{minipage}

\vfill
\begin{tcolorbox}[enhanced,nobeforeafter,width=7.1cm,height=2.1cm,valign=center,colback=csui!6,colframe=csui,arc=2.5mm,boxrule=0.9pt,left=3mm]{\Large\bfseries\color{csui}a)}\quad\begin{minipage}[c]{5.3cm}\raggedright\Large croissante\end{minipage}\end{tcolorbox}\hfill\begin{tcolorbox}[enhanced,nobeforeafter,width=7.1cm,height=2.1cm,valign=center,colback=csui!6,colframe=csui,arc=2.5mm,boxrule=0.9pt,left=3mm]{\Large\bfseries\color{csui}b)}\quad\begin{minipage}[c]{5.3cm}\raggedright\large ni croissante ni décroissante\end{minipage}\end{tcolorbox}\par\vspace{2.5mm}
\begin{tcolorbox}[enhanced,nobeforeafter,width=7.1cm,height=2.1cm,valign=center,colback=csui!6,colframe=csui,arc=2.5mm,boxrule=0.9pt,left=3mm]{\Large\bfseries\color{csui}c)}\quad\begin{minipage}[c]{5.3cm}\raggedright\Large constante\end{minipage}\end{tcolorbox}\hfill\begin{tcolorbox}[enhanced,nobeforeafter,width=7.1cm,height=2.1cm,valign=center,colback=csui!6,colframe=csui,arc=2.5mm,boxrule=0.9pt,left=3mm]{\Large\bfseries\color{csui}d)}\quad\begin{minipage}[c]{5.3cm}\raggedright\Large décroissante\end{minipage}\end{tcolorbox}
\end{tcolorbox}
\newpage
\begin{tcolorbox}[enhanced,width=16cm,height=14cm,colback=white,colframe=csui,boxrule=2pt,arc=0mm,
left=6mm,right=6mm,top=0mm,bottom=4mm,sharp corners,
overlay={\fill[csui] (frame.north west) rectangle ([yshift=-1.25cm]frame.north east);
\node[anchor=west,text=white,font=\sffamily\bfseries\LARGE] at ([xshift=5mm,yshift=-0.62cm]frame.north west) {Question 14 D};
\node[anchor=east,text=white,font=\sffamily\itshape\large] at ([xshift=-5mm,yshift=-0.62cm]frame.north east) {Suites};}]
\vspace*{1.75cm}
\begin{minipage}[t][6.9cm][c]{\linewidth}\LARGE\centering
\begin{minipage}{0.95\linewidth}\raggedright\linespread{1.25}\selectfont Quel est le sens de variation de la suite définie pour tout entier naturel $n$ par $u_n=7$~?\end{minipage}
\end{minipage}

\vfill
\begin{tcolorbox}[enhanced,nobeforeafter,width=7.1cm,height=2.1cm,valign=center,colback=csui!6,colframe=csui,arc=2.5mm,boxrule=0.9pt,left=3mm]{\Large\bfseries\color{csui}a)}\quad\begin{minipage}[c]{5.3cm}\raggedright\large ni croissante ni décroissante\end{minipage}\end{tcolorbox}\hfill\begin{tcolorbox}[enhanced,nobeforeafter,width=7.1cm,height=2.1cm,valign=center,colback=csui!6,colframe=csui,arc=2.5mm,boxrule=0.9pt,left=3mm]{\Large\bfseries\color{csui}b)}\quad\begin{minipage}[c]{5.3cm}\raggedright\Large décroissante\end{minipage}\end{tcolorbox}\par\vspace{2.5mm}
\begin{tcolorbox}[enhanced,nobeforeafter,width=7.1cm,height=2.1cm,valign=center,colback=csui!6,colframe=csui,arc=2.5mm,boxrule=0.9pt,left=3mm]{\Large\bfseries\color{csui}c)}\quad\begin{minipage}[c]{5.3cm}\raggedright\Large croissante\end{minipage}\end{tcolorbox}\hfill\begin{tcolorbox}[enhanced,nobeforeafter,width=7.1cm,height=2.1cm,valign=center,colback=csui!6,colframe=csui,arc=2.5mm,boxrule=0.9pt,left=3mm]{\Large\bfseries\color{csui}d)}\quad\begin{minipage}[c]{5.3cm}\raggedright\Large constante\end{minipage}\end{tcolorbox}
\end{tcolorbox}
\newpage
\begin{tcolorbox}[enhanced,width=16cm,height=14cm,colback=white,colframe=csui,boxrule=2pt,arc=0mm,
left=6mm,right=6mm,top=0mm,bottom=4mm,sharp corners,
overlay={\fill[csui] (frame.north west) rectangle ([yshift=-1.25cm]frame.north east);
\node[anchor=west,text=white,font=\sffamily\bfseries\LARGE] at ([xshift=5mm,yshift=-0.62cm]frame.north west) {Question 15 D};
\node[anchor=east,text=white,font=\sffamily\itshape\large] at ([xshift=-5mm,yshift=-0.62cm]frame.north east) {Suites};}]
\vspace*{1.75cm}
\begin{minipage}[t][6.9cm][c]{\linewidth}\LARGE\centering
\begin{minipage}{0.95\linewidth}\raggedright\linespread{1.25}\selectfont Pour tout entier naturel $n$, $u_n=-2n^2+3n+7$. Que vaut $u_{3}$~?\end{minipage}
\end{minipage}

\vfill
\begin{tcolorbox}[enhanced,nobeforeafter,width=7.1cm,height=2.1cm,valign=center,colback=csui!6,colframe=csui,arc=2.5mm,boxrule=0.9pt,left=3mm]{\Large\bfseries\color{csui}a)}\quad\begin{minipage}[c]{5.3cm}\raggedright\Large $-13$\end{minipage}\end{tcolorbox}\hfill\begin{tcolorbox}[enhanced,nobeforeafter,width=7.1cm,height=2.1cm,valign=center,colback=csui!6,colframe=csui,arc=2.5mm,boxrule=0.9pt,left=3mm]{\Large\bfseries\color{csui}b)}\quad\begin{minipage}[c]{5.3cm}\raggedright\Large $-2$\end{minipage}\end{tcolorbox}\par\vspace{2.5mm}
\begin{tcolorbox}[enhanced,nobeforeafter,width=7.1cm,height=2.1cm,valign=center,colback=csui!6,colframe=csui,arc=2.5mm,boxrule=0.9pt,left=3mm]{\Large\bfseries\color{csui}c)}\quad\begin{minipage}[c]{5.3cm}\raggedright\Large $5$\end{minipage}\end{tcolorbox}\hfill\begin{tcolorbox}[enhanced,nobeforeafter,width=7.1cm,height=2.1cm,valign=center,colback=csui!6,colframe=csui,arc=2.5mm,boxrule=0.9pt,left=3mm]{\Large\bfseries\color{csui}d)}\quad\begin{minipage}[c]{5.3cm}\raggedright\Large $2$\end{minipage}\end{tcolorbox}
\end{tcolorbox}
\newpage
\begin{tcolorbox}[enhanced,width=16cm,height=14cm,colback=white,colframe=csui,boxrule=2pt,arc=0mm,
left=6mm,right=6mm,top=0mm,bottom=4mm,sharp corners,
overlay={\fill[csui] (frame.north west) rectangle ([yshift=-1.25cm]frame.north east);
\node[anchor=west,text=white,font=\sffamily\bfseries\LARGE] at ([xshift=5mm,yshift=-0.62cm]frame.north west) {Question 16 D};
\node[anchor=east,text=white,font=\sffamily\itshape\large] at ([xshift=-5mm,yshift=-0.62cm]frame.north east) {Suites};}]
\vspace*{1.75cm}
\begin{minipage}[t][6.9cm][c]{\linewidth}\LARGE\centering
\begin{minipage}{0.95\linewidth}\raggedright\linespread{1.25}\selectfont Pour tout entier naturel $n$, $u_n=n^2+4n-1$. Que vaut $u_{2}$~?\end{minipage}
\end{minipage}

\vfill
\begin{tcolorbox}[enhanced,nobeforeafter,width=7.1cm,height=2.1cm,valign=center,colback=csui!6,colframe=csui,arc=2.5mm,boxrule=0.9pt,left=3mm]{\Large\bfseries\color{csui}a)}\quad\begin{minipage}[c]{5.3cm}\raggedright\Large $4$\end{minipage}\end{tcolorbox}\hfill\begin{tcolorbox}[enhanced,nobeforeafter,width=7.1cm,height=2.1cm,valign=center,colback=csui!6,colframe=csui,arc=2.5mm,boxrule=0.9pt,left=3mm]{\Large\bfseries\color{csui}b)}\quad\begin{minipage}[c]{5.3cm}\raggedright\Large $20$\end{minipage}\end{tcolorbox}\par\vspace{2.5mm}
\begin{tcolorbox}[enhanced,nobeforeafter,width=7.1cm,height=2.1cm,valign=center,colback=csui!6,colframe=csui,arc=2.5mm,boxrule=0.9pt,left=3mm]{\Large\bfseries\color{csui}c)}\quad\begin{minipage}[c]{5.3cm}\raggedright\Large $11$\end{minipage}\end{tcolorbox}\hfill\begin{tcolorbox}[enhanced,nobeforeafter,width=7.1cm,height=2.1cm,valign=center,colback=csui!6,colframe=csui,arc=2.5mm,boxrule=0.9pt,left=3mm]{\Large\bfseries\color{csui}d)}\quad\begin{minipage}[c]{5.3cm}\raggedright\Large $-11$\end{minipage}\end{tcolorbox}
\end{tcolorbox}
\newpage
\begin{tcolorbox}[enhanced,width=16cm,height=14cm,colback=white,colframe=cdeg,boxrule=2pt,arc=0mm,
left=6mm,right=6mm,top=0mm,bottom=4mm,sharp corners,
overlay={\fill[cdeg] (frame.north west) rectangle ([yshift=-1.25cm]frame.north east);
\node[anchor=west,text=white,font=\sffamily\bfseries\LARGE] at ([xshift=5mm,yshift=-0.62cm]frame.north west) {Question 17 D};
\node[anchor=east,text=white,font=\sffamily\itshape\large] at ([xshift=-5mm,yshift=-0.62cm]frame.north east) {Second degré};}]
\vspace*{1.75cm}
\begin{minipage}[t][6.9cm][c]{\linewidth}\LARGE\centering
\begin{minipage}{0.95\linewidth}\raggedright\linespread{1.25}\selectfont Quel est le discriminant du trinôme $x^2+x+1$~?\end{minipage}
\end{minipage}

\vfill
\begin{tcolorbox}[enhanced,nobeforeafter,width=7.1cm,height=2.1cm,valign=center,colback=cdeg!6,colframe=cdeg,arc=2.5mm,boxrule=0.9pt,left=3mm]{\Large\bfseries\color{cdeg}a)}\quad\begin{minipage}[c]{5.3cm}\raggedright\Large $5$\end{minipage}\end{tcolorbox}\hfill\begin{tcolorbox}[enhanced,nobeforeafter,width=7.1cm,height=2.1cm,valign=center,colback=cdeg!6,colframe=cdeg,arc=2.5mm,boxrule=0.9pt,left=3mm]{\Large\bfseries\color{cdeg}b)}\quad\begin{minipage}[c]{5.3cm}\raggedright\Large $-3$\end{minipage}\end{tcolorbox}\par\vspace{2.5mm}
\begin{tcolorbox}[enhanced,nobeforeafter,width=7.1cm,height=2.1cm,valign=center,colback=cdeg!6,colframe=cdeg,arc=2.5mm,boxrule=0.9pt,left=3mm]{\Large\bfseries\color{cdeg}c)}\quad\begin{minipage}[c]{5.3cm}\raggedright\Large $-1$\end{minipage}\end{tcolorbox}\hfill\begin{tcolorbox}[enhanced,nobeforeafter,width=7.1cm,height=2.1cm,valign=center,colback=cdeg!6,colframe=cdeg,arc=2.5mm,boxrule=0.9pt,left=3mm]{\Large\bfseries\color{cdeg}d)}\quad\begin{minipage}[c]{5.3cm}\raggedright\Large $3$\end{minipage}\end{tcolorbox}
\end{tcolorbox}
\newpage
\begin{tcolorbox}[enhanced,width=16cm,height=14cm,colback=white,colframe=cdeg,boxrule=2pt,arc=0mm,
left=6mm,right=6mm,top=0mm,bottom=4mm,sharp corners,
overlay={\fill[cdeg] (frame.north west) rectangle ([yshift=-1.25cm]frame.north east);
\node[anchor=west,text=white,font=\sffamily\bfseries\LARGE] at ([xshift=5mm,yshift=-0.62cm]frame.north west) {Question 18 D};
\node[anchor=east,text=white,font=\sffamily\itshape\large] at ([xshift=-5mm,yshift=-0.62cm]frame.north east) {Second degré};}]
\vspace*{1.75cm}
\begin{minipage}[t][6.9cm][c]{\linewidth}\LARGE\centering
\begin{minipage}{0.95\linewidth}\raggedright\linespread{1.25}\selectfont Quel est le discriminant du trinôme $4x^2-4x-3$~?\end{minipage}
\end{minipage}

\vfill
\begin{tcolorbox}[enhanced,nobeforeafter,width=7.1cm,height=2.1cm,valign=center,colback=cdeg!6,colframe=cdeg,arc=2.5mm,boxrule=0.9pt,left=3mm]{\Large\bfseries\color{cdeg}a)}\quad\begin{minipage}[c]{5.3cm}\raggedright\Large $-64$\end{minipage}\end{tcolorbox}\hfill\begin{tcolorbox}[enhanced,nobeforeafter,width=7.1cm,height=2.1cm,valign=center,colback=cdeg!6,colframe=cdeg,arc=2.5mm,boxrule=0.9pt,left=3mm]{\Large\bfseries\color{cdeg}b)}\quad\begin{minipage}[c]{5.3cm}\raggedright\Large $44$\end{minipage}\end{tcolorbox}\par\vspace{2.5mm}
\begin{tcolorbox}[enhanced,nobeforeafter,width=7.1cm,height=2.1cm,valign=center,colback=cdeg!6,colframe=cdeg,arc=2.5mm,boxrule=0.9pt,left=3mm]{\Large\bfseries\color{cdeg}c)}\quad\begin{minipage}[c]{5.3cm}\raggedright\Large $-32$\end{minipage}\end{tcolorbox}\hfill\begin{tcolorbox}[enhanced,nobeforeafter,width=7.1cm,height=2.1cm,valign=center,colback=cdeg!6,colframe=cdeg,arc=2.5mm,boxrule=0.9pt,left=3mm]{\Large\bfseries\color{cdeg}d)}\quad\begin{minipage}[c]{5.3cm}\raggedright\Large $64$\end{minipage}\end{tcolorbox}
\end{tcolorbox}
\newpage
\begin{tcolorbox}[enhanced,width=16cm,height=14cm,colback=white,colframe=cdeg,boxrule=2pt,arc=0mm,
left=6mm,right=6mm,top=0mm,bottom=4mm,sharp corners,
overlay={\fill[cdeg] (frame.north west) rectangle ([yshift=-1.25cm]frame.north east);
\node[anchor=west,text=white,font=\sffamily\bfseries\LARGE] at ([xshift=5mm,yshift=-0.62cm]frame.north west) {Question 19 D};
\node[anchor=east,text=white,font=\sffamily\itshape\large] at ([xshift=-5mm,yshift=-0.62cm]frame.north east) {Second degré};}]
\vspace*{1.75cm}
\begin{minipage}[t][6.9cm][c]{\linewidth}\LARGE\centering
\begin{minipage}{0.95\linewidth}\raggedright\linespread{1.25}\selectfont Combien de racines réelles le trinôme $5x^2-x-2$ possède-t-il~?\end{minipage}
\end{minipage}

\vfill
\begin{tcolorbox}[enhanced,nobeforeafter,width=7.1cm,height=2.1cm,valign=center,colback=cdeg!6,colframe=cdeg,arc=2.5mm,boxrule=0.9pt,left=3mm]{\Large\bfseries\color{cdeg}a)}\quad\begin{minipage}[c]{5.3cm}\raggedright\Large une infinité\end{minipage}\end{tcolorbox}\hfill\begin{tcolorbox}[enhanced,nobeforeafter,width=7.1cm,height=2.1cm,valign=center,colback=cdeg!6,colframe=cdeg,arc=2.5mm,boxrule=0.9pt,left=3mm]{\Large\bfseries\color{cdeg}b)}\quad\begin{minipage}[c]{5.3cm}\raggedright\Large $1$\end{minipage}\end{tcolorbox}\par\vspace{2.5mm}
\begin{tcolorbox}[enhanced,nobeforeafter,width=7.1cm,height=2.1cm,valign=center,colback=cdeg!6,colframe=cdeg,arc=2.5mm,boxrule=0.9pt,left=3mm]{\Large\bfseries\color{cdeg}c)}\quad\begin{minipage}[c]{5.3cm}\raggedright\Large $2$\end{minipage}\end{tcolorbox}\hfill\begin{tcolorbox}[enhanced,nobeforeafter,width=7.1cm,height=2.1cm,valign=center,colback=cdeg!6,colframe=cdeg,arc=2.5mm,boxrule=0.9pt,left=3mm]{\Large\bfseries\color{cdeg}d)}\quad\begin{minipage}[c]{5.3cm}\raggedright\Large $0$\end{minipage}\end{tcolorbox}
\end{tcolorbox}
\newpage
\begin{tcolorbox}[enhanced,width=16cm,height=14cm,colback=white,colframe=cdeg,boxrule=2pt,arc=0mm,
left=6mm,right=6mm,top=0mm,bottom=4mm,sharp corners,
overlay={\fill[cdeg] (frame.north west) rectangle ([yshift=-1.25cm]frame.north east);
\node[anchor=west,text=white,font=\sffamily\bfseries\LARGE] at ([xshift=5mm,yshift=-0.62cm]frame.north west) {Question 20 D};
\node[anchor=east,text=white,font=\sffamily\itshape\large] at ([xshift=-5mm,yshift=-0.62cm]frame.north east) {Second degré};}]
\vspace*{1.75cm}
\begin{minipage}[t][6.9cm][c]{\linewidth}\LARGE\centering
\begin{minipage}{0.95\linewidth}\raggedright\linespread{1.25}\selectfont Combien de racines réelles le trinôme $2x^2-8x+8$ possède-t-il~?\end{minipage}
\end{minipage}

\vfill
\begin{tcolorbox}[enhanced,nobeforeafter,width=7.1cm,height=2.1cm,valign=center,colback=cdeg!6,colframe=cdeg,arc=2.5mm,boxrule=0.9pt,left=3mm]{\Large\bfseries\color{cdeg}a)}\quad\begin{minipage}[c]{5.3cm}\raggedright\Large une infinité\end{minipage}\end{tcolorbox}\hfill\begin{tcolorbox}[enhanced,nobeforeafter,width=7.1cm,height=2.1cm,valign=center,colback=cdeg!6,colframe=cdeg,arc=2.5mm,boxrule=0.9pt,left=3mm]{\Large\bfseries\color{cdeg}b)}\quad\begin{minipage}[c]{5.3cm}\raggedright\Large $0$\end{minipage}\end{tcolorbox}\par\vspace{2.5mm}
\begin{tcolorbox}[enhanced,nobeforeafter,width=7.1cm,height=2.1cm,valign=center,colback=cdeg!6,colframe=cdeg,arc=2.5mm,boxrule=0.9pt,left=3mm]{\Large\bfseries\color{cdeg}c)}\quad\begin{minipage}[c]{5.3cm}\raggedright\Large $2$\end{minipage}\end{tcolorbox}\hfill\begin{tcolorbox}[enhanced,nobeforeafter,width=7.1cm,height=2.1cm,valign=center,colback=cdeg!6,colframe=cdeg,arc=2.5mm,boxrule=0.9pt,left=3mm]{\Large\bfseries\color{cdeg}d)}\quad\begin{minipage}[c]{5.3cm}\raggedright\Large $1$\end{minipage}\end{tcolorbox}
\end{tcolorbox}
\newpage
\begin{tcolorbox}[enhanced,width=16cm,height=14cm,colback=white,colframe=cdeg,boxrule=2pt,arc=0mm,
left=6mm,right=6mm,top=0mm,bottom=4mm,sharp corners,
overlay={\fill[cdeg] (frame.north west) rectangle ([yshift=-1.25cm]frame.north east);
\node[anchor=west,text=white,font=\sffamily\bfseries\LARGE] at ([xshift=5mm,yshift=-0.62cm]frame.north west) {Question 21 D};
\node[anchor=east,text=white,font=\sffamily\itshape\large] at ([xshift=-5mm,yshift=-0.62cm]frame.north east) {Second degré};}]
\vspace*{1.75cm}
\begin{minipage}[t][6.9cm][c]{\linewidth}\LARGE\centering
\begin{minipage}{0.95\linewidth}\raggedright\linespread{1.25}\selectfont Quelles sont les racines du trinôme $x^2-16$~?\end{minipage}
\end{minipage}

\vfill
\begin{tcolorbox}[enhanced,nobeforeafter,width=7.1cm,height=2.1cm,valign=center,colback=cdeg!6,colframe=cdeg,arc=2.5mm,boxrule=0.9pt,left=3mm]{\Large\bfseries\color{cdeg}a)}\quad\begin{minipage}[c]{5.3cm}\raggedright\Large $\{-3\,;\,5\}$\end{minipage}\end{tcolorbox}\hfill\begin{tcolorbox}[enhanced,nobeforeafter,width=7.1cm,height=2.1cm,valign=center,colback=cdeg!6,colframe=cdeg,arc=2.5mm,boxrule=0.9pt,left=3mm]{\Large\bfseries\color{cdeg}b)}\quad\begin{minipage}[c]{5.3cm}\raggedright\Large $\{-4\,;\,4\}$\end{minipage}\end{tcolorbox}\par\vspace{2.5mm}
\begin{tcolorbox}[enhanced,nobeforeafter,width=7.1cm,height=2.1cm,valign=center,colback=cdeg!6,colframe=cdeg,arc=2.5mm,boxrule=0.9pt,left=3mm]{\Large\bfseries\color{cdeg}c)}\quad\begin{minipage}[c]{5.3cm}\raggedright\Large $\{-3\,;\,4\}$\end{minipage}\end{tcolorbox}\hfill\begin{tcolorbox}[enhanced,nobeforeafter,width=7.1cm,height=2.1cm,valign=center,colback=cdeg!6,colframe=cdeg,arc=2.5mm,boxrule=0.9pt,left=3mm]{\Large\bfseries\color{cdeg}d)}\quad\begin{minipage}[c]{5.3cm}\raggedright\Large $\{-16\,;\,0\}$\end{minipage}\end{tcolorbox}
\end{tcolorbox}
\newpage
\begin{tcolorbox}[enhanced,width=16cm,height=14cm,colback=white,colframe=cdeg,boxrule=2pt,arc=0mm,
left=6mm,right=6mm,top=0mm,bottom=4mm,sharp corners,
overlay={\fill[cdeg] (frame.north west) rectangle ([yshift=-1.25cm]frame.north east);
\node[anchor=west,text=white,font=\sffamily\bfseries\LARGE] at ([xshift=5mm,yshift=-0.62cm]frame.north west) {Question 22 D};
\node[anchor=east,text=white,font=\sffamily\itshape\large] at ([xshift=-5mm,yshift=-0.62cm]frame.north east) {Second degré};}]
\vspace*{1.75cm}
\begin{minipage}[t][6.9cm][c]{\linewidth}\LARGE\centering
\begin{minipage}{0.95\linewidth}\raggedright\linespread{1.25}\selectfont Quelles sont les racines du trinôme $x^2-7x$~?\end{minipage}
\end{minipage}

\vfill
\begin{tcolorbox}[enhanced,nobeforeafter,width=7.1cm,height=2.1cm,valign=center,colback=cdeg!6,colframe=cdeg,arc=2.5mm,boxrule=0.9pt,left=3mm]{\Large\bfseries\color{cdeg}a)}\quad\begin{minipage}[c]{5.3cm}\raggedright\Large $\{1\,;\,8\}$\end{minipage}\end{tcolorbox}\hfill\begin{tcolorbox}[enhanced,nobeforeafter,width=7.1cm,height=2.1cm,valign=center,colback=cdeg!6,colframe=cdeg,arc=2.5mm,boxrule=0.9pt,left=3mm]{\Large\bfseries\color{cdeg}b)}\quad\begin{minipage}[c]{5.3cm}\raggedright\Large $\{-1\,;\,8\}$\end{minipage}\end{tcolorbox}\par\vspace{2.5mm}
\begin{tcolorbox}[enhanced,nobeforeafter,width=7.1cm,height=2.1cm,valign=center,colback=cdeg!6,colframe=cdeg,arc=2.5mm,boxrule=0.9pt,left=3mm]{\Large\bfseries\color{cdeg}c)}\quad\begin{minipage}[c]{5.3cm}\raggedright\Large $\{-7\,;\,0\}$\end{minipage}\end{tcolorbox}\hfill\begin{tcolorbox}[enhanced,nobeforeafter,width=7.1cm,height=2.1cm,valign=center,colback=cdeg!6,colframe=cdeg,arc=2.5mm,boxrule=0.9pt,left=3mm]{\Large\bfseries\color{cdeg}d)}\quad\begin{minipage}[c]{5.3cm}\raggedright\Large $\{0\,;\,7\}$\end{minipage}\end{tcolorbox}
\end{tcolorbox}
\newpage
\begin{tcolorbox}[enhanced,width=16cm,height=14cm,colback=white,colframe=cdeg,boxrule=2pt,arc=0mm,
left=6mm,right=6mm,top=0mm,bottom=4mm,sharp corners,
overlay={\fill[cdeg] (frame.north west) rectangle ([yshift=-1.25cm]frame.north east);
\node[anchor=west,text=white,font=\sffamily\bfseries\LARGE] at ([xshift=5mm,yshift=-0.62cm]frame.north west) {Question 23 D};
\node[anchor=east,text=white,font=\sffamily\itshape\large] at ([xshift=-5mm,yshift=-0.62cm]frame.north east) {Second degré};}]
\vspace*{1.75cm}
\begin{minipage}[t][6.9cm][c]{\linewidth}\LARGE\centering
\begin{minipage}{0.95\linewidth}\raggedright\linespread{1.25}\selectfont Quelle est la forme canonique de $x^2+4x+9$~?\end{minipage}
\end{minipage}

\vfill
\begin{tcolorbox}[enhanced,nobeforeafter,width=7.1cm,height=2.1cm,valign=center,colback=cdeg!6,colframe=cdeg,arc=2.5mm,boxrule=0.9pt,left=3mm]{\Large\bfseries\color{cdeg}a)}\quad\begin{minipage}[c]{5.3cm}\raggedright\Large $(x-2)^2+5$\end{minipage}\end{tcolorbox}\hfill\begin{tcolorbox}[enhanced,nobeforeafter,width=7.1cm,height=2.1cm,valign=center,colback=cdeg!6,colframe=cdeg,arc=2.5mm,boxrule=0.9pt,left=3mm]{\Large\bfseries\color{cdeg}b)}\quad\begin{minipage}[c]{5.3cm}\raggedright\Large $(x+2)^2+5$\end{minipage}\end{tcolorbox}\par\vspace{2.5mm}
\begin{tcolorbox}[enhanced,nobeforeafter,width=7.1cm,height=2.1cm,valign=center,colback=cdeg!6,colframe=cdeg,arc=2.5mm,boxrule=0.9pt,left=3mm]{\Large\bfseries\color{cdeg}c)}\quad\begin{minipage}[c]{5.3cm}\raggedright\Large $(x+4)^2+5$\end{minipage}\end{tcolorbox}\hfill\begin{tcolorbox}[enhanced,nobeforeafter,width=7.1cm,height=2.1cm,valign=center,colback=cdeg!6,colframe=cdeg,arc=2.5mm,boxrule=0.9pt,left=3mm]{\Large\bfseries\color{cdeg}d)}\quad\begin{minipage}[c]{5.3cm}\raggedright\Large $(x+2)^2+9$\end{minipage}\end{tcolorbox}
\end{tcolorbox}
\newpage
\begin{tcolorbox}[enhanced,width=16cm,height=14cm,colback=white,colframe=cdeg,boxrule=2pt,arc=0mm,
left=6mm,right=6mm,top=0mm,bottom=4mm,sharp corners,
overlay={\fill[cdeg] (frame.north west) rectangle ([yshift=-1.25cm]frame.north east);
\node[anchor=west,text=white,font=\sffamily\bfseries\LARGE] at ([xshift=5mm,yshift=-0.62cm]frame.north west) {Question 24 D};
\node[anchor=east,text=white,font=\sffamily\itshape\large] at ([xshift=-5mm,yshift=-0.62cm]frame.north east) {Second degré};}]
\vspace*{1.75cm}
\begin{minipage}[t][6.9cm][c]{\linewidth}\LARGE\centering
\begin{minipage}{0.95\linewidth}\raggedright\linespread{1.25}\selectfont Quelle est la forme canonique de $x^2-2x-8$~?\end{minipage}
\end{minipage}

\vfill
\begin{tcolorbox}[enhanced,nobeforeafter,width=7.1cm,height=2.1cm,valign=center,colback=cdeg!6,colframe=cdeg,arc=2.5mm,boxrule=0.9pt,left=3mm]{\Large\bfseries\color{cdeg}a)}\quad\begin{minipage}[c]{5.3cm}\raggedright\Large $(x-2)^2-9$\end{minipage}\end{tcolorbox}\hfill\begin{tcolorbox}[enhanced,nobeforeafter,width=7.1cm,height=2.1cm,valign=center,colback=cdeg!6,colframe=cdeg,arc=2.5mm,boxrule=0.9pt,left=3mm]{\Large\bfseries\color{cdeg}b)}\quad\begin{minipage}[c]{5.3cm}\raggedright\Large $(x+1)^2-9$\end{minipage}\end{tcolorbox}\par\vspace{2.5mm}
\begin{tcolorbox}[enhanced,nobeforeafter,width=7.1cm,height=2.1cm,valign=center,colback=cdeg!6,colframe=cdeg,arc=2.5mm,boxrule=0.9pt,left=3mm]{\Large\bfseries\color{cdeg}c)}\quad\begin{minipage}[c]{5.3cm}\raggedright\Large $(x-1)^2-8$\end{minipage}\end{tcolorbox}\hfill\begin{tcolorbox}[enhanced,nobeforeafter,width=7.1cm,height=2.1cm,valign=center,colback=cdeg!6,colframe=cdeg,arc=2.5mm,boxrule=0.9pt,left=3mm]{\Large\bfseries\color{cdeg}d)}\quad\begin{minipage}[c]{5.3cm}\raggedright\Large $(x-1)^2-9$\end{minipage}\end{tcolorbox}
\end{tcolorbox}
\newpage
\begin{tcolorbox}[enhanced,width=16cm,height=14cm,colback=white,colframe=cdeg,boxrule=2pt,arc=0mm,
left=6mm,right=6mm,top=0mm,bottom=4mm,sharp corners,
overlay={\fill[cdeg] (frame.north west) rectangle ([yshift=-1.25cm]frame.north east);
\node[anchor=west,text=white,font=\sffamily\bfseries\LARGE] at ([xshift=5mm,yshift=-0.62cm]frame.north west) {Question 25 D};
\node[anchor=east,text=white,font=\sffamily\itshape\large] at ([xshift=-5mm,yshift=-0.62cm]frame.north east) {Second degré};}]
\vspace*{1.75cm}
\begin{minipage}[t][6.9cm][c]{\linewidth}\LARGE\centering
\begin{minipage}{0.95\linewidth}\raggedright\linespread{1.25}\selectfont Quelles sont les coordonnées du sommet de la parabole d'équation $y=x^2-6x+10$~?\end{minipage}
\end{minipage}

\vfill
\begin{tcolorbox}[enhanced,nobeforeafter,width=7.1cm,height=2.1cm,valign=center,colback=cdeg!6,colframe=cdeg,arc=2.5mm,boxrule=0.9pt,left=3mm]{\Large\bfseries\color{cdeg}a)}\quad\begin{minipage}[c]{5.3cm}\raggedright\Large $(-3\,;\,1)$\end{minipage}\end{tcolorbox}\hfill\begin{tcolorbox}[enhanced,nobeforeafter,width=7.1cm,height=2.1cm,valign=center,colback=cdeg!6,colframe=cdeg,arc=2.5mm,boxrule=0.9pt,left=3mm]{\Large\bfseries\color{cdeg}b)}\quad\begin{minipage}[c]{5.3cm}\raggedright\Large $(3\,;\,1)$\end{minipage}\end{tcolorbox}\par\vspace{2.5mm}
\begin{tcolorbox}[enhanced,nobeforeafter,width=7.1cm,height=2.1cm,valign=center,colback=cdeg!6,colframe=cdeg,arc=2.5mm,boxrule=0.9pt,left=3mm]{\Large\bfseries\color{cdeg}c)}\quad\begin{minipage}[c]{5.3cm}\raggedright\Large $(3\,;\,10)$\end{minipage}\end{tcolorbox}\hfill\begin{tcolorbox}[enhanced,nobeforeafter,width=7.1cm,height=2.1cm,valign=center,colback=cdeg!6,colframe=cdeg,arc=2.5mm,boxrule=0.9pt,left=3mm]{\Large\bfseries\color{cdeg}d)}\quad\begin{minipage}[c]{5.3cm}\raggedright\Large $(3\,;\,-1)$\end{minipage}\end{tcolorbox}
\end{tcolorbox}
\newpage
\begin{tcolorbox}[enhanced,width=16cm,height=14cm,colback=white,colframe=cdeg,boxrule=2pt,arc=0mm,
left=6mm,right=6mm,top=0mm,bottom=4mm,sharp corners,
overlay={\fill[cdeg] (frame.north west) rectangle ([yshift=-1.25cm]frame.north east);
\node[anchor=west,text=white,font=\sffamily\bfseries\LARGE] at ([xshift=5mm,yshift=-0.62cm]frame.north west) {Question 26 D};
\node[anchor=east,text=white,font=\sffamily\itshape\large] at ([xshift=-5mm,yshift=-0.62cm]frame.north east) {Second degré};}]
\vspace*{1.75cm}
\begin{minipage}[t][6.9cm][c]{\linewidth}\LARGE\centering
\begin{minipage}{0.95\linewidth}\raggedright\linespread{1.25}\selectfont Quelles sont les coordonnées du sommet de la parabole d'équation $y=4x^2+8x-1$~?\end{minipage}
\end{minipage}

\vfill
\begin{tcolorbox}[enhanced,nobeforeafter,width=7.1cm,height=2.1cm,valign=center,colback=cdeg!6,colframe=cdeg,arc=2.5mm,boxrule=0.9pt,left=3mm]{\Large\bfseries\color{cdeg}a)}\quad\begin{minipage}[c]{5.3cm}\raggedright\Large $(-1\,;\,-5)$\end{minipage}\end{tcolorbox}\hfill\begin{tcolorbox}[enhanced,nobeforeafter,width=7.1cm,height=2.1cm,valign=center,colback=cdeg!6,colframe=cdeg,arc=2.5mm,boxrule=0.9pt,left=3mm]{\Large\bfseries\color{cdeg}b)}\quad\begin{minipage}[c]{5.3cm}\raggedright\Large $(-1\,;\,5)$\end{minipage}\end{tcolorbox}\par\vspace{2.5mm}
\begin{tcolorbox}[enhanced,nobeforeafter,width=7.1cm,height=2.1cm,valign=center,colback=cdeg!6,colframe=cdeg,arc=2.5mm,boxrule=0.9pt,left=3mm]{\Large\bfseries\color{cdeg}c)}\quad\begin{minipage}[c]{5.3cm}\raggedright\Large $(1\,;\,-5)$\end{minipage}\end{tcolorbox}\hfill\begin{tcolorbox}[enhanced,nobeforeafter,width=7.1cm,height=2.1cm,valign=center,colback=cdeg!6,colframe=cdeg,arc=2.5mm,boxrule=0.9pt,left=3mm]{\Large\bfseries\color{cdeg}d)}\quad\begin{minipage}[c]{5.3cm}\raggedright\Large $(-1\,;\,-1)$\end{minipage}\end{tcolorbox}
\end{tcolorbox}
\newpage
\begin{tcolorbox}[enhanced,width=16cm,height=14cm,colback=white,colframe=cdeg,boxrule=2pt,arc=0mm,
left=6mm,right=6mm,top=0mm,bottom=4mm,sharp corners,
overlay={\fill[cdeg] (frame.north west) rectangle ([yshift=-1.25cm]frame.north east);
\node[anchor=west,text=white,font=\sffamily\bfseries\LARGE] at ([xshift=5mm,yshift=-0.62cm]frame.north west) {Question 27 D};
\node[anchor=east,text=white,font=\sffamily\itshape\large] at ([xshift=-5mm,yshift=-0.62cm]frame.north east) {Second degré};}]
\vspace*{1.75cm}
\begin{minipage}[t][6.9cm][c]{\linewidth}\LARGE\centering
\begin{minipage}{0.95\linewidth}\raggedright\linespread{1.25}\selectfont Quel est l'ensemble des solutions de l'inéquation $x^2-8x+12<0$~?\end{minipage}
\end{minipage}

\vfill
\begin{tcolorbox}[enhanced,nobeforeafter,width=7.1cm,height=2.1cm,valign=center,colback=cdeg!6,colframe=cdeg,arc=2.5mm,boxrule=0.9pt,left=3mm]{\Large\bfseries\color{cdeg}a)}\quad\begin{minipage}[c]{5.3cm}\raggedright\large $]-\infty\,;\,2]\cup[6\,;\,+\infty[$\end{minipage}\end{tcolorbox}\hfill\begin{tcolorbox}[enhanced,nobeforeafter,width=7.1cm,height=2.1cm,valign=center,colback=cdeg!6,colframe=cdeg,arc=2.5mm,boxrule=0.9pt,left=3mm]{\Large\bfseries\color{cdeg}b)}\quad\begin{minipage}[c]{5.3cm}\raggedright\Large $]2\,;\,6[$\end{minipage}\end{tcolorbox}\par\vspace{2.5mm}
\begin{tcolorbox}[enhanced,nobeforeafter,width=7.1cm,height=2.1cm,valign=center,colback=cdeg!6,colframe=cdeg,arc=2.5mm,boxrule=0.9pt,left=3mm]{\Large\bfseries\color{cdeg}c)}\quad\begin{minipage}[c]{5.3cm}\raggedright\large $]-\infty\,;\,2[\cup]6\,;\,+\infty[$\end{minipage}\end{tcolorbox}\hfill\begin{tcolorbox}[enhanced,nobeforeafter,width=7.1cm,height=2.1cm,valign=center,colback=cdeg!6,colframe=cdeg,arc=2.5mm,boxrule=0.9pt,left=3mm]{\Large\bfseries\color{cdeg}d)}\quad\begin{minipage}[c]{5.3cm}\raggedright\Large $[2\,;\,6]$\end{minipage}\end{tcolorbox}
\end{tcolorbox}
\newpage
\begin{tcolorbox}[enhanced,width=16cm,height=14cm,colback=white,colframe=cdeg,boxrule=2pt,arc=0mm,
left=6mm,right=6mm,top=0mm,bottom=4mm,sharp corners,
overlay={\fill[cdeg] (frame.north west) rectangle ([yshift=-1.25cm]frame.north east);
\node[anchor=west,text=white,font=\sffamily\bfseries\LARGE] at ([xshift=5mm,yshift=-0.62cm]frame.north west) {Question 28 D};
\node[anchor=east,text=white,font=\sffamily\itshape\large] at ([xshift=-5mm,yshift=-0.62cm]frame.north east) {Second degré};}]
\vspace*{1.75cm}
\begin{minipage}[t][6.9cm][c]{\linewidth}\LARGE\centering
\begin{minipage}{0.95\linewidth}\raggedright\linespread{1.25}\selectfont Quel est l'ensemble des solutions de l'inéquation $-x^2+16\geqslant 0$~?\end{minipage}
\end{minipage}

\vfill
\begin{tcolorbox}[enhanced,nobeforeafter,width=7.1cm,height=2.1cm,valign=center,colback=cdeg!6,colframe=cdeg,arc=2.5mm,boxrule=0.9pt,left=3mm]{\Large\bfseries\color{cdeg}a)}\quad\begin{minipage}[c]{5.3cm}\raggedright\Large $]-4\,;\,4[$\end{minipage}\end{tcolorbox}\hfill\begin{tcolorbox}[enhanced,nobeforeafter,width=7.1cm,height=2.1cm,valign=center,colback=cdeg!6,colframe=cdeg,arc=2.5mm,boxrule=0.9pt,left=3mm]{\Large\bfseries\color{cdeg}b)}\quad\begin{minipage}[c]{5.3cm}\raggedright\large $]-\infty\,;\,-4]\cup[4\,;\,+\infty[$\end{minipage}\end{tcolorbox}\par\vspace{2.5mm}
\begin{tcolorbox}[enhanced,nobeforeafter,width=7.1cm,height=2.1cm,valign=center,colback=cdeg!6,colframe=cdeg,arc=2.5mm,boxrule=0.9pt,left=3mm]{\Large\bfseries\color{cdeg}c)}\quad\begin{minipage}[c]{5.3cm}\raggedright\large $]-\infty\,;\,-4[\cup]4\,;\,+\infty[$\end{minipage}\end{tcolorbox}\hfill\begin{tcolorbox}[enhanced,nobeforeafter,width=7.1cm,height=2.1cm,valign=center,colback=cdeg!6,colframe=cdeg,arc=2.5mm,boxrule=0.9pt,left=3mm]{\Large\bfseries\color{cdeg}d)}\quad\begin{minipage}[c]{5.3cm}\raggedright\Large $[-4\,;\,4]$\end{minipage}\end{tcolorbox}
\end{tcolorbox}
\newpage
\begin{tcolorbox}[enhanced,width=16cm,height=14cm,colback=white,colframe=cdeg,boxrule=2pt,arc=0mm,
left=6mm,right=6mm,top=0mm,bottom=4mm,sharp corners,
overlay={\fill[cdeg] (frame.north west) rectangle ([yshift=-1.25cm]frame.north east);
\node[anchor=west,text=white,font=\sffamily\bfseries\LARGE] at ([xshift=5mm,yshift=-0.62cm]frame.north west) {Question 29 D};
\node[anchor=east,text=white,font=\sffamily\itshape\large] at ([xshift=-5mm,yshift=-0.62cm]frame.north east) {Second degré};}]
\vspace*{1.75cm}
\begin{minipage}[t][6.9cm][c]{\linewidth}\LARGE\centering
\begin{minipage}{0.95\linewidth}\raggedright\linespread{1.25}\selectfont Le trinôme $x^2+10x+9$ a deux racines. Que vaut leur produit~?\end{minipage}
\end{minipage}

\vfill
\begin{tcolorbox}[enhanced,nobeforeafter,width=7.1cm,height=2.1cm,valign=center,colback=cdeg!6,colframe=cdeg,arc=2.5mm,boxrule=0.9pt,left=3mm]{\Large\bfseries\color{cdeg}a)}\quad\begin{minipage}[c]{5.3cm}\raggedright\Large $-9$\end{minipage}\end{tcolorbox}\hfill\begin{tcolorbox}[enhanced,nobeforeafter,width=7.1cm,height=2.1cm,valign=center,colback=cdeg!6,colframe=cdeg,arc=2.5mm,boxrule=0.9pt,left=3mm]{\Large\bfseries\color{cdeg}b)}\quad\begin{minipage}[c]{5.3cm}\raggedright\Large $10$\end{minipage}\end{tcolorbox}\par\vspace{2.5mm}
\begin{tcolorbox}[enhanced,nobeforeafter,width=7.1cm,height=2.1cm,valign=center,colback=cdeg!6,colframe=cdeg,arc=2.5mm,boxrule=0.9pt,left=3mm]{\Large\bfseries\color{cdeg}c)}\quad\begin{minipage}[c]{5.3cm}\raggedright\Large $9$\end{minipage}\end{tcolorbox}\hfill\begin{tcolorbox}[enhanced,nobeforeafter,width=7.1cm,height=2.1cm,valign=center,colback=cdeg!6,colframe=cdeg,arc=2.5mm,boxrule=0.9pt,left=3mm]{\Large\bfseries\color{cdeg}d)}\quad\begin{minipage}[c]{5.3cm}\raggedright\Large $-10$\end{minipage}\end{tcolorbox}
\end{tcolorbox}
\newpage
\begin{tcolorbox}[enhanced,width=16cm,height=14cm,colback=white,colframe=cdeg,boxrule=2pt,arc=0mm,
left=6mm,right=6mm,top=0mm,bottom=4mm,sharp corners,
overlay={\fill[cdeg] (frame.north west) rectangle ([yshift=-1.25cm]frame.north east);
\node[anchor=west,text=white,font=\sffamily\bfseries\LARGE] at ([xshift=5mm,yshift=-0.62cm]frame.north west) {Question 30 D};
\node[anchor=east,text=white,font=\sffamily\itshape\large] at ([xshift=-5mm,yshift=-0.62cm]frame.north east) {Second degré};}]
\vspace*{1.75cm}
\begin{minipage}[t][6.9cm][c]{\linewidth}\LARGE\centering
\begin{minipage}{0.95\linewidth}\raggedright\linespread{1.25}\selectfont Le trinôme $5x^2-20x+1$ a deux racines. Que vaut leur somme~?\end{minipage}
\end{minipage}

\vfill
\begin{tcolorbox}[enhanced,nobeforeafter,width=7.1cm,height=2.1cm,valign=center,colback=cdeg!6,colframe=cdeg,arc=2.5mm,boxrule=0.9pt,left=3mm]{\Large\bfseries\color{cdeg}a)}\quad\begin{minipage}[c]{5.3cm}\raggedright\Large $\dfrac{1}{5}$\end{minipage}\end{tcolorbox}\hfill\begin{tcolorbox}[enhanced,nobeforeafter,width=7.1cm,height=2.1cm,valign=center,colback=cdeg!6,colframe=cdeg,arc=2.5mm,boxrule=0.9pt,left=3mm]{\Large\bfseries\color{cdeg}b)}\quad\begin{minipage}[c]{5.3cm}\raggedright\Large $20$\end{minipage}\end{tcolorbox}\par\vspace{2.5mm}
\begin{tcolorbox}[enhanced,nobeforeafter,width=7.1cm,height=2.1cm,valign=center,colback=cdeg!6,colframe=cdeg,arc=2.5mm,boxrule=0.9pt,left=3mm]{\Large\bfseries\color{cdeg}c)}\quad\begin{minipage}[c]{5.3cm}\raggedright\Large $4$\end{minipage}\end{tcolorbox}\hfill\begin{tcolorbox}[enhanced,nobeforeafter,width=7.1cm,height=2.1cm,valign=center,colback=cdeg!6,colframe=cdeg,arc=2.5mm,boxrule=0.9pt,left=3mm]{\Large\bfseries\color{cdeg}d)}\quad\begin{minipage}[c]{5.3cm}\raggedright\Large $-4$\end{minipage}\end{tcolorbox}
\end{tcolorbox}
\newpage
\begin{tcolorbox}[enhanced,width=16cm,height=14cm,colback=white,colframe=cdeg,boxrule=2pt,arc=0mm,
left=6mm,right=6mm,top=0mm,bottom=4mm,sharp corners,
overlay={\fill[cdeg] (frame.north west) rectangle ([yshift=-1.25cm]frame.north east);
\node[anchor=west,text=white,font=\sffamily\bfseries\LARGE] at ([xshift=5mm,yshift=-0.62cm]frame.north west) {Question 31 D};
\node[anchor=east,text=white,font=\sffamily\itshape\large] at ([xshift=-5mm,yshift=-0.62cm]frame.north east) {Second degré};}]
\vspace*{1.75cm}
\begin{minipage}[t][6.9cm][c]{\linewidth}\LARGE\centering
\begin{minipage}{0.95\linewidth}\raggedright\linespread{1.25}\selectfont Le trinôme $-2x^2-6x+8$ a pour racines $-4$ et $1$. Quelle est sa forme factorisée~?\end{minipage}
\end{minipage}

\vfill
\begin{tcolorbox}[enhanced,nobeforeafter,width=7.1cm,height=2.1cm,valign=center,colback=cdeg!6,colframe=cdeg,arc=2.5mm,boxrule=0.9pt,left=3mm]{\Large\bfseries\color{cdeg}a)}\quad\begin{minipage}[c]{5.3cm}\raggedright\Large $-2(x-4)(x-1)$\end{minipage}\end{tcolorbox}\hfill\begin{tcolorbox}[enhanced,nobeforeafter,width=7.1cm,height=2.1cm,valign=center,colback=cdeg!6,colframe=cdeg,arc=2.5mm,boxrule=0.9pt,left=3mm]{\Large\bfseries\color{cdeg}b)}\quad\begin{minipage}[c]{5.3cm}\raggedright\Large $(x+4)(x-1)$\end{minipage}\end{tcolorbox}\par\vspace{2.5mm}
\begin{tcolorbox}[enhanced,nobeforeafter,width=7.1cm,height=2.1cm,valign=center,colback=cdeg!6,colframe=cdeg,arc=2.5mm,boxrule=0.9pt,left=3mm]{\Large\bfseries\color{cdeg}c)}\quad\begin{minipage}[c]{5.3cm}\raggedright\Large $-2(x+4)(x-1)$\end{minipage}\end{tcolorbox}\hfill\begin{tcolorbox}[enhanced,nobeforeafter,width=7.1cm,height=2.1cm,valign=center,colback=cdeg!6,colframe=cdeg,arc=2.5mm,boxrule=0.9pt,left=3mm]{\Large\bfseries\color{cdeg}d)}\quad\begin{minipage}[c]{5.3cm}\raggedright\Large $-2(x-4)(x+1)$\end{minipage}\end{tcolorbox}
\end{tcolorbox}
\newpage
\begin{tcolorbox}[enhanced,width=16cm,height=14cm,colback=white,colframe=cdeg,boxrule=2pt,arc=0mm,
left=6mm,right=6mm,top=0mm,bottom=4mm,sharp corners,
overlay={\fill[cdeg] (frame.north west) rectangle ([yshift=-1.25cm]frame.north east);
\node[anchor=west,text=white,font=\sffamily\bfseries\LARGE] at ([xshift=5mm,yshift=-0.62cm]frame.north west) {Question 32 D};
\node[anchor=east,text=white,font=\sffamily\itshape\large] at ([xshift=-5mm,yshift=-0.62cm]frame.north east) {Second degré};}]
\vspace*{1.75cm}
\begin{minipage}[t][6.9cm][c]{\linewidth}\LARGE\centering
\begin{minipage}{0.95\linewidth}\raggedright\linespread{1.25}\selectfont Le trinôme $x^2-49$ a pour racines $-7$ et $7$. Quelle est sa forme factorisée~?\end{minipage}
\end{minipage}

\vfill
\begin{tcolorbox}[enhanced,nobeforeafter,width=7.1cm,height=2.1cm,valign=center,colback=cdeg!6,colframe=cdeg,arc=2.5mm,boxrule=0.9pt,left=3mm]{\Large\bfseries\color{cdeg}a)}\quad\begin{minipage}[c]{5.3cm}\raggedright\Large $(x-7)(x+7)$\end{minipage}\end{tcolorbox}\hfill\begin{tcolorbox}[enhanced,nobeforeafter,width=7.1cm,height=2.1cm,valign=center,colback=cdeg!6,colframe=cdeg,arc=2.5mm,boxrule=0.9pt,left=3mm]{\Large\bfseries\color{cdeg}b)}\quad\begin{minipage}[c]{5.3cm}\raggedright\Large $-(x+7)(x-7)$\end{minipage}\end{tcolorbox}\par\vspace{2.5mm}
\begin{tcolorbox}[enhanced,nobeforeafter,width=7.1cm,height=2.1cm,valign=center,colback=cdeg!6,colframe=cdeg,arc=2.5mm,boxrule=0.9pt,left=3mm]{\Large\bfseries\color{cdeg}c)}\quad\begin{minipage}[c]{5.3cm}\raggedright\Large $2(x+7)(x-7)$\end{minipage}\end{tcolorbox}\hfill\begin{tcolorbox}[enhanced,nobeforeafter,width=7.1cm,height=2.1cm,valign=center,colback=cdeg!6,colframe=cdeg,arc=2.5mm,boxrule=0.9pt,left=3mm]{\Large\bfseries\color{cdeg}d)}\quad\begin{minipage}[c]{5.3cm}\raggedright\Large $(x+7)(x-7)$\end{minipage}\end{tcolorbox}
\end{tcolorbox}
\newpage
\begin{tcolorbox}[enhanced,width=16cm,height=14cm,colback=white,colframe=cder,boxrule=2pt,arc=0mm,
left=6mm,right=6mm,top=0mm,bottom=4mm,sharp corners,
overlay={\fill[cder] (frame.north west) rectangle ([yshift=-1.25cm]frame.north east);
\node[anchor=west,text=white,font=\sffamily\bfseries\LARGE] at ([xshift=5mm,yshift=-0.62cm]frame.north west) {Question 33 D};
\node[anchor=east,text=white,font=\sffamily\itshape\large] at ([xshift=-5mm,yshift=-0.62cm]frame.north east) {Dérivation};}]
\vspace*{1.75cm}
\begin{minipage}[t][6.9cm][c]{\linewidth}\LARGE\centering
\begin{minipage}{0.95\linewidth}\raggedright\linespread{1.25}\selectfont Soit $f(x)=-2x^3+x^2+5x-3$. Que vaut $f'(2)$~?\end{minipage}
\end{minipage}

\vfill
\begin{tcolorbox}[enhanced,nobeforeafter,width=7.1cm,height=2.1cm,valign=center,colback=cder!6,colframe=cder,arc=2.5mm,boxrule=0.9pt,left=3mm]{\Large\bfseries\color{cder}a)}\quad\begin{minipage}[c]{5.3cm}\raggedright\Large $-1$\end{minipage}\end{tcolorbox}\hfill\begin{tcolorbox}[enhanced,nobeforeafter,width=7.1cm,height=2.1cm,valign=center,colback=cder!6,colframe=cder,arc=2.5mm,boxrule=0.9pt,left=3mm]{\Large\bfseries\color{cder}b)}\quad\begin{minipage}[c]{5.3cm}\raggedright\Large $-5$\end{minipage}\end{tcolorbox}\par\vspace{2.5mm}
\begin{tcolorbox}[enhanced,nobeforeafter,width=7.1cm,height=2.1cm,valign=center,colback=cder!6,colframe=cder,arc=2.5mm,boxrule=0.9pt,left=3mm]{\Large\bfseries\color{cder}c)}\quad\begin{minipage}[c]{5.3cm}\raggedright\Large $-17$\end{minipage}\end{tcolorbox}\hfill\begin{tcolorbox}[enhanced,nobeforeafter,width=7.1cm,height=2.1cm,valign=center,colback=cder!6,colframe=cder,arc=2.5mm,boxrule=0.9pt,left=3mm]{\Large\bfseries\color{cder}d)}\quad\begin{minipage}[c]{5.3cm}\raggedright\Large $-15$\end{minipage}\end{tcolorbox}
\end{tcolorbox}
\newpage
\begin{tcolorbox}[enhanced,width=16cm,height=14cm,colback=white,colframe=cder,boxrule=2pt,arc=0mm,
left=6mm,right=6mm,top=0mm,bottom=4mm,sharp corners,
overlay={\fill[cder] (frame.north west) rectangle ([yshift=-1.25cm]frame.north east);
\node[anchor=west,text=white,font=\sffamily\bfseries\LARGE] at ([xshift=5mm,yshift=-0.62cm]frame.north west) {Question 34 D};
\node[anchor=east,text=white,font=\sffamily\itshape\large] at ([xshift=-5mm,yshift=-0.62cm]frame.north east) {Dérivation};}]
\vspace*{1.75cm}
\begin{minipage}[t][6.9cm][c]{\linewidth}\LARGE\centering
\begin{minipage}{0.95\linewidth}\raggedright\linespread{1.25}\selectfont Soit $f(x)=x^3-6x^2+9x$. Que vaut $f'(-1)$~?\end{minipage}
\end{minipage}

\vfill
\begin{tcolorbox}[enhanced,nobeforeafter,width=7.1cm,height=2.1cm,valign=center,colback=cder!6,colframe=cder,arc=2.5mm,boxrule=0.9pt,left=3mm]{\Large\bfseries\color{cder}a)}\quad\begin{minipage}[c]{5.3cm}\raggedright\Large $16$\end{minipage}\end{tcolorbox}\hfill\begin{tcolorbox}[enhanced,nobeforeafter,width=7.1cm,height=2.1cm,valign=center,colback=cder!6,colframe=cder,arc=2.5mm,boxrule=0.9pt,left=3mm]{\Large\bfseries\color{cder}b)}\quad\begin{minipage}[c]{5.3cm}\raggedright\Large $18$\end{minipage}\end{tcolorbox}\par\vspace{2.5mm}
\begin{tcolorbox}[enhanced,nobeforeafter,width=7.1cm,height=2.1cm,valign=center,colback=cder!6,colframe=cder,arc=2.5mm,boxrule=0.9pt,left=3mm]{\Large\bfseries\color{cder}c)}\quad\begin{minipage}[c]{5.3cm}\raggedright\Large $24$\end{minipage}\end{tcolorbox}\hfill\begin{tcolorbox}[enhanced,nobeforeafter,width=7.1cm,height=2.1cm,valign=center,colback=cder!6,colframe=cder,arc=2.5mm,boxrule=0.9pt,left=3mm]{\Large\bfseries\color{cder}d)}\quad\begin{minipage}[c]{5.3cm}\raggedright\Large $-16$\end{minipage}\end{tcolorbox}
\end{tcolorbox}
\newpage
\begin{tcolorbox}[enhanced,width=16cm,height=14cm,colback=white,colframe=cder,boxrule=2pt,arc=0mm,
left=6mm,right=6mm,top=0mm,bottom=4mm,sharp corners,
overlay={\fill[cder] (frame.north west) rectangle ([yshift=-1.25cm]frame.north east);
\node[anchor=west,text=white,font=\sffamily\bfseries\LARGE] at ([xshift=5mm,yshift=-0.62cm]frame.north west) {Question 35 D};
\node[anchor=east,text=white,font=\sffamily\itshape\large] at ([xshift=-5mm,yshift=-0.62cm]frame.north east) {Dérivation};}]
\vspace*{1.75cm}
\begin{minipage}[t][6.9cm][c]{\linewidth}\LARGE\centering
\begin{minipage}{0.95\linewidth}\raggedright\linespread{1.25}\selectfont Quelle est l'équation de la tangente à la courbe de $f(x)=-2x^2+5$ au point d'abscisse $1$~?\end{minipage}
\end{minipage}

\vfill
\begin{tcolorbox}[enhanced,nobeforeafter,width=7.1cm,height=2.1cm,valign=center,colback=cder!6,colframe=cder,arc=2.5mm,boxrule=0.9pt,left=3mm]{\Large\bfseries\color{cder}a)}\quad\begin{minipage}[c]{5.3cm}\raggedright\Large $y=-4x+3$\end{minipage}\end{tcolorbox}\hfill\begin{tcolorbox}[enhanced,nobeforeafter,width=7.1cm,height=2.1cm,valign=center,colback=cder!6,colframe=cder,arc=2.5mm,boxrule=0.9pt,left=3mm]{\Large\bfseries\color{cder}b)}\quad\begin{minipage}[c]{5.3cm}\raggedright\Large $y=-4x-1$\end{minipage}\end{tcolorbox}\par\vspace{2.5mm}
\begin{tcolorbox}[enhanced,nobeforeafter,width=7.1cm,height=2.1cm,valign=center,colback=cder!6,colframe=cder,arc=2.5mm,boxrule=0.9pt,left=3mm]{\Large\bfseries\color{cder}c)}\quad\begin{minipage}[c]{5.3cm}\raggedright\Large $y=3x-7$\end{minipage}\end{tcolorbox}\hfill\begin{tcolorbox}[enhanced,nobeforeafter,width=7.1cm,height=2.1cm,valign=center,colback=cder!6,colframe=cder,arc=2.5mm,boxrule=0.9pt,left=3mm]{\Large\bfseries\color{cder}d)}\quad\begin{minipage}[c]{5.3cm}\raggedright\Large $y=-4x+7$\end{minipage}\end{tcolorbox}
\end{tcolorbox}
\newpage
\begin{tcolorbox}[enhanced,width=16cm,height=14cm,colback=white,colframe=cder,boxrule=2pt,arc=0mm,
left=6mm,right=6mm,top=0mm,bottom=4mm,sharp corners,
overlay={\fill[cder] (frame.north west) rectangle ([yshift=-1.25cm]frame.north east);
\node[anchor=west,text=white,font=\sffamily\bfseries\LARGE] at ([xshift=5mm,yshift=-0.62cm]frame.north west) {Question 36 D};
\node[anchor=east,text=white,font=\sffamily\itshape\large] at ([xshift=-5mm,yshift=-0.62cm]frame.north east) {Dérivation};}]
\vspace*{1.75cm}
\begin{minipage}[t][6.9cm][c]{\linewidth}\LARGE\centering
\begin{minipage}{0.95\linewidth}\raggedright\linespread{1.25}\selectfont Quelle est l'équation de la tangente à la courbe de $f(x)=x^2-5x+6$ au point d'abscisse $4$~?\end{minipage}
\end{minipage}

\vfill
\begin{tcolorbox}[enhanced,nobeforeafter,width=7.1cm,height=2.1cm,valign=center,colback=cder!6,colframe=cder,arc=2.5mm,boxrule=0.9pt,left=3mm]{\Large\bfseries\color{cder}a)}\quad\begin{minipage}[c]{5.3cm}\raggedright\Large $y=3x-10$\end{minipage}\end{tcolorbox}\hfill\begin{tcolorbox}[enhanced,nobeforeafter,width=7.1cm,height=2.1cm,valign=center,colback=cder!6,colframe=cder,arc=2.5mm,boxrule=0.9pt,left=3mm]{\Large\bfseries\color{cder}b)}\quad\begin{minipage}[c]{5.3cm}\raggedright\Large $y=3x+2$\end{minipage}\end{tcolorbox}\par\vspace{2.5mm}
\begin{tcolorbox}[enhanced,nobeforeafter,width=7.1cm,height=2.1cm,valign=center,colback=cder!6,colframe=cder,arc=2.5mm,boxrule=0.9pt,left=3mm]{\Large\bfseries\color{cder}c)}\quad\begin{minipage}[c]{5.3cm}\raggedright\Large $y=2x-5$\end{minipage}\end{tcolorbox}\hfill\begin{tcolorbox}[enhanced,nobeforeafter,width=7.1cm,height=2.1cm,valign=center,colback=cder!6,colframe=cder,arc=2.5mm,boxrule=0.9pt,left=3mm]{\Large\bfseries\color{cder}d)}\quad\begin{minipage}[c]{5.3cm}\raggedright\Large $y=3x+14$\end{minipage}\end{tcolorbox}
\end{tcolorbox}
\newpage
\begin{tcolorbox}[enhanced,width=16cm,height=14cm,colback=white,colframe=cder,boxrule=2pt,arc=0mm,
left=6mm,right=6mm,top=0mm,bottom=4mm,sharp corners,
overlay={\fill[cder] (frame.north west) rectangle ([yshift=-1.25cm]frame.north east);
\node[anchor=west,text=white,font=\sffamily\bfseries\LARGE] at ([xshift=5mm,yshift=-0.62cm]frame.north west) {Question 37 D};
\node[anchor=east,text=white,font=\sffamily\itshape\large] at ([xshift=-5mm,yshift=-0.62cm]frame.north east) {Dérivation};}]
\vspace*{1.75cm}
\begin{minipage}[t][6.9cm][c]{\linewidth}\LARGE\centering
\begin{minipage}{0.95\linewidth}\raggedright\linespread{1.25}\selectfont Soit $f(x)=(5x+1)(2x)$. Que vaut $f'(1)$~?\end{minipage}
\end{minipage}

\vfill
\begin{tcolorbox}[enhanced,nobeforeafter,width=7.1cm,height=2.1cm,valign=center,colback=cder!6,colframe=cder,arc=2.5mm,boxrule=0.9pt,left=3mm]{\Large\bfseries\color{cder}a)}\quad\begin{minipage}[c]{5.3cm}\raggedright\Large $24$\end{minipage}\end{tcolorbox}\hfill\begin{tcolorbox}[enhanced,nobeforeafter,width=7.1cm,height=2.1cm,valign=center,colback=cder!6,colframe=cder,arc=2.5mm,boxrule=0.9pt,left=3mm]{\Large\bfseries\color{cder}b)}\quad\begin{minipage}[c]{5.3cm}\raggedright\Large $12$\end{minipage}\end{tcolorbox}\par\vspace{2.5mm}
\begin{tcolorbox}[enhanced,nobeforeafter,width=7.1cm,height=2.1cm,valign=center,colback=cder!6,colframe=cder,arc=2.5mm,boxrule=0.9pt,left=3mm]{\Large\bfseries\color{cder}c)}\quad\begin{minipage}[c]{5.3cm}\raggedright\Large $22$\end{minipage}\end{tcolorbox}\hfill\begin{tcolorbox}[enhanced,nobeforeafter,width=7.1cm,height=2.1cm,valign=center,colback=cder!6,colframe=cder,arc=2.5mm,boxrule=0.9pt,left=3mm]{\Large\bfseries\color{cder}d)}\quad\begin{minipage}[c]{5.3cm}\raggedright\Large $10$\end{minipage}\end{tcolorbox}
\end{tcolorbox}
\newpage
\begin{tcolorbox}[enhanced,width=16cm,height=14cm,colback=white,colframe=cder,boxrule=2pt,arc=0mm,
left=6mm,right=6mm,top=0mm,bottom=4mm,sharp corners,
overlay={\fill[cder] (frame.north west) rectangle ([yshift=-1.25cm]frame.north east);
\node[anchor=west,text=white,font=\sffamily\bfseries\LARGE] at ([xshift=5mm,yshift=-0.62cm]frame.north west) {Question 38 D};
\node[anchor=east,text=white,font=\sffamily\itshape\large] at ([xshift=-5mm,yshift=-0.62cm]frame.north east) {Dérivation};}]
\vspace*{1.75cm}
\begin{minipage}[t][6.9cm][c]{\linewidth}\LARGE\centering
\begin{minipage}{0.95\linewidth}\raggedright\linespread{1.25}\selectfont Soit $f(x)=(-2x+3)(3x+1)$. Que vaut $f'(0)$~?\end{minipage}
\end{minipage}

\vfill
\begin{tcolorbox}[enhanced,nobeforeafter,width=7.1cm,height=2.1cm,valign=center,colback=cder!6,colframe=cder,arc=2.5mm,boxrule=0.9pt,left=3mm]{\Large\bfseries\color{cder}a)}\quad\begin{minipage}[c]{5.3cm}\raggedright\Large $-6$\end{minipage}\end{tcolorbox}\hfill\begin{tcolorbox}[enhanced,nobeforeafter,width=7.1cm,height=2.1cm,valign=center,colback=cder!6,colframe=cder,arc=2.5mm,boxrule=0.9pt,left=3mm]{\Large\bfseries\color{cder}b)}\quad\begin{minipage}[c]{5.3cm}\raggedright\Large $7$\end{minipage}\end{tcolorbox}\par\vspace{2.5mm}
\begin{tcolorbox}[enhanced,nobeforeafter,width=7.1cm,height=2.1cm,valign=center,colback=cder!6,colframe=cder,arc=2.5mm,boxrule=0.9pt,left=3mm]{\Large\bfseries\color{cder}c)}\quad\begin{minipage}[c]{5.3cm}\raggedright\Large $3$\end{minipage}\end{tcolorbox}\hfill\begin{tcolorbox}[enhanced,nobeforeafter,width=7.1cm,height=2.1cm,valign=center,colback=cder!6,colframe=cder,arc=2.5mm,boxrule=0.9pt,left=3mm]{\Large\bfseries\color{cder}d)}\quad\begin{minipage}[c]{5.3cm}\raggedright\Large $-2$\end{minipage}\end{tcolorbox}
\end{tcolorbox}
\newpage
\begin{tcolorbox}[enhanced,width=16cm,height=14cm,colback=white,colframe=cder,boxrule=2pt,arc=0mm,
left=6mm,right=6mm,top=0mm,bottom=4mm,sharp corners,
overlay={\fill[cder] (frame.north west) rectangle ([yshift=-1.25cm]frame.north east);
\node[anchor=west,text=white,font=\sffamily\bfseries\LARGE] at ([xshift=5mm,yshift=-0.62cm]frame.north west) {Question 39 D};
\node[anchor=east,text=white,font=\sffamily\itshape\large] at ([xshift=-5mm,yshift=-0.62cm]frame.north east) {Dérivation};}]
\vspace*{1.75cm}
\begin{minipage}[t][6.9cm][c]{\linewidth}\LARGE\centering
\begin{minipage}{0.95\linewidth}\raggedright\linespread{1.25}\selectfont Soit $f(x)=\dfrac{8}{x}$ sur $\mathbb{R}^*$. Que vaut $f'(4)$~?\end{minipage}
\end{minipage}

\vfill
\begin{tcolorbox}[enhanced,nobeforeafter,width=7.1cm,height=2.1cm,valign=center,colback=cder!6,colframe=cder,arc=2.5mm,boxrule=0.9pt,left=3mm]{\Large\bfseries\color{cder}a)}\quad\begin{minipage}[c]{5.3cm}\raggedright\Large $-2$\end{minipage}\end{tcolorbox}\hfill\begin{tcolorbox}[enhanced,nobeforeafter,width=7.1cm,height=2.1cm,valign=center,colback=cder!6,colframe=cder,arc=2.5mm,boxrule=0.9pt,left=3mm]{\Large\bfseries\color{cder}b)}\quad\begin{minipage}[c]{5.3cm}\raggedright\Large $32$\end{minipage}\end{tcolorbox}\par\vspace{2.5mm}
\begin{tcolorbox}[enhanced,nobeforeafter,width=7.1cm,height=2.1cm,valign=center,colback=cder!6,colframe=cder,arc=2.5mm,boxrule=0.9pt,left=3mm]{\Large\bfseries\color{cder}c)}\quad\begin{minipage}[c]{5.3cm}\raggedright\Large $\dfrac{1}{2}$\end{minipage}\end{tcolorbox}\hfill\begin{tcolorbox}[enhanced,nobeforeafter,width=7.1cm,height=2.1cm,valign=center,colback=cder!6,colframe=cder,arc=2.5mm,boxrule=0.9pt,left=3mm]{\Large\bfseries\color{cder}d)}\quad\begin{minipage}[c]{5.3cm}\raggedright\Large $-\dfrac{1}{2}$\end{minipage}\end{tcolorbox}
\end{tcolorbox}
\newpage
\begin{tcolorbox}[enhanced,width=16cm,height=14cm,colback=white,colframe=cder,boxrule=2pt,arc=0mm,
left=6mm,right=6mm,top=0mm,bottom=4mm,sharp corners,
overlay={\fill[cder] (frame.north west) rectangle ([yshift=-1.25cm]frame.north east);
\node[anchor=west,text=white,font=\sffamily\bfseries\LARGE] at ([xshift=5mm,yshift=-0.62cm]frame.north west) {Question 40 D};
\node[anchor=east,text=white,font=\sffamily\itshape\large] at ([xshift=-5mm,yshift=-0.62cm]frame.north east) {Dérivation};}]
\vspace*{1.75cm}
\begin{minipage}[t][6.9cm][c]{\linewidth}\LARGE\centering
\begin{minipage}{0.95\linewidth}\raggedright\linespread{1.25}\selectfont Soit $f(x)=\dfrac{-2}{x}$ sur $\mathbb{R}^*$. Que vaut $f'(1)$~?\end{minipage}
\end{minipage}

\vfill
\begin{tcolorbox}[enhanced,nobeforeafter,width=7.1cm,height=2.1cm,valign=center,colback=cder!6,colframe=cder,arc=2.5mm,boxrule=0.9pt,left=3mm]{\Large\bfseries\color{cder}a)}\quad\begin{minipage}[c]{5.3cm}\raggedright\Large $2$\end{minipage}\end{tcolorbox}\hfill\begin{tcolorbox}[enhanced,nobeforeafter,width=7.1cm,height=2.1cm,valign=center,colback=cder!6,colframe=cder,arc=2.5mm,boxrule=0.9pt,left=3mm]{\Large\bfseries\color{cder}b)}\quad\begin{minipage}[c]{5.3cm}\raggedright\Large $3$\end{minipage}\end{tcolorbox}\par\vspace{2.5mm}
\begin{tcolorbox}[enhanced,nobeforeafter,width=7.1cm,height=2.1cm,valign=center,colback=cder!6,colframe=cder,arc=2.5mm,boxrule=0.9pt,left=3mm]{\Large\bfseries\color{cder}c)}\quad\begin{minipage}[c]{5.3cm}\raggedright\Large $-2$\end{minipage}\end{tcolorbox}\hfill\begin{tcolorbox}[enhanced,nobeforeafter,width=7.1cm,height=2.1cm,valign=center,colback=cder!6,colframe=cder,arc=2.5mm,boxrule=0.9pt,left=3mm]{\Large\bfseries\color{cder}d)}\quad\begin{minipage}[c]{5.3cm}\raggedright\Large $1$\end{minipage}\end{tcolorbox}
\end{tcolorbox}
\newpage
\begin{tcolorbox}[enhanced,width=16cm,height=14cm,colback=white,colframe=cder,boxrule=2pt,arc=0mm,
left=6mm,right=6mm,top=0mm,bottom=4mm,sharp corners,
overlay={\fill[cder] (frame.north west) rectangle ([yshift=-1.25cm]frame.north east);
\node[anchor=west,text=white,font=\sffamily\bfseries\LARGE] at ([xshift=5mm,yshift=-0.62cm]frame.north west) {Question 41 D};
\node[anchor=east,text=white,font=\sffamily\itshape\large] at ([xshift=-5mm,yshift=-0.62cm]frame.north east) {Dérivation};}]
\vspace*{1.75cm}
\begin{minipage}[t][6.9cm][c]{\linewidth}\LARGE\centering
\begin{minipage}{0.95\linewidth}\raggedright\linespread{1.25}\selectfont Soit $f(x)=\sqrt{x}$. Que vaut $f'(64)$~?\end{minipage}
\end{minipage}

\vfill
\begin{tcolorbox}[enhanced,nobeforeafter,width=7.1cm,height=2.1cm,valign=center,colback=cder!6,colframe=cder,arc=2.5mm,boxrule=0.9pt,left=3mm]{\Large\bfseries\color{cder}a)}\quad\begin{minipage}[c]{5.3cm}\raggedright\Large $\dfrac{1}{16}$\end{minipage}\end{tcolorbox}\hfill\begin{tcolorbox}[enhanced,nobeforeafter,width=7.1cm,height=2.1cm,valign=center,colback=cder!6,colframe=cder,arc=2.5mm,boxrule=0.9pt,left=3mm]{\Large\bfseries\color{cder}b)}\quad\begin{minipage}[c]{5.3cm}\raggedright\Large $4$\end{minipage}\end{tcolorbox}\par\vspace{2.5mm}
\begin{tcolorbox}[enhanced,nobeforeafter,width=7.1cm,height=2.1cm,valign=center,colback=cder!6,colframe=cder,arc=2.5mm,boxrule=0.9pt,left=3mm]{\Large\bfseries\color{cder}c)}\quad\begin{minipage}[c]{5.3cm}\raggedright\Large $\dfrac{1}{8}$\end{minipage}\end{tcolorbox}\hfill\begin{tcolorbox}[enhanced,nobeforeafter,width=7.1cm,height=2.1cm,valign=center,colback=cder!6,colframe=cder,arc=2.5mm,boxrule=0.9pt,left=3mm]{\Large\bfseries\color{cder}d)}\quad\begin{minipage}[c]{5.3cm}\raggedright\Large $16$\end{minipage}\end{tcolorbox}
\end{tcolorbox}
\newpage
\begin{tcolorbox}[enhanced,width=16cm,height=14cm,colback=white,colframe=cder,boxrule=2pt,arc=0mm,
left=6mm,right=6mm,top=0mm,bottom=4mm,sharp corners,
overlay={\fill[cder] (frame.north west) rectangle ([yshift=-1.25cm]frame.north east);
\node[anchor=west,text=white,font=\sffamily\bfseries\LARGE] at ([xshift=5mm,yshift=-0.62cm]frame.north west) {Question 42 D};
\node[anchor=east,text=white,font=\sffamily\itshape\large] at ([xshift=-5mm,yshift=-0.62cm]frame.north east) {Dérivation};}]
\vspace*{1.75cm}
\begin{minipage}[t][6.9cm][c]{\linewidth}\LARGE\centering
\begin{minipage}{0.95\linewidth}\raggedright\linespread{1.25}\selectfont Soit $f(x)=\sqrt{x}$. Que vaut $f'(100)$~?\end{minipage}
\end{minipage}

\vfill
\begin{tcolorbox}[enhanced,nobeforeafter,width=7.1cm,height=2.1cm,valign=center,colback=cder!6,colframe=cder,arc=2.5mm,boxrule=0.9pt,left=3mm]{\Large\bfseries\color{cder}a)}\quad\begin{minipage}[c]{5.3cm}\raggedright\Large $20$\end{minipage}\end{tcolorbox}\hfill\begin{tcolorbox}[enhanced,nobeforeafter,width=7.1cm,height=2.1cm,valign=center,colback=cder!6,colframe=cder,arc=2.5mm,boxrule=0.9pt,left=3mm]{\Large\bfseries\color{cder}b)}\quad\begin{minipage}[c]{5.3cm}\raggedright\Large $\dfrac{1}{10}$\end{minipage}\end{tcolorbox}\par\vspace{2.5mm}
\begin{tcolorbox}[enhanced,nobeforeafter,width=7.1cm,height=2.1cm,valign=center,colback=cder!6,colframe=cder,arc=2.5mm,boxrule=0.9pt,left=3mm]{\Large\bfseries\color{cder}c)}\quad\begin{minipage}[c]{5.3cm}\raggedright\Large $5$\end{minipage}\end{tcolorbox}\hfill\begin{tcolorbox}[enhanced,nobeforeafter,width=7.1cm,height=2.1cm,valign=center,colback=cder!6,colframe=cder,arc=2.5mm,boxrule=0.9pt,left=3mm]{\Large\bfseries\color{cder}d)}\quad\begin{minipage}[c]{5.3cm}\raggedright\Large $\dfrac{1}{20}$\end{minipage}\end{tcolorbox}
\end{tcolorbox}
\newpage
\begin{tcolorbox}[enhanced,width=16cm,height=14cm,colback=white,colframe=cder,boxrule=2pt,arc=0mm,
left=6mm,right=6mm,top=0mm,bottom=4mm,sharp corners,
overlay={\fill[cder] (frame.north west) rectangle ([yshift=-1.25cm]frame.north east);
\node[anchor=west,text=white,font=\sffamily\bfseries\LARGE] at ([xshift=5mm,yshift=-0.62cm]frame.north west) {Question 43 D};
\node[anchor=east,text=white,font=\sffamily\itshape\large] at ([xshift=-5mm,yshift=-0.62cm]frame.north east) {Dérivation};}]
\vspace*{1.75cm}
\begin{minipage}[t][6.9cm][c]{\linewidth}\LARGE\centering
\begin{minipage}{0.95\linewidth}\raggedright\linespread{1.25}\selectfont Soit $f(x)=(3x-1)e^x$. Que vaut $f'(0)$~?\end{minipage}
\end{minipage}

\vfill
\begin{tcolorbox}[enhanced,nobeforeafter,width=7.1cm,height=2.1cm,valign=center,colback=cder!6,colframe=cder,arc=2.5mm,boxrule=0.9pt,left=3mm]{\Large\bfseries\color{cder}a)}\quad\begin{minipage}[c]{5.3cm}\raggedright\Large $2$\end{minipage}\end{tcolorbox}\hfill\begin{tcolorbox}[enhanced,nobeforeafter,width=7.1cm,height=2.1cm,valign=center,colback=cder!6,colframe=cder,arc=2.5mm,boxrule=0.9pt,left=3mm]{\Large\bfseries\color{cder}b)}\quad\begin{minipage}[c]{5.3cm}\raggedright\Large $-1$\end{minipage}\end{tcolorbox}\par\vspace{2.5mm}
\begin{tcolorbox}[enhanced,nobeforeafter,width=7.1cm,height=2.1cm,valign=center,colback=cder!6,colframe=cder,arc=2.5mm,boxrule=0.9pt,left=3mm]{\Large\bfseries\color{cder}c)}\quad\begin{minipage}[c]{5.3cm}\raggedright\Large $-3$\end{minipage}\end{tcolorbox}\hfill\begin{tcolorbox}[enhanced,nobeforeafter,width=7.1cm,height=2.1cm,valign=center,colback=cder!6,colframe=cder,arc=2.5mm,boxrule=0.9pt,left=3mm]{\Large\bfseries\color{cder}d)}\quad\begin{minipage}[c]{5.3cm}\raggedright\Large $3$\end{minipage}\end{tcolorbox}
\end{tcolorbox}
\newpage
\begin{tcolorbox}[enhanced,width=16cm,height=14cm,colback=white,colframe=cder,boxrule=2pt,arc=0mm,
left=6mm,right=6mm,top=0mm,bottom=4mm,sharp corners,
overlay={\fill[cder] (frame.north west) rectangle ([yshift=-1.25cm]frame.north east);
\node[anchor=west,text=white,font=\sffamily\bfseries\LARGE] at ([xshift=5mm,yshift=-0.62cm]frame.north west) {Question 44 D};
\node[anchor=east,text=white,font=\sffamily\itshape\large] at ([xshift=-5mm,yshift=-0.62cm]frame.north east) {Dérivation};}]
\vspace*{1.75cm}
\begin{minipage}[t][6.9cm][c]{\linewidth}\LARGE\centering
\begin{minipage}{0.95\linewidth}\raggedright\linespread{1.25}\selectfont Soit $f(x)=(5x+2)e^x$. Que vaut $f'(0)$~?\end{minipage}
\end{minipage}

\vfill
\begin{tcolorbox}[enhanced,nobeforeafter,width=7.1cm,height=2.1cm,valign=center,colback=cder!6,colframe=cder,arc=2.5mm,boxrule=0.9pt,left=3mm]{\Large\bfseries\color{cder}a)}\quad\begin{minipage}[c]{5.3cm}\raggedright\Large $10$\end{minipage}\end{tcolorbox}\hfill\begin{tcolorbox}[enhanced,nobeforeafter,width=7.1cm,height=2.1cm,valign=center,colback=cder!6,colframe=cder,arc=2.5mm,boxrule=0.9pt,left=3mm]{\Large\bfseries\color{cder}b)}\quad\begin{minipage}[c]{5.3cm}\raggedright\Large $2$\end{minipage}\end{tcolorbox}\par\vspace{2.5mm}
\begin{tcolorbox}[enhanced,nobeforeafter,width=7.1cm,height=2.1cm,valign=center,colback=cder!6,colframe=cder,arc=2.5mm,boxrule=0.9pt,left=3mm]{\Large\bfseries\color{cder}c)}\quad\begin{minipage}[c]{5.3cm}\raggedright\Large $7$\end{minipage}\end{tcolorbox}\hfill\begin{tcolorbox}[enhanced,nobeforeafter,width=7.1cm,height=2.1cm,valign=center,colback=cder!6,colframe=cder,arc=2.5mm,boxrule=0.9pt,left=3mm]{\Large\bfseries\color{cder}d)}\quad\begin{minipage}[c]{5.3cm}\raggedright\Large $5$\end{minipage}\end{tcolorbox}
\end{tcolorbox}
\newpage
\begin{tcolorbox}[enhanced,width=16cm,height=14cm,colback=white,colframe=cder,boxrule=2pt,arc=0mm,
left=6mm,right=6mm,top=0mm,bottom=4mm,sharp corners,
overlay={\fill[cder] (frame.north west) rectangle ([yshift=-1.25cm]frame.north east);
\node[anchor=west,text=white,font=\sffamily\bfseries\LARGE] at ([xshift=5mm,yshift=-0.62cm]frame.north west) {Question 45 D};
\node[anchor=east,text=white,font=\sffamily\itshape\large] at ([xshift=-5mm,yshift=-0.62cm]frame.north east) {Dérivation};}]
\vspace*{1.75cm}
\begin{minipage}[t][6.9cm][c]{\linewidth}\LARGE\centering
\begin{minipage}{0.95\linewidth}\raggedright\linespread{1.25}\selectfont La fonction $f$ est dérivable sur $\mathbb{R}$ et $f'(x)=3x+9$. Sur quel intervalle $f$ est-elle décroissante~?\end{minipage}
\end{minipage}

\vfill
\begin{tcolorbox}[enhanced,nobeforeafter,width=7.1cm,height=2.1cm,valign=center,colback=cder!6,colframe=cder,arc=2.5mm,boxrule=0.9pt,left=3mm]{\Large\bfseries\color{cder}a)}\quad\begin{minipage}[c]{5.3cm}\raggedright\Large $[-3\,;\,+\infty[$\end{minipage}\end{tcolorbox}\hfill\begin{tcolorbox}[enhanced,nobeforeafter,width=7.1cm,height=2.1cm,valign=center,colback=cder!6,colframe=cder,arc=2.5mm,boxrule=0.9pt,left=3mm]{\Large\bfseries\color{cder}b)}\quad\begin{minipage}[c]{5.3cm}\raggedright\Large $]-\infty\,;\,3]$\end{minipage}\end{tcolorbox}\par\vspace{2.5mm}
\begin{tcolorbox}[enhanced,nobeforeafter,width=7.1cm,height=2.1cm,valign=center,colback=cder!6,colframe=cder,arc=2.5mm,boxrule=0.9pt,left=3mm]{\Large\bfseries\color{cder}c)}\quad\begin{minipage}[c]{5.3cm}\raggedright\Large $[3\,;\,+\infty[$\end{minipage}\end{tcolorbox}\hfill\begin{tcolorbox}[enhanced,nobeforeafter,width=7.1cm,height=2.1cm,valign=center,colback=cder!6,colframe=cder,arc=2.5mm,boxrule=0.9pt,left=3mm]{\Large\bfseries\color{cder}d)}\quad\begin{minipage}[c]{5.3cm}\raggedright\Large $]-\infty\,;\,-3]$\end{minipage}\end{tcolorbox}
\end{tcolorbox}
\newpage
\begin{tcolorbox}[enhanced,width=16cm,height=14cm,colback=white,colframe=cder,boxrule=2pt,arc=0mm,
left=6mm,right=6mm,top=0mm,bottom=4mm,sharp corners,
overlay={\fill[cder] (frame.north west) rectangle ([yshift=-1.25cm]frame.north east);
\node[anchor=west,text=white,font=\sffamily\bfseries\LARGE] at ([xshift=5mm,yshift=-0.62cm]frame.north west) {Question 46 D};
\node[anchor=east,text=white,font=\sffamily\itshape\large] at ([xshift=-5mm,yshift=-0.62cm]frame.north east) {Dérivation};}]
\vspace*{1.75cm}
\begin{minipage}[t][6.9cm][c]{\linewidth}\LARGE\centering
\begin{minipage}{0.95\linewidth}\raggedright\linespread{1.25}\selectfont La fonction $f$ est dérivable sur $\mathbb{R}$ et $f'(x)=-4x-12$. Sur quel intervalle $f$ est-elle décroissante~?\end{minipage}
\end{minipage}

\vfill
\begin{tcolorbox}[enhanced,nobeforeafter,width=7.1cm,height=2.1cm,valign=center,colback=cder!6,colframe=cder,arc=2.5mm,boxrule=0.9pt,left=3mm]{\Large\bfseries\color{cder}a)}\quad\begin{minipage}[c]{5.3cm}\raggedright\Large $[-3\,;\,+\infty[$\end{minipage}\end{tcolorbox}\hfill\begin{tcolorbox}[enhanced,nobeforeafter,width=7.1cm,height=2.1cm,valign=center,colback=cder!6,colframe=cder,arc=2.5mm,boxrule=0.9pt,left=3mm]{\Large\bfseries\color{cder}b)}\quad\begin{minipage}[c]{5.3cm}\raggedright\Large $[3\,;\,+\infty[$\end{minipage}\end{tcolorbox}\par\vspace{2.5mm}
\begin{tcolorbox}[enhanced,nobeforeafter,width=7.1cm,height=2.1cm,valign=center,colback=cder!6,colframe=cder,arc=2.5mm,boxrule=0.9pt,left=3mm]{\Large\bfseries\color{cder}c)}\quad\begin{minipage}[c]{5.3cm}\raggedright\Large $]-\infty\,;\,-3]$\end{minipage}\end{tcolorbox}\hfill\begin{tcolorbox}[enhanced,nobeforeafter,width=7.1cm,height=2.1cm,valign=center,colback=cder!6,colframe=cder,arc=2.5mm,boxrule=0.9pt,left=3mm]{\Large\bfseries\color{cder}d)}\quad\begin{minipage}[c]{5.3cm}\raggedright\Large $]-\infty\,;\,3]$\end{minipage}\end{tcolorbox}
\end{tcolorbox}
\newpage
\begin{tcolorbox}[enhanced,width=16cm,height=14cm,colback=white,colframe=cder,boxrule=2pt,arc=0mm,
left=6mm,right=6mm,top=0mm,bottom=4mm,sharp corners,
overlay={\fill[cder] (frame.north west) rectangle ([yshift=-1.25cm]frame.north east);
\node[anchor=west,text=white,font=\sffamily\bfseries\LARGE] at ([xshift=5mm,yshift=-0.62cm]frame.north west) {Question 47 D};
\node[anchor=east,text=white,font=\sffamily\itshape\large] at ([xshift=-5mm,yshift=-0.62cm]frame.north east) {Dérivation};}]
\vspace*{1.75cm}
\begin{minipage}[t][6.9cm][c]{\linewidth}\LARGE\centering
\begin{minipage}{0.95\linewidth}\raggedright\linespread{1.25}\selectfont Soit $f(x)=3x^2-6x+5$. Quel est le minimum de $f$ sur $\mathbb{R}$~?\end{minipage}
\end{minipage}

\vfill
\begin{tcolorbox}[enhanced,nobeforeafter,width=7.1cm,height=2.1cm,valign=center,colback=cder!6,colframe=cder,arc=2.5mm,boxrule=0.9pt,left=3mm]{\Large\bfseries\color{cder}a)}\quad\begin{minipage}[c]{5.3cm}\raggedright\Large $1$\end{minipage}\end{tcolorbox}\hfill\begin{tcolorbox}[enhanced,nobeforeafter,width=7.1cm,height=2.1cm,valign=center,colback=cder!6,colframe=cder,arc=2.5mm,boxrule=0.9pt,left=3mm]{\Large\bfseries\color{cder}b)}\quad\begin{minipage}[c]{5.3cm}\raggedright\Large $2$\end{minipage}\end{tcolorbox}\par\vspace{2.5mm}
\begin{tcolorbox}[enhanced,nobeforeafter,width=7.1cm,height=2.1cm,valign=center,colback=cder!6,colframe=cder,arc=2.5mm,boxrule=0.9pt,left=3mm]{\Large\bfseries\color{cder}c)}\quad\begin{minipage}[c]{5.3cm}\raggedright\Large $5$\end{minipage}\end{tcolorbox}\hfill\begin{tcolorbox}[enhanced,nobeforeafter,width=7.1cm,height=2.1cm,valign=center,colback=cder!6,colframe=cder,arc=2.5mm,boxrule=0.9pt,left=3mm]{\Large\bfseries\color{cder}d)}\quad\begin{minipage}[c]{5.3cm}\raggedright\Large $-2$\end{minipage}\end{tcolorbox}
\end{tcolorbox}
\newpage
\begin{tcolorbox}[enhanced,width=16cm,height=14cm,colback=white,colframe=cder,boxrule=2pt,arc=0mm,
left=6mm,right=6mm,top=0mm,bottom=4mm,sharp corners,
overlay={\fill[cder] (frame.north west) rectangle ([yshift=-1.25cm]frame.north east);
\node[anchor=west,text=white,font=\sffamily\bfseries\LARGE] at ([xshift=5mm,yshift=-0.62cm]frame.north west) {Question 48 D};
\node[anchor=east,text=white,font=\sffamily\itshape\large] at ([xshift=-5mm,yshift=-0.62cm]frame.north east) {Dérivation};}]
\vspace*{1.75cm}
\begin{minipage}[t][6.9cm][c]{\linewidth}\LARGE\centering
\begin{minipage}{0.95\linewidth}\raggedright\linespread{1.25}\selectfont Soit $f(x)=x^2-10x+20$. Quel est le minimum de $f$ sur $\mathbb{R}$~?\end{minipage}
\end{minipage}

\vfill
\begin{tcolorbox}[enhanced,nobeforeafter,width=7.1cm,height=2.1cm,valign=center,colback=cder!6,colframe=cder,arc=2.5mm,boxrule=0.9pt,left=3mm]{\Large\bfseries\color{cder}a)}\quad\begin{minipage}[c]{5.3cm}\raggedright\Large $-3$\end{minipage}\end{tcolorbox}\hfill\begin{tcolorbox}[enhanced,nobeforeafter,width=7.1cm,height=2.1cm,valign=center,colback=cder!6,colframe=cder,arc=2.5mm,boxrule=0.9pt,left=3mm]{\Large\bfseries\color{cder}b)}\quad\begin{minipage}[c]{5.3cm}\raggedright\Large $20$\end{minipage}\end{tcolorbox}\par\vspace{2.5mm}
\begin{tcolorbox}[enhanced,nobeforeafter,width=7.1cm,height=2.1cm,valign=center,colback=cder!6,colframe=cder,arc=2.5mm,boxrule=0.9pt,left=3mm]{\Large\bfseries\color{cder}c)}\quad\begin{minipage}[c]{5.3cm}\raggedright\Large $-5$\end{minipage}\end{tcolorbox}\hfill\begin{tcolorbox}[enhanced,nobeforeafter,width=7.1cm,height=2.1cm,valign=center,colback=cder!6,colframe=cder,arc=2.5mm,boxrule=0.9pt,left=3mm]{\Large\bfseries\color{cder}d)}\quad\begin{minipage}[c]{5.3cm}\raggedright\Large $5$\end{minipage}\end{tcolorbox}
\end{tcolorbox}
\newpage
\begin{tcolorbox}[enhanced,width=16cm,height=14cm,colback=white,colframe=ctri,boxrule=2pt,arc=0mm,
left=6mm,right=6mm,top=0mm,bottom=4mm,sharp corners,
overlay={\fill[ctri] (frame.north west) rectangle ([yshift=-1.25cm]frame.north east);
\node[anchor=west,text=white,font=\sffamily\bfseries\LARGE] at ([xshift=5mm,yshift=-0.62cm]frame.north west) {Question 49 D};
\node[anchor=east,text=white,font=\sffamily\itshape\large] at ([xshift=-5mm,yshift=-0.62cm]frame.north east) {Trigonométrie et exponentielle};}]
\vspace*{1.75cm}
\begin{minipage}[t][6.9cm][c]{\linewidth}\LARGE\centering
\begin{minipage}{0.95\linewidth}\raggedright\linespread{1.25}\selectfont Quelle est la mesure en radians d'un angle de $90^\circ$~?\end{minipage}
\end{minipage}

\vfill
\begin{tcolorbox}[enhanced,nobeforeafter,width=7.1cm,height=2.1cm,valign=center,colback=ctri!6,colframe=ctri,arc=2.5mm,boxrule=0.9pt,left=3mm]{\Large\bfseries\color{ctri}a)}\quad\begin{minipage}[c]{5.3cm}\raggedright\Large $\dfrac{\pi}{4}$\end{minipage}\end{tcolorbox}\hfill\begin{tcolorbox}[enhanced,nobeforeafter,width=7.1cm,height=2.1cm,valign=center,colback=ctri!6,colframe=ctri,arc=2.5mm,boxrule=0.9pt,left=3mm]{\Large\bfseries\color{ctri}b)}\quad\begin{minipage}[c]{5.3cm}\raggedright\Large $2\pi$\end{minipage}\end{tcolorbox}\par\vspace{2.5mm}
\begin{tcolorbox}[enhanced,nobeforeafter,width=7.1cm,height=2.1cm,valign=center,colback=ctri!6,colframe=ctri,arc=2.5mm,boxrule=0.9pt,left=3mm]{\Large\bfseries\color{ctri}c)}\quad\begin{minipage}[c]{5.3cm}\raggedright\Large $\dfrac{\pi}{2}$\end{minipage}\end{tcolorbox}\hfill\begin{tcolorbox}[enhanced,nobeforeafter,width=7.1cm,height=2.1cm,valign=center,colback=ctri!6,colframe=ctri,arc=2.5mm,boxrule=0.9pt,left=3mm]{\Large\bfseries\color{ctri}d)}\quad\begin{minipage}[c]{5.3cm}\raggedright\Large $\pi$\end{minipage}\end{tcolorbox}
\end{tcolorbox}
\newpage
\begin{tcolorbox}[enhanced,width=16cm,height=14cm,colback=white,colframe=ctri,boxrule=2pt,arc=0mm,
left=6mm,right=6mm,top=0mm,bottom=4mm,sharp corners,
overlay={\fill[ctri] (frame.north west) rectangle ([yshift=-1.25cm]frame.north east);
\node[anchor=west,text=white,font=\sffamily\bfseries\LARGE] at ([xshift=5mm,yshift=-0.62cm]frame.north west) {Question 50 D};
\node[anchor=east,text=white,font=\sffamily\itshape\large] at ([xshift=-5mm,yshift=-0.62cm]frame.north east) {Trigonométrie et exponentielle};}]
\vspace*{1.75cm}
\begin{minipage}[t][6.9cm][c]{\linewidth}\LARGE\centering
\begin{minipage}{0.95\linewidth}\raggedright\linespread{1.25}\selectfont Quelle est la mesure en radians d'un angle de $270^\circ$~?\end{minipage}
\end{minipage}

\vfill
\begin{tcolorbox}[enhanced,nobeforeafter,width=7.1cm,height=2.1cm,valign=center,colback=ctri!6,colframe=ctri,arc=2.5mm,boxrule=0.9pt,left=3mm]{\Large\bfseries\color{ctri}a)}\quad\begin{minipage}[c]{5.3cm}\raggedright\Large $3\pi$\end{minipage}\end{tcolorbox}\hfill\begin{tcolorbox}[enhanced,nobeforeafter,width=7.1cm,height=2.1cm,valign=center,colback=ctri!6,colframe=ctri,arc=2.5mm,boxrule=0.9pt,left=3mm]{\Large\bfseries\color{ctri}b)}\quad\begin{minipage}[c]{5.3cm}\raggedright\Large $\dfrac{3\pi}{4}$\end{minipage}\end{tcolorbox}\par\vspace{2.5mm}
\begin{tcolorbox}[enhanced,nobeforeafter,width=7.1cm,height=2.1cm,valign=center,colback=ctri!6,colframe=ctri,arc=2.5mm,boxrule=0.9pt,left=3mm]{\Large\bfseries\color{ctri}c)}\quad\begin{minipage}[c]{5.3cm}\raggedright\Large $\dfrac{2\pi}{3}$\end{minipage}\end{tcolorbox}\hfill\begin{tcolorbox}[enhanced,nobeforeafter,width=7.1cm,height=2.1cm,valign=center,colback=ctri!6,colframe=ctri,arc=2.5mm,boxrule=0.9pt,left=3mm]{\Large\bfseries\color{ctri}d)}\quad\begin{minipage}[c]{5.3cm}\raggedright\Large $\dfrac{3\pi}{2}$\end{minipage}\end{tcolorbox}
\end{tcolorbox}
\newpage
\begin{tcolorbox}[enhanced,width=16cm,height=14cm,colback=white,colframe=ctri,boxrule=2pt,arc=0mm,
left=6mm,right=6mm,top=0mm,bottom=4mm,sharp corners,
overlay={\fill[ctri] (frame.north west) rectangle ([yshift=-1.25cm]frame.north east);
\node[anchor=west,text=white,font=\sffamily\bfseries\LARGE] at ([xshift=5mm,yshift=-0.62cm]frame.north west) {Question 51 D};
\node[anchor=east,text=white,font=\sffamily\itshape\large] at ([xshift=-5mm,yshift=-0.62cm]frame.north east) {Trigonométrie et exponentielle};}]
\vspace*{1.75cm}
\begin{minipage}[t][6.9cm][c]{\linewidth}\LARGE\centering
\begin{minipage}{0.95\linewidth}\raggedright\linespread{1.25}\selectfont Quelle est la mesure en degrés d'un angle de $\dfrac{\pi}{9}$ radians~?\end{minipage}
\end{minipage}

\vfill
\begin{tcolorbox}[enhanced,nobeforeafter,width=7.1cm,height=2.1cm,valign=center,colback=ctri!6,colframe=ctri,arc=2.5mm,boxrule=0.9pt,left=3mm]{\Large\bfseries\color{ctri}a)}\quad\begin{minipage}[c]{5.3cm}\raggedright\Large $20^\circ$\end{minipage}\end{tcolorbox}\hfill\begin{tcolorbox}[enhanced,nobeforeafter,width=7.1cm,height=2.1cm,valign=center,colback=ctri!6,colframe=ctri,arc=2.5mm,boxrule=0.9pt,left=3mm]{\Large\bfseries\color{ctri}b)}\quad\begin{minipage}[c]{5.3cm}\raggedright\Large $40^\circ$\end{minipage}\end{tcolorbox}\par\vspace{2.5mm}
\begin{tcolorbox}[enhanced,nobeforeafter,width=7.1cm,height=2.1cm,valign=center,colback=ctri!6,colframe=ctri,arc=2.5mm,boxrule=0.9pt,left=3mm]{\Large\bfseries\color{ctri}c)}\quad\begin{minipage}[c]{5.3cm}\raggedright\Large $340^\circ$\end{minipage}\end{tcolorbox}\hfill\begin{tcolorbox}[enhanced,nobeforeafter,width=7.1cm,height=2.1cm,valign=center,colback=ctri!6,colframe=ctri,arc=2.5mm,boxrule=0.9pt,left=3mm]{\Large\bfseries\color{ctri}d)}\quad\begin{minipage}[c]{5.3cm}\raggedright\Large $10^\circ$\end{minipage}\end{tcolorbox}
\end{tcolorbox}
\newpage
\begin{tcolorbox}[enhanced,width=16cm,height=14cm,colback=white,colframe=ctri,boxrule=2pt,arc=0mm,
left=6mm,right=6mm,top=0mm,bottom=4mm,sharp corners,
overlay={\fill[ctri] (frame.north west) rectangle ([yshift=-1.25cm]frame.north east);
\node[anchor=west,text=white,font=\sffamily\bfseries\LARGE] at ([xshift=5mm,yshift=-0.62cm]frame.north west) {Question 52 D};
\node[anchor=east,text=white,font=\sffamily\itshape\large] at ([xshift=-5mm,yshift=-0.62cm]frame.north east) {Trigonométrie et exponentielle};}]
\vspace*{1.75cm}
\begin{minipage}[t][6.9cm][c]{\linewidth}\LARGE\centering
\begin{minipage}{0.95\linewidth}\raggedright\linespread{1.25}\selectfont Quelle est la mesure en degrés d'un angle de $\dfrac{4\pi}{3}$ radians~?\end{minipage}
\end{minipage}

\vfill
\begin{tcolorbox}[enhanced,nobeforeafter,width=7.1cm,height=2.1cm,valign=center,colback=ctri!6,colframe=ctri,arc=2.5mm,boxrule=0.9pt,left=3mm]{\Large\bfseries\color{ctri}a)}\quad\begin{minipage}[c]{5.3cm}\raggedright\Large $150^\circ$\end{minipage}\end{tcolorbox}\hfill\begin{tcolorbox}[enhanced,nobeforeafter,width=7.1cm,height=2.1cm,valign=center,colback=ctri!6,colframe=ctri,arc=2.5mm,boxrule=0.9pt,left=3mm]{\Large\bfseries\color{ctri}b)}\quad\begin{minipage}[c]{5.3cm}\raggedright\Large $120^\circ$\end{minipage}\end{tcolorbox}\par\vspace{2.5mm}
\begin{tcolorbox}[enhanced,nobeforeafter,width=7.1cm,height=2.1cm,valign=center,colback=ctri!6,colframe=ctri,arc=2.5mm,boxrule=0.9pt,left=3mm]{\Large\bfseries\color{ctri}c)}\quad\begin{minipage}[c]{5.3cm}\raggedright\Large $480^\circ$\end{minipage}\end{tcolorbox}\hfill\begin{tcolorbox}[enhanced,nobeforeafter,width=7.1cm,height=2.1cm,valign=center,colback=ctri!6,colframe=ctri,arc=2.5mm,boxrule=0.9pt,left=3mm]{\Large\bfseries\color{ctri}d)}\quad\begin{minipage}[c]{5.3cm}\raggedright\Large $240^\circ$\end{minipage}\end{tcolorbox}
\end{tcolorbox}
\newpage
\begin{tcolorbox}[enhanced,width=16cm,height=14cm,colback=white,colframe=ctri,boxrule=2pt,arc=0mm,
left=6mm,right=6mm,top=0mm,bottom=4mm,sharp corners,
overlay={\fill[ctri] (frame.north west) rectangle ([yshift=-1.25cm]frame.north east);
\node[anchor=west,text=white,font=\sffamily\bfseries\LARGE] at ([xshift=5mm,yshift=-0.62cm]frame.north west) {Question 53 D};
\node[anchor=east,text=white,font=\sffamily\itshape\large] at ([xshift=-5mm,yshift=-0.62cm]frame.north east) {Trigonométrie et exponentielle};}]
\vspace*{1.75cm}
\begin{minipage}[t][6.9cm][c]{\linewidth}\LARGE\centering
\begin{minipage}{0.95\linewidth}\raggedright\linespread{1.25}\selectfont Que vaut $\cos\left(\dfrac{\pi}{2}\right)$~?\end{minipage}
\end{minipage}

\vfill
\begin{tcolorbox}[enhanced,nobeforeafter,width=7.1cm,height=2.1cm,valign=center,colback=ctri!6,colframe=ctri,arc=2.5mm,boxrule=0.9pt,left=3mm]{\Large\bfseries\color{ctri}a)}\quad\begin{minipage}[c]{5.3cm}\raggedright\Large $0$\end{minipage}\end{tcolorbox}\hfill\begin{tcolorbox}[enhanced,nobeforeafter,width=7.1cm,height=2.1cm,valign=center,colback=ctri!6,colframe=ctri,arc=2.5mm,boxrule=0.9pt,left=3mm]{\Large\bfseries\color{ctri}b)}\quad\begin{minipage}[c]{5.3cm}\raggedright\Large $\dfrac12$\end{minipage}\end{tcolorbox}\par\vspace{2.5mm}
\begin{tcolorbox}[enhanced,nobeforeafter,width=7.1cm,height=2.1cm,valign=center,colback=ctri!6,colframe=ctri,arc=2.5mm,boxrule=0.9pt,left=3mm]{\Large\bfseries\color{ctri}c)}\quad\begin{minipage}[c]{5.3cm}\raggedright\Large $1$\end{minipage}\end{tcolorbox}\hfill\begin{tcolorbox}[enhanced,nobeforeafter,width=7.1cm,height=2.1cm,valign=center,colback=ctri!6,colframe=ctri,arc=2.5mm,boxrule=0.9pt,left=3mm]{\Large\bfseries\color{ctri}d)}\quad\begin{minipage}[c]{5.3cm}\raggedright\Large $-\dfrac12$\end{minipage}\end{tcolorbox}
\end{tcolorbox}
\newpage
\begin{tcolorbox}[enhanced,width=16cm,height=14cm,colback=white,colframe=ctri,boxrule=2pt,arc=0mm,
left=6mm,right=6mm,top=0mm,bottom=4mm,sharp corners,
overlay={\fill[ctri] (frame.north west) rectangle ([yshift=-1.25cm]frame.north east);
\node[anchor=west,text=white,font=\sffamily\bfseries\LARGE] at ([xshift=5mm,yshift=-0.62cm]frame.north west) {Question 54 D};
\node[anchor=east,text=white,font=\sffamily\itshape\large] at ([xshift=-5mm,yshift=-0.62cm]frame.north east) {Trigonométrie et exponentielle};}]
\vspace*{1.75cm}
\begin{minipage}[t][6.9cm][c]{\linewidth}\LARGE\centering
\begin{minipage}{0.95\linewidth}\raggedright\linespread{1.25}\selectfont Que vaut $\cos\left(\dfrac{5\pi}{6}\right)$~?\end{minipage}
\end{minipage}

\vfill
\begin{tcolorbox}[enhanced,nobeforeafter,width=7.1cm,height=2.1cm,valign=center,colback=ctri!6,colframe=ctri,arc=2.5mm,boxrule=0.9pt,left=3mm]{\Large\bfseries\color{ctri}a)}\quad\begin{minipage}[c]{5.3cm}\raggedright\Large $-\dfrac{\sqrt2}{2}$\end{minipage}\end{tcolorbox}\hfill\begin{tcolorbox}[enhanced,nobeforeafter,width=7.1cm,height=2.1cm,valign=center,colback=ctri!6,colframe=ctri,arc=2.5mm,boxrule=0.9pt,left=3mm]{\Large\bfseries\color{ctri}b)}\quad\begin{minipage}[c]{5.3cm}\raggedright\Large $-\dfrac{\sqrt3}{2}$\end{minipage}\end{tcolorbox}\par\vspace{2.5mm}
\begin{tcolorbox}[enhanced,nobeforeafter,width=7.1cm,height=2.1cm,valign=center,colback=ctri!6,colframe=ctri,arc=2.5mm,boxrule=0.9pt,left=3mm]{\Large\bfseries\color{ctri}c)}\quad\begin{minipage}[c]{5.3cm}\raggedright\Large $\dfrac{\sqrt3}{2}$\end{minipage}\end{tcolorbox}\hfill\begin{tcolorbox}[enhanced,nobeforeafter,width=7.1cm,height=2.1cm,valign=center,colback=ctri!6,colframe=ctri,arc=2.5mm,boxrule=0.9pt,left=3mm]{\Large\bfseries\color{ctri}d)}\quad\begin{minipage}[c]{5.3cm}\raggedright\Large $-1$\end{minipage}\end{tcolorbox}
\end{tcolorbox}
\newpage
\begin{tcolorbox}[enhanced,width=16cm,height=14cm,colback=white,colframe=ctri,boxrule=2pt,arc=0mm,
left=6mm,right=6mm,top=0mm,bottom=4mm,sharp corners,
overlay={\fill[ctri] (frame.north west) rectangle ([yshift=-1.25cm]frame.north east);
\node[anchor=west,text=white,font=\sffamily\bfseries\LARGE] at ([xshift=5mm,yshift=-0.62cm]frame.north west) {Question 55 D};
\node[anchor=east,text=white,font=\sffamily\itshape\large] at ([xshift=-5mm,yshift=-0.62cm]frame.north east) {Trigonométrie et exponentielle};}]
\vspace*{1.75cm}
\begin{minipage}[t][6.9cm][c]{\linewidth}\LARGE\centering
\begin{minipage}{0.95\linewidth}\raggedright\linespread{1.25}\selectfont Que vaut $\sin\left(\dfrac{3\pi}{2}\right)$~?\end{minipage}
\end{minipage}

\vfill
\begin{tcolorbox}[enhanced,nobeforeafter,width=7.1cm,height=2.1cm,valign=center,colback=ctri!6,colframe=ctri,arc=2.5mm,boxrule=0.9pt,left=3mm]{\Large\bfseries\color{ctri}a)}\quad\begin{minipage}[c]{5.3cm}\raggedright\Large $-1$\end{minipage}\end{tcolorbox}\hfill\begin{tcolorbox}[enhanced,nobeforeafter,width=7.1cm,height=2.1cm,valign=center,colback=ctri!6,colframe=ctri,arc=2.5mm,boxrule=0.9pt,left=3mm]{\Large\bfseries\color{ctri}b)}\quad\begin{minipage}[c]{5.3cm}\raggedright\Large $1$\end{minipage}\end{tcolorbox}\par\vspace{2.5mm}
\begin{tcolorbox}[enhanced,nobeforeafter,width=7.1cm,height=2.1cm,valign=center,colback=ctri!6,colframe=ctri,arc=2.5mm,boxrule=0.9pt,left=3mm]{\Large\bfseries\color{ctri}c)}\quad\begin{minipage}[c]{5.3cm}\raggedright\Large $-\dfrac{\sqrt3}{2}$\end{minipage}\end{tcolorbox}\hfill\begin{tcolorbox}[enhanced,nobeforeafter,width=7.1cm,height=2.1cm,valign=center,colback=ctri!6,colframe=ctri,arc=2.5mm,boxrule=0.9pt,left=3mm]{\Large\bfseries\color{ctri}d)}\quad\begin{minipage}[c]{5.3cm}\raggedright\Large $\dfrac{\sqrt3}{2}$\end{minipage}\end{tcolorbox}
\end{tcolorbox}
\newpage
\begin{tcolorbox}[enhanced,width=16cm,height=14cm,colback=white,colframe=ctri,boxrule=2pt,arc=0mm,
left=6mm,right=6mm,top=0mm,bottom=4mm,sharp corners,
overlay={\fill[ctri] (frame.north west) rectangle ([yshift=-1.25cm]frame.north east);
\node[anchor=west,text=white,font=\sffamily\bfseries\LARGE] at ([xshift=5mm,yshift=-0.62cm]frame.north west) {Question 56 D};
\node[anchor=east,text=white,font=\sffamily\itshape\large] at ([xshift=-5mm,yshift=-0.62cm]frame.north east) {Trigonométrie et exponentielle};}]
\vspace*{1.75cm}
\begin{minipage}[t][6.9cm][c]{\linewidth}\LARGE\centering
\begin{minipage}{0.95\linewidth}\raggedright\linespread{1.25}\selectfont Que vaut $\sin\left(\dfrac{5\pi}{6}\right)$~?\end{minipage}
\end{minipage}

\vfill
\begin{tcolorbox}[enhanced,nobeforeafter,width=7.1cm,height=2.1cm,valign=center,colback=ctri!6,colframe=ctri,arc=2.5mm,boxrule=0.9pt,left=3mm]{\Large\bfseries\color{ctri}a)}\quad\begin{minipage}[c]{5.3cm}\raggedright\Large $0$\end{minipage}\end{tcolorbox}\hfill\begin{tcolorbox}[enhanced,nobeforeafter,width=7.1cm,height=2.1cm,valign=center,colback=ctri!6,colframe=ctri,arc=2.5mm,boxrule=0.9pt,left=3mm]{\Large\bfseries\color{ctri}b)}\quad\begin{minipage}[c]{5.3cm}\raggedright\Large $\dfrac12$\end{minipage}\end{tcolorbox}\par\vspace{2.5mm}
\begin{tcolorbox}[enhanced,nobeforeafter,width=7.1cm,height=2.1cm,valign=center,colback=ctri!6,colframe=ctri,arc=2.5mm,boxrule=0.9pt,left=3mm]{\Large\bfseries\color{ctri}c)}\quad\begin{minipage}[c]{5.3cm}\raggedright\Large $-\dfrac12$\end{minipage}\end{tcolorbox}\hfill\begin{tcolorbox}[enhanced,nobeforeafter,width=7.1cm,height=2.1cm,valign=center,colback=ctri!6,colframe=ctri,arc=2.5mm,boxrule=0.9pt,left=3mm]{\Large\bfseries\color{ctri}d)}\quad\begin{minipage}[c]{5.3cm}\raggedright\Large $\dfrac{\sqrt2}{2}$\end{minipage}\end{tcolorbox}
\end{tcolorbox}
\newpage
\begin{tcolorbox}[enhanced,width=16cm,height=14cm,colback=white,colframe=ctri,boxrule=2pt,arc=0mm,
left=6mm,right=6mm,top=0mm,bottom=4mm,sharp corners,
overlay={\fill[ctri] (frame.north west) rectangle ([yshift=-1.25cm]frame.north east);
\node[anchor=west,text=white,font=\sffamily\bfseries\LARGE] at ([xshift=5mm,yshift=-0.62cm]frame.north west) {Question 57 D};
\node[anchor=east,text=white,font=\sffamily\itshape\large] at ([xshift=-5mm,yshift=-0.62cm]frame.north east) {Trigonométrie et exponentielle};}]
\vspace*{1.75cm}
\begin{minipage}[t][6.9cm][c]{\linewidth}\LARGE\centering
\begin{minipage}{0.95\linewidth}\raggedright\linespread{1.25}\selectfont On sait que $\cos x=-0{,}5$. Que vaut $\cos(x+2\pi)$~?\end{minipage}
\end{minipage}

\vfill
\begin{tcolorbox}[enhanced,nobeforeafter,width=7.1cm,height=2.1cm,valign=center,colback=ctri!6,colframe=ctri,arc=2.5mm,boxrule=0.9pt,left=3mm]{\Large\bfseries\color{ctri}a)}\quad\begin{minipage}[c]{5.3cm}\raggedright\Large $1{,}5$\end{minipage}\end{tcolorbox}\hfill\begin{tcolorbox}[enhanced,nobeforeafter,width=7.1cm,height=2.1cm,valign=center,colback=ctri!6,colframe=ctri,arc=2.5mm,boxrule=0.9pt,left=3mm]{\Large\bfseries\color{ctri}b)}\quad\begin{minipage}[c]{5.3cm}\raggedright\Large $-0{,}5$\end{minipage}\end{tcolorbox}\par\vspace{2.5mm}
\begin{tcolorbox}[enhanced,nobeforeafter,width=7.1cm,height=2.1cm,valign=center,colback=ctri!6,colframe=ctri,arc=2.5mm,boxrule=0.9pt,left=3mm]{\Large\bfseries\color{ctri}c)}\quad\begin{minipage}[c]{5.3cm}\raggedright\Large $0{,}75$\end{minipage}\end{tcolorbox}\hfill\begin{tcolorbox}[enhanced,nobeforeafter,width=7.1cm,height=2.1cm,valign=center,colback=ctri!6,colframe=ctri,arc=2.5mm,boxrule=0.9pt,left=3mm]{\Large\bfseries\color{ctri}d)}\quad\begin{minipage}[c]{5.3cm}\raggedright\Large $0{,}5$\end{minipage}\end{tcolorbox}
\end{tcolorbox}
\newpage
\begin{tcolorbox}[enhanced,width=16cm,height=14cm,colback=white,colframe=ctri,boxrule=2pt,arc=0mm,
left=6mm,right=6mm,top=0mm,bottom=4mm,sharp corners,
overlay={\fill[ctri] (frame.north west) rectangle ([yshift=-1.25cm]frame.north east);
\node[anchor=west,text=white,font=\sffamily\bfseries\LARGE] at ([xshift=5mm,yshift=-0.62cm]frame.north west) {Question 58 D};
\node[anchor=east,text=white,font=\sffamily\itshape\large] at ([xshift=-5mm,yshift=-0.62cm]frame.north east) {Trigonométrie et exponentielle};}]
\vspace*{1.75cm}
\begin{minipage}[t][6.9cm][c]{\linewidth}\LARGE\centering
\begin{minipage}{0.95\linewidth}\raggedright\linespread{1.25}\selectfont On sait que $\cos x=0{,}9$. Que vaut $\sin\left(\dfrac{\pi}{2}-x\right)$~?\end{minipage}
\end{minipage}

\vfill
\begin{tcolorbox}[enhanced,nobeforeafter,width=7.1cm,height=2.1cm,valign=center,colback=ctri!6,colframe=ctri,arc=2.5mm,boxrule=0.9pt,left=3mm]{\Large\bfseries\color{ctri}a)}\quad\begin{minipage}[c]{5.3cm}\raggedright\Large $-0{,}9$\end{minipage}\end{tcolorbox}\hfill\begin{tcolorbox}[enhanced,nobeforeafter,width=7.1cm,height=2.1cm,valign=center,colback=ctri!6,colframe=ctri,arc=2.5mm,boxrule=0.9pt,left=3mm]{\Large\bfseries\color{ctri}b)}\quad\begin{minipage}[c]{5.3cm}\raggedright\Large $0{,}9$\end{minipage}\end{tcolorbox}\par\vspace{2.5mm}
\begin{tcolorbox}[enhanced,nobeforeafter,width=7.1cm,height=2.1cm,valign=center,colback=ctri!6,colframe=ctri,arc=2.5mm,boxrule=0.9pt,left=3mm]{\Large\bfseries\color{ctri}c)}\quad\begin{minipage}[c]{5.3cm}\raggedright\Large $0{,}1$\end{minipage}\end{tcolorbox}\hfill\begin{tcolorbox}[enhanced,nobeforeafter,width=7.1cm,height=2.1cm,valign=center,colback=ctri!6,colframe=ctri,arc=2.5mm,boxrule=0.9pt,left=3mm]{\Large\bfseries\color{ctri}d)}\quad\begin{minipage}[c]{5.3cm}\raggedright\Large $0{,}19$\end{minipage}\end{tcolorbox}
\end{tcolorbox}
\newpage
\begin{tcolorbox}[enhanced,width=16cm,height=14cm,colback=white,colframe=ctri,boxrule=2pt,arc=0mm,
left=6mm,right=6mm,top=0mm,bottom=4mm,sharp corners,
overlay={\fill[ctri] (frame.north west) rectangle ([yshift=-1.25cm]frame.north east);
\node[anchor=west,text=white,font=\sffamily\bfseries\LARGE] at ([xshift=5mm,yshift=-0.62cm]frame.north west) {Question 59 D};
\node[anchor=east,text=white,font=\sffamily\itshape\large] at ([xshift=-5mm,yshift=-0.62cm]frame.north east) {Trigonométrie et exponentielle};}]
\vspace*{1.75cm}
\begin{minipage}[t][6.9cm][c]{\linewidth}\LARGE\centering
\begin{minipage}{0.95\linewidth}\raggedright\linespread{1.25}\selectfont Quelle est la mesure principale de l'angle $-\dfrac{10\pi}{3}$~?\end{minipage}
\end{minipage}

\vfill
\begin{tcolorbox}[enhanced,nobeforeafter,width=7.1cm,height=2.1cm,valign=center,colback=ctri!6,colframe=ctri,arc=2.5mm,boxrule=0.9pt,left=3mm]{\Large\bfseries\color{ctri}a)}\quad\begin{minipage}[c]{5.3cm}\raggedright\Large $-\dfrac{2\pi}{3}$\end{minipage}\end{tcolorbox}\hfill\begin{tcolorbox}[enhanced,nobeforeafter,width=7.1cm,height=2.1cm,valign=center,colback=ctri!6,colframe=ctri,arc=2.5mm,boxrule=0.9pt,left=3mm]{\Large\bfseries\color{ctri}b)}\quad\begin{minipage}[c]{5.3cm}\raggedright\Large $-\dfrac{\pi}{3}$\end{minipage}\end{tcolorbox}\par\vspace{2.5mm}
\begin{tcolorbox}[enhanced,nobeforeafter,width=7.1cm,height=2.1cm,valign=center,colback=ctri!6,colframe=ctri,arc=2.5mm,boxrule=0.9pt,left=3mm]{\Large\bfseries\color{ctri}c)}\quad\begin{minipage}[c]{5.3cm}\raggedright\Large $\dfrac{\pi}{3}$\end{minipage}\end{tcolorbox}\hfill\begin{tcolorbox}[enhanced,nobeforeafter,width=7.1cm,height=2.1cm,valign=center,colback=ctri!6,colframe=ctri,arc=2.5mm,boxrule=0.9pt,left=3mm]{\Large\bfseries\color{ctri}d)}\quad\begin{minipage}[c]{5.3cm}\raggedright\Large $\dfrac{2\pi}{3}$\end{minipage}\end{tcolorbox}
\end{tcolorbox}
\newpage
\begin{tcolorbox}[enhanced,width=16cm,height=14cm,colback=white,colframe=ctri,boxrule=2pt,arc=0mm,
left=6mm,right=6mm,top=0mm,bottom=4mm,sharp corners,
overlay={\fill[ctri] (frame.north west) rectangle ([yshift=-1.25cm]frame.north east);
\node[anchor=west,text=white,font=\sffamily\bfseries\LARGE] at ([xshift=5mm,yshift=-0.62cm]frame.north west) {Question 60 D};
\node[anchor=east,text=white,font=\sffamily\itshape\large] at ([xshift=-5mm,yshift=-0.62cm]frame.north east) {Trigonométrie et exponentielle};}]
\vspace*{1.75cm}
\begin{minipage}[t][6.9cm][c]{\linewidth}\LARGE\centering
\begin{minipage}{0.95\linewidth}\raggedright\linespread{1.25}\selectfont Quelle est la mesure principale de l'angle $\dfrac{15\pi}{4}$~?\end{minipage}
\end{minipage}

\vfill
\begin{tcolorbox}[enhanced,nobeforeafter,width=7.1cm,height=2.1cm,valign=center,colback=ctri!6,colframe=ctri,arc=2.5mm,boxrule=0.9pt,left=3mm]{\Large\bfseries\color{ctri}a)}\quad\begin{minipage}[c]{5.3cm}\raggedright\Large $-\dfrac{\pi}{4}$\end{minipage}\end{tcolorbox}\hfill\begin{tcolorbox}[enhanced,nobeforeafter,width=7.1cm,height=2.1cm,valign=center,colback=ctri!6,colframe=ctri,arc=2.5mm,boxrule=0.9pt,left=3mm]{\Large\bfseries\color{ctri}b)}\quad\begin{minipage}[c]{5.3cm}\raggedright\Large $-\dfrac{\pi}{8}$\end{minipage}\end{tcolorbox}\par\vspace{2.5mm}
\begin{tcolorbox}[enhanced,nobeforeafter,width=7.1cm,height=2.1cm,valign=center,colback=ctri!6,colframe=ctri,arc=2.5mm,boxrule=0.9pt,left=3mm]{\Large\bfseries\color{ctri}c)}\quad\begin{minipage}[c]{5.3cm}\raggedright\Large $\dfrac{3\pi}{4}$\end{minipage}\end{tcolorbox}\hfill\begin{tcolorbox}[enhanced,nobeforeafter,width=7.1cm,height=2.1cm,valign=center,colback=ctri!6,colframe=ctri,arc=2.5mm,boxrule=0.9pt,left=3mm]{\Large\bfseries\color{ctri}d)}\quad\begin{minipage}[c]{5.3cm}\raggedright\Large $\dfrac{\pi}{4}$\end{minipage}\end{tcolorbox}
\end{tcolorbox}
\newpage
\begin{tcolorbox}[enhanced,width=16cm,height=14cm,colback=white,colframe=ctri,boxrule=2pt,arc=0mm,
left=6mm,right=6mm,top=0mm,bottom=4mm,sharp corners,
overlay={\fill[ctri] (frame.north west) rectangle ([yshift=-1.25cm]frame.north east);
\node[anchor=west,text=white,font=\sffamily\bfseries\LARGE] at ([xshift=5mm,yshift=-0.62cm]frame.north west) {Question 61 D};
\node[anchor=east,text=white,font=\sffamily\itshape\large] at ([xshift=-5mm,yshift=-0.62cm]frame.north east) {Trigonométrie et exponentielle};}]
\vspace*{1.75cm}
\begin{minipage}[t][6.9cm][c]{\linewidth}\LARGE\centering
\begin{minipage}{0.95\linewidth}\raggedright\linespread{1.25}\selectfont On sait que $x\in[0\,;\,\pi]$ et $\cos x=\dfrac{12}{13}$. Que vaut $\sin x$~?\end{minipage}
\end{minipage}

\vfill
\begin{tcolorbox}[enhanced,nobeforeafter,width=7.1cm,height=2.1cm,valign=center,colback=ctri!6,colframe=ctri,arc=2.5mm,boxrule=0.9pt,left=3mm]{\Large\bfseries\color{ctri}a)}\quad\begin{minipage}[c]{5.3cm}\raggedright\Large $\dfrac{25}{169}$\end{minipage}\end{tcolorbox}\hfill\begin{tcolorbox}[enhanced,nobeforeafter,width=7.1cm,height=2.1cm,valign=center,colback=ctri!6,colframe=ctri,arc=2.5mm,boxrule=0.9pt,left=3mm]{\Large\bfseries\color{ctri}b)}\quad\begin{minipage}[c]{5.3cm}\raggedright\Large $\dfrac{1}{13}$\end{minipage}\end{tcolorbox}\par\vspace{2.5mm}
\begin{tcolorbox}[enhanced,nobeforeafter,width=7.1cm,height=2.1cm,valign=center,colback=ctri!6,colframe=ctri,arc=2.5mm,boxrule=0.9pt,left=3mm]{\Large\bfseries\color{ctri}c)}\quad\begin{minipage}[c]{5.3cm}\raggedright\Large $-\dfrac{5}{13}$\end{minipage}\end{tcolorbox}\hfill\begin{tcolorbox}[enhanced,nobeforeafter,width=7.1cm,height=2.1cm,valign=center,colback=ctri!6,colframe=ctri,arc=2.5mm,boxrule=0.9pt,left=3mm]{\Large\bfseries\color{ctri}d)}\quad\begin{minipage}[c]{5.3cm}\raggedright\Large $\dfrac{5}{13}$\end{minipage}\end{tcolorbox}
\end{tcolorbox}
\newpage
\begin{tcolorbox}[enhanced,width=16cm,height=14cm,colback=white,colframe=ctri,boxrule=2pt,arc=0mm,
left=6mm,right=6mm,top=0mm,bottom=4mm,sharp corners,
overlay={\fill[ctri] (frame.north west) rectangle ([yshift=-1.25cm]frame.north east);
\node[anchor=west,text=white,font=\sffamily\bfseries\LARGE] at ([xshift=5mm,yshift=-0.62cm]frame.north west) {Question 62 D};
\node[anchor=east,text=white,font=\sffamily\itshape\large] at ([xshift=-5mm,yshift=-0.62cm]frame.north east) {Trigonométrie et exponentielle};}]
\vspace*{1.75cm}
\begin{minipage}[t][6.9cm][c]{\linewidth}\LARGE\centering
\begin{minipage}{0.95\linewidth}\raggedright\linespread{1.25}\selectfont On sait que $x\in[0\,;\,\pi]$ et $\cos x=-\dfrac{20}{29}$. Que vaut $\sin x$~?\end{minipage}
\end{minipage}

\vfill
\begin{tcolorbox}[enhanced,nobeforeafter,width=7.1cm,height=2.1cm,valign=center,colback=ctri!6,colframe=ctri,arc=2.5mm,boxrule=0.9pt,left=3mm]{\Large\bfseries\color{ctri}a)}\quad\begin{minipage}[c]{5.3cm}\raggedright\Large $-\dfrac{21}{29}$\end{minipage}\end{tcolorbox}\hfill\begin{tcolorbox}[enhanced,nobeforeafter,width=7.1cm,height=2.1cm,valign=center,colback=ctri!6,colframe=ctri,arc=2.5mm,boxrule=0.9pt,left=3mm]{\Large\bfseries\color{ctri}b)}\quad\begin{minipage}[c]{5.3cm}\raggedright\Large $\dfrac{21}{29}$\end{minipage}\end{tcolorbox}\par\vspace{2.5mm}
\begin{tcolorbox}[enhanced,nobeforeafter,width=7.1cm,height=2.1cm,valign=center,colback=ctri!6,colframe=ctri,arc=2.5mm,boxrule=0.9pt,left=3mm]{\Large\bfseries\color{ctri}c)}\quad\begin{minipage}[c]{5.3cm}\raggedright\Large $\dfrac{441}{841}$\end{minipage}\end{tcolorbox}\hfill\begin{tcolorbox}[enhanced,nobeforeafter,width=7.1cm,height=2.1cm,valign=center,colback=ctri!6,colframe=ctri,arc=2.5mm,boxrule=0.9pt,left=3mm]{\Large\bfseries\color{ctri}d)}\quad\begin{minipage}[c]{5.3cm}\raggedright\Large $\dfrac{49}{29}$\end{minipage}\end{tcolorbox}
\end{tcolorbox}
\newpage
\begin{tcolorbox}[enhanced,width=16cm,height=14cm,colback=white,colframe=ctri,boxrule=2pt,arc=0mm,
left=6mm,right=6mm,top=0mm,bottom=4mm,sharp corners,
overlay={\fill[ctri] (frame.north west) rectangle ([yshift=-1.25cm]frame.north east);
\node[anchor=west,text=white,font=\sffamily\bfseries\LARGE] at ([xshift=5mm,yshift=-0.62cm]frame.north west) {Question 63 D};
\node[anchor=east,text=white,font=\sffamily\itshape\large] at ([xshift=-5mm,yshift=-0.62cm]frame.north east) {Trigonométrie et exponentielle};}]
\vspace*{1.75cm}
\begin{minipage}[t][6.9cm][c]{\linewidth}\LARGE\centering
\begin{minipage}{0.95\linewidth}\raggedright\linespread{1.25}\selectfont Simplifier $\left(e^{5}\right)^{2}\times e^{-4}$.\end{minipage}
\end{minipage}

\vfill
\begin{tcolorbox}[enhanced,nobeforeafter,width=7.1cm,height=2.1cm,valign=center,colback=ctri!6,colframe=ctri,arc=2.5mm,boxrule=0.9pt,left=3mm]{\Large\bfseries\color{ctri}a)}\quad\begin{minipage}[c]{5.3cm}\raggedright\Large $e^{-10}$\end{minipage}\end{tcolorbox}\hfill\begin{tcolorbox}[enhanced,nobeforeafter,width=7.1cm,height=2.1cm,valign=center,colback=ctri!6,colframe=ctri,arc=2.5mm,boxrule=0.9pt,left=3mm]{\Large\bfseries\color{ctri}b)}\quad\begin{minipage}[c]{5.3cm}\raggedright\Large $e^{3}$\end{minipage}\end{tcolorbox}\par\vspace{2.5mm}
\begin{tcolorbox}[enhanced,nobeforeafter,width=7.1cm,height=2.1cm,valign=center,colback=ctri!6,colframe=ctri,arc=2.5mm,boxrule=0.9pt,left=3mm]{\Large\bfseries\color{ctri}c)}\quad\begin{minipage}[c]{5.3cm}\raggedright\Large $e^{6}$\end{minipage}\end{tcolorbox}\hfill\begin{tcolorbox}[enhanced,nobeforeafter,width=7.1cm,height=2.1cm,valign=center,colback=ctri!6,colframe=ctri,arc=2.5mm,boxrule=0.9pt,left=3mm]{\Large\bfseries\color{ctri}d)}\quad\begin{minipage}[c]{5.3cm}\raggedright\Large $e^{21}$\end{minipage}\end{tcolorbox}
\end{tcolorbox}
\newpage
\begin{tcolorbox}[enhanced,width=16cm,height=14cm,colback=white,colframe=ctri,boxrule=2pt,arc=0mm,
left=6mm,right=6mm,top=0mm,bottom=4mm,sharp corners,
overlay={\fill[ctri] (frame.north west) rectangle ([yshift=-1.25cm]frame.north east);
\node[anchor=west,text=white,font=\sffamily\bfseries\LARGE] at ([xshift=5mm,yshift=-0.62cm]frame.north west) {Question 64 D};
\node[anchor=east,text=white,font=\sffamily\itshape\large] at ([xshift=-5mm,yshift=-0.62cm]frame.north east) {Trigonométrie et exponentielle};}]
\vspace*{1.75cm}
\begin{minipage}[t][6.9cm][c]{\linewidth}\LARGE\centering
\begin{minipage}{0.95\linewidth}\raggedright\linespread{1.25}\selectfont Simplifier $\left(e^{3}\right)^{3}\times e^{-5}$.\end{minipage}
\end{minipage}

\vfill
\begin{tcolorbox}[enhanced,nobeforeafter,width=7.1cm,height=2.1cm,valign=center,colback=ctri!6,colframe=ctri,arc=2.5mm,boxrule=0.9pt,left=3mm]{\Large\bfseries\color{ctri}a)}\quad\begin{minipage}[c]{5.3cm}\raggedright\Large $e^{-6}$\end{minipage}\end{tcolorbox}\hfill\begin{tcolorbox}[enhanced,nobeforeafter,width=7.1cm,height=2.1cm,valign=center,colback=ctri!6,colframe=ctri,arc=2.5mm,boxrule=0.9pt,left=3mm]{\Large\bfseries\color{ctri}b)}\quad\begin{minipage}[c]{5.3cm}\raggedright\Large $e^{22}$\end{minipage}\end{tcolorbox}\par\vspace{2.5mm}
\begin{tcolorbox}[enhanced,nobeforeafter,width=7.1cm,height=2.1cm,valign=center,colback=ctri!6,colframe=ctri,arc=2.5mm,boxrule=0.9pt,left=3mm]{\Large\bfseries\color{ctri}c)}\quad\begin{minipage}[c]{5.3cm}\raggedright\Large $e^{4}$\end{minipage}\end{tcolorbox}\hfill\begin{tcolorbox}[enhanced,nobeforeafter,width=7.1cm,height=2.1cm,valign=center,colback=ctri!6,colframe=ctri,arc=2.5mm,boxrule=0.9pt,left=3mm]{\Large\bfseries\color{ctri}d)}\quad\begin{minipage}[c]{5.3cm}\raggedright\Large $e$\end{minipage}\end{tcolorbox}
\end{tcolorbox}
\newpage
\begin{tcolorbox}[enhanced,width=16cm,height=14cm,colback=white,colframe=cpsc,boxrule=2pt,arc=0mm,
left=6mm,right=6mm,top=0mm,bottom=4mm,sharp corners,
overlay={\fill[cpsc] (frame.north west) rectangle ([yshift=-1.25cm]frame.north east);
\node[anchor=west,text=white,font=\sffamily\bfseries\LARGE] at ([xshift=5mm,yshift=-0.62cm]frame.north west) {Question 65 D};
\node[anchor=east,text=white,font=\sffamily\itshape\large] at ([xshift=-5mm,yshift=-0.62cm]frame.north east) {Produit scalaire et géométrie repérée};}]
\vspace*{1.75cm}
\begin{minipage}[t][6.9cm][c]{\linewidth}\LARGE\centering
\begin{minipage}{0.95\linewidth}\raggedright\linespread{1.25}\selectfont Dans un repère orthonormé, que vaut $\vec u\cdot\vec v$ pour $\vec u\begin{pmatrix}0\\4\end{pmatrix}$ et $\vec v\begin{pmatrix}7\\-2\end{pmatrix}$~?\end{minipage}
\end{minipage}

\vfill
\begin{tcolorbox}[enhanced,nobeforeafter,width=7.1cm,height=2.1cm,valign=center,colback=cpsc!6,colframe=cpsc,arc=2.5mm,boxrule=0.9pt,left=3mm]{\Large\bfseries\color{cpsc}a)}\quad\begin{minipage}[c]{5.3cm}\raggedright\Large $-8$\end{minipage}\end{tcolorbox}\hfill\begin{tcolorbox}[enhanced,nobeforeafter,width=7.1cm,height=2.1cm,valign=center,colback=cpsc!6,colframe=cpsc,arc=2.5mm,boxrule=0.9pt,left=3mm]{\Large\bfseries\color{cpsc}b)}\quad\begin{minipage}[c]{5.3cm}\raggedright\Large $8$\end{minipage}\end{tcolorbox}\par\vspace{2.5mm}
\begin{tcolorbox}[enhanced,nobeforeafter,width=7.1cm,height=2.1cm,valign=center,colback=cpsc!6,colframe=cpsc,arc=2.5mm,boxrule=0.9pt,left=3mm]{\Large\bfseries\color{cpsc}c)}\quad\begin{minipage}[c]{5.3cm}\raggedright\Large $-28$\end{minipage}\end{tcolorbox}\hfill\begin{tcolorbox}[enhanced,nobeforeafter,width=7.1cm,height=2.1cm,valign=center,colback=cpsc!6,colframe=cpsc,arc=2.5mm,boxrule=0.9pt,left=3mm]{\Large\bfseries\color{cpsc}d)}\quad\begin{minipage}[c]{5.3cm}\raggedright\Large $28$\end{minipage}\end{tcolorbox}
\end{tcolorbox}
\newpage
\begin{tcolorbox}[enhanced,width=16cm,height=14cm,colback=white,colframe=cpsc,boxrule=2pt,arc=0mm,
left=6mm,right=6mm,top=0mm,bottom=4mm,sharp corners,
overlay={\fill[cpsc] (frame.north west) rectangle ([yshift=-1.25cm]frame.north east);
\node[anchor=west,text=white,font=\sffamily\bfseries\LARGE] at ([xshift=5mm,yshift=-0.62cm]frame.north west) {Question 66 D};
\node[anchor=east,text=white,font=\sffamily\itshape\large] at ([xshift=-5mm,yshift=-0.62cm]frame.north east) {Produit scalaire et géométrie repérée};}]
\vspace*{1.75cm}
\begin{minipage}[t][6.9cm][c]{\linewidth}\LARGE\centering
\begin{minipage}{0.95\linewidth}\raggedright\linespread{1.25}\selectfont Dans un repère orthonormé, que vaut $\vec u\cdot\vec v$ pour $\vec u\begin{pmatrix}6\\-1\end{pmatrix}$ et $\vec v\begin{pmatrix}2\\4\end{pmatrix}$~?\end{minipage}
\end{minipage}

\vfill
\begin{tcolorbox}[enhanced,nobeforeafter,width=7.1cm,height=2.1cm,valign=center,colback=cpsc!6,colframe=cpsc,arc=2.5mm,boxrule=0.9pt,left=3mm]{\Large\bfseries\color{cpsc}a)}\quad\begin{minipage}[c]{5.3cm}\raggedright\Large $16$\end{minipage}\end{tcolorbox}\hfill\begin{tcolorbox}[enhanced,nobeforeafter,width=7.1cm,height=2.1cm,valign=center,colback=cpsc!6,colframe=cpsc,arc=2.5mm,boxrule=0.9pt,left=3mm]{\Large\bfseries\color{cpsc}b)}\quad\begin{minipage}[c]{5.3cm}\raggedright\Large $8$\end{minipage}\end{tcolorbox}\par\vspace{2.5mm}
\begin{tcolorbox}[enhanced,nobeforeafter,width=7.1cm,height=2.1cm,valign=center,colback=cpsc!6,colframe=cpsc,arc=2.5mm,boxrule=0.9pt,left=3mm]{\Large\bfseries\color{cpsc}c)}\quad\begin{minipage}[c]{5.3cm}\raggedright\Large $22$\end{minipage}\end{tcolorbox}\hfill\begin{tcolorbox}[enhanced,nobeforeafter,width=7.1cm,height=2.1cm,valign=center,colback=cpsc!6,colframe=cpsc,arc=2.5mm,boxrule=0.9pt,left=3mm]{\Large\bfseries\color{cpsc}d)}\quad\begin{minipage}[c]{5.3cm}\raggedright\Large $26$\end{minipage}\end{tcolorbox}
\end{tcolorbox}
\newpage
\begin{tcolorbox}[enhanced,width=16cm,height=14cm,colback=white,colframe=cpsc,boxrule=2pt,arc=0mm,
left=6mm,right=6mm,top=0mm,bottom=4mm,sharp corners,
overlay={\fill[cpsc] (frame.north west) rectangle ([yshift=-1.25cm]frame.north east);
\node[anchor=west,text=white,font=\sffamily\bfseries\LARGE] at ([xshift=5mm,yshift=-0.62cm]frame.north west) {Question 67 D};
\node[anchor=east,text=white,font=\sffamily\itshape\large] at ([xshift=-5mm,yshift=-0.62cm]frame.north east) {Produit scalaire et géométrie repérée};}]
\vspace*{1.75cm}
\begin{minipage}[t][6.9cm][c]{\linewidth}\LARGE\centering
\begin{minipage}{0.95\linewidth}\raggedright\linespread{1.25}\selectfont $\|\vec u\|=3$, $\|\vec v\|=6$ et l'angle entre $\vec u$ et $\vec v$ mesure $120^\circ$. Que vaut $\vec u\cdot\vec v$~?\end{minipage}
\end{minipage}

\vfill
\begin{tcolorbox}[enhanced,nobeforeafter,width=7.1cm,height=2.1cm,valign=center,colback=cpsc!6,colframe=cpsc,arc=2.5mm,boxrule=0.9pt,left=3mm]{\Large\bfseries\color{cpsc}a)}\quad\begin{minipage}[c]{5.3cm}\raggedright\Large $-8$\end{minipage}\end{tcolorbox}\hfill\begin{tcolorbox}[enhanced,nobeforeafter,width=7.1cm,height=2.1cm,valign=center,colback=cpsc!6,colframe=cpsc,arc=2.5mm,boxrule=0.9pt,left=3mm]{\Large\bfseries\color{cpsc}b)}\quad\begin{minipage}[c]{5.3cm}\raggedright\Large $18$\end{minipage}\end{tcolorbox}\par\vspace{2.5mm}
\begin{tcolorbox}[enhanced,nobeforeafter,width=7.1cm,height=2.1cm,valign=center,colback=cpsc!6,colframe=cpsc,arc=2.5mm,boxrule=0.9pt,left=3mm]{\Large\bfseries\color{cpsc}c)}\quad\begin{minipage}[c]{5.3cm}\raggedright\Large $-9$\end{minipage}\end{tcolorbox}\hfill\begin{tcolorbox}[enhanced,nobeforeafter,width=7.1cm,height=2.1cm,valign=center,colback=cpsc!6,colframe=cpsc,arc=2.5mm,boxrule=0.9pt,left=3mm]{\Large\bfseries\color{cpsc}d)}\quad\begin{minipage}[c]{5.3cm}\raggedright\Large $9$\end{minipage}\end{tcolorbox}
\end{tcolorbox}
\newpage
\begin{tcolorbox}[enhanced,width=16cm,height=14cm,colback=white,colframe=cpsc,boxrule=2pt,arc=0mm,
left=6mm,right=6mm,top=0mm,bottom=4mm,sharp corners,
overlay={\fill[cpsc] (frame.north west) rectangle ([yshift=-1.25cm]frame.north east);
\node[anchor=west,text=white,font=\sffamily\bfseries\LARGE] at ([xshift=5mm,yshift=-0.62cm]frame.north west) {Question 68 D};
\node[anchor=east,text=white,font=\sffamily\itshape\large] at ([xshift=-5mm,yshift=-0.62cm]frame.north east) {Produit scalaire et géométrie repérée};}]
\vspace*{1.75cm}
\begin{minipage}[t][6.9cm][c]{\linewidth}\LARGE\centering
\begin{minipage}{0.95\linewidth}\raggedright\linespread{1.25}\selectfont $\|\vec u\|=7$, $\|\vec v\|=2$ et l'angle entre $\vec u$ et $\vec v$ mesure $60^\circ$. Que vaut $\vec u\cdot\vec v$~?\end{minipage}
\end{minipage}

\vfill
\begin{tcolorbox}[enhanced,nobeforeafter,width=7.1cm,height=2.1cm,valign=center,colback=cpsc!6,colframe=cpsc,arc=2.5mm,boxrule=0.9pt,left=3mm]{\Large\bfseries\color{cpsc}a)}\quad\begin{minipage}[c]{5.3cm}\raggedright\Large $-7$\end{minipage}\end{tcolorbox}\hfill\begin{tcolorbox}[enhanced,nobeforeafter,width=7.1cm,height=2.1cm,valign=center,colback=cpsc!6,colframe=cpsc,arc=2.5mm,boxrule=0.9pt,left=3mm]{\Large\bfseries\color{cpsc}b)}\quad\begin{minipage}[c]{5.3cm}\raggedright\Large $7$\end{minipage}\end{tcolorbox}\par\vspace{2.5mm}
\begin{tcolorbox}[enhanced,nobeforeafter,width=7.1cm,height=2.1cm,valign=center,colback=cpsc!6,colframe=cpsc,arc=2.5mm,boxrule=0.9pt,left=3mm]{\Large\bfseries\color{cpsc}c)}\quad\begin{minipage}[c]{5.3cm}\raggedright\Large $14$\end{minipage}\end{tcolorbox}\hfill\begin{tcolorbox}[enhanced,nobeforeafter,width=7.1cm,height=2.1cm,valign=center,colback=cpsc!6,colframe=cpsc,arc=2.5mm,boxrule=0.9pt,left=3mm]{\Large\bfseries\color{cpsc}d)}\quad\begin{minipage}[c]{5.3cm}\raggedright\Large $9$\end{minipage}\end{tcolorbox}
\end{tcolorbox}
\newpage
\begin{tcolorbox}[enhanced,width=16cm,height=14cm,colback=white,colframe=cpsc,boxrule=2pt,arc=0mm,
left=6mm,right=6mm,top=0mm,bottom=4mm,sharp corners,
overlay={\fill[cpsc] (frame.north west) rectangle ([yshift=-1.25cm]frame.north east);
\node[anchor=west,text=white,font=\sffamily\bfseries\LARGE] at ([xshift=5mm,yshift=-0.62cm]frame.north west) {Question 69 D};
\node[anchor=east,text=white,font=\sffamily\itshape\large] at ([xshift=-5mm,yshift=-0.62cm]frame.north east) {Produit scalaire et géométrie repérée};}]
\vspace*{1.75cm}
\begin{minipage}[t][6.9cm][c]{\linewidth}\LARGE\centering
\begin{minipage}{0.95\linewidth}\raggedright\linespread{1.25}\selectfont Pour quelle valeur de $m$ les vecteurs $\vec u\begin{pmatrix}3\\m\end{pmatrix}$ et $\vec v\begin{pmatrix}5\\-3\end{pmatrix}$ sont-ils orthogonaux~?\end{minipage}
\end{minipage}

\vfill
\begin{tcolorbox}[enhanced,nobeforeafter,width=7.1cm,height=2.1cm,valign=center,colback=cpsc!6,colframe=cpsc,arc=2.5mm,boxrule=0.9pt,left=3mm]{\Large\bfseries\color{cpsc}a)}\quad\begin{minipage}[c]{5.3cm}\raggedright\Large $m=-15$\end{minipage}\end{tcolorbox}\hfill\begin{tcolorbox}[enhanced,nobeforeafter,width=7.1cm,height=2.1cm,valign=center,colback=cpsc!6,colframe=cpsc,arc=2.5mm,boxrule=0.9pt,left=3mm]{\Large\bfseries\color{cpsc}b)}\quad\begin{minipage}[c]{5.3cm}\raggedright\Large $m=5$\end{minipage}\end{tcolorbox}\par\vspace{2.5mm}
\begin{tcolorbox}[enhanced,nobeforeafter,width=7.1cm,height=2.1cm,valign=center,colback=cpsc!6,colframe=cpsc,arc=2.5mm,boxrule=0.9pt,left=3mm]{\Large\bfseries\color{cpsc}c)}\quad\begin{minipage}[c]{5.3cm}\raggedright\Large $m=\dfrac{1}{5}$\end{minipage}\end{tcolorbox}\hfill\begin{tcolorbox}[enhanced,nobeforeafter,width=7.1cm,height=2.1cm,valign=center,colback=cpsc!6,colframe=cpsc,arc=2.5mm,boxrule=0.9pt,left=3mm]{\Large\bfseries\color{cpsc}d)}\quad\begin{minipage}[c]{5.3cm}\raggedright\Large $m=-5$\end{minipage}\end{tcolorbox}
\end{tcolorbox}
\newpage
\begin{tcolorbox}[enhanced,width=16cm,height=14cm,colback=white,colframe=cpsc,boxrule=2pt,arc=0mm,
left=6mm,right=6mm,top=0mm,bottom=4mm,sharp corners,
overlay={\fill[cpsc] (frame.north west) rectangle ([yshift=-1.25cm]frame.north east);
\node[anchor=west,text=white,font=\sffamily\bfseries\LARGE] at ([xshift=5mm,yshift=-0.62cm]frame.north west) {Question 70 D};
\node[anchor=east,text=white,font=\sffamily\itshape\large] at ([xshift=-5mm,yshift=-0.62cm]frame.north east) {Produit scalaire et géométrie repérée};}]
\vspace*{1.75cm}
\begin{minipage}[t][6.9cm][c]{\linewidth}\LARGE\centering
\begin{minipage}{0.95\linewidth}\raggedright\linespread{1.25}\selectfont Pour quelle valeur de $m$ les vecteurs $\vec u\begin{pmatrix}-2\\m\end{pmatrix}$ et $\vec v\begin{pmatrix}-4\\8\end{pmatrix}$ sont-ils orthogonaux~?\end{minipage}
\end{minipage}

\vfill
\begin{tcolorbox}[enhanced,nobeforeafter,width=7.1cm,height=2.1cm,valign=center,colback=cpsc!6,colframe=cpsc,arc=2.5mm,boxrule=0.9pt,left=3mm]{\Large\bfseries\color{cpsc}a)}\quad\begin{minipage}[c]{5.3cm}\raggedright\Large $m=-8$\end{minipage}\end{tcolorbox}\hfill\begin{tcolorbox}[enhanced,nobeforeafter,width=7.1cm,height=2.1cm,valign=center,colback=cpsc!6,colframe=cpsc,arc=2.5mm,boxrule=0.9pt,left=3mm]{\Large\bfseries\color{cpsc}b)}\quad\begin{minipage}[c]{5.3cm}\raggedright\Large $m=-1$\end{minipage}\end{tcolorbox}\par\vspace{2.5mm}
\begin{tcolorbox}[enhanced,nobeforeafter,width=7.1cm,height=2.1cm,valign=center,colback=cpsc!6,colframe=cpsc,arc=2.5mm,boxrule=0.9pt,left=3mm]{\Large\bfseries\color{cpsc}c)}\quad\begin{minipage}[c]{5.3cm}\raggedright\Large $m=1$\end{minipage}\end{tcolorbox}\hfill\begin{tcolorbox}[enhanced,nobeforeafter,width=7.1cm,height=2.1cm,valign=center,colback=cpsc!6,colframe=cpsc,arc=2.5mm,boxrule=0.9pt,left=3mm]{\Large\bfseries\color{cpsc}d)}\quad\begin{minipage}[c]{5.3cm}\raggedright\Large $m=0$\end{minipage}\end{tcolorbox}
\end{tcolorbox}
\newpage
\begin{tcolorbox}[enhanced,width=16cm,height=14cm,colback=white,colframe=cpsc,boxrule=2pt,arc=0mm,
left=6mm,right=6mm,top=0mm,bottom=4mm,sharp corners,
overlay={\fill[cpsc] (frame.north west) rectangle ([yshift=-1.25cm]frame.north east);
\node[anchor=west,text=white,font=\sffamily\bfseries\LARGE] at ([xshift=5mm,yshift=-0.62cm]frame.north west) {Question 71 D};
\node[anchor=east,text=white,font=\sffamily\itshape\large] at ([xshift=-5mm,yshift=-0.62cm]frame.north east) {Produit scalaire et géométrie repérée};}]
\vspace*{1.75cm}
\begin{minipage}[t][6.9cm][c]{\linewidth}\LARGE\centering
\begin{minipage}{0.95\linewidth}\raggedright\linespread{1.25}\selectfont Dans le triangle $ABC$, $AB=1$, $AC=1$ et $\widehat{BAC}=60^\circ$. Combien mesure $BC$~?\end{minipage}
\end{minipage}

\vfill
\begin{tcolorbox}[enhanced,nobeforeafter,width=7.1cm,height=2.1cm,valign=center,colback=cpsc!6,colframe=cpsc,arc=2.5mm,boxrule=0.9pt,left=3mm]{\Large\bfseries\color{cpsc}a)}\quad\begin{minipage}[c]{5.3cm}\raggedright\Large $\sqrt{2}$\end{minipage}\end{tcolorbox}\hfill\begin{tcolorbox}[enhanced,nobeforeafter,width=7.1cm,height=2.1cm,valign=center,colback=cpsc!6,colframe=cpsc,arc=2.5mm,boxrule=0.9pt,left=3mm]{\Large\bfseries\color{cpsc}b)}\quad\begin{minipage}[c]{5.3cm}\raggedright\Large $2$\end{minipage}\end{tcolorbox}\par\vspace{2.5mm}
\begin{tcolorbox}[enhanced,nobeforeafter,width=7.1cm,height=2.1cm,valign=center,colback=cpsc!6,colframe=cpsc,arc=2.5mm,boxrule=0.9pt,left=3mm]{\Large\bfseries\color{cpsc}c)}\quad\begin{minipage}[c]{5.3cm}\raggedright\Large $1$\end{minipage}\end{tcolorbox}\hfill\begin{tcolorbox}[enhanced,nobeforeafter,width=7.1cm,height=2.1cm,valign=center,colback=cpsc!6,colframe=cpsc,arc=2.5mm,boxrule=0.9pt,left=3mm]{\Large\bfseries\color{cpsc}d)}\quad\begin{minipage}[c]{5.3cm}\raggedright\Large $\sqrt{3}$\end{minipage}\end{tcolorbox}
\end{tcolorbox}
\newpage
\begin{tcolorbox}[enhanced,width=16cm,height=14cm,colback=white,colframe=cpsc,boxrule=2pt,arc=0mm,
left=6mm,right=6mm,top=0mm,bottom=4mm,sharp corners,
overlay={\fill[cpsc] (frame.north west) rectangle ([yshift=-1.25cm]frame.north east);
\node[anchor=west,text=white,font=\sffamily\bfseries\LARGE] at ([xshift=5mm,yshift=-0.62cm]frame.north west) {Question 72 D};
\node[anchor=east,text=white,font=\sffamily\itshape\large] at ([xshift=-5mm,yshift=-0.62cm]frame.north east) {Produit scalaire et géométrie repérée};}]
\vspace*{1.75cm}
\begin{minipage}[t][6.9cm][c]{\linewidth}\LARGE\centering
\begin{minipage}{0.95\linewidth}\raggedright\linespread{1.25}\selectfont Dans le triangle $ABC$, $AB=7$, $AC=3$ et $\widehat{BAC}=60^\circ$. Combien mesure $BC$~?\end{minipage}
\end{minipage}

\vfill
\begin{tcolorbox}[enhanced,nobeforeafter,width=7.1cm,height=2.1cm,valign=center,colback=cpsc!6,colframe=cpsc,arc=2.5mm,boxrule=0.9pt,left=3mm]{\Large\bfseries\color{cpsc}a)}\quad\begin{minipage}[c]{5.3cm}\raggedright\Large $\sqrt{58}$\end{minipage}\end{tcolorbox}\hfill\begin{tcolorbox}[enhanced,nobeforeafter,width=7.1cm,height=2.1cm,valign=center,colback=cpsc!6,colframe=cpsc,arc=2.5mm,boxrule=0.9pt,left=3mm]{\Large\bfseries\color{cpsc}b)}\quad\begin{minipage}[c]{5.3cm}\raggedright\Large $\sqrt{37}$\end{minipage}\end{tcolorbox}\par\vspace{2.5mm}
\begin{tcolorbox}[enhanced,nobeforeafter,width=7.1cm,height=2.1cm,valign=center,colback=cpsc!6,colframe=cpsc,arc=2.5mm,boxrule=0.9pt,left=3mm]{\Large\bfseries\color{cpsc}c)}\quad\begin{minipage}[c]{5.3cm}\raggedright\Large $\sqrt{79}$\end{minipage}\end{tcolorbox}\hfill\begin{tcolorbox}[enhanced,nobeforeafter,width=7.1cm,height=2.1cm,valign=center,colback=cpsc!6,colframe=cpsc,arc=2.5mm,boxrule=0.9pt,left=3mm]{\Large\bfseries\color{cpsc}d)}\quad\begin{minipage}[c]{5.3cm}\raggedright\Large $37$\end{minipage}\end{tcolorbox}
\end{tcolorbox}
\newpage
\begin{tcolorbox}[enhanced,width=16cm,height=14cm,colback=white,colframe=cpsc,boxrule=2pt,arc=0mm,
left=6mm,right=6mm,top=0mm,bottom=4mm,sharp corners,
overlay={\fill[cpsc] (frame.north west) rectangle ([yshift=-1.25cm]frame.north east);
\node[anchor=west,text=white,font=\sffamily\bfseries\LARGE] at ([xshift=5mm,yshift=-0.62cm]frame.north west) {Question 73 D};
\node[anchor=east,text=white,font=\sffamily\itshape\large] at ([xshift=-5mm,yshift=-0.62cm]frame.north east) {Produit scalaire et géométrie repérée};}]
\vspace*{1.75cm}
\begin{minipage}[t][6.9cm][c]{\linewidth}\LARGE\centering
\begin{minipage}{0.95\linewidth}\raggedright\linespread{1.25}\selectfont Lequel de ces vecteurs est un vecteur normal à la droite d'équation $-2x-5y+4=0$~?\end{minipage}
\end{minipage}

\vfill
\begin{tcolorbox}[enhanced,nobeforeafter,width=7.1cm,height=2.1cm,valign=center,colback=cpsc!6,colframe=cpsc,arc=2.5mm,boxrule=0.9pt,left=3mm]{\Large\bfseries\color{cpsc}a)}\quad\begin{minipage}[c]{5.3cm}\raggedright\large $\begin{pmatrix}-2\\4\end{pmatrix}$\end{minipage}\end{tcolorbox}\hfill\begin{tcolorbox}[enhanced,nobeforeafter,width=7.1cm,height=2.1cm,valign=center,colback=cpsc!6,colframe=cpsc,arc=2.5mm,boxrule=0.9pt,left=3mm]{\Large\bfseries\color{cpsc}b)}\quad\begin{minipage}[c]{5.3cm}\raggedright\large $\begin{pmatrix}-5\\-2\end{pmatrix}$\end{minipage}\end{tcolorbox}\par\vspace{2.5mm}
\begin{tcolorbox}[enhanced,nobeforeafter,width=7.1cm,height=2.1cm,valign=center,colback=cpsc!6,colframe=cpsc,arc=2.5mm,boxrule=0.9pt,left=3mm]{\Large\bfseries\color{cpsc}c)}\quad\begin{minipage}[c]{5.3cm}\raggedright\large $\begin{pmatrix}-2\\-5\end{pmatrix}$\end{minipage}\end{tcolorbox}\hfill\begin{tcolorbox}[enhanced,nobeforeafter,width=7.1cm,height=2.1cm,valign=center,colback=cpsc!6,colframe=cpsc,arc=2.5mm,boxrule=0.9pt,left=3mm]{\Large\bfseries\color{cpsc}d)}\quad\begin{minipage}[c]{5.3cm}\raggedright\large $\begin{pmatrix}5\\-2\end{pmatrix}$\end{minipage}\end{tcolorbox}
\end{tcolorbox}
\newpage
\begin{tcolorbox}[enhanced,width=16cm,height=14cm,colback=white,colframe=cpsc,boxrule=2pt,arc=0mm,
left=6mm,right=6mm,top=0mm,bottom=4mm,sharp corners,
overlay={\fill[cpsc] (frame.north west) rectangle ([yshift=-1.25cm]frame.north east);
\node[anchor=west,text=white,font=\sffamily\bfseries\LARGE] at ([xshift=5mm,yshift=-0.62cm]frame.north west) {Question 74 D};
\node[anchor=east,text=white,font=\sffamily\itshape\large] at ([xshift=-5mm,yshift=-0.62cm]frame.north east) {Produit scalaire et géométrie repérée};}]
\vspace*{1.75cm}
\begin{minipage}[t][6.9cm][c]{\linewidth}\LARGE\centering
\begin{minipage}{0.95\linewidth}\raggedright\linespread{1.25}\selectfont Lequel de ces vecteurs est un vecteur normal à la droite d'équation $6x+y-3=0$~?\end{minipage}
\end{minipage}

\vfill
\begin{tcolorbox}[enhanced,nobeforeafter,width=7.1cm,height=2.1cm,valign=center,colback=cpsc!6,colframe=cpsc,arc=2.5mm,boxrule=0.9pt,left=3mm]{\Large\bfseries\color{cpsc}a)}\quad\begin{minipage}[c]{5.3cm}\raggedright\large $\begin{pmatrix}6\\1\end{pmatrix}$\end{minipage}\end{tcolorbox}\hfill\begin{tcolorbox}[enhanced,nobeforeafter,width=7.1cm,height=2.1cm,valign=center,colback=cpsc!6,colframe=cpsc,arc=2.5mm,boxrule=0.9pt,left=3mm]{\Large\bfseries\color{cpsc}b)}\quad\begin{minipage}[c]{5.3cm}\raggedright\large $\begin{pmatrix}1\\6\end{pmatrix}$\end{minipage}\end{tcolorbox}\par\vspace{2.5mm}
\begin{tcolorbox}[enhanced,nobeforeafter,width=7.1cm,height=2.1cm,valign=center,colback=cpsc!6,colframe=cpsc,arc=2.5mm,boxrule=0.9pt,left=3mm]{\Large\bfseries\color{cpsc}c)}\quad\begin{minipage}[c]{5.3cm}\raggedright\large $\begin{pmatrix}6\\-3\end{pmatrix}$\end{minipage}\end{tcolorbox}\hfill\begin{tcolorbox}[enhanced,nobeforeafter,width=7.1cm,height=2.1cm,valign=center,colback=cpsc!6,colframe=cpsc,arc=2.5mm,boxrule=0.9pt,left=3mm]{\Large\bfseries\color{cpsc}d)}\quad\begin{minipage}[c]{5.3cm}\raggedright\large $\begin{pmatrix}-1\\6\end{pmatrix}$\end{minipage}\end{tcolorbox}
\end{tcolorbox}
\newpage
\begin{tcolorbox}[enhanced,width=16cm,height=14cm,colback=white,colframe=cpsc,boxrule=2pt,arc=0mm,
left=6mm,right=6mm,top=0mm,bottom=4mm,sharp corners,
overlay={\fill[cpsc] (frame.north west) rectangle ([yshift=-1.25cm]frame.north east);
\node[anchor=west,text=white,font=\sffamily\bfseries\LARGE] at ([xshift=5mm,yshift=-0.62cm]frame.north west) {Question 75 D};
\node[anchor=east,text=white,font=\sffamily\itshape\large] at ([xshift=-5mm,yshift=-0.62cm]frame.north east) {Produit scalaire et géométrie repérée};}]
\vspace*{1.75cm}
\begin{minipage}[t][6.9cm][c]{\linewidth}\LARGE\centering
\begin{minipage}{0.95\linewidth}\raggedright\linespread{1.25}\selectfont Quelle est une équation cartésienne de la droite passant par $A(1\,;\,0)$ et de vecteur normal $\vec n\begin{pmatrix}5\\2\end{pmatrix}$~?\end{minipage}
\end{minipage}

\vfill
\begin{tcolorbox}[enhanced,nobeforeafter,width=7.1cm,height=2.1cm,valign=center,colback=cpsc!6,colframe=cpsc,arc=2.5mm,boxrule=0.9pt,left=3mm]{\Large\bfseries\color{cpsc}a)}\quad\begin{minipage}[c]{5.3cm}\raggedright\Large $5x+2y-5=0$\end{minipage}\end{tcolorbox}\hfill\begin{tcolorbox}[enhanced,nobeforeafter,width=7.1cm,height=2.1cm,valign=center,colback=cpsc!6,colframe=cpsc,arc=2.5mm,boxrule=0.9pt,left=3mm]{\Large\bfseries\color{cpsc}b)}\quad\begin{minipage}[c]{5.3cm}\raggedright\Large $-2x+5y+2=0$\end{minipage}\end{tcolorbox}\par\vspace{2.5mm}
\begin{tcolorbox}[enhanced,nobeforeafter,width=7.1cm,height=2.1cm,valign=center,colback=cpsc!6,colframe=cpsc,arc=2.5mm,boxrule=0.9pt,left=3mm]{\Large\bfseries\color{cpsc}c)}\quad\begin{minipage}[c]{5.3cm}\raggedright\Large $x-5=0$\end{minipage}\end{tcolorbox}\hfill\begin{tcolorbox}[enhanced,nobeforeafter,width=7.1cm,height=2.1cm,valign=center,colback=cpsc!6,colframe=cpsc,arc=2.5mm,boxrule=0.9pt,left=3mm]{\Large\bfseries\color{cpsc}d)}\quad\begin{minipage}[c]{5.3cm}\raggedright\Large $5x+2y+5=0$\end{minipage}\end{tcolorbox}
\end{tcolorbox}
\newpage
\begin{tcolorbox}[enhanced,width=16cm,height=14cm,colback=white,colframe=cpsc,boxrule=2pt,arc=0mm,
left=6mm,right=6mm,top=0mm,bottom=4mm,sharp corners,
overlay={\fill[cpsc] (frame.north west) rectangle ([yshift=-1.25cm]frame.north east);
\node[anchor=west,text=white,font=\sffamily\bfseries\LARGE] at ([xshift=5mm,yshift=-0.62cm]frame.north west) {Question 76 D};
\node[anchor=east,text=white,font=\sffamily\itshape\large] at ([xshift=-5mm,yshift=-0.62cm]frame.north east) {Produit scalaire et géométrie repérée};}]
\vspace*{1.75cm}
\begin{minipage}[t][6.9cm][c]{\linewidth}\LARGE\centering
\begin{minipage}{0.95\linewidth}\raggedright\linespread{1.25}\selectfont Quelle est une équation cartésienne de la droite passant par $A(2\,;\,2)$ et de vecteur normal $\vec n\begin{pmatrix}-1\\3\end{pmatrix}$~?\end{minipage}
\end{minipage}

\vfill
\begin{tcolorbox}[enhanced,nobeforeafter,width=7.1cm,height=2.1cm,valign=center,colback=cpsc!6,colframe=cpsc,arc=2.5mm,boxrule=0.9pt,left=3mm]{\Large\bfseries\color{cpsc}a)}\quad\begin{minipage}[c]{5.3cm}\raggedright\Large $-3x-y+8=0$\end{minipage}\end{tcolorbox}\hfill\begin{tcolorbox}[enhanced,nobeforeafter,width=7.1cm,height=2.1cm,valign=center,colback=cpsc!6,colframe=cpsc,arc=2.5mm,boxrule=0.9pt,left=3mm]{\Large\bfseries\color{cpsc}b)}\quad\begin{minipage}[c]{5.3cm}\raggedright\Large $2x+2y-4=0$\end{minipage}\end{tcolorbox}\par\vspace{2.5mm}
\begin{tcolorbox}[enhanced,nobeforeafter,width=7.1cm,height=2.1cm,valign=center,colback=cpsc!6,colframe=cpsc,arc=2.5mm,boxrule=0.9pt,left=3mm]{\Large\bfseries\color{cpsc}c)}\quad\begin{minipage}[c]{5.3cm}\raggedright\Large $-x+3y+4=0$\end{minipage}\end{tcolorbox}\hfill\begin{tcolorbox}[enhanced,nobeforeafter,width=7.1cm,height=2.1cm,valign=center,colback=cpsc!6,colframe=cpsc,arc=2.5mm,boxrule=0.9pt,left=3mm]{\Large\bfseries\color{cpsc}d)}\quad\begin{minipage}[c]{5.3cm}\raggedright\Large $-x+3y-4=0$\end{minipage}\end{tcolorbox}
\end{tcolorbox}
\newpage
\begin{tcolorbox}[enhanced,width=16cm,height=14cm,colback=white,colframe=cpsc,boxrule=2pt,arc=0mm,
left=6mm,right=6mm,top=0mm,bottom=4mm,sharp corners,
overlay={\fill[cpsc] (frame.north west) rectangle ([yshift=-1.25cm]frame.north east);
\node[anchor=west,text=white,font=\sffamily\bfseries\LARGE] at ([xshift=5mm,yshift=-0.62cm]frame.north west) {Question 77 D};
\node[anchor=east,text=white,font=\sffamily\itshape\large] at ([xshift=-5mm,yshift=-0.62cm]frame.north east) {Produit scalaire et géométrie repérée};}]
\vspace*{1.75cm}
\begin{minipage}[t][6.9cm][c]{\linewidth}\LARGE\centering
\begin{minipage}{0.95\linewidth}\raggedright\linespread{1.25}\selectfont Quels sont le centre et le rayon du cercle d'équation $(x+1)^2+(y-3)^2=1$~?\end{minipage}
\end{minipage}

\vfill
\begin{tcolorbox}[enhanced,nobeforeafter,width=7.1cm,height=2.1cm,valign=center,colback=cpsc!6,colframe=cpsc,arc=2.5mm,boxrule=0.9pt,left=3mm]{\Large\bfseries\color{cpsc}a)}\quad\begin{minipage}[c]{5.3cm}\raggedright\large $\Omega(1\,;\,-3)\text{, rayon }1$\end{minipage}\end{tcolorbox}\hfill\begin{tcolorbox}[enhanced,nobeforeafter,width=7.1cm,height=2.1cm,valign=center,colback=cpsc!6,colframe=cpsc,arc=2.5mm,boxrule=0.9pt,left=3mm]{\Large\bfseries\color{cpsc}b)}\quad\begin{minipage}[c]{5.3cm}\raggedright\large $\Omega(-1\,;\,3)\text{, rayon }2$\end{minipage}\end{tcolorbox}\par\vspace{2.5mm}
\begin{tcolorbox}[enhanced,nobeforeafter,width=7.1cm,height=2.1cm,valign=center,colback=cpsc!6,colframe=cpsc,arc=2.5mm,boxrule=0.9pt,left=3mm]{\Large\bfseries\color{cpsc}c)}\quad\begin{minipage}[c]{5.3cm}\raggedright\large $\Omega(3\,;\,-1)\text{, rayon }1$\end{minipage}\end{tcolorbox}\hfill\begin{tcolorbox}[enhanced,nobeforeafter,width=7.1cm,height=2.1cm,valign=center,colback=cpsc!6,colframe=cpsc,arc=2.5mm,boxrule=0.9pt,left=3mm]{\Large\bfseries\color{cpsc}d)}\quad\begin{minipage}[c]{5.3cm}\raggedright\large $\Omega(-1\,;\,3)\text{, rayon }1$\end{minipage}\end{tcolorbox}
\end{tcolorbox}
\newpage
\begin{tcolorbox}[enhanced,width=16cm,height=14cm,colback=white,colframe=cpsc,boxrule=2pt,arc=0mm,
left=6mm,right=6mm,top=0mm,bottom=4mm,sharp corners,
overlay={\fill[cpsc] (frame.north west) rectangle ([yshift=-1.25cm]frame.north east);
\node[anchor=west,text=white,font=\sffamily\bfseries\LARGE] at ([xshift=5mm,yshift=-0.62cm]frame.north west) {Question 78 D};
\node[anchor=east,text=white,font=\sffamily\itshape\large] at ([xshift=-5mm,yshift=-0.62cm]frame.north east) {Produit scalaire et géométrie repérée};}]
\vspace*{1.75cm}
\begin{minipage}[t][6.9cm][c]{\linewidth}\LARGE\centering
\begin{minipage}{0.95\linewidth}\raggedright\linespread{1.25}\selectfont Quels sont le centre et le rayon du cercle d'équation $(x-5)^2+(y+5)^2=36$~?\end{minipage}
\end{minipage}

\vfill
\begin{tcolorbox}[enhanced,nobeforeafter,width=7.1cm,height=2.1cm,valign=center,colback=cpsc!6,colframe=cpsc,arc=2.5mm,boxrule=0.9pt,left=3mm]{\Large\bfseries\color{cpsc}a)}\quad\begin{minipage}[c]{5.3cm}\raggedright\large $\Omega(5\,;\,-5)\text{, rayon }36$\end{minipage}\end{tcolorbox}\hfill\begin{tcolorbox}[enhanced,nobeforeafter,width=7.1cm,height=2.1cm,valign=center,colback=cpsc!6,colframe=cpsc,arc=2.5mm,boxrule=0.9pt,left=3mm]{\Large\bfseries\color{cpsc}b)}\quad\begin{minipage}[c]{5.3cm}\raggedright\large $\Omega(-5\,;\,5)\text{, rayon }6$\end{minipage}\end{tcolorbox}\par\vspace{2.5mm}
\begin{tcolorbox}[enhanced,nobeforeafter,width=7.1cm,height=2.1cm,valign=center,colback=cpsc!6,colframe=cpsc,arc=2.5mm,boxrule=0.9pt,left=3mm]{\Large\bfseries\color{cpsc}c)}\quad\begin{minipage}[c]{5.3cm}\raggedright\large $\Omega(5\,;\,-5)\text{, rayon }6$\end{minipage}\end{tcolorbox}\hfill\begin{tcolorbox}[enhanced,nobeforeafter,width=7.1cm,height=2.1cm,valign=center,colback=cpsc!6,colframe=cpsc,arc=2.5mm,boxrule=0.9pt,left=3mm]{\Large\bfseries\color{cpsc}d)}\quad\begin{minipage}[c]{5.3cm}\raggedright\large $\Omega(-5\,;\,5)\text{, rayon }36$\end{minipage}\end{tcolorbox}
\end{tcolorbox}
\newpage
\begin{tcolorbox}[enhanced,width=16cm,height=14cm,colback=white,colframe=cpsc,boxrule=2pt,arc=0mm,
left=6mm,right=6mm,top=0mm,bottom=4mm,sharp corners,
overlay={\fill[cpsc] (frame.north west) rectangle ([yshift=-1.25cm]frame.north east);
\node[anchor=west,text=white,font=\sffamily\bfseries\LARGE] at ([xshift=5mm,yshift=-0.62cm]frame.north west) {Question 79 D};
\node[anchor=east,text=white,font=\sffamily\itshape\large] at ([xshift=-5mm,yshift=-0.62cm]frame.north east) {Produit scalaire et géométrie repérée};}]
\vspace*{1.75cm}
\begin{minipage}[t][6.9cm][c]{\linewidth}\LARGE\centering
\begin{minipage}{0.95\linewidth}\raggedright\linespread{1.25}\selectfont On a $\|\vec u\|=2$, $\|\vec v\|=7$ et $\vec u\cdot\vec v=1$. Que vaut $\|\vec u+\vec v\|^2$~?\end{minipage}
\end{minipage}

\vfill
\begin{tcolorbox}[enhanced,nobeforeafter,width=7.1cm,height=2.1cm,valign=center,colback=cpsc!6,colframe=cpsc,arc=2.5mm,boxrule=0.9pt,left=3mm]{\Large\bfseries\color{cpsc}a)}\quad\begin{minipage}[c]{5.3cm}\raggedright\Large $81$\end{minipage}\end{tcolorbox}\hfill\begin{tcolorbox}[enhanced,nobeforeafter,width=7.1cm,height=2.1cm,valign=center,colback=cpsc!6,colframe=cpsc,arc=2.5mm,boxrule=0.9pt,left=3mm]{\Large\bfseries\color{cpsc}b)}\quad\begin{minipage}[c]{5.3cm}\raggedright\Large $53$\end{minipage}\end{tcolorbox}\par\vspace{2.5mm}
\begin{tcolorbox}[enhanced,nobeforeafter,width=7.1cm,height=2.1cm,valign=center,colback=cpsc!6,colframe=cpsc,arc=2.5mm,boxrule=0.9pt,left=3mm]{\Large\bfseries\color{cpsc}c)}\quad\begin{minipage}[c]{5.3cm}\raggedright\Large $54$\end{minipage}\end{tcolorbox}\hfill\begin{tcolorbox}[enhanced,nobeforeafter,width=7.1cm,height=2.1cm,valign=center,colback=cpsc!6,colframe=cpsc,arc=2.5mm,boxrule=0.9pt,left=3mm]{\Large\bfseries\color{cpsc}d)}\quad\begin{minipage}[c]{5.3cm}\raggedright\Large $55$\end{minipage}\end{tcolorbox}
\end{tcolorbox}
\newpage
\begin{tcolorbox}[enhanced,width=16cm,height=14cm,colback=white,colframe=cpsc,boxrule=2pt,arc=0mm,
left=6mm,right=6mm,top=0mm,bottom=4mm,sharp corners,
overlay={\fill[cpsc] (frame.north west) rectangle ([yshift=-1.25cm]frame.north east);
\node[anchor=west,text=white,font=\sffamily\bfseries\LARGE] at ([xshift=5mm,yshift=-0.62cm]frame.north west) {Question 80 D};
\node[anchor=east,text=white,font=\sffamily\itshape\large] at ([xshift=-5mm,yshift=-0.62cm]frame.north east) {Produit scalaire et géométrie repérée};}]
\vspace*{1.75cm}
\begin{minipage}[t][6.9cm][c]{\linewidth}\LARGE\centering
\begin{minipage}{0.95\linewidth}\raggedright\linespread{1.25}\selectfont On a $\|\vec u\|=3$, $\|\vec v\|=3$ et $\vec u\cdot\vec v=-4$. Que vaut $\|\vec u+\vec v\|^2$~?\end{minipage}
\end{minipage}

\vfill
\begin{tcolorbox}[enhanced,nobeforeafter,width=7.1cm,height=2.1cm,valign=center,colback=cpsc!6,colframe=cpsc,arc=2.5mm,boxrule=0.9pt,left=3mm]{\Large\bfseries\color{cpsc}a)}\quad\begin{minipage}[c]{5.3cm}\raggedright\Large $18$\end{minipage}\end{tcolorbox}\hfill\begin{tcolorbox}[enhanced,nobeforeafter,width=7.1cm,height=2.1cm,valign=center,colback=cpsc!6,colframe=cpsc,arc=2.5mm,boxrule=0.9pt,left=3mm]{\Large\bfseries\color{cpsc}b)}\quad\begin{minipage}[c]{5.3cm}\raggedright\Large $14$\end{minipage}\end{tcolorbox}\par\vspace{2.5mm}
\begin{tcolorbox}[enhanced,nobeforeafter,width=7.1cm,height=2.1cm,valign=center,colback=cpsc!6,colframe=cpsc,arc=2.5mm,boxrule=0.9pt,left=3mm]{\Large\bfseries\color{cpsc}c)}\quad\begin{minipage}[c]{5.3cm}\raggedright\Large $10$\end{minipage}\end{tcolorbox}\hfill\begin{tcolorbox}[enhanced,nobeforeafter,width=7.1cm,height=2.1cm,valign=center,colback=cpsc!6,colframe=cpsc,arc=2.5mm,boxrule=0.9pt,left=3mm]{\Large\bfseries\color{cpsc}d)}\quad\begin{minipage}[c]{5.3cm}\raggedright\Large $36$\end{minipage}\end{tcolorbox}
\end{tcolorbox}
\newpage
\begin{tcolorbox}[enhanced,width=16cm,height=14cm,colback=white,colframe=ccvp,boxrule=2pt,arc=0mm,
left=6mm,right=6mm,top=0mm,bottom=4mm,sharp corners,
overlay={\fill[ccvp] (frame.north west) rectangle ([yshift=-1.25cm]frame.north east);
\node[anchor=west,text=white,font=\sffamily\bfseries\LARGE] at ([xshift=5mm,yshift=-0.62cm]frame.north west) {Question 81 D};
\node[anchor=east,text=white,font=\sffamily\itshape\large] at ([xshift=-5mm,yshift=-0.62cm]frame.north east) {Probabilités conditionnelles et variables aléatoires};}]
\vspace*{1.75cm}
\begin{minipage}[t][6.9cm][c]{\linewidth}\LARGE\centering
\begin{minipage}{0.95\linewidth}\raggedright\linespread{1.25}\selectfont On a $P(A)=0{,}4$ et $P(A\cap B)=0{,}32$. Que vaut $P_A(B)$~?\end{minipage}
\end{minipage}

\vfill
\begin{tcolorbox}[enhanced,nobeforeafter,width=7.1cm,height=2.1cm,valign=center,colback=ccvp!6,colframe=ccvp,arc=2.5mm,boxrule=0.9pt,left=3mm]{\Large\bfseries\color{ccvp}a)}\quad\begin{minipage}[c]{5.3cm}\raggedright\Large $0{,}72$\end{minipage}\end{tcolorbox}\hfill\begin{tcolorbox}[enhanced,nobeforeafter,width=7.1cm,height=2.1cm,valign=center,colback=ccvp!6,colframe=ccvp,arc=2.5mm,boxrule=0.9pt,left=3mm]{\Large\bfseries\color{ccvp}b)}\quad\begin{minipage}[c]{5.3cm}\raggedright\Large $0{,}08$\end{minipage}\end{tcolorbox}\par\vspace{2.5mm}
\begin{tcolorbox}[enhanced,nobeforeafter,width=7.1cm,height=2.1cm,valign=center,colback=ccvp!6,colframe=ccvp,arc=2.5mm,boxrule=0.9pt,left=3mm]{\Large\bfseries\color{ccvp}c)}\quad\begin{minipage}[c]{5.3cm}\raggedright\Large $0{,}128$\end{minipage}\end{tcolorbox}\hfill\begin{tcolorbox}[enhanced,nobeforeafter,width=7.1cm,height=2.1cm,valign=center,colback=ccvp!6,colframe=ccvp,arc=2.5mm,boxrule=0.9pt,left=3mm]{\Large\bfseries\color{ccvp}d)}\quad\begin{minipage}[c]{5.3cm}\raggedright\Large $0{,}8$\end{minipage}\end{tcolorbox}
\end{tcolorbox}
\newpage
\begin{tcolorbox}[enhanced,width=16cm,height=14cm,colback=white,colframe=ccvp,boxrule=2pt,arc=0mm,
left=6mm,right=6mm,top=0mm,bottom=4mm,sharp corners,
overlay={\fill[ccvp] (frame.north west) rectangle ([yshift=-1.25cm]frame.north east);
\node[anchor=west,text=white,font=\sffamily\bfseries\LARGE] at ([xshift=5mm,yshift=-0.62cm]frame.north west) {Question 82 D};
\node[anchor=east,text=white,font=\sffamily\itshape\large] at ([xshift=-5mm,yshift=-0.62cm]frame.north east) {Probabilités conditionnelles et variables aléatoires};}]
\vspace*{1.75cm}
\begin{minipage}[t][6.9cm][c]{\linewidth}\LARGE\centering
\begin{minipage}{0.95\linewidth}\raggedright\linespread{1.25}\selectfont On a $P(A)=0{,}3$ et $P(A\cap B)=0{,}27$. Que vaut $P_A(B)$~?\end{minipage}
\end{minipage}

\vfill
\begin{tcolorbox}[enhanced,nobeforeafter,width=7.1cm,height=2.1cm,valign=center,colback=ccvp!6,colframe=ccvp,arc=2.5mm,boxrule=0.9pt,left=3mm]{\Large\bfseries\color{ccvp}a)}\quad\begin{minipage}[c]{5.3cm}\raggedright\Large $0{,}57$\end{minipage}\end{tcolorbox}\hfill\begin{tcolorbox}[enhanced,nobeforeafter,width=7.1cm,height=2.1cm,valign=center,colback=ccvp!6,colframe=ccvp,arc=2.5mm,boxrule=0.9pt,left=3mm]{\Large\bfseries\color{ccvp}b)}\quad\begin{minipage}[c]{5.3cm}\raggedright\Large $0{,}081$\end{minipage}\end{tcolorbox}\par\vspace{2.5mm}
\begin{tcolorbox}[enhanced,nobeforeafter,width=7.1cm,height=2.1cm,valign=center,colback=ccvp!6,colframe=ccvp,arc=2.5mm,boxrule=0.9pt,left=3mm]{\Large\bfseries\color{ccvp}c)}\quad\begin{minipage}[c]{5.3cm}\raggedright\Large $0{,}03$\end{minipage}\end{tcolorbox}\hfill\begin{tcolorbox}[enhanced,nobeforeafter,width=7.1cm,height=2.1cm,valign=center,colback=ccvp!6,colframe=ccvp,arc=2.5mm,boxrule=0.9pt,left=3mm]{\Large\bfseries\color{ccvp}d)}\quad\begin{minipage}[c]{5.3cm}\raggedright\Large $0{,}9$\end{minipage}\end{tcolorbox}
\end{tcolorbox}
\newpage
\begin{tcolorbox}[enhanced,width=16cm,height=14cm,colback=white,colframe=ccvp,boxrule=2pt,arc=0mm,
left=6mm,right=6mm,top=0mm,bottom=4mm,sharp corners,
overlay={\fill[ccvp] (frame.north west) rectangle ([yshift=-1.25cm]frame.north east);
\node[anchor=west,text=white,font=\sffamily\bfseries\LARGE] at ([xshift=5mm,yshift=-0.62cm]frame.north west) {Question 83 D};
\node[anchor=east,text=white,font=\sffamily\itshape\large] at ([xshift=-5mm,yshift=-0.62cm]frame.north east) {Probabilités conditionnelles et variables aléatoires};}]
\vspace*{1.75cm}
\begin{minipage}[t][6.9cm][c]{\linewidth}\LARGE\centering
\begin{minipage}{0.95\linewidth}\raggedright\linespread{1.25}\selectfont On a $P(A)=0{,}5$ et $P_A(B)=0{,}5$. Que vaut $P(A\cap B)$~?\end{minipage}
\end{minipage}

\vfill
\begin{tcolorbox}[enhanced,nobeforeafter,width=7.1cm,height=2.1cm,valign=center,colback=ccvp!6,colframe=ccvp,arc=2.5mm,boxrule=0.9pt,left=3mm]{\Large\bfseries\color{ccvp}a)}\quad\begin{minipage}[c]{5.3cm}\raggedright\Large $1$\end{minipage}\end{tcolorbox}\hfill\begin{tcolorbox}[enhanced,nobeforeafter,width=7.1cm,height=2.1cm,valign=center,colback=ccvp!6,colframe=ccvp,arc=2.5mm,boxrule=0.9pt,left=3mm]{\Large\bfseries\color{ccvp}b)}\quad\begin{minipage}[c]{5.3cm}\raggedright\Large $0{,}5$\end{minipage}\end{tcolorbox}\par\vspace{2.5mm}
\begin{tcolorbox}[enhanced,nobeforeafter,width=7.1cm,height=2.1cm,valign=center,colback=ccvp!6,colframe=ccvp,arc=2.5mm,boxrule=0.9pt,left=3mm]{\Large\bfseries\color{ccvp}c)}\quad\begin{minipage}[c]{5.3cm}\raggedright\Large $0{,}25$\end{minipage}\end{tcolorbox}\hfill\begin{tcolorbox}[enhanced,nobeforeafter,width=7.1cm,height=2.1cm,valign=center,colback=ccvp!6,colframe=ccvp,arc=2.5mm,boxrule=0.9pt,left=3mm]{\Large\bfseries\color{ccvp}d)}\quad\begin{minipage}[c]{5.3cm}\raggedright\Large $0{,}75$\end{minipage}\end{tcolorbox}
\end{tcolorbox}
\newpage
\begin{tcolorbox}[enhanced,width=16cm,height=14cm,colback=white,colframe=ccvp,boxrule=2pt,arc=0mm,
left=6mm,right=6mm,top=0mm,bottom=4mm,sharp corners,
overlay={\fill[ccvp] (frame.north west) rectangle ([yshift=-1.25cm]frame.north east);
\node[anchor=west,text=white,font=\sffamily\bfseries\LARGE] at ([xshift=5mm,yshift=-0.62cm]frame.north west) {Question 84 D};
\node[anchor=east,text=white,font=\sffamily\itshape\large] at ([xshift=-5mm,yshift=-0.62cm]frame.north east) {Probabilités conditionnelles et variables aléatoires};}]
\vspace*{1.75cm}
\begin{minipage}[t][6.9cm][c]{\linewidth}\LARGE\centering
\begin{minipage}{0.95\linewidth}\raggedright\linespread{1.25}\selectfont On a $P(A)=0{,}9$ et $P_A(B)=0{,}3$. Que vaut $P(A\cap B)$~?\end{minipage}
\end{minipage}

\vfill
\begin{tcolorbox}[enhanced,nobeforeafter,width=7.1cm,height=2.1cm,valign=center,colback=ccvp!6,colframe=ccvp,arc=2.5mm,boxrule=0.9pt,left=3mm]{\Large\bfseries\color{ccvp}a)}\quad\begin{minipage}[c]{5.3cm}\raggedright\Large $0{,}27$\end{minipage}\end{tcolorbox}\hfill\begin{tcolorbox}[enhanced,nobeforeafter,width=7.1cm,height=2.1cm,valign=center,colback=ccvp!6,colframe=ccvp,arc=2.5mm,boxrule=0.9pt,left=3mm]{\Large\bfseries\color{ccvp}b)}\quad\begin{minipage}[c]{5.3cm}\raggedright\Large $0{,}2$\end{minipage}\end{tcolorbox}\par\vspace{2.5mm}
\begin{tcolorbox}[enhanced,nobeforeafter,width=7.1cm,height=2.1cm,valign=center,colback=ccvp!6,colframe=ccvp,arc=2.5mm,boxrule=0.9pt,left=3mm]{\Large\bfseries\color{ccvp}c)}\quad\begin{minipage}[c]{5.3cm}\raggedright\Large $\dfrac{1}{3}$\end{minipage}\end{tcolorbox}\hfill\begin{tcolorbox}[enhanced,nobeforeafter,width=7.1cm,height=2.1cm,valign=center,colback=ccvp!6,colframe=ccvp,arc=2.5mm,boxrule=0.9pt,left=3mm]{\Large\bfseries\color{ccvp}d)}\quad\begin{minipage}[c]{5.3cm}\raggedright\Large $0{,}63$\end{minipage}\end{tcolorbox}
\end{tcolorbox}
\newpage
\begin{tcolorbox}[enhanced,width=16cm,height=14cm,colback=white,colframe=ccvp,boxrule=2pt,arc=0mm,
left=6mm,right=6mm,top=0mm,bottom=4mm,sharp corners,
overlay={\fill[ccvp] (frame.north west) rectangle ([yshift=-1.25cm]frame.north east);
\node[anchor=west,text=white,font=\sffamily\bfseries\LARGE] at ([xshift=5mm,yshift=-0.62cm]frame.north west) {Question 85 D};
\node[anchor=east,text=white,font=\sffamily\itshape\large] at ([xshift=-5mm,yshift=-0.62cm]frame.north east) {Probabilités conditionnelles et variables aléatoires};}]
\vspace*{1.75cm}
\begin{minipage}[t][6.9cm][c]{\linewidth}\LARGE\centering
\begin{minipage}{0.95\linewidth}\raggedright\linespread{1.25}\selectfont On a $P(A)=0{,}1$, $P_A(B)=0{,}5$ et $P_{\overline{A}}(B)=0{,}1$. Que vaut $P(B)$~?\end{minipage}
\end{minipage}

\vfill
\begin{tcolorbox}[enhanced,nobeforeafter,width=7.1cm,height=2.1cm,valign=center,colback=ccvp!6,colframe=ccvp,arc=2.5mm,boxrule=0.9pt,left=3mm]{\Large\bfseries\color{ccvp}a)}\quad\begin{minipage}[c]{5.3cm}\raggedright\Large $0{,}14$\end{minipage}\end{tcolorbox}\hfill\begin{tcolorbox}[enhanced,nobeforeafter,width=7.1cm,height=2.1cm,valign=center,colback=ccvp!6,colframe=ccvp,arc=2.5mm,boxrule=0.9pt,left=3mm]{\Large\bfseries\color{ccvp}b)}\quad\begin{minipage}[c]{5.3cm}\raggedright\Large $0{,}05$\end{minipage}\end{tcolorbox}\par\vspace{2.5mm}
\begin{tcolorbox}[enhanced,nobeforeafter,width=7.1cm,height=2.1cm,valign=center,colback=ccvp!6,colframe=ccvp,arc=2.5mm,boxrule=0.9pt,left=3mm]{\Large\bfseries\color{ccvp}c)}\quad\begin{minipage}[c]{5.3cm}\raggedright\Large $0{,}6$\end{minipage}\end{tcolorbox}\hfill\begin{tcolorbox}[enhanced,nobeforeafter,width=7.1cm,height=2.1cm,valign=center,colback=ccvp!6,colframe=ccvp,arc=2.5mm,boxrule=0.9pt,left=3mm]{\Large\bfseries\color{ccvp}d)}\quad\begin{minipage}[c]{5.3cm}\raggedright\Large $0{,}06$\end{minipage}\end{tcolorbox}
\end{tcolorbox}
\newpage
\begin{tcolorbox}[enhanced,width=16cm,height=14cm,colback=white,colframe=ccvp,boxrule=2pt,arc=0mm,
left=6mm,right=6mm,top=0mm,bottom=4mm,sharp corners,
overlay={\fill[ccvp] (frame.north west) rectangle ([yshift=-1.25cm]frame.north east);
\node[anchor=west,text=white,font=\sffamily\bfseries\LARGE] at ([xshift=5mm,yshift=-0.62cm]frame.north west) {Question 86 D};
\node[anchor=east,text=white,font=\sffamily\itshape\large] at ([xshift=-5mm,yshift=-0.62cm]frame.north east) {Probabilités conditionnelles et variables aléatoires};}]
\vspace*{1.75cm}
\begin{minipage}[t][6.9cm][c]{\linewidth}\LARGE\centering
\begin{minipage}{0.95\linewidth}\raggedright\linespread{1.25}\selectfont On a $P(A)=0{,}25$, $P_A(B)=0{,}4$ et $P_{\overline{A}}(B)=0{,}8$. Que vaut $P(B)$~?\end{minipage}
\end{minipage}

\vfill
\begin{tcolorbox}[enhanced,nobeforeafter,width=7.1cm,height=2.1cm,valign=center,colback=ccvp!6,colframe=ccvp,arc=2.5mm,boxrule=0.9pt,left=3mm]{\Large\bfseries\color{ccvp}a)}\quad\begin{minipage}[c]{5.3cm}\raggedright\Large $0{,}3$\end{minipage}\end{tcolorbox}\hfill\begin{tcolorbox}[enhanced,nobeforeafter,width=7.1cm,height=2.1cm,valign=center,colback=ccvp!6,colframe=ccvp,arc=2.5mm,boxrule=0.9pt,left=3mm]{\Large\bfseries\color{ccvp}b)}\quad\begin{minipage}[c]{5.3cm}\raggedright\Large $0{,}32$\end{minipage}\end{tcolorbox}\par\vspace{2.5mm}
\begin{tcolorbox}[enhanced,nobeforeafter,width=7.1cm,height=2.1cm,valign=center,colback=ccvp!6,colframe=ccvp,arc=2.5mm,boxrule=0.9pt,left=3mm]{\Large\bfseries\color{ccvp}c)}\quad\begin{minipage}[c]{5.3cm}\raggedright\Large $0{,}7$\end{minipage}\end{tcolorbox}\hfill\begin{tcolorbox}[enhanced,nobeforeafter,width=7.1cm,height=2.1cm,valign=center,colback=ccvp!6,colframe=ccvp,arc=2.5mm,boxrule=0.9pt,left=3mm]{\Large\bfseries\color{ccvp}d)}\quad\begin{minipage}[c]{5.3cm}\raggedright\Large $0{,}1$\end{minipage}\end{tcolorbox}
\end{tcolorbox}
\newpage
\begin{tcolorbox}[enhanced,width=16cm,height=14cm,colback=white,colframe=ccvp,boxrule=2pt,arc=0mm,
left=6mm,right=6mm,top=0mm,bottom=4mm,sharp corners,
overlay={\fill[ccvp] (frame.north west) rectangle ([yshift=-1.25cm]frame.north east);
\node[anchor=west,text=white,font=\sffamily\bfseries\LARGE] at ([xshift=5mm,yshift=-0.62cm]frame.north west) {Question 87 D};
\node[anchor=east,text=white,font=\sffamily\itshape\large] at ([xshift=-5mm,yshift=-0.62cm]frame.north east) {Probabilités conditionnelles et variables aléatoires};}]
\vspace*{1.75cm}
\begin{minipage}[t][6.9cm][c]{\linewidth}\LARGE\centering
\begin{minipage}{0.95\linewidth}\raggedright\linespread{1.25}\selectfont On a $P(A)=0{,}9$, $P(B)=0{,}1$ et $P(A\cap B)=0{,}09$. Les événements $A$ et $B$ sont-ils indépendants~?\end{minipage}
\end{minipage}

\vfill
\begin{tcolorbox}[enhanced,nobeforeafter,width=7.1cm,height=2.1cm,valign=center,colback=ccvp!6,colframe=ccvp,arc=2.5mm,boxrule=0.9pt,left=3mm]{\Large\bfseries\color{ccvp}a)}\quad\begin{minipage}[c]{5.3cm}\raggedright\Large on ne peut pas savoir\end{minipage}\end{tcolorbox}\hfill\begin{tcolorbox}[enhanced,nobeforeafter,width=7.1cm,height=2.1cm,valign=center,colback=ccvp!6,colframe=ccvp,arc=2.5mm,boxrule=0.9pt,left=3mm]{\Large\bfseries\color{ccvp}b)}\quad\begin{minipage}[c]{5.3cm}\raggedright\Large non\end{minipage}\end{tcolorbox}\par\vspace{2.5mm}
\begin{tcolorbox}[enhanced,nobeforeafter,width=7.1cm,height=2.1cm,valign=center,colback=ccvp!6,colframe=ccvp,arc=2.5mm,boxrule=0.9pt,left=3mm]{\Large\bfseries\color{ccvp}c)}\quad\begin{minipage}[c]{5.3cm}\raggedright\large seulement s'ils sont incompatibles\end{minipage}\end{tcolorbox}\hfill\begin{tcolorbox}[enhanced,nobeforeafter,width=7.1cm,height=2.1cm,valign=center,colback=ccvp!6,colframe=ccvp,arc=2.5mm,boxrule=0.9pt,left=3mm]{\Large\bfseries\color{ccvp}d)}\quad\begin{minipage}[c]{5.3cm}\raggedright\Large oui\end{minipage}\end{tcolorbox}
\end{tcolorbox}
\newpage
\begin{tcolorbox}[enhanced,width=16cm,height=14cm,colback=white,colframe=ccvp,boxrule=2pt,arc=0mm,
left=6mm,right=6mm,top=0mm,bottom=4mm,sharp corners,
overlay={\fill[ccvp] (frame.north west) rectangle ([yshift=-1.25cm]frame.north east);
\node[anchor=west,text=white,font=\sffamily\bfseries\LARGE] at ([xshift=5mm,yshift=-0.62cm]frame.north west) {Question 88 D};
\node[anchor=east,text=white,font=\sffamily\itshape\large] at ([xshift=-5mm,yshift=-0.62cm]frame.north east) {Probabilités conditionnelles et variables aléatoires};}]
\vspace*{1.75cm}
\begin{minipage}[t][6.9cm][c]{\linewidth}\LARGE\centering
\begin{minipage}{0.95\linewidth}\raggedright\linespread{1.25}\selectfont On a $P(A)=0{,}4$, $P(B)=0{,}4$ et $P(A\cap B)=0{,}2$. Les événements $A$ et $B$ sont-ils indépendants~?\end{minipage}
\end{minipage}

\vfill
\begin{tcolorbox}[enhanced,nobeforeafter,width=7.1cm,height=2.1cm,valign=center,colback=ccvp!6,colframe=ccvp,arc=2.5mm,boxrule=0.9pt,left=3mm]{\Large\bfseries\color{ccvp}a)}\quad\begin{minipage}[c]{5.3cm}\raggedright\large seulement s'ils sont incompatibles\end{minipage}\end{tcolorbox}\hfill\begin{tcolorbox}[enhanced,nobeforeafter,width=7.1cm,height=2.1cm,valign=center,colback=ccvp!6,colframe=ccvp,arc=2.5mm,boxrule=0.9pt,left=3mm]{\Large\bfseries\color{ccvp}b)}\quad\begin{minipage}[c]{5.3cm}\raggedright\Large on ne peut pas savoir\end{minipage}\end{tcolorbox}\par\vspace{2.5mm}
\begin{tcolorbox}[enhanced,nobeforeafter,width=7.1cm,height=2.1cm,valign=center,colback=ccvp!6,colframe=ccvp,arc=2.5mm,boxrule=0.9pt,left=3mm]{\Large\bfseries\color{ccvp}c)}\quad\begin{minipage}[c]{5.3cm}\raggedright\Large non\end{minipage}\end{tcolorbox}\hfill\begin{tcolorbox}[enhanced,nobeforeafter,width=7.1cm,height=2.1cm,valign=center,colback=ccvp!6,colframe=ccvp,arc=2.5mm,boxrule=0.9pt,left=3mm]{\Large\bfseries\color{ccvp}d)}\quad\begin{minipage}[c]{5.3cm}\raggedright\Large oui\end{minipage}\end{tcolorbox}
\end{tcolorbox}
\newpage
\begin{tcolorbox}[enhanced,width=16cm,height=14cm,colback=white,colframe=ccvp,boxrule=2pt,arc=0mm,
left=6mm,right=6mm,top=0mm,bottom=4mm,sharp corners,
overlay={\fill[ccvp] (frame.north west) rectangle ([yshift=-1.25cm]frame.north east);
\node[anchor=west,text=white,font=\sffamily\bfseries\LARGE] at ([xshift=5mm,yshift=-0.62cm]frame.north west) {Question 89 D};
\node[anchor=east,text=white,font=\sffamily\itshape\large] at ([xshift=-5mm,yshift=-0.62cm]frame.north east) {Probabilités conditionnelles et variables aléatoires};}]
\vspace*{1.75cm}
\begin{minipage}[t][6.9cm][c]{\linewidth}\Large\centering
\begin{minipage}{0.95\linewidth}\raggedright\linespread{1.25}\selectfont Voici la loi de probabilité d'une variable aléatoire $X$. Que vaut $E(X)$~? $\begin{array}{|l|c|c|c|}\hline x_i & -3 & 1 & 5\\\hline P(X=x_i) & 0{,}5 & 0{,}25 & 0{,}25\\\hline\end{array}$\end{minipage}
\end{minipage}

\vfill
\begin{tcolorbox}[enhanced,nobeforeafter,width=7.1cm,height=2.1cm,valign=center,colback=ccvp!6,colframe=ccvp,arc=2.5mm,boxrule=0.9pt,left=3mm]{\Large\bfseries\color{ccvp}a)}\quad\begin{minipage}[c]{5.3cm}\raggedright\Large $-1$\end{minipage}\end{tcolorbox}\hfill\begin{tcolorbox}[enhanced,nobeforeafter,width=7.1cm,height=2.1cm,valign=center,colback=ccvp!6,colframe=ccvp,arc=2.5mm,boxrule=0.9pt,left=3mm]{\Large\bfseries\color{ccvp}b)}\quad\begin{minipage}[c]{5.3cm}\raggedright\Large $1$\end{minipage}\end{tcolorbox}\par\vspace{2.5mm}
\begin{tcolorbox}[enhanced,nobeforeafter,width=7.1cm,height=2.1cm,valign=center,colback=ccvp!6,colframe=ccvp,arc=2.5mm,boxrule=0.9pt,left=3mm]{\Large\bfseries\color{ccvp}c)}\quad\begin{minipage}[c]{5.3cm}\raggedright\Large $0$\end{minipage}\end{tcolorbox}\hfill\begin{tcolorbox}[enhanced,nobeforeafter,width=7.1cm,height=2.1cm,valign=center,colback=ccvp!6,colframe=ccvp,arc=2.5mm,boxrule=0.9pt,left=3mm]{\Large\bfseries\color{ccvp}d)}\quad\begin{minipage}[c]{5.3cm}\raggedright\Large $11$\end{minipage}\end{tcolorbox}
\end{tcolorbox}
\newpage
\begin{tcolorbox}[enhanced,width=16cm,height=14cm,colback=white,colframe=ccvp,boxrule=2pt,arc=0mm,
left=6mm,right=6mm,top=0mm,bottom=4mm,sharp corners,
overlay={\fill[ccvp] (frame.north west) rectangle ([yshift=-1.25cm]frame.north east);
\node[anchor=west,text=white,font=\sffamily\bfseries\LARGE] at ([xshift=5mm,yshift=-0.62cm]frame.north west) {Question 90 D};
\node[anchor=east,text=white,font=\sffamily\itshape\large] at ([xshift=-5mm,yshift=-0.62cm]frame.north east) {Probabilités conditionnelles et variables aléatoires};}]
\vspace*{1.75cm}
\begin{minipage}[t][6.9cm][c]{\linewidth}\Large\centering
\begin{minipage}{0.95\linewidth}\raggedright\linespread{1.25}\selectfont Voici la loi de probabilité d'une variable aléatoire $X$. Que vaut $E(X)$~? $\begin{array}{|l|c|c|c|}\hline x_i & 10 & 20 & -4\\\hline P(X=x_i) & 0{,}1 & 0{,}2 & 0{,}7\\\hline\end{array}$\end{minipage}
\end{minipage}

\vfill
\begin{tcolorbox}[enhanced,nobeforeafter,width=7.1cm,height=2.1cm,valign=center,colback=ccvp!6,colframe=ccvp,arc=2.5mm,boxrule=0.9pt,left=3mm]{\Large\bfseries\color{ccvp}a)}\quad\begin{minipage}[c]{5.3cm}\raggedright\Large $\dfrac{26}{3}$\end{minipage}\end{tcolorbox}\hfill\begin{tcolorbox}[enhanced,nobeforeafter,width=7.1cm,height=2.1cm,valign=center,colback=ccvp!6,colframe=ccvp,arc=2.5mm,boxrule=0.9pt,left=3mm]{\Large\bfseries\color{ccvp}b)}\quad\begin{minipage}[c]{5.3cm}\raggedright\Large $2{,}2$\end{minipage}\end{tcolorbox}\par\vspace{2.5mm}
\begin{tcolorbox}[enhanced,nobeforeafter,width=7.1cm,height=2.1cm,valign=center,colback=ccvp!6,colframe=ccvp,arc=2.5mm,boxrule=0.9pt,left=3mm]{\Large\bfseries\color{ccvp}c)}\quad\begin{minipage}[c]{5.3cm}\raggedright\Large $-2{,}2$\end{minipage}\end{tcolorbox}\hfill\begin{tcolorbox}[enhanced,nobeforeafter,width=7.1cm,height=2.1cm,valign=center,colback=ccvp!6,colframe=ccvp,arc=2.5mm,boxrule=0.9pt,left=3mm]{\Large\bfseries\color{ccvp}d)}\quad\begin{minipage}[c]{5.3cm}\raggedright\Large $3{,}2$\end{minipage}\end{tcolorbox}
\end{tcolorbox}
\newpage
\begin{tcolorbox}[enhanced,width=16cm,height=14cm,colback=white,colframe=ccvp,boxrule=2pt,arc=0mm,
left=6mm,right=6mm,top=0mm,bottom=4mm,sharp corners,
overlay={\fill[ccvp] (frame.north west) rectangle ([yshift=-1.25cm]frame.north east);
\node[anchor=west,text=white,font=\sffamily\bfseries\LARGE] at ([xshift=5mm,yshift=-0.62cm]frame.north west) {Question 91 D};
\node[anchor=east,text=white,font=\sffamily\itshape\large] at ([xshift=-5mm,yshift=-0.62cm]frame.north east) {Probabilités conditionnelles et variables aléatoires};}]
\vspace*{1.75cm}
\begin{minipage}[t][6.9cm][c]{\linewidth}\Large\centering
\begin{minipage}{0.95\linewidth}\raggedright\linespread{1.25}\selectfont Voici la loi de probabilité d'une variable aléatoire $X$. Que vaut sa variance $V(X)$~? $\begin{array}{|l|c|c|}\hline x_i & 0 & 5\\\hline P(X=x_i) & 0{,}2 & 0{,}8\\\hline\end{array}$\end{minipage}
\end{minipage}

\vfill
\begin{tcolorbox}[enhanced,nobeforeafter,width=7.1cm,height=2.1cm,valign=center,colback=ccvp!6,colframe=ccvp,arc=2.5mm,boxrule=0.9pt,left=3mm]{\Large\bfseries\color{ccvp}a)}\quad\begin{minipage}[c]{5.3cm}\raggedright\Large $8$\end{minipage}\end{tcolorbox}\hfill\begin{tcolorbox}[enhanced,nobeforeafter,width=7.1cm,height=2.1cm,valign=center,colback=ccvp!6,colframe=ccvp,arc=2.5mm,boxrule=0.9pt,left=3mm]{\Large\bfseries\color{ccvp}b)}\quad\begin{minipage}[c]{5.3cm}\raggedright\Large $4$\end{minipage}\end{tcolorbox}\par\vspace{2.5mm}
\begin{tcolorbox}[enhanced,nobeforeafter,width=7.1cm,height=2.1cm,valign=center,colback=ccvp!6,colframe=ccvp,arc=2.5mm,boxrule=0.9pt,left=3mm]{\Large\bfseries\color{ccvp}c)}\quad\begin{minipage}[c]{5.3cm}\raggedright\Large $20$\end{minipage}\end{tcolorbox}\hfill\begin{tcolorbox}[enhanced,nobeforeafter,width=7.1cm,height=2.1cm,valign=center,colback=ccvp!6,colframe=ccvp,arc=2.5mm,boxrule=0.9pt,left=3mm]{\Large\bfseries\color{ccvp}d)}\quad\begin{minipage}[c]{5.3cm}\raggedright\Large $5$\end{minipage}\end{tcolorbox}
\end{tcolorbox}
\newpage
\begin{tcolorbox}[enhanced,width=16cm,height=14cm,colback=white,colframe=ccvp,boxrule=2pt,arc=0mm,
left=6mm,right=6mm,top=0mm,bottom=4mm,sharp corners,
overlay={\fill[ccvp] (frame.north west) rectangle ([yshift=-1.25cm]frame.north east);
\node[anchor=west,text=white,font=\sffamily\bfseries\LARGE] at ([xshift=5mm,yshift=-0.62cm]frame.north west) {Question 92 D};
\node[anchor=east,text=white,font=\sffamily\itshape\large] at ([xshift=-5mm,yshift=-0.62cm]frame.north east) {Probabilités conditionnelles et variables aléatoires};}]
\vspace*{1.75cm}
\begin{minipage}[t][6.9cm][c]{\linewidth}\Large\centering
\begin{minipage}{0.95\linewidth}\raggedright\linespread{1.25}\selectfont Voici la loi de probabilité d'une variable aléatoire $X$. Que vaut sa variance $V(X)$~? $\begin{array}{|l|c|c|}\hline x_i & 0 & 10\\\hline P(X=x_i) & 0{,}9 & 0{,}1\\\hline\end{array}$\end{minipage}
\end{minipage}

\vfill
\begin{tcolorbox}[enhanced,nobeforeafter,width=7.1cm,height=2.1cm,valign=center,colback=ccvp!6,colframe=ccvp,arc=2.5mm,boxrule=0.9pt,left=3mm]{\Large\bfseries\color{ccvp}a)}\quad\begin{minipage}[c]{5.3cm}\raggedright\Large $1$\end{minipage}\end{tcolorbox}\hfill\begin{tcolorbox}[enhanced,nobeforeafter,width=7.1cm,height=2.1cm,valign=center,colback=ccvp!6,colframe=ccvp,arc=2.5mm,boxrule=0.9pt,left=3mm]{\Large\bfseries\color{ccvp}b)}\quad\begin{minipage}[c]{5.3cm}\raggedright\Large $9$\end{minipage}\end{tcolorbox}\par\vspace{2.5mm}
\begin{tcolorbox}[enhanced,nobeforeafter,width=7.1cm,height=2.1cm,valign=center,colback=ccvp!6,colframe=ccvp,arc=2.5mm,boxrule=0.9pt,left=3mm]{\Large\bfseries\color{ccvp}c)}\quad\begin{minipage}[c]{5.3cm}\raggedright\Large $10$\end{minipage}\end{tcolorbox}\hfill\begin{tcolorbox}[enhanced,nobeforeafter,width=7.1cm,height=2.1cm,valign=center,colback=ccvp!6,colframe=ccvp,arc=2.5mm,boxrule=0.9pt,left=3mm]{\Large\bfseries\color{ccvp}d)}\quad\begin{minipage}[c]{5.3cm}\raggedright\Large $18$\end{minipage}\end{tcolorbox}
\end{tcolorbox}
\newpage
\begin{tcolorbox}[enhanced,width=16cm,height=14cm,colback=white,colframe=ccvp,boxrule=2pt,arc=0mm,
left=6mm,right=6mm,top=0mm,bottom=4mm,sharp corners,
overlay={\fill[ccvp] (frame.north west) rectangle ([yshift=-1.25cm]frame.north east);
\node[anchor=west,text=white,font=\sffamily\bfseries\LARGE] at ([xshift=5mm,yshift=-0.62cm]frame.north west) {Question 93 D};
\node[anchor=east,text=white,font=\sffamily\itshape\large] at ([xshift=-5mm,yshift=-0.62cm]frame.north east) {Probabilités conditionnelles et variables aléatoires};}]
\vspace*{1.75cm}
\begin{minipage}[t][6.9cm][c]{\linewidth}\Large\centering
\begin{minipage}{0.95\linewidth}\raggedright\linespread{1.25}\selectfont On choisit au hasard une fille de ce groupe. Quelle est la probabilité que ce soit un demi-pensionnaire~? $\begin{array}{|l|c|c|}\hline & \text{Demi-pens.} & \text{Externes}\\\hline \text{Filles} & 5 & 15\\\hline \text{Garçons} & 25 & 5\\\hline\end{array}$\end{minipage}
\end{minipage}

\vfill
\begin{tcolorbox}[enhanced,nobeforeafter,width=7.1cm,height=2.1cm,valign=center,colback=ccvp!6,colframe=ccvp,arc=2.5mm,boxrule=0.9pt,left=3mm]{\Large\bfseries\color{ccvp}a)}\quad\begin{minipage}[c]{5.3cm}\raggedright\Large $\dfrac{1}{4}$\end{minipage}\end{tcolorbox}\hfill\begin{tcolorbox}[enhanced,nobeforeafter,width=7.1cm,height=2.1cm,valign=center,colback=ccvp!6,colframe=ccvp,arc=2.5mm,boxrule=0.9pt,left=3mm]{\Large\bfseries\color{ccvp}b)}\quad\begin{minipage}[c]{5.3cm}\raggedright\Large $\dfrac{3}{4}$\end{minipage}\end{tcolorbox}\par\vspace{2.5mm}
\begin{tcolorbox}[enhanced,nobeforeafter,width=7.1cm,height=2.1cm,valign=center,colback=ccvp!6,colframe=ccvp,arc=2.5mm,boxrule=0.9pt,left=3mm]{\Large\bfseries\color{ccvp}c)}\quad\begin{minipage}[c]{5.3cm}\raggedright\Large $\dfrac{1}{6}$\end{minipage}\end{tcolorbox}\hfill\begin{tcolorbox}[enhanced,nobeforeafter,width=7.1cm,height=2.1cm,valign=center,colback=ccvp!6,colframe=ccvp,arc=2.5mm,boxrule=0.9pt,left=3mm]{\Large\bfseries\color{ccvp}d)}\quad\begin{minipage}[c]{5.3cm}\raggedright\Large $\dfrac{1}{10}$\end{minipage}\end{tcolorbox}
\end{tcolorbox}
\newpage
\begin{tcolorbox}[enhanced,width=16cm,height=14cm,colback=white,colframe=ccvp,boxrule=2pt,arc=0mm,
left=6mm,right=6mm,top=0mm,bottom=4mm,sharp corners,
overlay={\fill[ccvp] (frame.north west) rectangle ([yshift=-1.25cm]frame.north east);
\node[anchor=west,text=white,font=\sffamily\bfseries\LARGE] at ([xshift=5mm,yshift=-0.62cm]frame.north west) {Question 94 D};
\node[anchor=east,text=white,font=\sffamily\itshape\large] at ([xshift=-5mm,yshift=-0.62cm]frame.north east) {Probabilités conditionnelles et variables aléatoires};}]
\vspace*{1.75cm}
\begin{minipage}[t][6.9cm][c]{\linewidth}\Large\centering
\begin{minipage}{0.95\linewidth}\raggedright\linespread{1.25}\selectfont On choisit au hasard un garçon de ce groupe. Quelle est la probabilité que ce soit un demi-pensionnaire~? $\begin{array}{|l|c|c|}\hline & \text{Demi-pens.} & \text{Externes}\\\hline \text{Filles} & 18 & 2\\\hline \text{Garçons} & 10 & 10\\\hline\end{array}$\end{minipage}
\end{minipage}

\vfill
\begin{tcolorbox}[enhanced,nobeforeafter,width=7.1cm,height=2.1cm,valign=center,colback=ccvp!6,colframe=ccvp,arc=2.5mm,boxrule=0.9pt,left=3mm]{\Large\bfseries\color{ccvp}a)}\quad\begin{minipage}[c]{5.3cm}\raggedright\Large $\dfrac{1}{2}$\end{minipage}\end{tcolorbox}\hfill\begin{tcolorbox}[enhanced,nobeforeafter,width=7.1cm,height=2.1cm,valign=center,colback=ccvp!6,colframe=ccvp,arc=2.5mm,boxrule=0.9pt,left=3mm]{\Large\bfseries\color{ccvp}b)}\quad\begin{minipage}[c]{5.3cm}\raggedright\Large $\dfrac{5}{14}$\end{minipage}\end{tcolorbox}\par\vspace{2.5mm}
\begin{tcolorbox}[enhanced,nobeforeafter,width=7.1cm,height=2.1cm,valign=center,colback=ccvp!6,colframe=ccvp,arc=2.5mm,boxrule=0.9pt,left=3mm]{\Large\bfseries\color{ccvp}c)}\quad\begin{minipage}[c]{5.3cm}\raggedright\Large $\dfrac{1}{4}$\end{minipage}\end{tcolorbox}\hfill\begin{tcolorbox}[enhanced,nobeforeafter,width=7.1cm,height=2.1cm,valign=center,colback=ccvp!6,colframe=ccvp,arc=2.5mm,boxrule=0.9pt,left=3mm]{\Large\bfseries\color{ccvp}d)}\quad\begin{minipage}[c]{5.3cm}\raggedright\Large $\dfrac{7}{10}$\end{minipage}\end{tcolorbox}
\end{tcolorbox}
\newpage
\begin{tcolorbox}[enhanced,width=16cm,height=14cm,colback=white,colframe=ccvp,boxrule=2pt,arc=0mm,
left=6mm,right=6mm,top=0mm,bottom=4mm,sharp corners,
overlay={\fill[ccvp] (frame.north west) rectangle ([yshift=-1.25cm]frame.north east);
\node[anchor=west,text=white,font=\sffamily\bfseries\LARGE] at ([xshift=5mm,yshift=-0.62cm]frame.north west) {Question 95 D};
\node[anchor=east,text=white,font=\sffamily\itshape\large] at ([xshift=-5mm,yshift=-0.62cm]frame.north east) {Probabilités conditionnelles et variables aléatoires};}]
\vspace*{1.75cm}
\begin{minipage}[t][6.9cm][c]{\linewidth}\Large\centering
\begin{minipage}{0.95\linewidth}\raggedright\linespread{1.25}\selectfont Voici la loi de probabilité d'une variable aléatoire $X$. Que vaut $P(X\geqslant 1)$~? $\begin{array}{|l|c|c|c|c|c|}\hline x_i & 0 & 1 & 2 & 3 & 4\\\hline P(X=x_i) & 0{,}4 & 0{,}3 & 0{,}2 & 0{,}05 & 0{,}05\\\hline\end{array}$\end{minipage}
\end{minipage}

\vfill
\begin{tcolorbox}[enhanced,nobeforeafter,width=7.1cm,height=2.1cm,valign=center,colback=ccvp!6,colframe=ccvp,arc=2.5mm,boxrule=0.9pt,left=3mm]{\Large\bfseries\color{ccvp}a)}\quad\begin{minipage}[c]{5.3cm}\raggedright\Large $0{,}6$\end{minipage}\end{tcolorbox}\hfill\begin{tcolorbox}[enhanced,nobeforeafter,width=7.1cm,height=2.1cm,valign=center,colback=ccvp!6,colframe=ccvp,arc=2.5mm,boxrule=0.9pt,left=3mm]{\Large\bfseries\color{ccvp}b)}\quad\begin{minipage}[c]{5.3cm}\raggedright\Large $0{,}4$\end{minipage}\end{tcolorbox}\par\vspace{2.5mm}
\begin{tcolorbox}[enhanced,nobeforeafter,width=7.1cm,height=2.1cm,valign=center,colback=ccvp!6,colframe=ccvp,arc=2.5mm,boxrule=0.9pt,left=3mm]{\Large\bfseries\color{ccvp}c)}\quad\begin{minipage}[c]{5.3cm}\raggedright\Large $1$\end{minipage}\end{tcolorbox}\hfill\begin{tcolorbox}[enhanced,nobeforeafter,width=7.1cm,height=2.1cm,valign=center,colback=ccvp!6,colframe=ccvp,arc=2.5mm,boxrule=0.9pt,left=3mm]{\Large\bfseries\color{ccvp}d)}\quad\begin{minipage}[c]{5.3cm}\raggedright\Large $0{,}3$\end{minipage}\end{tcolorbox}
\end{tcolorbox}
\newpage
\begin{tcolorbox}[enhanced,width=16cm,height=14cm,colback=white,colframe=ccvp,boxrule=2pt,arc=0mm,
left=6mm,right=6mm,top=0mm,bottom=4mm,sharp corners,
overlay={\fill[ccvp] (frame.north west) rectangle ([yshift=-1.25cm]frame.north east);
\node[anchor=west,text=white,font=\sffamily\bfseries\LARGE] at ([xshift=5mm,yshift=-0.62cm]frame.north west) {Question 96 D};
\node[anchor=east,text=white,font=\sffamily\itshape\large] at ([xshift=-5mm,yshift=-0.62cm]frame.north east) {Probabilités conditionnelles et variables aléatoires};}]
\vspace*{1.75cm}
\begin{minipage}[t][6.9cm][c]{\linewidth}\Large\centering
\begin{minipage}{0.95\linewidth}\raggedright\linespread{1.25}\selectfont Voici la loi de probabilité d'une variable aléatoire $X$. Que vaut $P(X\geqslant 2)$~? $\begin{array}{|l|c|c|c|c|c|}\hline x_i & 0 & 1 & 2 & 3 & 4\\\hline P(X=x_i) & 0{,}12 & 0{,}18 & 0{,}3 & 0{,}3 & 0{,}1\\\hline\end{array}$\end{minipage}
\end{minipage}

\vfill
\begin{tcolorbox}[enhanced,nobeforeafter,width=7.1cm,height=2.1cm,valign=center,colback=ccvp!6,colframe=ccvp,arc=2.5mm,boxrule=0.9pt,left=3mm]{\Large\bfseries\color{ccvp}a)}\quad\begin{minipage}[c]{5.3cm}\raggedright\Large $0{,}4$\end{minipage}\end{tcolorbox}\hfill\begin{tcolorbox}[enhanced,nobeforeafter,width=7.1cm,height=2.1cm,valign=center,colback=ccvp!6,colframe=ccvp,arc=2.5mm,boxrule=0.9pt,left=3mm]{\Large\bfseries\color{ccvp}b)}\quad\begin{minipage}[c]{5.3cm}\raggedright\Large $0{,}7$\end{minipage}\end{tcolorbox}\par\vspace{2.5mm}
\begin{tcolorbox}[enhanced,nobeforeafter,width=7.1cm,height=2.1cm,valign=center,colback=ccvp!6,colframe=ccvp,arc=2.5mm,boxrule=0.9pt,left=3mm]{\Large\bfseries\color{ccvp}c)}\quad\begin{minipage}[c]{5.3cm}\raggedright\Large $0{,}3$\end{minipage}\end{tcolorbox}\hfill\begin{tcolorbox}[enhanced,nobeforeafter,width=7.1cm,height=2.1cm,valign=center,colback=ccvp!6,colframe=ccvp,arc=2.5mm,boxrule=0.9pt,left=3mm]{\Large\bfseries\color{ccvp}d)}\quad\begin{minipage}[c]{5.3cm}\raggedright\Large $0{,}88$\end{minipage}\end{tcolorbox}
\end{tcolorbox}
\newpage
\end{document}